\clearpage
\thispagestyle{plain}
\markboth{预备知识}{预备知识}
\oldpage[IV-1]{5}
\phantomsection
\label{chapter:0IV1-fr}

\begin{center}
{\large 第 0 章（续）\footnotemark\par}
\vspace{1.5em}
{\Large\bfseries 预备知识\par}
\vspace{2em}
{\large\bfseries 概要\par}
\end{center}

\footnotetext{请注意，符号 $0_N$（或 $\operatorname{Err}_N$）分别指随第 $N$ 章附载的第 0 章部分和勘误表。}

\begin{center}
\begin{tabular}{@{}r l@{}}
\S~14. & 拓扑空间的组合维数。\\
\S~15. & M-正则序列与 $\mathscr F$-正则序列。\\
\S~16. & Noether 局部环中的维数与深度。\\
\S~17. & 正则环。\\
\S~18. & 代数扩张的补充。\\
\S~19. & 形式平滑代数与 Cohen 环。\\
\S~20. & 导射与微分。\\
\S~21. & 特征 $p$ 的环里的微分。\\
\S~22. & 形式平滑性与正则性的微分判别准则。\\
\S~23. & 日本环。
\end{tabular}
\end{center}

前面各节几乎全部用于阐述第 IV 章中将要用到的交换代数概念。
这些概念中相当一部分虽已见于若干著作
（[1]、[12]、[13]、[17]、[30]），我们仍认为，为读者提供一套连贯而又大体自足的论述会更为方便。
这些节与第 IV 章的第 5、6、7 节（其中采用概形语言）合在一起，
在本论著内部构成一部大体独立于第 I 至 III 章的小型专论；其宗旨，是以统一的形式阐述
那些对完备化、整闭包等运算“表现良好”的环的性质，并系统地把这些性质同
若干一般观念联系起来\footnote{这里所论的大多数性质由 Chevalley、Zariski、Nagata 和 Serre 发现。
这里采用的方法最初由 A.~Grothendieck 于 1961 年秋在哈佛大学的一门课程中发展出来。}。

\setcounter{section}{13}
\section{拓扑空间的组合维数}
\label{section:0.14-fr}

\setcounter{subsection}{0}
\subsection{拓扑空间的组合维数}
\label{subsection:0.14.1-fr}

\oldpage[0\textsubscript{IV-1}]{6}
\begin{env}[14.1.1]
\label{0.14.1.1-fr}
设 $I$ 是一个有序集；按定义，$I$ 中元素的一条\emph{链}是一个有限严格递增序列
$i_0<i_1<\cdots<i_n$，其各项均属于 $I$（$n\geq 0$）；按定义，这条链的\emph{长度}为 $n$。
若 $X$ 是一个拓扑空间，则其不可约闭子集所成的集合按包含关系有序，
由此得到 $X$ 的不可约闭子集链的概念。
\end{env}

\begin{definition}[14.1.2]
\label{0.14.1.2-fr}
设 $X$ 是一个拓扑空间。把\emph{$X$ 的组合维数}（如果不会混淆，也简称为
\emph{$X$ 的维数}），记为 $\dim_c(X)$，或简称为 $\dim(X)$，定义为 $X$
的不可约闭子集链的长度的上确界。对每个 $x\in X$，把\emph{$X$ 在 $x$ 处的组合维数}
（或简称为\emph{$X$ 在 $x$ 处的维数}），记为 $\dim_x(X)$，定义为数
$\inf_U(\dim(U))$；其中 $U$ 遍历 $x$ 在 $X$ 中的所有开邻域。
\end{definition}

由此定义可得
\[
  \dim(\varnothing)=-\infty
\]
（空集在 $\overline{\mathbb R}$ 中的上确界为 $-\infty$）。若 $(X_\alpha)$ 是
$X$ 的各不可约分支所成的族，则
\[
  \dim(X)=\sup_\alpha\dim(X_\alpha)
  \tag{14.1.2.1}\label{0.14.1.2.1-fr}
\]
因为按定义，$X$ 的每条不可约闭子集链都包含于 $X$ 的某个不可约分支中；反过来，
这些分支在 $X$ 中是闭的，所以 $X_\alpha$ 的每个不可约闭子集也是 $X$ 的不可约闭子集。

\begin{definition}[14.1.3]
\label{0.14.1.3-fr}
若拓扑空间 $X$ 的所有不可约分支都有相同的维数，则称这个空间是\emph{等维的}
（由 (14.1.2.1)，这个共同维数等于 $\dim(X)$）。
\end{definition}

\begin{proposition}[14.1.4]
\label{0.14.1.4-fr}
\begin{enumerate}
  \item[\rm(i)] 设 $Y$ 是拓扑空间 $X$ 的任一子集，则有
  $\dim(Y)\leq\dim(X)$。
  \item[\rm(ii)] 若拓扑空间 $X$ 是有限多个闭子集 $X_i$ 的并，则
  $\dim(X)=\sup_i\dim(X_i)$。
\end{enumerate}
\end{proposition}

\begin{proof}
任取不可约闭子集 $Z$，它位于 $Y$ 中；其闭包 $\overline Z$，即 $Z$ 在
$X$ 中的闭包，是不可约的（$0_{\mathrm I}$, 2.1.2），并且
$\overline Z\cap Y=Z$，由此得到 (i)。另一方面，若
$X=\bigcup_{i=1}^n X_i$，其中各 $X_i$ 均为闭集，则 $X$ 的每个不可约闭子集都包含于
某个 $X_i$ 中（$0_{\mathrm I}$, 2.1.1）；因此，$X$ 的每条不可约闭子集链都包含于
某个 $X_i$ 中，由此得到 (ii)。
\end{proof}

由 (14.1.4, (i)) 可知，对每个 $x\in X$，还可以写成
\[
  \dim_x(X)=\lim_U\dim(U)
  \tag{14.1.4.1}\label{0.14.1.4.1-fr}
\]
其中的极限沿 $x$ 在 $X$ 中的开邻域按反向包含关系所成的滤过有序集取值。

\oldpage[0\textsubscript{IV-1}]{7}
\begin{corollary}[14.1.5]
\label{0.14.1.5-fr}
设 $X$ 是一个拓扑空间，$x$ 是 $X$ 的一点，$U$ 是 $x$ 的一个邻域，
$Y_i$（$1\leq i\leq n$）是 $U$ 的闭子集，并且 $x\in Y_i$ 对每个 $i$ 都成立，
且 $U$ 是各 $Y_i$ 的并。则
\[
  \dim_x(X)=\sup_i\bigl(\dim_x(Y_i)\bigr).
  \tag{14.1.5.1}\label{0.14.1.5.1-fr}
\]
\end{corollary}

\begin{proof}
由 (14.1.4, (ii)) 可得
$\dim_x(X)=\inf_V(\sup_i(\dim(Y_i\cap V)))$，其中 $V$ 遍历 $x$ 的、包含于 $U$ 的
开邻域；同样，$\dim_x(Y_i)=\inf_V(\dim(Y_i\cap V))$ 对每个 $i$ 都成立。因此，若
\[
  \sup_i\bigl(\dim_x(Y_i)\bigr)=+\infty~;
\]
则推论显然成立；否则，存在一个开邻域 $V_0\subset U$，它含有 $x$，使得
$\dim(Y_i\cap V)=\dim_x(Y_i)$ 对 $1\leq i\leq n$ 以及每个 $V\subset V_0$ 都成立，
结论随即成立。
\end{proof}

\begin{proposition}[14.1.6]
\label{0.14.1.6-fr}
对每个拓扑空间 $X$，都有
$\dim(X)=\sup_{x\in X}\dim_x(X)$。
\end{proposition}

\begin{proof}
由定义 (14.1.2) 和命题 (14.1.4) 可知，
$\dim_x(X)\leq\dim(X)$ 对每个 $x\in X$ 都成立。另一方面，设
$Z_0\subset Z_1\subset\cdots\subset Z_n$ 是 $X$ 的一条不可约闭子集链，并取
$x\in Z_0$；对每个开集 $U\subset X$，若它包含 $x$，则 $U\cap Z_i$ 是不可约的
（$0_{\mathrm I}$, 2.1.6），并且在 $U$ 中是闭的；又因为有
$\overline{U\cap Z_i}=Z_i$（在 $X$ 中），所以各 $U\cap Z_i$ 两两不同；故
$\dim(U)\geq n$，证明完毕。
\end{proof}

\begin{corollary}[14.1.7]
\label{0.14.1.7-fr}
若 $(X_\alpha)$ 是 $X$ 的一个开覆盖，或 $X$ 的一个局部有限闭覆盖，则
$\dim(X)=\sup_\alpha(\dim(X_\alpha))$。
\end{corollary}

\begin{proof}
若 $X_\alpha$ 是 $x\in X$ 的一个邻域，则
$\dim_x(X)\leq\dim(X_\alpha)$，由此得到关于开覆盖的断言。另一方面，若
各 $X_\alpha$ 都是闭的，而 $U$ 是 $x\in X$ 的一个只与有限多个 $X_\alpha$ 相交的邻域，则
\[
  \dim_x(X)\leq\dim(U)=\sup_\alpha\bigl(\dim(U\cap X_\alpha)\bigr)
  \leq\sup_\alpha\bigl(\dim(X_\alpha)\bigr)
\]
由 (14.1.4) 即得第二个断言。
\end{proof}

\begin{corollary}[14.1.8]
\label{0.14.1.8-fr}
设 $X$ 是一个 Noether Kolmogorov 空间（$0_{\mathrm I}$, 2.1.3），并设
$F$ 是 $X$ 的全体闭点所成的集合。则
$\dim(X)=\sup_{x\in F}\dim_x(X)$。
\end{corollary}

\begin{proof}
沿用 (14.1.6) 的证明中的记号，只须注意 $Z_0$ 中存在一个闭点
（$0_{\mathrm I}$, 2.1.3）。
\end{proof}

\begin{proposition}[14.1.9]
\label{0.14.1.9-fr}
设 $X$ 是一个非空 Noether Kolmogorov 空间。$\dim(X)=0$ 的充分必要条件是
$X$ 有限且离散。
\end{proposition}

\begin{proof}
若空间 $X$ 是 Hausdorff 的（特别地，若 $X$ 是离散空间），则 $X$ 的不可约闭子集
只能是单点集，故 $\dim(X)=0$。反过来，设 $X$ 是满足 $\dim(X)=0$ 的 Noether
Kolmogorov 空间；由于 $X$ 的每个不可约分支都包含一个闭点
（$0_{\mathrm I}$, 2.1.3），该分支必等于这个单点。又因 $X$ 只有有限多个不可约分支，
所以它有限且离散。
\end{proof}

\begin{corollary}[14.1.10]
\label{0.14.1.10-fr}
设 $X$ 是一个 Noether Kolmogorov 空间。点 $x\in X$ 是孤立点的充分必要条件是
$\dim_x(X)=0$。
\end{corollary}

\begin{proof}
必要性显然（无需对 $X$ 作任何假设）。该条件也充分，
\oldpage[0\textsubscript{IV-1}]{8}
因为由此可知存在一个满足 $\dim(U)=0$ 的开邻域 $U$，它含有 $x$；
而 $U$ 是 Noether Kolmogorov 空间，所以 $U$ 有限且离散。
\end{proof}

\begin{proposition}[14.1.11]
\label{0.14.1.11-fr}
函数 $x\mapsto\dim_x(X)$ 在 $X$ 上是上半连续的。
\end{proposition}

\begin{proof}
在该函数取值为 $+\infty$ 的每一点，它显然是上半连续的。因此，设
$\dim_x(X)=n<+\infty$；于是公式 (14.1.4.1) 表明，存在一个开邻域 $U_0$，它含有 $x$，
使得 $\dim(U)=n$ 对每个开邻域 $U\subset U_0$（它含有 $x$）都成立。既然如此，
对每个 $y\in U_0$ 以及每个开邻域 $V\subset U_0$（它含有 $y$），由 (14.1.4) 有
$\dim(V)\leq\dim(U_0)=n$；再由 (14.1.4.1) 得到 $\dim_y(X)\leq n$。
\end{proof}

\begin{remark}[14.1.12]
\label{0.14.1.12-fr}
设 $X$、$Y$ 是两个拓扑空间，$f:X\to Y$ 是一个连续映射，则可能有
$\dim(f(X))>\dim(X)$；以下给出一个例子：取 $X$ 为具有两个元素 $a$、$b$ 的离散空间，
取 $Y$ 为集合 $\{a,b\}$，并赋予它以闭集恰为 $\varnothing$、$\{a\}$、$\{a,b\}$
的拓扑；若 $f:X\to Y$ 是恒等映射，则 $\dim(Y)=1$ 且 $\dim(X)=0$。
请注意，$Y$ 是一个离散赋值环 $A$ 的谱，其中 $a$ 是唯一的闭点，$b$ 是泛点；
若 $K$ 和 $k$ 分别是 $A$ 的分式域和剩余域，则 $X$ 是环 $k\times K$ 的谱，而 $f$
是对应于同态 $(\varphi,\psi):A\to k\times K$ 的连续映射，其中
$\varphi:A\to k$ 和 $\psi:A\to K$ 是典范同态（参见 IV, 5.4.3）。
\end{remark}


\subsection{闭子空间的余维数}
\label{subsection:0.14.2-fr}

\begin{definition}[14.2.1]
\label{0.14.2.1-fr}
设 $Y$ 是拓扑空间 $X$ 的一个不可约闭子集。$Y$ 在 $X$ 中的
\emph{组合余维数}（或简称\emph{余维数}），记为 $\operatorname{codim}(Y,X)$，定义为
$X$ 的不可约闭子集链的长度的上确界，其中 $Y$ 是最小元素。若 $Y$ 是 $X$ 的
任一闭子集，则\emph{$Y$ 在 $X$ 中的余维数}，仍记为 $\operatorname{codim}(Y,X)$，定义为
$X$ 中由 $Y$ 的各不可约分支给出的余维数的下确界。若 $X$ 的所有极小不可约闭子集
在 $X$ 中的余维数都相同，则称 $X$ 是\emph{等余维的}。
\end{definition}

由此定义可得 $\operatorname{codim}(\varnothing,X)=+\infty$，因为空集在
$\overline{\mathbb R}$ 中的下确界为 $+\infty$。若 $Y$ 在 $X$ 中闭，并且
$(X_\alpha)$（分别地，$(Y_\beta)$）是 $X$（分别地，$Y$）的各不可约分支所成的族，
则每个 $Y_\beta$ 都包含于某个 $X_\alpha$ 中；更一般地，$X$ 的每条不可约闭子集链，
若以 $Y_\beta$ 为最小元素，则其各项都是某个 $X_\alpha$ 的子集。因此
\[
  \operatorname{codim}(Y,X)=
  \inf_\beta\left(\sup_\alpha
  \bigl(\operatorname{codim}(Y_\beta,X_\alpha)\bigr)\right)
  \tag{14.2.1.1}\label{0.14.2.1.1-fr}
\]
其中，对每个 $\beta$，$\alpha$ 遍历满足 $Y_\beta\subset X_\alpha$ 的指标集。

\begin{proposition}[14.2.2]
\label{0.14.2.2-fr}
设 $X$ 是一个拓扑空间。
\begin{enumerate}
  \item[\rm(i)] 若 $\Phi$ 是 $X$ 的全体不可约闭子集所成的集合，则
  \[
    \dim(X)=\sup_{Y\in\Phi}\bigl(\operatorname{codim}(Y,X)\bigr).
    \tag{14.2.2.1}\label{0.14.2.2.1-fr}
  \]
  \oldpage[0\textsubscript{IV-1}]{9}
  \item[\rm(ii)] 对任一非空闭子集 $Y$（视为 $X$ 的闭子集），都有
  \[
    \dim(Y)+\operatorname{codim}(Y,X)\leq\dim(X).
    \tag{14.2.2.2}\label{0.14.2.2.2-fr}
  \]
  \item[\rm(iii)] 若 $Y$、$Z$、$T$ 是 $X$ 的三个闭子集，并且
  $Y\subset Z\subset T$，则
  \[
    \operatorname{codim}(Y,Z)+\operatorname{codim}(Z,T)
    \leq\operatorname{codim}(Y,T).
    \tag{14.2.2.3}\label{0.14.2.2.3-fr}
  \]
  \item[\rm(iv)] 闭子集 $Y$（视为 $X$ 的子集）满足
  $\operatorname{codim}(Y,X)=0$ 的充分必要条件是 $Y$ 包含 $X$ 的一个不可约分支。
\end{enumerate}
\end{proposition}

\begin{proof}
断言 (i) 与 (iv) 是定义 (14.2.1) 的直接推论。为证明 (ii)，只须考虑 $Y$ 不可约的情形；
这时公式直接由定义 (14.1.1) 和 (14.2.1) 得出。最后，为证明 (iii)，由定义 (14.2.1)，
可以先只考虑 $Y$ 不可约的情形；于是
$\operatorname{codim}(Y,Z)=
\sup_\alpha(\operatorname{codim}(Y,Z_\alpha))$，其中 $Z_\alpha$ 遍历 $Z$ 的、包含 $Y$
的不可约分支。显然有
$\operatorname{codim}(Y,T)\geq\operatorname{codim}(Y,Z)$，所以当
$\operatorname{codim}(Y,Z)=+\infty$ 时不等式成立；否则，存在某个 $\alpha$，使
$\operatorname{codim}(Y,Z)=\operatorname{codim}(Y,Z_\alpha)$；再由定义 (14.2.1)，
可以只考虑 $Z$ 也不可约的情形；这时，不等式 (14.2.2.3) 显然由定义 (14.2.1) 得出。
\end{proof}

\begin{proposition}[14.2.3]
\label{0.14.2.3-fr}
设 $X$ 是一个拓扑空间，$Y$ 是 $X$ 的一个闭子集。对任一开集 $U$（视为 $X$ 的开子集），都有
\[
  \operatorname{codim}(Y\cap U,U)\geq\operatorname{codim}(Y,X).
  \tag{14.2.3.1}\label{0.14.2.3.1-fr}
\]
此外，(14.2.3.1) 两端相等的充分必要条件是：若 $(Y_\alpha)$ 是 $Y$ 的、与 $U$ 相交的
各不可约分支所成的族，则有
$\operatorname{codim}(Y,X)=
\inf_\alpha(\operatorname{codim}(Y_\alpha,X))$。
\end{proposition}

\begin{proof}
我们知道（$0_{\mathrm I}$, 2.1.6），$Z\mapsto\overline Z$ 是从 $U$ 的全体不可约闭子集
到 $X$ 的与 $U$ 相交的全体不可约闭子集的双射；特别地，它使 $Y\cap U$ 的不可约分支
与 $Y$ 的、同 $U$ 相交的不可约分支对应。若 $Y_\alpha$ 是后一类分支之一，则
$\operatorname{codim}(Y_\alpha,X)=
\operatorname{codim}(Y_\alpha\cap U,U)$，命题于是由定义 (14.2.1) 得出。
\end{proof}

\begin{definition}[14.2.4]
\label{0.14.2.4-fr}
设 $X$ 是一个拓扑空间，$Y$ 是 $X$ 的一个闭子集，$x$ 是 $X$ 的一点。把
\emph{$Y$ 在 $X$ 中、点 $x$ 处的余维数}记为 $\operatorname{codim}_x(Y,X)$，并定义为数
$\sup_U(\operatorname{codim}(Y\cap U,U))$，其中 $U$ 遍历 $x$ 在 $X$ 中的所有开邻域。
\end{definition}

由 (14.2.3)，还可以写成
\[
  \operatorname{codim}_x(Y,X)=
  \lim_U\bigl(\operatorname{codim}(Y\cap U,U)\bigr)
  \tag{14.2.4.1}\label{0.14.2.4.1-fr}
\]
其中的极限沿 $x$ 在 $X$ 中的开邻域按反向包含关系所成的滤过有序集取值。还应注意
\[
  \operatorname{codim}_x(Y,X)=+\infty\quad\text{若}\quad x\in X-Y.
\]

\oldpage[0\textsubscript{IV-1}]{10}

\begin{proposition}[14.2.5]
\label{0.14.2.5-fr}
若 $(Y_i)_{1\leq i\leq n}$ 是拓扑空间 $X$ 的有限个闭子集所成的族，且 $Y$ 是这个族的并，则
\[
  \operatorname{codim}(Y,X)=
  \inf_i\bigl(\operatorname{codim}(Y_i,X)\bigr).
  \tag{14.2.5.1}\label{0.14.2.5.1-fr}
\]
\end{proposition}

\begin{proof}
事实上，某个 $Y_i$ 的每个不可约分支都包含于 $Y$ 的一个不可约分支中；反过来，
$Y$ 的每个不可约分支也都是某个 $Y_i$ 的不可约分支（$0_{\mathrm I}$, 2.1.1）。
因此，结论由定义 (14.2.1) 和不等式 (14.2.2.3) 得出。
\end{proof}

\begin{corollary}[14.2.6]
\label{0.14.2.6-fr}
设 $X$ 是一个拓扑空间，$Y$ 是 $X$ 的一个局部 Noether 的闭子空间。
\begin{enumerate}
  \item[\rm(i)] 对每个 $x\in X$，仅有有限个不可约分支 $Y_i$
  （$1\leq i\leq n$）既是 $Y$ 的分支又包含 $x$，并且
  $\operatorname{codim}_x(Y,X)=
  \inf_i(\operatorname{codim}(Y_i,X))$。
  \item[\rm(ii)] 函数
  $x\mapsto\operatorname{codim}_x(Y,X)$ 在 $X$ 上是下半连续的。
\end{enumerate}
\end{corollary}

\begin{proof}
事实上，由假设可取开邻域 $U_0$；它是 $x$ 在 $X$ 中的邻域，并且 $Y\cap U_0$ 是 Noether 空间，
因而只有有限个不可约分支；这些分支在 $U_0$ 上都是 $Y$ 的某些不可约分支的迹。
因而，尤其只有有限个不可约分支 $Y_i$（$1\leq i\leq n$）既是 $Y$ 的分支又包含 $x$；
再把 $U_0$ 换成开邻域 $U\subset U_0$，并使它含 $x$，且不与任何 $Y_j$ 相交，
只要该分支不含 $x$；这样便可设 $Y_i\cap U$ 恰为 $Y\cap U$ 的各不可约分支。
于是，对每个开邻域 $V\subset U$，若它是 $x$ 在 $X$ 中的邻域，则各 $Y_i\cap V$ 恰为 $Y\cap V$
的不可约分支；此时 (14.2.3) 表明
$\operatorname{codim}(Y_i,X)=\operatorname{codim}(Y_i\cap V,V)$，从而证明了 (i)。
此外，对每个 $x'\in U$，$Y$ 的包含 $x'$ 的不可约分支是某些 $Y_i$，所以
$\operatorname{codim}_{x'}(Y,X)\geq\operatorname{codim}_x(Y,X)$，这就证明了 (ii)。
\end{proof}

\subsection{链条件}
\label{subsection:0.14.3-fr}

\begin{env}[14.3.1]
\label{0.14.3.1-fr}
在拓扑空间 $X$ 中，称不可约闭子集链
$Z_0\subset Z_1\subset\cdots\subset Z_n$ 是\emph{饱和的}，如果不存在不可约闭子集 $Z'$，
它异于诸 $Z_i$，并使得 $Z_k\subset Z'\subset Z_{k+1}$ 对某个指标 $k$ 成立。
\end{env}

\begin{proposition}[14.3.2]
\label{0.14.3.2-fr}
设 $X$ 是一个拓扑空间，并假定任取不可约闭子集 $Y$、$Z$（二者均为 $X$ 的子集），只要 $Y\subset Z$，就有
$\operatorname{codim}(Y,Z)<+\infty$。下列两个条件等价：
\begin{enumerate}
  \item[\rm(a)] $X$ 的任意两条具有相同端点的不可约闭子集饱和链都有相同长度。
  \item[\rm(b)] 若 $Y$、$Z$、$T$ 是 $X$ 的三个不可约闭子集，并且
  $Y\subset Z\subset T$，则
  \[
    \operatorname{codim}(Y,T)=
    \operatorname{codim}(Y,Z)+\operatorname{codim}(Z,T).
    \tag{14.3.2.1}\label{0.14.3.2.1-fr}
  \]
\end{enumerate}
\end{proposition}

\begin{proof}
(a) 蕴含 (b) 是显然的。反过来，设 (b) 成立；若两条具有相同端点的饱和链的长度分别为
$m$ 和 $n\leq m$，我们来证明必有 $m=n$。对 $n$ 作归纳；当 $n=1$ 时命题显然。
因而设 $n>1$、$n<m$，并取一条饱和链
$Z_0\subset Z_1\subset\cdots\subset Z_n$，
\oldpage[0\textsubscript{IV-1}]{11}
使得还存在另一条以
$Z_0$、$Z_n$ 为端点且长度为 $m$ 的饱和链。由于
$\operatorname{codim}(Z_0,Z_n)\geq m>n$ 且
$\operatorname{codim}(Z_0,Z_1)=1$，由 (b) 得
$\operatorname{codim}(Z_1,Z_n)=\operatorname{codim}(Z_0,Z_n)-1>n-1$，
这与归纳假设矛盾。
\end{proof}

当 (14.3.2) 的条件满足时，称 $X$ 满足\emph{链条件}，或称它为\emph{悬链空间}。
显然，悬链空间的每个闭子空间仍是悬链空间。

\begin{proposition}[14.3.3]
\label{0.14.3.3-fr}
设 $X$ 是一个有限维 Noether Kolmogorov 空间。下列条件等价：
\begin{enumerate}
  \item[\rm(a)] $X$ 的任意两条不可约闭子集极大链都有相同长度。
  \item[\rm(b)] $X$ 是等维、等余维且悬链的。
  \item[\rm(c)] $X$ 是等维的，并且对任意两个不可约闭子集 $Y$、$Z$，若二者均为 $X$ 的子集且 $Y\subset Z$，则
  \[
    \dim(Z)=\dim(Y)+\operatorname{codim}(Y,Z).
    \tag{14.3.3.1}\label{0.14.3.3.1-fr}
  \]
  \item[\rm(d)] $X$ 是等余维的，并且对任意两个不可约闭子集 $Y$、$Z$，若二者均为 $X$ 的子集且 $Y\subset Z$，则
  \[
    \operatorname{codim}(Y,X)=
    \operatorname{codim}(Y,Z)+\operatorname{codim}(Z,X).
    \tag{14.3.3.2}\label{0.14.3.3.2-fr}
  \]
\end{enumerate}
\end{proposition}

\begin{proof}
关于 $X$ 的假设蕴涵：$X$ 的不可约闭子集的任一极大链，其端点必为一个闭点和 $X$
的一个不可约分支（$0_{\mathrm I}$, 2.1.3）。此外，任一以 $Y$、$Z$ 为端点
（其中 $Y\subset Z$）的饱和链，都包含于一条极大链中；该极大链里不属于给定链的元素，
或者包含于 $Y$，或者包含 $Z$。这些观察立即建立了 (a) 与 (b) 的等价性，并且还证明：
若 (a) 成立，则对每个不可约闭子集 $Y$（视为 $X$ 的子集）都有
\[
  \dim(Y)+\operatorname{codim}(Y,X)=\dim(X)~;
  \tag{14.3.3.3}\label{0.14.3.3.3-fr}
\]
又因 (14.3.2.1) 成立，由 (14.3.3.3) 立即推出 (14.3.3.1) 和 (14.3.3.2)。
反过来，(14.3.3.1) 蕴涵 (14.3.2.1)，故由 (14.3.2)，(14.3.3.1) 蕴涵链条件。
此外，在 (14.3.3.1) 中取 $Y$ 为闭点 $x$ 所成的单点集（其中该闭点属于 $X$），
并取 $Z$ 为 $X$ 的一个不可约分支，得到
$\operatorname{codim}(\{x\},X)=\dim(Z)$；由此可见 (c) 蕴涵 (b)。
同样，(14.3.3.2) 蕴涵 (14.3.2.1)，因而蕴涵链条件；此外，对 $Y$、$Z$
作与上文相同的选择，(14.3.3.2) 仍给出
$\operatorname{codim}(\{x\},X)=\dim(Z)$。于是，由于 $X$ 的每个不可约分支都包含一个闭点
（$0_{\mathrm I}$, 2.1.3），(d) 蕴涵 (b)。
\end{proof}

若一个 Noether Kolmogorov 空间具有\emph{有限维数}并满足 (14.3.3) 的等价条件，
则称它为\emph{双等维的}。

\begin{corollary}[14.3.4]
\label{0.14.3.4-fr}
设 $X$ 是一个双等维的 Noether Kolmogorov 空间；则对每个闭点 $x$（作为 $X$ 的点）和
每个不可约分支 $Z$（作为 $X$ 的分支），都有
\[
  \dim(X)=\dim(Z)=\operatorname{codim}(\{x\},X)=\dim_x(X).
  \tag{14.3.4.1}\label{0.14.3.4.1-fr}
\]
\end{corollary}

\oldpage[0\textsubscript{IV-1}]{12}
\begin{proof}
最后一个等式源于以下事实：若
$Y_0=\{x\}\subset Y_1\subset\cdots\subset Y_m$ 是 $X$ 的不可约闭子集的一条极大链，
且 $U$ 是 $x$ 的一个开邻域，则各 $U\cap Y_i$ 是 $U$ 的两两不同的不可约闭子集
（因为 $\overline{U\cap Y_i}=Y_i$）；于是由 (14.1.4) 得
$\dim(U)=\dim(X)$。
\end{proof}

\begin{corollary}[14.3.5]
\label{0.14.3.5-fr}
设 $X$ 是一个 Noether Kolmogorov 空间；若 $X$ 是双等维的，则 $X$ 的任意若干不可约分支之并
以及 $X$ 的任意不可约闭子集也都是双等维的。此外，对每个闭子集 $Y$（视为 $X$ 的子集），都有
\[
  \dim(Y)+\operatorname{codim}(Y,X)=\dim(X).
  \tag{14.3.5.1}\label{0.14.3.5.1-fr}
\]
\end{corollary}

\begin{proof}
由于 $X$ 的每条不可约闭子集链都包含于 $X$ 的一个不可约分支中，第一个断言立即由
(14.3.3) 得出。此外，若 $X'$ 是 $X$ 的一个不可约闭子集，则 $X'$ 显然满足条件
(14.3.3(c))，从而得到第二个断言。

最后，为证明 (14.3.5.1)，注意我们在 (14.3.3) 的证明中已经看到：当 $Y$ 不可约时，
该等式成立。若 $Y_i$（$1\leq i\leq m$）是 $Y$ 的各不可约分支，则诸 $Y_i$ 中
使 $\dim(Y_i)$ 取最大值的那些，恰是使 $\operatorname{codim}(Y_i,X)$ 取最小值的那些；
所以 (14.3.5.1) 由 $\dim(Y)$ 和 $\operatorname{codim}(Y,X)$ 的定义得出。
\end{proof}

\begin{remark}[14.3.6]
\label{0.14.3.6-fr}
读者会注意到，(14.3.2) 的证明适用于任意有序集；该证明并未使用所讨论的元素是某个拓扑空间
的不可约闭子集这一事实。对有序集 $E$，若每个 $x\in E$ 都存在一个 $z\leq x$，
它在 $E$ 中是\emph{极小的}，并且 $E$ 的元素链的长度有界，则 (14.3.3) 的证明也同样适用。
\end{remark}

\section{\texorpdfstring{$M$}{M}-正则序列与
\texorpdfstring{$\mathscr F$}{F}-正则序列}
\label{section:0.15-fr}

\subsection{\texorpdfstring{$M$}{M}-正则序列与
\texorpdfstring{$M$}{M}-拟正则序列}
\label{subsection:0.15.1-fr}
\label{0.15.1-fr}

\begin{env}[15.1.1]
\label{0.15.1.1-fr}
\oldpage[0\textsubscript{IV-1}]{12}
设 $A$ 是一个环，$N$ 是一个 $A$-模；以下为简便起见，以
$N[T_1,\ldots,T_n]$ 表示分次 $A$-模
$N\otimes_A A[T_1,\ldots,T_n]$（$T_i$ 为不定元），其中 $N$ 取平凡分次，
$A[T_1,\ldots,T_n]$ 按总次数分次。

设 $M$ 是一个 $A$-模，$f_i$（$1\leq i\leq n$）是 $A$ 的元素，
$\mathfrak{J}=\sum_{i=1}^n f_iA$ 是 $A$ 中由诸 $f_i$ 生成的理想，并赋予
$M$ 以 $\mathfrak{J}$-预进制滤过；于是定义一个\emph{分次 $A$-模的满同态}，
其次数为 $0$，
\[
  \varphi : (\gr_0(M))[T_1,\ldots,T_n]\longrightarrow\gr_\bullet(M)
  \tag{15.1.1.1}\label{0.15.1.1.1-fr}
\]
如下：若 $\xi\in\gr_0(M)$ 是一个模 $\mathfrak{J}M$ 的类，并由某个元素
$x\in M$ 表出，则对由 $\xi$ 与一个齐次多项式 $P[T_1,\ldots,T_n]$
所成的偶，其中后者的次数为 $k$，使之对应于模 $\mathfrak{J}^{k+1}M$ 的类，
该类以 $P(f_1,\ldots,f_n)x$ 为代表元；立即可见，这个类
\oldpage[0\textsubscript{IV-1}]{13}
只依赖于 $\xi$ 和 $P$。若 $S_k$ 表示次数为 $k$、属于
$A[T_1,\ldots,T_n]$ 的齐次多项式所成的集合，则由此定义了一个 $A$-线性映射
\[
  \varphi_k:\gr_0(M)\otimes_A S_k\longrightarrow\gr_k(M),
\]
它由 $\mathfrak{J}^kM$ 的定义即为满射。

以下称元素 $f\in A$ 是\emph{$A$-模 $M$ 中的非零因子}，如果位似
$f_M:x\mapsto fx$（定义在 $M$ 上）是\emph{单射}。
\end{env}

\begin{lemma}[15.1.2]
\label{0.15.1.2-fr}
设 $M$ 是一个 $A$-模，$f$ 是 $A$ 的一个元素，并赋予 $M$ 以
$(fA)$-预进制滤过。若 $f$ 是 $M$ 中的非零因子，则 (15.1.1) 中定义的典范同态
\[
  \varphi:(M/fM)[T]\longrightarrow\gr_\bullet(M)
  \tag{15.1.2.1}\label{0.15.1.2.1-fr}
\]
是双射。若 $M$ 对 $(fA)$-预进制拓扑是分离的，则逆命题也成立。
\end{lemma}

\begin{proof}
这里 $\gr_k(M)=f^kM/f^{k+1}M$；说 $\varphi$ 是单射，意即对每个
$x\in M$ 和每个 $k$，关系 $f^kx\in f^{k+1}M$ 都蕴涵 $x\in fM$；
若 $f$ 是 $M$ 中的非零因子，这显然成立。反过来，注意位似
$f_M:x\mapsto fx$ 通过取商定义了同态
$\gr_k(M)\to\gr_{k+1}(M)$，从而定义了一个分次同态 $g$，其次数为 $1$，
并从 $\gr_\bullet(M)$ 映到自身；而说 $\varphi$ 是单射，也正是说 $g$ 是单射。
若再假设 $M$ 的滤过是分离的，则由此得到 $f_M$ 是单射
（Bourbaki，\emph{Alg.~comm.}，chap.~III，\S~2，no~8，cor.~1 du th.~1）。
\end{proof}

\begin{corollary}[15.1.3]
\label{0.15.1.3-fr}
设 $A$ 是一个 Noether 环，$M$ 是一个有限型 $A$-模。元素 $f$ 是
$M$ 中的非零因子，当且仅当 $\varphi$ 是双射。
\end{corollary}

\begin{proof}
只需证明该条件是充分的；为证明 $f_M$ 是单射，只需证明：对每个素理想
$\mathfrak p$（作为 $A$ 的素理想），比为 $f'=f/1\in A'=A_{\mathfrak p}$ 的位似
在 $A_{\mathfrak p}$-模 $M_{\mathfrak p}=M'$ 上是单射
（Bourbaki，\emph{Alg.~comm.}，chap.~II，\S~3，no~4，th.~1）。不过有
$f'^kM'/f'^{k+1}M'=(f^kM/f^{k+1}M)\otimes_A A'$，而且立即可见，典范同态
$\varphi':(M'/f'M')[T]\to\gr_\bullet(M')$ 等于
$\varphi\otimes 1_{A'}$；所以，如果 $\varphi$ 是单射，$\varphi'$ 也是单射
（$0_{\mathrm I}$，1.3.2）。于是归结到 $A$ 是 Noether \emph{局部}环的情形，
其极大理想记为 $\mathfrak m$。若 $f\notin\mathfrak m$，则 $f$ 可逆，无须证明。
若 $f\in\mathfrak m$，则已知 $M$ 对 $(fA)$-预进制拓扑是分离的
（$0_{\mathrm I}$，7.3.5），所以只需应用 (15.1.2)。
\end{proof}

\begin{definition}[15.1.4]
\label{0.15.1.4-fr}
设 $A$ 是一个环，$M$ 是一个 $A$-模。称元素 $f\in A$ 是
\emph{$M$-正则的}（相应地，\emph{$M$-拟正则的}），如果 $f$ 是
$M$ 中的非零因子（相应地，如果同态 (15.1.2.1) 是双射）。
\end{definition}

所以，我们刚刚证明了每个 $M$-正则元素都是 $M$-拟正则的；当 $A$ 为 Noether 环且
$M$ 是有限型 $A$-模时，逆命题也成立。当不再假设 $M$ 为有限型时，逆命题可能不成立：
例如，若 $A$ 是一个不是域的主理想整环，$K$ 是它的分式域，$M$ 是 $A$-模
$K/A$，则 $fM=M$ 对每个非零元 $f\neq0$（取自 $A$）都成立；因而 (15.1.2.1) 的两端均化为
$0$，可是此时 $f$ 是 $M$ 中的零因子。

\begin{env}[15.1.5]
\label{0.15.1.5-fr}
\oldpage[0\textsubscript{IV-1}]{14}
现在假设 $M$ 赋有一个\emph{滤过} $(M_k)$，满足 $M_0=M$；以
$\gr_\bullet(M)=\bigoplus_{k\geq0}M_k/M_{k+1}$ 表示相伴分次 $A$-模。
记号仍如 (15.1.1)；对每个 $k\geq0$，置
\[
  M'_k=M_k+\mathfrak{J}M_{k-1}+\cdots+
  \mathfrak{J}^{k-1}M_1+\mathfrak{J}^kM_0.
  \tag{15.1.5.1}\label{0.15.1.5.1-fr}
\]
显然，$(M'_k)$ 是 $M$ 上的一个滤过。当
$M_k=\mathfrak{Q}^kM$，其中 $\mathfrak Q$ 是 $A$ 的一个理想时，滤过
$(M'_k)$ 就是 $(\mathfrak{J}+\mathfrak{Q})$-预进制滤过。以
$\gr'_\bullet(M)$ 表示分次 $A$-模，它与 $(M'_k)$ 相伴。再以
\[
  (\gr_\bullet(M)\otimes_A(A/\mathfrak{J}))[T_1,\ldots,T_n]
\]
表示分次 $A$-模的张量积（$\mathrm{II}$，2.1.2）
\[
  (\gr_\bullet(M)\otimes_A(A/\mathfrak{J}))
  \otimes_A A[T_1,\ldots,T_n].
\]
定义一个分次 $A$-模的\emph{满同态}，其次数为 $0$：
\[
  \psi:(\gr_\bullet(M)\otimes_A(A/\mathfrak{J}))[T_1,\ldots,T_n]
  \longrightarrow\gr'_\bullet(M)
  \tag{15.1.5.2}\label{0.15.1.5.2-fr}
\]
如下：若 $\alpha\in A/\mathfrak{J}$ 是元素 $a\in A$ 的类，
$\xi\in M_k/M_{k+1}$ 是元素 $x\in M_k$ 的类，而
$P[T_1,\ldots,T_n]$ 是一个次数为 $h$ 的齐次多项式，则使三元组
$(\alpha,\xi,P)$ 对应于模 $M'_{h+k+1}$ 的类，该类以
$P(f_1,\ldots,f_n)ax$ 为代表元；仍可立即验证，这个元素只依赖于 $\xi$、
$\alpha$ 和 $P$，并且这就定义了分次同态 (15.1.5.2)，它由 (15.1.5.1)
的定义即为满射。对 $M$ 考虑满足 $M_1=0$ 的滤过，就重新得到同态 (15.1.1.1)。
\end{env}

\begin{lemma}[15.1.6]
\label{0.15.1.6-fr}
沿用 (15.1.5) 的记号，并假设 $n=1$，置 $f=f_1$。若 $f$ 是
$\gr_\bullet(M)$-正则的，则同态 (15.1.5.2) 是双射。
\end{lemma}

\begin{proof}
以 $Q_k$（相应地，$Q'_k$）表示 (15.1.5.2) 左端（相应地，右端）的
$A$-子模，其元素均由次数为 $k$ 的项组成。赋予 $Q_k$ 滤过
\[
  (Q_k)_i=\sum_{j\leq k-i}
  (\gr_{k-j}(M)\otimes_A(A/fA))T^j
  \tag{15.1.6.1}\label{0.15.1.6.1-fr}
\]
使得 $(Q_k)_0=Q_k$ 且 $(Q_k)_{k+1}=0$；对这个滤过，有
\[
  \gr_i(Q_k)=(\gr_i(M)\otimes_A(A/fA))T^{k-i}.
\]
赋予 $Q'_k$ 由 $\psi((Q_k)_i)=(Q'_k)_i$ 所成的滤过。由于这些滤过是
\emph{有限的}，只需证明同态
$\psi_{ki}:\gr_i(Q_k)\to\gr_i(Q'_k)$ 是单射；该同态由 $\psi_k$ 诱导
（Bourbaki，\emph{Alg.~comm.}，chap.~III，\S~3，no~8，cor.~1 du th.~1）。不过有
\[
  \gr_i(Q_k)=((M_i/M_{i+1})\otimes_A(A/fA))T^{k-i}~;
\]
另一方面，$(Q'_k)_{i+1}$ 是
\[
  M_k+fM_{k-1}+\cdots+f^{k-i-1}M_{i+1}
\]
\oldpage[0\textsubscript{IV-1}]{15}
在 $M'_k/M'_{k+1}$ 中的像。因此归结为说明：对 $x\in M_i$，关系
\[
  f^{k-i}x\in M_k+fM_{k-1}+\cdots+
  f^{k-i-1}M_{i+1}+M'_{k+1}
  \tag{15.1.6.2}\label{0.15.1.6.2-fr}
\]
蕴涵 $x\in fM_i+M_{i+1}$。不过，(15.1.6.2) 的右端包含于
$M_{i+1}+M'_{k+1}$；而由 (15.1.5.1)，$M'_{k+1}$ 包含于
$M_{i+1}+f^{k+1-i}M$。关于 $f$ 的假设蕴涵 $f$ 在 $M/M_h$ 中
不是零因子，不论 $h$ 取何值（Bourbaki，\emph{loc.~cit.}）。所以，关系
$f^{k-i}x\in f^{k+1-i}M+M_{i+1}$ 蕴涵（在 $M/M_{i+1}$ 中推理）存在
$y\in M$，使 $x-fy\in M_{i+1}$；又因 $x\in M_i$，故 $fy\in M_i$，
从而（在 $M/M_i$ 中推理）$y\in M_i$，最后得到 $x\in fM_i+M_{i+1}$。
证毕。
\end{proof}

\begin{definition}[15.1.7]
\label{0.15.1.7-fr}
设 $(f_i)_{1\leq i\leq n}$ 是 $A$ 的元素序列。称它是
\emph{$M$-正则的}，是指对每个 $1\leq i\leq n$，元素 $f_i$ 都是
$\left(M/(\sum_{1\leq j\leq i-1}f_jM)\right)$-正则的。称序列 $(f_i)$ 是
\emph{$M$-拟正则的}，是指 (15.1.1.1) 中定义的典范同态 $\varphi$ 为双射。

当 $M=A$ 时，简称为“正则序列”和“拟正则序列”。
\end{definition}

\begin{proposition}[15.1.8]
\label{0.15.1.8-fr}
沿用 (15.1.5) 的记号，假设序列 $(f_i)_{1\leq i\leq n}$ 是
$\gr_\bullet(M)$-正则序列。则典范同态 (15.1.5.2) 是双射。
\end{proposition}

\begin{proof}
当 $n=1$ 时，这个命题就是 (15.1.6)；现对 $n$ 作归纳。置
\[
  \mathfrak{J}''=\sum_{i=1}^{n-1}f_iA,
\]
于是 $\mathfrak{J}=\mathfrak{J}''+f_nA$；再置
\[
  M''_k=M_k+\mathfrak{J}''M_{k-1}+\cdots+
  \mathfrak{J}''^{k-1}M_1+\mathfrak{J}''^kM_0~;
\]
因而可以写成
\[
  M'_k=M''_k+f_nM''_{k-1}+\cdots+f_n^{k-1}M''_1+f_n^kM''_0.
\]
以 $\gr''_\bullet(M)$ 表示分次 $A$-模，它与滤过 $(M''_k)$（在 $M$ 上）相伴。
在典范同构的意义下，可以写成
\[
  (\gr_\bullet(M)\otimes(A/\mathfrak{J}))\otimes_A A[T_1,\ldots,T_n]
  =
  (\gr_\bullet(M)\otimes(A/\mathfrak{J}''))[T_1,\ldots,T_{n-1}]
  \otimes(A/f_nA)\otimes_A A[T_n]
\]
从而 (15.1.5.2) 分解为
\[
  (\gr''_\bullet(M)\otimes(A/f_nA))[T_n]\longrightarrow\gr'_\bullet(M)
  \tag{15.1.8.1}\label{0.15.1.8.1-fr}
\]
以及
\[
  ((\gr_\bullet(M)\otimes(A/\mathfrak{J}''))[T_1,\ldots,T_{n-1}])
  \otimes(A/f_nA)\otimes_A A[T_n]
  \xrightarrow{\psi'\otimes1}
  \gr''_\bullet(M)\otimes(A/f_nA)\otimes_A A[T_n]
\]
其中 $\psi'$ 是 (15.1.5.2) 中把 $n$ 换成 $n-1$、把
$\mathfrak{J}$ 换成 $\mathfrak{J}''$ 并把 $\gr'_\bullet(M)$ 换成
$\gr''_\bullet(M)$ 所得的映射。归纳假设蕴涵 $\psi'$ 是双射，故只须证明
(15.1.8.1) 也是双射。应用 (15.1.6)，只须证明 $f_n$ 是
$\gr''_\bullet(M)$-正则的。不过 $\gr''_\bullet(M)$ 等同于
\[
  (\gr_\bullet(M)\otimes(A/\mathfrak{J}''))
  [T_1,\ldots,T_{n-1}]~;
\]
依假设，比为 $f_n$ 的位似在 $\gr_\bullet(M)\otimes(A/\mathfrak{J}'')$ 中是
单射；因此，比为 $f_n$ 的位似在 $\gr''_\bullet(M)$ 中也同样是单射，因为
$A[T_1,\ldots,T_{n-1}]$ 是自由 $A$-模
\SourceCorrectionNote{法文印本及外交式转录先误作“双射”，再以“同样”沿用； 两处均为“单射”。正则性仅给出 单射性，而自由多项式扩张保持单射性。}。
\end{proof}

\begin{proposition}[15.1.9]
\label{0.15.1.9-fr}
设 $M$ 是 $A$-模，$(f_i)_{1\leq i\leq n}$ 是 $A$ 的元素序列。
若序列 $(f_i)$ 是 $M$-正则的，则它是 $M$-拟正则的。反之，若还假设
$A$-模
\[
  M,\quad M/f_1M,\quad\ldots,\quad
  M/\left(\sum_{1\leq j\leq n-1}f_jM\right)
\]
对于 $\mathfrak{J}$-预进拓扑是分离的（其中
$\mathfrak{J}=\sum_{i=1}^n f_iA$），则逆命题也成立。
\end{proposition}

\begin{proof}
第一个断言由 (15.1.8) 得出：在其中取 $M_1=0$，于是同态
(15.1.5.2) 化为 (15.1.1.1)。反过来，假设 (15.1.1.1) 的同态
$\varphi$ 是双射。注意，若置
\[
  \mathfrak{J}''=\sum_{i=1}^{n-1}f_iA,
\]
\oldpage[0\textsubscript{IV-1}]{16}
并以 $\gr''_\bullet(M)$ 表示分次 $A$-模，它与 $M$（赋有
$\mathfrak{J}''$-预进制滤过）相伴，则 $\varphi$ 分解为
\[
  (\gr''_\bullet(M)\otimes(A/f_nA))[T_n]\longrightarrow\gr_\bullet(M)
  \tag{15.1.9.1}\label{0.15.1.9.1-fr}
\]
以及 $\varphi'\otimes1$，其中 $\varphi'$ 是典范同态
\[
  (\gr''_0(M))[T_1,\ldots,T_{n-1}]
  \longrightarrow\gr''_\bullet(M).
\]
由于这两个同态都是满射，在假设 $\varphi$ 为双射时，它们也都是双射；
所以 $\varphi'$ 必为单射，并且由于它是满射，因而是双射。于是可以对 $n$
作归纳（$n=1$ 的情形由 (15.1.2) 得出），因为在
\[
  M,\quad M/f_1M,\quad\ldots,\quad
  M/\left(\sum_{j=1}^{n-2}f_jM\right)
\]
上，$\mathfrak{J}''$-预进拓扑细于 $\mathfrak{J}$-预进拓扑，故是分离的。
由此可见，序列 $(f_i)_{1\leq i\leq n-1}$ 是 $M$-正则的。另一方面，
我们来证明：该假设蕴涵，如果 $y\in M/\mathfrak{J}''M$ 满足
\[
  f_n^ky\in f_n^{k+1}(M/\mathfrak{J}''M),
\]
那么 $y\in f_n(M/\mathfrak{J}''M)$。事实上，若 $x$ 是 $y$ 的一个代表元，
我们来证明 $f_n^kx\in\mathfrak{J}^{k+1}M$；这样，我们的断言便由
(15.1.9.1) 的单射性得出。该假设蕴涵
\[
  f_n^kx\in f_n^{k+1}M+\mathfrak{J}''M\subset\mathfrak{J}M,
\]
所以，由于 $\varphi$ 是单射，
$f_n^{k-1}x\in\mathfrak{J}^kM$；如此归纳便得到
$x\in\mathfrak{J}M$，从而最终得到
$f_n^kx\in\mathfrak{J}^{k+1}M$。最后像 (15.1.2) 中那样结束论证：
在 $M/\mathfrak{J}''M$ 上，$\mathfrak{J}$-预进制滤过与
$(f_nA)$-预进制滤过相同，故依假设是分离的；于是得出比为 $f_n$ 的位似在
$M/\mathfrak{J}''M$ 中是单射，这就证明了序列
$(f_i)_{1\leq i\leq n}$ 是 $M$-正则的。
\end{proof}

\begin{corollary}[15.1.10]
\label{0.15.1.10-fr}
沿用 (15.1.9) 的记号。假设对每个有限型 $A$-模 $N$，
$M\otimes_A N$ 对于 $\mathfrak{J}$-预进拓扑都是分离的。则序列
$(f_i)$ 是 $M$-正则的，当且仅当它是 $M$-拟正则的。
\end{corollary}

应当指出，(15.1.10) 的假设蕴涵：若序列
$(f_i)_{1\leq i\leq n}$ 是 $M$-正则的，则序列 $(f_{\sigma(i)})$ 也同样如此，
其中 $\sigma$ 是区间 $[1,n]$ 的任意置换。

\begin{corollary}[15.1.11]
\label{0.15.1.11-fr}
设 $A$ 是 Noether 环，$M$ 是有限型 $A$-模，
$(f_i)_{1\leq i\leq n}$ 是由 $A$ 的 Jacobson 根中元素组成的有限序列。
则序列 $(f_i)$ 是 $M$-正则的，当且仅当它是 $M$-拟正则的。
\end{corollary}

\begin{proof}
这是 (15.1.10) 的一个特殊情形，因为这时每个有限型 $A$-模对于
$\mathfrak{J}$-预进拓扑都是分离的（$0_{\mathrm I}$，7.3.5）。
\end{proof}

\begin{remarks}[15.1.12]
\label{0.15.1.12-fr}
\mbox{}
\begin{enumerate}
  \item[(i)] 当 $A$ 是 Noether 局部环、$M$ 是有限型 $A$-模，并且各
  $f_i$ 都属于 $A$ 的极大理想时，J.-P.~Serre 以 $M$-\emph{序列}
  之名引入了 $M$-正则序列的概念 [17]。
  \item[(ii)] 即使 $A$ 是 Noether 的且 $M=A$，拟正则序列 $(f,g)$
  （由 $A$ 的两个元素组成）也未必正则。例如，序列 $(0,1)$ 是拟正则的，因为这时
  $\mathfrak{J}=fA+gA=A$，故 $\gr_\bullet(A)=0$；但若 $A\neq0$，它并不正则。
  不过应当指出，在这种情况下序列 $(1,0)$ 是正则的，因为 $1$ 不能把任何非零元乘成
  $0$（这里是在 $A$ 中而言），所以它不是零因子；并且由于 $A/A=0$，$0$ 在这个模中不是零因子
  （参见 (15.1.1)）。
  \item[(iii)]
  \oldpage[0\textsubscript{IV-1}]{17}
  在 (15.1.6) 中，一般不能把 $f$ 是 $\gr_\bullet(M)$-正则的这一假设换成
  它是 $\gr_\bullet(M)$-拟正则的这一假设。为说明这一点，仍取 (15.1.4)
  的例子，其中 $A$ 是分式域为 $K\neq A$ 的主理想整环，并取 $M=K$；滤过
  $(M_k)$（在 $M$ 上）定义如下：
  $M_0=M=K$，$M_1=A$，$M_k=0$（$k\geq2$），所以
  $\gr_0(M)=K/A$，$\gr_1(M)=A$，$\gr_k(M)=0$（$k\geq2$）。若
  $f\neq0$（在 $A$ 中），我们已经看到 $f$ 不是
  $\gr_\bullet(M)$-正则的；但由于
  \[
    f^k\gr_\bullet(M)/f^{k+1}\gr_\bullet(M)=f^kA/f^{k+1}A,
  \]
  可立即看出同态 (15.1.2.1)（其中把 $M$ 换成 $\gr_\bullet(M)$）是单射，
  所以 $f$ 是 $\gr_\bullet(M)$-拟正则的。不过，沿用 (15.1.6) 的记号，
  $M'_k=K$ 对每个 $k\geq0$ 都成立，因而 $\gr'_\bullet(M)=0$，而
  (15.1.5.2) 的左端并不化为 $0$。
\end{enumerate}
\end{remarks}



\begin{proposition}[15.1.13]
\label{0.15.1.13-fr}
设 $A$ 是一个环，$M$、$N$ 是两个 $A$-模，$(f_i)_{1\leq i\leq n}$
是 $A$ 的元素序列。若序列 $(f_i)$ 是 $M$-正则的，并且 $N$ 是平坦
$A$-模，则序列 $(f_i)$ 是 $M\otimes_A N$-正则的。若 $N$ 是忠实平坦
$A$-模，则逆命题也成立。
\end{proposition}

\begin{proof}
事实上（不对 $N$ 作任何假设），
\[
  (M\otimes_A N)/
  \left(\sum_{j=1}^{i-1}f_j(M\otimes_A N)\right)
\]
同构于
\[
  \left(M/\left(\sum_{j=1}^{i-1}f_jM\right)\right)\otimes_A N~;
\]
于是只须把（$0_{\mathrm I}$，6.1.1）和（$0_{\mathrm I}$，6.4.1）
应用于比为 $f_i$、作用在 $M/(\sum_{j=1}^{i-1}f_jM)$ 中的位似。
\end{proof}

\begin{corollary}[15.1.14]
\label{0.15.1.14-fr}
设 $A$ 是一个环，$M$ 是 $A$-模，$(f_i)_{1\leq i\leq n}$ 是 $A$
的有限元素序列。设 $\rho:A\to B$ 是环同态，$M'=M_{(B)}=M\otimes_A B$，
且置 $f'_i=\rho(f_i)$（$1\leq i\leq n$）。若 $B$ 是平坦 $A$-模且序列
$(f_i)$ 是 $M$-正则的，则序列 $(f'_i)$ 是 $M'$-正则的；若 $B$ 是忠实平坦
$A$-模，则逆命题也成立。
\end{corollary}

\begin{proof}
事实上，
\[
  M'/\left(\sum_{j=1}^{i-1}f'_jM'\right)
\]
等同于
\[
  \left(M/\left(\sum_{j=1}^{i-1}f_jM\right)\right)\otimes_A B,
\]
而比为 $f'_i$ 的位似作用在 $B$-模 $M'/(\sum_{j=1}^{i-1}f'_jM')$ 上；它等同于
比为 $f_i$ 的位似作用在如下 $A$-模上：
\[
  \left(M/\left(\sum_{j=1}^{i-1}f_jM\right)\right)\otimes_A B~;
\]
所以断言由 (15.1.13) 得出。
\end{proof}

\begin{proposition}[15.1.15]
\label{0.15.1.15-fr}
设 $A$ 是一个环，$B$ 是交换 $A$-代数，
$(f_i)_{1\leq i\leq n}$ 是 $B$ 的有限元素序列，$M$ 是 $B$-模。
假设序列 $(f_i)$ 是 $M$-正则的，并且每个 $A$-模
\[
  M/\left(\sum_{j=1}^{i-1}f_jM\right)\quad(1\leq i\leq n)
\]
都是平坦的。那么，对于任意环同态 $\rho:A\to A'$，若置
$B'=B\otimes_A A'$、$M'=M\otimes_A A'$ 和 $f'_i=f_i\otimes1$
（$1\leq i\leq n$），则序列 $(f'_i)$ 是 $M'$-正则的，并且各
$A'$-模
\[
  M'/\left(\sum_{j=1}^{i-1}f'_jM'\right)\quad(1\leq i\leq n)
\]
都是平坦的。
\end{proposition}

\begin{proof}
考虑下列序列（依假设它是正合的）：
\[
  0\longrightarrow M\xrightarrow{f_1}M\longrightarrow M/f_1M
  \longrightarrow0
\]
（箭头 $M\xrightarrow{f_1}M$ 是比为 $f_1$、作用在 $M$ 中的位似）。假设
$M/f_1M$ 是平坦 $A$-模，这蕴涵序列
\[
  0\longrightarrow M\otimes_A A'
  \xrightarrow{f_1\otimes1}M\otimes_A A'
  \longrightarrow(M/f_1M)\otimes_A A'\longrightarrow0
\]
\oldpage[0\textsubscript{IV-1}]{18}
是正合的（$0_{\mathrm I}$，6.1.2）。这说明比为 $f'_1$ 的位似在 $M'$ 中是单射。

其次，对 $1\leq i\leq n-1$ 置
\[
  M_i=M/\left(\sum_{j=1}^i f_jM\right),\qquad
  M'_i=M'/\left(\sum_{j=1}^i f'_jM'\right)~;
\]
有 $M'_i=M_i\otimes_A A'$、$M_i/f_{i+1}M_i=M_{i+1}$、
$M'_i/f'_{i+1}M'_i=M'_{i+1}$；只须在上面的论证中把 $M$、$f_1$
分别换成 $M_i$、$f_{i+1}$，便可得出比为 $f'_{i+1}$ 的位似在
$M'_i$ 中是单射，这是由于 $M_i/f_{i+1}M_i$ 是平坦的。最后一个断言由
（$0_{\mathrm I}$，6.2.1）得出。
\end{proof}

\begin{proposition}[15.1.16]
\label{0.15.1.16-fr}
设 $A$、$B$ 是两个 Noether 局部环，$k$ 是 $A$ 的剩余域，
$\varphi:A\to B$ 是局部同态，$M$ 是有限型 $B$-模。
设 $(f_i)_{1\leq i\leq n}$ 是由 $B$ 的极大理想中元素组成的序列；对每个
$i$，设 $g_i$ 是 $f_i$ 在 $B\otimes_A k$ 中的像。下列条件等价：
\begin{enumerate}
  \item[\rm(a)] 序列 $(f_i)$ 是 $M$-正则的，并且各商模
  \[
    M_i=M/\left(\sum_{j=1}^i f_jM\right)
  \]
  都是平坦 $A$-模（$1\leq i\leq n$）。
  \item[\rm(b)] 序列 $(f_i)$ 是 $M$-正则的，并且 $A$-模 $M_n$ 是平坦的。
  \item[\rm(c)] $A$-模 $M$ 是平坦的，并且序列
  $(g_i)_{1\leq i\leq n}$ 是 $(M\otimes_A k)$-正则的。
  \item[\rm(d)] $A$-模 $M$ 是平坦的；并且对于任意环同态
  $\rho:A\to A'$，若置 $B'=B\otimes_A A'$、$M'=M\otimes_A A'$、
  $f'_i=f_i\otimes1$（$1\leq i\leq n$），则序列 $(f'_i)$ 是
  $M'$-正则的。
\end{enumerate}
\end{proposition}

\begin{proof}
显然 a) 蕴涵 b)。由 (15.1.15)，要证明 b) 蕴涵 d)，只须证明 b) 蕴涵各
$M_i$ 都是平坦 $A$-模（$0\leq i\leq n-1$；置 $M_0=M$）。由于
$M_{i+1}=M_i/f_{i+1}M_i$，用下降归纳法，问题归结为证明下述事实：
若 $h$ 是 $B$ 的极大理想中的一个元素，比为 $h$ 的位似在 $M$ 中是单射，
并且 $M/hM$ 是平坦 $A$-模，则 $M$ 是平坦 $A$-模；这恰是
（$0_{\mathrm{III}}$，10.2.7）。显然，c) 是 d) 在取 $A'=k$ 时所得的特殊情形。
所以只须证明 c) 蕴涵 a)。由于 $M$ 是平坦 $A$-模且是有限型 $B$-模，
由 c) 和（$0_{\mathrm{III}}$，10.2.4）可知，比为 $f_1$ 的位似在 $M$ 中是单射，
且其余核 $M/f_1M=M_1$ 是平坦 $A$-模。对 $i$ 作归纳，并假设已经证明
$M_i$ 是平坦 $A$-模；依假设，比为 $g_{i+1}$ 的位似在
$M_i\otimes_A k$ 中是单射，再由（$0_{\mathrm{III}}$，10.2.4）可知，
比为 $f_{i+1}$ 的位似在 $M_i$ 中是单射，并且
$M_i/f_{i+1}M_i=M_{i+1}$ 是平坦 $A$-模，这就完成了证明。应当指出，
在证明 c) 蕴涵 a) 时，并未使用各 $f_i$ 属于 $B$ 的极大理想这一事实。
\end{proof}

\begin{corollary}[15.1.17]
\label{0.15.1.17-fr}
在 (15.1.16) 关于 $A$、$B$ 的假设下，设 $M$ 是有限型 $B$-模且在 $A$
上平坦，$R$ 是一个 $B$-子模，包含于 $M$，并且 $N=M/R$ 是平坦 $A$-模。
设 $(f_i)_{1\leq i\leq n}$ 是由 $B$ 的极大理想中元素组成的序列，满足下列性质：
\begin{enumerate}
  \item[(i)] 有 $f_iM\subset R$（$1\leq i\leq n$）。
  \item[(ii)] 若 $g_i$ 是 $f_i$ 在 $B\otimes_A k$ 中的像，则序列
  $(g_i)_{1\leq i\leq n}$ 是 $(M\otimes_A k)$-正则的。
  \item[(iii)] 各 $(B\otimes_A k)$-子模
  $g_i(M\otimes_A k)$（它们都是 $M\otimes_A k$ 的子模）之和，
  等于 $R\otimes_A k$ 在 $M\otimes_A k$
  中的典范像（即典范同态 $M\otimes_A k\to N\otimes_A k$ 的核）。
\end{enumerate}
则序列 $(f_i)$ 是 $M$-正则的，并且
$R=\sum_{i=1}^n f_iM$。
\end{corollary}

\begin{proof}
\oldpage[0\textsubscript{IV-1}]{19}
由于 $N$ 是平坦 $A$-模，由正合序列 $0\to R\to M\to N\to0$ 的正合性可知，
序列
\[
  0\longrightarrow R\otimes_A k\longrightarrow M\otimes_A k
  \longrightarrow N\otimes_A k\longrightarrow0
\]
是正合的（$0_{\mathrm I}$，6.1.2）；所以
$R\otimes_A k=\sum_i g_i(M\otimes_A k)$。由于 $M$ 是有限型 $B$-模且
$B$ 是 Noether 环，$R$ 也是有限型 $B$-模；故由 Nakayama 引理得到
$R=\sum_i f_iM$，这里各 $f_i$ 属于 $B$ 的极大理想。
\end{proof}

\begin{lemma}[15.1.18]
\label{0.15.1.18-fr}
设 $A$ 是一个环，$M$ 是 $A$-模，赋有分离且穷竭滤过
$(M_k)_{k\geq0}$；并设 $f\in A$ 是一个 $\gr_\bullet(M)$-正则元。则：
\begin{enumerate}
  \item[(i)] $f$ 是 $M$-正则的。
  \item[(ii)] 对每个 $k$，有 $M_k\cap fM=fM_k$；因而典范像
  $M'_k$（即 $M_k$ 在 $M'=M/fM$ 中的像）典范同构于 $M_k/fM_k$。
  \item[(iii)] 序列
  \[
    0\longrightarrow M_{k+1}/fM_{k+1}\longrightarrow M_k/fM_k
    \longrightarrow\gr_k(M)/f\cdot\gr_k(M)\longrightarrow0
  \]
  是正合的；因此，若在 $M'$ 上赋予商滤过 $(M'_k)$，则
  $\gr_\bullet(M')$ 典范同构于
  $\gr_\bullet(M)/f\cdot\gr_\bullet(M)$。
\end{enumerate}
\end{lemma}

\begin{proof}
断言 (i) 由关于 $f$ 的假设以及 Bourbaki，
\emph{Alg.~comm.}，chap.~III，\S~2，no~8，cor.~1 du th.~1 得出。
由于 $\gr_\bullet(M/M_k)$ 是 $\gr_\bullet(M)$ 的直和因子，
$f$ 是 $\gr_\bullet(M/M_k)$-正则的（对每个 $k$），故由 (i) 可知它是
$(M/M_k)$-正则的，这就证明了 (ii)。为证明 (iii)，只须说明映射
$M_{k+1}/fM_{k+1}\to M_k/fM_k$ 是单射的：事实上，若
$x\in M_{k+1}$ 且 $x\in fM_k$，则由 (ii) 可知
$x\in fM_{k+1}$，引理得证。
\end{proof}

\begin{proposition}[15.1.19]
\label{0.15.1.19-fr}
设 $A$ 是一个环，$M$ 是 $A$-模，赋有滤过 $(M_k)_{k\geq0}$，其中
$M_0=M$；设 $(f_i)_{1\leq i\leq n}$ 是一个
$\gr_\bullet(M)$-正则的、由 $A$ 的元素组成的序列。还假设对
$0\leq i\leq n-1$，滤过 $(M_k)$
在
\[
  M/\left(\sum_{j\leq i}f_jM\right)
\]
上所诱导的商滤过是分离的。则序列 $(f_i)$ 是 $M$-正则的；若置
\[
  M'=M/\left(\sum_{i\leq n}f_iM\right)
\]
并在 $M'$ 上赋予 $(M_k)$ 的商滤过，则 $\gr_\bullet(M')$ 典范同构于
\[
  \gr_\bullet(M)/
  \left(\sum_{i\leq n}f_i\gr_\bullet(M)\right).
\]
\end{proposition}

\begin{proof}
当 $n=1$ 时，命题由引理 (15.1.18) 得出。对 $n$ 作归纳，置
\[
  N=M/\left(\sum_{i\leq n-1}f_iM\right)
\]
并用 $(N_k)$ 表示 $(M_k)$ 在 $N$ 上的商滤过；依假设，该滤过是分离的。
此外，$\gr_\bullet(N)$ 同构于
\[
  \gr_\bullet(M)/
  \left(\sum_{i\leq n-1}f_i\gr_\bullet(M)\right).
\]
所以依假设，$f_n$ 是 $\gr_\bullet(N)$-正则的；把引理 (15.1.18) 应用于
$N$，可知 $f_n$ 是 $N$-正则的，并且 $\gr_\bullet(M')$ 同构于
$\gr_\bullet(N)/f_n\gr_\bullet(N)$，这就证明了命题。
\end{proof}

命题 (15.1.19) 尤其适用于满足下列任一假设的滤过：
\begin{enumerate}
  \item[$1^\circ$] 滤过 $(M_k)$ 是有限且分离的（因为这蕴涵
  $M_k=0$，其中 $k$ 足够大）；
  \item[$2^\circ$] $A$ 是 Noether 环，$\mathfrak{J}$ 是包含于 $A$ 的
  Jacobson 根中的理想，$M$ 是有限型 $A$-模，并且 $(M_k)$ 是
  $\mathfrak{J}$-预进制滤过（$0_{\mathrm I}$，7.3.5）。
\end{enumerate}

\begin{corollary}[15.1.20]
\label{0.15.1.20-fr}
\oldpage[0\textsubscript{IV-1}]{20}
设 $A$ 是一个环，$M$ 是 $A$-模，$(f_i)_{1\leq i\leq n}$ 是一个
$M$-正则序列，其各项属于 $A$。则对于任意整数 $\nu_i\geq1$，序列
$(f_i^{\nu_i})$ 是 $M$-正则的，并且
$M/(\sum_{i\leq n}f_i^{\nu_i}M)$ 有一条合成列，其各商模同构于
$M/(\sum_{i\leq n}f_iM)$。
\end{corollary}

\begin{proof}
对 $n$ 作归纳。当 $n=1$ 时，断言是立即的，因为
$f^jM/f^{j+1}M$ 同构于 $M/fM$，而依假设，比为 $f$ 的位似在 $M$ 中是单射。
置
\[
  N=M/\left(\sum_{i\leq n-1}f_i^{\nu_i}M\right)~;
\]
由归纳假设，$N$ 有一个有限滤过 $(N_k)$，使各 $\gr_k(N)$ 同构于
$M/(\sum_{i\leq n-1}f_iM)$；依假设，$f_n$ 因而是
$\gr_\bullet(N)$-正则的。所以由 (15.1.19) 可知，$f_n$ 是 $N$-正则的。
应用开头已经证明的 $n=1$ 情形，便得 $f_n^{\nu_n}$ 是 $N$-正则的，并且
$N/f_n^{\nu_n}N$ 有一条合成列，其各商模同构于 $N/f_nN$；而对于
$N/f_nN$ 上由滤过 $(N_k)$ 诱导的商滤过，由 (15.1.19)，
$\gr_k(N/f_nN)$ 同构于 $\gr_k(N)/f_n\gr_k(N)$（对每个 $k$），从而同构于
$M/(\sum_{i\leq n}f_iM)$，这就证明了推论。
\end{proof}

\begin{proposition}[15.1.21]
\label{0.15.1.21-fr}
设 $A$、$B$ 是两个 Noether 局部环，$\varphi:A\to B$ 是局部同态，
$M$ 是有限型 $B$-模，$(f_i)_{1\leq i\leq n}$ 是由 $A$ 的极大理想中
元素组成的 $A$-正则序列。下列性质等价：
\begin{enumerate}
  \item[a)] $M$ 是平坦 $A$-模。
  \item[b)] 若 $\mathfrak{J}$ 是理想 $\sum_{i\leq n}f_iA$，则
  $M/\mathfrak{J}M$ 是平坦 $A/\mathfrak{J}$-模，并且序列 $(f_i)$ 是
  $M$-正则的。
\end{enumerate}
\end{proposition}

\begin{proof}
a) 蕴涵 b) 由 (15.1.13) 和（$0_{\mathrm{III}}$，10.2.2）得出。反过来，
若序列 $(f_i)$ 是 $M$-正则的，则它是 $M$-拟正则的 (15.1.9)，故同态
(15.1.1.1) 是双射。不过，由于还假设序列 $(f_i)_{1\leq i\leq n}$ 是
$A$-正则的，所以它是 $A$-拟正则的 (15.1.9)，相应的典范同态 (15.1.1.1)
$(A/\mathfrak{J})[T_1,\ldots,T_n]\to\gr_\bullet(A)$
是双射。由此可知，典范同态（$0_{\mathrm{III}}$，10.1.1.2）
\[
  \gr_0(M)\otimes_{A/\mathfrak{J}}\gr_\bullet(A)
  \longrightarrow\gr_\bullet(M)
\]
是双射。于是结论由（$0_{\mathrm{III}}$，10.2.2）得出，因为各 $f_i$
属于 $A$ 的极大理想，并且同态 $\varphi$ 是局部的。
\end{proof}

\subsection{\texorpdfstring{$\mathscr{F}$}{F}-正则序列}
\label{subsection:0.15.2-fr}

\begin{env}[15.2.1]
\label{0.15.2.1-fr}
设 $X$ 是赋环空间，$\mathscr{F}$ 是 $\mathcal{O}_X$-模层，
$f_i$（$1\leq i\leq n$）是 $\mathcal{O}_X$ 在 $X$ 上的截面所成的序列，
$\mathscr{J}$ 是 $\mathcal{O}_X$ 的理想层，由各 $f_i$ 生成
（即 $\mathcal{O}_X^n$ 在典范同态
$\mathcal{O}_X^n\to\mathcal{O}_X$ 下的像；该同态由各 $f_i$ 定义
（$0_{\mathrm I}$，5.1.1））。子-$\mathcal{O}_X$-模层
$\mathscr{J}^k\mathscr{F}$ 在 $\mathscr{F}$ 上定义一个滤过，并且仍置
$\gr_k(\mathscr{F})=\mathscr{J}^k\mathscr{F}/\mathscr{J}^{k+1}\mathscr{F}$；
若 $S_k$ 是总次数为 $k$ 的齐次多项式的集合，这些多项式取自环
$\mathbb{Z}[T_1,\ldots,T_n]$，则可将 $S_k$ 看作 $X$ 上的常值层，并且如
(15.1.1) 那样定义一个
典范同态
\[
\label{0.15.2.1.1-fr}
  \varphi_k:\gr_0(\mathscr{F})\otimes_{\mathbb Z}S_k
  \longrightarrow\gr_k(\mathscr{F}).
\tag{15.2.1.1}
\]
\oldpage[0\textsubscript{IV-1}]{21}
其定义如下：对任一开集 $U$（它是 $X$ 的开集；也可只取 $X$ 的一个
拓扑基中的开集），考虑同态 (15.1.1.1)
\[
  \varphi_{k,U}:\gr_0(\Gamma(U,\mathscr{F}))\otimes_{\mathbb Z}S_k
  \longrightarrow\gr_k(\Gamma(U,\mathscr{F}))
\]
它由 $n$ 个截面 $f_i|U\in\Gamma(U,\mathcal{O}_X)$ 定义，其中
$\Gamma(U,\mathcal{O}_X)$-模 $\Gamma(U,\mathscr{F})$ 带有由子模
$(\Gamma(U,\mathscr{J}))^k\Gamma(U,\mathscr{F})$ 所定义的滤过。
$\mathcal{O}_X$-模层 $\gr_k(\mathscr{F})$ 与预层
$U\mapsto\gr_k(\Gamma(U,\mathscr{F}))$ 相伴（对每个 $k$），所以各 $\varphi_{k,U}$
定义了 $\mathcal{O}_X$-模层的同态 (15.2.1.1)。由模的归纳极限的正合性以及
归纳极限与张量积可交换，以上说明还表明：对于每个 $x\in X$，
(15.2.1.1) 两端的茎之间的同态 $(\varphi_k)_x$（这些茎取在点 $x$），可等同于
(15.1.1) 中的同态
\[
\label{0.15.2.1.2-fr}
  \gr_0(\mathscr{F}_x)\otimes_{\mathbb Z}S_k
  \longrightarrow\gr_k(\mathscr{F}_x)
\tag{15.2.1.2}
\]
它由序列 $((f_i)_x)_{1\leq i\leq n}$ 定义，该序列由
$\mathcal{O}_x$ 中的元素组成。
\end{env}

\begin{env}[15.2.2]
\label{0.15.2.2-fr}
记号如 (15.2.1)。我们称序列 $(f_i)_{1\leq i\leq n}$ 是
$\mathscr{F}$-正则的，如果如下条件成立：以 $f_j\mathscr{F}$ 表示
$\mathcal{O}_X$-模层 $\mathscr{F}$ 的一个子模层，即 $\mathscr{F}$ 的一个自同态
之像（该自同态由 $f_j$ 定义），则在
$\mathscr{F}/(\sum_{j=1}^{i-1}f_j\mathscr{F})$ 上由 $f_i$ 定义的自同态，
对 $1\leq i\leq n$ 是单射；若同态 (15.2.1.1) 都是单射，则称序列
$(f_i)$ 是 $\mathscr{F}$-拟正则的。由这些定义和
(15.2.1) 的说明可知，序列 $(f_i)$ 是 $\mathscr{F}$-正则的（相应地，
$\mathscr{F}$-拟正则的），当且仅当对于每个 $x\in X$，序列
$((f_i)_x)_{1\leq i\leq n}$ 是 $\mathscr{F}_x$-正则的（相应地，
$\mathscr{F}_x$-拟正则的）。所以每个 $\mathscr{F}$-正则序列都是
$\mathscr{F}$-拟正则的 (15.1.9)；若 $\mathscr{F}$ 是有限型
$\mathcal{O}_X$-模层、$\mathcal{O}_x$ 对每个 $x\in X$ 都是 Noether
环，并且各 $(f_i)_x$ 还属于 $\mathcal{O}_x$ 的 Jacobson 根（对每个
$x\in X$），则逆命题也成立 (15.1.11)。

若 $X=\operatorname{Spec}(A)$ 是仿射概形且
$\mathscr{F}=\widetilde{M}$，则由（I，1.3）可知，序列 $(f_i)$ 是
$\mathscr{F}$-正则的（相应地，$\mathscr{F}$-拟正则的），当且仅当它是
$M$-正则的（相应地，$M$-拟正则的）。

当 $\mathscr{F}=\mathcal{O}_X$ 时，在上述定义中省略 $\mathscr{F}$。
\end{env}

\begin{remark}[15.2.3]
\label{0.15.2.3-fr}
当 $X=\operatorname{Spec}(A)$ 且
$\mathfrak{J}=\sum_{i=1}^n f_iA$ 时，序列 $(f_j)$ 对于任意 $A$-模
$M$，在任一开集 $U$ 上都是拟正则的，只要该开集不与
$V(\mathfrak{J})$ 相交；
这是因为对于每个 $x\in U$，都有 $\mathfrak{J}_x=\mathcal{O}_x$，并且
(15.2.1.2) 的两端都为零。因此，$\widetilde{M}$-正则序列的概念只有在
$V(\mathfrak{J})$ 的各点附近才有意义。
\end{remark}

\begin{proposition}[15.2.4]
\label{0.15.2.4-fr}
设 $X$ 是赋环空间，$\mathscr{F}$ 是凝聚 $\mathcal{O}_X$-模层，$x$ 是
$X$ 的一点，$(f_i)_{1\leq i\leq n}$ 是
$\Gamma(X,\mathcal{O}_X)$ 中元素所成的序列。若序列
$((f_i)_x)_{1\leq i\leq n}$ 是 $\mathscr{F}_x$-正则的，则存在一个开邻域
$U$，它包含 $x$，并使得序列 $(f_i|U)_{1\leq i\leq n}$ 是
$(\mathscr{F}|U)$-正则的。
\end{proposition}

\begin{proof}
事实上，对 $1\leq i\leq n$，$\mathcal{O}_X$-模层
$\mathscr{F}/(\sum_{j=1}^{i-1}f_j\mathscr{F})$ 是凝聚的
（$0_{\mathrm I}$，5.3.4），所以这个 $\mathcal{O}_X$-模层的自同态
（由 $f_i$ 定义）之核也是凝聚的（$0_{\mathrm I}$，5.3.4）；
由于该核在点 $x$ 处的茎为零，它在 $x$ 的某个邻域上的限制也为零
（$0_{\mathrm I}$，5.2.2）。
\end{proof}

\begin{proposition}[15.2.5]
\label{0.15.2.5-fr}
\oldpage[0\textsubscript{IV-1}]{22}
设 $u=(\psi,\theta):X\to Y$ 是赋环空间的平坦态射
（$0_{\mathrm I}$，6.7.1）；设 $\mathscr{F}$ 是 $\mathcal{O}_Y$-模层，
$\mathscr{G}=u^*(\mathscr{F})$，$(f_i)_{1\leq i\leq n}$ 是
$\Gamma(Y,\mathcal{O}_Y)$ 中元素所成的序列，而
$(g_i)_{1\leq i\leq n}$ 是它们的像所成的序列，其中
$g_i=\theta(f_i)\in\Gamma(Y,u_*(\mathcal{O}_X))=\Gamma(X,\mathcal{O}_X)$。
若序列 $(f_i)$ 是 $\mathscr{F}$-正则的，则序列 $(g_i)$ 是
$\mathscr{G}$-正则的。当 $u$ 是忠实平坦态射时，逆命题也成立
（$0_{\mathrm I}$，6.7.8）。
\end{proposition}

\begin{proof}
设 $x$ 是 $X$ 的一点，$y=u(x)=\psi(x)$；有
\[
  \mathscr{G}_x=\psi_x^*(\mathscr{F}_y)
  \otimes_{\psi_x^*(\mathcal{O}_y)}\mathcal{O}_x,
  \qquad\text{且}\qquad
  (g_i)_x=\theta_x^\#(\psi_x^*((f_i)_y))~;
\]
依假设，$\theta_x^\#:\psi_x^*(\mathcal{O}_y)\to\mathcal{O}_x$ 使
$\mathcal{O}_x$ 成为平坦 $\psi_x^*(\mathcal{O}_y)$-模，并且函子
$\psi^*$ 是正合的，所以第一个断言由 (15.1.14) 的相应断言得出；第二个
断言同理，因为 $u$ 忠实平坦这一假设意指 $\psi$ 是满射，并且
$\theta_x^\#$ 使 $\mathcal{O}_x$ 成为忠实平坦
$\psi_x^*(\mathcal{O}_y)$-模。
\end{proof}

\section{Noether 局部环中的维数与深度}
\label{section:0.16-fr}

\subsection{环的维数}
\label{subsection:0.16.1-fr}

\begin{env}[16.1.1]
\label{0.16.1.1-fr}
把环 $A$ 的谱的（组合）维数称为该环的\emph{维数}（或 \emph{Krull 维数}），
记为 $\dim A$；这里所说的谱是 $\operatorname{Spec}(A)$ (14.2.1)。由于
$\operatorname{Spec}(A)$ 的不可约闭子集都是
$V(\mathfrak{p})$，其中 $\mathfrak{p}$ 是 $A$ 的素理想（Bourbaki，
\emph{Alg.~comm.}，第 II 章，\S~4，第 3 号，命题 14 的推论 1），
这也就是说，$\dim(A)$ 是 $A$ 的素理想链 (14.1.1)
\[
  \mathfrak{p}_0\subset\mathfrak{p}_1\subset\cdots\subset\mathfrak{p}_n
\]
之长度的上确界。若 $A\neq0$，则 $\dim A\geq0$。非零（$\neq0$）Artin 环
（特别地，域）的维数为 0（Bourbaki，\emph{Alg.~comm.}，第 IV 章，
\S~2，第 5 号，命题 7）；不是域的 Dedekind 整环（特别地，
$\mathbb{Z}$ 或一元多项式环 $k[T]$，其基域为 $k$）的维数为 1
（Bourbaki，\emph{Alg.~comm.}，第 VII 章，\S~2）。Noether 环不必是
有限维的 [30]。

若 $(\mathfrak{p}_\alpha)$ 是 $A$ 的极小素理想族，则由 (14.1.2.1) 有
\[
\label{0.16.1.1.1-fr}
  \dim(A)=\sup_\alpha(\dim(A/\mathfrak{p}_\alpha)).
\tag{16.1.1.1}
\]
\end{env}

\begin{env}[16.1.2]
\label{0.16.1.2-fr}
设 $\mathfrak{a}$ 是 $A$ 的任意理想；$A/\mathfrak{a}$ 的素理想与 $A$ 中
包含 $\mathfrak{a}$ 的素理想一一对应，故
\[
\label{0.16.1.2.1-fr}
  \dim(A/\mathfrak{a})\leq\dim(A).
\tag{16.1.2.1}
\]
若再假设 $\mathfrak{a}$ 不包含在任何极小素理想
$\mathfrak{p}_\alpha$ 中（这里的极小素理想取自 $A$），则 $A$ 中包含
$\mathfrak{a}$ 的任一素理想链，
都可在前端补上某个 $\mathfrak{p}_\alpha$，因而
\[
\label{0.16.1.2.2-fr}
  1+\dim(A/\mathfrak{a})\leq\dim(A).
\tag{16.1.2.2}
\]
\end{env}

\begin{env}[16.1.3]
\label{0.16.1.3-fr}
\oldpage[0\textsubscript{IV-1}]{23}
若 $S$ 是 $A$ 的乘法子集，则 $S^{-1}A$ 的素理想与 $A$ 中不与 $S$
相交的素理想一一对应，故
\[
\label{0.16.1.3.1-fr}
  \dim(S^{-1}A)\leq\dim(A).
\tag{16.1.3.1}
\]
对于任意素理想 $\mathfrak{p}$（在 $A$ 中），依定义 (14.2.1) 有
\[
\label{0.16.1.3.2-fr}
  \dim(A_{\mathfrak{p}})
  =\operatorname{codim}(V(\mathfrak{p}),\operatorname{Spec}(A)).
\tag{16.1.3.2}
\]
这个数也称为素理想 $\mathfrak{p}$ 的\emph{高度}，记作
$\operatorname{ht}(\mathfrak{p})$。更一般地，设 $\mathfrak{a}$ 是 $A$
的任意理想；称 $\mathfrak{a}$ 的\emph{高度}为下列数：
\[
\label{0.16.1.3.3-fr}
  \operatorname{ht}(\mathfrak{a})
  =\operatorname{codim}(V(\mathfrak{a}),\operatorname{Spec}(A)).
\tag{16.1.3.3}
\]
若 $(\mathfrak{p}_\lambda)$ 是所有包含 $\mathfrak{a}$ 的素理想中极小者所成
的族，则由 (14.2.1) 有
\[
\label{0.16.1.3.4-fr}
  \operatorname{ht}(\mathfrak{a})
  =\inf_\lambda(\operatorname{ht}(\mathfrak{p}_\lambda)).
\tag{16.1.3.4}
\]
显然有
\[
\label{0.16.1.3.5-fr}
  \dim(A)=\sup_{\mathfrak{m}}(\dim(A_{\mathfrak{m}}))
\tag{16.1.3.5}
\]
其中 $\mathfrak{m}$ 遍历 $A$ 的全部极大理想。
\end{env}

\begin{env}[16.1.4]
\label{0.16.1.4-fr}
对于任意素理想 $\mathfrak{p}$（在 $A$ 中），由 (14.2.2.2) 有
\[
\label{0.16.1.4.1-fr}
  \dim(A_{\mathfrak{p}})+\dim(A/\mathfrak{p})\leq\dim A.
\tag{16.1.4.1}
\]
称 $A$ 是\emph{悬链的}（分别地，\emph{等维的}、\emph{等余维的}、
\emph{双等维的}），如果 $\operatorname{Spec}(A)$ 是悬链的（分别地，
等维的、等余维的、双等维的）。当 $A$ 是 Noether 且双等维的时，
(16.1.4.1) 的两边对于 $A$ 的每个素理想都相等 (14.3.5)。

设 $\mathfrak{a}$ 是 $A$ 的任意理想，则 $A/\mathfrak{a}$ 的素理想与 $A$
中包含 $\mathfrak{a}$ 的素理想一一对应；因此，如果 $A$ 是悬链的，
$A/\mathfrak{a}$ 也同样如此。类似地，设 $S$ 是 $A$ 的任意乘法子集，
则 $S^{-1}A$ 的素理想与 $A$ 中不与 $S$ 相交的素理想一一对应；因此，
如果 $A$ 是悬链的，$S^{-1}A$ 也同样如此。

此外，由于 $A$ 中包含在素理想 $\mathfrak{p}$ 内的素理想与
$A_{\mathfrak{p}}$ 的素理想一一对应，要使 $A$ 是悬链的，充要条件是
$A_{\mathfrak{p}}$ 对每个 $\mathfrak{p}$ 都是悬链的。所以，称 $A$ 是
悬链的，意指由 (14.3.2) 和 (16.1.3.2)，对于任意一对素理想
$\mathfrak{p}$、$\mathfrak{q}$（均在 $A$ 中），只要
$\mathfrak{q}\subset\mathfrak{p}$，就有
\[
\label{0.16.1.4.2-fr}
  \dim(A_{\mathfrak{q}})
  +\dim(A_{\mathfrak{p}}/\mathfrak{q}A_{\mathfrak{p}})
  =\dim(A_{\mathfrak{p}}).
\tag{16.1.4.2}
\]
若 $A$ 是双等维的，则 $A_{\mathfrak{p}}$ 也同样如此，其中
$\mathfrak{p}$ 是 $A$ 的任意素理想。
\end{env}

\begin{proposition}[16.1.5]
\label{0.16.1.5-fr}
设 $\rho:A\to B$ 是环同态，使 $B$ 成为 $A$ 上的整代数。则
$\dim(B)\leq\dim(A)$；若此外 $\rho$ 是单射，则
$\dim(B)=\dim(A)$。
\end{proposition}

\begin{proof}
\oldpage[0\textsubscript{IV-1}]{24}
若 $\mathfrak{n}$ 是 $\rho$ 的核，则 $\rho(A)$ 同构于
$A/\mathfrak{n}$，故由 (16.1.2.1) 有
$\dim\rho(A)\leq\dim(A)$；又因 $B$ 整于 $\rho(A)$，所以只须在
$A\subset B$ 时证明第二个断言。若
$\mathfrak{p}'_0\subset\mathfrak{p}'_1\subset\cdots\subset\mathfrak{p}'_n$
是 $B$ 的一个素理想链，则 $\mathfrak{p}'_i\cap A$ 是 $A$ 中彼此不同的
素理想（Bourbaki，\emph{Alg.~comm.}，第 V 章，\S~2，第 1 号，命题 1
的推论 1），因而组成一个素理想链。反之，任取链
$\mathfrak{p}_0\subset\mathfrak{p}_1\subset\cdots\subset\mathfrak{p}_m$，
其各项为 $A$ 的素理想；已知对每个 $i$，存在素理想
$\mathfrak{p}'_i$（在 $B$ 中），使得 $\mathfrak{p}'_i\cap A=\mathfrak{p}_i$，并且
$\mathfrak{p}'_i\subset\mathfrak{p}'_{i+1}$ 对
$0\leq i\leq m-1$ 成立（\emph{同上}，定理 1 的推论 2）。结论随即得出。
\end{proof}

\begin{proposition}[16.1.6]
\label{0.16.1.6-fr}
设 $B$ 是整环，$A$ 是 $B$ 的局部子环，$\mathfrak{m}$ 是它的极大理想。
假设 $A$ 整闭且 $B$ 整于 $A$；则对任意极大理想 $\mathfrak{n}$
（在 $B$ 中），有 $\dim(B_{\mathfrak{n}})=\dim(A)$。
\end{proposition}

\begin{proof}
(16.1.5) 论证的第一部分表明
$\dim(B_{\mathfrak{n}})\leq\dim(A)$。反之，考虑素理想链
$\mathfrak{p}_0\subset\mathfrak{p}_1\subset\cdots\subset
\mathfrak{p}_k=\mathfrak{m}$，这是 $A$ 中的一条链；已知
$\mathfrak{n}\cap A=\mathfrak{m}$（Bourbaki，\emph{Alg.~comm.}，第 V 章，
\S~2，第 1 号，命题 1），并且由所作假设，可以应用 Cohen--Seidenberg
第二定理（Bourbaki，\emph{Alg.~comm.}，第 V 章，\S~2，第 4 号，定理 3），
它保证存在素理想链
$\mathfrak{p}'_0\subset\mathfrak{p}'_1\subset\cdots\subset
\mathfrak{p}'_k=\mathfrak{n}$，这是 $B$ 中的一条链，并且
$\mathfrak{p}'_i\cap A=\mathfrak{p}_i$ 对每个 $i$ 成立；从而立即得到
$\dim(B_{\mathfrak{n}})\geq\dim(A)$，命题得证。
\end{proof}

\begin{env}[16.1.7]
\label{0.16.1.7-fr}
设 $M$ 是 $A$-模；把 $M$ 的\emph{维数}记为
$\dim_A(M)$ 或 $\dim(M)$，并定义为环 $A/\operatorname{Ann}(M)$ 的维数
（该环同构于位似环 $A_M$，即 $M$ 的位似环）。因而有
$\dim(M^k)=\dim(M)$，这对每个 $k>0$ 成立。若 $M$ 是有限型的，则 $A$ 中包含
$\operatorname{Ann}(M)$ 的素理想，恰好是属于 $\operatorname{Supp}(M)$
的那些素理想（$0_{\mathrm I}$，1.7.4）；而后者在
$\operatorname{Spec}(A)$ 中是闭的，所以此时
\[
\label{0.16.1.7.1-fr}
  \dim(M)=\dim(\operatorname{Supp}(M))\leq\dim(A).
\tag{16.1.7.1}
\]
若 $M$、$N$ 是两个 $A$-模且 $N\subset M$，则
\[
\label{0.16.1.7.2-fr}
  \dim(N)\leq\dim(M),\qquad \dim(M/N)\leq\dim(M).
\tag{16.1.7.2}
\]
若 $S$ 是 $A$ 的乘法子集且 $M$ 是有限型的，则
\[
\label{0.16.1.7.3-fr}
  \dim_{S^{-1}A}(S^{-1}M)\leq\dim_A(M).
\tag{16.1.7.3}
\]
事实上，有
$\operatorname{Ann}(S^{-1}M)=S^{-1}\operatorname{Ann}(M)$（Bourbaki，
\emph{Alg.~comm.}，第 II 章，\S~2，第 4 号），并且
$S^{-1}A/S^{-1}\operatorname{Ann}(M)=S^{-1}(A/\operatorname{Ann}(M))$。

另一方面，包含 $\operatorname{Ann}(M)$ 的任一素理想极大链都终止于某个
极大理想 $\mathfrak{m}$，所以它的长度等于 $A_{\mathfrak{m}}$ 中相应素理想链
的长度（这些素理想包含
$\operatorname{Ann}(M_{\mathfrak{m}})=(\operatorname{Ann}(M))_{\mathfrak{m}}$）；
由这个观察和 (16.1.7.3) 得到
\[
\label{0.16.1.7.4-fr}
  \dim_A(M)=\sup_{\mathfrak{m}}
  (\dim_{A_{\mathfrak{m}}}(M_{\mathfrak{m}}))
\tag{16.1.7.4}
\]
其中 $\mathfrak{m}$ 遍历 $A$ 的全部极大理想。

称 $M$ 是悬链的（分别地，等维的、等余维的、双等维的），如果
$A/\operatorname{Ann}(M)$ 具有相应性质。
\end{env}

\begin{proposition}[16.1.8]
\label{0.16.1.8-fr}
\oldpage[0\textsubscript{IV-1}]{25}
设 $A$ 是环，$M$ 是有限型 $A$-模，$\mathfrak{p}$ 是 $A$ 的素理想。则
\[
\label{0.16.1.8.1-fr}
  \dim_{A_{\mathfrak{p}}}(M_{\mathfrak{p}})
  +\dim_A(M/\mathfrak{p}M)\leq\dim_A(M).
\tag{16.1.8.1}
\]
\end{proposition}

\begin{proof}
有
$\operatorname{Ann}(M_{\mathfrak{p}})=(\operatorname{Ann}(M))_{\mathfrak{p}}$，
故 $A_{\mathfrak{p}}$ 中包含 $\operatorname{Ann}(M_{\mathfrak{p}})$ 的素理想，
与 $A$ 中包含 $\operatorname{Ann}(M)$ 且包含于 $\mathfrak{p}$ 的素理想
一一对应。另一方面，若 $\operatorname{Ann}(M)=\mathfrak{n}$，则
$V(\operatorname{Ann}(M/\mathfrak{p}M))=V(\mathfrak{p}+\mathfrak{n})$
（$0_{\mathrm I}$，1.7.5），所以 $A$ 中包含
$\operatorname{Ann}(M/\mathfrak{p}M)$ 的素理想都包含
$\mathfrak{p}+\mathfrak{n}$；结论由此得出。

最后这个观察还表明
\[
\label{0.16.1.8.2-fr}
  \dim_{A/\mathfrak{p}}(M/\mathfrak{p}M)=\dim_A(M/\mathfrak{p}M).
\tag{16.1.8.2}
\]
\end{proof}

\begin{proposition}[16.1.9]
\label{0.16.1.9-fr}
设 $A$ 是 Noether 环，$\rho:A\to B$ 是环同态，$M$ 是 $B$-模，并且
$M_{[\rho]}$ 是有限型 $A$-模。则
\[
\label{0.16.1.9.1-fr}
  \dim_A(M_{[\rho]})=\dim_B(M).
\tag{16.1.9.1}
\]
\end{proposition}

\begin{proof}
事实上，$B_M$（$B$-模 $M$ 的位似环）可视为
$C=\operatorname{End}_A(M_{[\rho]})$ 的子环，而 $C$ 是有限型 $A$-模
（Bourbaki，\emph{Alg.~comm.}，第 III 章，\S~3，第 1 号，引理 2），
所以 $B_M$ 也是有限型的；此外，$A_M$（$A$-模 $M_{[\rho]}$ 的位似环）
是 $B_M$ 的子环，并且是 $A$ 在 $C$ 中的典范像，故 $B_M$ 在 $A_M$
上有限；结论遂由 (16.1.5) 得出。
\end{proof}

\begin{env}[16.1.10]
\label{0.16.1.10-fr}
设 $A$ 是 Noether 环；若 $M$ 是有限长度的 $A$-模，则已知
（Bourbaki，\emph{Alg.~comm.}，第 IV 章，\S~2，第 5 号，命题 7 及其
推论 1），$\operatorname{Supp}(M)$ 是有限离散空间，所以
$\dim(M)=0$（只要 $M\neq0$）；反之，若 $M$ 是有限型 $A$-模且
$\dim(M)=0$，则
$\operatorname{Supp}(M)$ 的每一点都是闭点，也就是说，它是 $A$ 的极大理想，
因而（\emph{同上}，命题 7）$M$ 具有有限长度。
\end{env}

\subsection[Noether 半局部环的维数]{Noether 半局部环的维数\protect\footnote{本号重录
Serre 于 1958 年在法兰西公学院课程中的讲演。}}
\label{subsection:0.16.2-fr}

\begin{env}[16.2.1]
\label{0.16.2.1-fr}
设 $A$ 是 Noether 半局部环，$\mathfrak{r}$ 是它的 Jacobson 根。回忆，
若理想 $\mathfrak{q}$ 是 $A$ 的定义理想，则它满足
$\mathfrak{q}\subset\mathfrak{r}$，并且 $\mathfrak{q}$ 包含
$\mathfrak{r}$ 的某个幂；
等价地说，包含 $\mathfrak{q}$ 的素理想全是极大理想，或者
$A/\mathfrak{q}$ 是 Artin 环。设 $M$ 是有限型 $A$-模；对每个整数
$n>0$，$M/\mathfrak{q}^nM$ 是有限长度 $A$-模，其长度等于
$\sum_{i=0}^{n-1}\operatorname{long}(\gr_i(M))$；其中
$\gr_\bullet(M)$ 是 $(A/\mathfrak{q})$-分次模，它是由 $M$ 上的
$\mathfrak{q}$-预进制滤过所得的相伴分次模。已知存在有理系数多项式
$Q(n)$，使等式 $\operatorname{long}(\gr_n(M))=Q(n)$ 对充分大的 $n$
成立（III，2.5.4）；由此立即推出存在有理系数多项式
$P_{\mathfrak{q}}(M,n)$（以 $n$ 为变量），使
\[
  \operatorname{long}(M/\mathfrak{q}^nM)=P_{\mathfrak{q}}(M,n)
\]
对充分大的 $n$ 成立；显然这个多项式是唯一的，并且其首项系数
$>0$（当 $M\neq0$ 时）；称它为 $M$ 对于 $\mathfrak{q}$-预进制滤过的
Hilbert--Samuel 多项式。
\end{env}

\begin{lemma}[16.2.2.1]
\label{0.16.2.2.1-fr}
\oldpage[0\textsubscript{IV-1}]{26}
设 $(M_n)$ 是一个 $\mathfrak{q}$-良滤过（在 $M$ 上；Bourbaki，
\emph{Alg.~comm.}，第 III 章，\S~3，第 1 号）。则
$\operatorname{long}(M/M_n)$ 对充分大的 $n$ 等于一个关于 $n$ 的多项式，
其次数和首项系数分别与 $P_{\mathfrak{q}}(M,n)$ 的次数和首项系数相同。
\end{lemma}

\begin{proof}
事实上，存在 $n_0$，使得 $M_{n+1}=\mathfrak{q}M_n$ 对
$n\geq n_0$ 成立；由此有包含关系
$\mathfrak{q}^{n+n_0}M\subset M_{n+n_0}=\mathfrak{q}^nM_{n_0}
\subset\mathfrak{q}^nM\subset M_n$，这对充分大的 $n$ 成立，从而
\[
  P_{\mathfrak{q}}(M,n+n_0)\geq
  \operatorname{long}(M/M_{n+n_0})\geq
  P_{\mathfrak{q}}(M,n)\geq\operatorname{long}(M/M_n)
\]
并由此得到结论。
\end{proof}

若 $\mathfrak{q}'$ 是第二个定义理想，则有
$\mathfrak{q}'\supset\mathfrak{q}^m$，其中 $m$ 是某个整数，因而
$P_{\mathfrak{q}'}(M,n)\leq P_{\mathfrak{q}}(M,mn)$；所以
$P_{\mathfrak{q}}(M,n)$ 的次数与所取的定义理想 $\mathfrak{q}$ 无关；
记这个次数为 $d(M)$。

\begin{lemma}[16.2.2.2]
\label{0.16.2.2.2-fr}
设 $0\to M'\to M\to M''\to0$ 是有限型 $A$-模的正合列。则多项式
\[
  P_{\mathfrak{q}}(M,n)-P_{\mathfrak{q}}(M',n)
  -P_{\mathfrak{q}}(M'',n)
\]
的次数满足 $\leq d(M')-1\leq d(M)-1$。
\end{lemma}

\begin{proof}
事实上，由
$M'_n=M'\cap\mathfrak{q}^nM'$ 给出的滤过是 $\mathfrak{q}$-良滤过
（Bourbaki，\emph{Alg.~comm.}，第 III 章，\S~3，第 1 号，命题 1）；由于
\[
  \operatorname{long}(M/\mathfrak{q}^nM)
  =\operatorname{long}(M''/\mathfrak{q}^nM'')
  +\operatorname{long}(M'/M'_n),
\]
把 (16.2.2.1) 应用于 $M'$，便立即得到本引理。
\end{proof}

\begin{theorem}[16.2.3]
\label{0.16.2.3-fr}
\textup{(Krull--Chevalley--Samuel).}
设 $A$ 是 Noether 半局部环，$\mathfrak{r}$ 是它的 Jacobson 根，
$M\neq0$ 是有限型 $A$-模，$d(M)$ 是 $M$ 的 Hilbert--Samuel 多项式
（对于 $A$ 的任意定义理想）的次数；另一方面，令 $s(M)$ 为满足下述条件的
整数 $n$ 的下确界：存在元素 $x_1,\ldots,x_n$（它们属于
$\mathfrak{r}$），使 $M/(x_1M+\cdots+x_nM)$ 具有有限长度。则
\[
  d(M)=s(M)=\dim(M).
\]
\end{theorem}

\begin{lemma}[16.2.3.1]
\label{0.16.2.3.1-fr}
对每个 $x\in\mathfrak{r}$，令 ${}_xM$ 表示 $M$ 中由被 $x$ 零化的元素
组成的子模。则：
\begin{enumerate}
  \item[\rm{(i)}] 有 $s(M)\leq s(M/xM)+1$。
  \item[\rm{(ii)}] 设 $\mathfrak{p}_i$ 是
  $\operatorname{Supp}(M)$ 中满足
  \[
    \dim(A/\mathfrak{p}_i)=\dim(M)\quad(1\leq i\leq m).
  \]
  的素理想。
\end{enumerate}
若 $x\notin\mathfrak{p}_i$ 对每个 $i$ 都成立，则
$\dim(M/xM)\leq\dim(M)-1$。
\begin{enumerate}
  \item[\rm{(iii)}] 若 $\mathfrak{q}$ 是 $A$ 的定义理想，则多项式
  \[
    P_{\mathfrak{q}}({}_xM,n)-P_{\mathfrak{q}}(M/xM,n)
  \]
  的次数 $\leq d(M)-1$。
\end{enumerate}
\end{lemma}

\begin{proof}
断言 (i) 是平凡的：若 $(x_i)_{1\leq i\leq m}$ 使模
$N=M/xM$ 的商 $N/(x_1N+\cdots+x_mN)$ 具有有限长度，则只须注意，该商
同构于 $M/(xM+x_1M+\cdots+x_mM)$。

为证明 (ii)，注意，若 $\mathfrak{a}$ 是 $M$ 的零化子，则包含
$M/xM$ 的零化子的素理想，恰好是包含 $\mathfrak{a}+Ax$
的素理想（$0_{\mathrm I}$，1.7.5）；因此没有一个 $\mathfrak{p}_i$
能够包含 $\mathfrak{a}+Ax$，而依 $\dim(M)$ 的定义（16.1.7），
任何包含 $\mathfrak{a}+Ax$ 的素理想链，其长度都
$\leq\dim(M)-1$。

\oldpage[0\textsubscript{IV-1}]{27}

最后，为证明 (iii)，注意有两个正合列
\[
  0\to{}_xM\to M\xrightarrow{x}xM\to0,\qquad
  0\to xM\to M\to M/xM\to0,
\]
而断言立即由引理 (16.2.2.2) 得出。
\end{proof}

(16.2.3) 的证明分为三步：

\begin{env}[16.2.3.2]
\label{0.16.2.3.2-fr}
有 $\dim(M)\leq d(M)$。
\end{env}

对 $d(M)$ 作归纳；$d(M)=0$ 的情形是平凡的。于是设
$d(M)\geq1$，并取 $\mathfrak{p}_0\in\operatorname{Supp}(M)$，使得
$\dim(A/\mathfrak{p}_0)=\dim(M)$。这说明 $\mathfrak{p}_0$ 是
$\operatorname{Supp}(M)$ 的一个\emph{极小}元，故它是 $M$ 的相伴素理想
（Bourbaki，\emph{Alg.~comm.}，第 IV 章，\S~1，第 4 号，定理 2），因而
$M$ 含有一个子模 $N$，它同构于 $A/\mathfrak{p}_0$
（\emph{同上}，\S~1，第 1 号）；由 (16.2.2.2) 有
$d(M)\geq d(N)$，所以问题归结为证明 $\dim(N)\leq d(N)$。于是取
$\mathfrak{p}_0\subset\mathfrak{p}_1\subset\cdots\subset\mathfrak{p}_n$
为 $A$ 中彼此不同的素理想链，并证明 $n\leq d(N)$。若 $n=0$，这显然成立。
否则，存在 $x\in\mathfrak{p}_1\cap\mathfrak{r}$，使得
$x\notin\mathfrak{p}_0$；事实上，关系
$\mathfrak{p}_0\supset\mathfrak{p}_1\cap\mathfrak{r}$ 会推出
$\mathfrak{p}_0\supset\mathfrak{r}$，因为 $\mathfrak{p}_0$ 不包含
$\mathfrak{p}_1$；于是 $\mathfrak{p}_0$ 将是极大理想，这就产生矛盾。
有 $\mathfrak{p}_i\in\operatorname{Supp}(N/xN)$ 对 $i\geq1$ 成立，故
$n-1\leq\dim(N/xN)$（16.1.7）。由 ${}_xN=0$（这是所选 $x$ 的缘故），
再由 (16.2.3.1, (iii)) 得到 $d(N/xN)\leq d(N)-1$。归纳假设因而给出
$d(N/xN)=\dim(N/xN)\geq n-1$，所以 $n-1\leq d(N)-1$。

\begin{env}[16.2.3.3]
\label{0.16.2.3.3-fr}
有 $d(M)\leq s(M)$。
\end{env}

事实上，取一族元素 $(x_i)_{1\leq i\leq n}$，它们属于
$\mathfrak{r}$，并使得置
$\mathfrak{a}=x_1A+\cdots+x_nA$ 时，$M/\mathfrak{a}M$ 具有有限长度。
于是 $M/\mathfrak{a}M$ 的相伴素理想都是极大理想（Bourbaki，
\emph{Alg.~comm.}，第 IV 章，\S~2，第 5 号，命题 7）；又因为它们恰是
包含 $\mathfrak{q}=\mathfrak{a}+(\mathfrak{r}\cap\operatorname{Ann}(M))$
的全部素理想，所以 $A/\mathfrak{q}$ 是 Artin 环（\emph{同上}，命题 9），
而 $\mathfrak{q}$ 是 $A$ 的定义理想。但有
$\mathfrak{q}^mM/\mathfrak{q}^{m+1}M=
\mathfrak{a}^mM/\mathfrak{a}^{m+1}M$；若 $M$ 由 $r$ 个元素 $z_j$ 生成，
则 $\mathfrak{a}^mM/\mathfrak{a}^{m+1}M$ 是一个
$(A/\mathfrak{q})$-模，它由诸
$x_1^{\nu_1}x_2^{\nu_2}\cdots x_n^{\nu_n}z_j$ 的典范像生成，其中
$\sum_i\nu_i=m$；因而它由 $r\binom{m+n-1}{n-1}$ 个元素生成，其长度
$\leq ar\binom{m+n-1}{n-1}$，其中 $a$ 是常数。由此立即推出
$P_{\mathfrak{q}}(M,n)$ 的次数 $\leq n$，故 $d(M)\leq s(M)$。

\begin{env}[16.2.3.4]
\label{0.16.2.3.4-fr}
有 $s(M)\leq\dim(M)$。
\end{env}

首先注意，由 (16.2.3.2)，$n=\dim(M)$ 是有限的。对 $n$ 作归纳；
当 $n=0$ 时，命题是显然的，因为这时 $M$ 具有有限长度（16.1.10）。
设 $n\geq1$，并令 $\mathfrak{p}_i$ $(1\leq i\leq m)$ 为
$\operatorname{Supp}(M)$ 中满足 $\dim(A/\mathfrak{p}_i)=n$ 的素理想；
$\mathfrak{p}_i$ 的定义表明，它们是 $\operatorname{Supp}(M)$ 的极小元，
所以其数目有限（Bourbaki，\emph{Alg.~comm.}，第 IV 章，\S~1，第 4 号，
定理 2）。由于 $n\geq1$，诸 $\mathfrak{p}_i$ 都不是极大理想，故存在
$x\in\mathfrak{r}$，使 $x\notin\mathfrak{p}_i$ 对每个 $i$ 都成立
（Bourbaki，\emph{Alg.~comm.}，第 II 章，\S~1，第 1 号，命题 2）。
根据 (16.2.3.1)，有 $s(M)\leq s(M/xM)+1$ 和
$\dim(M)\geq\dim(M/xM)+1$；归纳假设给出
$s(M/xM)\leq\dim(M/xM)$，结论随即得出。

\begin{corollary}[16.2.4]
\label{0.16.2.4-fr}
在 (16.2.3) 的假设下，有
\begin{equation}
\label{0.16.2.4.1-fr}
\tag{16.2.4.1}
  \dim_{\widehat{A}}(\widehat{M})=\dim_A(M).
\end{equation}
\end{corollary}

\oldpage[0\textsubscript{IV-1}]{28}

\begin{proof}
事实上，$M/\mathfrak{r}^nM$ 与
$\widehat{M}/\mathfrak{r}^n\widehat{M}$ 同构，故
$d(\widehat{M})=d(M)$。
\end{proof}

\begin{corollary}[16.2.5]
\label{0.16.2.5-fr}
Noether 半局部环 $A$ 的维数有限，并且等于 $A$ 的 Jacobson 根中生成一个
定义理想所需元素的最少数目。
\end{corollary}

\begin{proof}
这就是等式 $s(M)=\dim(M)$ 在 $M=A$ 时的情形。
\end{proof}

特别地：

\begin{corollary}[16.2.6]
\label{0.16.2.6-fr}
设 $A$ 是 Noether 局部环，$\mathfrak{m}$ 是它的极大理想，
$k=A/\mathfrak{m}$ 是它的剩余域；则
\begin{equation}
\label{0.16.2.6.1-fr}
\tag{16.2.6.1}
  \dim(A)\leq\operatorname{rg}_k(\mathfrak{m}/\mathfrak{m}^2).
\end{equation}
\end{corollary}

\begin{proof}
事实上，已知若 $x_i$ $(1\leq i\leq n)$ 是 $\mathfrak{m}$ 的元素，
其模 $\mathfrak{m}^2$ 的类构成 $k$-向量空间
$\mathfrak{m}/\mathfrak{m}^2$ 的一组基，则诸 $x_i$ 生成
$\mathfrak{m}$（Bourbaki，\emph{Alg.~comm.}，第 II 章，\S~3，第 2 号，
命题 4 的推论 2）。
\end{proof}

\begin{proposition}[16.2.7]
\label{0.16.2.7-fr}
设 $k$ 是域，$B$ 是一个按正次数分次的有限型 $k$-代数，它由次数为 1 的齐次元生成，
并且 $B_0=k$；设 $E$ 是有限型分次 $B$-模，因而（III，2.5.4），对充分大的
$n$，$\operatorname{rg}_k(E_n)=r(n)$ 具有
$\Phi(n)-\Phi(n-1)$ 的形式，其中 $\Phi$ 是关于 $n$ 的次数为 $d$ 的多项式。
则 $B$ 中存在一族 $(t_i)_{1\leq i\leq d}$，由 $d$ 个次数
$\geq1$ 的齐次元组成，使得 $E/(\sum_i t_iE)$ 作为 $k$-向量空间是有限维的。
\end{proposition}

\begin{proof}
令 $\mathfrak{q}$ 为 $\bigoplus_{n\geq1}B_n$；这是 $B$ 的极大理想。对每个
$n$，令 $E'_n$ 为 $B$-子模 $\bigoplus_{h\geq n+1}E_h$，它是 $E$ 的子模。
$B-\mathfrak{q}$ 的每个元素都在 Artin $B$-模 $E/E'_n$ 上确定一个单射位似，
故这个位似是双射（Bourbaki，\emph{Alg.}，第 VIII 章，\S~2，第 2 号，引理 3）；
这说明 $(E/E'_n)_{\mathfrak{q}}$ 可与 $E/E'_n$ 等同。由于
$B_{\mathfrak{q}}$ 的剩余域是 $k$，由假设可知
$\operatorname{long}_{B_{\mathfrak{q}}}((E/E'_n)_{\mathfrak{q}})$ 是关于
$n$ 的次数为 $d$ 的多项式。另一方面，由于 $E$ 是有限型 $B$-模，而且
$B$ 由其次数为 1 的齐次元生成，所以滤过 $(E'_n)$ 在 $E$ 上是
$\mathfrak{q}$-良滤过（Bourbaki，\emph{Alg.~comm.}，第 III 章，\S~1，
第 3 号，引理 1）；因而滤过 $((E'_n)_{\mathfrak{q}})$ 在
$E_{\mathfrak{q}}$ 上是 $\mathfrak{q}B_{\mathfrak{q}}$-良滤过。
由于 $B_{\mathfrak{q}}$ 是 Noether 局部环，由 (16.2.3) 和 (16.2.2.1)
得到 $\dim(E_{\mathfrak{q}})=d$。现在对 $d$ 作归纳；当 $d=0$ 时断言显然。
设 $d\geq1$，并注意，极小素理想 $\mathfrak{p}_i$ $(1\leq i\leq r)$
（即 $\operatorname{Ass}(E)$ 的极小元）都是\emph{分次的}（Bourbaki，
\emph{Alg.~comm.}，第 IV 章，\S~3，第 1 号，命题 1），而
$\mathfrak{q}$ 不等于诸 $\mathfrak{p}_i$ 的并集，因为 $d\geq1$；所以存在一个
\emph{齐次}元 $t_1\in\mathfrak{q}$，它不属于任何 $\mathfrak{p}_i$
（Bourbaki，\emph{Alg.~comm.}，第 III 章，\S~1，
第 4 号，命题 8）。因为素理想 $(\mathfrak{p}_i)_{\mathfrak{q}}$ 是
$\operatorname{Ass}(E_{\mathfrak{q}})$ 的极小元（Bourbaki，
\emph{Alg.~comm.}，第 IV 章，\S~1，第 2 号，命题 5），所以
$\dim(E_{\mathfrak{q}}/t_1E_{\mathfrak{q}})=d-1$（16.3.4）。此时只须把归纳
假设应用于分次 $B$-模 $E/t_1E$，便得到结论。
\end{proof}

\subsection{Noether 局部环中的参数系}
\label{subsection:0.16.3-fr}

\begin{proposition}[16.3.1]
\label{0.16.3.1-fr}
设 $A$ 是 Noether 环，$\mathfrak{p}$ 是 $A$ 的素理想，$n$ 是整数。下列条件等价：
\begin{enumerate}
  \item[$a)$] $\operatorname{ht}(\mathfrak{p})=
  \dim(A_{\mathfrak{p}})\leq n$；
  \item[$b)$] 存在 $n$ 个元素 $x_i$，它们属于 $A$ $(1\leq i\leq n)$，使得
  $\mathfrak{p}$ 在包含理想 $\mathfrak{a}$ 的素理想集合中是极小的，其中该理想
  由诸 $x_i$ 生成
  （换言之，$V(\mathfrak{p})$ 是 $V(\mathfrak{a})$ 的一个不可约分支）。
\end{enumerate}
\end{proposition}

\begin{proof}
条件 $b)$ 蕴含 $\mathfrak{a}A_{\mathfrak{p}}$ 是 $A_{\mathfrak{p}}$ 的一个定义理想，
故由 (16.2.5) 得到 $a)$。反过来，若 $a)$ 成立，则存在一个定义理想
$\mathfrak{b}$（属于 $A_{\mathfrak{p}}$），它由 $n$ 个元素 $x_i/s$ 生成，其中
$s\in A-\mathfrak{p}$。
根据 $A_{\mathfrak{p}}$ 的素理想与 $A$ 中包含于 $\mathfrak{p}$ 的素理想之间的一一对应，
理想 $\mathfrak{a}$ 是 $A$ 的理想，由诸 $x_i$ 生成，并且满足 $b)$。
\end{proof}

\oldpage[0\textsubscript{IV-1}]{29}

当 $n=1$ 时，得到如下特殊情形：

\begin{corollary}[16.3.2]
\label{0.16.3.2-fr}
\textup{（“Krull 主理想定理”（Hauptidealsatz））。}
一个 Noether 环的素理想的高度为 $\leq1$，当且仅当它在包含某个适当主理想的
素理想集合中是极小的。
\end{corollary}

\begin{corollary}[16.3.3]
\label{0.16.3.3-fr}
\textup{（Artin--Tate）。}
设 $A$ 是 Noether 整环。下列条件等价：
\begin{enumerate}
  \item[$a)$] $A$ 是维数 $\leq1$ 的半局部环。
  \item[$b)$] 存在 $f\neq0$，它是 $A$ 中的元素，并使 $A_f$ 是域。
\end{enumerate}
\end{corollary}

\begin{proof}
条件 $a)$ 等价于说，$A$ 只有有限多个素理想 $\mathfrak{m}_i\neq0$，而且这些理想
全都是极大的（$0$ 是素理想）。由于诸 $\mathfrak{m}_i$ 的乘积非零（因为 $A$ 是整环），
故该乘积的一个元素 $f\neq0$ 属于每个 $\mathfrak{m}_i$；因而，$(0)$ 是整环 $A_f$ 的
唯一素理想，换言之，$A_f$ 是域。（请注意，这一部分证明没有用到 $A$ 是
Noether 环这一事实。）现在证明 $b)$ 蕴含 $a)$；在条件 $b)$ 下，每个素理想
$\mathfrak{p}\neq0$（属于 $A$）都包含 $f$。设 $\mathfrak{p}_i$ $(1\leq i\leq n)$ 是包含 $f$ 的
素理想中的极小元（它们只有有限多个，因为 $A/fA$ 是 Noether 环）；只须证明诸
$\mathfrak{p}_i$ 是\emph{极大}理想，因为这样每个素理想 $\neq0$ 必包含某个
$\mathfrak{p}_i$，因而就等于它。现在假设存在一个极大理想 $\mathfrak{m}$，它包含
某个 $\mathfrak{p}_i$ 而不同于该 $\mathfrak{p}_i$；于是 $\mathfrak{m}$ 不包含于诸
$\mathfrak{p}_i$ 的并集（Bourbaki，\emph{Alg.~comm.}，第 II 章，\S~1，第 1 号，
命题 2），换言之，存在 $g\in\mathfrak{m}$，它不属于任何 $\mathfrak{p}_i$。
令 $\mathfrak{q}$ 为包含 $g$ 的素理想中的一个极小元；由“Krull 主理想定理”
(16.3.2) 可知，$\mathfrak{q}$ 的高度为 1，因而 $\neq0$；所以按假设它包含 $f$，
进而包含某个 $\mathfrak{p}_i$。又因为 $g\in\mathfrak{q}$ 而 $g$ 不属于任何
$\mathfrak{p}_i$，所以 $\mathfrak{q}$ 不同于诸 $\mathfrak{p}_i$；于是它的高度将
为 $\geq2$，这就矛盾。
\end{proof}

\begin{proposition}[16.3.4]
\label{0.16.3.4-fr}
设 $A$ 是 Noether 半局部环，$\mathfrak{r}$ 是它的 Jacobson 根，$M$ 是有限型
$A$-模，而 $\mathfrak{p}_i$ 是 $\operatorname{Supp}(M)$ 中满足
$\dim(A/\mathfrak{p}_i)=\dim(M)$ 的元素。则诸 $\mathfrak{p}_i$ 是
$\operatorname{Ass}(M)$ 的极小元，并且
$\dim(M/\mathfrak{p}_iM)=\dim(M)$。对每个 $x\in\mathfrak{r}$，都有
\begin{equation}
\label{0.16.3.4.1-fr}
\tag{16.3.4.1}
  \dim(M/xM)\geq\dim(M)-1.
\end{equation}
此外，(16.3.4.1) 两边相等，当且仅当 $x$ 不属于任何 $\mathfrak{p}_i$。
\end{proposition}

\begin{proof}
第一个断言是立即的，因为诸 $\mathfrak{p}_i$ 按定义是
$\operatorname{Supp}(M)$ 的极小元（16.1.6），而后者也都是
$\operatorname{Ass}(M)$ 的极小元（Bourbaki，\emph{Alg.~comm.}，第 IV 章，
\S~1，第 4 号，定理 2）。此外，$\operatorname{Supp}(M/\mathfrak{p}_iM)$
是包含 $\mathfrak{p}_i$ 的素理想集合（$0_{\mathrm I}$, 1.7.5），故
$\dim(M/\mathfrak{p}_iM)=\dim(A/\mathfrak{p}_i)=\dim(M)$。

令 $N=M/xM$。若 $x_1,\ldots,x_n$ 是 $\mathfrak{r}$ 的元素，并且
$N/(x_1N+\cdots+x_nN)$ 长度有限，这就意味着
$M/(xM+x_1M+\cdots+x_nM)$ 长度有限，故由 (16.2.3) 得到不等式 (16.3.4.1)。

若 $x$ 不属于任何 $\mathfrak{p}_i$，则 (16.3.4.1) 两边相等；这是 (16.3.4.1)
与 (16.2.3.1) 的结果。反过来，若
\oldpage[0\textsubscript{IV-1}]{30}
$\dim(M/xM)=\dim(M)-1$，则没有一个 $\mathfrak{p}_i$ 属于
$\operatorname{Supp}(M/xM)$；而这个支撑中的素理想正是包含
$\operatorname{Ann}(M)+Ax$ 的素理想。因为诸 $\mathfrak{p}_i$ 包含
$\operatorname{Ann}(M)$，所以它们不能包含 $x$。
\end{proof}

\begin{corollary}[16.3.5]
\label{0.16.3.5-fr}
沿用 (16.3.4) 的记号，对每个理想 $\mathfrak{a}\subset\mathfrak{p}_i$，有
$\dim(M/\mathfrak{a}M)=\dim(M)$；若令 $N=M/\mathfrak{a}M$，则关系
$\dim(M/xM)=\dim(M)-1$ 蕴含
$\dim(N/xN)=\dim(N)-1$。
\end{corollary}

\begin{proof}
事实上，$\operatorname{Supp}(N)$ 是包含 $\mathfrak{a}+\operatorname{Ann}(M)$
的素理想集合，因而 $\mathfrak{p}_i$ 属于 $\operatorname{Supp}(N)$；又因为
$\dim(N)\leq\dim(M)=\dim(A/\mathfrak{p}_i)$，所以
$\dim(A/\mathfrak{p}_i)=\dim(N)$。此外，满足
$\mathfrak{q}\in\operatorname{Supp}(N)$ 和
$\dim(A/\mathfrak{q})=\dim(N)$ 的素理想是诸 $\mathfrak{p}_j$ 中的某些；
若 $x$ 不属于任何 $\mathfrak{p}_j$，则由 (16.3.4) 有
$\dim(N/xN)=\dim(N)-1$。
\end{proof}

\begin{definition}[16.3.6]
\label{0.16.3.6-fr}
设 $A$ 是 Noether 半局部环，$\mathfrak{r}$ 是它的 Jacobson 根，$M$ 是有限型
$A$-模，并令 $n=\dim(M)$。所谓 $M$ 的一个\emph{参数系}，是指任一系
$(x_i)_{1\leq i\leq n}$，它由 $n$ 个 $\mathfrak{r}$ 中的元素组成，并使得
$M/(x_1M+\cdots+x_nM)$ 具有有限长度。
\end{definition}

我们已在 (16.2.3) 中看到，这样的参数系总是存在。

在 (16.2.3.3) 的证明过程中已经看到，若
$\mathfrak{J}=\sum_{i=1}^n x_iA$ 使 $M/\mathfrak{J}M$ 具有有限长度，则
$\mathfrak{q}=\mathfrak{J}+\operatorname{Ann}(M)$ 是 $A$ 的一个\emph{定义理想}；
反过来，若是这种情形，则 $M/\mathfrak{J}M$ 是 Artin 环 $A/\mathfrak{q}$ 上的
有限型模，故具有有限长度。因此，说 $(x_i)_{1\leq i\leq n}$ 是\emph{$M$ 的}参数系，
或说它是\emph{$A/\operatorname{Ann}(M)$ 的}参数系，是同一回事（也可说，它们在
$A/\operatorname{Ann}(M)$ 中的像构成这个环的一个参数系）。若
$(x_i)_{1\leq i\leq n}$ 是 $M$ 的参数系，则系 $(x_i^k)$ 对每个整数
$k\geq1$ 也同样是；事实上，由上文只须考察 $M=A$ 的情形，而若
$\mathfrak{J}$ 是 $A$ 的定义理想，则理想 $\sum_{i=1}^n x_i^kA$ 包含
$\mathfrak{J}^{kn}$，故它也是定义理想。

\begin{proposition}[16.3.7]
\label{0.16.3.7-fr}
设 $A$ 是 Noether 半局部环，$\mathfrak{r}$ 是它的 Jacobson 根，
$x_1,\ldots,x_k$ 是 $\mathfrak{r}$ 的元素，而 $M$ 是有限型 $A$-模。则
\begin{equation}
\label{0.16.3.7.1-fr}
\tag{16.3.7.1}
  \dim(M/(x_1M+\cdots+x_kM))\geq\dim(M)-k.
\end{equation}
此外，(16.3.7.1) 两边相等，当且仅当存在元素 $x_i$
$(k+1\leq i\leq n=\dim(M))$，它们属于 $\mathfrak{r}$，并使
$(x_i)_{1\leq i\leq n}$ 是 $M$ 的参数系。
\end{proposition}

\begin{proof}
不等式 (16.3.7.1) 由 (16.3.4.1) 对 $k$ 作归纳而得。若 (16.3.7.1) 两边相等，
而 $x_{k+1},\ldots,x_n$ 是 $M/(x_1M+\cdots+x_kM)$ 的参数系，则显然
$M/(x_1M+\cdots+x_kM+x_{k+1}M+\cdots+x_nM)$ 具有有限长度，因而
$(x_i)_{1\leq i\leq n}$ 是 $M$ 的参数系。反过来，若存在具有这一性质的
$x_i$ $(k+1\leq i\leq n)$，并令 $N=M/(x_1M+\cdots+x_kM)$，则显然
$N/(x_{k+1}N+\cdots+x_nN)$ 具有有限长度，故 $\dim(N)\leq n-k$；再由
(16.3.7.1)，两边相等。
\end{proof}

\oldpage[0\textsubscript{IV-1}]{31}

\begin{proposition}[16.3.8]
\label{0.16.3.8-fr}
设 $A$ 是 Noether 半局部环，$X=\operatorname{Spec}(A)$ 是它的谱，
$\mathfrak{a}$ 是 $A$ 的一个不同于 $A$ 的理想，并且
$\operatorname{codim}(V(\mathfrak{a}),X)=r>0$。则存在 $r$ 个元素
$x_1,\ldots,x_r$（属于 $\mathfrak{a}$），它们构成 $A$ 的某个参数系的一部分，且
\[
  \operatorname{codim}(V(\mathfrak{a}),
  V(Ax_1+\cdots+Ax_r))=0.
\]
\end{proposition}

\begin{proof}
设 $\mathfrak{p}_i$ $(1\leq i\leq m)$ 是环 $A$ 的极小素理想，
$\mathfrak{q}_j$ $(1\leq j\leq n)$ 是包含 $\mathfrak{a}$ 的素理想中的极小元；
由 (14.2.1)，有
\[
  \operatorname{codim}(V(\mathfrak{a}),X)
  =\inf_j(\operatorname{codim}(V(\mathfrak{q}_j),X))
  =\inf_j(\dim(A_{\mathfrak{q}_j}))
\]
（16.1.3.2）。条件 $r>0$ 意味着 $\mathfrak{a}$ 不包含于任何
$\mathfrak{p}_i$；因而存在 $x\in\mathfrak{a}$，它不属于任何
$\mathfrak{p}_i$（Bourbaki，\emph{Alg.~comm.}，第 II 章，\S~1，第 1 号，命题 2）。
如上可得
\[
  \operatorname{codim}(V(\mathfrak{a}),V(Ax))
  =\inf_j(\dim((A/Ax)_{\mathfrak{q}_j/Ax}))
  =\inf_j(\dim(A_{\mathfrak{q}_j}/A_{\mathfrak{q}_j}x))~;
\]
但因为 $x/1$ 不属于 $A_{\mathfrak{q}_j}$ 的任何极小素理想，所以
$\dim(A_{\mathfrak{q}_j}/A_{\mathfrak{q}_j}x)
=\dim(A_{\mathfrak{q}_j})-1$（16.3.4），从而
$\operatorname{codim}(V(\mathfrak{a}),V(Ax))=r-1$。出于同一理由，有
$\dim(A/Ax)=\dim(A)-1$。现在对 $r$ 作归纳；若 $r>1$，则存在
$r-1$ 个元素 $x'_i$ $(2\leq i\leq r)$（属于 $A'=A/Ax$），它们构成这个环的
某个参数系的一部分，属于 $\mathfrak{a}/Ax$，并且
\[
  \operatorname{codim}(V(\mathfrak{a}/Ax),
  V(A'x'_2+\cdots+A'x'_r))=0.
\]
令 $(x'_i)_{r+1\leq i\leq s}$ 为 $A'$ 中的一族元素，使得
$(x'_i)_{2\leq i\leq s}$ 是 $A'$ 的一个参数系；对每个 $i$，若
$2\leq i\leq s$，令 $x_i$ 为类 $x'_i$ 在 $A$ 中的一个代表，并要求
$x_i\in\mathfrak{a}$（当 $2\leq i\leq r$ 时）。显然，
$A/(Ax+Ax_2+\cdots+Ax_s)$ 长度有限；又因 $\dim(A)=s$，可见
$x_1=x$ 以及诸 $x_i$（其指标 $i\geq2$）满足命题的条件。
\end{proof}

\begin{proposition}[16.3.9]
\label{0.16.3.9-fr}
设 $A$、$B$ 是两个 Noether 局部环，$\mathfrak{m}$ 是 $A$ 的极大理想，
$k$ 是它的剩余域，而 $\varphi:A\to B$ 是局部同态。
\begin{enumerate}
  \item[\rm{(i)}] 有
  \begin{equation}
  \label{0.16.3.9.1-fr}
  \tag{16.3.9.1}
    \dim(B)\leq\dim(A)+\dim(B\otimes_A k).
  \end{equation}
  \item[\rm{(ii)}] 设 $M\neq0$ 是有限型 $A$-模，$N\neq0$ 是有限型
  $B$-模。则有
  \begin{equation}
  \label{0.16.3.9.2-fr}
  \tag{16.3.9.2}
    \dim_B(M\otimes_A N)
    \leq\dim_A(M)+\dim_{B\otimes_A k}(N\otimes_A k).
  \end{equation}
\end{enumerate}
\end{proposition}

\begin{proof}
设 $m$ 是 $A$ 的维数，$(s_i)_{1\leq i\leq m}$ 是 $A$ 的一个参数系
（16.3.6），从而若 $\mathfrak{a}=\sum_i s_iA$，则 $A/\mathfrak{a}$ 是长度有限的
$A$-模。于是环 $A/\mathfrak{a}$ 是 Artin 环，所以这个环的极大理想
$\mathfrak{m}/\mathfrak{a}$ 是幂零的；理想 $\mathfrak{m}B/\mathfrak{a}B$ 是
$B/\mathfrak{a}B$ 的理想，并且是 $(\mathfrak{m}/\mathfrak{a})\otimes_A B$
的像，故它也幂零；因此
$\dim(B/\mathfrak{a}B)=\dim(B/\mathfrak{m}B)$（这两个环的谱具有相同的底空间）。
另一方面，$B/\mathfrak{a}B=B/(s_1B+\cdots+s_mB)$，而诸 $s_i$ 在 $B$ 中的像
属于 $B$ 的极大理想。因而由 (16.3.7)，有
\[
  \dim(B)\leq m+\dim(B\otimes_A k),
\]
这正是 (16.3.9.1)。

为证明 (16.3.9.2)，注意，若 $\mathfrak{r}$（相应地，$\mathfrak{s}$）是 $M$
（相应地，$N$）
\oldpage[0\textsubscript{IV-1}]{32}
在 $A$（相应地，$B$）中的零化子，则按定义 (16.1.7)，有
$\dim_A(M)=\dim(A/\mathfrak{r})$，$\dim_B(N)=\dim(B/\mathfrak{s})$；另一方面，
$\operatorname{Supp}(M\otimes_A N)$ 是 $\operatorname{Spec}(B)$ 的闭子集，
并且可与下列空间等同（I，9.1.13）：
\[
  \operatorname{Supp}(M)\times_{\operatorname{Spec}(A)}\operatorname{Supp}(N)
  =\operatorname{Spec}((B/\mathfrak{s})\otimes_A(A/\mathfrak{r}))
  =\operatorname{Spec}(B/(\mathfrak{s}+\mathfrak{r}B)).
\]
因为 $N\otimes_A k$ 是有限型 $B$-模，所以
$\dim_{B\otimes_A k}(N\otimes_A k)=\dim_B(N\otimes_A k)$（16.1.9）；而由同上
推理，$\operatorname{Supp}(N\otimes_A k)$ 等于
$\operatorname{Spec}(B/(\mathfrak{s}+\mathfrak{m}B))$。把 (16.3.9.1) 应用于
$A/\mathfrak{r}$ 和 $B/(\mathfrak{s}+\mathfrak{r}B)$，得到
\[
  \dim(B/(\mathfrak{s}+\mathfrak{r}B))
  \leq\dim(A/\mathfrak{r})+
  \dim(B/(\mathfrak{s}+\mathfrak{m}B)),
\]
这正是 (16.3.9.2)。
\end{proof}

\begin{corollary}[16.3.10]
\label{0.16.3.10-fr}
在 (16.3.9) 的假设下，若 $\mathfrak{m}B$ 是 $B$ 的定义理想，则
$\dim(B)\leq\dim(A)$；若进一步 $A$ 是整环且 $\dim(A)=\dim(B)$，则
$\varphi$ 是单射。
\end{corollary}

\begin{proof}
第一个断言由 (16.3.9) 得出，因为此时 $\dim(B/\mathfrak{m}B)=0$。此外，可把
$A$ 替换为 $\varphi(A)$，故 $\dim(B)\leq\dim(\varphi(A))$；若
$\mathfrak{a}=\operatorname{Ker}(\varphi)\neq0$，而 $A$ 是整环，则
$\dim(\varphi(A))=\dim(A/\mathfrak{a})<\dim(A)$（16.1.2.2），因而
$\dim(A)=\dim(B)$ 只能在 $\mathfrak{a}=0$ 时成立。
\end{proof}

\subsection[深度和余深度]{深度和余深度\protect\footnote{在第三章第三部分中，我们将从
上同调观点并在更一般的框架下发展深度概念。}}
\label{subsection:0.16.4-fr}

\begin{proposition}[16.4.1]
\label{0.16.4.1-fr}
设 $A$ 是 Noether 局部环，$\mathfrak{m}$ 是它的极大理想，而 $M$ 是有限型
$A$-模。则每个 $M$-正则序列 $(x_i)_{1\leq i\leq r}$（15.1.7），若由
$\mathfrak{m}$ 中的元素组成，就是 $M$ 的某个参数系的一部分。
\end{proposition}

\begin{proof}
对 $r$ 作归纳；考察 $A$-模
\[
  N=M/(x_1M+\cdots+x_{r-1}M)~;
\]
按归纳假设，$x_1,\ldots,x_{r-1}$ 构成 $M$ 的某个参数系的一部分，所以
$\dim(N)=\dim(M)-(r-1)$（16.3.7）。按假设，以 $x_r$ 为比的位似在 $N$ 中
是单射，故 $x_r$ 不属于 $N$ 的任何相伴素理想，因而由 (16.3.4)，有
$\dim(N/x_rN)=\dim(N)-1$，这也写作
\[
  \dim(M/(x_1M+\cdots+x_rM))=\dim(M)-r~;
\]
所以结论由 (16.3.7) 得出。
\end{proof}

因此，对 $A$-模 $M\neq0$，一个 $M$-正则序列，若其元素属于
$\mathfrak{m}$，则\emph{至多}有 $\dim(M)$ 个元素。

\begin{lemma}[16.4.2]
\label{0.16.4.2-fr}
设 $A$ 是 Noether 局部环，$\mathfrak{m}$ 是它的极大理想，
$k=A/\mathfrak{m}$ 是它的剩余域，而 $M$ 是 $A$-模。一个 $M$-正则序列
$(x_i)_{1\leq i\leq n}$，若由 $\mathfrak{m}$ 中的元素组成，则它是极大的，
当且仅当
\[
  \operatorname{Hom}_A(k,M/(x_1M+\cdots+x_nM))\neq0.
\]
\end{lemma}

\begin{proof}
说 $(x_i)$ 是极大的，意味着对任何 $x\in\mathfrak{m}$，以 $x$ 为比的位似在
$N=M/(x_1M+\cdots+x_nM)$ 中都不是单射；也就是说，
$\mathfrak{m}$ 包含于 $N$ 的相伴素理想之并（Bourbaki，\emph{Alg.~comm.}，
第 IV 章，\S~1，第 1 号，命题 2 的推论 2），因而它等于后者之一，因为
前者是极大理想（Bourbaki，\emph{Alg.~comm.}，第 II 章，\S~1，
第 1 号，命题 2）。
\oldpage[0\textsubscript{IV-1}]{33}
换言之，存在一个元素 $z\in N$，其零化子是 $\mathfrak{m}$；所以子模 $Az$
是 $N$ 的子模，并同构于 $A/\mathfrak{m}=k$，引理得证。
\end{proof}

\begin{lemma}[16.4.3]
\label{0.16.4.3-fr}
设 $A$ 是 Noether 局部环，$\mathfrak{m}$ 是它的极大理想，
$k=A/\mathfrak{m}$ 是它的剩余域，$M$ 是 $A$-模，而
$(x_i)_{1\leq i\leq n}$ 是一个 $M$-正则序列，其元素属于 $\mathfrak{m}$。
则 $A$-模
\[
  \operatorname{Hom}_A(k,M/(x_1M+\cdots+x_nM))
  \quad\text{和}\quad \operatorname{Ext}^n_A(k,M)
\]
同构；当 $n\geq1$ 时，它们还同构于
$\operatorname{Ext}^{n-1}_A(k,M/x_1M)$。
\end{lemma}

\begin{proof}
对 $n$ 作归纳；当 $n=0$ 时，断言显然。令 $N=M/x_1M$；序列
$(x_i)_{2\leq i\leq n}$ 是 $N$-正则的，故
\[
  \operatorname{Hom}_A(k,M/(x_1M+\cdots+x_nM))
  =\operatorname{Hom}_A(k,N/(x_2N+\cdots+x_nN))
\]
由归纳假设同构于 $\operatorname{Ext}^{n-1}_A(k,N)$。现在考察正合序列
$0\to M\xrightarrow{x_1}M\to M/x_1M=N\to0$；由此得到 Ext 正合序列
\begin{equation}
\label{0.16.4.3.1-fr}
\tag{16.4.3.1}
  \operatorname{Ext}^{n-1}_A(k,M)
  \to\operatorname{Ext}^{n-1}_A(k,N)
  \to\operatorname{Ext}^n_A(k,M)
  \xrightarrow{x_1}\operatorname{Ext}^n_A(k,M).
\end{equation}
由归纳假设，$\operatorname{Ext}^{n-1}_A(k,M)$ 同构于
\[
  \operatorname{Hom}_A(k,M/(x_1M+\cdots+x_{n-1}M)),
\]
而这个模为零，因为 $(x_i)_{1\leq i\leq n-1}$ 不是极大的 $M$-正则序列
（16.4.2）。另一方面，因为 $x_1\in\mathfrak{m}$，以 $x_1$ 为比的位似在
$A$-模 $\operatorname{Hom}_A(k,T)$ 中对任意 $A$-模 $T$ 都为零；因而以 $x_1$
为比的位似在每个 $A$-模 $\operatorname{Ext}^n_A(k,T)$ 中也为零。于是引理的
断言立即由正合序列 (16.4.3.1) 得出。
\end{proof}

\begin{corollary}[16.4.4]
\label{0.16.4.4-fr}
设 $A$ 是 Noether 局部环，$\mathfrak{m}$ 是它的极大理想，
$k=A/\mathfrak{m}$ 是它的剩余域，而 $M$ 是有限型 $A$-模，且该模 $\neq0$。
所有极大 $M$-正则序列，若由 $\mathfrak{m}$ 中的元素组成，都具有相同的元素数；
这个数是最小整数 $n$，满足 $\operatorname{Ext}^n_A(k,M)\neq0$。
\end{corollary}

\begin{definition}[16.4.5]
\label{0.16.4.5-fr}
设 $A$ 是 Noether 局部环，$\mathfrak{m}$ 是它的极大理想，而 $M$ 是有限型
$A$-模。把 $M$ 的\emph{深度}记作 $\operatorname{prof}_A M$ 或
$\operatorname{prof}(M)$；它是 $M$-正则序列的元素数的上确界，其中序列由
$\mathfrak{m}$ 中的元素组成。因此，$\operatorname{prof}(M)=+\infty$ 当且仅当
$M=0$；并且由 (16.4.1)，
\begin{equation}
\label{0.16.4.5.1-fr}
\tag{16.4.5.1}
  \operatorname{prof}(M)\leq\dim(M)\quad\text{若}\quad M\neq0.
\end{equation}
显然，$\operatorname{prof}(M^k)=\operatorname{prof}(M)$ 对每个 $k>0$ 成立。
\end{definition}

\begin{proposition}[16.4.6]
\label{0.16.4.6-fr}
设 $A$ 是 Noether 局部环，$\mathfrak{m}$ 是它的极大理想，而 $M$ 是有限型
$A$-模。
\begin{enumerate}
  \item[\rm{(i)}] $\operatorname{prof}(M)=0$ 当且仅当
  $\mathfrak{m}\in\operatorname{Ass}(M)$（也就是说，$M$ 有一个子模同构于
  剩余域 $k=A/\mathfrak{m}$，即 $A$ 的剩余域）。
  \item[\rm{(ii)}] 对每个 $M$-正则元素 $x$，若它属于 $\mathfrak{m}$，则有
  \begin{equation}
  \label{0.16.4.6.1-fr}
  \tag{16.4.6.1}
    \operatorname{prof}(M/xM)=\operatorname{prof}(M)-1.
  \end{equation}
  \item[\rm{(iii)}] 若 $M\neq0$，则
  \begin{equation}
  \label{0.16.4.6.2-fr}
  \tag{16.4.6.2}
    \operatorname{prof}(M)
    \leq\inf_{\mathfrak{p}\in\operatorname{Ass}(M)}\dim(A/\mathfrak{p})
    \leq\dim(M).
  \end{equation}
  \oldpage[0\textsubscript{IV-1}]{34}
  \item[\rm{(iv)}] 有
  $\operatorname{prof}_A(M)=
  \operatorname{prof}_{\widehat{A}}(\widehat{M})$。
\end{enumerate}
\end{proposition}

\begin{proof}
(i) 说 $\operatorname{prof}(M)=0$，意味着不存在 $x\in\mathfrak{m}$，使得以
$x$ 为比的位似在 $M$ 中是单射；换言之，每个 $x\in\mathfrak{m}$ 都属于 $M$
的某个相伴素理想（Bourbaki，\emph{Alg.~comm.}，第 IV 章，\S~1，第 1 号，
命题 2 的推论 2）。但是，$\mathfrak{m}$ 只有在等于 $M$ 的某个相伴素理想时，
才可能是这些理想的并（Bourbaki，\emph{Alg.~comm.}，第 II 章，\S~1，第 1 号，
命题 2）。

(ii) 若 $(x_i)_{1\leq i\leq m}$ 是一个 $(M/xM)$-正则序列，且由
$\mathfrak{m}$ 中的元素组成，则由 $x$ 与诸 $x_i$ 组成的序列是 $M$-正则的；
由定义可知，当序列 $(x_i)_{1\leq i\leq m}$ 是极大的 $(M/xM)$-正则序列时，
前一序列也是极大的，结论由此得出。

(iii) 对 $n=\operatorname{prof}(M)$ 作归纳；当 $n=0$ 时断言显然。我们将使用
下述引理：
\end{proof}

\begin{lemma}[16.4.6.3]
\label{0.16.4.6.3-fr}
设 $A$ 是 Noether 环，$M$ 是有限型 $A$-模，而 $t\in A$ 是 $M$-正则元素。
则对每个 $\mathfrak{p}\in\operatorname{Ass}(M)$，在所有包含
$\mathfrak{p}+At$ 的素理想中，任取一个极小元，它都是 $M/tM$ 的相伴素理想。
\end{lemma}

\begin{proof}
我们知道存在正合序列
\[
  0\to M'\to M\to M''\to0
\]
使得 $\operatorname{Ass}(M')=\{\mathfrak{p}\}$，且
$\operatorname{Ass}(M'')=\operatorname{Ass}(M)-\{\mathfrak{p}\}$（Bourbaki，
\emph{Alg.~comm.}，第 IV 章，\S~1，第 1 号，命题 4）。因为 $t$ 是
$M$-正则的，它也是 $M'$-正则和 $M''$-正则的
（\emph{同上}，\S~1，第 1 号，命题 2 的推论 2）。由此得出序列
\[
  0\to M'/tM'\to M/tM\to M''/tM''\to0
\]
是正合的（15.1.18）。因此有
$\operatorname{Ass}(M'/tM')\subset\operatorname{Ass}(M/tM)$。但是，$A$ 中
包含 $\mathfrak{p}$ 的素理想正是 $\operatorname{Supp}(M')$ 的点
（\emph{同上}，\S~1，第 3 号，命题 7），所以包含 $\mathfrak{p}+At$ 的素理想
正是 $M'/tM'=M'\otimes_A(A/tA)$ 的支撑的点（$0_{\mathrm I}$，1.7.5）；
这些素理想之中的一个极小元因而属于 $\operatorname{Ass}(M'/tM')$
（Bourbaki，\emph{同上}，\S~1，第 3 号，命题 7 的推论 1），因而当然也属于
$\operatorname{Ass}(M/tM)$。
\end{proof}

\begin{proof}[（16.4.6）证明的续篇]
引理既已证明，现在回到 (iii) 的证明。取 $x\in\mathfrak{m}$ 为 $M$-正则元素，
于是由 (ii) 有 $\operatorname{prof}(M/xM)=n-1$；对每个
$\mathfrak{p}\in\operatorname{Ass}(M)$，由引理可知，存在一个素理想
$\mathfrak{p}'$，它与 $M/xM$ 相伴且包含 $\mathfrak{p}+xA$。但是，因为 $x$ 是 $M$-正则的，
有 $x\notin\mathfrak{p}$，所以 $\mathfrak{p}'\neq\mathfrak{p}$，从而
$\dim(A/\mathfrak{p}')\leq\dim(A/\mathfrak{p})-1$。而归纳假设表明
$n-1\leq\dim(A/\mathfrak{p}')$；故有 $n\leq\dim(A/\mathfrak{p})$，对每个
$\mathfrak{p}\in\operatorname{Ass}(M)$ 都成立。

(iv) 只需考虑 $M\neq0$ 的情形。回忆
$\widehat{M}=M\otimes_A\widehat{A}$，这里的等式是在典范同构意义下成立的（$0_{\mathrm I}$，7.3.3）；
因为 $\widehat{A}$ 是平坦 $A$-模，有典范同构
\[
  \operatorname{Ext}^i_A(k,M)\otimes_A\widehat{A}
  \xrightarrow{\sim}
  \operatorname{Ext}^i_{\widehat{A}}
  (k,M\otimes_A\widehat{A})
\]
这是因为 $k=k\otimes_A\widehat{A}$ 在典范同构意义下成立
（$0_{\mathrm{III}}$，12.3.5）。因此，(iv) 由 (16.4.4) 以及 $\widehat{A}$ 是
忠实平坦 $A$-模这一事实得出（$0_{\mathrm I}$，7.3.5）。
\end{proof}

\begin{env}[16.4.6.4]
\label{0.16.4.6.4-fr}
注意，如果 $\mathfrak{p}$ 是 $A$ 的一个素理想，则既可能有
$\operatorname{prof}_{A_{\mathfrak{p}}}(M_{\mathfrak{p}})
<\operatorname{prof}_A(M)$，也可能有
$\operatorname{prof}_{A_{\mathfrak{p}}}(M_{\mathfrak{p}})
>\operatorname{prof}_A(M)$。第一种情形的例子，
\oldpage[0\textsubscript{IV-1}]{35}
只需取 $A$ 为维数
$\geq1$ 的局部整环；此时 $\operatorname{prof}(A)\geq1$，但
$\operatorname{prof}(K)=0$，其中 $K$ 是 $A$ 的分式域。

第二种情形的例子如下。考察一个 Noether 局部整环 $A_0$，其维数 $\geq2$，
它的极大理想为 $\mathfrak{m}_0$，而 $k=A_0/\mathfrak{m}_0$ 是它的剩余域；令
$A=A_0\oplus\mathfrak{J}$ 为 $A_0$ 由 $A_0$-模 $k$ 所作的平凡扩张（18.2.3），
其中理想 $\mathfrak{J}$ 与 $A_0$-模 $k$ 同构（于是 $A$ 中的乘法由
$(x,z)(x',z')=(xx',xz'+x'z)$ 给出）。显然，$\mathfrak{J}$ 是 $A$ 的幂零根，
而 $A/\mathfrak{J}$ 同构于 $A_0$；取 $\mathfrak{p}$ 为 $A$ 的任一素理想，则它形如
$\mathfrak{p}_0\oplus\mathfrak{J}$，其中 $\mathfrak{p}_0$ 是 $A_0$ 的一个素理想；
特别地，$\mathfrak{m}=\mathfrak{m}_0\oplus\mathfrak{J}$ 是 $A$ 的唯一极大理想。
因为 $\mathfrak{m}_0$ 的每个元素都零化 $\mathfrak{J}$，有
$\operatorname{prof}(A)=\operatorname{prof}(A_{\mathfrak{m}})=0$；相反，若
$\mathfrak{p}_0\neq\mathfrak{m}_0$，则 $\mathfrak{J}$ 的元素被
$\mathfrak{m}_0-\mathfrak{p}_0$ 的元素零化，所以 $A_{\mathfrak{p}}$ 典范同构于
$(A_0)_{\mathfrak{p}_0}$；又因 $A_0$ 是整环，有
$\operatorname{prof}(A_{\mathfrak{p}})
=\operatorname{prof}((A_0)_{\mathfrak{p}_0})\geq1$，只要 $\mathfrak{p}_0\neq0$。
\end{env}

\begin{corollary}[16.4.7]
\label{0.16.4.7-fr}
一个既约 Noether 局部环 $A$ 的深度为 $0$，当且仅当 $A$ 是域。
\end{corollary}

\begin{proof}
事实上，因为 $A$ 的极小素理想之交为 $0$，所以 $A$ 没有嵌入相伴素理想；
由 (16.4.6, (i))，若 $\operatorname{prof}(A)=0$，则极大理想
$\mathfrak{m}$ 是 $A$ 的极大理想，同时也是极小素理想，故它是 $A$ 的唯一素理想；又因 $A$ 既约，
所以 $\mathfrak{m}=0$。
\end{proof}

\begin{proposition}[16.4.8]
\label{0.16.4.8-fr}
设 $A$、$B$ 是两个 Noether 局部环，$\rho:A\to B$ 是局部同态，而 $M$ 是
$B$-模，使得 $M_{[\rho]}$ 是有限型 $A$-模。则有
\begin{equation}
\label{0.16.4.8.1-fr}
\tag{16.4.8.1}
  \operatorname{prof}_A(M_{[\rho]})=\operatorname{prof}_B(M).
\end{equation}
\end{proposition}

\begin{proof}
只需考虑 $M\neq0$ 的情形；设 $\operatorname{prof}_A(M_{[\rho]})=n$。若
$(x_i)_{1\leq i\leq n}$ 是一个极大的 $(M_{[\rho]})$-正则序列，由 $A$ 的极大理想
中的元素组成，则诸 $\rho(x_i)$ 构成一个 $M$-正则序列，由 $B$ 的极大理想中的
元素组成。令
$N=M/(\rho(x_1)M+\cdots+\rho(x_n)M)$，则有
$\operatorname{prof}_B(N)=\operatorname{prof}_B(M)-n$（16.4.6, (ii)）；同样，
$N_{[\rho]}=M_{[\rho]}/(x_1M_{[\rho]}+\cdots+x_nM_{[\rho]})$，且
$\operatorname{prof}_A(N_{[\rho]})=0$。这样就归结为证明 $n=0$ 时的命题。

设 $k$ 是 $A$ 的剩余域。因为 $k$ 是 $A$ 的商，故 $A$-模
$P=\operatorname{Hom}_A(k,M_{[\rho]})$ 是
$\operatorname{Hom}_A(A,M_{[\rho]})=M_{[\rho]}$ 的子模，且假设推出
$P\neq0$（16.4.6, (i)）。此外，$P$ 是有限型 $A$-模（因为 $M_{[\rho]}$ 是有限型
$A$-模），所以它是有限维 $k$-向量空间（这里用到它被 $A$ 的极大理想零化），进而
是有限长度 $A$-模。另一方面，$P$ 也是 $B$-模 $M$ 的子模；既然它作为
$A$-模具有有限长度，作为 $B$-模也就具有有限长度。因为它 $\neq0$，所以它含有
一个单 $B$-子模，也就是同构于 $B$ 的剩余域的子模；于是由 (16.4.6, (i)) 得出
$\operatorname{prof}_B(M)=0$。
\end{proof}

\begin{definition}[16.4.9]
\label{0.16.4.9-fr}
设 $A$ 是 Noether 局部环，而 $M$ 是有限型 $A$-模。若 $M\neq0$，把 $M$ 的
\emph{余深度}记作 $\operatorname{coprof}_A(M)$ 或 $\operatorname{coprof}(M)$，并定义为
有限整数 $\dim_A(M)-\operatorname{prof}_A(M)\geq0$（16.4.5.1）。若 $M=0$，约定
$\operatorname{coprof}(M)=0$。

有 $\operatorname{coprof}(M^k)=\operatorname{coprof}(M)$，对每个 $k>0$ 成立。
\end{definition}

\begin{proposition}[16.4.10]
\label{0.16.4.10-fr}
设 $A$ 是 Noether 局部环，而 $M$ 是有限型 $A$-模。
\oldpage[0\textsubscript{IV-1}]{36}
\begin{enumerate}
  \item[(i)] 对每个 $M$-正则元素 $x$，若它属于 $A$ 的极大理想，则有
  $\operatorname{coprof}(M/xM)=\operatorname{coprof}(M)$。
  \item[(ii)] 有
  $\operatorname{coprof}_A(M)=\operatorname{coprof}_{\widehat A}(\widehat M)$。
\end{enumerate}
\end{proposition}

\begin{proof}
事实上，(i) 由 (16.3.4) 和 (16.4.6, (ii)) 得出，而 (ii) 由 (16.2.4) 和
(16.4.6, (iv)) 得出。
\end{proof}

稍后我们将证明（IV，6.11.5），对每个素理想 $\mathfrak{p}$（在 $A$ 中），都有
$\operatorname{coprof}_{A_{\mathfrak{p}}}(M_{\mathfrak{p}})
\leq\operatorname{coprof}_A(M)$。

\begin{proposition}[16.4.11]
\label{0.16.4.11-fr}
在 (16.4.8) 的假设下，有
\begin{equation}
\label{0.16.4.11.1-fr}
\tag{16.4.11.1}
  \operatorname{coprof}_A(M_{[\rho]})=\operatorname{coprof}_B(M).
\end{equation}
\end{proposition}

\begin{proof}
这由 (16.1.9) 和 (16.4.8) 得出。
\end{proof}

\subsection{Cohen--Macaulay 模}
\label{subsection:0.16.5-fr}

\begin{definition}[16.5.1]
\label{0.16.5.1-fr}
设 $A$ 是 Noether 局部环。若有限型 $A$-模 $M$ 满足
$\operatorname{coprof}(M)=0$，即 $M=0$，或 $M\neq0$ 且
$\dim(M)=\operatorname{prof}(M)$，则称它是一个\emph{Cohen--Macaulay 模}。
如果 $A$-模 $A$ 是 Cohen--Macaulay 模，则称 $A$ 是
\emph{Cohen--Macaulay 环}。
\end{definition}

有限长度的 $A$-模是 Cohen--Macaulay 模，因为它的维数为 $0$（16.1.10）。
既约局部环 $A$ 若维数为 $1$，就是 Cohen--Macaulay 环，因为这时它的极大理想
$\mathfrak{m}$ 不可能相伴于 $A$（Bourbaki，\emph{Alg.~comm.}，第 IV 章，
\S~2，第 5 号，命题 10），所以由 (16.4.6, (i))，
$\operatorname{prof}(A)=1$。若 $A$ 是一个\emph{维数 $\geq2$ 的整闭局部整环}，则
$\operatorname{prof}(A)\geq2$。事实上，若 $x\neq0$ 是 $\mathfrak{m}$ 的元素，
我们知道（Bourbaki，\emph{Alg.~comm.}，第 VII 章，\S~1，第 4 号，命题 8），
$A/xA$ 的相伴素理想高度为 $1$，因而不同于 $\mathfrak{m}$；所以它们的并也不同于
$\mathfrak{m}$（Bourbaki，\emph{Alg.~comm.}，第 II 章，\S~1，第 1 号，命题 2）。
因此，存在序列 $(x,y)$，它是 $A$-正则的，且由 $\mathfrak{m}$ 中元素组成，断言得证。
于是，维数为 $2$ 的整闭局部环是 Cohen--Macaulay 环。

\begin{proposition}[16.5.2]
\label{0.16.5.2-fr}
设 $A$ 是 Noether 局部环，而 $M$ 是有限型 $A$-模。$M$ 是 Cohen--Macaulay 模，
当且仅当 $\widehat M$ 作为 $\widehat A$-模是 Cohen--Macaulay 模。
\end{proposition}

\begin{proof}
这立即由 (16.4.10, (ii)) 得出。
\end{proof}

\begin{proposition}[16.5.3]
\label{0.16.5.3-fr}
在 (16.4.8) 的假设下，$M$ 作为 $B$-模是 Cohen--Macaulay 模，当且仅当
$M_{[\rho]}$ 作为 $A$-模是 Cohen--Macaulay 模。
\end{proposition}

\begin{proof}
这由 (16.4.11) 得出。
\end{proof}

\begin{proposition}[16.5.4]
\label{0.16.5.4-fr}
设 $A$ 是 Noether 局部环，而 $M$ 是有限型 $A$-模，而且它 $\neq0$。
若 $M$ 是 Cohen--Macaulay $A$-模，则有
$\dim(A/\mathfrak{p})=\dim(M)$，对每个素理想
$\mathfrak{p}\in\operatorname{Ass}(M)$ 都成立；特别地，$M$ 没有嵌入相伴素理想。
\end{proposition}

\begin{proof}
这立即由不等式 (16.4.6.2) 得出。
\end{proof}

还可以说 $M$（或 $\operatorname{Supp}(M)$）是\emph{等维的}（16.1.7）。

\oldpage[0\textsubscript{IV-1}]{37}

\begin{proposition}[16.5.5]
\label{0.16.5.5-fr}
设 $A$ 是 Noether 局部环，$\mathfrak{m}$ 是它的极大理想，而 $M$ 是有限型
$A$-模，而且它 $\neq0$；设 $x$ 是 $\mathfrak{m}$ 的一个元素。假设 $M$ 是
Cohen--Macaulay 模。则下列条件等价：
\begin{enumerate}
  \item[a)] $x$ 是 $M$-正则的；
  \item[b)] $\dim(M/xM)=\dim(M)-1$；
  \item[c)] $x$ 不属于 $\operatorname{Supp}(M)$ 的任何极小元。
\end{enumerate}
此外，这时 $M/xM$ 是 Cohen--Macaulay 模。
\end{proposition}

\begin{proof}
我们已经看到（16.5.4），诸素理想 $\mathfrak{p}$（即
$\operatorname{Ass}(M)$ 的元素）都是 $\operatorname{Supp}(M)$ 的极小元，并且
$\dim(A/\mathfrak{p})=\dim(M)$；
a) 与 c) 的等价性因而由 Bourbaki，\emph{Alg.~comm.}，第 IV 章，\S~1，
第 1 号，命题 2 的推论 2 得出；b) 与 c) 的等价性由 (16.3.4) 得出。最后，
$M/xM$ 是 Cohen--Macaulay 模这一事实由 (16.4.10, (i)) 得出。
\end{proof}

\begin{corollary}[16.5.6]
\label{0.16.5.6-fr}
在 (16.5.5) 的假设下，设 $(x_i)_{1\leq i\leq r}$ 是由 $\mathfrak{m}$ 中元素
组成的序列。下列条件等价：
\begin{enumerate}
  \item[a)] 序列 $(x_i)$ 是 $M$-正则的。
  \item[b)] 有
  $\dim(M/(x_1M+\cdots+x_rM))=\dim(M)-r$。
  \item[c)] 诸 $x_i$ 构成 $M$ 的某个参数系的一部分。
\end{enumerate}
此外，当这些条件满足时，$N=M/(x_1M+\cdots+x_rM)$ 是 Cohen--Macaulay 模，
所以对每个理想 $\mathfrak{p}\in\operatorname{Ass}(N)$，有
$\dim(A/\mathfrak{p})=\dim(M)-r$。
\end{corollary}

\begin{proof}
我们已经知道 b) 与 c) 等价（16.3.6），并且 a) 推出 c)（16.4.1），无须对 $M$
另加假设。为证明 c) 推出 a) 并证明陈述的最后断言，对 $r$ 作归纳；当 $r=0$ 时
断言是平凡的。因为 $x_1$ 构成某个参数系的一部分，故
$\dim(M/x_1M)=\dim(M)-1$（16.3.6），所以 $x_1$ 是 $M$-正则的，并且由 (16.5.5)，
$M/x_1M$ 是 Cohen--Macaulay 模。因为序列 $(x_i)_{2\leq i\leq r}$ 构成
$P=M/x_1M$ 的某个参数系的一部分，归纳假设表明 $(x_i)_{2\leq i\leq r}$ 是
$P$-正则序列，而且 $P/(x_2P+\cdots+x_rP)=N$ 是 Cohen--Macaulay 模；此外，
$(x_i)_{1\leq i\leq r}$ 是 $M$-正则的。
\end{proof}

\begin{remark}[16.5.7]
\label{0.16.5.7-fr}
由 (16.5.6)，当 $M$ 是 Cohen--Macaulay 模时，\emph{$M$ 的参数系}与
\emph{由 $\mathfrak{m}$ 中元素组成的极大 $M$-正则序列}是同一类对象。反过来，
由 (16.4.1)，如果一个非零有限型 $A$-模有一个 $M$-正则序列，同时又是 $M$
的参数系，则 $M$ 是 Cohen--Macaulay 模。(16.5.6) 中陈述的最后一个性质也
\emph{刻画} Cohen--Macaulay 模：
\end{remark}

\begin{proposition}[16.5.8]
\label{0.16.5.8-fr}
设 $A$ 是 Noether 局部环，而 $M$ 是有限型 $A$-模。假设对每个序列
$(x_i)_{1\leq i\leq r}$，只要它由 $A$ 中元素组成并构成 $M$ 的某个参数系的一部分，
以及对每个素理想 $\mathfrak{p}$，只要它相伴于
$N=M/(x_1M+\cdots+x_rM)$，都有
$\dim(A/\mathfrak{p})=\dim(N)$。则 $M$ 是 Cohen--Macaulay 模。
\end{proposition}

\begin{proof}
设 $(y_i)_{1\leq i\leq n}$ 是 $M$ 的一个参数系；只需证明序列 $(y_i)$ 是
$M$-正则的。对 $n=\dim(M)$ 作归纳；当 $n=0$ 时命题是平凡的。根据假设
（取 $r=0$），每个素理想 $\mathfrak{p}_j$ 若相伴于 $M$，就满足
$\dim(A/\mathfrak{p}_j)=\dim(M)$。因为 $y_1$ 是
\oldpage[0\textsubscript{IV-1}]{38}
$M$ 的某个参数系的一部分，由 (16.3.7) 和 (16.3.4) 可知，$y_1$ 不属于任何 $\mathfrak{p}_j$，
所以是 $M$-正则的。若令 $P=M/y_1M$，则 $(y_i)_{2\leq i\leq n}$ 是 $P$ 的
一个参数系，而且 $P$ 显然满足命题的假设；应用归纳假设可知，
$(y_i)_{2\leq i\leq n}$ 是 $P$-正则序列，命题得证。
\end{proof}

\begin{proposition}[16.5.9]
\label{0.16.5.9-fr}
设 $A$ 是 Noether 局部环，而 $M$ 是有限型 $A$-模；假设 $M$ 是
Cohen--Macaulay 模。设 $\mathfrak{p}\in\operatorname{Supp}(M)$，并令
$r=\dim(M)-\dim(M/\mathfrak{p}M)$。则存在一个 $M$-正则序列
$(x_i)_{1\leq i\leq r}$，其元素属于 $\mathfrak{p}$；对每个具有这些性质的序列，都有
\[
  \dim(M/(x_1M+\cdots+x_rM))
  =\dim(M/\mathfrak{p}M)=\dim(A/\mathfrak{p}),
\]
而 $\mathfrak{p}$ 是 $\operatorname{Ass}(M/(x_1M+\cdots+x_rM))$ 的一个极小元。
\end{proposition}

\begin{proof}
先证明诸 $x_i$ 的存在性；当 $r=0$ 时无须证明。当 $r>0$ 时，对 $r$ 作归纳。
因为每个素理想 $\mathfrak{q}_i$ 若相伴于 $M$，就有
$\dim(A/\mathfrak{q}_i)=\dim(M)$（16.5.4），假设
$\dim(M/\mathfrak{p}M)<\dim(M)$ 推出 $\mathfrak{p}$ 不可能包含在任何
$\mathfrak{q}_i$；后者是包含 $\operatorname{Ann}(M)$ 的素理想之中的极小元
（16.1.2.2 与 16.1.7）。由此推出 $\mathfrak{p}$ 也不包含在诸
$\mathfrak{q}_i$ 的并（Bourbaki，\emph{Alg.~comm.}，第 II 章，\S~1，第 1 号，
命题 2），因而存在 $x_1\in\mathfrak{p}$，它是 $M$-正则的（Bourbaki，
\emph{Alg.~comm.}，第 IV 章，\S~1，第 1 号，命题 2 的推论 2）。所以由
(16.5.5)，$N=M/x_1M$ 是维数为 $\dim(M)-1$ 的 Cohen--Macaulay 模；又因为
$x_1\in\mathfrak{p}$，有 $N/\mathfrak{p}N=M/\mathfrak{p}M$，从而
$\dim(N)-\dim(N/\mathfrak{p}N)=r-1$。于是可以对 $N$ 应用归纳假设；若
$(x_i)_{2\leq i\leq r}$ 是 $N$-正则序列，其元素属于 $\mathfrak{p}$，显然
$(x_i)_{1\leq i\leq r}$ 是 $M$-正则序列，其元素属于 $\mathfrak{p}$。

若令 $P=M/(x_1M+\cdots+x_rM)$，则有
$P/\mathfrak{p}P=M/\mathfrak{p}M$，因为诸 $x_i$ 属于 $\mathfrak{p}$；并由 (16.5.6) 有
$\dim(P)=\dim(P/\mathfrak{p}P)$。但是，因为 $P$ 是 Cohen--Macaulay 模，也有
$\dim(P)=\dim(A/\mathfrak{p}')$，对每个
$\mathfrak{p}'\in\operatorname{Ass}(P)$ 都成立（16.5.4）。因为
$\mathfrak{p}\in\operatorname{Supp}(P)$（$0_{\mathrm I}$，1.7.5），
$\mathfrak{p}$ 包含某个 $\mathfrak{p}'\in\operatorname{Ass}(P)$
（Bourbaki，\emph{Alg.~comm.}，第 IV 章，\S~1，第 4 号，定理 2）；又因为
\[
  \dim(P)=\dim(P/\mathfrak{p}P)
  \leq\dim(A/\mathfrak{p})
  \leq\dim(A/\mathfrak{p}')=\dim(P),
\]
必有 $\mathfrak{p}=\mathfrak{p}'$（16.1.2.2），证明完毕。
\end{proof}

\begin{corollary}[16.5.10]
\label{0.16.5.10-fr}
在 (16.5.9) 的同一假设下：
\begin{enumerate}
  \item[(i)] $M_{\mathfrak{p}}$ 作为 $A_{\mathfrak{p}}$-模是 Cohen--Macaulay 模。
  \item[(ii)] 有
  \begin{equation}
  \label{0.16.5.10.1-fr}
  \tag{16.5.10.1}
    \dim(M)=\dim_A(M/\mathfrak{p}M)
    +\dim_{A_{\mathfrak{p}}}(M_{\mathfrak{p}}).
  \end{equation}
\end{enumerate}
\end{corollary}

\begin{proof}
沿用 (16.5.9) 的记号，设 $x_i'$ 是 $x_i$ 在 $A_{\mathfrak{p}}$ 中的典范像
$(1\leq i\leq r)$。因为诸 $x_i'$ 属于 $A_{\mathfrak{p}}$ 的极大理想，并且由
平坦性（15.1.14），它们构成 $M_{\mathfrak{p}}$-正则序列，所以有
\[
  \operatorname{prof}_{A_{\mathfrak{p}}}(M_{\mathfrak{p}})
  \geq r=\dim(M)-\dim(M/\mathfrak{p}M)
  \geq\dim_{A_{\mathfrak{p}}}(M_{\mathfrak{p}})
\]
（由 (16.1.8.1)）。但是，考虑到 (16.4.5.1)，这个不等式的三项必然相等，
推论得证。
\end{proof}

\oldpage[0\textsubscript{IV-1}]{39}

\begin{corollary}[16.5.11]
\label{0.16.5.11-fr}
设 $A$ 是 Noether 局部环；假设存在有限型 $A$-模 $M$，其支撑为
$\operatorname{Spec}(A)$，并且它是 Cohen--Macaulay 模。则对每个素理想
$\mathfrak{p}$（在 $A$ 中），都有
\begin{equation}
\label{0.16.5.11.1-fr}
\tag{16.5.11.1}
  \dim(A)=\dim(A/\mathfrak{p})+\dim(A_{\mathfrak{p}}).
\end{equation}
\end{corollary}

\begin{proof}
事实上，$\operatorname{Supp}(M/\mathfrak{p}M)
=\operatorname{Spec}(A/\mathfrak{p})$（$0_{\mathrm I}$，1.7.5），且
$\operatorname{Supp}(M_{\mathfrak{p}})
=\operatorname{Spec}(A_{\mathfrak{p}})$；关系式 (16.5.11.1) 因而是
(16.5.10.1) 的一个特例。
\end{proof}

\begin{corollary}[16.5.12]
\label{0.16.5.12-fr}
满足 (16.5.11) 条件的 Noether 局部环 $A$ 的每个商环（特别是
Cohen--Macaulay 局部环的每个商环）都是\emph{悬链的}。
\end{corollary}

\begin{proof}
显然，只需证明 $A$ 本身是悬链的；也就是说，对两个素理想
$\mathfrak{p}\subset\mathfrak{q}$（均在 $A$ 中），有 (16.1.4.2)
\[
  \dim(A_{\mathfrak{q}})
  =\dim(A_{\mathfrak{p}})
  +\dim(A_{\mathfrak{q}}/\mathfrak{p}A_{\mathfrak{q}}).
\]
而由 (16.5.10, (i))，$A_{\mathfrak{q}}$ 满足与 $A$ 相同的假设，所以前一关系
就是把 (16.5.11.1) 应用于 $A_{\mathfrak{q}}$ 以及素理想
$\mathfrak{p}A_{\mathfrak{q}}$；后者是 $A_{\mathfrak{q}}$ 的素理想。
\end{proof}

\begin{env}[16.5.13]
\label{0.16.5.13-fr}
设 $A$ 是 Noether 环，而 $M$ 是有限型 $A$-模。称 $M$ 是一个
\emph{Cohen--Macaulay $A$-模}，如果对每个素理想 $\mathfrak{p}$（在 $A$ 中），
$M_{\mathfrak{p}}$ 作为 $A_{\mathfrak{p}}$-模都是 Cohen--Macaulay 模；由 (16.5.10)，
当 $A$ 是\emph{局部}环时，这个定义与 (16.5.1) 相同。称 $A$ 是
\emph{Cohen--Macaulay 环}，如果所有 $A_{\mathfrak{p}}$ 都是 Cohen--Macaulay 环。
显然，如果 $M$（相应地，$A$）是 Cohen--Macaulay
$A$-模（相应地，Cohen--Macaulay 环），则 $S^{-1}M$ 是 Cohen--Macaulay
$S^{-1}A$-模（相应地，$S^{-1}A$ 是 Cohen--Macaulay 环）；这里 $S$ 是 $A$ 的任一乘法子集。
\end{env}

\section{正则环}
\label{section:0.17-fr}

\subsection{正则环的定义}
\label{subsection:0.17.1-fr}

\begin{proposition}[17.1.1]
\label{0.17.1.1-fr}
设 $A$ 是维数为 $n$ 的 Noether 局部环，$\mathfrak{m}$ 是它的极大理想，
$k=A/\mathfrak{m}$ 是它的剩余域。下列条件等价：
\begin{enumerate}
  \item[\rm a)] 分次 $k$-模的典范满同态 (15.1.1.1)
  \[
    \varphi:\mathbf{S}^*_k(\mathfrak{m}/\mathfrak{m}^2)
    \longrightarrow\operatorname{gr}^*_{\mathfrak{m}}(A)
  \]
  是双射（这里后者是分次 $k$-模，即 $A$ 关于 $\mathfrak{m}$-预进制滤过
  的相伴分次模）。
  \item[\rm b)] 有
  $\operatorname{rg}_k(\mathfrak{m}/\mathfrak{m}^2)=n$。
  \item[\rm c)] 理想 $\mathfrak{m}$ 有一组由 $n$ 个元素构成的生成元。
  \item[\rm d)] 理想 $\mathfrak{m}$ 有一组构成 $A$-正则序列的生成元。
\end{enumerate}
\end{proposition}


\begin{proof}
由 Nakayama 引理，$\operatorname{rg}_k(\mathfrak{m}/\mathfrak{m}^2)$ 等于
$\mathfrak{m}$ 的生成元系所含元素的最小数目，故 b) 与 c) 等价。另一方面，
若 $(x_i)_{1\leq i\leq r}$ 是 $\mathfrak{m}$ 的一组生成元，并且它们的类
$\overline{x}_i$（模 $\mathfrak{m}^2$）构成 $k$-向量空间
$\mathfrak{m}/\mathfrak{m}^2$ 的一组基，则
$\mathbf{S}^*(\mathfrak{m}/\mathfrak{m}^2)$ 同构于
\oldpage[0\textsubscript{IV-1}]{40}
多项式环 $k[\mathrm{T}_1,\ldots,\mathrm{T}_r]$；考虑到每个有限型 $A$-模关于
$\mathfrak{m}$-预进制滤过都是分离的（$0_{\mathrm I}$，7.3.5），由 (15.1.9)
可知条件 a) 与 d) 等价；又因为每个 $A$-正则序列，只要各项都属于 $\mathfrak{m}$，
便至多含有 $n$ 个元素（16.4.1），所以 d) 蕴含 c)。还须证明 c) 蕴含 a)。

为简便起见，置
$\mathrm{S}=\mathbf{S}^*(\mathfrak{m}/\mathfrak{m}^2)
=k[\mathrm{T}_1,\ldots,\mathrm{T}_n]$，
$\mathrm{G}=\operatorname{gr}^*_{\mathfrak{m}}(A)$，并考虑正合列
$0\longrightarrow\mathfrak{J}\longrightarrow\mathrm{S}
\xrightarrow{\varphi}\mathrm{G}\longrightarrow0$；这里 $\mathfrak{J}$ 是
$\varphi$ 的核，因而是 $\mathrm{S}$ 的一个分次理想。于是，对每个整数 $s>0$，有
\begin{equation}
\label{0.17.1.1.1-fr}
\tag{17.1.1.1}
  \binom{s+n-1}{n-1}
  =\operatorname{long}(\mathrm{S}_s)
  =\operatorname{long}(\mathrm{G}_s)
  +\operatorname{long}(\mathfrak{J}_s).
\end{equation}

假设 $\mathfrak{J}\neq0$，于是存在齐次元 $u\in\mathfrak{J}$，其次数
$h>0$；$\mathfrak{J}_s$ 包含 $u\mathrm{S}_{s-h}$（$s\geq h$），而因为
$u\neq0$ 且 $\mathrm{S}$ 是整环，$u\mathrm{S}_{s-h}$ 同构于
$\mathrm{S}_{s-h}$（作为 $k$-向量空间）；因而将有
$\operatorname{long}(\mathfrak{J}_s)\geq
\operatorname{long}(\mathrm{S}_{s-h})=\binom{s-h+n-1}{n-1}$，从而当
$s\geq h$ 时，
\begin{equation}
\label{0.17.1.1.2-fr}
\tag{17.1.1.2}
  \operatorname{long}(\mathrm{G}_s)
  \leq\binom{s+n-1}{n-1}-\binom{s-h+n-1}{n-1}.
\end{equation}
但是，(17.1.1.2) 的右端是关于 $s$ 的次数为 $\leq n-2$ 的多项式；另一方面，
$\mathrm{G}_s=\mathfrak{m}^s/\mathfrak{m}^{s+1}$，并且当 $s$ 充分大时，
$\operatorname{long}(\mathrm{G}_s)$ 是关于 $s$ 的次数恰为 $n-1$、首项系数
$>0$ 的多项式（16.2.1），这与不等式 (17.1.1.2) 矛盾。证毕。
\end{proof}

\begin{definition}[17.1.2]
\label{0.17.1.2-fr}
称满足 (17.1.1) 中诸等价条件的 Noether 局部环为\emph{正则局部环}。
\end{definition}

\begin{corollary}[17.1.3]
\label{0.17.1.3-fr}
正则局部环是整环、整闭环，并且是 Cohen--Macaulay 环。
\end{corollary}

\begin{proof}
最后一个断言由 (17.1.1, d)) 得出；另一方面，若 $A$ 是正则局部环，则
$\operatorname{gr}^*_{\mathfrak{m}}(A)$ 同构于
$k[\mathrm{T}_1,\ldots,\mathrm{T}_n]$，因而是整环且完全整闭；由此可知 $A$
也具有这两个性质（Bourbaki，\emph{Alg.~comm.}，第~III~章，\S~2，
第~3~号，命题~1 的推论，以及第~V~章，\S~1，第~5~号，命题~15）。
\end{proof}

\begin{examples}[17.1.4]
\label{0.17.1.4-fr}
\begin{enumerate}
  \item[\rm (i)] 维数为 $0$ 的正则局部环依 (17.1.3) 为整环，因而必为一个
  \emph{域}；反之亦然。
  \item[\rm (ii)] 维数为 $1$ 的 Noether 局部环是正则的，当且仅当它是一个
  \emph{离散赋值环}：事实上，说维数为 $1$ 的 Noether 局部环是正则的，
  就是说它的极大理想是\emph{主理想}；结论由 Bourbaki，\emph{Alg.~comm.}，
  第~VI~章，\S~3，第~6~号，命题~9 得出。
  \item[\rm (iii)] 设 $k$ 是一个域；形式幂级数环
  $A=k[[\mathrm{T}_1,\ldots,\mathrm{T}_n]]$ 是一个\emph{维数为 $n$ 的正则环}；
  事实上，显然诸 $\mathrm{T}_i$ $(1\leq i\leq n)$ 生成极大理想
  $\mathfrak{m}$（在 $A$ 中），并构成一个 $A$-正则序列，因为
  $A/(\mathrm{T}_1A+\cdots+\mathrm{T}_iA)$ 同构于
  $k[[\mathrm{T}_{i+1},\ldots,\mathrm{T}_n]]$。
\end{enumerate}
\end{examples}

\begin{proposition}[17.1.5]
\label{0.17.1.5-fr}
Noether 局部环 $A$ 是正则的，当且仅当它的完备化 $\widehat{A}$ 是正则的。
\end{proposition}

\begin{proof}
\oldpage[0\textsubscript{IV-1}]{41}
事实上，$\widehat{A}$ 的极大理想是 $\mathfrak{m}\widehat{A}$，并且我们知道
$\mathfrak{m}^h/\mathfrak{m}^{h+1}$ 与
$(\mathfrak{m}\widehat{A})^h/(\mathfrak{m}\widehat{A})^{h+1}$ 是同构的
$k$-向量空间；所以 (17.1.1) 的条件 a) 对 $A$ 和 $\widehat{A}$ 是同一个条件。
\end{proof}

\begin{definition}[17.1.6]
\label{0.17.1.6-fr}
给定 Noether 局部环 $A$，其极大理想为 $\mathfrak{m}$，把
\emph{$A$ 的正则参数系}定义为 $A$ 的一个生成 $\mathfrak{m}$ 的参数系
（16.3.5）。
\end{definition}

这样的参数系存在，当且仅当 $A$ 是正则的（17.1.1）；若
$(x_i)_{1\leq i\leq n}$ 是一个正则参数系，则它是 $\mathfrak{m}$ 的一个极小
生成元系；它的各元素模 $\mathfrak{m}^2$ 的类构成 $\mathfrak{m}/\mathfrak{m}^2$
在 $k$ 上的一组基；由 (17.1.1, a)) 和 (15.1.9)，$(x_i)$ 是一个极大
$A$-正则序列（术语由此而来）。

不过应当注意，在正则环中，$A$ 的一个参数系，即使同时也是 $A$-正则序列，
也\emph{不一定}是一个正则参数系；把正则参数系的每个元素都取 $k$ 次幂所得的例子说明了这一点
（15.1.20）。

\begin{proposition}[17.1.7]
\label{0.17.1.7-fr}
设 $A$ 是 Noether 局部环，$\mathfrak{m}$ 是它的极大理想，
$k=A/\mathfrak{m}$ 是它的剩余域，$(x_i)_{1\leq i\leq r}$ 是
$\mathfrak{m}$ 中的元素序列，而 $\mathfrak{J}$ 是理想
$x_1A+\cdots+x_rA$。下列条件等价：
\begin{enumerate}
  \item[\rm a)] $A$ 是正则的，并且诸 $x_i$ 是 $A$ 的某个正则参数系的一部分。
  \item[\rm a')] $A$ 是正则的，并且诸类 $\overline{x}_i$（即 $x_i$ 模
  $\mathfrak{m}^2$ 的类）在 $k$ 上线性无关。
  \item[\rm b)] 诸 $x_i$ 是 $A$ 的某个参数系的一部分，并且 $A/\mathfrak{J}$
  是正则环。
\end{enumerate}
此外，当这些条件成立时，$\mathfrak{J}$ 是素理想。
\end{proposition}

\begin{proof}
最后一个断言立即由如下事实得出：$A/\mathfrak{J}$ 是正则的，因而是整环
（17.1.3）。置 $n=\dim(A)$。a) 与 a') 的等价性是显然的；事实上，正则参数系
中各元素模 $\mathfrak{m}^2$ 的类构成 $\mathfrak{m}/\mathfrak{m}^2$ 在 $k$
上的一组基；反过来，若诸 $\overline{x}_i$ $(1\leq i\leq r)$ 在 $k$ 上
线性无关，则可以找到 $n-r$ 个元素 $x_i\in\mathfrak{m}$
$(r+1\leq i\leq n)$，使得诸 $\overline{x}_i$ $(1\leq i\leq n)$ 构成
$\mathfrak{m}/\mathfrak{m}^2$ 的一组基，因而 $(x_i)_{1\leq i\leq n}$ 是一个
正则参数系。为说明 a) 蕴含 b)，注意，因为诸 $x_i$ 是 $A$ 的某个参数系的一部分，
所以
$\dim(A/\mathfrak{J})=n-r$（16.3.6）；令
$\mathfrak{n}=\mathfrak{m}/\mathfrak{J}$ 为 $A/\mathfrak{J}$ 的极大理想。
有正合列
\begin{equation}
\label{0.17.1.7.1-fr}
\tag{17.1.7.1}
  0\longrightarrow(\mathfrak{m}^2+\mathfrak{J})/\mathfrak{m}^2
  \longrightarrow\mathfrak{m}/\mathfrak{m}^2
  \longrightarrow\mathfrak{n}/\mathfrak{n}^2\longrightarrow0
\end{equation}
因为 $\mathfrak{n}^2=(\mathfrak{m}^2+\mathfrak{J})/\mathfrak{J}$；条件 a) 蕴含
$\operatorname{rg}_k((\mathfrak{m}^2+\mathfrak{J})/\mathfrak{m}^2)=r$，故
$\operatorname{rg}_k(\mathfrak{n}/\mathfrak{n}^2)=n-r
=\dim(A/\mathfrak{J})$；由 (17.1.1, b))，这就证明了 b)。

反过来，为证明 b) 蕴含 a)，注意诸 $x_i$ 是某个参数系的一部分这一事实蕴含
$\dim(A/\mathfrak{J})=n-r$（16.3.6）；另一方面，$A/\mathfrak{J}$ 是正则的这一
假设蕴含
$\operatorname{rg}_k(\mathfrak{n}/\mathfrak{n}^2)=n-r$（17.1.1）；此外，显然有
$\operatorname{rg}_k((\mathfrak{m}^2+\mathfrak{J})/\mathfrak{m}^2)\leq r$，
所以由 (17.1.7.1) 得出
$\operatorname{rg}_k(\mathfrak{m}/\mathfrak{m}^2)\leq r+(n-r)=n$；因而由
(16.2.6) 和 (17.1.1)，$A$ 是正则的。此外，关系式
$\operatorname{rg}_k(\mathfrak{m}/\mathfrak{m}^2)=n$ 又蕴含
$\operatorname{rg}_k((\mathfrak{m}^2+\mathfrak{J})/\mathfrak{m}^2)=r$，从而诸
$\overline{x}_i$ 在 $k$ 上线性无关。证毕。
\end{proof}

\oldpage[0\textsubscript{IV-1}]{42}

\begin{corollary}[17.1.8]
\label{0.17.1.8-fr}
设 $A$ 是 Noether 局部环，$\mathfrak{m}$ 是它的极大理想，$t$ 是
$\mathfrak{m}$ 的一个元素。下列条件等价：
\begin{enumerate}
  \item[\rm a)] $A/tA$ 是正则的，并且 $t$ 在 $A$ 中是非零因子。
  \item[\rm b)] $A$ 是正则的，并且 $t\notin\mathfrak{m}^2$。
\end{enumerate}
\end{corollary}

\begin{proof}
这是 (17.1.7) 中 $r=1$ 的特殊情形（须考虑到 $\mathfrak{m}$ 中的非零因子
是某个参数系的一部分（16.4.1））。
\end{proof}

\begin{corollary}[17.1.9]
\label{0.17.1.9-fr}
设 $A$ 是正则局部环，$\mathfrak{m}$ 是它的极大理想，而 $\mathfrak{J}$ 是
$A$ 的一个包含于 $\mathfrak{m}$ 的理想。下列条件等价：
\begin{enumerate}
  \item[\rm a)] 环 $A/\mathfrak{J}$ 是正则的。
  \item[\rm b)] $\mathfrak{J}$ 由序列 $(x_i)_{1\leq i\leq r}$ 生成，而这个
  序列是 $A$ 的某个正则参数系的一部分。
\end{enumerate}
\end{corollary}

\begin{proof}
我们已经看到（17.1.7），b) 蕴含 a)。反过来，假设 $A/\mathfrak{J}$ 是正则的
（这蕴含 $\mathfrak{J}$ 是素理想），并置 $n=\dim(A)$、
$n-r=\dim(A/\mathfrak{J})$；沿用 (17.1.7) 的记号，正合列 (17.1.7.1)
给出
$\operatorname{rg}_k((\mathfrak{m}^2+\mathfrak{J})/\mathfrak{m}^2)=r$，
所以存在序列 $(x_i)_{1\leq i\leq r}$，它由 $\mathfrak{J}$ 中的元素构成，
并且是 $A$ 的某个正则参数系的一部分。置
$\mathfrak{J}'=x_1A+\cdots+x_rA$；由 (17.1.7)，$\mathfrak{J}'$ 是 $A$ 的
素理想，并且 $A/\mathfrak{J}'$ 是维数为 $n-r$ 的正则环；但是，由于
$A/\mathfrak{J}=(A/\mathfrak{J}')/(\mathfrak{J}/\mathfrak{J}')$，且
$\mathfrak{J}/\mathfrak{J}'$ 是整环 $A/\mathfrak{J}'$ 中的素理想，
$A/\mathfrak{J}$ 与 $A/\mathfrak{J}'$ 的维数只有在
$\mathfrak{J}=\mathfrak{J}'$ 时才可能相等（16.1.2.2）。
\end{proof}

\subsection{模的投射维数与内射维数回顾}
\label{subsection:0.17.2-fr}

\begin{env}[17.2.1]
\label{0.17.2.1-fr}
设 $A$ 是环，$M$ 是 $A$-模。回忆（M, VI, 2），$M$ 的\emph{投射维数}
（相应地，\emph{内射维数}），记作 $\operatorname{dim.proj}(M)$ 或
$\operatorname{dim.proj}_A(M)$（相应地，$\operatorname{dim.inj}(M)$ 或
$\operatorname{dim.inj}_A(M)$），是满足下述条件的最小数 $n$（它可以是整数，也可以等于
$+\infty$）：存在 $M$ 的一个长度为 $n$ 的左投射消解（相应地，右内射消解）。
等价地（\emph{同上}），$n$ 是具有下述性质的最小数：对于每个 $A$-模 $N$，
都有 $\operatorname{Ext}^i_A(M,N)=0$（相应地，
$\operatorname{Ext}^i_A(N,M)=0$）对每个 $i>n$ 成立，或者只需对 $i=n+1$
要求这一等式成立。因此，若 $M'$ 是 $M$ 的任一直和项，则有
$\operatorname{dim.proj}(M')\leq\operatorname{dim.proj}(M)$，这是因为
$\operatorname{Ext}^i_A(M',N)$ 是 $\operatorname{Ext}^i_A(M,N)$ 的直和项；
同样地，$\operatorname{dim.inj}(M')\leq\operatorname{dim.inj}(M)$。
$\operatorname{dim.proj}(M)=n$（相应地，$\operatorname{dim.inj}(M)=n$）的一个
等价条件是：$n$ 是具有下述性质的最小数，即对于每个正合列
\[
  0\longrightarrow R\longrightarrow P_{n-1}\longrightarrow\cdots
  \longrightarrow P_0\longrightarrow M\longrightarrow0
\]
（相应地，
$0\longrightarrow M\longrightarrow Q_0\longrightarrow\cdots
\longrightarrow Q_{n-1}\longrightarrow C\longrightarrow0$），其中所有 $P_i$
（$0\leq i<n$）都是投射的（相应地，所有 $Q_i$（$0\leq i<n$）都是内射的），
则 $R$ 是投射的（相应地，$C$ 是内射的）。
\end{env}

\begin{remarks}[17.2.2]
\label{0.17.2.2-fr}
\begin{enumerate}
  \item[\rm (i)] 说 $M$ 是投射（相应地，内射）$A$-模，等价于说
  $\operatorname{dim.proj}(M)=0$（相应地，$\operatorname{dim.inj}(M)=0$）。
  \item[\rm (ii)] 设 $T$ 是从 $A$-模范畴到某个 Abel 范畴的\emph{协变}加性函子。
  由 (17.2.1) 的定义和导出函子的定义立即可知，
  \oldpage[0\textsubscript{IV-1}]{43}
  如果 $\operatorname{dim.proj}(M)\leq n$（相应地，
  $\operatorname{dim.inj}(M)\leq n$），则有 $L^iT(M)=0$（相应地，
  $R^iT(M)=0$），并且此式\emph{对于每个 $i>n$}成立。
  \item[\rm (iii)] 如果还假设 $A$ 是 Noether 环且 $M$ 是有限型的，则由
  (17.2.1) 中投射维数的最后一种刻画可知：若
  $\operatorname{dim.proj}(M)=n$，则 $M$ 有一个长度为 $n$ 的左消解，
  其中各项都是\emph{有限型投射模}。如果此外 $A$ 是 Noether 局部环，则 $M$
  有一个长度为 $n$ 的消解，其中各项都是\emph{有限型自由模}；事实上，此时
  有限型投射 $A$-模是自由的（Bourbaki，\emph{Alg.~comm.}，第 II 章，
  \S~3，第 2 号，命题 5 的推论 2）。
\end{enumerate}
\end{remarks}

\begin{lemma}[17.2.3]
\label{0.17.2.3-fr}
设 $A$ 是环，$M$ 是 $A$-模。$\operatorname{dim.inj}(M)\leq n$ 当且仅当
对于每个循环 $A$-模 $N$，都有 $\operatorname{Ext}^{n+1}_A(N,M)=0$。
\end{lemma}

\begin{proof}
沿用 (17.2.1) 的记号，须证明 $C$ 是内射的；对于每个 $A$-模 $N$，
$\operatorname{Ext}^{n+1}_A(N,M)$ 同构于 $\operatorname{Ext}^1_A(N,C)$
（M, V, 7），所以有 $\operatorname{Ext}^1_A(N,C)=0$，其中所考虑的是所有
$A$-模 $N$，其形式为 $A/\mathfrak{J}$，而 $\mathfrak{J}$ 是 $A$ 的任一理想。
于是，Ext 正合列表明，典范同态
$\operatorname{Hom}(A,C)\longrightarrow\operatorname{Hom}(\mathfrak{J},C)$ 对
每个理想 $\mathfrak{J}$（在 $A$ 中）都是满射；这蕴含 $C$ 是内射 $A$-模
（M, I, 3.2）。
\end{proof}

\begin{lemma}[17.2.4]
\label{0.17.2.4-fr}
设 $A$ 是 Noether 环，$M$ 是有限型 $A$-模。
$\operatorname{dim.proj}(M)\leq n$ 当且仅当对于每个循环 $A$-模 $N$，都有
$\operatorname{Ext}^{n+1}_A(M,N)=0$。
\end{lemma}

\begin{proof}
事实上，我们知道（M, VI, 2.5），条件 $\operatorname{dim.proj}(M)\leq n$
等价于如下条件：$\operatorname{Ext}^{n+1}_A(M,N)=0$ 对每个\emph{有限型}
$A$-模 $N$ 成立。为说明命题中的条件也充分，按生成元数 $m$ 对 $N$ 归纳：
存在子模 $N_1$，它包含于 $N$，由 $m-1$ 个元素生成，并使
$N_2=N/N_1$ 是循环模；由正合列
$0\longrightarrow N_1\longrightarrow N\longrightarrow N_2\longrightarrow0$，
得到正合列
$\operatorname{Ext}^{n+1}_A(M,N_1)\longrightarrow
\operatorname{Ext}^{n+1}_A(M,N)\longrightarrow
\operatorname{Ext}^{n+1}_A(M,N_2)$，而归纳假设表明，命题中的条件确实蕴含
$\operatorname{Ext}^{n+1}_A(M,N)=0$。
\end{proof}

\begin{corollary}[17.2.5]
\label{0.17.2.5-fr}
设 $A$ 是 Noether 环，$M$ 是 $A$-模。则
\begin{equation}
\label{0.17.2.5.1-fr}
\tag{17.2.5.1}
  \operatorname{dim.inj}_A(M)
  =\sup_{\mathfrak{m}}
  \bigl(\operatorname{dim.inj}_{A_{\mathfrak{m}}}(M_{\mathfrak{m}})\bigr)
\end{equation}
并且，如果 $M$ 是\emph{有限型}的，则
\begin{equation}
\label{0.17.2.5.2-fr}
\tag{17.2.5.2}
  \operatorname{dim.proj}_A(M)
  =\sup_{\mathfrak{m}}
  \bigl(\operatorname{dim.proj}_{A_{\mathfrak{m}}}(M_{\mathfrak{m}})\bigr)
\end{equation}
其中 $\mathfrak{m}$ 遍历 $A$ 的素理想集（或者极大理想集）。
\end{corollary}

\begin{proof}
事实上，若 $N$ 是有限型 $A$-模（因而是有限表示的），则对于每个乘法子集 $S$
（在 $A$ 中）以及每个 $i\geq0$，有
\[
  S^{-1}\operatorname{Ext}^i_A(N,M)
  =\operatorname{Ext}^i_{S^{-1}A}(S^{-1}N,S^{-1}M)
\]
这是由平坦性得到的：考察 $N$ 的一个自由消解，并利用上述关系在 $i=0$ 时成立这一
事实（Bourbaki，\emph{Alg.~comm.}，第 II 章，\S~2，第 7 号，命题 19）
\SourceCorrectionNote{法文印本及外交式转录作“考察 M 的一个自由消解”；应为 N。这里的有限表示假设施加在 Ext 的第一自变量 N 上， 局部化公式也通过它的自由消解计算。}。
特别地，有
$\operatorname{Ext}^i_{A_{\mathfrak{m}}}
(A_{\mathfrak{m}}/\mathfrak{J}A_{\mathfrak{m}},M_{\mathfrak{m}})
=(\operatorname{Ext}^i_A(A/\mathfrak{J},M))_{\mathfrak{m}}$，其中
$\mathfrak{m}$ 遍历所有素理想，$\mathfrak{J}$ 遍历 $A$ 的所有理想；考虑到
(17.2.3)，以及
\oldpage[0\textsubscript{IV-1}]{44}
$A_{\mathfrak{m}}$ 的每个理想都可写成 $\mathfrak{J}A_{\mathfrak{m}}$，其中
$\mathfrak{J}$ 是 $A$ 的某个适当理想，由前述结果及 Bourbaki，
\emph{Alg.~comm.}，第 II 章，\S~3，第 3 号，定理 1，即得公式 (17.2.5.1)。
对 (17.2.5.2) 作同样处理，这次使用 (17.2.4)，并交换 $M$ 与 $N$ 的角色。
\end{proof}

因此，对于 Noether 环和有限型模，由 (17.2.5)，投射维数或内射维数的研究归结为
\emph{局部}环的情形。这时有如下引理。

\begin{lemma}[17.2.6]
\label{0.17.2.6-fr}
设 $A$ 是 Noether 局部环，$k$ 是它的剩余域，$M$ 是有限型 $A$-模。
$\operatorname{dim.proj}(M)\leq n$ 的必要条件是
$\operatorname{Tor}^A_i(M,k)=0$ 对 $i>n$ 成立，而
$\operatorname{Tor}^A_{n+1}(M,k)=0$
则足以推出这一不等式。
\end{lemma}

\begin{proof}
必要性是把备注 (17.2.2 (ii)) 应用于协变函子 $M\mapsto k\otimes_A M$ 所得的
特殊情形。为证明该条件充分，沿用 (17.2.1) 的记号，只须证明 $R$ 是投射的，
这里假设各 $P_i$ 都是有限型的。然而，$\operatorname{Tor}^A_{n+1}(M,k)$
同构于 $\operatorname{Tor}^A_1(R,k)$（M, V, 7）；又已知，由于 $R$ 是有限型的，
条件 $\operatorname{Tor}^A_1(R,k)=0$ 蕴含 $R$ 是自由的（Bourbaki，
\emph{Alg.~comm.}，第 II 章，\S~3，第 2 号，命题 5 的推论 2）。
\end{proof}

\begin{corollary}[17.2.7]
\label{0.17.2.7-fr}
在 (17.2.6) 的假设下，设 $x$ 是一个 $M$-正则元，属于环 $A$，并且位于
$A$ 的极大理想；则有
\begin{equation}
\label{0.17.2.7.1-fr}
\tag{17.2.7.1}
  \operatorname{dim.proj}(M/xM)=\operatorname{dim.proj}(M)+1.
\end{equation}
\end{corollary}

\begin{proof}
事实上，正合列
$0\longrightarrow M\xrightarrow{x}M\longrightarrow M/xM\longrightarrow0$
由乘以 $x$ 的映射在 $M$ 上定义，并给出如下 Tor 正合列：
\[
  \operatorname{Tor}^A_i(M,k)\xrightarrow{x}
  \operatorname{Tor}^A_i(M,k)\longrightarrow
  \operatorname{Tor}^A_i(M/xM,k)\longrightarrow
  \operatorname{Tor}^A_{i-1}(M,k)\xrightarrow{x}
  \operatorname{Tor}^A_{i-1}(M,k)
\]
对每个 $i\geq1$ 都成立。现在，对每个 $A$-模 $N$，乘以 $x$ 的自同态在
$M\otimes_A N$ 上既可由乘以 $x$ 的自同态在 $M$ 上的作用诱导，也可由乘以 $x$
的自同态在 $N$ 上的作用诱导；由于按定义，乘以 $x$ 的自同态在 $k$ 上为零，因而立即得出：
对每个 $i$，
同态
$\operatorname{Tor}^A_i(M,k)\xrightarrow{x}\operatorname{Tor}^A_i(M,k)$
为\emph{零}；换言之，有正合列
\[
  0\longrightarrow\operatorname{Tor}^A_i(M,k)\longrightarrow
  \operatorname{Tor}^A_i(M/xM,k)\longrightarrow
  \operatorname{Tor}^A_{i-1}(M,k)\longrightarrow0.
\]
若 $\operatorname{dim.proj}(M)=n$，则由 (17.2.6)，有
$\operatorname{Tor}^A_n(M,k)\neq0$，且
$\operatorname{Tor}^A_{n+1}(M,k)=\operatorname{Tor}^A_{n+2}(M,k)=0$。
因此，由前述结果可知
$\operatorname{Tor}^A_{n+1}(M/xM,k)\neq0$，且
$\operatorname{Tor}^A_{n+2}(M/xM,k)=0$；再由 (17.2.6)，得到
$\operatorname{dim.proj}(M/xM)=n+1$。
\end{proof}

\begin{proposition}[17.2.8]
\label{0.17.2.8-fr}
\textnormal{（M.~Auslander）} 设 $A$ 是环。下列条件等价：
\begin{enumerate}
  \item[\rm a)] 每个 $A$-模的投射维数都 $\leq n$。
  \item[\rm a')] 每个有限型 $A$-模的投射维数都 $\leq n$。
  \item[\rm b)] 每个 $A$-模的内射维数都 $\leq n$。
  \item[\rm c)] 对于每对 $A$-模 $M$、$N$，都有
  $\operatorname{Ext}^{n+1}_A(M,N)=0$。
  \item[\rm c')] 对于每对 $A$-模 $M$、$N$，若 $M$ 是有限型的（甚至只须是
  循环模），则有 $\operatorname{Ext}^{n+1}_A(M,N)=0$。
\end{enumerate}
\end{proposition}

\begin{proof}
\oldpage[0\textsubscript{IV-1}]{45}
这立即由 (17.2.1) 和 (17.2.3) 得出。
\end{proof}

使 (17.2.8) 中各等价条件成立的最小数 $n$（整数或 $+\infty$）称为 $A$ 的
\emph{整体上同调维数}（或简称\emph{上同调维数}），记作
$\operatorname{dim.coh}(A)$。

\begin{proposition}[17.2.9]
\label{0.17.2.9-fr}
设 $A$ 是 Noether 环。下列条件等价：
\begin{enumerate}
  \item[\rm a)] $\operatorname{dim.coh}(A)\leq n$。
  \item[\rm b)] 每个有限型 $A$-模的内射维数都 $\leq n$。
  \item[\rm c)] 对于每对有限型 $A$-模 $M$、$N$，都有
  $\operatorname{Ext}^{n+1}_A(M,N)=0$。
\end{enumerate}
\end{proposition}

\begin{proof}
这立即由定义和 (17.2.4) 得出。
\end{proof}

\begin{corollary}[17.2.10]
\label{0.17.2.10-fr}
若 $A$ 是 Noether 环，则有
\begin{equation}
\label{0.17.2.10.1-fr}
\tag{17.2.10.1}
  \operatorname{dim.coh}(A)
  =\sup_{\mathfrak{m}}
  \bigl(\operatorname{dim.coh}(A_{\mathfrak{m}})\bigr)
\end{equation}
其中 $\mathfrak{m}$ 遍历 $A$ 的谱（或者 $A$ 的极大理想集）。
\end{corollary}

\begin{proof}
这由 (17.2.9) 和 (17.2.5) 得出。
\end{proof}

\begin{proposition}[17.2.11]
\label{0.17.2.11-fr}
设 $A$ 是 Noether 局部环，$k$ 是它的剩余域。要使
$\operatorname{dim.coh}(A)\leq n$，必须有
$\operatorname{Tor}^A_i(k,k)=0$（$i>n$）；而只须有
$\operatorname{Tor}^A_{n+1}(k,k)=0$ 即已充分。
\end{proposition}

\begin{proof}
考虑到 (17.2.6)，只须证明关系
$\operatorname{dim.coh}(A)\leq n$ 与
$\operatorname{dim.proj}_A(k)\leq n$ 等价。由定义，前者显然推出后者。反过来，若
$\operatorname{dim.proj}_A(k)\leq n$，则由 (17.2.2, (ii))，有
$\operatorname{Tor}^A_i(M,k)=0$；这里 $i>n$，且所取的 $A$-模 $M$ 任意。因而
$\operatorname{dim.proj}(M)\leq n$，再由 (17.2.8) 即得命题。
\end{proof}

\begin{corollary}[17.2.12]
\label{0.17.2.12-fr}
设 $A$ 是 Noether 环。要使
$\operatorname{dim.coh}(A)\leq n$，必须对每个极大理想
$\mathfrak{m}$（在 $A$ 中）都有
$\operatorname{Tor}^{A_{\mathfrak{m}}}_i(A/\mathfrak{m},A/\mathfrak{m})=0$
（$i>n$）；而只须这些关系在 $i=n+1$ 时成立即已充分。
\end{corollary}

\begin{proof}
这立即由 (17.2.11) 和 (17.2.10) 得出。
\end{proof}

\begin{proposition}[17.2.13]
\label{0.17.2.13-fr}
设 $A$、$B$ 是两个 Noether 局部环，
$\varphi:A\longrightarrow B$ 是使 $B$ 成为平坦 $A$-模的局部同态。则有
\begin{equation}
\label{0.17.2.13.1-fr}
\tag{17.2.13.1}
  \operatorname{dim.coh}(A)\leq\operatorname{dim.coh}(B).
\end{equation}
\end{proposition}

\begin{proof}
设 $\operatorname{dim.coh}(B)=n$ 为有限数；只须证明：对每对
$(M,N)$（均为有限型 $A$-模），都有
$\operatorname{Tor}^A_i(M,N)=0$（$i>n$）(17.2.11)。由于
$B$ 是忠实平坦 $A$-模（$0_{\rm I}$, 6.6.2），由（$0_{\rm I}$, 6.4.1），这等价于证明
\begin{equation}
\label{0.17.2.13.2-fr}
\tag{17.2.13.2}
  \operatorname{Tor}^A_i(M,N)\otimes_A B=0.
\end{equation}
现在，若 $L_\bullet=(L_j)$ 是 $M$ 的、由自由 $A$-模组成的右消解，则由 $B$ 是
平坦 $A$-模可知，$L_\bullet\otimes_A B=(L_j\otimes_A B)$ 是
$B$-模 $M\otimes_A B$ 的、由自由 $B$-模组成的右消解；此外，
$(L_j\otimes_A B)\otimes_B(N\otimes_A B)
=(L_j\otimes_A N)\otimes_A B$，由此立即推出 (17.2.13.2) 的左端等于
\oldpage[0\textsubscript{IV-1}]{46}
$\operatorname{Tor}^B_i(M\otimes_A B,N\otimes_A B)$；关于 $B$ 的假设蕴含这个
$B$-模在 $i>n$ 时为零（17.2.2, (ii)），结论随即得证。
\end{proof}

\begin{env}[17.2.14]
\label{0.17.2.14-fr}
设 $(X,\mathscr{O}_X)$ 是赋环空间，$\mathscr{F}$ 是
$\mathscr{O}_X$-模。将 $\mathscr{F}$ 的\emph{逐点投射维数}（相应地，
\emph{逐点内射维数}）记作
$\operatorname{dim.proj}(\mathscr{F})$（相应地，
$\operatorname{dim.inj}(\mathscr{F})$）；其定义为数
$\sup_{x\in X}(\operatorname{dim.proj}(\mathscr{F}_x))$（相应地，数
$\sup_{x\in X}(\operatorname{dim.inj}(\mathscr{F}_x))$）。将 $X$ 的
\emph{逐点上同调维数}记作 $\operatorname{dim.coh}(X)$；其定义为数
$\sup_{x\in X}(\operatorname{dim.coh}(\mathscr{O}_x))$。由 (17.2.5) 和
(17.2.10) 可知，当 $A$ 是 Noether 环，$M$ 是有限型 $A$-模，且
$X=\operatorname{Spec}(A)$ 时，有
$\operatorname{dim.proj}(\widetilde{M})=\operatorname{dim.proj}(M)$、
$\operatorname{dim.inj}(\widetilde{M})=\operatorname{dim.inj}(M)$ 以及
$\operatorname{dim.coh}(X)=\operatorname{dim.coh}(A)$。称 $\mathscr{F}$
\emph{在点 $x\in X$ 处的投射维数}（相应地，\emph{内射维数}）为
$\mathscr{F}_x$ 的投射维数（相应地，内射维数）；称 $X$
\emph{在点 $x$ 处的上同调维数}为 $\mathscr{O}_x$ 的上同调维数。
\end{env}

\begin{proposition}[17.2.15]
\label{0.17.2.15-fr}
设 $X$、$Y$ 是两个以 Noether 局部环为茎的赋环空间，
$f:X\longrightarrow Y$ 是平坦态射。若
$\operatorname{dim.coh}(X)\leq n$，则 $Y$ 在这些点处的上同调维数都
$\leq n$；这里的点取遍 $f(X)$。
\end{proposition}

\begin{proof}
这立即由 (17.2.13) 得出。
\end{proof}

\subsection{正则环的上同调理论}
\label{subsection:0.17.3-fr}

\begin{theorem}[17.3.1]
\label{0.17.3.1-fr}
\textnormal{(Hilbert--Serre)} 设 $A$ 是 Noether 局部环。$A$ 具有有限上同调维数，
当且仅当 $A$ 是正则环；在此情形下，有
\begin{equation}
\label{0.17.3.1.1-fr}
\tag{17.3.1.1}
  \operatorname{dim.coh}(A)=\dim(A).
\end{equation}
\end{theorem}

\begin{proof}
假设 $A$ 是正则环；设 $\mathfrak{m}$ 是它的极大理想，
$\mathbf{x}=(x_i)_{1\leq i\leq n}$ 是 $A$ 的一个正则参数系
(17.1.6)；考虑外代数复形 $K_\bullet(\mathbf{x})$（\textbf{III}, 1.1.1）。
它由自由 $A$-模组成，并且 $K_i(\mathbf{x})=0$（$i>n$）。由于
$(x_i)$ 是正则序列，故
$H_i(K_\bullet(\mathbf{x}))=0$（$i>0$）（\textbf{III}, 1.1.4 和 1.1.3.3），且
$H_0(K_\bullet(\mathbf{x}))=A/(x_1A+\cdots+x_nA)=A/\mathfrak{m}=k$；
因此，各 $K_i(\mathbf{x})$ 构成 $k$ 的一个长度为 $n$ 的\emph{自由消解}。
另一方面，由于各 $x_i$ 都属于 $\mathfrak{m}$，在复形
$K_\bullet(\mathbf{x},k)=K_\bullet(\mathbf{x})\otimes_A k$ 中，边界算子在
所有维数上均为零，因而按照定义，
$\operatorname{Tor}^A_i(k,k)=\bigwedge^i(k^n)$；于是等式 (17.3.1.1)
立即由 (17.2.11) 得出（这一结果本质上就是 Hilbert 的“合冲定理”）。

现在证明：若 $A$ 是 Noether 局部环，$\mathfrak{m}$ 是它的极大理想，并且
$\operatorname{dim.proj}_A(\mathfrak{m})$ \emph{有限}，则 $A$ 是正则环；这将完成
(17.3.1) 的证明。我们对
$n=\operatorname{rg}_k(\mathfrak{m}/\mathfrak{m}^2)$ 作归纳。当 $n=0$ 时，
$\mathfrak{m}=0$，断言显然成立。
\end{proof}

\begin{lemma}[17.3.1.2]
\label{0.17.3.1.2-fr}
\textnormal{(Nagata)} 若 $\mathfrak{m}-\mathfrak{m}^2$ 的每个元素都是 $A$
中的零因子，则存在 $c\neq0$，它属于 $A$，并使得 $c\mathfrak{m}=0$（换言之，
$\mathfrak{m}\in\operatorname{Ass}(A)$）。
\end{lemma}

\begin{proof}
只需考虑 $\mathfrak{m}\neq0$ 的情形，因而
$\mathfrak{m}\neq\mathfrak{m}^2$。上述假设意味着
$\mathfrak{m}-\mathfrak{m}^2$ 包含于各理想 $\mathfrak{p}_i$ 的并；这里各理想属于
$\operatorname{Ass}(A)$（Bourbaki，\emph{Alg.~comm.}，第 IV 章，\S~1，
第 1 号，命题 2 的推论 3）；所以 $\mathfrak{m}$ 包含于 $\mathfrak{m}^2$
\oldpage[0\textsubscript{IV-1}]{47}
与各 $\mathfrak{p}_i$ 的并，从而包含于其中某个 $\mathfrak{p}_i$，这是因为
$\mathfrak{m}\not\subset\mathfrak{m}^2$（Bourbaki，\emph{Alg.~comm.}，
第 II 章，\S~1，第 1 号，命题 2）；由于 $\mathfrak{m}$ 是极大理想，这便证明了引理。
\end{proof}

\begin{lemma}[17.3.1.3]
\label{0.17.3.1.3-fr}
若 $a\in\mathfrak{m}-\mathfrak{m}^2$，则 $\mathfrak{m}/Aa$ 同构于
$\mathfrak{m}/a\mathfrak{m}$ 的一个直和项。
\end{lemma}

\begin{proof}
存在 $\mathfrak{m}$ 的一个包含 $a$ 的极小生成元系（这些生成元在
$\mathfrak{m}/\mathfrak{m}^2$ 中的像构成这个 $k$-向量空间的一组基）；令
$\mathfrak{b}$ 为该生成元系中除 $a$ 以外各元素所生成的理想。由于关系
$xa\in\mathfrak{b}$ 蕴含 $x\in\mathfrak{m}$（只须考察 $xa$ 在
$\mathfrak{m}/\mathfrak{m}^2$ 中的像），故有
$\mathfrak{b}\cap Aa\subset a\mathfrak{m}$；于是得到同态
$\mathfrak{b}/(\mathfrak{b}\cap Aa)\longrightarrow
\mathfrak{m}/a\mathfrak{m}$；它由单射
$\mathfrak{b}\longrightarrow\mathfrak{m}$ 取商而得，并且是\emph{单射}。此外，
$\mathfrak{b}+Aa=\mathfrak{m}$，而复合同态
\[
  \mathfrak{m}/Aa=(\mathfrak{b}+Aa)/Aa
  \xrightarrow{\sim}\mathfrak{b}/(\mathfrak{b}\cap Aa)
  \longrightarrow\mathfrak{m}/a\mathfrak{m}
  \longrightarrow\mathfrak{m}/Aa
\]
是恒等同态；引理得证。
\end{proof}

\begin{lemma}[17.3.1.4]
\label{0.17.3.1.4-fr}
设 $A$ 是 Noether 局部环，$\mathfrak{m}$ 是它的极大理想，$E$ 是有限型
$A$-模，且具有有限投射维数。若 $a\in\mathfrak{m}$ 是 $A$-正则且
$E$-正则的，则 $E/aE$ 是具有有限投射维数的 $(A/aA)$-模，并且其投射维数
不超过 $\operatorname{dim.proj}_A(E)$。
\end{lemma}

\begin{proof}
对 $h=\operatorname{dim.proj}_A(E)$ 作归纳；当 $h=0$ 时，结论显然，因为此时
$E$ 是投射 $A$-模，因而 $E/aE$ 是投射 $(A/aA)$-模。存在正合列
\[
  0\longrightarrow N\longrightarrow L\longrightarrow E\longrightarrow0
\]
其中 $L$ 是自由模，且 $\operatorname{dim.proj}_A(N)=h-1$（17.2.2, (iii)），
$N$ 是有限型模。此外，序列
\[
  0\longrightarrow N/aN\longrightarrow L/aL\longrightarrow E/aE
  \longrightarrow0
\]
是正合的 (15.1.18)。由于 $a$ 是 $A$-正则的，它也是 $N$-正则的，这是因为
$L$ 是自由模；归纳假设蕴含 $N/aN$ 是 $(A/aA)$-模，其投射维数
$\leq h-1$；又因 $L/aL$ 是自由 $(A/aA)$-模，故 $E/aE$ 是
$(A/aA)$-模，其投射维数 $\leq h$。
\end{proof}

现在考察两种情形：

\medskip\noindent\textbf{I.} --- 首先假设 $\mathfrak{m}-\mathfrak{m}^2$
的每个元素都是 $A$ 中的零因子；此时由 (17.3.1.2)，存在 $c\neq0$，它属于 $A$，
使得 $c\mathfrak{m}=0$。我们来证明 $\mathfrak{m}=0$。若非如此，先注意
$\mathfrak{m}$ 不可能是投射 $A$-模，因为否则它将是自由模（Bourbaki，
\emph{Alg.~comm.}，第 II 章，\S~3，第 2 号，命题 5 的推论 2），这与关系
$c\mathfrak{m}=0$ 矛盾。因此将有
$n=\operatorname{dim.coh}(A)\geq1$。由于
$\mathfrak{m}\in\operatorname{Ass}(A)$，存在 $A$-同态的正合列
\[
  0\longrightarrow k\longrightarrow A\longrightarrow E\longrightarrow0.
\]
但这是不可能的，因为由关系 $n\geq1$，Tor 的长正合列将给出正合列
$0\longrightarrow\operatorname{Tor}^A_{n+1}(E,k)\longrightarrow
\operatorname{Tor}^A_n(k,k)\longrightarrow0$；然而
$\operatorname{Tor}^A_{n+1}(E,k)=0$（17.2.2, (ii)），且
$\operatorname{Tor}^A_n(k,k)\neq0$（17.2.11），于是得到矛盾。

\medskip\noindent\textbf{II.} --- 因此，只须考虑存在
$a\in\mathfrak{m}-\mathfrak{m}^2$ 是 $A$-正则元、从而也是
$\mathfrak{m}$-正则元的情形。考虑环 $A'=A/aA$ 及其极大理想
\oldpage[0\textsubscript{IV-1}]{48}
$\mathfrak{m}'=\mathfrak{m}/aA$；显然
$\operatorname{rg}_k(\mathfrak{m}'/\mathfrak{m}'^2)=n-1$。由 (17.3.1.4)，
$\mathfrak{m}/a\mathfrak{m}$ 是具有有限投射维数的 $A'$-模；因而它的直和项
$\mathfrak{m}'=\mathfrak{m}/aA$ 也具有这个性质（17.3.1.3 和 17.2.1）。归纳假设
于是蕴含 $A'=A/aA$ 是正则环，进而由 (17.1.8) 可知 $A$ 是正则环。

\begin{corollary}[17.3.2]
\label{0.17.3.2-fr}
若 $A$ 是正则局部环，则 $A_{\mathfrak{p}}$ 是正则环，其中
$\mathfrak{p}$ 可为 $A$ 的任一素理想。
\end{corollary}

\begin{proof}
事实上，(17.2.10) 已经说明
$\operatorname{dim.coh}(A_{\mathfrak{p}})\leq
\operatorname{dim.coh}(A)$，故结论立即由 (17.3.1) 得出。
\end{proof}

\begin{proposition}[17.3.3]
\label{0.17.3.3-fr}
设 $A$、$B$ 是两个 Noether 局部环，
$\varphi:A\longrightarrow B$ 是局部同态，$\mathfrak{m}$（相应地，
$\mathfrak{n}$）是 $A$（相应地，$B$）的极大理想，
$k=A/\mathfrak{m}$（相应地，$k'=B/\mathfrak{n}$）是 $A$（相应地，
$B$）的剩余域。于是有 $k'$-向量空间的典范同态
\begin{equation}
\label{0.17.3.3.1-fr}
\tag{17.3.3.1}
  \psi:(\mathfrak{m}/\mathfrak{m}^2)\otimes_k k'
  \longrightarrow\mathfrak{n}/\mathfrak{n}^2.
\end{equation}
\begin{enumerate}
  \item[\rm (i)] 若 $B$ 是正则环，且 $B$ 是平坦 $A$-模，则 $A$ 是正则环。
  \item[\rm (ii)] 下列条件等价：
  \begin{enumerate}
    \item[\rm a)] $B$ 是正则环，且同态 $\psi$ (17.3.3.1) 是单射。
    \item[\rm b)] $A$ 和 $B$ 都是正则环，并且若 $(x_i)_{1\leq i\leq m}$
    是 $A$ 的一个正则参数系，则各 $\varphi(x_i)$ 构成 $B$ 的某个正则参数系
    的一部分（此时，对 $A$ 的\emph{每个}正则参数系都具有这一性质）。
    \item[\rm c)] $B$ 和 $B\otimes_A k$ 都是正则环，且 $B$ 是平坦
    $A$-模。
    \item[\rm d)] $A$ 和 $B\otimes_A k$ 都是正则环，且 $B$ 是平坦
    $A$-模。
    \item[\rm e)] $A$ 和 $B\otimes_A k$ 都是正则环，并且有
    \begin{equation}
    \label{0.17.3.3.2-fr}
    \tag{17.3.3.2}
      \dim(B)=\dim(A)+\dim(B\otimes_A k).
    \end{equation}
  \end{enumerate}
\end{enumerate}
\end{proposition}

\begin{proof}
\textnormal{(i)} 由 (17.2.13)，有
$\operatorname{dim.coh}(A)\leq\operatorname{dim.coh}(B)$，因而只须应用
(17.3.1)。

\textnormal{(ii)} 假设 $A$ 是正则环，且 $(x_i)$ 是 $A$ 的正则参数系。由
(15.1.21)，说 $B$ 是\emph{平坦} $A$-模，等价于说序列
$y_i=\varphi(x_i)$ 是 $B$-正则的（因为 $B\otimes_A k$ 是平坦
$k$-模）。另一方面，由于 $\dim(A)=m$，并且序列 $(y_i)$ 生成理想
$\mathfrak{m}B$，关系 (17.3.3.2) 也可写成
$\dim(B/\mathfrak{m}B)=\dim(B)-m$；这一关系等价于说序列
$(y_i)_{1\leq i\leq m}$ 构成 $B$ 的某个参数系的一部分 (16.3.7)。因此可见，
条件 b)、d) 和 e) 分别等价于下列条件：
\begin{enumerate}
  \item[\rm b')] $A$ 和 $B$ 都是正则环，且 $(y_i)_{1\leq i\leq m}$
  构成 $B$ 的某个正则参数系的一部分。
  \item[\rm d')] $A$ 是正则环，$B/(\sum_{i=1}^m y_iB)$ 是正则环，且序列
  $(y_i)$ 是 $B$-正则的。
  \item[\rm e')] $A$ 是正则环，$B/(\sum_{i=1}^m y_iB)$ 是正则环，且序列
  $(y_i)$ 构成 $B$ 的某个参数系的一部分。
\end{enumerate}
由 (17.1.7)，b') 和 e') 等价；又因 d') 蕴含 e') (16.4.1)，且 b') 蕴含
d') (17.1.7)，所以 d') 与它们等价。条件 b) 和 d) 同时成立
\oldpage[0\textsubscript{IV-1}]{49}
显然蕴含 c)；而由 (i)，c) 蕴含 d)；所以已经证明 b)、c)、d) 和 e) 等价。
此外，b) 显然蕴含 a)：此时各 $x_i$ 在
$\mathfrak{m}/\mathfrak{m}^2$ 中的类构成这个 $k$-向量空间的一组基，而各
$y_i$ 在 $\mathfrak{n}/\mathfrak{n}^2$ 中的类构成这个 $k'$-向量空间的一个
线性无关组。因此，只须证明 a) 蕴含 e)。令
$V=\mathfrak{m}/\mathfrak{m}^2$、$W=\mathfrak{n}/\mathfrak{n}^2$；立即有
$\psi(V\otimes_k k')=(\mathfrak{n}^2+\mathfrak{m}B)/\mathfrak{n}^2$。
一方面，由 (16.3.9)，有
\begin{equation}
\label{0.17.3.3.3-fr}
\tag{17.3.3.3}
  \dim(B)\leq\dim(A)+\dim(B\otimes_A k)~;
\end{equation}
其次，由 (16.2.6)，有
\begin{equation}
\label{0.17.3.3.4-fr}
\tag{17.3.3.4}
  \dim(A)\leq\operatorname{rg}_{k'}(V\otimes_k k').
\end{equation}
最后，在局部环 $B\otimes_A k=B/\mathfrak{m}B$ 中，
$\mathfrak{n}'=\mathfrak{n}/\mathfrak{m}B$ 是极大理想，$k'$ 是剩余域，且
$\mathfrak{n}'/\mathfrak{n}'^2$ 同构于
$\mathfrak{n}/(\mathfrak{n}^2+\mathfrak{m}B)
=W/\psi(V\otimes_k k')$；因此，由 (16.2.6)，有
\begin{equation}
\label{0.17.3.3.5-fr}
\tag{17.3.3.5}
  \dim(B\otimes_A k)
  \leq\operatorname{rg}_{k'}W
  -\operatorname{rg}_{k'}\psi(V\otimes_k k').
\end{equation}
又因假设 $B$ 是正则环，故
$\dim(B)=\operatorname{rg}_{k'}W$ (17.1.1)；于是由 (17.3.3.3)、
(17.3.3.4) 和 (17.3.3.5) 得到
\[
  \operatorname{rg}_{k'}W
  \leq\dim(A)+\dim(B\otimes_A k)
  \leq\operatorname{rg}_{k'}W
  +\operatorname{rg}_{k'}(V\otimes_k k')
  -\operatorname{rg}_{k'}\psi(V\otimes_k k')
\]
说这个不等式的两个端项相等，正是说 $\psi$ \emph{是单射}。因此条件 a)
必然蕴含：在关系 (17.3.3.3)、(17.3.3.4) 和 (17.3.3.5) 中，两端都
\emph{相等}；而 (17.3.3.4)（相应地，(17.3.3.5)）中的等号意味着 $A$
（相应地，$B\otimes_A k$）是\emph{正则环} (17.1.1)。这样便证明了 a)
蕴含 e)。
\end{proof}

\begin{proposition}[17.3.4]
\label{0.17.3.4-fr}
设 $A$ 是维数为 $n$ 的正则局部环，$\mathfrak{m}$ 是它的极大理想。对于每个
非零有限型 $A$-模 $M$，有
\begin{equation}
\label{0.17.3.4.1-fr}
\tag{17.3.4.1}
  \operatorname{prof}(M)+\operatorname{dim.proj}(M)=n.
\end{equation}
\end{proposition}

\begin{proof}
对 $m=\operatorname{prof}(M)$ 作归纳。若 $m=0$，则由 (16.4.6, (i))，存在
一个子-$A$-模 $N$，它包含于 $M$，并同构于 $k=A/\mathfrak{m}$；对正合列
$0\longrightarrow N\longrightarrow M\longrightarrow M/N\longrightarrow0$
应用 Tor 的长正合列，得到正合列
\[
  \operatorname{Tor}^A_{n+1}(M/N,k)\longrightarrow
  \operatorname{Tor}^A_n(N,k)\longrightarrow
  \operatorname{Tor}^A_n(M,k),
\]
而 $A$ 为正则环的假设蕴含
$\operatorname{Tor}^A_{n+1}(M/N,k)=0$（(17.3.1.1) 和 (17.2.6)），故
$\operatorname{Tor}^A_n(k,k)$ 同构于某个子-$A$-模；这个子模嵌入
$\operatorname{Tor}^A_n(M,k)$。再次应用 (17.3.1.1) 和
(17.2.6)，可知 $\operatorname{Tor}^A_n(M,k)\neq0$，从而
$\operatorname{dim.proj}(M)\geq n$ (17.2.6)；但由 (17.3.1.1)，
$\operatorname{dim.proj}(M)\leq n$，所以确有
$\operatorname{dim.proj}(M)=n$。

现在假设 $m>0$，并设 $x$ 是一个 $M$-正则元，且它属于 $\mathfrak{m}$；由
(16.4.6, (ii))，此时有 $\operatorname{prof}(M/xM)=m-1$；另一方面，由
(17.2.7)，有
$\operatorname{dim.proj}(M/xM)=\operatorname{dim.proj}(M)+1$；归纳假设立即
给出关系 (17.3.4.1)。
\SourceCorrectionNote{法文印本在此深度降一处误引 (16.4.6, (i))；应为 (ii)。}
\end{proof}

\oldpage[0\textsubscript{IV-1}]{50}

\begin{corollary}[17.3.5]
\label{0.17.3.5-fr}
\begin{enumerate}
  \item[\rm(i)] 设 $A$ 是维数为 $n$ 的正则局部环；要使有限型
  $A$-模 $M\neq0$ 为自由模，必须且只须 $M$ 是 Cohen--Macaulay
  $A$-模，且维数为 $n$。
  \item[\rm(ii)] 设 $A$ 是正则局部环，$B$ 是局部环，
  $\rho:A\to B$ 是使 $B$ 成为有限型 $A$-模的局部同态。要使 $B$ 是自由
  $A$-模，必须且只须 $B$ 是 Cohen--Macaulay 环且 $\rho$ 是单射（或者，
  等价地 (16.1.5)，$\dim(B)=\dim(A)$）。
\end{enumerate}
\end{corollary}

\begin{proof}
(i) 这由 (17.3.4.1) 以及下述事实得出：对于有限型 $A$-模，说该模是投射模，
与说它是自由模等价（Bourbaki，\emph{Alg. comm.}，第 II 章，\S~3，
第 2 号，命题 5 的推论 2）；因而有限型自由 $A$-模 $M$
由关系 $\operatorname{dim.proj}(M)=0$ (17.2.2, (i)) 刻画。

(ii) 说 $B$ 是 Cohen--Macaulay 环，等价于说 $B$ 是 Cohen--Macaulay
$A$-模 (16.5.3)；因此，只须应用 (i)，因为 $\dim(B)=\dim(A)$ (16.1.5)。
\end{proof}

\begin{env}[17.3.6]
\label{0.17.3.6-fr}
称 Noether 环 $A$ 是\emph{正则的}，如果对于每个素理想 $\mathfrak{p}$（在
$A$ 中），局部环 $A_{\mathfrak{p}}$ 都是正则环；当 $A$ 本身是
局部环时，由 (17.3.2)，这一定义与 (17.1.2) 的定义等价。由 (17.3.2)，要使
$A$ 正则，只须 $A_{\mathfrak{m}}$ 对每个极大理想 $\mathfrak{m}$（在
$A$ 中）都正则。此外，由这个定义立即可知：对于任意乘法子集
$S$（在 $A$ 中），$S^{-1}A$ 都是正则环。
\end{env}

\begin{proposition}[17.3.7]
\label{0.17.3.7-fr}
若 $A$ 是正则 Noether 环，则每个多项式环 $A[T_1,\ldots,T_n]$ 都是正则环。
\end{proposition}

\begin{proof}
显然，只须证明多项式环 $B=A[T]$ 是正则环。由于 $B$ 是自由 $A$-模，故对
每个素理想 $\mathfrak{q}$（在 $B$ 中），$B_{\mathfrak{q}}$ 是平坦
$A_{\mathfrak{p}}$-模，其中 $\mathfrak{p}=\mathfrak{q}\cap A$
（$0_{\mathrm I}$, 6.3.1）；因而由 (17.3.3)，只须证明
$B_{\mathfrak{q}}/\mathfrak{p}B_{\mathfrak{q}}$ 是正则环。由于这个环是在
$A_{\mathfrak{p}}[T]/\mathfrak{p}A_{\mathfrak{p}}[T]=k[T]$ 的某个素理想处
的局部环（这里 $k$ 是剩余域
$A_{\mathfrak{p}}/\mathfrak{p}A_{\mathfrak{p}}$），只须证明
$C=k[T]$ 是正则环；而 $C$ 是主理想整环，故 $C$ 在各素理想处的局部环都是
离散赋值环，或是一个域（对于理想 $(0)$），所以是正则环 (17.1.4)，证明完毕。
\SourceCorrectionNote{法文印本作 (17.1.10)；《勘误与补遗》第 3 表在 “0\textsubscript{IV}, 17.3.7”项明确要求改为 (17.3.3)。}
\end{proof}

\begin{corollary}[17.3.8]
\label{0.17.3.8-fr}
若 $A$ 是正则环，则每个形式幂级数环
$A[[T_1,\ldots,T_n]]$ 都是正则环。
\end{corollary}

\begin{proof}
设 $\mathfrak{J}$ 是由各 $T_i$ 在多项式环 $A[T_1,\ldots,T_n]$ 中生成的理想；
由 (17.3.7)，后者是正则环；又因为 $A[[T_1,\ldots,T_n]]$ 是
$A[T_1,\ldots,T_n]$ 关于 $\mathfrak{J}$-预进制拓扑的完备化（Bourbaki，
\emph{Alg. comm.}，第 III 章，\S~2，第 12 号，
\emph{例 1}），结论由下述引理得出。
\SourceCorrectionNote{法文印本在此止于语法不完整的“résulte du”；补为“下述引理”。}
\end{proof}

\begin{lemma}[17.3.8.1]
\label{0.17.3.8.1-fr}
设 $A$ 是正则环，$\mathfrak{J}$ 是 $A$ 的理想，$\widehat{A}$ 是 $A$ 关于
$\mathfrak{J}$-预进制拓扑的分离完备化。则 $\widehat{A}$ 是正则环。
\end{lemma}

\begin{proof}
事实上，由 (17.3.6)，只须证明：对于每个极大理想 $\mathfrak{n}$（在
$\widehat{A}$ 中），局部环 $\widehat{A}_{\mathfrak{n}}$ 都是正则环。已知
（Bourbaki，\emph{Alg. comm.}，第 III 章，\S~3，第 4 号，
命题 8），$\mathfrak{n}$
\oldpage[0\textsubscript{IV-1}]{51}
具有形式 $\mathfrak{m}\widehat{A}$，其中 $\mathfrak{m}$ 是 $A$ 的、包含
$\mathfrak{J}$ 的极大理想；并且典范同态 $A\to\widehat{A}$ 给出\emph{单射}
$A_{\mathfrak{m}}\to\widehat{A}_{\mathfrak{n}}$，使得
$\mathfrak{n}\widehat{A}_{\mathfrak{n}}$-预进制拓扑（定义在
$\widehat{A}_{\mathfrak{n}}$ 上）在
$A_{\mathfrak{m}}$ 上诱导出
$\mathfrak{m}A_{\mathfrak{m}}$-预进制拓扑，并且 $A_{\mathfrak{m}}$ 在
$\widehat{A}_{\mathfrak{n}}$ 中关于
$\mathfrak{n}\widehat{A}_{\mathfrak{n}}$-预进制拓扑稠密；因此，Noether 局部环
$A_{\mathfrak{m}}$ 和 $\widehat{A}_{\mathfrak{n}}$ 的\emph{完备化}典范地
\emph{同构}。但依假设，$A_{\mathfrak{m}}$ 是正则环，故它的完备化也是正则环
(17.1.5)；而由于 $\widehat{A}_{\mathfrak{n}}$ 的完备化是正则环，由
(17.1.5) 可知 $\widehat{A}_{\mathfrak{n}}$ 也是正则环。
\end{proof}

\begin{corollary}[17.3.9]
\label{0.17.3.9-fr}
设 $A$ 是 Noether 环，并且是某个正则 Noether 环 $B$ 的商环。若 $C$ 是有限型
$A$-代数，且 $S^{-1}C$ 是 $C$ 的任一分式环，则它是某个正则环的商环。
\end{corollary}

\begin{proof}
事实上，$C$ 是多项式环 $A[T_1,\ldots,T_n]$ 的商环；由于
$A=B/\mathfrak{J}$，其中 $\mathfrak{J}$ 是 $B$ 的理想，故有
$A[T_1,\ldots,T_n]=B[T_1,\ldots,T_n]/\mathfrak{J}B[T_1,\ldots,T_n]$；由
(17.3.7)，因而可归结到 $C$ 是某个正则环 $B'$ 的商环的情形；但若 $S'$ 是
$S$ 在 $B'$ 中的逆像，则 $S'$ 是 $B'$ 的乘法子集，且 $S^{-1}C$ 是
$S'^{-1}B'$ 的商环。于是最终只须证明：当 $C$ 是\emph{正则环}时，
$S^{-1}C$ 也是正则环；这已经在 (17.3.6) 中证明。
\SourceCorrectionNote{法文印本把商环误写成未局部化的 C；这里应为局部化后的对象。}
\end{proof}

正则环之商的重要性，除了上述性质以外，还在于它们是\emph{悬链环}
(16.5.12)。

代数几何应用中的一切重要环，都是正则环的商环。

\section{代数扩张的补充}
\label{section:0.18-fr}

本节汇集若干关于环的函子性构造，这些构造将在 \S\S~19 和 20 中反复使用；
其中不含任何非平凡的结果。

\subsection{增广环的逆像}
\label{subsection:0.18.1-fr}

\begin{env}[18.1.1]
\label{0.18.1.1-fr}
给定一个环 $A$（不必交换），\emph{$A$-环}范畴的对象是偶 $(B,\rho)$，它由
一个环 $B$ 和一个环同态 $\rho:A\to B$ 组成；其态射（也称
\emph{$A$-同态}）是这样的环同态 $u:B\to C$：若
$\rho:A\to B$ 和 $\sigma:A\to C$ 分别是 $B$ 与 $C$ 的 $A$-环结构同态
（也称\emph{结构同态}），则图表
\begin{equation}
\label{0.18.1.1.1-fr}
\tag{18.1.1.1}
  \xymatrix{
    B\ar[r]^u & C\\
    & A\ar[lu]^\rho\ar[u]_\sigma
  }
\end{equation}
交换。$u$ 的\emph{核} $\mathfrak{J}$ 是一个双边理想，并且是一个
$A$-\emph{双模}（相对于 $\rho$）。

若 $f:A'\to A$ 是环同态，则对于每个 $A$-环 $(B,\rho)$，
$(B,\rho\circ f)$ 都是 $A'$-环；若 $u:B\to C$ 是一个 $A$-同态，从
$A$-环 $(B,\rho)$
\oldpage[0\textsubscript{IV-1}]{52}
到 $A$-环 $(C,\sigma)$，则它也是一个 $A'$-同态，从 $A'$-环
$(B,\rho\circ f)$ 到 $A'$-环 $(C,\sigma\circ f)$；由此定义了
从 $A$-环范畴到 $A'$-环范畴的一个典范函子。
\end{env}

\begin{env}[18.1.2]
\label{0.18.1.2-fr}
设 $B$、$E$、$F$ 是三个 $A$-环，$f:E\to B$、$g:F\to B$ 是两个
$A$-同态。回顾，$E$ 与 $F$ 在 $B$ 上（相对于 $f$ 和 $g$）的
\emph{纤维积} $E\times_BF$，是积环 $E\times F$ 的子环 $G$，由满足
$f(x)=g(y)$ 的偶 $(x,y)$ 组成；典范投影
$p_1:G\to E$、$p_2:G\to F$
分别是投影 $\operatorname{pr}_1$ 和 $\operatorname{pr}_2$ 在 $G$ 上的限制。
$A$-环 $G$ 的结构由同态
$\alpha\mapsto(\rho(\alpha),\sigma(\alpha))$ 定义（其中
$\rho:A\to E$ 和 $\sigma:A\to F$ 是结构同态）；由 (18.1.1)，这个同态确实把
$A$ 映入 $G$。$E\times_BF$ 的特征性质如下：对于任意一对 $A$-同态
$u:C\to E$、$v:C\to F$，若 $f\circ u=g\circ v$，则存在唯一的 $A$-同态
$w:C\to G$，使得 $u=p_1\circ w$、$v=p_2\circ w$。也可以说，
$E\times_BF$ 是由 $B$、$E$、$F$ 以及 $A$-同态 $f$、$g$ 构成的逆系统在
$A$-环范畴中的\emph{逆极限}（$0_{\mathrm{III}}$, 8.1.9）。
\end{env}

\begin{env}[18.1.3]
\label{0.18.1.3-fr}
设 $f$ 的核 $\mathfrak{J}$ 是 $E$ 的双边理想；显然，
$p_2:G\to F$ 的核 $\mathfrak{J}'$ 是由元素 $(x,0)$ 组成的理想，其中
$x\in\mathfrak{J}$；因而，映射 $i_1:\mathfrak{J}'\to\mathfrak{J}$，即
$p_1$ 的限制，是一个\emph{无单位元环}的同构，也是
$A$-双模的同构。同样，若 $\mathfrak{R}$ 是 $g$ 的核，则
$\mathfrak{R}'$，即 $p_1:G\to E$ 的核，是由元素 $(0,y)$ 组成的理想，其中
$y\in\mathfrak{R}$；而映射 $i_2:\mathfrak{R}'\to\mathfrak{R}$，即 $p_2$ 的
限制，是同样意义下的同构。最后，显然，$A$-同态
$q=f\circ p_1=g\circ p_2$（从 $E\times_BF$ 到 $B$）的核是
$\mathfrak{J}\times\mathfrak{R}=\mathfrak{J}'\oplus\mathfrak{R}'$，从而有交换图
\begin{equation}
\label{0.18.1.3.1-fr}
\tag{18.1.3.1}
  \xymatrix@C=2.2pc@R=1.8pc{
    0\ar[dr] & & 0\ar[d] & 0\ar[d]\\
    & \mathfrak{J}'\oplus\mathfrak{R}'\ar[r]\ar[d]\ar[dr]
      & \mathfrak{R}'\ar[r]^{i_2}_{\sim}\ar[d]^j
      & \mathfrak{R}\ar[d]^{p_2}\\
    0\ar[r] & \mathfrak{J}'\ar[r]\ar[d]^{i_1}_{\sim}
      & G\ar[r]\ar[d]^{p_1}\ar[dr]^q & F\ar[d]^g\\
    0\ar[r] & \mathfrak{J}\ar[r] & E\ar[r]_f & B
  }
\end{equation}

(18.1.2) 和 (18.1.3) 的定义与结果立即推广到任意一族
$(E_\lambda)_{\lambda\in I}$ 的 $A$-环的纤维积；这一纤维积由一族
$A$-同态 $f_\lambda:E_\lambda\to B$ 定义。我们把这些结果的表述留给读者。
\end{env}

\begin{env}[18.1.4]
\label{0.18.1.4-fr}
依照 (M, VIII) 的术语，所谓\emph{在 $B$ 上增广的 $A$-环}，是指配备一个
满射 $A$-同态 $f:E\to B$ 的 $A$-环 $E$；$f$ 称为 $E$ 的\emph{增广}，
其核 $\mathfrak{J}$ 称为\emph{增广理想}。若这个增广 $A$-环 $E$ 存在一个
$A$-同态 $s:B\to E$，它是增广
$f:E\to B$ 的\emph{右逆}（换言之，$f\circ s=1_B$），则称这个增广环是
\emph{平凡的}。此时，$A$-双模的正合列
\[
  0\longrightarrow\mathfrak{J}\longrightarrow E
  \xrightarrow{f}B\longrightarrow0
\]
\oldpage[0\textsubscript{IV-1}]{53}
是\emph{分裂的}；换言之，可以把 $A$-双模 $E$ 等同于
$B\times\mathfrak{J}$，而 $E$ 中的乘法由
\[
  (b,z)(b',z')=(bb',bz'+zb'+zz')
\]
给出，其中 $\mathfrak{J}$ 被视为一个 $B$-双模，其结构由 $s:B\to E$ 给出。
\end{env}

\begin{env}[18.1.5]
\label{0.18.1.5-fr}
沿用 (18.1.2) 的记号，假设 $A$-同态 $f$ 是\emph{满射}，换言之，$E$ 是
在 $B$ 上增广的 $A$-环。那么显然，$p_2:G\to F$ 也是\emph{满射}，换言之，
$G$ 是在 $F$ 上增广的 $A$-环；这个环称为
经由 $g:F\to B$ 所取的增广环 $E$ 的\emph{逆像}。
\end{env}

\begin{proposition}[18.1.6]
\label{0.18.1.6-fr}
在 $B$ 上增广的 $A$-环 $E$ 沿 $g:F\to B$ 的逆像
$G=E\times_BF$ 是平凡增广 $A$-环，当且仅当存在一个 $A$-同态
$u:F\to E$，使图表
\[
  \xymatrix{
    & F\ar[ld]_u\ar[d]^g\\
    E\ar[r]_f & B
  }
\]
交换。
\end{proposition}

\begin{proof}
取 $u=p_1\circ s$，其中 $s:F\to G$ 是一个 $A$-同态，并且是 $p_2$ 的右逆，
便知该条件显然是必要的。反之，若存在满足命题中条件的 $A$-同态 $u$，则
$s$ 作为 $p_2$ 的右逆的存在性，由纤维积的泛性质 (18.1.2) 应用于 $A$-同态
$u:F\to E$ 和 $1_F:F\to F$ 而得。

特别地，这一结果蕴含：若 $E$ 是平凡增广环，则它的所有逆像也都是平凡的。
\end{proof}

\begin{env}[18.1.7]
\label{0.18.1.7-fr}
重新考虑 (18.1.2) 和 (18.1.3) 中描述的情形；$\mathfrak{R}$ 显然具有一个
$G$-双模结构，它来自其 $F$-双模结构以及环同态 $p_2:G\to F$。由于 $i_2$
是双射，还得到一个单射 $\theta=j\circ i_2^{-1}:\mathfrak{R}\to G$，它是
$G$-双模同态。反过来，我们将看到，若还假设 $f$ 和 $g$ 是\emph{满射}，
换言之，$E$ 和 $F$ 都是在 $B$ 上增广的 $A$-环，那么给定这样的同态
$\theta$，便可由在 $F$ 上增广的 $A$-环 $G$ 重构在 $B$ 上增广的 $A$-环 $E$。

更确切地说，给定一个在 $B$ 上增广的 $A$-环 $F$，以及一个在 $F$ 上增广的
$A$-环 $G$，分别以 $\mathfrak{R}$ 和 $\mathfrak{J}'$ 表示其增广理想：
\begin{equation}
\label{0.18.1.7.1-fr}
\tag{18.1.7.1}
  \xymatrix@C=2pc@R=1.5pc{
    & & 0\ar[d] & & \\
    & & \mathfrak{R}\ar[ld]^\theta\ar[d] & & \\
    0\ar[r] & \mathfrak{J}'\ar[r] & G\ar[r]_h & F\ar[r]\ar[d]^g & 0\\
    & & & B\ar[d] & \\
    & & & 0 &
  }
\end{equation}

\oldpage[0\textsubscript{IV-1}]{54}
同态 $h$ 在 $F$ 的所有理想上（特别是在 $\mathfrak{R}$ 上）定义一个 $G$-双模
结构。设 $\theta:\mathfrak{R}\to G$ 是一个使图 (18.1.7.1) 交换的 $G$-双模
同态；这蕴含 $\theta$ 是\emph{单射}，并且
$\mathfrak{J}'\cap\theta(\mathfrak{R})=0$；由于 $\theta(\mathfrak{R})$ 是 $G$ 的
双边理想，并且 $h(\theta(\mathfrak{R}))=\mathfrak{R}$，复合同态
$g\circ h:G\to B$ 分解为
$G\xrightarrow{q}G/\theta(\mathfrak{R})\xrightarrow{f}B$，其中 $f$ 是满射；
此外，$\mathfrak{J}$（即 $\mathfrak{J}'$ 在 $E=G/\theta(\mathfrak{R})$ 中的像）
是 $f$ 的核（$g\circ h$ 的核为 $\mathfrak{J}'\oplus\theta(\mathfrak{R})$），而
在 $\mathfrak{J}'$ 上，典范同态 $q:G\to E$ 的限制是单射。因此可以作
\emph{纤维积} $G'=E\times_BF$；又由于两个 $A$-同态 $q:G\to E$ 和
$h:G\to F$ 满足 $f\circ q=g\circ h$，它们定义唯一的 $A$-同态
$u:G\to G'$，这是由 $G'$ 的泛性质 (18.1.2) 得到的；我们将看到 $u$ 是
\emph{双射}。只须在把 $u$ 看作 $A$-双模同态时证明这一点；这时注意到，
$u$ 与 $G$ 和 $G'$ 上的有限滤过相容，这两个滤过分别由 $G$ 与
$\mathfrak{J}'$，以及 $G'$ 与 $\mathfrak{J}''=\operatorname{Ker}(G'\to F)$ 组成；
此外，由 (18.1.3) 已知 $\operatorname{gr}_0u:F\to F$ 是恒等映射，而
$\operatorname{gr}_1u:\mathfrak{J}'\to\mathfrak{J}''$ 是双射，因而 $u$ 本身是
双射（Bourbaki，\emph{Alg. comm.}，第 III 章，\S~2，
n\textsuperscript{o}~8，定理 1 的推论 3）
\SourceCorrectionNote{法文印本及外交式转录把泛性质所给箭头的方向写反，并随之 倒置了两个滤过及次数一分次映射的方向；箭头应从原环指向纤维积，滤过应先取原环再取纤维积，次数一分次映射应从 前者的增广理想指向后者的核。}。
\end{env}

\subsection{环关于双模的扩张}
\label{subsection:0.18.2-fr}

\begin{env}[18.2.1]
\label{0.18.2.1-fr}
设 $E$ 是一个在 $B$ 上增广的 $A$-环，$f:E\to B$ 是增广同态，
$\mathfrak{J}=\operatorname{Ker}(f)$ 是增广理想。若 $\mathfrak{J}^2=0$，则
$\mathfrak{J}$ 不仅是 $E$-双模，而且还是 $B$-\emph{双模}，因为 $B$ 同构于
$E/\mathfrak{J}$。更确切地，每个 $b\in B$ 都可写成 $f(x)$，其中 $x\in E$；
若 $z,z'$ 属于 $\mathfrak{J}$，则
$(x+z)z'=xz'$（相应地，$z'(x+z)=z'x$），所以 $xz'$（相应地，$z'x$）的值
不依赖于元素 $x$ 在 $f^{-1}(b)$ 中的选取，并可写成 $bz'$（相应地，$z'b$）；
这就定义了所考察的 $B$-双模结构。反过来，若在 $\mathfrak{J}$ 上给定一个
$B$-双模结构，使得 $xz'=f(x)z'$ 且 $z'x=z'f(x)$，其中 $x\in E$、
$z'\in\mathfrak{J}$，则显然 $\mathfrak{J}^2=0$。
\end{env}

\begin{env}[18.2.2]
\label{0.18.2.2-fr}
所谓\emph{$A$-环 $B$ 关于 $B$-双模 $L$ 的一个 $A$-扩张}，是指一个
$A$-双模同态的\emph{正合列}
\[
  0\longrightarrow L\xrightarrow{j}E\xrightarrow{f}B\longrightarrow0
\]
其中 $E$ 是 $A$-环，$f$ 是环的 $A$-同态，并且对于 $x\in E$ 和 $z\in L$ 有
\[
  j(f(x)z)=xj(z),\qquad j(zf(x))=j(z)x.
\]
由此根据 (18.2.1)，$j(L)$ 是 $E$ 中平方为零的双边理想。略去严格措辞时，
也称 $E$ 是 $B$ 关于 $L$ 的一个扩张。对于 $B$ 关于 $L$ 的两个 $A$-扩张
$E$、$E'$，若存在一个使下图交换的 $A$-环同构
$u:E\xrightarrow{\sim}E'$（也称为\emph{$A$-扩张的 $A$-等价}）：
\begin{equation}
\label{0.18.2.2.1-fr}
\tag{18.2.2.1}
  \xymatrix@C=2.1pc@R=1.2pc{
    & & E\ar[dr]\ar[dd]_u\ar@{}[dd]|{\wr} & & \\
    0\ar[r] & L\ar[ur]\ar[dr] & & B\ar[r] & 0\\
    & & E'\ar[ur] & &
  }
\end{equation}
则称这两个扩张\emph{$A$-等价}。
\end{env}

\oldpage[0\textsubscript{IV-1}]{55}

\begin{env}[18.2.3]
\label{0.18.2.3-fr}
称 $B$ 关于 $B$-双模 $L$ 的一个 $A$-扩张 $E$ 是\emph{$A$-平凡的}，
如果 $E$ 是一个在 $B$ 上平凡增广的 $A$-环 (18.1.4)。在乘积 $A$-双模
$B\times L$ 上定义一个 $A$-环结构，规定
$(x,s)(y,t)=(xy,xt+sy)$；这是立即可验的。典范映射 $j:L\to B\times L$ 与
$f:B\times L\to B$ 显然定义了 $B$ 关于 $L$ 的一个 $A$-扩张，并且它是
$A$-平凡的；典范映射 $g:B\to B\times L$ 是一个 $A$-同态，并且是 $f$ 的
右逆。称这个扩张为 $B$ 关于 $L$ 的\emph{典型平凡扩张}，记为 $D_B(L)$；
显然，$B$ 关于 $L$ 的每个 $A$-平凡 $A$-扩张都 $A$-等价于 $D_B(L)$。

注意，$A$ 关于任一 $A$-双模的每个扩张都必然是 $A$-平凡的。
\end{env}

\begin{env}[18.2.4]
\label{0.18.2.4-fr}
给定两个 $A$-扩张
$L\xrightarrow{j}E\xrightarrow{f}B$、
$L'\xrightarrow{j'}E'\xrightarrow{f'}B'$，所谓从第一个到第二个的一个
\emph{态射}，按定义是一个 $A$-双模同态三元组 $(u,v,w)$，使得下图
\begin{equation}
\label{0.18.2.4.1-fr}
\tag{18.2.4.1}
  \xymatrix{
    0\ar[r] & L\ar[r]^j\ar[d]_w & E\ar[r]^f\ar[d]_u
      & B\ar[r]\ar[d]_v & 0\\
    0\ar[r] & L'\ar[r]_{j'} & E'\ar[r]_{f'} & B'\ar[r] & 0
  }
\end{equation}
交换，其中 $u$ 和 $v$ 是环的 $A$-同态，而 $w$ 满足
$w(bz)=v(b)w(z)$ 和 $w(zb)=w(z)v(b)$，这里 $z\in L$、$b\in B$（换言之，
二元组 $(v,w)$ 构成从 $B$-双模 $L$ 到 $B'$-双模 $L'$ 的一个
\emph{双同态}）。显然，若 $(u',v',w')$ 是从 $L'\to E'\to B'$ 到一个
$A$-扩张 $L''\to E''\to B''$ 的态射，则
$(u'\circ u,v'\circ v,w'\circ w)$ 仍是一个态射；这说明上述术语是合理的。

考察图 (18.2.4.1) 中的两个交换方块，将引出 $A$-环扩张上的两种运算。
\end{env}

\begin{env}[18.2.5]
\label{0.18.2.5-fr}
首先考察 $B'$ 关于 $L'$ 的一个 $A$-扩张 $E'$：
\[
  0\longrightarrow L'\xrightarrow{j'}E'\xrightarrow{f'}B'\longrightarrow0
\]
再给定一个环的 $A$-同态 $v:B\to B'$，并令
$F=E'\times_{B'}B$ 是在 $v$ 下取增广 $A$-环 $E'$ 的\emph{逆像} (18.1.5)，
于是有交换图
\[
  \xymatrix{
    0\ar[r] & L_0\ar[r]\ar[d]^i\ar@{}[d]|{\wr} & F\ar[r]^{p_2}\ar[d]^{p_1}
      & B\ar[r]\ar[d]^v & 0\\
    0\ar[r] & L'\ar[r]_{j'} & E'\ar[r]_{f'} & B'\ar[r] & 0
  }
\]
其两行正合，$p_1$ 和 $p_2$ 是典范同态；(18.1.3) 已说明 $i$ 是双射，而由
(18.1.2) 的定义还可得 $L_0^2=0$，所以可以把 $F$ 看作 $B$ 关于 $L'$ 的
一个 $A$-扩张；称它为扩张 $E'$ \emph{沿 $v$ 的逆像}，其中借助环同态 $v$
自然地把 $L'$ 视为 $B$-双模。纤维积对
每个因子的函子性还表明：若给定 $B'$ 的两个扩张之间的一个态射
\oldpage[0\textsubscript{IV-1}]{56}
\[
  \xymatrix{
    0\ar[r] & L'_1\ar[r]\ar[d]^h & E'_1\ar[r]\ar[d]^g
      & B'\ar[r]\ar[d]^{1_{B'}} & 0\\
    0\ar[r] & L'_2\ar[r] & E'_2\ar[r] & B'\ar[r] & 0
  }
\]
则由此得到一个\emph{在 $v$ 下取逆像}的态射
\[
  \xymatrix{
    0\ar[r] & L'_1\ar[r]\ar[d]^h & E'_1\times_{B'}B\ar[r]
      \ar[d]^{g\times_{B'}1_B} & B\ar[r]\ar[d]^{1_B} & 0\\
    0\ar[r] & L'_2\ar[r] & E'_2\times_{B'}B\ar[r] & B\ar[r] & 0.
  }
\]
特别地，若 $E'_1$ 和 $E'_2$ 是 $B'$ 关于 $L'$ 的两个彼此 $A$-等价的
$A$-扩张，则它们沿 $v$ 的逆像是 $B$ 关于 $L'$ 的两个彼此 $A$-等价的
$A$-扩张。

纤维积的定义表明，当给定 $A$-扩张的一个态射 (18.2.4.1) 时，它可
\emph{经由} $E'$ 沿 $v$ 的逆像\emph{分解}；更确切地，存在唯一的 $A$-同态
$u_0:E\to F=E'\times_{B'}B$，使下图交换：
\[
  \xymatrix{
    0\ar[r] & L\ar[r]^j\ar[d]^{w_0} & E\ar[r]^f\ar[d]^{u_0}
      & B\ar[r]\ar[d]^{1_B} & 0\\
    0\ar[r] & L_0\ar[r]\ar[d]^i\ar@{}[d]|{\wr} & F\ar[r]^{p_2}\ar[d]^{p_1}
      & B\ar[r]\ar[d]^v & 0\\
    0\ar[r] & L'\ar[r]_{j'} & E'\ar[r]_{f'} & B'\ar[r] & 0
  }
\]
其中 $w_0$ 是 $u_0$ 在 $L$ 上的限制，并且
$p_1\circ u_0=u$、$i\circ w_0=w$。
\end{env}

\begin{env}[18.2.6]
\label{0.18.2.6-fr}
我们特别研究沿\emph{满射} $A$-同态取扩张逆像的情形。设 $v:B\to B'$ 是
满射 $A$-同态，$F$ 是 $B$ 关于 $B$-双模 $L$ 的一个 $A$-扩张，并设
$\mathfrak{R}$ 是 $B$ 的双边理想、即 $v$ 的核（可把它看作
$F$-双模，所借助的增广为 $F\to B$）；由 (18.1.7) 已知，每个使下图交换的
$F$-双模同态 $\theta:\mathfrak{R}\to F$
\[
  \xymatrix{
    & \mathfrak{R}\ar[ld]^\theta\ar[d]\\
    F\ar[r] & B
  }
\]
都确定 $B'=B/\mathfrak{R}$ 关于 $L$ 的一个扩张 $E'$，并且 $E'$ 沿
$v:B\to B'$ 的逆像等价于 $F$；而具有后一性质的每个 $A$-扩张 $E'$ 也都可
如此得到（精确到 $A$-等价）。此外，要使两个
同态 $\theta_1,\theta_2$（作为 $F$-双模同态，从 $\mathfrak{R}$ 到 $F$）给出
两个 $B'$ 关于 $L$ 的彼此 $A$-等价的扩张，充要条件是存在一个
$A$-等价 $u$，它是 $A$-扩张 $F$ 到自身的等价，并且
\oldpage[0\textsubscript{IV-1}]{57}
$\theta_2=u\circ\theta_1$；这立即由 (18.2.5) 以及 $F$ 到纤维积
$E'\times_{B'}B$ 的典范双射的定义 (18.1.7) 得出。
\SourceCorrectionNote{法文印本及外交式转录在 18.2.6 的分类判据中把重构所得 扩张的底环写作未带撇的 B；这里应为带撇的 B，因为所比较的是由 商环重构的两个扩张，而它们在原环上的逆像才是后续对象。}
\end{env}

\begin{env}[18.2.7]
\label{0.18.2.7-fr}
现在考察 (18.2.4.1) 的左方块。先回顾 $A$-双模范畴中\emph{推出}的概念：
给定三个 $A$-双模 $X$、$Y$、$Z$ 以及两个 $A$-同态
$f:X\to Y$、$g:X\to Z$，推出 $Y\oplus_XZ$ 是归纳系统的\emph{归纳极限}；
该系统在 $A$-双模范畴中由 $X$、$Y$、$Z$ 以及 $A$-同态 $f$、$g$ 构成
（$0_{\mathrm{III}}$，8.1.11）。这个 $A$-双模定义为乘积 $Y\times Z$ 对如下
$A$-子双模 $M$ 的商：该子双模是 $X$ 在同态 $x\mapsto(f(x),-g(x))$ 下的像。
它的特征性质是：对于每一对 $A$-双模同态 $u:Y\to T$、$v:Z\to T$，只要
$u\circ f=v\circ g$，就存在唯一的同态 $w:Y\oplus_XZ\to T$，使得
$u=w\circ j_1$ 且 $v=w\circ j_2$；这里
$j_1:Y\to Y\oplus_XZ$ 和 $j_2:Z\to Y\oplus_XZ$ 是典范映射。
\end{env}

\begin{env}[18.2.8]
\label{0.18.2.8-fr}
现在考察 $B$ 关于 $B$-双模 $L$ 的一个 $A$-扩张 $E$：
\[
  0\longrightarrow L\xrightarrow{j}E\xrightarrow{f}B\longrightarrow0
\]
另一方面，设 $w:L\to L'$ 是一个 $B$-双模同态。令 $H$ 为 $A$-双模的推出
$E\oplus_LL'$；我们来说明如何在这个 $A$-双模上赋予 $A$-环结构，并定义一个
$A$-扩张
\[
  0\longrightarrow L'\longrightarrow H\longrightarrow B\longrightarrow0.
\]

为此注意，$L'$ 带有一个 $E$-\emph{双模}结构，所借助的是增广同态 $E\to B$；所以
可以构造典型平凡扩张 $G=D_E(L')$ (18.2.3)。再考察映射
$\theta:z\mapsto(j(z),-w(z))$，它从 $L$ 到 $G$；这是一个 $G$-双模同态
（这里把 $L$ 看作 $G$-双模，所借助的是典范同态 $\rho:G\to E$）。事实上，对于
$(x,z')\in G$ 和 $z\in L$，有
$(j(z),-w(z))(x,z')=(j(z)x,j(z)z'-w(z)f(x))$；但
$j(z)z'=f(j(z))z'=0$，这是由 $E$-双模在 $L'$ 上的结构定义；而
$w(z)f(x)=w(zf(x))$，这是由 $B$-双模在 $L'$ 上的结构定义；同样可验证
$\theta$ 是左 $G$-模同态。于是可以把
(18.1.7) 的结果应用于交换图
\[
  \xymatrix@C=2pc@R=1.4pc{
    & & 0\ar[d] & & \\
    & & L\ar[ld]^\theta\ar[d]^j & & \\
    0\ar[r] & L'\ar[r] & G\ar[r]_p & E\ar[r]\ar[d]^f & 0\\
    & & & B\ar[d] & \\
    & & & 0 &
  }
\]
由于按定义 $H=G/\theta(L)$，我们的断言立即由 (18.1.7) 得出。

\oldpage[0\textsubscript{IV-1}]{58}

称 $B$ 关于 $L'$ 的 $A$-扩张 $H$ 是\emph{由 $E$ 经同态
$w:L\to L'$ 推得的}。推出对每个求和项的函子性还表明：若给定扩张的一个态射
\[
  \xymatrix{
    0\ar[r] & L\ar[r]\ar[d]^{1_L}\ar@{}[d]|{\wr}
      & E_1\ar[r]\ar[d]^g & B_1\ar[r]\ar[d]^h & 0\\
    0\ar[r] & L\ar[r] & E_2\ar[r] & B_2\ar[r] & 0
  }
\]
则由此典范地得到推得扩张的一个态射
\[
  \xymatrix{
    0\ar[r] & L'\ar[r]\ar[d]^{1_{L'}}\ar@{}[d]|{\wr}
      & E_1\oplus_LL'\ar[r]\ar[d]^{g\oplus_L1_{L'}}
      & B_1\ar[r]\ar[d]^h & 0\\
    0\ar[r] & L'\ar[r] & E_2\oplus_LL'\ar[r] & B_2\ar[r] & 0.
  }
\]
特别地，若 $E_1$ 和 $E_2$ 是 $B$ 关于 $L$ 的两个彼此 $A$-等价的
$A$-扩张，则由它们经 $w$ 推得的 $B$ 关于 $L'$ 的扩张彼此 $A$-等价。

当给定 $A$-扩张的一个态射 (18.2.4.1) 时，它可\emph{经由}如下
$A$-扩张 $H$ \emph{分解}：该扩张是 $B$ 关于 $L'$ 的扩张，并由 $E$ 经同态
$w:L\to L'$ 推得（这里把 $L'$ 看作 $B$-双模，所借助的是同态
$v:B\to B'$）。事实上，推出的定义表明，存在唯一的 $A$-同态
$u_0:H\to E'$，它是 $A$-双模同态，
使下图交换：
\[
  \xymatrix{
    0\ar[r] & L\ar[r]^j\ar[d]^{w_0} & E\ar[r]^f\ar[d]^{j_1}
      & B\ar[r]\ar[d]^{1_B}\ar@{}[d]|{\wr} & 0\\
    0\ar[r] & L'_0\ar[r]^{j_2}\ar[d]^{j_0}\ar@{}[d]|{\wr}
      & H\ar[r]^{f_0}\ar[d]^{u_0} & B\ar[r]\ar[d]^v & 0\\
    0\ar[r] & L'\ar[r]_{j'} & E'\ar[r]_{f'} & B'\ar[r] & 0
  }
\]
其中 $u=u_0\circ j_1$、$w=j_0\circ w_0$，而 $j_1$ 和 $j_2$ 是典范同态；
立即可验 $u_0$ 也是环同态。

最后记下与平凡扩张有关的函子性质：
\end{env}

\begin{proposition}[18.2.9]
\label{0.18.2.9-fr}
设 $B$、$B'$ 是两个 $A$-环，$L$ 是 $B$-双模，$L'$ 是 $B'$-双模，
$v:B\to B'$ 是环的 $A$-同态，$w:L\to L'$ 是 $A$-双模同态，并且
$(v,w)$ 是双模的双同态。则存在唯一的环 $A$-同态
$u:D_B(L)\to D_{B'}(L')$，使下列图交换：
\[
  \xymatrix@C=2.2pc@R=1.6pc{
    D_B(L)\ar[r]^u & D_{B'}(L') & & D_B(L)\ar[r]^u & D_{B'}(L')\\
    B\ar[u]\ar[r]_v & B'\ar[u] & & L\ar[u]\ar[r]_w & L'\ar[u]
  }
\]
其中竖直箭头是典范单射。
\end{proposition}

\oldpage[0\textsubscript{IV-1}]{59}

\begin{proof}
事实上，$u$ 只能是映射 $(x,s)\mapsto(v(x),w(s))$；只须再验证它是环的
$A$-同态，而这由定义 (18.2.3) 立即得出。还应注意，$u$ 也使下图交换：
\[
  \xymatrix@C=2.2pc@R=1.5pc{
    D_B(L)\ar[r]^u\ar[d] & D_{B'}(L')\ar[d]\\
    B\ar[r]_v & B'
  }
\]
这一次，竖直箭头是增广同态。
\end{proof}

\begin{proposition}[18.2.10]
\label{0.18.2.10-fr}
设 $B$ 是 $A$-环，$L$ 是 $B$-双模，$E$ 是 $B$ 关于 $L$ 的一个 $A$-扩张。
令 $G$ 为从 $B$ 到 $E$ 的 $B$-环同态集合（换言之，为增广 $E\to B$ 的
$A$-同态右逆的集合），令 $G'$ 为从 $D_B(L)$ 到 $E$ 的 $A$-等价集合。映射
把 $g\in G$ 送到 $A$-等价 $g':(x,s)\mapsto g(x)+s$；这个映射是
双射，其逆映射把 $g'\in G'$ 送到 $B$-同态 $x\mapsto g'(x,0)$。
\end{proposition}

\begin{proof}
这立即由定义得出。
\end{proof}

\subsection{$A$-扩张类群}
\label{subsection:0.18.3-fr}

\begin{env}[18.3.1]
\label{0.18.3.1-fr}
固定一个 $A$-环 $B$ 和一个 $B$-双模 $L$。在 $B$ 关于 $L$ 的 $A$-扩张之间，
“$E$ 与 $E'$ 是 $A$-等价的”是一个等价关系；其等价类构成一个\emph{集合}。
要说明这一点，只须注意：对每个
$x\in B$，若 $c_x$ 是 $E$ 中一个在 $B$ 内的像为 $x$ 的元素，则每个
$z\in E$ 都能且只能写成 $c_x+t$ 的形式，其中 $t\in L$；并且可以写成
\[
  c_{x+y}=c_x+c_y+\varphi(x,y),\qquad
  c_{xy}=c_xc_y+\psi(x,y),
\]
其中 $\varphi(x,y)$ 和 $\psi(x,y)$ 是 $L$ 的元素；映射 $\varphi$ 和 $\psi$
从 $B\times B$ 到 $L$，并须满足表达 $E$ 是一个 $A$-环的若干条件，这里无须
把它们写出。因此，$B$ 关于 $L$ 的每个 $A$-扩张都 $A$-等价于一个以
$B\times L$ 为底层集合的 $A$-扩张；我们的结论立即由此得出。
\end{env}

\begin{env}[18.3.2]
\label{0.18.3.2-fr}
固定 $A$-环 $B$，暂以 $T(L)$ 表示 $B$ 关于 $L$ 的 $A$-扩张类的集合。
对每个 $B$-双模同态 $w:L\to L'$，典范地定义映射
$T(w):T(L)\to T(L')$：把扩张 $E$ 的类送到由 $E$ 经 $w$ 推得的扩张
$E\oplus_LL'$ 的类 (18.2.8)。若 $w':L'\to L''$ 是第二个 $B$-双模同态，则还有
\begin{equation}
\label{0.18.3.2.1-fr}
\tag{18.3.2.1}
  T(w'\circ w)=T(w')\circ T(w).
\end{equation}
事实上，存在 $A$-双模的典范同构
\[
  E\oplus_LL''\xrightarrow{\sim}(E\oplus_LL')\oplus_{L'}L''
\]
这是由归纳极限的一般性质（例如参见
($\mathrm{I}$, 3.3.9)）得出的；
立即可验，它确是这些 $A$-扩张的一个 $A$-等价；于是

\oldpage[0\textsubscript{IV-1}]{60}

得到关系 (18.3.2.1)。若 $\mathbf C_B$ 表示 $B$-双模范畴，则可见我们由此定义了
一个\emph{协变函子} $T:\mathbf C_B\to\mathbf{Ens}$。
\end{env}

\begin{env}[18.3.3]
\label{0.18.3.3-fr}
现取一族 $(L_\alpha)_{\alpha\in I}$ 的 $B$-双模，及其乘积
$L=\prod_{\alpha\in I}L_\alpha$；各投影
$\operatorname{pr}_\alpha:L\to L_\alpha$ 定义映射
\[
  T(\operatorname{pr}_\alpha):T(L)\to T(L_\alpha),
\]
从而定义一个典范映射
\begin{equation}
\label{0.18.3.3.1-fr}
\tag{18.3.3.1}
  \prod_\alpha T(\operatorname{pr}_\alpha):T(L)
  \longrightarrow\prod_{\alpha\in I}T(L_\alpha).
\end{equation}
我们来说明这个映射是\emph{双射}。事实上，对每个 $\alpha\in I$，设
$E_\alpha$ 是 $B$ 关于 $L_\alpha$ 的一个 $A$-扩张；并设 $F$ 为各
$E_\alpha$ 在 $B$ 上的\emph{纤维积} (18.1.3)。显然，$F$ 是 $B$ 关于
$L=\prod_{\alpha\in I}L_\alpha$ 的一个 $A$-扩张。若把每个 $E_\alpha$ 替换为
一个与它 $A$-等价的 $A$-扩张 $E'_\alpha$，则纤维积 $F'$——即各
$E'_\alpha$ 在 $B$ 上的纤维积——是 $A$-等价于 $F$ 的。这样就定义了一个映射
$\prod_{\alpha\in I}T(L_\alpha)\to T(L)$；显然，这个映射是 (18.3.3.1) 的逆映射，
断言由此得证。我们将通过映射 (18.3.3.1) 把 $T(L)$ 与
$\prod_{\alpha\in I}T(L_\alpha)$ 典范地等同起来。此外立即可验：若
$(L'_\alpha)_{\alpha\in I}$ 是第二族 $B$-双模，并且对每个 $\alpha$，
$w_\alpha:L_\alpha\to L'_\alpha$ 是一个 $B$-双模同态，那么，置
$w=\prod_\alpha w_\alpha:L\to L'=\prod_{\alpha\in I}L'_\alpha$ 后，在采用上述
等同的情形下，$T(w)$ 可与 $\prod_{\alpha\in I}T(w_\alpha)$ 等同。
\end{env}

\begin{env}[18.3.4]
\label{0.18.3.4-fr}
在此情形下，对于一个 $B$-双模 $L$，加法 $s:L\times L\to L$ 是一个
$B$-双模同态；$L$ 的加法取负映射 $t:L\to L$ 也是如此。由 (18.3.3) 得到合成律
\[
  T(s):T(L)\times T(L)\longrightarrow T(L)
\]
这个合成律使 $T(L)$ 成为一个\emph{交换群}，其取逆映射为 $T(t)$；这由用交换图
定义群对象的方式得出
($\mathbf 0_{\mathrm{III}}$, 8.2.5 和 8.2.6)。我们以
$\operatorname{Exan}_A(B,L)$ 表示这样定义的交换群，称它为\emph{$B$ 关于 $L$ 的
$A$-扩张类群}。
\end{env}

\begin{env}[18.3.5]
\label{0.18.3.5-fr}
以 $\mathbf K$ 表示如下范畴：其对象是三元组 $(A,B,L)$，其中 $A$ 是环，$B$ 是
$A$-环，而 $L$ 是 $B$-双模。这个范畴的\emph{态射}是三元组 $(u,v,w)$，其中
$u:A'\to A$ 和 $v:B'\to B$ 是两个环同态，使下图的左方块交换：
\[
  \xymatrix@C=2pc@R=1.5pc{
    A\ar[r] & B & L\ar[d]^w\\
    A'\ar[u]^u\ar[r] & B'\ar[u]^v & L'
  }
\]
这里水平箭头是结构同态；最后，$w:L\to L'$ 是交换群同态，并且满足
$w(v(b')z)=b'w(z)$ 以及

\oldpage[0\textsubscript{IV-1}]{61}

$w(zv(b'))=w(z)b'$，其中任意取 $z\in L$ 和 $b'\in B'$（换言之，$w$ 是一个
$B'$-双模同态，此处给 $L$ 所赋的 $B'$-双模结构由 $v$ 定义）。态射的复合规定为
$(u',v',w')\circ(u,v,w)=(u\circ u',v\circ v',w'\circ w)$；其合理性立即可见。
我们的目标是证明
\begin{equation}
\label{0.18.3.5.1-fr}
\tag{18.3.5.1}
  (A,B,L)\longmapsto\operatorname{Exan}_A(B,L)
\end{equation}
是从范畴 $\mathbf K$ 到交换群范畴 $\mathbf{Ab}$ 的一个\emph{协变函子}。因此，对
上述每个三元组 $(u,v,w)$，须定义一个交换群同态
\[
  (u,v,w)_*:\operatorname{Exan}_A(B,L)
  \longrightarrow\operatorname{Exan}_{A'}(B',L').
\]
由范畴 $\mathbf K$ 中态射的定义，可写成
\[
  (u,v,w)=(1_{A'},1_{B'},w)\circ(1_{A'},v,1_L)
  \circ(u,1_B,1_L)
\]
其中在第一个因子里，$L$ 的 $B'$-双模结构由 $v$ 定义；因此，我们先定义
$(u,v,w)_*$ 在 $u$、$v$、$w$ 中有两个为恒等映射的情形。
\end{env}

\begin{env}[18.3.6]
\label{0.18.3.6-fr}
首先，把 $(1_A,1_B,w)_*$ 取为映射
\begin{equation}
\label{0.18.3.6.1-fr}
\tag{18.3.6.1}
  w_*:\operatorname{Exan}_A(B,L)\longrightarrow
  \operatorname{Exan}_A(B,L')
\end{equation}
也就是 (18.3.2) 中记作 $T(w)$ 的映射。立即可验，这是一个群同态；该性质可用
一些交换图表示，而这些图是由 $T$ 将关于 $L$ 和 $L'$ 的类似图变换而来的。

映射 $(1_A,v,1_L)_*$ 取为如下映射：
\begin{equation}
\label{0.18.3.6.2-fr}
\tag{18.3.6.2}
  v^*:\operatorname{Exan}_A(B,L)\longrightarrow
  \operatorname{Exan}_A(B',L)
\end{equation}
其定义如下。若 $E$ 是 $B$ 关于 $L$ 的一个 $A$-扩张，则由 (18.2.5) 已知
$E\times_BB'$ 是 $B'$ 关于 $L$ 的一个 $A$-扩张；若把 $E$ 替换为一个与它
$A$-等价的 $A$-扩张 $E'$，则 $E'\times_BB'$ 在 $A$-等价意义下等同于
$E\times_BB'$。
$v^*$ 把 $E$ 的类映为 $E\times_BB'$ 的类。立即可验：若 $w:L\to L'$ 是一个
$B$-双模同态，则下图
\begin{equation}
\label{0.18.3.6.3-fr}
\tag{18.3.6.3}
  \xymatrix@C=3pc@R=2.4pc{
    \operatorname{Exan}_A(B,L)\ar[r]^{v^*}\ar[d]_{w_*}
      & \operatorname{Exan}_A(B',L)\ar[d]^{w_*}\\
    \operatorname{Exan}_A(B,L')\ar[r]_{v^*}
      & \operatorname{Exan}_A(B',L')
  }
\end{equation}
交换；在右列中，$L$ 和 $L'$ 都视为 $B'$-双模，其结构借助 $v$ 得到。把
$L$ 和 $L'$ 分别替换为 $L\times L$ 和 $L$，再把 $w$ 替换为加法 $s$
（定义在 $L$ 中），就可推出
$v^*$ 确是一个\emph{群同态}。
\end{env}
\oldpage[0\textsubscript{IV-1}]{62}

\begin{env}[18.3.6.4]
\label{0.18.3.6.4-fr}
最后，映射 $(u,1_B,1_L)_*$ 就是映射
\[
  u^*:\operatorname{Exan}_A(B,L)\longrightarrow
  \operatorname{Exan}_{A'}(B,L)
\]
其定义是：把 $B$ 关于 $L$ 的一个 $A$-扩张 $E$ 送到同一个环 $E$，
但借助 $u$ 将其视为 $A'$-\emph{环} (18.1.1)；这显然是 $B$ 关于 $L$ 的一个
$A'$-扩张，其中也借助 $u$ 将 $B$ 视为 $A'$-环。显然，一个
$A$-等价也是一个 $A'$-等价；由此得到映射 (18.3.6.4)。对每个同态
$w:L\to L'$，若它是 $B$-双模同态，则它仍使下图交换：
\begin{equation}
\label{0.18.3.6.5-fr}
\tag{18.3.6.5}
  \xymatrix@C=3pc@R=2.4pc{
    \operatorname{Exan}_A(B,L)\ar[r]^{u^*}\ar[d]_{w_*}
      & \operatorname{Exan}_{A'}(B,L)\ar[d]^{w_*}\\
    \operatorname{Exan}_A(B,L')\ar[r]_{u^*}
      & \operatorname{Exan}_{A'}(B,L')
  }
\end{equation}
因而同上可推出，$u^*$ 是群同态。

现在置
$(u,v,w)_*=(1_{A'},1_{B'},w)_*\circ(1_{A'},v,1_L)_*
\circ(u,1_B,1_L)_*$。由图 (18.3.6.3) 和 (18.3.6.5) 的交换性，容易验证
\[
  (u\circ u',v\circ v',w'\circ w)_*
  =(u',v',w')_*\circ(u,v,w)_*,
\]
这就完成了 (18.3.5.1) 是一个函子的证明。
\SourceCorrectionNote{法文印本及外交式转录在 18.3.6 末尾误引作 (18.3.4.1)；应为 (18.3.5.1)，因为刚刚验证复合律的正是该式所定义的函子。}

群同态 (18.3.6.1) 的存在尤其表明：若 $E$ 是 $B$ 关于 $L$ 的一个\emph{平凡}
$A$-扩张，则 (18.2.8) 中定义的 $B$ 关于 $L'$ 的扩张 $E\oplus_LL'$ 也是平凡的；
这当然也不难直接验证。

另一方面，设 $z$ 是 $B$ 的中心 $Z$ 的任意元素；右乘 $z$ 的映射
$h_z:y\mapsto yz$ 是 $B$-双模 $L$ 的一个自同态，故 $(h_z)_*$ 是交换群
$\operatorname{Exan}_A(B,L)$ 的一个自同态；由函子性，这些自同态在
$\operatorname{Exan}_A(B,L)$ 上定义了一个典范的 $Z$-\emph{模}结构。
\SourceCorrectionNote{法文印本及外交式转录在倍乘公式中写作 yz=zy； 应删去 =zy，因为底环中心元素在任意 双模上的左、右作用未必相同。这里的倍乘映射由右作用定义。}
\end{env}

\begin{env}[18.3.7]
\label{0.18.3.7-fr}
设 $A$、$A'$ 是两个环，$u:A'\to A$ 是一个同态，$B$ 是一个 $A$-环，而
$L$ 是一个 $B$-双模。群同态
\[
  u^*:\operatorname{Exan}_A(B,L)\longrightarrow
  \operatorname{Exan}_{A'}(B,L)
\]
的\emph{核}，按定义由 $B$ 关于 $L$ 的那些 $A$-扩张的类组成，这些扩张借助
$u$ 视为 $A'$-扩张时是 $A'$-\emph{平凡}的。不致混淆时，把这个核记作
$\operatorname{Exan}_{A/A'}(B,L)$。

设 $\Lambda$ 是一个环，并且有一个由 $\Lambda$-环的 $\Lambda$-\emph{同态}组成的
交换图
\begin{equation}
\label{0.18.3.7.1-fr}
\tag{18.3.7.1}
  \xymatrix@C=2pc@R=1.5pc{
    B'\ar[r] & B\\
    A'\ar[u]\ar[r] & A\ar[u]
  }
\end{equation}

\oldpage[0\textsubscript{IV-1}]{63}

则可由此典范地得到同态
\begin{equation}
\label{0.18.3.7.2-fr}
\tag{18.3.7.2}
  \operatorname{Exan}_{A/\Lambda}(B,L)\longrightarrow
  \operatorname{Exan}_{A'/\Lambda}(B,L)\longrightarrow
  \operatorname{Exan}_{A'/\Lambda}(B',L)
\end{equation}
它们来自下图的交换性：
\[
  \xymatrix@C=2.8pc@R=2.3pc{
    \operatorname{Exan}_A(B,L)\ar[r]\ar[d]
      & \operatorname{Exan}_{A'}(B,L)\ar[r]\ar[d]
      & \operatorname{Exan}_{A'}(B',L)\ar[d]\\
    \operatorname{Exan}_{\Lambda}(B,L)\ar[r]
      & \operatorname{Exan}_{\Lambda}(B,L)\ar[r]
      & \operatorname{Exan}_{\Lambda}(B',L)
  }
\]
其中各箭头都由 (18.3.7.1) 中的箭头借助函子性导出。
\end{env}

\begin{proposition}[18.3.8]
\label{0.18.3.8-fr}
设 $B$ 是一个环，$\mathfrak J$ 是 $B$ 的一个双边理想，而
$C=B/\mathfrak J$ 是商环；于是 $\mathfrak J/\mathfrak J^2$ 典范地具有一个
$C$-双模结构。对任意 $C$-双模 $L$，以
$\operatorname{Hom}_C(\mathfrak J/\mathfrak J^2,L)$ 表示 $C$-双模同态所成的
加法群，这些同态从 $\mathfrak J/\mathfrak J^2$ 到 $L$。于是可定义
交换群的一个典范同构
\begin{equation}
\label{0.18.3.8.1-fr}
\tag{18.3.8.1}
  \eta_L:\operatorname{Hom}_C(\mathfrak J/\mathfrak J^2,L)
  \xrightarrow{\sim}\operatorname{Exan}_B(C,L)
\end{equation}
如下：对每个 $C$-同态 $w:\mathfrak J/\mathfrak J^2\to L$（它当然也是一个
$B$-同态），让它对应于扩张
$(B/\mathfrak J^2)\oplus_{(\mathfrak J/\mathfrak J^2)}L$ 的类；该扩张由
$C$ 关于 $\mathfrak J/\mathfrak J^2$ 的扩张 $B/\mathfrak J^2$ 沿 $w$ 推得。
其逆同构把 $C$ 关于 $L$ 的一个 $B$-扩张 $E$ 的类送到同态 $w$，使复合
$\mathfrak J\to\mathfrak J/\mathfrak J^2\xrightarrow{w}L$ 等于结构同态 $B\to E$
在 $\mathfrak J$ 上的限制。
\end{proposition}

\begin{proof}
置 $F=B/\mathfrak J^2$；它是 $C$ 关于 $\mathfrak J/\mathfrak J^2$ 的一个
$B$-扩张。对任意 $C$ 关于 $L$ 的 $B$-扩张
\[
  0\longrightarrow L\xrightarrow{j}E\xrightarrow{p}C\longrightarrow0
\]
结构同态 $f:B\to E$ 满足：复合
$B\xrightarrow{f}E\xrightarrow{p}C$ 是典范同态
$B\to B/\mathfrak J=C$。由于 $\mathfrak J$ 在 $p\circ f$ 下的像为零，
$f(\mathfrak J)$ 含于 $p$ 的核，即 $j(L)$；又因 $j(L)$ 的
平方为零，故 $f(\mathfrak J^2)=0$。所以 $f$ 分解为
$B\to B/\mathfrak J^2=F\xrightarrow{u}E$。若以 $w$ 表示 $u$ 在
$\mathfrak J/\mathfrak J^2$ 上的限制，便有交换图
\[
  \xymatrix@C=2pc@R=1.5pc{
    0\ar[r] & \mathfrak J/\mathfrak J^2\ar[r]\ar[d]^w
      & F\ar[r]^g\ar[d]^u
      & C\ar[r]\ar[d]^{1_C}\ar@{}[d]|{\wr} & 0\\
    0\ar[r] & L\ar[r]_j & E\ar[r]_p & C\ar[r] & 0
  }
\]

换言之，$(u,1_C,w)$ 是一个扩张态射 (18.2.4)。由此得到一个扩张态射
\[
  \xymatrix@C=2pc@R=1.5pc{
    0\ar[r] & L\ar[r]\ar[d]^{1_L}\ar@{}[d]|{\wr}
      & E'\ar[r]\ar[d]^{u'}
      & C\ar[r]\ar[d]^{1_C}\ar@{}[d]|{\wr} & 0\\
    0\ar[r] & L\ar[r]_j & E\ar[r]_p & C\ar[r] & 0
  }
\]

\oldpage[0\textsubscript{IV-1}]{64}

其中 $E'=F\oplus_{(\mathfrak J/\mathfrak J^2)}L$ (18.2.8)。须证明 $u'$ 是一个
$B$-\emph{等价}；这将表明 $\eta_L$ 为双射。由 (18.2.8) 中的构造，只须证明
（并计及 (18.1.7)）：$C$ 的扩张 $E$ 沿典范同态 $g:F\to C$ 的逆像 $G$
是 $F$-平凡的。这是显然的，因为 $g$ 分解为
$F\xrightarrow{u}E\xrightarrow{p}C$ (18.1.6)。

还须验证 $\eta_L$ 是一个\emph{群}同态。对任意 $B$-同态 $h:L\to L'$，下图
显然交换：
\[
  \xymatrix@C=3pc@R=2.4pc{
    \operatorname{Hom}_C(\mathfrak J/\mathfrak J^2,L)
      \ar[r]^{\eta_L}\ar[d]_{\operatorname{Hom}(1,h)}
      & \operatorname{Exan}_B(C,L)\ar[d]^{h_*}\\
    \operatorname{Hom}_C(\mathfrak J/\mathfrak J^2,L')
      \ar[r]_{\eta_{L'}}^{\sim}
      & \operatorname{Exan}_B(C,L')
  }
\]
将这一点用于定义加法的同态 $L\times L\to L$，即可得出结论。
\end{proof}
\subsection{代数的扩张}
\label{subsection:0.18.4-fr}

\begin{env}[18.4.1]
\label{0.18.4.1-fr}
设 $A$ 是一个\emph{交换}环。于是，$A$-\emph{代数}范畴是 $A$-环范畴的一个
满子范畴；其特征是结构同态 $\rho:A\to B$ 的像含于中心。

若 $B$ 是一个 $A$-代数，$L$ 是一个 $B$-双模，且其经 $B$ 诱导的两个
$A$-模结构相同，则显然，典型平凡
$A$-扩张 $D_B(L)$ 是一个 $A$-代数。同样，在 (18.2.5)（相应地，(18.2.8)）
的构造中，若 $B$、$B'$ 和 $E'$（相应地，$B$ 和 $E$）都是 $A$-代数，则
$E'\times_{B'}B$（相应地，$E\oplus_LL'$）也同样如此。最后，显然，若
$B$ 是一个 $A$-代数，而 $E$ 是 $B$ 关于 $B$-双模 $L$ 的一个 $A$-扩张，
且本身是 $A$-代数，则 $B$ 关于 $L$ 的任意 $A$-扩张，只要
$A$-\emph{等价于} $E$，也都是 $A$-代数。于是，由定义 (18.3.4) 立即推出：
$B$ 关于 $L$ 的那些本身为 $A$-\emph{代数}的 $A$-扩张的等价类
构成一个\emph{子群}，记作 $\operatorname{Exal}_A(B,L)$；它是
$\operatorname{Exan}_A(B,L)$ 的子群。设 $\mathbf K'$ 是 (18.3.5) 所定义的范畴
$\mathbf K$ 的满子范畴；其对象 $(A,B,L)$ 满足：$A$ 交换，$B$ 是一个
$A$-代数，且 $L$ 经 $B$ 诱导的两个 $A$-模结构相同。上述结论表明
\[
  (A,B,L)\longmapsto\operatorname{Exal}_A(B,L)
\]
是从 $\mathbf K'$ 到 $\mathbf{Ab}$ 的一个\emph{协变函子}。若处处把
$\operatorname{Exan}$ 换成 $\operatorname{Exal}$，(18.3.7) 的结果保持不变。
\SourceCorrectionNote{法文印本及外交式转录在本段的典型平凡扩张及范畴对象 条件中，均遗漏双模经底代数诱导的两个 A-模结构相同这一假设；在这两个位置补入该假设。这是结构同态的像含于 平凡扩张的中心所需的条件。}
\end{env}

\begin{env}[18.4.2]
\label{0.18.4.2-fr}
仍设环 $A$ 交换。把“代数”换成“交换代数”，并把“双模”换成“模”，
(18.4.1) 的论述仍然成立。若 $B$ 是一个交换 $A$-代数，而 $L$ 是一个
$B$-模，

\oldpage[0\textsubscript{IV-1}]{65}

则 $B$ 关于 $L$ 的那些本身为 $A$-\emph{交换代数}（或等价地，为
$A$-\emph{交换环}）的 $A$-扩张的等价类构成
$\operatorname{Exal}_A(B,L)$ 的一个\emph{子群}，记作
$\operatorname{Exalcom}_A(B,L)$。若 $\mathbf K''$ 是 $\mathbf K'$ 的满子范畴，
由三元组 $(A,B,L)$ 构成，其中 $A$ 与 $B$ 交换，而 $L$ 是一个 $B$-模，则
\[
  (A,B,L)\longmapsto\operatorname{Exalcom}_A(B,L)
\]
仍是从 $\mathbf K''$ 到 $\mathbf{Ab}$ 的一个\emph{协变函子}。在 (18.3.7) 中，
也可以处处把 $\operatorname{Exan}$ 换成 $\operatorname{Exalcom}$。最后，在
(18.3.8) 中，若假设 $B$ 是一个\emph{交换}环，而 $L$ 是一个 $C$-\emph{模}，
则同样的论证给出典范同构
\begin{equation}
\label{0.18.4.2.1-fr}
\tag{18.4.2.1}
  \operatorname{Hom}_C(\mathfrak J/\mathfrak J^2,L)
  \xrightarrow{\sim}\operatorname{Exalcom}_B(C,L)
\end{equation}
其中左端是 $C$-\emph{模}同态所成的群。
\end{env}

\begin{env}[18.4.3]
\label{0.18.4.3-fr}
设 $A$ 是一个交换环。$A$-代数扩张的一类重要情形如下：$A$-代数 $E$ 是
$A$-代数 $B$ 关于 $B$-双模 $L$ 的扩张，且 $E$ 作为 $A$-\emph{模}，是
$B$ 关于 $L$ 的\emph{平凡}扩张；换言之，$A$-\emph{模}的正合列
$0\to L\to E\to B\to0$ 是\emph{分裂的}。此时称 $E$ 是 $B$ 关于 $L$ 的
一个 \emph{Hochschild 扩张}。当 $B$ 是一个\emph{投射} $A$-模时，总是如此；
特别地，当 $A$ 是一个交换\emph{域}时也是如此。作为 $A$-模，可把 $E$ 与
$B\times L$ 等同；$E$ 中的乘法由
$(x,s)(y,t)=(xy,xt+sy+f(x,y))$ 给出，其中 $f(x,y)\in L$。写出这个乘法在
$B\times L$ 上定义 $A$-代数结构的条件，便得到 (M, XIV, 2)：$f$ 必须是
一个 $A$-\emph{双线性}映射，从 $B\times B$ 到 $L$，并满足
\begin{equation}
\label{0.18.4.3.1-fr}
\tag{18.4.3.1}
  f(xy,z)+f(x,y)z=xf(y,z)+f(x,yz),
\end{equation}
换言之，$f$ 是 Hochschild 意义下的一个 $2$-\emph{上闭链}，它定义在 $B$
上并取值于 $L$；
扩张 $E$ 为 $A$-平凡的充要条件是
\begin{equation}
\label{0.18.4.3.2-fr}
\tag{18.4.3.2}
  f(x,y)=xg(y)-g(xy)+g(x)y,
\end{equation}
其中 $g$ 是一个 $A$-\emph{线性}映射，从 $B$ 到 $L$；换言之，$f$ 必须是
Hochschild 意义下的一个 $2$-\emph{上边缘}。由此立即推出：$B$ 关于 $L$ 的
Hochschild 扩张的等价类构成 $\operatorname{Exal}_A(B,L)$ 的一个子群，
它同构于 Hochschild \emph{上同调群} $\mathrm H_A^2(B,L)$。

若 $B$ 是一个\emph{交换} $A$-代数，而 $L$ 是一个 $B$-模，则 $B$ 关于 $L$ 的
\emph{交换} Hochschild 扩张对应于\emph{对称的} $2$-上闭链，也就是满足
$f(y,x)=f(x,y)$ 的上闭链。因此，$B$ 关于 $L$ 的交换 Hochschild 扩张的等价类
构成 $\operatorname{Exalcom}_A(B,L)$ 的一个子群；它同构于
$\mathrm H_A^2(B,L)$ 的如下子群：由对称 $2$-上闭链所成的群在其中的像。我们把
这个子群记作 $\mathrm H_A^2(B,L)^s$。
\end{env}

\begin{env}[18.4.4]
\label{0.18.4.4-fr}
我们限于 $A$ 和 $B$ 都\emph{交换}、而 $L$ 是一个 $B$-模的情形，并回顾此时
Hochschild 群 $\mathrm H_A^i(B,L)$ 的一个等价定义 (M, IX, 6)。考虑复形
$\mathbf P_\bullet=(\mathbf P_n)_{n\geqslant0}$，它由 $B$-模组成，其中
$\mathbf P_n=B^{\otimes(n+1)}=B\otimes_A B\otimes_A\cdots\otimes_A B$
（共 $n+1$ 个因子），其 $B$-模结构由下式定义：

\oldpage[0\textsubscript{IV-1}]{66}

\[
  x(y_1\otimes y_2\otimes\cdots\otimes y_{n+1})
  =(xy_1)\otimes y_2\otimes\cdots\otimes y_{n+1}\,;
\]
次数为 $-1$ 的微分 $d_n:\mathbf P_n\to\mathbf P_{n-1}$ 定义为
\begin{align*}
  d_n(x_1\otimes x_2\otimes\cdots\otimes x_{n+1})
    &=(x_1x_2)\otimes x_3\otimes\cdots\otimes x_{n+1}
      -x_1\otimes(x_2x_3)\otimes\cdots\otimes x_{n+1}+\cdots\\
    &\quad+(-1)^{n-1}x_1\otimes x_2\otimes\cdots\otimes x_{n-1}
      \otimes(x_nx_{n+1})
      +(-1)^n(x_{n+1}x_1)\otimes x_2\otimes\cdots\otimes x_n,
\end{align*}
由于 $B$ 交换，这的确是 $B$-线性的。

这个复形的一个 $2$-上闭链，取值于 $L$，是一个 $B$-线性映射 $h$，从
$B\otimes B\otimes B$ 到 $L$，并满足
$h(d_3(x\otimes y\otimes z\otimes t))=0$。但由于
$h(x\otimes y\otimes z)=xh(1\otimes y\otimes z)$，上闭链 $h$ 由一个
$A$-双线性映射 $(y,z)\longmapsto f(y,z)=h(1\otimes y\otimes z)$ 决定；
该映射从 $B\times B$ 到 $L$。对 $h$ 写出上述条件，
就重新得到 $f$ 的条件 (18.4.3.1)。同样，一个 $2$-上边缘是形如
$x\otimes y\otimes z\longmapsto h'(d_2(x\otimes y\otimes z))$ 的映射，其中
$h':B\otimes_A B\to L$ 是 $B$-线性的；这里，$h'$ 仍由 $A$-线性映射
$g:y\longmapsto h'(1\otimes y)$ 决定；该映射从 $B$ 到 $L$。于是得到
$h'(d_2(1\otimes x\otimes y))=xg(y)-g(xy)+yg(x)$，这就重新给出
(18.4.3.2)。
\SourceCorrectionNote{法文印本及外交式转录在 18.4.4 的这一步把映射 g 误作 B-线性；应为 A-线性，因为 g 由 B 与自身在 A 上的张量积上的 B-线性映射 h' 经映射 y↦1⊗y 诱导，而所显示的上边缘公式所需且所得正是 A-线性。}
对任意 $i$ 同样处理，就可看出
\begin{equation}
\label{0.18.4.4.1-fr}
\tag{18.4.4.1}
  \mathrm H_A^i(B,L)=\mathrm H^i(\operatorname{Hom}_B(\mathbf P_\bullet,L)).
\end{equation}
\end{env}

\begin{env}[18.4.5]
\label{0.18.4.5-fr}
在 (18.4.4) 的条件下，可以用同样的方法解释群
$\mathrm H_A^2(B,L)^s$ (18.4.3)。为此，把复形 $\mathbf P_\bullet$ 在次数 $3$
处加以修改，考虑新复形
\[
  \mathbf P'_\bullet:\mathbf P'_3\xrightarrow{d'_3}\mathbf P_2
  \xrightarrow{d_2}\mathbf P_1\,;
\]
我们取 $\mathbf P'_3=\mathbf P_3\oplus(B\otimes_A B\otimes_A B)$；
$d'_3$ 与 $d_3$ 在 $\mathbf P_3$ 上相同，而在
$B\otimes_A B\otimes_A B$ 上，它由下式给出：
\[
  d'_3(x\otimes y\otimes z)=x\otimes y\otimes z-x\otimes z\otimes y.
\]

关系 $d_2d'_3=0$ 由 $B$ 的交换性得到。采用上面引入的记号，一个
$2$-上闭链（属于 $\mathbf P'_\bullet$）现在对应于一个\emph{对称的}
$A$-双线性映射 $f:B\times B\to L$，并且它满足 (18.4.3.1)；因此
\[
  \mathrm H_A^2(B,L)^s
  =\mathrm H^2(\operatorname{Hom}_B(\mathbf P'_\bullet,L)).
\]
\end{env}

\begin{env}[18.4.6]
\label{0.18.4.6-fr}
在特殊情形下，设 $k$ 是一个交换\emph{域}，而 $K$ 是 $k$ 的一个
\emph{扩域}；恒有 $\operatorname{Ext}_K^1(M,L)=0$，这里的 $K$-向量空间
$L$、$M$ 可以任意选取；因此，由 (M, VI, 3.3 $a$)，有典范同构
\begin{equation}
\label{0.18.4.6.1-fr}
\tag{18.4.6.1}
  \mathrm H_k^i(K,L)\xrightarrow{\sim}
  \operatorname{Hom}_K(\mathrm H_i(\mathbf P_\bullet),L)
\end{equation}
以及同样的典范同构
\begin{equation}
\label{0.18.4.6.2-fr}
\tag{18.4.6.2}
  \mathrm H_k^2(K,L)^s\xrightarrow{\sim}
  \operatorname{Hom}_K(\mathrm H_2(\mathbf P'_\bullet),L).
\end{equation}
\SourceCorrectionNote{法文印本及外交式转录在本段文献引注末尾多出一个右圆括号； 删去该括号，原引注为 (M, VI, 3.3 a))。}
\end{env}

\subsection{拓扑环的情形}
\label{subsection:0.18.5-fr}

\begin{env}[18.5.1]
\label{0.18.5.1-fr}
设 $A$、$B$ 是两个具有\emph{线性}拓扑的\emph{拓扑}环，
$\rho:A\to B$ 是一个连续同态，而 $L$ 是一个拓扑 $B$-双模。假设存在
一个\emph{开}双边理想 $\mathfrak R_0$（属于 $B$），满足
$\mathfrak R_0.L=L.\mathfrak R_0=0$，从而 $L$ 可视作

\oldpage[0\textsubscript{IV-1}]{67}

一个 $(B/\mathfrak R)$-双模，只要开理想满足
$\mathfrak R\subset\mathfrak R_0$。现在，设
$\mathfrak R\subset\mathfrak R_0$ 是 $B$ 的一个开双边理想，$\mathfrak J$ 是
$A$ 的一个开双边理想，满足 $\rho(\mathfrak J)\subset\mathfrak R$，从而
$B/\mathfrak R$ 可视作一个离散的 $(A/\mathfrak J)$-环。前面的论述说明群
$\operatorname{Exan}_{A/\mathfrak J}(B/\mathfrak R,L)$ 有定义。此外，若
$\mathfrak R'\subset\mathfrak R$、$\mathfrak J'\subset\mathfrak J$ 分别是
$B$ 和 $A$ 的两个开双边理想，并且满足
$\rho(\mathfrak J')\subset\mathfrak R'$，则由 (18.3.5.1) 有典范同态
\begin{equation}
\label{0.18.5.1.1-fr}
\tag{18.5.1.1}
  \operatorname{Exan}_{A/\mathfrak J}(B/\mathfrak R,L)
  \longrightarrow
  \operatorname{Exan}_{A/\mathfrak J'}(B/\mathfrak R',L).
\end{equation}
\SourceCorrectionNote{法文印本及外交式转录在 18.5.1 中把带撇理想的所属环 互换，误称 R' 与 J' 分别属于 A 与 B； R' 属于 B、J' 属于 A， 因为相应商环正是 B/R' 与 A/J'，且结构同态要求 J' 的像包含在 R' 中。}

开理想偶 $(\mathfrak J,\mathfrak R)$ 所成的集合，若这些偶满足
$\rho(\mathfrak J)\subset\mathfrak R\subset\mathfrak R_0$，则关于关系“
$\mathfrak J\supset\mathfrak J'$ 且 $\mathfrak R\supset\mathfrak R'$”显然成为
一个向右滤过的有序集；诸映射 (18.5.1.1) 以这个集合为指标集，定义了加法群的
一个\emph{归纳系统}。虽然后面的群不再与某个扩张集合自然地一一对应，我们仍
滥用记号定义
\begin{equation}
\label{0.18.5.1.2-fr}
\tag{18.5.1.2}
  \operatorname{Exantop}_A(B,L)
  =\varinjlim\operatorname{Exan}_{A/\mathfrak J}(B/\mathfrak R,L).
\end{equation}

说 (18.5.1.2) 的右端为\emph{零}，因而意味着：对每一对开理想
$\mathfrak J\subset A$、$\mathfrak R\subset B$，若满足
$\rho(\mathfrak J)\subset\mathfrak R\subset\mathfrak R_0$，以及 $B/\mathfrak R$ 关于
$L$ 的每个 $(A/\mathfrak J)$-扩张 $E$，都存在两个
开理想 $\mathfrak J'\subset\mathfrak J$、$\mathfrak R'\subset\mathfrak R$，
满足 $\rho(\mathfrak J')\subset\mathfrak R'$，并且沿同态
$B/\mathfrak R'\to B/\mathfrak R$ 所取的 $E$ 的逆像为\emph{平凡}扩张。

我们把归纳极限 $\operatorname{Exaltop}_A(B,L)$、
$\operatorname{Exalcotop}_A(B,L)$ 的类似定义留给读者；它们分别由
$\operatorname{Exal}_{A/\mathfrak J}(B/\mathfrak R,L)$ 和
$\operatorname{Exalcom}_{A/\mathfrak J}(B/\mathfrak R,L)$ 得到，适用于 $A$
交换而 $B$ 是一个拓扑 $A$-代数（相应地，一个拓扑交换 $A$-代数）的情形。
\end{env}

\begin{env}[18.5.2]
\label{0.18.5.2-fr}
若有一个由连续环同态组成的交换图
\[
  \xymatrix@C=3pc@R=2pc{
    B'\ar[r] & B\\
    A'\ar[u]\ar[r] & A\ar[u]
  }
\]
则从 (18.3.5.1) 出发取归纳极限，可典范地得到加法群的两个同态
\[
  \operatorname{Exantop}_A(B,L)
  \longrightarrow\operatorname{Exantop}_{A'}(B,L)
  \longrightarrow\operatorname{Exantop}_{A'}(B',L)
\]

由于函子 $\varinjlim$ 在交换群范畴中是正合的，同态
\[
  \operatorname{Exantop}_A(B,L)
  \longrightarrow\operatorname{Exantop}_{A'}(B,L)
\]
的核，是下列同态之核的归纳极限：
\[
  \operatorname{Exan}_{A/\mathfrak J}(B/\mathfrak R,L)
  \longrightarrow
  \operatorname{Exan}_{A'/\mathfrak J'}(B/\mathfrak R,L)
\]
其中以 $\mathfrak J'$ 表示 $\mathfrak J$ 的逆像；这个核记作
$\operatorname{Exantop}_{A/A'}(B,L)$。同样定义
$\operatorname{Exaltop}_{A/A'}(B,L)$ 和
$\operatorname{Exalcotop}_{A/A'}(B,L)$。最后，若有一个

\oldpage[0\textsubscript{IV-1}]{68}

连续的 $B$-双模同态 $L\to L'$，则从 (18.3.6.1) 出发取归纳极限，可典范地
得到加法群同态
\[
  \operatorname{Exantop}_A(B,L)\longrightarrow
  \operatorname{Exantop}_A(B,L')
\]
\end{env}

\begin{env}[18.5.3]
\label{0.18.5.3-fr}
给定一个拓扑环 $C$ 和两个拓扑 $C$-双模 $M$、$N$，以
$\operatorname{Hom.\ cont}_C(M,N)$ 表示连续 $C$-同态（从 $M$ 到 $N$）所成的
加法群。
\end{env}

\begin{lemma}[18.5.3.1]
\label{0.18.5.3.1-fr}
设 $C$ 是一个拓扑环，而 $E$、$L$ 是两个拓扑 $C$-双模。假设这些拓扑是
线性的，$L$ 是离散的，并被 $C$ 的一个开双边理想零化。于是有典范同构
\begin{equation}
\label{0.18.5.3.2-fr}
\tag{18.5.3.2}
  \varinjlim\operatorname{Hom}_{C/\mathfrak R}(E/V,L)
  \xrightarrow{\sim}\operatorname{Hom.\ cont}_C(E,L)
\end{equation}
其中，左端的归纳极限取遍如下偶 $(\mathfrak R,V)$ 所成的滤过有序集：
$\mathfrak R$ 是 $C$ 的一个开双边理想，$V$ 是 $E$ 的一个开 $C$-子双模，
并满足 $\mathfrak R.L=L.\mathfrak R=0$、$\mathfrak R.E\subset V$ 和
$E.\mathfrak R\subset V$。
\end{lemma}

\begin{proof}
由于 $C/\mathfrak R$ 和 $E/V$ 都是离散的，存在典范同态
\[
  w_{\mathfrak R,V}:\operatorname{Hom}_{C/\mathfrak R}(E/V,L)
  \longrightarrow\operatorname{Hom.\ cont}_C(E,L)
\]
它们构成一个归纳系统；取归纳极限即得到同态 (18.5.3.2)。同态
$E/V'\to E/V$ 在 $V'\subset V$ 时是满射；由定义立即可知同态
\[
  \operatorname{Hom}_{C/\mathfrak R}(E/V,L)
  \longrightarrow\operatorname{Hom}_{C/\mathfrak R}(E/V',L)
\]
是单射（这里 $\mathfrak R.L=L.\mathfrak R=0$、
$\mathfrak R.E\subset V'$、$E.\mathfrak R\subset V'$）；下列同态
\[
  \operatorname{Hom}_{C/\mathfrak R}(E/V,L)
  \longrightarrow\operatorname{Hom}_{C/\mathfrak R'}(E/V,L)
\]
在 $\mathfrak R'\subset\mathfrak R$ 时显然也同样是单射。由此推出同态
(18.5.3.2) 是单射。另一方面，若 $u$ 是一个连续 $C$-同态（从 $E$ 到 $L$），
则其核是 $E$ 的一个开子双模 $V_0$；若 $\mathfrak R_0$ 是
$C$ 的一个开双边理想，满足 $\mathfrak R_0.L=L.\mathfrak R_0=0$、
$\mathfrak R_0.E\subset V_0$ 和 $E.\mathfrak R_0\subset V_0$，则显然，$u$ 是
某个 $(C/\mathfrak R_0)$-同态（从 $E/V_0$ 到 $L$）的典范像。因此
(18.5.3.2) 是满射。
\end{proof}

现在，命题 (18.3.8) 可如下推广到拓扑环：

\begin{proposition}[18.5.4]
\label{0.18.5.4-fr}
设 $B$ 是一个具有线性拓扑的拓扑环，$\mathfrak J$ 是 $B$ 的一个双边理想，
$C=B/\mathfrak J$ 是商拓扑环。于是 $\mathfrak J/\mathfrak J^2$（其中，
$\mathfrak J$ 具有从 $B$ 诱导的拓扑，而 $\mathfrak J/\mathfrak J^2$ 具有从
$\mathfrak J$ 得到的商拓扑）典范地具有拓扑 $C$-双模结构。对任意离散
$C$-双模 $L$，若它被 $C$ 的一个开理想零化，则存在典范同构
\begin{equation}
\label{0.18.5.4.1-fr}
\tag{18.5.4.1}
  \operatorname{Hom.\ cont}_C(\mathfrak J/\mathfrak J^2,L)
  \xrightarrow{\sim}\operatorname{Exantop}_B(C,L).
\end{equation}
\end{proposition}

\begin{proof}
事实上，对每个开理想 $\mathfrak R$（属于 $B$），若
$(\mathfrak J+\mathfrak R)/\mathfrak J$ 零化 $L$，则由 (18.3.8) 有典范同构
\[
  \operatorname{Hom}_{B/(\mathfrak J+\mathfrak R)}
    ((\mathfrak J+\mathfrak R)/(\mathfrak J^2+\mathfrak R),L)
  \xrightarrow{\sim}
  \operatorname{Exan}_{B/\mathfrak R}(B/(\mathfrak J+\mathfrak R),L)
\]
只须取归纳极限，并计及 (18.5.3.1)，即可得证。
\SourceCorrectionNote{法文印本及外交式转录在 18.5.4 的证明中引用不存在的 (18.3.9)；应为 (18.3.8)，因为该命题正给出此处逐层使用的典范同构。}
\end{proof}

\section{形式平滑代数与 Cohen 环}
\label{section:0.19-fr}
\label{0.19-fr}

\setcounter{subsection}{-1}
\subsection{引言}
\label{subsection:0.19.0-fr}
\label{0.19.0-fr}

\oldpage[0\textsubscript{IV-1}]{69}

\begin{env}[19.0.1]
\label{0.19.0.1-fr}
在第四章中，我们将引入并研究一类重要的预概形态射，即\emph{光滑态射}\footnote{
某些近期工作（仿照经典术语“简单点”）也把它们称为\emph{简单态射}；不过这一
术语容易引起混淆，尤其是在代数群理论中。}。它们的一项基本性质（与一个有限性
条件合在一起时，这项性质也刻画它们）是\emph{态射的提升性质}：若
$f:X\to Y$ 是光滑态射，$g:Y'\to Y$ 是一个态射，而 $h:Y''\to Y'$ 是一个
使 $Y''$ 成为与 $Y'$“相差很小”的预概形的态射，则\emph{每个 $Y$-态射
$Y''\to X$ 都可分解}为 $Y''\xrightarrow{h}Y'\to X$。更确切地说，我们只考察
$Y=\operatorname{Spec}(A)$、$X=\operatorname{Spec}(B)$ 都是\emph{仿射}的
情形。若满足下列条件，则称 $B$ 是一个\emph{形式平滑}的 $A$-代数：对任意
$A$-代数 $C$ 和任意\emph{幂零}理想 $\mathfrak J$（属于 $C$），每个 $A$-同态
$B\to C/\mathfrak J$ 都能\emph{提升}为 $B\to C\to C/\mathfrak J$。
换言之，映射
\[
  \operatorname{Hom}_A(B,C)\longrightarrow
  \operatorname{Hom}_A(B,C/\mathfrak J)
\]
是\emph{满射}。在许多应用中，$X$ 将作为一个\emph{可表反变函子}
$Y'\mapsto F(Y')$ 的表示对象；该函子从 $Y$-预概形范畴到集合范畴，因而
($\mathbf{0}_{\mathrm{III}}$, 8.1.8) 有
$F(Y')=\operatorname{Hom}_Y(Y',X)$。在仿射情形，若置
$F^0(C)=F(\operatorname{Spec}(C))$，那么证明在上述条件下
$F^0(C)\to F^0(C/\mathfrak J)$ 是\emph{满射}的（甚至可以在\emph{不知道}
$F$ 可表时完成这一验证），只要附加的有限性条件得到满足，就能确定 $f$ 是
光滑的。
\end{env}

\begin{env}[19.0.2]
\label{0.19.0.2-fr}
对于拓扑环 $A$ 上的\emph{拓扑代数}，有一个类似的\emph{形式平滑}代数概念，
这里暂不详述（参见 19.3.1）。我们首先在 §~19 中从初等观点研究这些概念，
然后利用 §§~20 和 21 中将要发展的微分性质，在 §~22 中证明更精细的定理。
现将形式平滑代数的主要结果概述如下：

\medskip\noindent\textbf{I.} --- 设 $A$、$B$ 是两个 Noether 局部环，
$\varphi:A\to B$ 是局部同态，$k$ 是 $A$ 的剩余域，并置
$B_0=B\otimes_A k$；在 $A$、$B$ 和 $B_0$ 上取其预赋值拓扑，在 $k$ 上取
离散拓扑。于是，$B$ 是形式平滑的 $A$-代数，当且仅当 $B$ 是平坦的
$A$-模，并且 $B_0$ 是形式平滑的 $k$-代数 (19.7.1)。因而，对 Noether
局部环而言，这一定理把形式平滑性归结为同一问题在域上代数这一情形中的研究。

\medskip\noindent\textbf{II.} --- 设 $k$ 是一个域，$A$ 是一个作为 $k$-代数的
Noether 局部环。于是，$A$ 形式平滑，当且仅当 $A$ 在 $k$ 上\emph{几何正则}，
也就是说，对每个\emph{有限}扩域 $k'$（它是 $k$ 的扩域），半局部环
$A\otimes_k k'$ 都是

\oldpage[0\textsubscript{IV-1}]{70}

\noindent 正则的 (22.5.8)；此时，\emph{完备化} $\widehat A$（属于 $A$）同构于
一个形式幂级数环 $K[[T_1,\ldots,T_n]]$ (19.6.5)。此外，$k$-代数 $A$ 的结构
（这里假设 $A$ \emph{完备}且在 $k$ 上形式平滑）完全由剩余域 $K$（属于 $A$）
以及 $A$ 的\emph{维数}确定；只要满足不等式
\[
  \dim(A)\geq \operatorname{rg}_k(\Gamma_{K/k}),
\]
后者可以任意取值，其中 $\Gamma_{K/k}$ 是 $K$ 的“不完美模” (22.2.6)。

特别地，$K$ 作为 $k$ 的一个\emph{扩域}是形式平滑的，当且仅当 $K$ 是 $k$ 的
一个\emph{可分}扩域 (19.6.1)。

\medskip\noindent\textbf{III.} --- 设 $A$ 是一个 Noether 局部环，
$\mathfrak J$ 是 $A$ 的一个不等于 $A$ 的理想，$A_0=A/\mathfrak J$，而
$B_0$ 是一个\emph{完备} Noether 局部环；设 $A_0\to B_0$ 是一个局部同态，
使 $B_0$ 成为形式平滑的 $A_0$-代数。于是，存在一个\emph{完备} Noether
局部环 $B$、一个局部同态 $A\to B$，使 $B$ 成为平坦的 $A$-模，并且存在
一个 $A_0$-\emph{同构}
\[
  B\otimes_A A_0\xrightarrow{\sim}B_0
\]
（因此，由 I，$B$ 是形式平滑的 $A$-代数）；此外，这些性质在相差一个同构的
意义下确定 $B$ (19.7.2)。这个定理尤其包含关于\emph{完备} Noether 局部环
结构的\emph{Cohen 定理} (19.8)，这些定理将在第四章 §§~6 和 7 中起重要作用。

\medskip\noindent\textbf{IV.} --- 研究一个 Noether 局部环相对于另一个环形式
光滑的意义，来自光滑态射的如下“逐点”刻画：若 $X$ 和 $Y$ 是局部 Noether
预概形，$f:X\to Y$ 是局部有限型态射，则 $f$ 光滑，当且仅当对每个
$x\in X$，环 $\mathcal O_x$ 在 $\mathcal O_{f(x)}$ 上形式平滑。特别地，若
$Y=\operatorname{Spec}(k)$，其中 $k$ 是完美域，则说 $f$ 光滑，等价于说 $X$
是一个\emph{正则}预概形。

\medskip\noindent\textbf{V.} --- 最后，我们将在 §§~20、21 和 22 中看到，形式
光滑代数的概念自然地出现在\emph{Kähler 微分}理论中，而且这两个理论彼此阐明。
\end{env}

\begin{env}[19.0.3]
\label{0.19.0.3-fr}
在本节以及以下各节中，总假设拓扑环和拓扑模是\emph{线性拓扑化的}
($\mathbf{0}_{\mathrm I}$, 7.1.1)；除非另有明言，所考虑的拓扑环总假设为
\emph{交换}的。回顾：若 $A$ 和 $B$ 是两个拓扑环，$\rho:A\to B$ 是一个环
同态，并在 $B$ 上定义一个 $A$-代数结构，则称 $B$ 是一个\emph{拓扑
$A$-代数}，如果 $\rho$ 对所考虑的拓扑是\emph{连续}的。

为简便起见，在拓扑环 $A$（相应地，在拓扑 $A$-模 $M$）中，我们说“开理想
（相应地，开子模）的\emph{邻域基}”，而不说“由理想（相应地，由子模）组成的
零点邻域基”。

给定一个拓扑环 $A$ 和一个 $A$-模 $M$，诸集合 $\mathfrak J M$（其中
$\mathfrak J$ 遍历开理想的一个邻域基）构成 $M$ 上某个拓扑的开子模邻域基；
该拓扑使 $M$ 成为拓扑 $A$-模，并称为从 $A$ 的拓扑\emph{导出}的拓扑。

设 $M$ 是一个拓扑 $A$-模，其拓扑比从 $A$ 的拓扑导出的拓扑\emph{更粗}。
于是，若 $N$ 是 $M$ 的一个\emph{开}子模，则离散 $A$-模 $M/N$ 被 $A$ 的
一个开理想\emph{零化}，因为由假设存在这样的理想 $\mathfrak R$，满足
$\mathfrak R.M\subset N$。

\oldpage[0\textsubscript{IV-1}]{71}

若 $M$ 和 $N$ 是两个拓扑 $A$-模，而且它们的拓扑都从 $A$ 的拓扑导出，则
\emph{每个} $A$-同态 $u:M\to N$ 都是\emph{连续}的；事实上，对每个零点邻域
$V$（位于 $N$ 中），由假设存在一个开理想 $\mathfrak J$（属于 $A$），使
$\mathfrak JN\subset V$，从而 $u(\mathfrak JM)\subset\mathfrak JN\subset V$。

当 $B$ 是一个拓扑 $A$-代数时，$B$ 上从 $A$ 的拓扑导出的拓扑比给定拓扑
\emph{更细}；事实上，对每个开理想 $\mathfrak R$（属于 $B$），由假设存在
一个开理想 $\mathfrak J$（属于 $A$），使 $\mathfrak JB\subset\mathfrak R$。
\end{env}

\subsection{形式满态射与形式单态射}
\label{subsection:0.19.1-fr}

\begin{proposition}[19.1.1]
\label{0.19.1.1-fr}
设 $A$ 是一个拓扑环，$M$、$N$ 是两个拓扑 $A$-模，$(W_\lambda)$ 是 $N$
中开子模的一个邻域基，$u:M\to N$ 是一个连续 $A$-同态，$\widehat M$、
$\widehat N$ 分别是 $M$、$N$ 的分离完备化，而
$\widehat u:\widehat M\to\widehat N$ 是 $u$ 在这些分离完备化上的连续延拓。
\begin{enumerate}
  \item[(i)] 下列条件等价：
  \begin{enumerate}
    \item[a)] $u(M)$ 在 $N$ 中\emph{稠密}。
    \item[a$'$)] $\widehat u(\widehat M)$ 在 $\widehat N$ 中\emph{稠密}。
    \item[a$''$)] 对每个 $\lambda$，复合同态
    $M\xrightarrow{u}N\to N/W_\lambda$ 都是\emph{满射}。
  \end{enumerate}
  \item[(ii)] 下列条件等价：
  \begin{enumerate}
    \item[b)] 由 $u$ 从 $N$ 的拓扑拉回所得的拓扑等于 $M$ 的拓扑。
    \item[b$'$)] $\widehat u$ 是从拓扑 $\widehat A$-模 $\widehat M$ 到拓扑
    $\widehat A$-子模 $\widehat u(\widehat M)$ 上的一个\emph{同构}；后者是
    $\widehat N$ 的子模（且必为\emph{闭}子模）。
    \item[b$''$)] 各 $u^{-1}(W_\lambda)$ 构成 $M$ 中零点邻域的一个邻域基。
  \end{enumerate}
\end{enumerate}
这立即由分离完备化的定义以及 $N/W_\lambda$ 为离散模这一事实得出。
\end{proposition}

\begin{env}[19.1.2]
\label{0.19.1.2-fr}
当 (19.1.1) 中 (i)（相应地，(ii)）的等价条件成立时，称 $u$ 为一个
\emph{形式满态射}（相应地，一个\emph{形式单态射}）。若它同时是形式
单态射和形式满态射，则称 $u$ 为一个\emph{形式双态射}；由 (19.1.1)，这等价于
说 $\widehat u$ 是从拓扑 $\widehat A$-模 $\widehat M$ 到拓扑
$\widehat A$-模 $\widehat N$ 上的一个\emph{同构}。
\end{env}

\begin{proposition}[19.1.3]
\label{0.19.1.3-fr}
设 $A$ 是一个拓扑环，$M$、$N$ 是两个拓扑 $A$-模，$u:M\to N$ 是一个
连续 $A$-同态。假设存在开子模的两个邻域基 $(V_\lambda)$、$(W_\lambda)$，
它们分别位于 $M$ 和 $N$ 中，具有同一个指标集，并且
$u(V_\lambda)\subset W_\lambda$ 对每个 $\lambda$ 都成立；设
$u_\lambda:M/V_\lambda\to N/W_\lambda$ 是由 $u$ 取商所得的同态。则：
\begin{enumerate}
  \item[(i)] $u$ 为形式满态射，当且仅当对每个 $\lambda$，$u_\lambda$ 都是
  满射。
  \item[(ii)] $u$ 为形式单态射，当且仅当对每个 $\lambda$，存在一个 $\mu$，
  使 $V_\mu\subset V_\lambda$，并且 $\operatorname{Ker}(u_\mu)$ 包含于核
  $V_\lambda/V_\mu$ 中；这个核属于典范映射 $M/V_\mu\to M/V_\lambda$。
  \item[(iii)] $u$ 为形式双态射，当且仅当对每个 $\lambda$，$u_\lambda$ 都是
  满射，并且存在一个 $\mu$，使 $V_\mu\subset V_\lambda$，而典范同态
  $M/V_\mu\to M/V_\lambda$

  \oldpage[0\textsubscript{IV-1}]{72}

  \noindent 可分解为
  $M/V_\mu\xrightarrow{u_\mu}N/W_\mu\to M/V_\lambda$（其中同态
  $N/W_\mu\to M/V_\lambda$ 必唯一）。
\end{enumerate}
\end{proposition}

\begin{proof}
断言 (i) 是立即的；另一方面，
$\operatorname{Ker}(u_\mu)=u^{-1}(W_\mu)/V_\mu$，而典范映射
$M/V_\mu\to M/V_\lambda$ 的核是 $V_\lambda/V_\mu$；因此，(ii) 的条件等价于
$u^{-1}(W_\mu)\subset V_\lambda$，断言 (ii) 随即得出；最后，当 $u_\mu$ 是
满射时，说 $\operatorname{Ker}(u_\mu)$ 包含于 $V_\lambda/V_\mu$，等价于说
$M/V_\mu\to M/V_\lambda$ 可分解为 $v\circ u_\mu$，其中 $v$ 是一个同态
$N/W_\mu\to M/V_\lambda$，因为这时 $N/W_\mu$ 可与
$(M/V_\mu)/\operatorname{Ker}(u_\mu)$ 等同。
\end{proof}

\begin{corollary}[19.1.4]
\label{0.19.1.4-fr}
设 $A$ 是一个拓扑环，$M$、$N$ 是两个其拓扑都从 $A$ 的拓扑导出的拓扑
$A$-模，而 $u:M\to N$ 是一个形式满态射。设 $(\mathfrak J_\lambda)$ 是 $A$
中开理想的一个邻域基。则 $u$ 为形式双态射，当且仅当对每个 $\lambda$，同态
$u_\lambda:M/\mathfrak J_\lambda M\to N/\mathfrak J_\lambda N$ 都是双射；这里
该同态由 $u$ 取商得到。
\end{corollary}

\begin{proof}
事实上，$u(\mathfrak J_\lambda M)\subset\mathfrak J_\lambda N$，故可应用判据
(19.1.3, (iii))；但若有一个分解
\[
  M/\mathfrak J_\mu M\xrightarrow{u_\mu}
  N/\mathfrak J_\mu N\xrightarrow{v}M/\mathfrak J_\lambda M,
\]
则 $\mathfrak J_\lambda M/\mathfrak J_\mu M$ 在 $u_\mu$ 下的像是
$\mathfrak J_\lambda N/\mathfrak J_\mu N$，因而
$\mathfrak J_\lambda N/\mathfrak J_\mu N$ 在 $v$ 下的像为 0，并且 $v$ 可分解为
\[
  N/\mathfrak J_\mu N\longrightarrow N/\mathfrak J_\lambda N
  \xrightarrow{v'}M/\mathfrak J_\lambda M~;
\]
但此时 $M/\mathfrak J_\lambda M$ 的恒等自同态可分解为
\[
  M/\mathfrak J_\lambda M\xrightarrow{u_\lambda}
  N/\mathfrak J_\lambda N\xrightarrow{v'}M/\mathfrak J_\lambda M,
\]
这表明 $u_\lambda$ 是单射；而我们已经知道它是满射，故推论得证。
\end{proof}

\begin{proposition}[19.1.5]
\label{0.19.1.5-fr}
设 $A$ 是一个拓扑环，$M$、$N$ 是两个拓扑 $A$-模，而 $u:M\to N$ 是一个
连续 $A$-同态。下列条件等价：
\begin{enumerate}
  \item[a)] 对每个离散拓扑 $A$-模 $E$ 和每个连续 $A$-同态
  $v:M\to E$，存在一个连续 $A$-同态 $w:N\to E$，使 $v=w\circ u$。
  \item[b)] 对每个开子模 $V'$（属于 $M$），存在一个开子模 $W''$（属于 $N$）、
  一个开子模 $V''\subset V'\cap u^{-1}(W'')$ 以及一个同态
  $h:N/W''\to M/V'$，使典范同态 $M/V''\to M/V'$ 可分解为
  \[
    M/V''\xrightarrow{u''}N/W''\xrightarrow{h}M/V',
  \]
  其中 $u''$ 由 $u$ 取商得到。
\end{enumerate}
\end{proposition}

\begin{proof}
要证明 a) 蕴涵 b)，只须把 a) 应用于 $E=M/V'$ 和典范同态
$v:M\to M/V'$；于是 $W''=w^{-1}(0)$ 是 $N$ 的开子模，而
$V''=u^{-1}(W'')$ 是 $M$ 的一个包含于 $V'$ 的开子模；这些子模以及同态
$h:N/W''\to M/V'$（由 $w$ 取商得到）即满足要求。反过来，为证明 b) 蕴涵
a)，只须考虑 $v$ 为满射的情形；这是把 $E$ 换成 $v(M)$ 即可做到。于是
$V'=\operatorname{Ker}(v)$ 是 $M$ 的一个开子模，故 $E$ 同构于离散 $A$-模
$M/V'$；应用 b)，并取复合

\oldpage[0\textsubscript{IV-1}]{73}

\noindent 所求的连续 $A$-同态 $w$ 可取为复合
$N\to N/W''\xrightarrow{h}M/V'$；
这里交换图为
\[
  \xymatrix{
    M\ar[r]^u\ar[d] & N\ar[d]\\
    M/V''\ar[r]_{u''} & N/W''
  }
\]
。
\end{proof}

当 (19.1.5) 的等价条件成立时，我们称 $u$ \emph{形式左可逆}；由于 (19.1.5)
中的条件 b) 蕴涵 $u^{-1}(W'')\subset V'$，此时 $u$ 是一个形式单态射。下列
推论进一步说明了这一术语的理由：

\begin{corollary}[19.1.6]
\label{0.19.1.6-fr}
若存在一个连续 $\widehat A$-同态 $s:\widehat N\to\widehat M$，使
$s\circ\widehat u=1_{\widehat M}$，则 $u$ 形式左可逆。
\end{corollary}

注意，此时 $\widehat u$ 是从 $\widehat M$ 到 $\widehat N$ 的一个
\emph{拓扑直因子}上的拓扑同构。为证明 (19.1.6)，只须注意：沿用
(19.1.5, a)) 的记号，$v:M\to E$ 延拓成一个连续 $\widehat A$-同态
$\widehat v:\widehat M\to E$，因为 $E$ 是离散的；设
$j:M\to\widehat M$ 和 $j':N\to\widehat N$ 为典范同态；于是，令
$w=\widehat v\circ s\circ j'$，就有
$w\circ u=\widehat v\circ s\circ\widehat u\circ j=\widehat v\circ j=v$。

\begin{corollary}[19.1.7]
\label{0.19.1.7-fr}
假设 $M$ 和 $N$ 的拓扑都从 $A$ 的拓扑导出，并设 $(\mathfrak J_\lambda)$ 是
$A$ 中开理想的一个邻域基。则 $u$ 形式左可逆，当且仅当对每个 $\lambda$，同态
$u_\lambda:M/\mathfrak J_\lambda M\to N/\mathfrak J_\lambda N$（由 $u$ 取商
得到）左可逆（换言之，
它是从 $M/\mathfrak J_\lambda M$ 到 $N/\mathfrak J_\lambda N$ 的一个直因子
上的同构）。
\end{corollary}

事实上，该条件是充分的，因为在 (19.1.5) 的条件 b) 中取
$V'=\mathfrak J_\lambda M$，再取 $W''=\mathfrak J_\lambda N$、$V''=V'$，并取
$h$ 使 $h\circ u_\lambda$ 为 $M/\mathfrak J_\lambda M$ 上的恒等映射，即满足
要求。反过来，若 $u$ 形式左可逆，则由 (19.1.5, b))，对每个 $\lambda$，
存在一个 $\mu$，使 $\mathfrak J_\mu\subset\mathfrak J_\lambda$，并存在同态
$h:N/\mathfrak J_\mu N\to M/\mathfrak J_\lambda M$，使典范同态
$M/\mathfrak J_\mu M\to M/\mathfrak J_\lambda M$ 可分解为
\[
  M/\mathfrak J_\mu M\xrightarrow{u_\mu}N/\mathfrak J_\mu N
  \xrightarrow{h}M/\mathfrak J_\lambda M~;
\]
但由于
$h(\mathfrak J_\lambda(N/\mathfrak J_\mu N))
\subset\mathfrak J_\lambda(M/\mathfrak J_\lambda M)=0$，$h$ 典范地分解为
\[
  N/\mathfrak J_\mu N\longrightarrow N/\mathfrak J_\lambda N
  \xrightarrow{s}M/\mathfrak J_\lambda M,
\]
而 $s$ 显然是 $u_\lambda$ 的左逆。

\begin{proposition}[19.1.8]
\label{0.19.1.8-fr}
设 $A$ 是一个拓扑环，它有一个由开理想组成的可数递减邻域基
$(\mathfrak J_n)$。设 $M$、$N$ 是两个其拓扑都从 $A$ 的拓扑导出的拓扑
$A$-模，$u:M\to N$ 是一个 $A$-同态。置
$M_n=M/\mathfrak J_nM$、$N_n=N/\mathfrak J_nN$，并令
$u_n:M_n\to N_n$ 为一个 $(A/\mathfrak J_n)$-同态；它由 $u$ 取商得到。假设对每个
$n$，$P_n=\operatorname{Coker}(u_n)$ 都是一个投射 $(A/\mathfrak J_n)$-模，
并且 $M$ 是分离且完备的。则 $u$ 形式左可逆，当且仅当 $u$ 左可逆（而且此时
$u$ 是从 $M$ 到 $N$ 的一个拓扑直因子上的拓扑同构）。
\end{proposition}

\oldpage[0\textsubscript{IV-1}]{74}

\begin{proof}
由 (19.1.7)，有交换图
\[
  \xymatrix{
    0\ar[r] & M_n\ar[r]^{u_n} & N_n\ar[r]^{p_n} & P_n\ar[r] & 0\\
    0\ar[r] & M_{n+1}\ar[u]^f\ar[r]_{u_{n+1}} &
      N_{n+1}\ar[u]^g\ar[r]_{p_{n+1}} & P_{n+1}\ar[u]^h\ar[r] & 0
  }
\]
其中各行都是\emph{分裂}正合列；换言之，对每个 $n$，存在一个同态
$s_n:P_n\to N_n$，使 $p_n\circ s_n=1_{P_n}$。我们将证明，可以对 $n$
作归纳而定义一个同态 $s'_n:P_n\to N_n$，使
$p_n\circ s'_n=1_{P_n}$，并使各图
\[
  \xymatrix{
    P_n\ar[r]^{s'_n} & N_n\\
    P_{n+1}\ar[u]^h\ar[r]_{s'_{n+1}} & N_{n+1}\ar[u]^g
  }
\]
交换。事实上，$g\circ s_{n+1}-s'_n\circ h$ 是从 $P_{n+1}$ 到
\[
  u_n(M_n)=u_{n+1}(M_{n+1})/
  \mathfrak J_nu_{n+1}(M_{n+1})~;
\]
的一个同态；由于 $P_{n+1}$ 是一个\emph{投射} $(A/\mathfrak J_{n+1})$-模，
这个同态可分解为
\[
  P_{n+1}\xrightarrow{t_{n+1}}u_{n+1}(M_{n+1})
  \longrightarrow
  u_{n+1}(M_{n+1})/\mathfrak J_nu_{n+1}(M_{n+1}),
\]
由此立即得出 $s'_{n+1}=s_{n+1}-t_{n+1}$ 满足要求。这样，由 $N_n$ 分解为
$u_n(M_n)$ 与 $s'_n(P_n)$ 的直和，立即得到一个同态 $w_n:N_n\to M_n$；
它是 $u_n$ 的左逆，并使各图
\[
  \xymatrix{
    N_n\ar[r]^{w_n} & M_n\\
    N_{n+1}\ar[u]^g\ar[r]_{w_{n+1}} & M_{n+1}\ar[u]^f
  }
\]
交换。于是逆系统 $(w_n)$ 有逆极限 $\widehat w:\widehat N\to\widehat M=M$；
把它与典范同态 $N\to\widehat N$ 复合，得到一个同态 $w:N\to M$，使得对
每个 $n$，自同态 $(w\circ u)_n$（属于 $M_n=M/\mathfrak J_nM$，并由
$w\circ u$ 取商得到）都是恒等映射；这蕴涵 $w\circ u=1_M$，因为 $M$ 分离且
完备。
\end{proof}

\begin{proposition}[19.1.9]
\label{0.19.1.9-fr}
设 $A$ 是一个预可容拓扑环（$\mathbf{0}_{\mathrm I}$, 7.1.2），
$\mathfrak Q$ 是 $A$ 的一个定义理想，$(\mathfrak J_\lambda)$ 是 $A$ 中开理想
的一个邻域基。设 $M$、$N$ 是两个其拓扑都从 $A$ 的拓扑导出的拓扑 $A$-模，
并且对每个 $\lambda$，$N/\mathfrak J_\lambda N$ 都是一个投射
$(A/\mathfrak J_\lambda)$-模（参见 (19.2.3)）。设 $u:M\to N$ 是一个
$A$-同态。则下列条件等价：
\begin{enumerate}
  \item[a)] $u$ 形式左可逆。

  \oldpage[0\textsubscript{IV-1}]{75}

  \item[b)] 同态 $u_0:M/\mathfrak QM\to N/\mathfrak QN$（由 $u$ 取商得到）
  左可逆。
\end{enumerate}
\end{proposition}

\begin{proof}
我们在 (19.1.7) 中已经看到，条件 a) 等价于 $u_\lambda$ 对每个 $\lambda$
都左可逆；由于 $\mathfrak Q$ 是一个开理想，因而包含某个
$\mathfrak J_\lambda$，立即得出 $u_0$ 左可逆。反过来，为证明 b) 蕴涵 a)，
注意，对每个 $\lambda$，$\mathfrak Q/\mathfrak J_\lambda$ 依假设都是
$A/\mathfrak J_\lambda$ 的一个幂零理想。我们的断言将由下列命题得出。
\end{proof}

\begin{proposition}[19.1.10]
\label{0.19.1.10-fr}
设 $A$ 是一个环，$M$、$N$ 是两个 $A$-模，其中 $N$ 是投射模，而
$u:M\to N$ 是一个 $A$-同态。设 $\mathfrak J$ 是 $A$ 的一个理想；假设下列
条件之一成立：
\begin{enumerate}
  \item[(i)] $\mathfrak J$ 是幂零的。
  \item[(ii)] $\mathfrak J$ 包含于 $A$ 的根中，并且 $M$ 是有限型的。
\end{enumerate}
则 $u$ 左可逆，当且仅当同态 $u_0:M/\mathfrak JM\to N/\mathfrak JN$ 作为
$(A/\mathfrak J)$-模同态左可逆；该同态由 $u$ 取商得到。
\end{proposition}

\begin{proof}
该条件显然是必要的；下面证明它是充分的。设 $v_0$ 是 $u_0$ 的一个左逆；
复合同态 $N\to N/\mathfrak JN\xrightarrow{v_0}M/\mathfrak JM$ 可分解为
$N\xrightarrow{v}M\to M/\mathfrak JM$，因为 $N$ 是投射模；于是
$w=v\circ u$ 是 $M$ 的一个自同态，使得自同态 $w_0$（属于
$M/\mathfrak JM$，并由取商得到）是恒等映射；只须证明 $w$ 本身是双射
（因为这时 $w^{-1}\circ v$ 将是 $u$ 的一个左逆）。现在分别考虑两种情形。

(i) 对每个 $n$，有交换图
\[
  \xymatrix{
    (\mathfrak J^n/\mathfrak J^{n+1})\otimes_{A/\mathfrak J}
      (M/\mathfrak JM)\ar[r]\ar[d]_{1\otimes\operatorname{gr}^0(w)} &
      \mathfrak J^nM/\mathfrak J^{n+1}M\ar[d]^{\operatorname{gr}^n(w)}\\
    (\mathfrak J^n/\mathfrak J^{n+1})\otimes_{A/\mathfrak J}
      (M/\mathfrak JM)\ar[r] & \mathfrak J^nM/\mathfrak J^{n+1}M
  }
\]
其中水平箭头都是\emph{满射}；由于 $\operatorname{gr}^0(w)=w_0$ 是恒等映射，
$\operatorname{gr}^n(w)$ 也如此，因而它当然是双射。由于 $\mathfrak J$-预进制
滤过（位于 $M$ 上）是有限的（因为 $\mathfrak J$ 是幂零的），由此得出
$w$ 是双射（Bourbaki，\emph{Alg. comm.}，chap.~III，§~2，
n\textsuperscript{o}~8，th.~1 的 cor.~3）。

(ii) 只须证明，对每个极大理想 $\mathfrak m$（属于 $A$），自同态
$w_{\mathfrak m}$（属于 $M_{\mathfrak m}$）都是双射（Bourbaki，
\emph{Alg. comm.}，chap.~II，
§~3，n\textsuperscript{o}~3，th.~1）；由于
$\mathfrak JA_{\mathfrak m}\subset\mathfrak mA_{\mathfrak m}$ 且
$A_{\mathfrak m}/\mathfrak JA_{\mathfrak m}=(A/\mathfrak J)_{\mathfrak m}$，
问题归结为 $A$ 是局部环的情形。此外，可以假设 $\mathfrak J$ 就是极大理想
$\mathfrak J_0$（属于 $A$），因为若 $u_0$ 左可逆，

\oldpage[0\textsubscript{IV-1}]{76}

则同态 $u_{00}:M/\mathfrak J_0M\to N/\mathfrak J_0N$ 也左可逆；该同态由
$u_0$ 与 $1_{A/\mathfrak J_0}$ 作张量积得到，因为
\[
  M/\mathfrak J_0M=(M/\mathfrak JM)\otimes_{A/\mathfrak J}(A/\mathfrak J_0).
\]
\SourceCorrectionNote{法文印本及外交式转录把 u00 的陪域写成 N/J0M；但由 u0 与 A/J0 作张量积所得的陪域必为 N/J0N。此处采用后一读法；英文比较本也采用该 读法并明确标注法文误植。}
因此，假设 $\mathfrak J$ 是极大理想，从而 $A/\mathfrak J$ 是一个域。显然，
只须证明在命题的条件下 $M$ 是一个\emph{自由} $A$-模；事实上，
$\det(w_0)$ 这时是 $\det(w)$ 在 $A/\mathfrak J$ 中的典范像，故
$\det(w)$ 不属于理想 $\mathfrak J$，因而可逆。另一方面，由于
$(A/\mathfrak J)$-向量空间 $M/\mathfrak JM$ 是有限维自由的，存在一个有限型
自由 $A$-模 $L$ 和一个 $A$-同态 $f:L\to M$，使同态
$f_0:L/\mathfrak JL\to M/\mathfrak JM$（由 $f$ 取商得到）是双射。
\SourceCorrectionNote{法文印本及外交式转录此处只说 L 为有限型；应为有限型自由模，因为构造提升的是剩余向量空间的 一组基，而且下文明确使用 L 的自由性。}
由于 $M$ 是有限型的，首先由 Nakayama 引理得出 $f$ 是满射（Bourbaki，
\emph{Alg. comm.}，chap.~II，§~3，n\textsuperscript{o}~2，prop.~4 的
cor.~1）；此外，若 $g=u\circ f$，则取商所得的同态
$g_0:L/\mathfrak JL\to N/\mathfrak JN$ 左可逆；而这里 $L$ 是自由的，开头的
说明便证明 $g$ 本身左可逆；显然，这蕴涵 $f$ 是\emph{单射}，证明完毕。
\end{proof}

顺便指出 (19.1.10) 的下列推论：

\begin{corollary}[19.1.11]
\label{0.19.1.11-fr}
设 $A$ 是一个局部环，$k$ 是其剩余域，$M$ 是一个有限型 $A$-模，$N$ 是一个
投射 $A$-模，而 $u:M\to N$ 是一个同态。则 $u$ 左可逆，当且仅当存在一个
生成元系 $(x_i)_{1\leq i\leq m}$（属于 $M$），使同态
$u\otimes 1:M\otimes_Ak\to N\otimes_Ak$ 对各 $x_i\otimes1$ 的像，在
$N\otimes_Ak$ 中线性
无关；此时各 $x_i$ 构成 $M$ 的一组基。
\end{corollary}

该条件显然是必要的，因为若 $u$ 左可逆，则 $M$ 是有限型投射 $A$-模，因而
自由。反过来，若该条件成立，则各 $x_i\otimes1$ 显然构成 $M\otimes_Ak$
的一组基，并且 $u\otimes1$ 左可逆；所以只须把 (19.1.10) 应用于极大理想
$\mathfrak J$（属于 $A$）。

\begin{corollary}[19.1.12]
\label{0.19.1.12-fr}
设 $A$ 是一个环，$M$ 是一个有限型 $A$-模，$N$ 是一个投射 $A$-模，而
$u:M\to N$ 是一个同态。对每个素理想
$\mathfrak p\in\operatorname{Spec}(A)$，下列条件等价：
\begin{enumerate}
  \item[a)] 同态 $u_{\mathfrak p}:M_{\mathfrak p}\to N_{\mathfrak p}$ 左可逆。
  \item[b)] 同态
  $u\otimes1:M\otimes_Ak(\mathfrak p)\to N\otimes_Ak(\mathfrak p)$ 是单射
  （回顾 $k(\mathfrak p)$ 表示 $A$ 在 $\mathfrak p$ 处的剩余域）。
  \item[c)] 存在有限个元素 $x_i\in M$ $(1\leq i\leq m)$，使各 $x_i$ 在
  $M_{\mathfrak p}$ 中的像生成 $M_{\mathfrak p}$；还存在 $m$ 个线性形式
  $y_i^*$（定义在 $N$ 上）$(1\leq i\leq m)$，使
  $\det(\langle y_i^*,u(x_j)\rangle)\notin\mathfrak p$。
  \item[d)] 存在 $f\in A-\mathfrak p$，使同态 $u_f:M_f\to N_f$ 左可逆。
\end{enumerate}
此外，满足这些条件的 $\mathfrak p\in\operatorname{Spec}(A)$ 所成的集合是
$\operatorname{Spec}(A)$ 的开集。
\end{corollary}

最后一个断言由 d) 立即得出。
\SourceCorrectionNote{法文印本及外交式转录写作“最后一个条件”，这会错误地把 d) 说成由自身推出；此处指列举后 关于开集的“最后一个断言”。}
由于 $N$ 是投射模，它是一个自由 $A$-模
$A^{(I)}$ 的直因子；又由于 $M$ 是有限型的，$u(M)$ 包含在 $A^{(I)}$ 的某个
形如 $A^n$ 的子模中，其中 $n$ 有限；由于 a)、b)、c)、d) 中的每个断言都
等价于把 $N$ 换成 $N\oplus P$ 后的相应断言（仍假设 $u$ 把 $M$ 映到 $N$
中），问题归结为 $N$ 是\emph{有限型自由模}的情形。d) 蕴涵 a) 是立即的，
而由 (19.1.11)，a) 与 b) 等价；此外，a) 蕴涵 $M_{\mathfrak p}$ 是自由的
(19.1.11)，所以 a) 蕴涵 c)：取一族 $x_i$，使它们在 $M_{\mathfrak p}$ 中的
像构成这个

\oldpage[0\textsubscript{IV-1}]{77}

$A_{\mathfrak p}$-模的一组基，并注意（因为 $N$ 是有限型自由模）
$A_{\mathfrak p}$-模 $N_{\mathfrak p}$ 上的每个线性形式都可写成 $u=v/f$，
其中 $f\in A-\mathfrak p$，而 $v$ 是 $A$-模 $N$ 上的一个线性形式。显然，
c) 蕴涵 b)，故只须证明 a) 蕴涵 d)。由于 $N$ 是有限表示的，
$(\operatorname{Hom}_A(N,M))_{\mathfrak p}$ 典范地等同于
$\operatorname{Hom}_{A_{\mathfrak p}}(N_{\mathfrak p},M_{\mathfrak p})$
（Bourbaki，\emph{Alg. comm.}，chap.~II，§~2，n\textsuperscript{o}~7，
prop.~19）。若 $w_{\mathfrak p}$ 是 $u_{\mathfrak p}$ 的左逆，则存在同态
$w:N\to M$ 和元素 $f\in A-\mathfrak p$，使
$w_{\mathfrak p}=w\otimes(1/f)1_{A_{\mathfrak p}}$。
\SourceCorrectionNote{法文印本及外交式转录把上式两边倒置；应由全局同态局部化并缩放而得到局部左逆。}
因此，等式 $w_{\mathfrak p}\circ u_{\mathfrak p}=1_{M_{\mathfrak p}}$ 也可写成
$(w\circ u)\otimes1_{A_{\mathfrak p}}=f.1_{M_{\mathfrak p}}$。但由于 $M$
是有限型 $A$-模，存在 $g\in A-\mathfrak p$，使自同态
$g((w\circ u)-f.1_M)$ 在 $M$ 的所有生成元上为零，因而本身为零。若置
$h=fg$，以及 $u_h=u\otimes1_{A_h}$、$w_h=w\otimes1_{A_h}$，则有
$w_h(u_h(z))=(f/1)z$ 对每个 $z\in M_h$ 成立；但由于 $h/1$ 在 $A_h$ 中可逆，
$f/1=f'$ 也可逆，故 $f'^{-1}w_h$ 是 $u_h$ 的左逆，证明完毕。

\begin{proposition}[19.1.13]
\label{0.19.1.13-fr}
假设 (19.1.9) 的各项假设成立，并且进一步假设 $(\mathfrak J_\lambda)$ 是一个
递减序列 $(\mathfrak J_n)$，而 $M$ 是分离且完备的。则 (19.1.9) 的条件 a)
和 b) 还等价于：
\begin{enumerate}
  \item[c)] $u$ 左可逆。
\end{enumerate}
\end{proposition}

我们已经知道 c) 蕴涵 a) (19.1.6)。反过来，若 a) 成立，则已知（沿用
(19.1.8) 的记号）$M_n$ 是投射 $(A/\mathfrak J_n)$-模 $N_n$ 的一个直因子，
故 $P_n$ 也同构于 $N_n$ 的一个直因子，因而是投射模；所以只须应用
(19.1.8)。

最后指出下列命题：

\begin{proposition}[19.1.14]
\label{0.19.1.14-fr}
设 $A$ 是一个环，$M$ 是一个有限型 $A$-模，$N$ 是一个投射 $A$-模，而
$u:M\to N$ 是一个 $A$-同态。
\begin{enumerate}
  \item[(i)] $u$ 左可逆，当且仅当对每个极大理想 $\mathfrak m$（属于 $A$），
  $u_{\mathfrak m}:M_{\mathfrak m}\to N_{\mathfrak m}$ 都左可逆。
  \item[(ii)] 设 $A'$ 是一个 $A$-代数，并且作为 $A$-模忠实平坦。则 $u$
  左可逆，当且仅当
  $u\otimes1:M\otimes_AA'\to N\otimes_AA'$ 左可逆。
\end{enumerate}
\end{proposition}

与 (19.1.12) 中一样，可以只考虑 $N$ 是有限型自由模的情形；说 $u$ 左可逆，
就意味着 $u$ 是单射，而且商模 $P=N/u(M)$ 是\emph{投射}的，因为这时
$u(M)$ 将是 $N$ 的一个直因子。还要注意，由于 $M$ 是有限型的，$P$ 是
\emph{有限表示}的。于是：

(i) 该条件显然是必要的。反过来，若该条件成立，则已知 $u$ 是单射
（Bourbaki，\emph{Alg. comm.}，chap.~II，§~3，n\textsuperscript{o}~3，
th.~1）；而
$P_{\mathfrak m}=N_{\mathfrak m}/u_{\mathfrak m}(M_{\mathfrak m})$ 对每个
$\mathfrak m$ 都是投射的，
故已知这蕴涵 $P$ 是投射的（\emph{loc. cit.}，§~5，
n\textsuperscript{o}~2，th.~1）。

(ii) 这里，该条件仍显然是必要的。反过来，若它成立，则已知 $u$ 是单射
（$\mathbf{0}_{\mathrm I}$, 6.4.1）；又由于
$P\otimes_AA'=\operatorname{Coker}(u\otimes1)$ 是投射的，因而是平坦的，
故得出 $P$ 是一个平坦 $A$-模（$\mathbf{0}_{\mathrm I}$, 6.6.3），并因其
有限表示而为投射模（Bourbaki，\emph{Alg. comm.}，chap.~II，§~5，
n\textsuperscript{o}~2，th.~1 的 cor.~2）。

\oldpage[0\textsubscript{IV-1}]{78}

\begin{remark}[19.1.15]
\label{0.19.1.15-fr}
形式左可逆同态的“对偶”概念是\emph{形式右可逆}同态；按定义，这样的连续
$A$-同态 $u:M\to N$ 满足下列条件：对每个开子模 $W'$（属于 $N$），存在
一个开子模 $V'\subset u^{-1}(W')$（属于 $M$）、一个开子模
$W''\subset W'$（属于 $N$）以及
一个同态 $h:N/W''\to M/V'$，使典范同态 $N/W''\to N/W'$ 可分解为
\[
  N/W''\xrightarrow{h}M/V'\xrightarrow{u'}N/W',
\]
其中 $u'$ 由 $u$ 取商得到。这蕴涵 $u$ 是一个形式满态射；若存在一个连续
$\widehat A$-同态 $r:\widehat N\to\widehat M$，使
$\widehat u\circ r=1_{\widehat N}$，则立即验证 $u$ 形式右可逆。若 (19.1.7)
的条件成立，则 $u$ 形式右可逆，当且仅当各 $u_\lambda$ 都右可逆（也就是说，
$u_\lambda$ 的核是 $M/\mathfrak J_\lambda M$ 的一个直因子，而且
$u_\lambda$ 是从 $\operatorname{Ker}(u_\lambda)$ 的一个补模到
$N/\mathfrak J_\lambda N$ 上的同构）。我们把陈述并证明与 (19.1.8) 和
(19.1.9) 相对应的命题留给读者（在 (19.1.8) 的类似命题中，必须假设 $M$
分离且完备，并假设 $M_n$ 是一个投射 $(A/\mathfrak J_n)$-模；在 (19.1.9)
的类似命题中，必须删去关于各 $N/\mathfrak J_\lambda N$ 的假设，反而假设对
每个 $\lambda$，$M/\mathfrak J_\lambda M$ 是一个投射
$(A/\mathfrak J_\lambda)$-模）。
\end{remark}

\subsection{形式投射模}
\label{subsection:0.19.2-fr}

\begin{definition}[19.2.1]
\label{0.19.2.1-fr}
设 $A$ 是一个拓扑环。若一个拓扑 $A$-模 $M$ 满足下列条件，则称它是
\emph{形式投射的}：对任意开理想 $\mathfrak J$（属于 $A$）、任意一对离散
$(A/\mathfrak J)$-模 $P$、$Q$、任意满 $A$-同态 $u:P\to Q$ 以及任意连续
$A$-同态 $v:M\to Q$，存在连续 $A$-同态 $w:M\to P$，使得
$v=u\circ w$。
\end{definition}

\begin{env}[19.2.2]
\label{0.19.2.2-fr}
为了验证 (19.2.1) 的条件，显然可以把 $Q$ 换成 $v(M)$、把 $P$ 换成
$u^{-1}(v(M))$，从而只考虑 $v$ 本身为满射的情形；此时 $Q$ 同构于
$M/V$，其中 $V$ 是 $M$ 的一个开子模，并且 $\mathfrak JM\subset V$。
于是，(19.2.1) 的条件等价于：对每个离散 $(A/\mathfrak J)$-模 $P$ 以及
每个满 $A$-同态 $u:P\to M/V$，存在一个开子模 $V'\subset V$（属于
$M$）和一个 $A$-同态 $w:M/V'\to P$，使典范同态 $M/V'\to M/V$
可分解为
\[
  M/V'\xrightarrow{w}P\xrightarrow{u}M/V.
\]
应当注意，只须在 $\mathfrak J$ 遍历由理想组成的 $A$ 中 0 的一个基本邻域系时
验证这一条件。
\end{env}

\begin{env}[19.2.3]
\label{0.19.2.3-fr}
假设存在一个开理想基本系 $(\mathfrak J_\lambda)$（属于 $A$），以及一个
开子模基本系 $(V_\lambda)$（属于 $M$）；后者与 $(\mathfrak J_\lambda)$ 具有
\emph{同一个}指标集，并且对每个 $\lambda$，$M/V_\lambda$ 都是一个投射
$(A/\mathfrak J_\lambda)$-模。则 $M$ 是形式投射的：事实上，沿用
(19.2.2) 的记号，只须取 $\lambda$，使 $\mathfrak J_\lambda\subset\mathfrak J$
且 $V_\lambda\subset V$；由于 $P$ 也是一个 $(A/\mathfrak J_\lambda)$-模，
典范同态 $M/V_\lambda\to M/V$ 可分解为
$M/V_\lambda\to P\to M/V$，这是因为它是一个
$(A/\mathfrak J_\lambda)$-同态，并且 $M/V_\lambda$ 按假设是一个投射
$(A/\mathfrak J_\lambda)$-模。

若本条中这个较强的条件成立，则称 $M$ 是\emph{严格形式投射的}。
\end{env}

\begin{proposition}[19.2.4]
\label{0.19.2.4-fr}
设 $A$ 是一个拓扑环，$M$ 是一个拓扑由 $A$ 的拓扑导出的拓扑 $A$-模。
则 $M$ 是形式投射的，当且仅当它是严格形式投射的。
\end{proposition}

\oldpage[0\textsubscript{IV-1}]{79}

\begin{proof}
刚才已经看到，该条件是充分的。反过来，假设 $M$ 是形式投射的，并设
$(\mathfrak J_\lambda)$ 是 $A$ 中的一个开理想基本系。对每个 $\lambda$，
取一个模 $P$，它是自由 $(A/\mathfrak J_\lambda)$-模；再取一个满同态
$u:P\to M/\mathfrak J_\lambda M$，它是
$(A/\mathfrak J_\lambda)$-线性的。于是存在
$\mathfrak J_\mu\subset\mathfrak J_\lambda$，使典范同态
$M/\mathfrak J_\mu M\to M/\mathfrak J_\lambda M$ 可分解为
\[
  M/\mathfrak J_\mu M\xrightarrow{w}P\xrightarrow{u}M/\mathfrak J_\lambda M;
\]
但是，由于
$w(\mathfrak J_\lambda(M/\mathfrak J_\mu M))\subset\mathfrak J_\lambda P=0$，
$w$ 可分解为 $M/\mathfrak J_\lambda M\xrightarrow{v}P$；显然，
$u\circ v$ 是 $M/\mathfrak J_\lambda M$ 上的恒等映射。这就证明了这个
$(A/\mathfrak J_\lambda)$-模是投射的。
\end{proof}

\begin{proposition}[19.2.5]
\label{0.19.2.5-fr}
设 $A$ 是一个拓扑环，$M$ 是一个拓扑 $A$-模。
\begin{enumerate}
  \item[(i)] $M$ 是形式投射的（相应地，严格形式投射的），当且仅当拓扑
  $\hat A$-模 $\hat M$ 具有相同性质。
  \item[(ii)] 设 $A'$ 是一个拓扑 $A$-代数。若 $M$ 是形式投射的（相应地，
  严格形式投射的），则赋予张量积拓扑的 $M\otimes_A A'$ 是一个形式投射的
  （相应地，严格形式投射的）拓扑 $A'$-模。
\end{enumerate}
\end{proposition}

\begin{proof}
\begin{enumerate}
  \item[(i)] 只须注意：当 $\mathfrak J$（相应地，$V$）遍历 $A$ 的所有开理想
  （相应地，$M$ 的所有开子模）时，分离完备化 $\hat{\mathfrak J}$
  （相应地，$\hat V$）遍历 $\hat A$ 的所有开理想（相应地，
  $\hat M$ 的所有开子模），而且在相差一个典范同构的意义下，
  $\hat A/\hat{\mathfrak J}=A/\mathfrak J$（相应地，
  $\hat M/\hat V=M/V$）。由于形式投射模（相应地，严格形式投射模）
  的概念只涉及各 $A/\mathfrak J$ 和各 $M/V$，立即得到 (i)。
  \item[(ii)] 首先假设 $M$ 是形式投射的，并置 $M'=M\otimes_A A'$。设
  $\mathfrak J'$ 是 $A'$ 的一个开理想，$P'$、$Q'$ 是两个离散
  $(A'/\mathfrak J')$-模，$u':P'\to Q'$ 是一个满 $A'$-同态，而
  $v':M'\to Q'$ 是一个连续 $A'$-同态。存在一个开理想 $\mathfrak J$
  （属于 $A$），使 $\mathfrak JA'\subset\mathfrak J'$，因而 $P'$ 和
  $Q'$ 也都是离散 $(A/\mathfrak J)$-模。考察连续的复合 $A$-同态
  $v:M\xrightarrow{j}M'\xrightarrow{v'}Q'$；由假设，存在连续
  $A$-同态 $w:M\to P'$，使 $v=u'\circ w$。但由于 $P'$ 是一个拓扑
  $A'$-模，$w$ 可分解为
  $M\xrightarrow{j}M'\xrightarrow{w'}P'$，其中 $w'$ 是一个连续
  $A'$-同态；又因 $v'\circ j=(u'\circ w')\circ j$，可得
  $v'=u'\circ w'$。

  其次假设 $M$ 是严格形式投射的。设 $(\mathfrak J_\lambda)$（相应地，
  $(V_\lambda)$）是 $A$（相应地，$M$）中的一个开理想基本系（相应地，
  开子模基本系），并且 $(M/V_\lambda)$ 是一个投射
  $(A/\mathfrak J_\lambda)$-模；再设 $(\mathfrak J'_\mu)$ 是 $A'$ 中的一个
  开理想基本系。取下列子模，可得关于 $0$ 的、位于 $M'$ 中的一个基本邻域系：
  \[
    \operatorname{Im}(M\otimes_A\mathfrak J'_\mu)
    +\operatorname{Im}(V_\lambda\otimes_A A')=W_{\lambda\mu},
  \]
  其中 $\lambda$ 和 $\mu$ 满足
  $\mathfrak J_\lambda A'\subset\mathfrak J'_\mu$。由于
  \[
    (M\otimes_A A')/W_{\lambda\mu}
    =(M/V_\lambda)\otimes_{A/\mathfrak J_\lambda}
      (A'/\mathfrak J'_\mu)
  \]
  并且 $(M/V_\lambda)$ 是一个投射 $(A/\mathfrak J_\lambda)$-模，故
  $M'/W_{\lambda\mu}$ 是一个投射 $(A'/\mathfrak J'_\mu)$-模。
\end{enumerate}
\end{proof}

\subsection{形式平滑代数}
\label{subsection:0.19.3-fr}

\begin{definition}[19.3.1]
\label{0.19.3.1-fr}
设 $A$ 是一个拓扑环，$B$ 是一个拓扑 $A$-代数。若 $B$ 作为 $A$-代数满足
下列条件，则称它是\emph{形式平滑的}：对每个离散拓扑 $A$-代数 $C$ 和每个
幂零理想 $\mathfrak J$（属于 $C$），每个连续 $A$-同态
$u:B\to C/\mathfrak J$ 都可分解为
\[
  B\xrightarrow{v}C\xrightarrow{\varphi}C/\mathfrak J,
\]
其中 $v$ 是一个连续同态，而 $\varphi$ 是典范同态。
\end{definition}

\oldpage[0\textsubscript{IV-1}]{80}

定义 (19.3.1) 等价于下列性质成立：
\begin{enumerate}
  \item[(P)] 对每个开理想 $\mathfrak R$（属于 $A$-代数 $B$）以及每个
  $A$-同态 $u':B/\mathfrak R\to C/\mathfrak J$，存在一个开理想
  $\mathfrak R'\subset\mathfrak R$（属于 $B$），使同态
  $B\to B/\mathfrak R\xrightarrow{u'}C/\mathfrak J$ 可分解为
  \[
    B\to B/\mathfrak R'\xrightarrow{v'}C\to C/\mathfrak J,
  \]
  其中 $v'$ 是一个 $A$-同态。
\end{enumerate}
事实上，若 $u:B\to C/\mathfrak J$ 是一个连续 $A$-同态，则它的核
$\mathfrak R$ 是 $B$ 的一个开理想，故 $u$ 可分解为
$B\to B/\mathfrak R\xrightarrow{u'}C/\mathfrak J$；若 (P) 成立，只须把它
应用于 $u'$，并把 $v$ 取为复合 $B\to B/\mathfrak R'\xrightarrow{v'}C$，
即可满足 (19.3.1) 的条件。反过来，假设 $B$ 是一个形式平滑的 $A$-代数；
给定一个开理想 $\mathfrak R$（属于 $B$）以及一个
$A$-同态 $u':B/\mathfrak R\to C/\mathfrak J$，把定义 (19.3.1) 应用于
$u:B\to B/\mathfrak R\xrightarrow{u'}C/\mathfrak J$。若
$v:B\to C$ 是一个连续 $A$-同态，并且 $u$ 可分解为
$B\xrightarrow{v}C\to C/\mathfrak J$，则理想
$\mathfrak R'=\operatorname{Ker}(v)\cap\mathfrak R$（属于 $B$）是开理想且包含于
$\mathfrak R$；于是 $v$ 可分解为
$B\to B/\mathfrak R'\xrightarrow{v'}C$，而 $v'$ 确实满足条件 (P)。

\begin{proposition}[19.3.2]
\label{0.19.3.2-fr}
设 $A$ 是一个离散环，$V$ 是一个投射 $A$-模；赋予离散拓扑的对称代数
$B=\mathbf S_A^\bullet(V)$ 是一个形式平滑的 $A$-代数。
\end{proposition}

\begin{proof}
事实上，记号 $C$ 和 $\mathfrak J$ 的含义与 (19.3.1) 相同，并设
$u:B\to C/\mathfrak J$ 是一个 $A$-代数同态；把它限制到
$V=\mathbf S_A^1(V)$ 上，就得到一个 $A$-模同态
$u_1:V\to C/\mathfrak J$。由于 $V$ 是投射的，$u_1$ 可分解为
$V\xrightarrow{v_1}C\xrightarrow{\varphi}C/\mathfrak J$；又可把 $v_1$
延拓为 $A$-代数同态 $v:\mathbf S_A^\bullet(V)\to C$，使复合
$\varphi\circ v$ 与 $u$ 在 $V$ 上重合，因而 $\varphi\circ v=u$；命题得证。
\end{proof}

\begin{corollary}[19.3.3]
\label{0.19.3.3-fr}
若 $A$ 是一个离散环，则赋予离散拓扑的任意多项式代数
$B=A[T_\alpha]_{\alpha\in I}$ 都是一个形式平滑的 $A$-代数。
\end{corollary}

\begin{proposition}[19.3.4]
\label{0.19.3.4-fr}
设 $A$ 是一个拓扑环，并设
$B=A[[T_\alpha]]_{\alpha\in I}$ 是形式幂级数
$\sum_{\lambda\in\mathbb{N}^{(I)}}c_\lambda T^\lambda$ 所成的环（$A$ 上的
幺半群 $\mathbb{N}^{(I)}$ 的大代数；把它作为 $A$-模，与乘积
$A^{\mathbb{N}^{(I)}}$ 等同）。若赋予 $B$ 乘积拓扑，则 $B$ 是一个形式平滑的
$A$-代数。
\end{proposition}

设 $(\mathfrak L_\mu)_{\mu\in M}$ 是一个开理想基本系（属于 $A$）。对每个
有限子集 $H$（属于 $I$）、每个 $\mu\in M$ 以及每个整数 $n$，以
$\mathfrak K_{H,\mu,n}$ 表示 $0$ 在 $B$ 中的如下邻域：它由满足下列条件的
$(c_\lambda)$ 所成。对每个 $\lambda=(\lambda_\alpha)_{\alpha\in I}$，若
$\lambda_\alpha=0$ 对 $\alpha\notin H$ 成立，并且
$\sum_{\alpha\in H}\lambda_\alpha\leq n$，则
$c_\lambda\in\mathfrak L_\mu$。立即验证，各 $\mathfrak K_{H,\mu,n}$ 组成
$0$ 在 $B$ 中的一个基本邻域系，而且都是 $B$ 的理想；因此乘积拓扑与
$A$-代数 $B$ 的结构相容。

仍沿用这些记号，先指出下列引理。
\begin{lemma}[19.3.4.1]
\label{0.19.3.4.1-fr}
设 $E$ 是一个离散 $A$-代数。
\begin{enumerate}
  \item[(i)] 若 $f:B\to E$ 是一个连续 $A$-同态，则存在一个有限子集 $H$
  （属于 $I$），使得 $f(T_\alpha)=0$ 对 $\alpha\notin H$ 成立，并且
  $f(T_\alpha)$ 是 $E$ 中的幂零元，对每个 $\alpha\in H$ 都成立。
  \item[(ii)] 反过来，设 $H$ 是 $I$ 的一个有限子集，
  $(z_\alpha)_{\alpha\in H}$ 是 $E$ 中的一族幂零元。则存在唯一的连续
  $A$-同态 $g:B\to E$，使得 $g(T_\alpha)=z_\alpha$ 对 $\alpha\in H$
  成立，而 $g(T_\alpha)=0$ 对 $\alpha\notin H$ 成立。
\end{enumerate}
\end{lemma}

\oldpage[0\textsubscript{IV-1}]{81}

\begin{proof}
(i) 立即由下述事实得出：$f^{-1}(0)$ 是 $0$ 在乘积 $A$-模
$A^{\mathbb{N}^{(I)}}$ 中的一个邻域；因而 $f(T^\lambda)=0$，只有有限多个
$\lambda\in\mathbb{N}^{(I)}$ 例外。为证明
(ii)，只须注意多项式环 $B'=A[T_\alpha]_{\alpha\in I}$ 在 $B$ 中稠密；
限制 $g|B'$ 的存在性和唯一性都是显然的，而其连续性来自下述假设：各
$z_\alpha$ 对 $\alpha\in H$ 都是幂零的。事实上，若
$z_\alpha^n=0$ 对每个 $\alpha\in H$ 成立，则对每个
$g(T^\lambda)=0$；这里 $\lambda=(\lambda_\alpha)_{\alpha\in I}$ 具有有限支撑，
但以下情形除外：$\lambda_\alpha=0$ 对 $\alpha\notin H$ 成立，并且
$\lambda_\alpha\leq n$ 对 $\alpha\in H$ 成立；也就是说，这种例外的
$\lambda$ 只有有限多个取值。
\SourceCorrectionNote{法文印本和外交式转录在本证明的上述两个位置误用 $f(T_\alpha)$；该符号属于 (i)，在 (ii) 中并未绑定。这里应使用已给定的幂零元 $z_\alpha$（等价地，使用 $g(T_\alpha)$）。}
\end{proof}

\begin{proof}[命题~{\rm(19.3.4)}的证明]
该引理既已证明，并且记号 $C$ 和 $\mathfrak J$ 的含义与 (19.3.1) 相同，则
$u(T_\alpha)=0$ 对 $\alpha\in H$ 以外的指标成立，其中 $H$ 是 $I$ 的一个
有限子集；而 $z_\alpha=u(T_\alpha)$ 对 $\alpha\in H$ 而言是
$C/\mathfrak J$ 中的幂零元。由于 $\mathfrak J$ 是幂零的，存在一族幂零元
$(x_\alpha)_{\alpha\in H}$（属于 $C$），它们在 $C/\mathfrak J$ 中的典范像
就是 $z_\alpha$。若 $v$ 是一个连续 $A$-同态，它从 $B$ 映入 $C$，并满足
$v(T_\alpha)=0$ 对 $\alpha\notin H$ 成立，而
$v(T_\alpha)=x_\alpha$ 对 $\alpha\in H$ 成立，则显然 $u$ 可分解为
$B\xrightarrow{v}C\to C/\mathfrak J$。
\end{proof}

\begin{proposition}[19.3.5]
\label{0.19.3.5-fr}
设 $A$ 是一个拓扑环。
\begin{enumerate}
  \item[(i)] $A$ 是一个形式平滑的 $A$-代数。
  \item[(ii)] 若 $B$ 是一个形式平滑的 $A$-代数，而 $C$ 是一个形式平滑的
  $B$-代数，则 $C$ 是一个形式平滑的 $A$-代数。
  \item[(iii)] 设 $B$ 是一个形式平滑的 $A$-代数，而 $A'$ 是一个拓扑
  $A$-代数；则拓扑 $A'$-代数 $B\otimes_A A'$（$\mathbf 0_{\rm I}$，
  7.7.5 和 7.7.6）是形式平滑的。
  \item[(iv)] 设 $B$ 是一个拓扑 $A$-代数，而 $S$（相应地，$T$）是
  $A$（相应地，$B$）的一个乘法子集，并且 $S$ 在 $B$ 中的典范像包含于
  $T$。若 $B$ 是一个形式平滑的 $A$-代数，则 $T^{-1}B$ 是一个形式平滑的
  $S^{-1}A$-代数。
  \item[(v)] 设 $B_i$（$1\leq i\leq n$）是一些拓扑 $A$-代数。乘积
  $\prod_{i=1}^n B_i$ 是一个形式平滑的 $A$-代数，当且仅当每个 $B_i$
  都是如此。
\end{enumerate}
\end{proposition}

\oldpage[0\textsubscript{IV-1}]{82}

\begin{proof}
\begin{enumerate}
  \item[(i)] 若 $C$ 是一个离散 $A$-代数，而
  $\varphi:C\to C/\mathfrak J$ 是从 $C$ 到任意一个商 $A$-代数（属于 $C$）的典范同态，
  则\emph{唯一}的 $A$-同态，从 $A$ 映入 $C/\mathfrak J$ 的那个，是复合
  $A\xrightarrow{\psi}C\xrightarrow{\varphi}C/\mathfrak J$，其中 $\psi$
  是定义 $A$-代数 $C$ 结构的同态；由于 $\psi$ 连续，(19.3.1) 的条件显然
  成立。
  \item[(ii)] 设 $\alpha:A\to B$、$\beta:B\to C$ 是分别定义 $A$-代数
  $B$ 的结构和 $B$-代数 $C$ 的结构的连续同态，因而
  $\beta\circ\alpha$ 定义 $A$-代数 $C$ 的结构。设 $E$ 是一个离散
  $A$-代数，$\mathfrak L$ 是 $E$ 的一个幂零理想，而
  $u:C\to E/\mathfrak L$ 是一个连续 $A$-同态，因而
  $u\circ\beta\circ\alpha$ 是定义 $A$-代数 $E/\mathfrak L$ 的结构的同态。
  由于 $B$ 是一个形式平滑的 $A$-代数，连续 $A$-同态
  $u\circ\beta:B\to E/\mathfrak L$ 可分解为
  $B\xrightarrow{v}E\to E/\mathfrak L$，其中 $v$ 是一个连续 $A$-同态；
  此时，$v$ 和 $u\circ\beta$ 分别在 $E$ 和 $E/\mathfrak L$ 上定义拓扑
  $B$-代数结构，在这些结构之下，$E/\mathfrak L$ 仍是 $B$-代数 $E$ 的
  （离散）商代数。此外，$u$ 是一个连续 $B$-同态，故可分解为
  $C\xrightarrow{w}E\to E/\mathfrak L$，其中 $w$ 是一个连续 $B$-同态；
  由于 $w\circ\beta\circ\alpha$ 是一个 $A$-同态，它定义 $A$-代数 $E$ 的结构，$w$
  确实是一个连续 $A$-同态，断言得证。
  \SourceCorrectionNote{法文印本和外交式转录在上句写作不可复合的 $w\circ\alpha$。应在其中插入 $\circ\beta$，即结构同态为 $w\circ\beta\circ\alpha$。}
  \item[(iii)] 设 $C$ 是一个离散拓扑 $A'$-代数，$\mathfrak J$ 是 $C$ 的一个
  幂零理想，而 $u:B\otimes_A A'\to C/\mathfrak J$ 是一个连续 $A'$-同态。
  把 $u$ 与典范同态 $\rho:B\to B\otimes_A A'$ 复合；由于复合
  $A\to B\xrightarrow{\rho}B\otimes_A A'$ 等于复合
  $A\to A'\xrightarrow{\sigma}B\otimes_A A'$，所得连续同态是 $A$-线性的，
  故由假设可分解为 $B\xrightarrow{v}C\to C/\mathfrak J$，其中 $v$ 是一个
  连续 $A$-同态。复合 $A\to B\xrightarrow{v}C$ 与复合
  $A\to A'\xrightarrow{w}C$ 相等，因而存在一个连续环同态
  $f:B\otimes_A A'\to C$，使得 $v=f\circ\rho$ 且 $w=f\circ\sigma$
  （$\mathbf 0_{\rm I}$，7.7.6）。由于两个复合
  $B\xrightarrow{\rho}B\otimes_A A'\xrightarrow{f}C\to C/\mathfrak J$
  和 $B\xrightarrow{\rho}B\otimes_A A'\xrightarrow{u}C/\mathfrak J$（相应地，
  两个复合 $A'\xrightarrow{\sigma}B\otimes_A A'\xrightarrow{f}C\to C/\mathfrak J$
  和 $A'\xrightarrow{\sigma}B\otimes_A A'\xrightarrow{u}C/\mathfrak J$）相等，
  确有分解 $u:B\otimes_A A'\xrightarrow{f}C\to C/\mathfrak J$，这就证明了
  (iii)。
  \item[(iv)] 所考虑的 $S^{-1}A$（相应地，$T^{-1}B$）上的拓扑自然是如下
  拓扑：$0$ 的一个基本邻域系由 $S^{-1}\mathfrak J$（相应地，
  $T^{-1}\mathfrak R$）组成，其中 $\mathfrak J$（相应地，$\mathfrak R$）遍历
  $A$（相应地，$B$）中的一个开理想基本系（$\mathbf 0_{\rm I}$，7.6.1）。若
  $\alpha:A\to B$ 是典范同态，则显然，典范同态
  $\alpha':S^{-1}A\to T^{-1}B$（由 $\alpha$ 导出）连续。现在设 $C$ 是一个离散拓扑
  $S^{-1}A$-代数，$\mathfrak J$ 是这个代数的一个幂零理想，而
  $u:T^{-1}B\to C/\mathfrak J$ 是一个连续 $S^{-1}A$-同态；于是复合
  $B\to T^{-1}B\xrightarrow{u}C/\mathfrak J$ 是一个连续 $A$-同态，故由
  假设可分解为 $B\xrightarrow{v}C\to C/\mathfrak J$，其中 $v$ 是一个连续
  $A$-同态。对每个 $t\in T$，$t/1$ 在 $T^{-1}B$ 中可逆，所以 $u(t/1)$
  在 $C/\mathfrak J$ 中可逆。由于 $\mathfrak J$ 幂零，$u(t/1)$ 在 $C$ 中的
  类中的每个元素，特别是 $v(t)$，都在 $C$ 中可逆（$\mathbf 0_{\rm I}$，
  7.1.12）；因此 $v$ 可分解为 $B\to T^{-1}B\xrightarrow{w}C$。由于 $v$
  连续，$w$ 也连续（$\mathbf 0_{\rm I}$，7.6.6）；它还是一个
  $S^{-1}A$-同态，因为复合 $A\to B\to T^{-1}B\xrightarrow{w}C$ 等于
  \[
    A\to S^{-1}A\xrightarrow{\alpha'}T^{-1}B\xrightarrow{w}C,
  \]
  故 $S^{-1}A\xrightarrow{\alpha'}T^{-1}B\xrightarrow{w}C$ 正是定义 $C$ 上
  $S^{-1}A$-代数结构的典范同态。最后，由于两个复合
  $B\to T^{-1}B\xrightarrow{w}C\to C/\mathfrak J$ 和
  $B\to T^{-1}B\xrightarrow{u}C/\mathfrak J$ 相等，同一论证表明 $u$ 确实等于
  复合 $T^{-1}B\xrightarrow{w}C\to C/\mathfrak J$；由此得到 (iv)。
  \item[(v)] 必要性由 (iv) 得出，因为每个 $B_i$ 都是
  $B=\prod_{i=1}^n B_i$ 关于其典范幂等元的局部化。反过来，设 $C$ 是一个
  离散 $A$-代数，$\mathfrak J$ 是 $C$ 的一个幂零理想，而
  $u:B\to C/\mathfrak J$ 是一个连续 $A$-同态。在 $u$ 下，$B$ 的各典范幂等元之像
  在余域中组成一族有限的两两正交幂等元，其和为 $1$；把这一族提升为一族
  幂等元 $e_i$（属于 $C$），它们两两正交且和为 $1$。置 $C_i=e_iC$ 以及
  $\mathfrak J_i=e_i\mathfrak J$，则 $u$ 在 $B_i$ 上的限制是一个连续
  $A$-同态 $B_i\to C_i/\mathfrak J_i$。由 $B_i$ 的形式平滑性，它可提升为
  一个连续 $A$-同态 $B_i\to C_i$。这些提升的乘积就是一个连续
  $A$-同态 $B\to C$，并且它提升 $u$。
  \SourceCorrectionNote{法文印本和外交式转录错误地把从有限乘积出发的同态 当成任意一族分量同态，并给出了无效证明。以幂等元提升和逐因子形式平滑性替换该证明。}
\end{enumerate}
\end{proof}

\oldpage[0\textsubscript{IV-1}]{83}

\begin{proposition}[19.3.6]
\label{0.19.3.6-fr}
设 $A$ 是一个拓扑环，$B$ 是一个拓扑 $A$-代数，$\hat A$ 和 $\hat B$
分别是 $A$ 和 $B$ 的分离完备化，因而 $\hat B$ 是一个拓扑 $\hat A$-代数。
下列条件等价：
\begin{enumerate}
  \item[a)] $B$ 是一个形式平滑的 $A$-代数。
  \item[b)] $\hat B$ 是一个形式平滑的 $A$-代数。
  \item[c)] $\hat B$ 是一个形式平滑的 $\hat A$-代数。
\end{enumerate}
\end{proposition}

\begin{proof}
当然，作为 $\hat A$-代数的 $\hat B$，其结构由同态 $\hat\varphi$ 定义；这里
$\varphi:A\to B$ 是给出 $A$-代数 $B$ 结构的同态。由于每个离散
$A$-代数 $C$ 都分离且完备，它也成为一个 $\hat A$-代数（这是把从 $A$ 到
$C$ 的同态典范延拓所得）；而一个连续 $A$-同态，即从 $B$ 到 $C$ 的同态，
经过典范延拓，给出一个 $\hat A$-同态（从而更是一个 $A$-同态），它从
$\hat B$ 映入 $C$（换言之，每个连续 $A$-同态 $B\to C$ 都唯一地分解为
$B\to\hat B\to C$）。这些观察和定义 (19.3.1) 立即推出本命题。
\SourceCorrectionNote{法文印本和外交式转录在上述括注中把“连续 $A$-同态”误写成了任意的 $A$-同态。完备化的泛性质只延拓连续同态；补入连续性条件。}
\end{proof}

\begin{corollary}[19.3.7]
\label{0.19.3.7-fr}
在 (19.3.5), (iv) 的假设下，拓扑 $A\{S^{-1}\}$-代数
$B\{T^{-1}\}$ 是形式平滑的。
\end{corollary}

这由定义（$\mathbf 0_{\rm I}$，7.6.1 和 7.6.7）以及 (19.3.5), (iv) 和
(19.3.6) 得出。

\begin{proposition}[19.3.8]
\label{0.19.3.8-fr}
设 $A$ 是一个拓扑环，并设 $B$ 是一个拓扑 $A$-代数；假设对每个开理想
$\mathfrak R$（属于 $B$），$\mathfrak R^2$ 也开。设 $A'$、$B'$ 是分别在 $A$ 和
$B$ 上赋予比原拓扑更细的拓扑所得的环，并假设典范同态
$\varphi:A\to B$ 作为从 $A'$ 到 $B'$ 的同态仍连续。于是，若 $B'$ 是一个
形式平滑的 $A'$-代数，则 $B$ 是一个形式平滑的 $A$-代数。
\end{proposition}

只须应用下列引理：

\begin{lemma}[19.3.8.1]
\label{0.19.3.8.1-fr}
设 $C$ 是一个离散 $A$-代数，$\mathfrak J$ 是 $C$ 的一个幂零理想。假设
$B$ 的每个开理想的平方都开。于是，若复合同态
$v:B\xrightarrow{u}C\to C/\mathfrak J$ 连续，则同态 $u$ 连续。
\end{lemma}

\begin{proof}
事实上，$\mathfrak R=u^{-1}(\mathfrak J)$ 是 $v$ 的核，并且依假设它开；
由于存在一个整数 $n$ 使 $\mathfrak J^n=0$，$\mathfrak R^n$ 包含于
$\operatorname{Ker}(u)$。但由假设对 $h$ 作归纳可知，$\mathfrak R^{2^h}$
是开理想，因而 $\mathfrak R^n$ 对每个 $n$ 也开；于是
$\operatorname{Ker}(u)$ 也开，这就证明了断言。
\end{proof}

应注意，关于 $B$ 的这个假设意味着 $B$ 的拓扑是各 $\mathfrak R$-预进拓扑的
上确界，其中 $\mathfrak R$ 遍历 $B$ 的所有开理想。若 $B$ 是一个预进环
（$\mathbf 0_{\rm I}$，7.1.9），则这个条件总是成立。

\begin{env}[19.3.9]
\label{0.19.3.9-fr}
现在我们将看到，形式平滑性蕴涵同态的若干“提升”性质，其中的条件比定义
(19.3.1) 中的条件更为一般。
\end{env}

\oldpage[0\textsubscript{IV-1}]{84}

\begin{proposition}[19.3.10]
\label{0.19.3.10-fr}
设 $A$ 是一个拓扑环，$B$ 是一个形式平滑的 $A$-代数。设 $C$ 是一个拓扑
$A$-代数，$\mathfrak J$ 是 $C$ 的一个理想，并且满足下列条件：
\begin{enumerate}
  \item[$1^\circ$] $C$ 可度量且完备。
  \item[$2^\circ$] $\mathfrak J$ 闭，并且序列 $(\mathfrak J^n)$ 趋于 $0$。
\end{enumerate}
于是，每个连续 $A$-同态 $u:B\to C/\mathfrak J$ 都可分解为
$B\xrightarrow{v}C\to C/\mathfrak J$，其中 $v$ 是一个连续 $A$-同态。
\end{proposition}

\begin{proof}
设 $(\mathfrak Q_n)$ 是 $C$ 的一个递减理想序列，并组成 $0$ 的一个基本邻域系。
对每个 $n$，考虑离散 $A$-代数 $C/\mathfrak Q_n$ 及其中的理想
$(\mathfrak J+\mathfrak Q_n)/\mathfrak Q_n$；由于存在 $k$ 使得
$\mathfrak J^k\subset\mathfrak Q_n$，理想
$(\mathfrak J+\mathfrak Q_n)/\mathfrak Q_n$ 是幂零的。对每个 $n$，令 $u_n$
为复合同态
\[
  B\xrightarrow{u}C/\mathfrak J\to C/(\mathfrak J+\mathfrak Q_n)
  =(C/\mathfrak Q_n)/((\mathfrak J+\mathfrak Q_n)/\mathfrak Q_n)~;
\]
依假设，$u_n$ 可分解为
\[
  B\xrightarrow{v_n}C/\mathfrak Q_n
  \xrightarrow{\varphi_n}C/(\mathfrak J+\mathfrak Q_n),
\]
其中 $v_n$ 是一个连续 $A$-同态。我们来证明，可以逐个选择 $v_n$，使得
$v_n$ 分解为
\[
  B\xrightarrow{v_{n+1}}C/\mathfrak Q_{n+1}\to C/\mathfrak Q_n
\]
对每个 $n$ 都成立（换言之，各 $v_n$ 组成一个同态的\emph{逆系统}）。这将由
下列引理得出。
\end{proof}

\begin{lemma}[19.3.10.1]
\label{0.19.3.10.1-fr}
设 $B$ 是一个形式平滑的 $A$-代数，$E$、$E'$ 是两个离散拓扑
$A$-代数，$\mathfrak R$（相应地，$\mathfrak R'$）是 $E$（相应地，$E'$）的
一个幂零理想；设 $f:E\to E'$ 是一个满 $A$-同态，满足
$f(\mathfrak R)\subset\mathfrak R'$，并设
$f':E/\mathfrak R\to E'/\mathfrak R'$ 是由 $f$ 取商所得的同态。设
$u:B\to E/\mathfrak R$ 是一个连续 $A$-同态，$u'$ 是复合同态
$B\xrightarrow{u}E/\mathfrak R\xrightarrow{f'}E'/\mathfrak R'$；再设
$v':B\to E'$ 是一个连续 $A$-同态，使得 $u'$ 可分解为
$B\xrightarrow{v'}E'\to E'/\mathfrak R'$。则存在一个连续 $A$-同态
$v:B\to E$，使得 $v'$ 可分解为
$B\xrightarrow{v}E\xrightarrow{f}E'$。
\end{lemma}

在交换图
\[
\xymatrix{
  E \ar[r]^f \ar[d] & E' \ar[d] \\
  E/\mathfrak R \ar[r]_{f'} & E'/\mathfrak R'
}
\]
中，所有同态都是满的；若 $\mathfrak L$ 是 $f$ 的核，则 $E'$ 等同于
$E/\mathfrak L$，而 $\mathfrak R'$ 等同于
$(\mathfrak R+\mathfrak L)/\mathfrak L$。现在使用下列初等引理。

\begin{lemma}[19.3.10.2]
\label{0.19.3.10.2-fr}
设 $F$ 是一个环（不必交换），$\mathfrak a$、$\mathfrak b$ 是 $F$ 的两个
双边理想；则纤维积
\[
  (F/\mathfrak a)\times_{F/(\mathfrak a+\mathfrak b)}(F/\mathfrak b)
\]
典范地等同于 $F/(\mathfrak a\cap\mathfrak b)$。
\end{lemma}

\begin{proof}
这正是 (18.1.7) 的一个特殊应用：在图 (18.1.7.1) 中，把 $G$ 替换为
$F/(\mathfrak a\cap\mathfrak b)$，把 $\mathfrak J'$ 替换为
$\mathfrak b/(\mathfrak a\cap\mathfrak b)$，把 $F$ 替换为
$F/\mathfrak b$，把 $\mathfrak R$ 替换为
$(\mathfrak a+\mathfrak b)/\mathfrak b$，把 $B$ 替换为
$F/(\mathfrak a+\mathfrak b)$，并把 $\theta$ 替换为一个 $F$-同态，即典范双射
$(\mathfrak a+\mathfrak b)/\mathfrak b\to
\mathfrak a/(\mathfrak a\cap\mathfrak b)$。
\end{proof}

\oldpage[0\textsubscript{IV-1}]{85}

\begin{proof}[引理~{\rm(19.3.10.1)}的证明（续）]
把这个引理应用于 (19.3.10.1) 的情形，交换图
\[
\xymatrix{
  B \ar[r]^{v'} \ar[d]_u & E'=E/\mathfrak L \ar[d] \\
  E/\mathfrak R \ar[r]_{f'} & E'/\mathfrak R'
}
\]
的交换性表明，存在同态
$w:B\to E/(\mathfrak R\cap\mathfrak L)$，使得 $v'$ 和 $u$ 分别分解为
$B\xrightarrow{w}E/(\mathfrak R\cap\mathfrak L)\to E/\mathfrak L$
和
$B\xrightarrow{w}E/(\mathfrak R\cap\mathfrak L)\to E/\mathfrak R$；由于
$w$ 的核包含 $v'$ 和 $u$ 的核的交，它在 $B$ 中开，故 $w$ 连续。最后，
$\mathfrak R\cap\mathfrak L$ 显然是幂零的，因而 $w$ 可分解为
$B\xrightarrow{v}E\to E/(\mathfrak R\cap\mathfrak L)$，其中 $v$ 是所求的
连续 $A$-同态。
\end{proof}

\begin{proof}[命题~{\rm(19.3.10)}的证明（续）]
$v_{n+1}$ 的存在性由引理 (19.3.10.1) 得出：在该引理中，分别把
$E$、$E'$、$\mathfrak R$、$\mathfrak R'$ 替换为
$C/\mathfrak Q_{n+1}$、$C/\mathfrak Q_n$、
$(\mathfrak J+\mathfrak Q_{n+1})/\mathfrak Q_{n+1}$ 和
$(\mathfrak J+\mathfrak Q_n)/\mathfrak Q_n$，并分别把 $u$、$u'$ 和 $v'$
替换为 $u_{n+1}$、$u_n$ 和 $v_n$。应注意，由于 $C$ 分离且完备，
它等于 $\varprojlim_n(C/\mathfrak Q_n)$；并且
$v=\varprojlim_n v_n$ 是一个连续 $A$-同态，它从 $B$ 映入 $C$。由于
$\mathfrak J$ 在 $C$ 中闭，有
$\mathfrak J=\bigcap_n(\mathfrak J+\mathfrak Q_n)$；又由于
$C/\mathfrak J$ 可度量且完备，并且
$(\mathfrak J+\mathfrak Q_n)/\mathfrak J$ 组成 $0$ 的一个基本邻域系（属于
$C/\mathfrak J$），有
\[
  C/\mathfrak J=\varprojlim_n C/(\mathfrak J+\mathfrak Q_n)
\]
而 $\varprojlim_n\varphi_n$ 是典范映射 $\varphi:C\to C/\mathfrak J$。由于
$\varprojlim_n u_n=u$，确有 $u=\varphi\circ v$。证毕。
\SourceCorrectionNote{法文印本和外交式转录在本证明末句把复合次序误写为 不可定型的相反次序。应先取提升再取商，即典范商映射与提升的复合。}
\end{proof}

\begin{corollary}[19.3.11]
\label{0.19.3.11-fr}
设 $A$ 是一个拓扑环，$B$ 是一个形式平滑的 $A$-代数，$C$ 是一个完备
Noether 局部环，$\mathfrak J$ 是 $C$ 的一个真理想，而
$\varphi:A\to C$ 是一个连续同态，使 $C$ 成为一个拓扑 $A$-代数。于是，
每个连续 $A$-同态 $u_0:B\to C_0=C/\mathfrak J$ 都可分解为
$B\xrightarrow{u}C\to C/\mathfrak J$，其中 $u$ 是一个连续 $A$-同态。
\end{corollary}

事实上，(19.3.10) 的全部条件均满足（$\mathbf 0_{\rm I}$，7.3.5）。

\begin{proposition}[19.3.12]
\label{0.19.3.12-fr}
设 $A$ 是一个拓扑环，$B$ 是一个形式平滑的 $A$-代数，$C$ 是一个拓扑
$A$-代数，而 $\mathfrak J$ 是 $C$ 的一个理想，并且满足下列条件：
\begin{enumerate}
  \item[$1^\circ$] 存在一个开理想基本系 $\mathfrak Q_\lambda$（属于 $C$），使得
  $C_\lambda=C/\mathfrak Q_\lambda$ 都是 Artin 环，并且典范同态
  $C\to\varprojlim_\lambda C_\lambda$ 是拓扑环同构。
  \item[$2^\circ$] 理想 $\mathfrak J$ 在 $C$ 中闭且拓扑幂零。
  \item[$3^\circ$] $B$ 的每个开理想的平方都是开理想。
\end{enumerate}
在这些条件下，每个连续 $A$-同态 $u:B\to C/\mathfrak J$ 都可分解为
$B\xrightarrow{v}C\to C/\mathfrak J$，其中 $v$ 是一个连续 $A$-同态。
\end{proposition}

\begin{proof}
令 $\mathfrak J_\lambda=(\mathfrak J+\mathfrak Q_\lambda)/\mathfrak Q_\lambda$
为 $\mathfrak J$ 在 $C_\lambda$ 中的典范像；它是一个幂零理想。令
$u_\lambda$ 为复合同态
$B\xrightarrow{u}C/\mathfrak J\to C_\lambda/\mathfrak J_\lambda
=C/(\mathfrak J+\mathfrak Q_\lambda)$。我们先证明，每个 $A$-同态
$w:B\to C$，只要 $u$ 可分解为
$B\xrightarrow{w}C\to C/\mathfrak J$，就必定连续。事实上，对每个
$\lambda$，存在一个开理想 $\mathfrak R$（属于 $B$），使得
$u(\mathfrak R)\subset (\mathfrak J+\mathfrak Q_\lambda)/\mathfrak J$，从而
$w(\mathfrak R)\subset\mathfrak J+\mathfrak Q_\lambda$；于是存在
$\mathfrak R^{2^m}$，它是 $\mathfrak R$ 的一个幂，并且满足
$w(\mathfrak R^{2^m})\subset\mathfrak Q_\lambda$。由于
$\mathfrak R^{2^m}$ 在 $B$ 中开，这就证明了断言。因此，只须找到具有上述
性质的一个 $A$-同态 $w$，而无须顾虑其连续性。现在，集合
$\operatorname{Hom}(B,C)$，即全体 $A$-代数同态（从 $B$ 映入 $C$）所成的集合，等于
$\varprojlim_\lambda\operatorname{Hom}(B,C_\lambda)$；后者在
$C_\lambda$-模 $C_\lambda^B$ 中是\emph{闭的}，该模赋予乘积拓扑，并且由于
$C_\lambda$ 是 Artin 环，它是\emph{线性紧的}。对每个 $\lambda$，令
$W_\lambda$

\oldpage[0\textsubscript{IV-1}]{86}

为所有满足下列条件的
$w_\lambda\in\operatorname{Hom}(B,C_\lambda)$ 所成的集合：$u_\lambda$ 可分解为
$B\xrightarrow{w_\lambda}C_\lambda\to
C_\lambda/\mathfrak J_\lambda$。$W_\lambda$ 是 $C_\lambda$-模
$\operatorname{Hom}(B,C_\lambda)$ 中的一个闭线性簇，并且非空，因为 $B$
形式平滑。在乘积
$\prod_{\mu\leq\lambda}\operatorname{Hom}(B,C_\mu)=E_\lambda$ 中，考虑线性
子簇 $U_\lambda$，它由各 $(w_\mu)_{\mu\leq\lambda}$ 所成，其中
$w_\mu=\varphi_{\mu\lambda}\circ w_\lambda$ 对 $\mu\leq\lambda$ 成立，并且
$w_\lambda\in W_\lambda$；这里
$\varphi_{\mu\lambda}:C_\lambda\to C_\mu$ 是典范同态。这个子簇也是闭的。
最后，令 $V_\lambda$ 为乘积
$\prod_\mu\operatorname{Hom}(B,C_\mu)$ 中的线性簇，它是 $U_\lambda$ 与各
$\operatorname{Hom}(B,C_\mu)$（$\mu\not\leq\lambda$）的乘积；它仍然闭。
问题归结为证明各 $V_\lambda$ 的交\emph{非空}，因为这个交的任一元素 $w$
依定义都属于 $\varprojlim_\lambda\operatorname{Hom}(B,C_\lambda)$。此外，
由于 $\mathfrak J$ 闭，并且 $C=\varprojlim_\lambda C_\lambda$ 线性紧，
$C/\mathfrak J$ 等同于
$\varprojlim_\lambda C_\lambda/\mathfrak J_\lambda$，而典范映射
$\psi:C\to C/\mathfrak J$ 等同于 $\varprojlim_\lambda\psi_\lambda$；这里
$\psi_\lambda$ 是典范映射
$C_\lambda\to C_\lambda/\mathfrak J_\lambda$。因此如同在 (19.3.10) 中那样，
得出 $\psi\circ w=u$。
\SourceCorrectionNote{本段原证明存在 实质缺口：代数同态集合在逐点运算下并非一般的模，而沿幂零理想的提升纤维也 未必是仿射线性子空间，所述线性紧性论证因而没有证明所需的非空交。此处保留原证明，并指出这一未解决的缺口。}
\end{proof}

\subsection{形式平滑性的一些基础性的判别法}
\label{subsection:0.19.4-fr}

\begin{proposition}[19.4.1]
\label{0.19.4.1-fr}
设 $A$ 是一个拓扑环，$B$ 是一个拓扑 $A$-代数；假设存在两个递减有向族
$(\mathfrak J_\alpha)_{\alpha\in I}$、
$(\mathfrak R_\alpha)_{\alpha\in I}$，分别由 $A$ 和 $B$ 的理想组成，并且满足：
\begin{enumerate}
  \item[$1^\circ$] $(\mathfrak J_\alpha)$ 趋于 $0$（在 $A$ 中），而
  $(\mathfrak R_\alpha)$ 趋于 $0$（在 $B$ 中）；
  \item[$2^\circ$] 对每个 $\alpha\in I$，有
  $\mathfrak J_\alpha B\subset\mathfrak R_\alpha$（因而
  $B/\mathfrak R_\alpha$ 是一个拓扑
  $(A/\mathfrak J_\alpha)$-代数）；
  \item[$3^\circ$] 对每个 $\alpha\in I$，
  $B/\mathfrak R_\alpha$ 是一个形式平滑的
  $(A/\mathfrak J_\alpha)$-代数。
\end{enumerate}
则 $B$ 是一个形式平滑的 $A$-代数。
\end{proposition}

\begin{proof}
事实上，设 $C$ 是一个离散 $A$-代数，$\mathfrak L$ 是 $C$ 的一个幂零理想，
$\mathfrak R$ 是 $B$ 的一个开理想；由假设，存在 $\alpha\in I$，使得
$\mathfrak R_\alpha\subset\mathfrak R$，故 $B/\mathfrak R$ 是
$B/\mathfrak R_\alpha$ 关于一个开理想的商。每个 $A$-同态
$u':B/\mathfrak R\to C/\mathfrak L$ 也是一个
$(A/\mathfrak J_\alpha)$-同态；因此存在一个开理想
$\mathfrak R'$（属于 $B$），满足
$\mathfrak R_\alpha\subset\mathfrak R'\subset\mathfrak R$，并且复合
$B/\mathfrak R_\alpha\to B/\mathfrak R\xrightarrow{u'}C/\mathfrak L$
可分解为
$B/\mathfrak R_\alpha\to B/\mathfrak R'\to C\to C/\mathfrak L$，
其中第二个同态是一个 $(A/\mathfrak J_\alpha)$-同态。由 (19.3.1) 后的
说明即得结论。
\end{proof}

\begin{corollary}[19.4.2]
\label{0.19.4.2-fr}
设 $A$ 是一个拓扑环，$B$ 是一个拓扑 $A$-代数，而
$(\mathfrak J_\alpha)_{\alpha\in I}$ 是由 $A$ 的理想组成、趋于 $0$ 的一个
递减有向族。$B$ 是一个形式平滑的 $A$-代数，当且仅当对每个
$\alpha\in I$，置 $A_\alpha=A/\mathfrak J_\alpha$、
$B_\alpha=B/\mathfrak J_\alpha B$ 时，$B_\alpha$ 都是一个形式平滑的
$A_\alpha$-代数。
\end{corollary}

充分性由 (19.4.1) 得出，必要性则由 (19.3.5), (iii) 得出。

\begin{proposition}[19.4.3]
\label{0.19.4.3-fr}
设 $A$ 是一个拓扑环，$B$ 是一个拓扑 $A$-代数。假设对每个离散
$A$-代数 $C$ 以及每个理想 $\mathfrak J$（属于 $C$），若
$\mathfrak J^2=0$，则每个连续 $A$-同态 $u:B\to C/\mathfrak J$ 都可分解为
$B\xrightarrow{v}C\to C/\mathfrak J$，其中 $v$ 是一个连续 $A$-同态。
则 $B$ 是一个形式平滑的 $A$-代数。
\end{proposition}

\begin{proof}
事实上，仍沿用相同的记号，设 $\mathfrak R$ 是 $C$ 的任意一个幂零理想，
并考虑一个连续 $A$-同态 $u':B\to C/\mathfrak R$。假设
$\mathfrak R^m=0$，并置 $C_i=C/\mathfrak R^i$，其中
$1\leq i\leq m$；于是 $C_m=C$，而
$\mathfrak R^{i-1}/\mathfrak R^i$ 是一个平方为零的理想
$\mathfrak J_i$（属于 $C_i$），并且 $C_i=C_{i+1}/\mathfrak J_{i+1}$。此时，假设以及对
$i$ 的归纳给出一族连续 $A$-同态 $u_i:B\to C_i$，使得
$u_1=u'$，并且 $u_i$ 可分解为
$B\xrightarrow{u_{i+1}}C_{i+1}\to
C_{i+1}/\mathfrak J_{i+1}=C_i$；结论随即得出。
\SourceCorrectionNote{法文印本和外交式转录在归纳起点误写为 $u_1=u$，但本证明引入并需要提升的同态是 $u':B\to C/\mathfrak R=C_1$。起点为 $u_1=u'$。}
\end{proof}

\begin{proposition}[19.4.4]
\label{0.19.4.4-fr}
设 $A$ 是一个拓扑环，$B$ 是一个（交换）拓扑 $A$-代数。$B$ 是一个形式平滑的
$A$-代数，当且仅当对每个

\oldpage[0\textsubscript{IV-1}]{87}

离散拓扑 $B$-模 $L$，只要它被 $B$ 的一个开理想消去，就有（参见 (18.5.1)）
\[
  \operatorname{Exalcotop}_A(B,L)=0.
\]
\end{proposition}

\begin{proof}
设 $(\mathfrak R_\lambda)$ 是 $B$ 的一个递减开理想基本系，并置
$B_\lambda=B/\mathfrak R_\lambda$。考虑一个离散拓扑 $A$-代数 $C$ 和
一个理想 $\mathfrak J$（属于 $C$），满足 $\mathfrak J^2=0$；因而 $C$ 是一个
$A$-扩张，即 $C/\mathfrak J$ 被 $\mathfrak J$ 扩张所得的扩张。设给定一个
$A$-同态 $u:B_\lambda\to C/\mathfrak J$，并作一个 $A$-扩张
$E_\lambda$，即 $B_\lambda$ 被 $\mathfrak J$ 扩张所得的扩张；它是 $C$ 在同态 $u$ 下的
逆像 (18.2.5)；赋予离散拓扑之后，它是一个拓扑 $A$-代数。若
$\operatorname{Exalcotop}_A(B,L)=0$，则由定义 (18.5.1)，存在一个 $\mu$，
使得 $\mathfrak R_\mu\subset\mathfrak R_\lambda$，并且逆像扩张
$E_\mu$ 是 $A$-平凡的；但这意味着 (18.1.6)，存在一个连续 $A$-同态
$v:B_\mu\to E_\lambda$，使得典范同态 $B_\mu\to B_\lambda$ 可分解为
$B_\mu\xrightarrow{v}E_\lambda\to B_\lambda$。立即得出，复合
$B_\mu\to B_\lambda\xrightarrow{u}C/\mathfrak J$ 可分解为
$B_\mu\xrightarrow{v}E_\lambda\to C\to C/\mathfrak J$；由 (19.3.1) 和
(19.4.3)，这就证明了 $B$ 是一个形式平滑的 $A$-代数。反过来，取 $C$ 为
一个拓扑 $A$-代数，它是一个 $A$-扩张，即 $B_\lambda$ 被 $L$ 扩张所得的扩张，
并把 (19.3.1) 应用于恒等同态 $B_\lambda\to B_\lambda=C/L$，即可立即得到
逆命题。
\end{proof}

当 $A$ 和 $B$ 都是离散环时，形式平滑性判别法 (19.4.4) 化为
\begin{equation}
\tag{19.4.4.1}
\label{0.19.4.4.1-fr}
\operatorname{Exalcom}_A(B,L)=0\qquad\text{对每个 $B$-模 $L$}；
\end{equation}
换言之，\emph{每个交换 $A$-扩张，即 $B$ 被一个 $B$-模扩张所得的
扩张，都必须是 $A$-平凡的}。

\begin{corollary}[19.4.5]
\label{0.19.4.5-fr}
设 $A$ 是一个拓扑环，$B$ 是一个（交换）拓扑 $A$-代数。
\begin{enumerate}
  \item[(i)] 假设 $B$ 是一个形式平滑的 $A$-代数。于是，对每个开理想
  $\mathfrak R$（属于 $B$）、每个 $(B/\mathfrak R)$-模 $L$，以及每个
  $2$-次 Hochschild $A$-双线性对称上闭链 $f$，它从
  $(B/\mathfrak R)\times(B/\mathfrak R)$ 映入 $L$ (18.4.3)，存在一个开理想
  $\mathfrak R'\subset\mathfrak R$，使得当
  $\varphi:B/\mathfrak R'\to B/\mathfrak R$ 是典范同态时，
  $f\circ(\varphi\times\varphi)$ 是一个 $2$-次 Hochschild $A$-双线性上边缘，
  它从 $(B/\mathfrak R')\times(B/\mathfrak R')$ 映入 $L$。
  \item[(ii)] 若 $B$ 是一个形式投射的 $A$-模 (19.2.1)，并且满足 (i) 的条件，
  则 $B$ 是一个形式平滑的 $A$-代数。
\end{enumerate}
\end{corollary}

\begin{proof}
\begin{enumerate}
  \item[(i)] $2$-上闭链 $f$ 定义一个 Hochschild $A$-扩张 $C$，即
  $B/\mathfrak R$ 被 $L$ 扩张所得的扩张 (18.4.3)。把 (19.3.1) 应用于
  $C$、平方为零的理想 $L$（属于 $C$）以及恒等同态
  $B/\mathfrak R\to B/\mathfrak R=C/L$，再由 (18.4.3)，即得 (i) 的条件。
  \item[(ii)] 应用判别法 (19.4.3)。考虑一个开理想 $\mathfrak R$（属于 $B$）、
  一个开理想 $\mathfrak J$（属于 $A$），满足 $\mathfrak JB\subset\mathfrak R$，
  以及一个 $(A/\mathfrak J)$-扩张 $E$，即 $B/\mathfrak R$ 被 $L$ 扩张所得的
  扩张。由于 $B$ 是一个形式投射的 $A$-模，典范的
  连续 $A$-线性映射 $\rho:B\to B/\mathfrak R$ 可分解为
  $B\xrightarrow{w}E\to B/\mathfrak R$，其中 $w$ 是一个连续 $A$-线性映射
  (19.2.1)；因而该映射本身又可分解为
  $B\to B/\mathfrak R'\xrightarrow{w'}E$，其中 $\mathfrak R'$ 是 $B$ 的一个
  包含于 $\mathfrak R$ 的开理想。把 $\mathfrak J$ 换成一个更小的理想之后，
  可以假设 $\mathfrak JB\subset\mathfrak R'$。于是，在典范同态
  $B/\mathfrak R'\to B/\mathfrak R$ 下，扩张 $E$（即
  $B/\mathfrak R$ 被 $L$ 扩张所得的扩张）的逆像，是一个
  $(A/\mathfrak J)$-扩张 $E'$，即 $B/\mathfrak R'$ 被 $L$ 扩张所得的
  \emph{Hochschild} 扩张。

  \oldpage[0\textsubscript{IV-1}]{88}

  把 (i) 应用于定义这个扩张的一个上闭链 (18.4.3)，便得出：存在一个开理想
  $\mathfrak R''\subset\mathfrak R'$（属于 $B$），使得逆像 $E''$，即 $E$ 在
  $B/\mathfrak R''\to B/\mathfrak R$ 下的逆像，是 $A$-平凡的。证毕。
\end{enumerate}
\end{proof}

\begin{corollary}[19.4.6]
\label{0.19.4.6-fr}
设 $A$ 是一个拓扑环，$B$ 是一个拓扑 $A$-代数，并且作为 $A$-模是形式投射的。
设 $A'$ 是一个 $A$-代数，\emph{赋予由 $A$ 的拓扑导出的拓扑}。再假设
$A'$ 是一个忠实平坦 $A$-模，并且下列条件之一成立：
\begin{enumerate}
  \item[$1^\circ$] 存在一个开理想基本系 $(\mathfrak J_\lambda)$（属于 $A$）以及
  一个开理想基本系 $(\mathfrak M_\lambda)$（属于 $B$），它们具有同一个指标集，并且对
  每个 $\lambda$，都有 $\mathfrak J_\lambda B\subset\mathfrak M_\lambda$，而
  $B/\mathfrak M_\lambda$ 是一个有限型投射
  $(A/\mathfrak J_\lambda)$-模。
  \item[$2^\circ$] $A'$ 是一个有限型投射 $A$-模。
\end{enumerate}
于是，赋予张量积拓扑的 $B'=B\otimes_A A'$ 是一个形式平滑的 $A'$-代数，
当且仅当 $B$ 是一个形式平滑的 $A$-代数。
\end{corollary}

\begin{proof}
充分性已包含在 (19.3.5), (iii) 中，并且无须对 $B$ 或 $A'$ 加任何附加假设。
为证明逆命题，我们应用判别法 (19.4.5)。在假设 $2^\circ$ 下，仍以
$(\mathfrak M_\lambda)$ 表示 $B$ 的一个开理想基本系，并且对每个 $\lambda$，
以 $\mathfrak J_\lambda$ 表示 $A$ 的一个满足
$\mathfrak J_\lambda B\subset\mathfrak M_\lambda$ 的开理想。在两种情形中，都置
$A_\lambda=A/\mathfrak J_\lambda$、
$B_\lambda=B/\mathfrak M_\lambda$、
$A'_\lambda=A_\lambda\otimes_A A'$、
$B'_\lambda=B_\lambda\otimes_A A'$。设 $f$ 是一个 $2$-次 Hochschild
$A_\lambda$-双线性对称上闭链，它从 $B_\lambda\times B_\lambda$ 映入一个
$B_\lambda$-模 $L$；经扩张标量，得到一个 $2$-次 Hochschild 上闭链
$f'=f\otimes1$，它是 $A'_\lambda$-双线性对称的，并从
$B'_\lambda\times B'_\lambda$ 映入 $L'=L\otimes_A A'$。由于按假设，$B'$ 是一个
形式平滑的 $A'$-代数，由 (19.4.5), (i)，存在指标 $\mu$，使得
$\mathfrak M_\mu\subset\mathfrak M_\lambda$；并且，若
$\varphi':B'_\mu\to B'_\lambda$ 是典范映射，则
$f'\circ(\varphi'\times\varphi')$ 是一个 $2$-次 Hochschild 上边缘，它从
$B'_\mu\times B'_\mu$ 映入 $L'$；换言之，它的像 $c'$（属于 Hochschild 群
$\mathrm H^2_{A'_\mu}(B'_\mu,L')^s$）为零。显然，
若 $\varphi:B_\mu\to B_\lambda$ 是典范同态，而 $c$ 是
$f\circ(\varphi\times\varphi)$ 在 Hochschild 群
$\mathrm H^2_{A_\mu}(B_\mu,L)^s$ 中的类，则 $c'$ 是 $c$ 的典范像。
现在，若 $P_\bullet$ 是 (18.4.5) 中定义的、相对于环 $A_\mu$ 和 $B_\mu$
的复形，用来计算 $\mathrm H^2_{A_\mu}(B_\mu,L)^s$，则相对于环
$A'_\mu$ 和 $B'_\mu$ 的类似复形显然为 $P_\bullet\otimes_A A'$。在假设
$1^\circ$ 下，$P_\bullet$ 的构造表明，它是一个有限型投射 $A_\mu$-模。
因此，由 Bourbaki，\emph{Alg.}，第 II 章，第 3 版，§~5，n\textsuperscript{o}~3，
命题 7，得出在两种假设下，均有典范同构
\[
  \operatorname{Hom}_{A'_\mu}(P_\bullet\otimes_A A',L\otimes_A A')
  =\operatorname{Hom}_{A_\mu}(P_\bullet,L)\otimes_A A'
\]
（相差一个典范同构）；由于 $A'$ 是一个平坦 $A$-模，继而由 (18.4.5) 有
\[
  \mathrm H^2_{A'_\mu}(B'_\mu,L')^s
  =\mathrm H^2_{A_\mu}(B_\mu,L)^s\otimes_A A'
\]
因而可以写成 $c'=c\otimes1$；但由于 $A'$ 是一个忠实平坦 $A$-模，假设
$c'=0$ 蕴涵 $c=0$，证明完毕。
\end{proof}

\begin{proposition}[19.4.7]
\label{0.19.4.7-fr}
设 $A$ 是一个预可容拓扑环，$\mathfrak J$ 是 $A$ 的一个定义理想
（$\mathbf 0_{\rm I}$，7.1.2），而 $B$ 是一个拓扑 $A$-代数，并且作为
$A$-模是形式投射的。考虑拓扑商环 $A_0=A/\mathfrak J$、
$B_0=B/\mathfrak JB$。则 $B$ 是一个形式平滑的 $A$-代数，当且仅当
$B_0$ 是一个形式平滑的 $A_0$-代数。
\end{proposition}

\oldpage[0\textsubscript{IV-1}]{89}

\begin{proof}
必要性由 (19.4.2) 得出。为证明充分性，先注意：考虑一个 $0$ 的开邻域基本系
（属于 $A$），它由理想 $\mathfrak J_\alpha$ 组成，并且这些理想包含于
$\mathfrak J$；由
(19.4.2)，可以归结到 $A$ 离散的情形，因为 $B/\mathfrak J_\alpha B$ 是一个
形式投射的 $(A/\mathfrak J_\alpha)$-模 (19.2.5)。此时，由预可容环的定义，
$\mathfrak J$ 是\emph{幂零的}。此外，只须在 $\mathfrak J^2=0$ 时证明本命题；
事实上，若 $\mathfrak J^m=0$，就把命题依次应用于环
$A_k=A/\mathfrak J^k$、$B_k=B/\mathfrak J^kB$，以及理想
$\mathfrak J^{k-1}/\mathfrak J^k$（属于 $A_k$），其中 $k=2,3,\ldots,m$；并由 (19.2.5)
注意到 $B_k=B\otimes_A A_k$ 是一个形式投射的 $A_k$-模。

应用判别法 (19.4.5), (ii)。考虑一个开理想 $\mathfrak R$（属于 $B$），以及
一个 $2$-次 Hochschild $A$-双线性对称上闭链 $f$，它从
$(B/\mathfrak R)\times(B/\mathfrak R)$ 映入一个
$(B/\mathfrak R)$-模 $L$。先考虑特殊情形 $\mathfrak JL=0$；于是 $L$ 也可以
视为一个 $(B_0/\mathfrak RB_0)$-模，而 $f$ 可分解为
\[
  (B/\mathfrak R)\times(B/\mathfrak R)\to
  (B_0/\mathfrak RB_0)\times(B_0/\mathfrak RB_0)
  \xrightarrow{f_0}L,
\]
其中 $f_0$ 是一个双线性对称 Hochschild $2$-上闭链。因而由假设，存在一个
开理想 $\mathfrak R'\subset\mathfrak R$（属于 $B$）以及一个 $A_0$-线性映射
$g_0:B_0/\mathfrak R'B_0\to L$，使得
\[
  f_0(\varphi_0(x),\varphi_0(y))
  =xg_0(y)-g_0(xy)+g_0(x)y
\]
对 $x,y$（属于 $B_0/\mathfrak R'B_0$）成立；这里 $\varphi_0$ 是典范同态
$B_0/\mathfrak R'B_0\to B_0/\mathfrak RB_0$。立即得出，复合的
$A$-线性映射 $g:B/\mathfrak R'\to
B_0/\mathfrak R'B_0\xrightarrow{g_0}L$ 满足
$dg=f\circ(\varphi\times\varphi)$，其中
$\varphi:B/\mathfrak R'\to B/\mathfrak R$ 是典范同态。

现在转到一般情形。先考虑 $(B/\mathfrak R)$-模
$L/\mathfrak JL=L'$；它满足 $\mathfrak JL'=0$。若 $f'$ 是复合的
$A$-双线性映射
$(B/\mathfrak R)\times(B/\mathfrak R)\xrightarrow{f}L\to L'$，则 $f'$ 仍是
一个对称 Hochschild $2$-上闭链；由上述结果，存在一个开理想
$\mathfrak R'\subset\mathfrak R$（属于 $B$）以及一个 $A$-线性映射
$g':B/\mathfrak R'\to L'$，使得
$dg'=f'\circ(\varphi'\times\varphi')$；这里
$\varphi':B/\mathfrak R'\to B/\mathfrak R$ 是典范映射。由于 $B$ 是一个
形式投射的 $A$-模，存在一个开理想
$\mathfrak R'_1\subset\mathfrak R'$ 以及一个 $A$-线性映射
$g_1:B/\mathfrak R'_1\to L$，使得同态
$B/\mathfrak R'_1\to B/\mathfrak R'\xrightarrow{g'}L'$ 可分解为
$B/\mathfrak R'_1\xrightarrow{g_1}L\to L'$。于是考虑 Hochschild
$2$-上闭链
\[
  f_1(x,y)=f(\varphi_1(x),\varphi_1(y))-xg_1(y)+g_1(xy)-g_1(x)y,
\]
它是一个 $A$-双线性对称映射，从
$(B/\mathfrak R'_1)\times(B/\mathfrak R'_1)$ 映入 $L$。$g_1$ 的选取使得 $f_1$ 的取值落在
$\mathfrak JL$ 中。由于 $\mathfrak J(\mathfrak JL)=0$，可以再次应用第一种情形；
因而存在一个开理想 $\mathfrak R''\subset\mathfrak R'_1$ 以及一个
$A$-线性映射 $g_2:B/\mathfrak R''\to\mathfrak JL$，使得
\[
  f_1(\varphi_2(x),\varphi_2(y))=xg_2(y)-g_2(xy)+g_2(x)y
\]
在 $(B/\mathfrak R'')\times(B/\mathfrak R'')$ 中成立；这里
$\varphi_2:B/\mathfrak R''\to B/\mathfrak R'_1$ 是典范映射。因此，
$A$-线性映射 $g=g_2+g_1\circ\varphi_2:B/\mathfrak R''\to L$ 满足
$dg=f\circ(\varphi\times\varphi)$，其中
$\varphi:B/\mathfrak R''\to B/\mathfrak R$ 是典范映射。证毕。
\end{proof}

\subsection{形式平滑性与相伴分次环}
\label{subsection:0.19.5-fr}

\begin{env}[19.5.1]
\label{0.19.5.1-fr}
\oldpage[0\textsubscript{IV-1}]{90}
设 $C$ 是一个（交换）拓扑环，$V$ 是一个拓扑 $C$-模，并考虑对称代数
$\mathbf S_C^*(V)=\bigoplus_{n\geq0}\mathbf S_C^n(V)$；我们将典范地赋予它
一个与其 $C$-代数结构相容的\emph{线性拓扑}。为此，考虑一个开子模
$U$（属于 $V$），以及由它生成的\emph{分次理想}
$U.\mathbf S_C^*(V)$（属于 $\mathbf S_C^*(V)$）；作为 $0$ 在
$\mathbf S_C^*(V)$ 中的一个基本邻域系，我们取所有和
$\mathfrak a.\mathbf S_C^*(V)+U.\mathbf S_C^*(V)$，其中 $\mathfrak a$（相应地，
$U$）遍历 $C$ 的开理想基本系（相应地，$V$ 的开子模基本系）。应注意，
如果 $V$ 的拓扑比由 $C$ 的拓扑导出的拓扑粗，就可以只取偶对
$(\mathfrak a,U)$，并要求 $\mathfrak aV\subset U$，从而
$\mathfrak a.\mathbf S_C^*(V)+U.\mathbf S_C^*(V)
=\mathfrak a.1_C+U.\mathbf S_C^*(V)$；在这种情形下，每个
$\mathbf S_C^n(V)$（$n\geq1$）上的诱导拓扑以下列各项作为 $0$ 的基本邻域：
$U.\mathbf S_C^{n-1}(V)$，其中 $U$ 遍历 $V$ 的开子模；
特别地，在 $V=\mathbf S_C^1(V)$ 上，这个拓扑与给定拓扑一致（一般说来，
前者比后者粗）。此外，在所有情形下，如此定义的\emph{拓扑代数}
$\mathbf S_C^*(V)$，在拓扑 $C$-模范畴与拓扑 $C$-代数范畴中，给出与
$C$-模范畴和 $C$-代数范畴中的情形相同的泛问题之解。事实上，设 $E$ 是
一个拓扑 $C$-代数，$u:V\to E$ 是一个 $C$-模同态，而
$\mathbf S(u)=\widetilde u$ 是它到 $\mathbf S_C^*(V)$ 上的典范延拓。假设
$u$ 连续；若 $\mathfrak b$ 是 $E$ 的一个开理想，则它的原像
$U=u^{-1}(\mathfrak b)$ 是一个开 $C$-子模（属于 $V$），因而
$\widetilde u$ 将 $U.\mathbf S_C^*(V)$ 映入 $\mathfrak b$；另一方面，存在
一个开理想 $\mathfrak a$（属于 $C$），使得
$\mathfrak a.E\subset\mathfrak b$，从而
$\widetilde u(\mathfrak a.\mathbf S_C^*(V))\subset\mathfrak b$；这证明
$\widetilde u$ 连续。反过来，若 $\widetilde u$ 连续，而 $\mathfrak b$ 是
$E$ 的一个开理想，则存在一个开子模 $U$（属于 $V$），使得
$\widetilde u(U.\mathbf S_C^*(V))\subset\mathfrak b$；特别地，
$\widetilde u(U.1_C)\subset\mathfrak b$，即
$u(U)\subset\mathfrak b$，故 $u$ 连续。还要回顾，有拓扑 $C$-模（离散的）
的典范同构
\begin{equation}
\label{0.19.5.1.1-fr}
  \mathbf S_C^n(V)/U.\mathbf S_C^{n-1}(V)
  \xrightarrow{\sim}\mathbf S_C^n(V/U).
  \tag{19.5.1.1}
\end{equation}
\end{env}

\begin{env}[19.5.2]
\label{0.19.5.2-fr}
设 $A$ 是一个拓扑环，$B$ 是一个（交换）拓扑 $A$-代数，而
$\mathfrak J$ 是 $B$ 的一个理想（未必开，也未必闭）；在以下全部讨论中，
给 $\mathfrak J$ 以及 $\mathfrak J^n$（$n\geq2$）赋予由 $B$ 的拓扑诱导的
拓扑，并给商
$C=B/\mathfrak J$、
$\mathfrak J^n/\mathfrak J^{n+1}=\gr_{\mathfrak J}^n(B)$ 赋予商拓扑；于是
$\mathfrak J^n/\mathfrak J^{n+1}$ 是拓扑 $C$-模。典范单射
$\mathfrak J/\mathfrak J^2\to\gr_{\mathfrak J}^*(B)$ 延拓为一个（起初不考虑
拓扑的）$C$-代数同态
\[
  \varphi:\mathbf S_C^*(\mathfrak J/\mathfrak J^2)
  \longrightarrow\gr_{\mathfrak J}^*(B)
\]
它对每个 $n$ 给出一个\emph{满的} $C$-模同态
\begin{equation}
\label{0.19.5.2.1-fr}
  \varphi_n:\mathbf S_C^n(\mathfrak J/\mathfrak J^2)
  \longrightarrow\mathfrak J^n/\mathfrak J^{n+1}
  =\gr_{\mathfrak J}^n(B).
  \tag{19.5.2.1}
\end{equation}

\oldpage[0\textsubscript{IV-1}]{91}

当 $\mathbf S_C^n(\mathfrak J/\mathfrak J^2)$ 被赋予 (19.5.1) 中定义的拓扑时，
同态 $\varphi_n$ 是连续的。这可由
$\mathbf S_C^*(\mathfrak J/\mathfrak J^2)$ 的泛性质 (19.5.1) 得出：把它应用于
拓扑 $C$-代数
$E=\gr_{\mathfrak J/\mathfrak J^{n+1}}(B/\mathfrak J^{n+1})$（该代数被赋予
$\mathfrak J^k/\mathfrak J^{k+1}$ 的拓扑的乘积拓扑），以及典范单射
$\mathfrak J/\mathfrak J^2\to E$。注意，此处 $\mathfrak J/\mathfrak J^2$
的拓扑比由 $C$ 的拓扑导出的拓扑粗（后一个在
$\mathfrak J/\mathfrak J^2$ 上的拓扑也就是由 $B$ 的拓扑导出的）。
\SourceCorrectionNote{法文印本和外交式转录在上述分次代数中把两个截断指数都 写成 $n$；这样所得相伴分次代数的 $n$ 次分量为零，不能承接 $\varphi_n$ 的值域 $\mathfrak J^n/\mathfrak J^{n+1}$。两个指数都应为 $n+1$。}
\end{env}

\begin{theorem}[19.5.3]
\label{0.19.5.3-fr}
设 $A$ 是一个拓扑环，$B$ 是一个拓扑 $A$-代数，$\mathfrak J$ 是 $B$ 的一个
理想，而 $C=B/\mathfrak J$ 是商拓扑 $A$-代数。假设离散 $A$-代数 $C$
是形式平滑的。则：
\begin{enumerate}
  \item[\rm(i)] 若 $B$ 是一个形式平滑的 $A$-代数，则
  $\mathfrak J/\mathfrak J^2$ 是一个形式投射的拓扑 $C$-模 (19.2.1)。
  \item[\rm(ii)] 若 $B$ 是一个形式平滑的 $A$-代数，并且是一个预可容环
  （$\mathbf 0_{\rm I}$，7.1.2），则同态 $\varphi_n$ (19.5.2) 是形式双态射
  (19.1.2)。
  \item[\rm(iii)] 反过来，假设 $B$ 预可容，序列 $(\mathfrak J^n)$ 趋于
  $0$（在 $B$ 中），$\mathfrak J/\mathfrak J^2$ 是一个形式投射的 $C$-模，并且
  $\varphi_n$ 是形式双态射。则 $B$ 是一个形式平滑的 $A$-代数。
\end{enumerate}
\end{theorem}

这个定理的证明很长，并且充斥技术细节；因此我们先证明一个较简单的推论，
在其中主导思想表现得更为清楚；此外，这个推论正是定理 (19.5.3) 中往后最常
使用的特殊情形。

\begin{corollary}[19.5.4]
\label{0.19.5.4-fr}
设 $A$ 是一个拓扑环，$B$ 是一个拓扑 $A$-代数，$\mathfrak J$ 是 $B$ 的一个
理想，使得 $B$ 的拓扑是 $\mathfrak J$-预进拓扑。假设离散 $A$-代数
$C=B/\mathfrak J$ 是形式平滑的。则下列三个条件等价：
\begin{enumerate}
  \item[a)] $B$ 是一个形式平滑的 $A$-代数。
  \item[b)] $\mathfrak J/\mathfrak J^2$ 是一个投射 $C$-模，并且典范同态
  \[
    \varphi:\mathbf S_C^*(\mathfrak J/\mathfrak J^2)
    \longrightarrow\gr_{\mathfrak J}^*(B)
  \]
  是双射。
  \item[c)] 分离完备化 $\widehat B$（属于 $B$）是一个拓扑 $\widehat A$-代数，
  并且同构于一个 $\widehat A$-代数，其形式为 $\widehat{B'}$，其中
  $B'=\mathbf S_C^*(V)$，$V$ 是一个投射 $C$-模；这里给 $B'$ 赋予
  ${B'}^+$-预进拓扑，而 ${B'}^+$ 是增广理想。
\end{enumerate}
\end{corollary}

\begin{env}[19.5.4.1]
\label{0.19.5.4.1-fr}
先证明 $a)$ 蕴含 $\mathfrak J/\mathfrak J^2$ 是一个投射 $C$-模。为此，设
$P$ 和 $Q$ 是两个 $C$-模，$u:P\to Q$ 是一个满 $C$-同态，而
$v:\mathfrak J/\mathfrak J^2\to Q$ 是一个 $C$-同态。把环
$B/\mathfrak J^2$ 看作一个 $B$-扩张：它以 $C$ 为商、以
$\mathfrak J/\mathfrak J^2$ 为核；由 $v$ 从中推出一个 $B$-扩张
$E=(B/\mathfrak J^2)\oplus_{(\mathfrak J/\mathfrak J^2)}Q$，它以 $C$ 为商、
以 $Q$ 为核 (18.2.8)。因为
$C$ 是一个形式平滑的离散 $A$-代数，扩张 $E$ 是 $A$-\emph{平凡的}
(19.4.4)，因而可以等同于 $\mathrm D_C(Q)$ (18.2.3)。满同态
$u:P\to Q$ 随即典范地定义一个满 $A$-同态
$\psi:\mathrm D_C(P)\to\mathrm D_C(Q)$ (18.2.9)；它的核是一个理想
$\mathfrak L$（属于扩张 $E'=\mathrm D_C(P)$），包含于 $P$，因而当然平方为零。
设 $f:B\to E=E'/\mathfrak L$ 是定义 $B$-代数 $E$ 的结构的同态：由于
$\mathfrak L$ 平方为零，并且 $B$ 是一个形式平滑的 $A$-代数，$f$ 可分解为
$B\xrightarrow{g}E'\to E'/\mathfrak L$，其中 $g$ 是一个连续 $A$-同态。

\oldpage[0\textsubscript{IV-1}]{92}

图表
\[
  \xymatrix{
    E' \ar[r]^{\psi} & E \\
    B \ar[u]^{g} \ar[r]_{\theta} & B/\mathfrak J^2 \ar[u]_{v'}
  }
\]
是交换的，其中 $v'$ 由 $v$ 推出 (18.2.8)。此外，$\mathfrak J$ 在
$\psi\circ g$ 下的像包含于 $Q$，因此 $\mathfrak J$ 在 $g$ 下的像包含于
$P+\mathfrak L=P$，而 $\mathfrak J^2$ 在 $g$ 下的像为零。把 $g$ 和
$\theta$ 限制到 $\mathfrak J$，得到交换图表
\[
  \xymatrix{
    P \ar[r]^{u} & Q \\
    \mathfrak J/\mathfrak J^2 \ar[u]^{w} \ar[ur]_{v}
  }
\]
这就证明了 $C$-模 $\mathfrak J/\mathfrak J^2$ 是投射的。
\end{env}

\begin{env}[19.5.4.2]
\label{0.19.5.4.2-fr}
其次证明 $a)$ 蕴含 $\varphi$ 是双射的，这将完成 $a)$ 蕴含 $b)$ 的证明。置
$E_n=B/\mathfrak J^{n+1}$，并用 $F_n$ 表示
$\mathbf S_C^*(\mathfrak J/\mathfrak J^2)$ 对其增广理想的第 $(n+1)$ 次幂的商。
因为 $\mathfrak J/\mathfrak J^{n+1}$ 在 $E_n$ 中是幂零的，而 $C$ 是一个
形式平滑的离散 $A$-代数，$C$ 的恒等自同态可分解为
$C\xrightarrow{f}E_n\to C=E_n/(\mathfrak J/\mathfrak J^{n+1})$，其中 $f$ 是
一个 $A$-同态；另一方面，由 (19.5.4.1)，$\mathfrak J/\mathfrak J^2$ 是一个
投射 $C$-模，故 $\mathfrak J/\mathfrak J^2$ 的恒等自同态可分解为
$\mathfrak J/\mathfrak J^2\xrightarrow{g}\mathfrak J/\mathfrak J^{n+1}
\to\mathfrak J/\mathfrak J^2$，其中 $g$ 是一个 $C$-线性映射；由 $f$ 和 $g$
典范地得到一个 $C$-代数同态
$\mathbf S_C^*(\mathfrak J/\mathfrak J^2)\to E_n$；此外（由 $g$ 的定义），
它在第 $(n+1)$ 次幂上为零，其中该幂取自
$\mathbf S_C^*(\mathfrak J/\mathfrak J^2)$ 的增广理想；因而取商后得到一个满的
$A$-代数同态
$v:F_n\to E_n$，使得 $\gr^0(v)$ 和 $\gr^1(v)$ 分别是 $C$ 和
$\mathfrak J/\mathfrak J^2$ 的恒等自同态。由典范同态 $\varphi$ 的定义可见，
$\gr^j(v)=\varphi_j$ 对每个 $j\leq n$ 都成立。现在注意，核
$\mathfrak R$（属于 $v$）是 $F_n$ 的一个幂零理想，所以 $E_n$ 等同于
$F_n/\mathfrak R$。因为 $B$ 是一个形式平滑的 $A$-代数，连续的典范
$A$-同态 $p_n:B\to E_n=B/\mathfrak J^{n+1}$ 可分解为
$B\xrightarrow{w}F_n\xrightarrow{v}E_n$，其中 $w$ 是一个连续 $A$-同态；
因为 $\gr^0(v)$ 是恒等映射，$w(\mathfrak J)$ 包含于 $F_n$ 的增广理想，
从而 $w(\mathfrak J^{n+1})=0$；于是 $w$ 可分解为
$B\xrightarrow{p_n}B/\mathfrak J^{n+1}=E_n\xrightarrow{w'}F_n$，并且复合
$E_n\xrightarrow{w'}F_n\xrightarrow{v}E_n$ 是恒等映射。此外，由于
$\gr^0(v)$ 和 $\gr^1(v)$ 分别是 $C$ 和 $\mathfrak J/\mathfrak J^2$ 的恒等
自同态，$\gr^0(w')$ 和 $\gr^1(w')$ 也分别如此。因为 $\gr^*(E_n)$ 由
$\gr^1(E_n)$ 生成，复合同态
\[
  \gr^j(F_n)\xrightarrow{\varphi_j}\gr^j(E_n)
  \xrightarrow{\gr^j(w')}\gr^j(F_n)
\]
对每个 $j\leq n$ 都是恒等映射，因为对 $j=0$ 和 $j=1$ 已经如此；取
$j=n$，便证明了 $\varphi_n$ 是单射，从而完成 $a)$ 蕴含 $b)$ 的证明。
\SourceCorrectionNote{法文印本和外交式转录在上述 $w$ 的分解中写作 $B/\mathfrak J^n=E_n$；但本段已定义 $E_n=B/\mathfrak J^{n+1}$，并且 $w(\mathfrak J^{n+1})=0$ 正好给出经过该商的分解。指数为 $n+1$。}
\end{env}

\begin{env}[19.5.4.3]
\label{0.19.5.4.3-fr}
\oldpage[0\textsubscript{IV-1}]{93}
下面证明 $b)$ 蕴含 $a)$。与 (19.5.4.2) 开头相同的推理证明，存在一个满的
$A$-代数同态 $v:F_n\to E_n$，使得 $\gr^j(v)=\varphi_j$ 对每个 $j$ 都成立；
因为 $\varphi_j$ 对每个 $j$ 都是双射，并且 $F_n$ 和
$E_n$ 的滤过是有限的，所以 $v$ 是\emph{双射的}（Bourbaki，
\emph{Alg. comm.}，第三章，§~2，n\textsuperscript{o}~8，定理~1 的推论~3）。
现在设 $G$ 是一个离散拓扑 $A$-代数，$\mathfrak R$ 是 $G$ 中一个平方为零的
理想，而 $f:B\to G/\mathfrak R$ 是一个连续 $A$-代数同态。由于 $G$ 离散，
存在一个整数 $m$，使得 $f$ 在 $\mathfrak J^m$ 上为零，因而 $f$ 可分解为
$B\to E_n\xrightarrow{f_n}G/\mathfrak R$，这里取 $n=2m$。于是由复合得到一个
连续 $A$-代数同态
$r:C\to F_n\xrightarrow{v}E_n\xrightarrow{f_n}G/\mathfrak R$；因为 $C$
是一个形式平滑的离散 $A$-代数，$r$ 可分解为
$C\xrightarrow{r'}G\to G/\mathfrak R$，从而 $G$ 被 $r'$ 赋予一个 $C$-代数结构。
另一方面，限制到 $\mathfrak J/\mathfrak J^2$ 的同态
$f'_n:F_n\xrightarrow{v}E_n\xrightarrow{f_n}G/\mathfrak R$ 是一个 $C$-线性映射
$g:\mathfrak J/\mathfrak J^2\to G/\mathfrak R$。由于
$\mathfrak J/\mathfrak J^2$ 是一个投射 $C$-模，$g$ 可分解为
$\mathfrak J/\mathfrak J^2\xrightarrow{h}G\to G/\mathfrak R$，其中 $h$ 是
一个 $C$-线性映射；将其延拓，得到一个 $C$-代数同态
$w:\mathbf S_C^*(\mathfrak J/\mathfrak J^2)\to G$。由构造，每个次数为 $m$
的元素在 $w$ 下的像属于 $\mathfrak R$；因而，每个次数为 $n=2m$ 的元素之像
为零，因为 $\mathfrak R^2=0$；换言之，$w$ 可分解为
$\mathbf S_C^*(\mathfrak J/\mathfrak J^2)\to F_n\xrightarrow{w'}G$。由构造，
复合 $F_n\xrightarrow{w'}G\xrightarrow{\theta}G/\mathfrak R$ 与 $f'_n$ 一致：
这在 $C$ 和 $\mathfrak J/\mathfrak J^2$ 上成立，故该复合等于 $f'_n$。最后，复合
$B\to E_n\xrightarrow{v^{-1}}F_n\xrightarrow{w'}G\to G/\mathfrak R$ 等于
$f$，从而可见 $B$ 是一个形式平滑的 $A$-代数。
\end{env}

\begin{env}[19.5.4.4]
\label{0.19.5.4.4-fr}
还须证明 $a)$ 与 $c)$ 等价。首先，$c)$ 蕴含 $a)$：事实上，$B'$ 对离散拓扑
而言是一个形式平滑的 $C$-代数 (19.3.2)，因而对 ${B'}^+$-预进拓扑而言也
如此 (19.3.8)；由于 $C$ 是一个形式平滑的 $A$-代数，$B'$ 也是一个形式平滑的
$A$-代数 (19.3.5, (ii))，最终 $\widehat{B'}$ 是一个形式平滑的
$\widehat A$-代数 (19.3.6)，故 $B$ 是一个形式平滑的 $A$-代数 (19.3.6)。
还须验证 $a)$ 蕴含 $c)$。先注意，由于 $C$ 是一个形式平滑的 $A$-代数，
恒等同态 $C\to B/\mathfrak J$ 可分解为
$C\to\widehat B\to B/\mathfrak J$，其中 $C\to\widehat B$ 是一个 $A$-同态
(19.3.10)；因此，$\widehat B$ 以及随之所有 $B/\mathfrak J^{n+1}$ 都被赋予
$C$-代数结构。另一方面，因为 $\mathfrak J/\mathfrak J^2$ 是一个 $C$-模，
并且由 $b)$ 它是投射的，典范单射 $\mathfrak J/\mathfrak J^2\to B/\mathfrak J^2$ 使我们
能够构成一个 $C$-同态的投射系
$u_n:\mathfrak J/\mathfrak J^2\to B/\mathfrak J^{n+1}$（$n\geq1$），从而也
能够（借助对称代数的泛性质）构成一个 $C$-代数同态的投射系
$v_n:\mathbf S_C^*(\mathfrak J/\mathfrak J^2)=B'\to
B/\mathfrak J^{n+1}$；此外，显然 $v_n$ 在 $({B'}^+)^{n+1}$ 上为零，而由于
$\gr^0(v_n)=\varphi_0$ 且 $\gr^1(v_n)=\varphi_1$，有
$\gr^j(v_n)=\varphi_j$（$0\leq j\leq n$）。因为各 $\varphi_j$ 由 $b)$ 都是
同构，并且 $B'/({B'}^+)^{n+1}$ 与 $B/\mathfrak J^{n+1}$ 的滤过是有限的，
所以 $v_n$ 对

\oldpage[0\textsubscript{IV-1}]{94}

每个 $n$ 都是\emph{双射的}（Bourbaki，\emph{Alg. comm.}，第三章，§~2，
n\textsuperscript{o}~8，定理~1 的推论~3）；取投射极限便得到 $c)$。
\end{env}

\begin{remark}[19.5.5]
\label{0.19.5.5-fr}
在 (19.5.4) 的一般假设下，进一步假设 $\mathfrak J$ 是 $B$ 的一个
\emph{极大}理想，使得 $C=B/\mathfrak J$ 是一个域，并且
$\mathfrak J/\mathfrak J^2$ 是一个\emph{有限}维 $C$-向量空间。于是 (19.5.4)
中的条件 $a)$、$b)$ 和 $c)$ 还等价于下列条件：
\begin{enumerate}
  \item[d)] \emph{给定一个多项式代数} $E=C[T_1,\ldots,T_n]$
  \emph{（定义在域} $C$ \emph{上），并用} $\mathfrak n$ \emph{表示} $E$
  \emph{的由诸} $T_i$ $(1\leq i\leq n)$ \emph{生成的极大理想，则对每个理想}
  $\mathfrak b$ \emph{（属于} $E$\emph{），只要它包含某个幂}
  $\mathfrak n^k$，\emph{每个同态} $v:B\to E/\mathfrak b$ \emph{（它是
  $A$-代数的 $C$-增广同态）都可分解为}
  $B\xrightarrow{w}E/\mathfrak n^k\to E/\mathfrak b$，\emph{其中} $w$
  \emph{是一个 $A$-同态}。
\end{enumerate}
事实上，此处 $\mathfrak J/\mathfrak J^2$ 是一个自由 $C$-模，因此只须验证条件
$d)$ 蕴含典范同态
$\varphi:\mathbf S_C^*(\mathfrak J/\mathfrak J^2)\to\gr_{\mathfrak J}^*(B)$
是双射。这里 $\mathbf S_C^*(\mathfrak J/\mathfrak J^2)$ 等同于多项式代数
$E=C[T_1,\ldots,T_n]$，其中
$n=\operatorname{rg}_C(\mathfrak J/\mathfrak J^2)$，而 $E$ 的增广理想就是理想
$\mathfrak n$，它由诸 $T_i$ 生成。对每个整数 $k\geq0$，$C$ 是一个形式平滑的
离散 $A$-代数这一假设如同在 (19.5.4.2) 中一样，给出一个满 $A$-同态
$f:E/\mathfrak n^{k+1}\to B/\mathfrak J^{k+1}$，使得
$\gr^j(f)=\varphi_j$ 对每个 $j\leq k$ 都成立。若
$\mathfrak b/\mathfrak n^{k+1}=\operatorname{Ker}(f)$，则
$B/\mathfrak J^{k+1}$ 等同于 $E/\mathfrak b$，而假设 $d)$ 使典范同态
$B\to B/\mathfrak J^{k+1}$ 可分解为
$B\xrightarrow{w}E/\mathfrak n^{k+1}\xrightarrow{f}B/\mathfrak J^{k+1}$；
因为 $\gr^0(f)$ 是恒等映射，有
$w(\mathfrak J)\subset\mathfrak n/\mathfrak n^{k+1}$，从而
$w(\mathfrak J^{k+1})=0$；如同在 (19.5.4.2) 中那样，得出
$\varphi_k$ 对每个 $k$ 都是单射的，这就完成了断言的证明。
\end{remark}

\begin{env}[19.5.6]
\label{0.19.5.6-fr}
\emph{定理} (19.5.3) \emph{的证明}。--- 设
$(\mathfrak b_\alpha)$ 是 $B$ 的一个递减的开理想基本系。置
\[
  B_\alpha=B/\mathfrak b_\alpha,\qquad
  C_\alpha=B/(\mathfrak b_\alpha+\mathfrak J),\qquad
  \mathfrak J_\alpha=(\mathfrak b_\alpha+\mathfrak J)/\mathfrak b_\alpha,
\]
从而
\[
  C_\alpha=B_\alpha/\mathfrak J_\alpha\qquad\text{且}\qquad
  \mathfrak J_\alpha^n/\mathfrak J_\alpha^{n+1}
  =(\mathfrak b_\alpha+\mathfrak J^n)/(\mathfrak b_\alpha+\mathfrak J^{n+1}).
\]
诸 $C$-模
$((\mathfrak b_\alpha\cap\mathfrak J^n)+\mathfrak J^{n+1})/
\mathfrak J^{n+1}$，作为 $\mathfrak J^n/\mathfrak J^{n+1}$ 的子模，构成拓扑 $C$-模
$\mathfrak J^n/\mathfrak J^{n+1}$ 的一个开子模基本系；由于
$(\mathfrak b_\alpha\cap\mathfrak J^n)+\mathfrak J^{n+1}
=\mathfrak J^n\cap(\mathfrak b_\alpha+\mathfrak J^{n+1})$，有
\[
  \mathfrak J^n/((\mathfrak b_\alpha\cap\mathfrak J^n)+\mathfrak J^{n+1})
  =\mathfrak J^n/(\mathfrak J^n\cap(\mathfrak b_\alpha+\mathfrak J^{n+1}))
  =(\mathfrak b_\alpha+\mathfrak J^n)/(\mathfrak b_\alpha+\mathfrak J^{n+1})
  =\mathfrak J_\alpha^n/\mathfrak J_\alpha^{n+1}.
\]
\end{env}

\begin{env}[19.5.6.1]
\label{0.19.5.6.1-fr}
设 $P$、$Q$ 是两个离散 $C_\alpha$-模，而
$u:P\to Q$ 是一个满 $C_\alpha$-同态。离散环
$B/(\mathfrak b_\alpha+\mathfrak J^2)$ 是一个 $B$-扩张：它以
$C_\alpha$ 为商、以平方为零理想
$\mathfrak J_\alpha/\mathfrak J_\alpha^2$ 为核。设
$v:\mathfrak J/\mathfrak J^2\to Q$ 是一个连续 $C$-同态；必要时把
$\mathfrak b_\alpha$ 替换为 $B$ 的一个更小开理想，便可假设 $v$ 的核包含开
$C$-子模
$((\mathfrak b_\alpha\cap\mathfrak J)+\mathfrak J^2)/\mathfrak J^2$；取商后，
由此得到离散模的一个 $C_\alpha$-同态
$v':\mathfrak J_\alpha/\mathfrak J_\alpha^2\to Q$。设 $E_\alpha$ 是一个
$B_\alpha$-扩张，即 $C_\alpha$ 被 $Q$ 扩张所得的扩张；它由
$B/(\mathfrak b_\alpha+\mathfrak J^2)$ 借助 $v'$ 推出。并设
$v'_*:B/(\mathfrak b_\alpha+\mathfrak J^2)\to E_\alpha$ 是相应的
$B_\alpha$-同态 (18.2.8)。$E_\alpha$ 是一个离散的拓扑 $A$-代数，而典范同构
$C_\alpha\xrightarrow{\sim}E_\alpha/Q$ 通过复合给出一个连续 $A$-同态
$f:C\to C_\alpha\to E_\alpha/Q$。然而，由于 $Q$ 在 $E_\alpha$ 中平方为零，
并且 $C$ 是一个形式平滑的 $A$-代数，$f$ 可分解为
$C\xrightarrow{g}E_\alpha\to E_\alpha/Q$，其中 $g$ 是一个连续 $A$-同态。
由于 $g$ 连续而 $E_\alpha$ 离散，存在 $\beta\geq\alpha$，使得 $g$ 可分解为
$C\to C_\beta\xrightarrow{g'}E_\alpha$。另一方面，设 $E_\beta$ 是一个
$A$-扩张，即 $C_\beta$ 被 $Q$ 扩张所得的扩张；它是 $E_\alpha$ 在典范同态
$C_\beta\to C_\alpha$ 下的逆像。$g'$ 的存在（其中复合
$C_\beta\to E_\alpha\to C_\alpha$ 是典范同态）等价于 $E_\beta$ 是一个
$A$-\emph{平凡的}扩张，因而可以等同于 $\mathrm D_{C_\beta}(Q)$。于是，满同态
$u:P\to Q$ 典范地定义一个满同态
$\mathrm D_{C_\beta}(P)\to\mathrm D_{C_\beta}(Q)$ (18.2.9)；它的核是一个理想
$\mathfrak L$，属于扩张 $E'_\beta=\mathrm D_{C_\beta}(P)$，并且包含于 $P$，
因而当然平方为零。设
$f':B\to E_\beta=E'_\beta/\mathfrak L$ 是一个连续 $A$-同态，它定义拓扑
$B$-代数 $E_\beta$ 的结构；由于 $\mathfrak L$ 平方为零，并且 $B$ 是一个
形式平滑的 $A$-代数，$f'$ 可分解为
$B\xrightarrow{h'}E'_\beta\to E'_\beta/\mathfrak L$，其中 $h'$ 是一个连续
$A$-同态。$E'_\beta$ 的构造通过复合给出一个 $A$-同态
$\psi:E'_\beta\to E_\beta\to E_\alpha$，并且图表
\[
  \xymatrix{
    E'_\beta \ar[r]^{\psi} & E_\alpha \\
    B \ar[u]^{h'} \ar[r]_{\theta} &
    B/(\mathfrak b_\alpha+\mathfrak J^2) \ar[u]_{v'_*}
  }
\]
显然是交换的。此外，在 $\psi\circ h'$ 下，$\mathfrak J$ 的像包含于 $Q$，
所以在 $h'$ 下，$\mathfrak J$ 的像包含于 $P$，而 $\mathfrak J^2$ 在 $h'$ 下的
像为零。把 $h'$ 和 $\theta$ 限制到 $\mathfrak J$，便得到交换图表
\[
  \xymatrix{
    P \ar[r]^{u} & Q \\
    \mathfrak J/\mathfrak J^2 \ar[u]^{w} \ar[ur]_{v}
  }
\]
其中 $w$ 是连续的；这就证明了 $\mathfrak J/\mathfrak J^2$ 是一个形式投射的
$C$-模。
\end{env}

\oldpage[0\textsubscript{IV-1}]{95}

\begin{env}[19.5.6.2]
\label{0.19.5.6.2-fr}
对每个整数 $n\geq0$，置
\[
  E_n=B/\mathfrak J^{n+1},
  \qquad\text{从而}\qquad E_0=C~;
\]
诸理想 $(\mathfrak b_\alpha+\mathfrak J^{n+1})/\mathfrak J^{n+1}$ 构成
$E_n$ 中的一个开理想基本系，并置
\[
  E_{\alpha,n}=B/(\mathfrak b_\alpha+\mathfrak J^{n+1})
  =B_\alpha/\mathfrak J_\alpha^{n+1},
\]
这是一个离散商环。同样，在
$\gr_{\mathfrak J}^n(B)=\mathfrak J^n/\mathfrak J^{n+1}$ 中，已经看到诸
$((\mathfrak b_\alpha\cap\mathfrak J^n)+\mathfrak J^{n+1})/
\mathfrak J^{n+1}$ 构成 $0$ 的一个邻域基本系，而
$\mathfrak J^n/\mathfrak J^{n+1}$ 对这个子模的商典范地等同于
$\mathfrak J_\alpha^n/\mathfrak J_\alpha^{n+1}$。考虑对称代数
$\mathbf S_{C_\alpha}^*(\mathfrak J_\alpha/\mathfrak J_\alpha^2)$；用
$F_{\alpha,n}$ 表示
$\mathbf S_{C_\alpha}^*(\mathfrak J_\alpha/\mathfrak J_\alpha^2)$ 对其增广理想的
第 $(n+1)$ 次幂的商。对固定的 $n\geq1$，由 (19.5.1.1) 可知，诸
$\mathbf S_{C_\alpha}^n(\mathfrak J_\alpha/\mathfrak J_\alpha^2)$ 是
$\mathbf S_C^n(\mathfrak J/\mathfrak J^2)$ 对该拓扑 $C$-模中的一个开子模基本系
取商所得的商模。

为简化用语，若 $\alpha\leq\beta$，则称典范同态
$B_\beta\to B_\alpha$、$C_\beta\to C_\alpha$、
$\mathfrak J_\beta/\mathfrak J_\beta^2\to
\mathfrak J_\alpha/\mathfrak J_\alpha^2$、$E_{\beta,n}\to E_{\alpha,n}$、
$F_{\beta,n}\to F_{\alpha,n}$ 等为\emph{转移同态}。
\end{env}

\begin{lemma}[19.5.6.3]
\label{0.19.5.6.3-fr}
假设 $A$-代数 $C$ 是形式平滑的，环 $B$ 是预可容的，并且 $C$-模
$\mathfrak J/\mathfrak J^2$ 是形式投射的。于是：
\begin{enumerate}
  \item[\rm(i)] 对每个 $\alpha$，存在 $\beta\geq\alpha$ 以及一个满
  $A$-代数同态
  \[
    v_{\alpha\beta}:F_{\beta,n}\longrightarrow E_{\alpha,n}
  \]
  使得 $\gr^0(v_{\alpha\beta}):C_\beta\to C_\alpha$ 和
  $\gr^1(v_{\alpha\beta}):\mathfrak J_\beta/\mathfrak J_\beta^2
  \to\mathfrak J_\alpha/\mathfrak J_\alpha^2$ 是转移同态。
  \item[\rm(ii)] 若 $\beta$ 和 $v_{\alpha\beta}$ 满足 \emph{(i)} 的条件，
  则对每个 $\gamma\geq\beta$，存在 $\delta\geq\gamma$ 以及一个满
  $A$-代数同态 $v_{\gamma\delta}:F_{\delta,n}\to E_{\gamma,n}$，它满足
  \emph{(i)} 的条件（对于 $\gamma$ 和 $\delta$），并使图表
  \begin{equation}
  \label{0.19.5.6.4-fr}
  \xymatrix@C=3em{
    F_{\beta,n} \ar[r]^{v_{\alpha\beta}} & E_{\alpha,n} \\
    F_{\delta,n} \ar[u] \ar[r]_{v_{\gamma\delta}} & E_{\gamma,n} \ar[u]
  }
  \tag{19.5.6.4}
  \end{equation}
  交换；其中竖直箭头是转移同态。
\end{enumerate}
\end{lemma}

\begin{env}[19.5.6.4]
\label{0.19.5.6.4-proof-fr}
\begin{enumerate}
  \item[\rm(i)] 在离散拓扑 $A$-代数 $E_{\alpha,n}$ 中，理想
  $\mathfrak J_\alpha/\mathfrak J_\alpha^{n+1}$ 是幂零的，而恒等同构
  $C_\alpha\xrightarrow{\sim}
  E_{\alpha,n}/(\mathfrak J_\alpha/\mathfrak J_\alpha^{n+1})$ 通过复合给出一个
  连续 $A$-同态
  $C\to C_\alpha\to
  E_{\alpha,n}/(\mathfrak J_\alpha/\mathfrak J_\alpha^{n+1})$；

\oldpage[0\textsubscript{IV-1}]{96}

  由于 $C$ 是一个形式平滑的 $A$-代数，这个同态可分解为
  $C\xrightarrow{f_\alpha}E_{\alpha,n}\to
  E_{\alpha,n}/(\mathfrak J_\alpha/\mathfrak J_\alpha^{n+1})$，其中 $f_\alpha$
  是一个连续 $A$-同态；因此，
  $\mathfrak J_\alpha/\mathfrak J_\alpha^{n+1}$ 借助 $f_\alpha$ 成为一个离散拓扑 $C$-模，并被
  $C$ 的一个开理想零化。$\mathfrak J/\mathfrak J^2$ 是形式投射的 $C$-模这一
  假设于是蕴含：存在一个连续 $C$-线性映射
  $g_\alpha:\mathfrak J/\mathfrak J^2\to
  \mathfrak J_\alpha/\mathfrak J_\alpha^{n+1}$，使图表
  \[
    \xymatrix{
      \mathfrak J_\alpha/\mathfrak J_\alpha^{n+1} \ar[r] &
      \mathfrak J_\alpha/\mathfrak J_\alpha^2 \\
      & \mathfrak J/\mathfrak J^2 \ar[ul]^{g_\alpha} \ar[u]
    }
  \]
  交换。由于 $f_\alpha$ 和 $g_\alpha$ 连续，存在 $\beta\geq\alpha$，使这些
  同态分别可分解为
  \[
    C\to C_\beta\xrightarrow{f'_{\alpha\beta}}E_{\alpha,n},
  \]
  \[
    \mathfrak J/\mathfrak J^2\to
    \mathfrak J_\beta/\mathfrak J_\beta^2
    \xrightarrow{g'_{\alpha\beta}}
    \mathfrak J_\alpha/\mathfrak J_\alpha^{n+1}
  \]
  因而由 $f'_{\alpha\beta}$ 和 $g'_{\alpha\beta}$ 典范地得到一个
  $C_\beta$-代数同态
  \[
    \mathbf S_{C_\beta}^*(\mathfrak J_\beta/\mathfrak J_\beta^2)
    \longrightarrow E_{\alpha,n}
  \]
  此外，按 $g'_{\alpha\beta}$ 的定义，这个同态在第 $(n+1)$ 次幂上为零；
  这里该幂取自 $\mathbf S_{C_\beta}^*(\mathfrak J_\beta/\mathfrak J_\beta^2)$
  的增广理想。对这个第 $(n+1)$ 次幂取商，便得到所求同态
  $v_{\alpha\beta}$；这里已顾及 $f_\alpha$ 和 $g_\alpha$ 的定义。事实上，
  $v_{\alpha\beta}$ 的满性由两个同态 $\gr^0(v_{\alpha\beta})$ 和
  $\gr^1(v_{\alpha\beta})$ 的满性得出，因为这蕴含 $\gr(v_{\alpha\beta})$
  是满的（分次代数 $\gr^*(E_{\alpha,n})$ 由 $\gr^0(E_{\alpha,n})$ 和
  $\gr^1(E_{\alpha,n})$ 生成）；而所考虑的滤过都是有限的，所以可以应用
  Bourbaki，\emph{Alg. comm.}，第三章，§~2，n\textsuperscript{o}~8，
  定理~1 的推论~3。

  \item[\rm(ii)] $B$ 是\emph{预可容的}这一假设意味着，可以假设所有
  $\mathfrak b_\alpha$ 都包含于同一个 $\mathfrak b_{\alpha_0}$，后者的幂
  \emph{趋于} $0$。这特别蕴含：转移同态
  $E_{\gamma,n}\to E_{\alpha,n}$ 的核等于
  $(\mathfrak b_\alpha+\mathfrak J^{n+1})/
  (\mathfrak b_\gamma+\mathfrak J^{n+1})$，并且是\emph{幂零的}；应用引理
  (19.3.10.1)，可见可以选择 $f_\gamma$，使图表
  \[
    \xymatrix{
      & E_{\alpha,n} \\
      C \ar[ur]^{f_\alpha} \ar[r]_{f_\gamma} & E_{\gamma,n} \ar[u]
    }
  \]
  交换（其中竖直箭头是转移同态）。另一方面，
  $\mathfrak J/\mathfrak J^2$ 是形式投射的 $C$-模这一假设，使得可以选择
  $g_\gamma:\mathfrak J/\mathfrak J^2\to
  \mathfrak J_\gamma/\mathfrak J_\gamma^{n+1}$，使图表
  \[
    \xymatrix@C=5em@R=3em{
      \mathfrak J_\alpha/\mathfrak J_\alpha^{n+1} \ar[r] &
      \mathfrak J_\alpha/\mathfrak J_\alpha^2 \\
      \mathfrak J_\gamma/\mathfrak J_\gamma^{n+1} \ar[u] \ar[r] &
      \mathfrak J_\gamma/\mathfrak J_\gamma^2 \ar[u] \\
      & \mathfrak J/\mathfrak J^2
        \ar[ul]_{g_\gamma} \ar[u]_{p_\gamma} \ar[uul]^{g_\alpha}
    }
  \]
  交换（只须注意：在 $(g_\alpha,p_\gamma)$ 下，$\mathfrak J/\mathfrak J^2$
  在乘积模 $(\mathfrak J_\alpha/\mathfrak J_\alpha^{n+1})
  \times(\mathfrak J_\gamma/\mathfrak J_\gamma^2)$ 中的像，鉴于
  $g_\alpha$ 的定义和关系 $\beta\leq\gamma$，包含于典范像 $Q$；这个像由模
  $P=\mathfrak J_\gamma/\mathfrak J_\gamma^{n+1}$ 映入该乘积模所得；再把定义
  (19.2.1) 应用于满同态 $P\to Q$ 以及同态 $(g_\alpha,p_\gamma)$，后者从
  $\mathfrak J/\mathfrak J^2$ 映到 $Q$ 即可）。如此选择 $f_\gamma$ 和
  $g_\gamma$ 后，如 \emph{(i)} 中那样构造 $v_{\gamma\delta}$，还可使图表
  (19.5.6.4) 交换。
\end{enumerate}
\end{env}

\oldpage[0\textsubscript{IV-1}]{97}

\begin{lemma}[19.5.6.5]
\label{0.19.5.6.5-fr}
假设 $A$-代数 $B$ 和 $C$ 都是形式平滑的，并且 $B$ 是预可容的（从而根据
(19.5.6.1)，(19.5.6.3) 的条件得到满足）。对于每一对指标
$\alpha\leq\beta$ 和每一个满足 (19.5.6.3, (i)) 条件的同态
$v_{\alpha\beta}$，存在一个指标 $\lambda\geq\beta$ 以及一个 $A$-代数同态
\[
  w_{\beta\lambda}:E_{\lambda,n}\longrightarrow F_{\beta,n}
\]
使得：$1^\circ$ $\gr^0(w_{\beta\lambda}):C_\lambda\to C_\beta$ 和
$\gr^1(w_{\beta\lambda}):
\mathfrak J_\lambda/\mathfrak J_\lambda^2\to
\mathfrak J_\beta/\mathfrak J_\beta^2$ 是转移同态；$2^\circ$ 复合
$E_{\lambda,n}\xrightarrow{w_{\beta\lambda}}F_{\beta,n}
\xrightarrow{v_{\alpha\beta}}E_{\alpha,n}$ 是转移同态。
\end{lemma}

\begin{proof}
把引理 (19.5.6.3, (ii)) 应用于 $\gamma=\beta$，便得到一个
$\delta\geq\beta$ 以及一个
$v_{\beta\delta}:F_{\delta,n}\to E_{\beta,n}$。回顾一下，
$v_{\beta\delta}$ 是满的；另一方面，它的核 $\mathfrak R$ 是\emph{幂零的}：
事实上，增广理想 $F_{\delta,n}^+$（属于 $F_{\delta,n}$）按定义是幂零的，而
$\gr^0(v_{\beta\delta}):C_\delta\to C_\beta$ 的核也由于 $B$ 是预可容的这一
假设而幂零。因此可见 $E_{\beta,n}$ 等同于
$F_{\delta,n}/\mathfrak R$。由于 $B$ 是一个形式平滑的 $A$-代数，连续的
典范 $A$-同态 $p_{\beta,n}:B\to E_{\beta,n}$ 可分解为
$B\xrightarrow{w}F_{\delta,n}\xrightarrow{v_{\beta\delta}}E_{\beta,n}$，其中
$w$ 是一个连续 $A$-同态。因为 $F_{\delta,n}$ 是离散的，存在
$\lambda\geq\delta$，使 $w$ 在 $\mathfrak b_\lambda$ 上为零，因而 $w$ 可分解为
$B\to B/\mathfrak b_\lambda\xrightarrow{w'_\lambda}F_{\delta,n}$。考虑复合同态
$B/\mathfrak b_\lambda\xrightarrow{w'_\lambda}F_{\delta,n}
\xrightarrow{q_{\beta\delta}}F_{\beta,n}$，其中 $q_{\beta\delta}$ 是转移同态；
注意，$\gr^0(q_{\beta\delta})=\gr^0(v_{\beta\delta}):C_\delta\to C_\beta$
是转移同态。因为复合
$v_{\alpha\beta}\circ q_{\beta\delta}\circ w'_\lambda$ 是典范同态
$B/\mathfrak b_\lambda\to
B/(\mathfrak b_\alpha+\mathfrak J^{n+1})$，这表明
$\mathfrak J_\lambda=(\mathfrak b_\lambda+\mathfrak J)/\mathfrak b_\lambda$
在 $q_{\beta\delta}\circ w'_\lambda$ 下的像包含于增广理想
$F_{\beta,n}^+$，从而
$\mathfrak J_\lambda^{n+1}=(\mathfrak b_\lambda+\mathfrak J^{n+1})/
\mathfrak b_\lambda$ 的像为零。换言之，$q_{\beta\delta}\circ w'_\lambda$
可分解为
\[
  B/\mathfrak b_\lambda\to
  B/(\mathfrak b_\lambda+\mathfrak J^{n+1})=E_{\lambda,n}
  \xrightarrow{w_{\beta\lambda}}F_{\beta,n}
\]
使得复合
$E_{\lambda,n}\xrightarrow{w_{\beta\lambda}}F_{\beta,n}
\xrightarrow{v_{\alpha\beta}}E_{\alpha,n}$ 是转移同态。前面的论证还表明，
$\gr^0(w_{\beta\lambda})$，即复合
$B/(\mathfrak b_\lambda+\mathfrak J)\to\gr^0(F_{\delta,n})
\xrightarrow{\gr^0(q_{\beta\delta})}\gr^0(F_{\beta,n})$，是转移同态
$C_\lambda\to C_\beta$，因为 $v_{\beta\delta}\circ w$ 是典范同态。此外还有
$\gr^1(q_{\beta\delta})=\gr^1(v_{\beta\delta})$ (19.5.6.3)，所以同一论证证明
$\gr^1(w_{\beta\lambda})$ 是转移同态。
\end{proof}

\begin{env}[19.5.6.6]
\label{0.19.5.6.6-fr}
\emph{(19.5.3, (ii)) 的证明}。--- 为证明 $\varphi_n$ 是一个形式双态射，
将应用 (19.1.3, (iii)) 的判别法。由于按假设，引理 (19.5.6.5) 的条件得到满足，
故对每个指标 $\alpha$，取满足该引理结论的 $v_{\alpha\beta}$ 和
$w_{\beta\lambda}$。同态
\[
  \gr^n(v_{\alpha\beta}):\gr^n(F_{\beta,n})
  \longrightarrow\gr^n(E_{\alpha,n})
\]
无非是同态
\[
  \varphi_{\alpha\beta,n}:
  \mathbf S_{C_\beta}^n(\mathfrak J_\beta/\mathfrak J_\beta^2)
  \longrightarrow\mathfrak J_\alpha^n/\mathfrak J_\alpha^{n+1}
\]
它由典范同态 $\varphi_n$ (19.5.2.1) 取商得到；事实上，由 (19.5.6.3) 可知
$\gr^0(v_{\alpha\beta})$ 和 $\gr^1(v_{\alpha\beta})$ 分别与
$\varphi_{\alpha\beta,0}$ 和 $\varphi_{\alpha\beta,1}$ 一致，而典范同态
$\varphi$ 的定义于是通过对 $j\leq n$ 的归纳表明，
$\gr^j(v_{\alpha\beta})$ 和 $\varphi_{\alpha\beta,j}$ 对每个 $j$ 都相等。既然如此，由于
$\varphi_n$ 是满的，它当然是一个形式满态射；此外，对于 $j\leq n$，复合同态
\[
  \gr^j(F_{\lambda,n})\xrightarrow{\varphi_{\lambda\lambda,j}}
  \gr^j(E_{\lambda,n})\xrightarrow{\gr^j(w_{\beta\lambda})}
  \gr^j(F_{\beta,n})
\]
是转移同态，因为根据 (19.5.6.5)，这对 $j=0$ 和 $j=1$ 成立，而
$\gr^*(E_{\lambda,n})$ 由 $\gr^1(E_{\lambda,n})$ 生成，所以对 $j$ 的归纳证明了
这个断言。再与转移同态
$\gr^j(F_{\beta,n})\to\gr^j(F_{\alpha,n})$ 复合，便可看出，对于 $j=n$，
转移同态 $\gr^n(F_{\lambda,n})\to\gr^n(F_{\alpha,n})$ 被分解为
\[
  \gr^n(F_{\lambda,n})\xrightarrow{\varphi_{\lambda\lambda,n}}
  \gr^n(E_{\lambda,n})\longrightarrow\gr^n(F_{\alpha,n}),
\]
这正是判别法 (19.1.3) 中使 $\varphi_n$ 成为形式双态射的条件。
\end{env}

\oldpage[0\textsubscript{IV-1}]{98}

\begin{lemma}[19.5.6.7]
\label{0.19.5.6.7-fr}
假设 $B$ 是预可容的，$\mathfrak J/\mathfrak J^2$ 是一个形式投射的 $C$-模，
$C$ 是一个形式平滑的 $A$-代数，并且诸 $\varphi_n$ 是形式双态射。于是，对于
每一对指标 $\alpha$、$\beta$，只要 $\alpha\leq\beta$，以及每个满足
(19.5.6.3, (i)) 条件的同态
$v_{\alpha\beta}:F_{\beta,n}\to E_{\alpha,n}$，都存在一个指标
$\gamma\geq\beta$，使得对于每个指标 $\mu\geq\gamma$，都有一个由 $A$-同态组成的
交换图表
\[
\xymatrix@C=4em@R=3em{
  F_{\beta,n}\ar[r]^{v_{\alpha\beta}} & E_{\alpha,n} \\
  & E_{\mu,n}\ar[ul]_{u_{\beta\mu}}\ar[u]_{p_{\alpha\mu}}
}
\]
其中 $p_{\alpha\mu}$ 是转移同态。
\end{lemma}

\begin{proof}
把判别法 (19.1.3) 应用于每个 $\varphi_j$（$0\leq j\leq n$），可见存在一个
指标 $\gamma\geq\beta$ 以及唯一确定的（并且是满的）$C_\gamma$-模同态
\[
  w_{\beta\gamma}^j:\gr^j(E_{\gamma,n})\longrightarrow
  \gr^j(F_{\beta,n})
\]
使诸复合
\[
  \gr^j(F_{\gamma,n})
  \xrightarrow{\varphi_{\gamma\gamma,j}}\gr^j(E_{\gamma,n})
  \xrightarrow{w_{\beta\gamma}^j}\gr^j(F_{\beta,n})
\]
为转移同态（可以选择同一个指标 $\gamma$ 供所有 $\varphi_j$ 使用，是因为它们只有
有限多个）。此外，诸 $w_{\beta\gamma}^j$ 的唯一性证明（因为 $\varphi$ 是一个
分次代数同态）：
$w_{\beta\gamma}=(w_{\beta\gamma}^j)_{0\leq j\leq n}$ 是一个分次
$C_\gamma$-代数同态
\[
  \gr^*(E_{\gamma,n})\longrightarrow
  \gr^*(F_{\beta,n})=F_{\beta,n}.
\]
此外，因为 $\varphi_{\gamma\gamma,0}$ 和 $\varphi_{\gamma\gamma,1}$ 是恒等同态，
$w_{\beta\gamma}^0:C_\gamma\to C_\beta$ 和
$w_{\beta\gamma}^1:\mathfrak J_\gamma/\mathfrak J_\gamma^2
\to\mathfrak J_\beta/\mathfrak J_\beta^2$ 是转移同态，因而
$\gr^0(v_{\alpha\beta}\circ w_{\beta\gamma})$ 和
$\gr^1(v_{\alpha\beta}\circ w_{\beta\gamma})$ 也是转移同态；由于
$\gr^*(E_{\gamma,n})$ 由 $\gr^1(E_{\gamma,n})$ 生成，于是可知
$\gr^j(v_{\alpha\beta}\circ w_{\beta\gamma})$ 对 $0\leq j\leq n$ 也是转移同态。

现在把引理 (19.5.6.3, (ii)) 应用于 $\alpha$、$\beta$ 和 $\gamma$，便得到
(19.5.6.4) 的图表，其中 $\delta\geq\gamma$；然后把 $\alpha$ 和 $\beta$ 分别
替换为 $\gamma$ 和 $\delta$，重复开头的论证。由此得到一个指标
$\lambda\geq\delta$ 以及同态的交换图表
\[
\xymatrix@C=4em@R=4em{
  \gr^*(E_{\gamma,n})\ar[r]^{w_{\beta\gamma}} &
  F_{\beta,n}\ar[r]^{v_{\alpha\beta}} & E_{\alpha,n} \\
  \gr^*(E_{\lambda,n})\ar[u]_{\gr^*(p_{\gamma\lambda})}
    \ar[r]_{w_{\delta\lambda}} &
  F_{\delta,n}\ar[u]_{q_{\beta\delta}}\ar[r]_{v_{\gamma\delta}} &
  E_{\gamma,n}\ar[u]_{p_{\alpha\gamma}}
    \ar@{-->}[ul]_{u_{\beta\gamma}}
}
\]
其中竖直箭头是转移同态。问题归结为证明存在使该图表交换的同态
$u_{\beta\gamma}$；为此，显然只须证明
\[
  \operatorname{Ker}(v_{\gamma\delta})
  \subset\operatorname{Ker}(q_{\beta\delta}),
\]
因为 $v_{\gamma\delta}$ 和 $q_{\beta\delta}$ 都是满的。由于
$w_{\delta\lambda}$ 是满的，最后这个关系等价于
\[
\begin{split}
  \operatorname{Ker}(v_{\gamma\delta}\circ w_{\delta\lambda})
  &\subset\operatorname{Ker}(q_{\beta\delta}\circ w_{\delta\lambda})\\
  &=\operatorname{Ker}\bigl(w_{\beta\gamma}\circ
  (\gr(p_{\gamma\lambda}))\bigr).
\end{split}
\]
但上面已经看到
$\gr(p_{\gamma\lambda})=\gr(v_{\gamma\delta}\circ w_{\delta\lambda})$，而且显然
\[
\begin{split}
  \operatorname{Ker}(v_{\gamma\delta}\circ w_{\delta\lambda})
  &\subset\operatorname{Ker}\bigl(\gr(v_{\gamma\delta}\circ
  w_{\delta\lambda})\bigr)\\
  &=\operatorname{Ker}(\gr(p_{\gamma\lambda}))
  \subset\operatorname{Ker}\bigl(w_{\beta\gamma}\circ
  (\gr(p_{\gamma\lambda}))\bigr),
\end{split}
\]
这就完成了引理的证明，因为对于每个 $\mu\geq\gamma$，只须把
$u_{\beta\mu}$ 定义为复合
\[
  E_{\mu,n}\xrightarrow{p_{\gamma\mu}}E_{\gamma,n}
  \xrightarrow{u_{\beta\gamma}}F_{\beta,n}.
\]
\end{proof}

\begin{env}[19.5.6.8]
\label{0.19.5.6.8-fr}
\emph{(19.5.3, (iii)) 的证明}。--- 设 $G$ 是一个离散拓扑 $A$-代数，
$\mathfrak R$ 是 $G$ 中一个平方为零的理想，而
$f:B\to G/\mathfrak R$ 是一个连续的 $A$-代数同态。由于
$G/\mathfrak R$ 是离散的，存在一个指标 $\alpha$，使得 $f$ 在
$\mathfrak b_\alpha$ 上为零；依假设，存在一个整数 $m$，使得
$\mathfrak J^m\subset\mathfrak b_\alpha$，因而 $f$ 分解为
\[
  B\xrightarrow{p_\alpha}E_{\alpha,n}
  \xrightarrow{f_{\alpha,n}}G/\mathfrak R,
\]
其中取 $n=2m$。由于引理 (19.5.6.3) 的假设得到满足，首先存在
$\beta\geq\alpha$，并有复合 $A$-同态
\begin{equation}
\label{0.19.5.6.9-fr}
  f_{\alpha,n}\circ v_{\alpha\beta}:F_{\beta,n}
  \xrightarrow{v_{\alpha\beta}}E_{\alpha,n}
  \xrightarrow{f_{\alpha,n}}G/\mathfrak R.
  \tag{19.5.6.9}
\end{equation}
因为 $F_{\beta,n}$ 是一个 $C_\beta$-代数，从而当然也是一个（离散的）
拓扑 $C$-代数，故通过复合得到一个连续的 $A$-同态
$r:C\to F_{\beta,n}\to G/\mathfrak R$。另一方面，因为 $C$ 是一个
形式平滑的 $A$-代数，

\oldpage[0\textsubscript{IV-1}]{99}

这个同态分解为
$r:C\xrightarrow{r'}G\to G/\mathfrak R$，其中 $r'$ 是一个连续的
$A$-同态；这样便赋予 $G$ 一个（离散的）拓扑 $C$-代数结构。另一方面，
与典范同态
\[
  \mathfrak J/\mathfrak J^2\longrightarrow
  \mathfrak J_\beta/\mathfrak J_\beta^2\longrightarrow F_{\beta,n},
\]
复合后，由 (19.5.6.9) 得到一个连续的 $C$-同态
$g:\mathfrak J/\mathfrak J^2\to G/\mathfrak R$。因为
$\mathfrak J/\mathfrak J^2$ 是一个形式投射的 $C$-模，$g$ 分解为
\[
  \mathfrak J/\mathfrak J^2\xrightarrow{h}G\longrightarrow G/\mathfrak R,
\]
其中 $h$ 是一个连续的 $C$-同态。由于 $G$ 是离散的，存在一个指标
$\gamma\geq\beta$，使得 $h$ 在
$((\mathfrak b_\gamma\cap\mathfrak J)+\mathfrak J^2)/\mathfrak J^2$ 上为零，
而 $r'$ 在 $(\mathfrak b_\gamma+\mathfrak J)/\mathfrak J$ 上也为零；因而
$h$ 分解为
\[
  \mathfrak J/\mathfrak J^2\longrightarrow
  \mathfrak J_\gamma/\mathfrak J_\gamma^2
  \xrightarrow{h'}G,
\]
其中 $h'$ 是一个 $C_\gamma$-同态。通过延拓，遂由 $h'$ 得到一个
$C_\gamma$-代数同态
\[
  w:\mathbf S_{C_\gamma}^*(\mathfrak J_\gamma/\mathfrak J_\gamma^2)
  \longrightarrow G
\]
并且按构造，每个 $m$ 次齐次元素在 $w$ 下的像都属于 $\mathfrak R$；所以每个
$n=2m$ 次齐次元素的像都为零，因为 $\mathfrak R^2=0$。换言之，$w$ 分解为
\[
  \mathbf S_{C_\gamma}^*(\mathfrak J_\gamma/\mathfrak J_\gamma^2)
  \longrightarrow F_{\gamma,n}\xrightarrow{w'_\gamma}G.
\]
现在把引理 (19.5.6.7) 应用于
\[
  v_{\alpha\gamma}:F_{\gamma,n}
  \xrightarrow{q_{\beta\gamma}}F_{\beta,n}
  \xrightarrow{v_{\alpha\beta}}E_{\alpha,n}
\]
其假设已经得到满足；于是存在 $\delta\geq\gamma$ 以及一个同态
$u_{\gamma\delta}:E_{\delta,n}\to F_{\gamma,n}$，使得复合
$v_{\alpha\gamma}\circ u_{\gamma\delta}$ 是转移同态
$p_{\alpha\delta}:E_{\delta,n}\to E_{\alpha,n}$。最终得到由连续
$A$-同态组成的交换图表
\[
\xymatrix@C=3em@R=2.5em{
  && E_{\alpha,n}\ar[r]^{f_{\alpha,n}} & G/\mathfrak R \\
  && F_{\beta,n}\ar[u]_{v_{\alpha\beta}} & \\
  B\ar[r]_{p_\delta} & E_{\delta,n}\ar[uur]^{p_{\alpha\delta}}
    \ar[r]_{u_{\gamma\delta}} &
  F_{\gamma,n}\ar[u]_{q_{\beta\gamma}}\ar[r]^{w'_\gamma} & G\ar[uu]
}
\]
因此，下方一行诸同态的复合 $f':B\to G$ 使 $f$ 分解为
$B\xrightarrow{f'}G\to G/\mathfrak R$，这就证明了 $B$ 是一个形式平滑的
$A$-代数。

定理 (19.5.3) 的证明至此完成。
\end{env}

\begin{corollary}[19.5.7]
\label{0.19.5.7-fr}
设 $A$ 是一个拓扑环，$B$ 是一个拓扑 $A$-代数，
$(\mathfrak b_\lambda)$ 是 $B$ 中开理想的一个基本系，$\mathfrak J$ 是 $B$ 的一个
理想，而 $C=B/\mathfrak J$ 是商拓扑 $A$-代数。置
$C_\lambda=B/(\mathfrak b_\lambda+\mathfrak J)$。假设：$1^\circ$ 对每个 $n$，
$\mathfrak J^n$ 上由 $B$ 的拓扑诱导出的拓扑，也就是 $B$-模
$\mathfrak J^n$ 上由 $B$ 的拓扑所导出的拓扑 (19.0.2)（特别地，若 $B$ 是
Noether 环且其
拓扑是准进拓扑，则这个条件得到满足 ($\mathbf 0_{\mathrm I}$, 7.3.2)）；
$2^\circ$ $C$ 是一个形式平滑的 $A$-代数。在这些条件下：
\begin{enumerate}
  \item[\rm(i)] 若 $B$ 是一个形式平滑的 $A$-代数，则对每个 $\lambda$，
  $(\mathfrak J/\mathfrak J^2)\otimes_C C_\lambda$ 是一个投射
  $C_\lambda$-模。
  \item[\rm(ii)] 若 $B$ 是一个形式平滑的 $A$-代数并且是一个预可容环，则对每个
  $\lambda$，典范同态
  \begin{equation}
  \label{0.19.5.7.1-fr}
    \varphi_\lambda=\varphi\otimes 1_{C_\lambda}:
    \mathbf S_{C_\lambda}^*((\mathfrak J/\mathfrak J^2)\otimes_C C_\lambda)
    \longrightarrow\gr_{\mathfrak J}^*(B)\otimes_C C_\lambda
    \tag{19.5.7.1}
  \end{equation}
  是双射。

  \oldpage[0\textsubscript{IV-1}]{100}

  \item[\rm(iii)] 反过来，假设 $B$ 是预可容的，序列 $(\mathfrak J^n)$ 趋于
  $0$，并且对每个 $\lambda$，
  $(\mathfrak J/\mathfrak J^2)\otimes_C C_\lambda$ 是一个投射
  $C_\lambda$-模，而同态 (19.5.4.1) 是双射。则 $B$ 是一个形式平滑的
  $A$-代数。
\end{enumerate}
\end{corollary}

\begin{proof}
事实上，$\mathfrak J/\mathfrak J^2$ 的拓扑以及诸
$\mathbf S_C^n(\mathfrak J/\mathfrak J^2)$ 的拓扑，这时也都是由 $B$ 的拓扑
导出的 (19.5.1)；各结论于是立即由 (19.5.3) 以及 (19.1.4) 和 (19.2.4) 中
专用于这种拓扑的判别准则得到。
\end{proof}

\begin{remark}[19.5.8]
\label{0.19.5.8-fr}
假设 $\mathfrak J/\mathfrak J^2$ 是一个有限型 $C$-模，并且对于 $B$ 的拓扑所
诱导的商拓扑，$C$ 是一个 Zariski 环；设 $\mathfrak r$ 是 $C$ 的一个定义理想，
从而诸 $\mathfrak r^n$ 构成 $0$ 在 $C$ 中的一个邻域基本系。这时，(19.5.7) 中
(i) 和 (ii) 的结论可以分别替换为：
\begin{enumerate}
  \item[\rm(i')] $\mathfrak J/\mathfrak J^2$ 是一个投射 $C$-模。
  \item[\rm(ii')] 典范同态
  $\varphi:\mathbf S_C^*(\mathfrak J/\mathfrak J^2)
  \to\gr_{\mathfrak J}^*(B)$ 是双射。
\end{enumerate}
事实上，(i') 蕴涵 (19.5.7) 中 (i) 的结论，这是显然的。反过来，若
$(\mathfrak J/\mathfrak J^2)\otimes_C C_\lambda$ 是一个投射
$C_\lambda$-模，而这对每个 $\lambda$ 都成立，则
$(\mathfrak J/\mathfrak J^2)\otimes_C(C/\mathfrak r^n)$ 是一个投射（从而
平坦的）$(C/\mathfrak r^n)$-模，而这对每个 $n$ 都成立；由此可知
$\mathfrak J/\mathfrak J^2$ 是一个
平坦 $C$-模 ($\mathbf 0_{\mathrm{III}}$, 10.2.2)，而且因为它是有限表示的，
所以它是投射的（Bourbaki，\emph{Alg. comm.}，chap.~II，\textsection5，
n\textsuperscript{o}~2，th.~1 的 cor.~2）。另一方面，$C$-模
$\mathbf S_C^n(\mathfrak J/\mathfrak J^2)$ 和 $\gr_{\mathfrak J}^n(B)$ 都是
有限型的；已知当 $C$ 是一个 Zariski 环时，说 $\varphi_n$ 是双射与说
$\widehat\varphi_n$ 是双射是等价的（Bourbaki，\emph{Alg. comm.}，chap.~III，
\textsection3，n\textsuperscript{o}~5，prop.~9），因此 (ii) 与 (ii') 等价。
\end{remark}

\subsection{域上代数的情形}
\label{subsection:0.19.6-fr}

\begin{theorem}[19.6.1 (Cohen)]
\label{0.19.6.1-fr}
设 $k$ 是一个域，$K$ 是 $k$ 的一个扩域，并且 $k$ 和 $K$ 都赋予离散拓扑。
$K$ 是一个形式平滑的 $k$-代数，当且仅当 $K$ 是 $k$ 的一个可分扩域。
\end{theorem}
\SourceCorrectionNote{法文印本及外交式转录的定理标题括号失衡，作 “19.6.1) (Cohen”；应为平衡标题“19.6.1 (Cohen)”。}。

\begin{proof}
条件的必要性将在 (19.6.5.1) 中证明（而且在此之前当然不会使用）；这里我们只证明
这个条件是\emph{充分的}。分两种情形：

\emph{I. --- $K$ 是 $k$ 的一个有限型可分扩域。} 这时已知（Bourbaki，
\emph{Alg.}，chap.~V，\textsection9，n\textsuperscript{o}~3，th.~2），
存在一个纯超越子扩域 $K'=k(T_1,\ldots,T_n)$（属于 $K$），使得 $K$ 是
$K'$ 的一个有限可分代数扩域。考虑到 (19.3.5, (iii))，因此只须分别处理
$K=K'$ 的情形，以及 $K$ 是 $k$ 的有限代数扩域的情形。在第一种情形中，
已知 $A=k[T_1,\ldots,T_n]$ 是一个形式平滑的 $k$-代数 (19.3.3)，所以
$k(T_1,\ldots,T_n)=S^{-1}A$ 也是如此，其中 $S=A-\{0\}$
(19.3.5, (iv))。在第二种情形中，\emph{所有} Hochschild 上同调群
$\mathrm H_k^n(K,L)$，对于任意 $(K,K)$-双模 $L$，都为零：事实上，若考察 $k$-代数
张量积 $C=K\otimes_k K$，则 $L$ 是一个左 $C$-模，并且上同调群
$\mathrm H_k^n(K,L)$ 等于 $\operatorname{Ext}_C^n(K,L)$，这里 $K$ 也视为
一个 $(K,K)$-双模，

\oldpage[0\textsubscript{IV-1}]{101}

因而视为一个左 $C$-模（M, IX, 4）。但因为 $K$ 是 $k$ 的一个有限可分扩域，
已知 $K\otimes_k K$ 是若干 $k$ 的扩域的直积，其中一个就是 $K$ 本身
（Bourbaki，\emph{Alg.}，chap.~VIII，\textsection8，prop.~3）；所以 $K$ 是
一个\emph{投射} $C$-模，由此得到我们的断言。于是，每个 $k$-扩张（它是
$K$ 的扩张且核为 $L$）都是 $k$-平凡的，更不必说交换的 $k$-扩张也是如此；
因而定理在这种情形成立 (19.4.4)。

\emph{II. --- 一般情形。} --- 沿用 (18.4.5) 的记号，需要证明
$\operatorname{Hom}_K(\mathrm H_2(P'_\bullet),L)=0$ 对\emph{每个}
$K$-向量空间 $L$ 都成立，这显然就是说 $\mathrm H_2(P'_\bullet)=0$。若
$K$ 是一族有向子扩域 $K^{(\alpha)}$ 的并，其中每个都是 $k$ 的扩域，则
$P'_\bullet$ 是相应复形 $P'^{(\alpha)}_\bullet$ 的\emph{归纳极限}，因为函子
$\varinjlim$ 在 $k$-模范畴中与张量积可交换；由这个范畴中函子 $\varinjlim$ 的
正合性，遂有
$\mathrm H_2(P'_\bullet)=\varinjlim\mathrm H_2(P'^{(\alpha)}_\bullet)$。
由于 $K$ 可分这一假设蕴涵 $K$ 的每个子扩域也可分，并且 $K$ 是其有限型
子扩域的并，证明的第一部分确实推出 $K$ 是一个形式平滑的 $k$-代数。
证毕。
\end{proof}

\begin{corollary}[19.6.2]
\label{0.19.6.2-fr}
设 $A$ 是一个分离完备局部环，$K$ 是它的剩余域，$k$ 是 $K$ 的一个子域，
并且 $K$ 是 $k$ 的一个可分扩域。则存在一个子域 $K'$（属于 $A$），它包含
$k$，使限制在 $K'$ 上的典范同态 $A\to K$ 是 $K'$ 到 $K$ 的一个同构。
\end{corollary}

\begin{proof}
设 $\mathfrak m$ 是 $A$ 的极大理想。由假设及 (19.6.1)，$K$ 是一个形式平滑的
$k$-代数；在 (19.3.10) 中把 $C$ 替换为 $A$、把 $\mathfrak J$ 替换为
$\mathfrak m$，可见 $K=A/\mathfrak m$ 的恒等自同构分解为
$K\xrightarrow{u}A\to A/\mathfrak m$，其中 $u$ 是一个 $k$-同态，因而必为
一个 $k$-同构，从 $K$ 到一个子域 $K'$，而后者包含 $k$。
\end{proof}

\begin{corollary}[19.6.3]
\label{0.19.6.3-fr}
设 $A$ 是一个完备 Noether 局部环，$\mathfrak m$ 是它的极大理想，$k$ 是它的
剩余域。下列条件等价：
\begin{enumerate}
  \item[a)] $A$ 包含一个子域。
  \item[b)] 若 $p$ 是 $k$ 的特征，则 $pA=0$。
  \item[c)] 存在一个域 $K$，使 $A$ 同构于形式幂级数环
  $B=K[[T_1,\ldots,T_n]]$ 的一个商环。
\end{enumerate}
当这些条件成立时，可以假设 $A$ 同构于 $B/\mathfrak b$，其中
$\mathfrak b$ 包含在 $B$ 的极大理想的平方中。
\end{corollary}

\begin{proof}
$c)$ 蕴涵 $a)$ 是立即的，因为 $A$ 是 $\neq0$ 的，这是由于它是局部环；又因
$K$ 是一个域，并且复合同态
$K\xrightarrow{\varphi}B\to A$（其中 $\varphi$ 是典范嵌入）不为零，所以它是
单射。为了看出 $a)$ 蕴涵 $b)$，只须注意，若 $K$ 是 $A$ 的一个子域，则 $K$
同构于 $k$ 的一个子域，因而有同一个特征 $p$；因为 $p.1=0$ 在 $K$ 中成立，
所以也在 $A$ 中成立，故 $pA=0$。反过来，$b)$ 蕴涵 $a)$，因为诸典范同态的复合
$\mathbf Z\xrightarrow{j}A\to k$ 的核是 $p\mathbf Z$，而且有
$j(p\mathbf Z)=0$，故 $\operatorname{Ker}(j)=p\mathbf Z$；于是 $A$ 包含一个环
$A'$，它同构于 $\mathbf Z/p\mathbf Z$ 且与 $\mathfrak m$ 不相交；若 $p>0$，
$A'$ 是一个域；否则，

\oldpage[0\textsubscript{IV-1}]{102}

$A'$ 的每个非零元素都在 $A$ 中可逆，所以 $A'$ 的分式域（同构于 $\mathbf Q$）
也包含在 $A$ 中。
\SourceCorrectionNote{法文印本及外交式转录误作“每个元素都可逆”，其中零元当然 不可逆；应为“每个非零元素”。}。

最后证明 $a)$ 蕴涵 $c)$；条件 $a)$ 蕴涵 $A$ 包含一个素子域 $k_0$，它同构于
$k$ 的素子域，因而 $k$ 是它的一个可分扩域。应用 (19.6.2)，首先可知存在
一个同构 $f$，从 $k$ 到一个子域 $K$（属于 $A$）。另一方面，设
$(x_i)_{1\leq i\leq n}$ 是 $\mathfrak m$ 中的一组元素，使得模
$\mathfrak m^2$ 取的诸 $x_i$ 的类构成一个基（在 $k$ 上），即
$\mathfrak m/\mathfrak m^2$ 的一个基。
由于 $A$ 是完备的，遂存在一个连续环同态
$u:k[[T_1,\ldots,T_n]]\to A$，使 $u$ 等于 $f$（在 $k$ 上），并且
$u(T_i)=x_i$ 对每个 $i$ 都成立；由于诸 $x_i$ 的选择，这个同态是满的（Bourbaki，
\emph{Alg. comm.}，chap.~III，\textsection2，n\textsuperscript{o}~9，
prop.~11）。
\end{proof}

\begin{theorem}[19.6.4]
\label{0.19.6.4-fr}
设 $k$ 是一个域，$A$ 是一个 Noether 局部 $k$-代数，$\mathfrak m$ 是它的
极大理想，$K=A/\mathfrak m$ 是它的剩余域，并赋予 $A$ 以
$\mathfrak m$-准进拓扑。假设 $K$ 是 $k$ 的一个可分扩域。则下列条件等价：
\begin{enumerate}
  \item[a)] $A$ 是一个形式平滑的 $k$-代数。
  \item[b)] 完备化 $\widehat A$（即 $A$ 的完备化）作为 $k$-代数同构于
  形式幂级数环 $K[[T_1,\ldots,T_n]]$（其中 $k$-代数结构由复合同态
  \[
    k\xrightarrow{\varphi}K\xrightarrow{\psi}K[[T_1,\ldots,T_n]],
  \]
  定义；这里 $\psi$ 是典范嵌入，$\varphi$ 是定义 $K$ 的扩域结构的同态）。
  \item[b')] $\widehat A$ 作为环同构于形式幂级数环
  $K[[T_1,\ldots,T_n]]$。
  \item[c)] $A$ 是一个正则局部环。
\end{enumerate}
\end{theorem}

\begin{proof}
$a)$ 蕴涵 $b)$ 是 (19.5.4)（$a)$ 与 $c)$ 的等价性）的特殊情形；应用时把
$A$ 替换为 $k$、把 $B$ 替换为 $A$、把 $\mathfrak J$ 替换为 $\mathfrak m$，
并考虑 (19.6.1)。显然 $b)$ 蕴涵 $b')$，而由 (17.1.1)，$b')$ 蕴涵 $c)$。
最后，若 $c)$ 成立，则由 (17.1.1, $a)$) 以及 (19.5.4)（$b)$ 与 $a)$ 的
等价性）可知：应用时以 $K=A/\mathfrak m$ 替换 $C$，$A$ 是一个形式平滑的
$k$-代数。
\end{proof}

\begin{corollary}[19.6.5]
\label{0.19.6.5-fr}
设 $k$ 是一个域，$A$ 是一个 Noether 局部 $k$-代数，$\mathfrak m$ 是它的
极大理想，并赋予 $A$ 以 $\mathfrak m$-准进拓扑。假设 $A$ 是一个形式平滑的
$k$-代数；则 $A$ 在 $k$ 上\emph{几何正则}，换言之（参见 $\mathbf{IV}$,
6.7.6），对于每个有限扩域 $k'$（它是 $k$ 的扩域），半局部环
$A'=A\otimes_k k'$ 是正则的 (17.3.6)。此外，若 $K$ 是 $A$ 的剩余域，则
$\widehat A$ 同构于一个形式幂级数环 $K[[T_1,\ldots,T_n]]$。
\end{corollary}

\begin{proof}
由于 $\widehat A'=\widehat A\otimes_k k'$，由 (19.3.6) 和 (17.1.5)
（应用于 $A'$ 在其极大理想处的局部环），可以归结到 $A$ 完备的情形；这时
$A'$ 是一个形式平滑的 $k'$-代数 (19.3.5, (iii))，并且又是若干局部
$k'$-代数的直积，而这些局部代数也都是形式平滑的 (19.3.5, (v))。因此归结为
证明 $A$ 是\emph{正则的}。设 $k_0$ 是 $k$ 的素子域；因为 $k$ 是 $k_0$ 的
一个可分扩域，所以它是一个形式平滑的 $k_0$-代数 (19.6.1)，并且由假设，
$A$ 也是如此 (19.3.5, (ii))。由于剩余域 $K$（即 $A$ 的剩余域）是 $k_0$ 的一个可分
扩域，于是可以把 (19.6.4) 应用于 $A$ 和 $k_0$，从而得到结论。
\end{proof}

\oldpage[0\textsubscript{IV-1}]{103}

\begin{corollary}[19.6.5.1]
\label{0.19.6.5.1-fr}
设 $A$ 是一个 Noether 局部环，$k$ 是 $A$ 的一个子域，并且 $A$（关于其
准进拓扑）是一个形式平滑的 $k$-代数。则每个域 $K$，若满足
$k\subset K\subset A$，都是 $k$ 的一个可分扩域。
\end{corollary}

\begin{proof}
事实上，对于每个有限扩域 $k'$（它是 $k$ 的扩域），环 $K\otimes_k k'$ 可视为
$A'=A\otimes_k k'$ 的一个子环；由于 $A'$ 是正则环，因而是既约的，所以
$K\otimes_k k'$ 也是既约的。这就证明了 $K$ 是 $k$ 的一个可分扩域
（Bourbaki，\emph{Alg.}，chap.~VIII，\textsection7，
n\textsuperscript{o}~3，th.~1）。
\SourceCorrectionNote{法文印本及外交式转录在第二句把应为 $K\otimes_k k'$ 的既约代数误写成 $K$；应采用上述张量积。}。
\end{proof}

请注意，这也证明了 (19.6.1) 的陈述中所列条件是\emph{必要的}。

\begin{remark}[19.6.5.2]
\label{0.19.6.5.2-fr}
若 $K$ 是 $k$ 的一个\emph{不可分}扩域，则形式幂级数环
$K[[T_1,\ldots,T_n]]$ \emph{不是}一个形式平滑的 $k$-代数（这里采用
(19.6.4) 中回顾的通常 $k$-代数结构）。另一方面，确实存在形式平滑的
$k$-代数 $A$，它们是完备 Noether 局部环，而其剩余域 $K$ 是 $k$ 的一个
\emph{不可分}扩域（例如，把 $B=k[T_1,\ldots,T_n]$ 在极大理想
$\mathfrak n$ 处局部化后再完备，其中 $B/\mathfrak n$ 是 $k$ 的一个有限不可分
代数扩域）；因此，这样的环不可能作为 $k$-代数同构于
$K[[T_1,\ldots,T_n]]$，尽管由 (19.6.5)，它作为环同构于后者。

在某些情形下，可以把推论 (19.6.5) 说得更精确：
\end{remark}

\begin{proposition}[19.6.6]
\label{0.19.6.6-fr}
设 $k$ 是一个域，$A$ 是一个 Noether 局部 $k$-代数，$\mathfrak m$ 是它的
极大理想，并赋予 $A$ 以 $\mathfrak m$-准进拓扑。假设
$K=A/\mathfrak m$（即 $A$ 的剩余域）具有如下性质：存在一个有限纯不可分扩域
$k_0$（它是 $k$ 的扩域），
使局部环 $K\otimes_k k_0$ 的剩余域是 $k_0$ 的一个可分扩域（以后我们把这个
性质表述为 $K$ 在 $k$ 上具有\emph{有限纯不可分重数}
（$\mathbf{IV}$，4.7.4）；并且我们将看到，当 $K$ 是 $k$ 的一个\emph{有限型扩域}
时，这个条件总是成立）。则下列条件等价：
\begin{enumerate}
  \item[a)] $A$ 是一个形式平滑的 $k$-代数。
  \item[b)] 存在一个有限纯不可分扩域 $k'$（它是 $k$ 的扩域），使
  $\widehat A\otimes_k k'$ 作为 $k'$-代数同构于形式幂级数代数
  $K'[[T_1,\ldots,T_n]]$，其中 $K'$ 是 $k'$ 的一个可分扩域。
  \item[b')] 存在一个有限扩域 $k'$（它是 $k$ 的扩域），使半局部环
  $\widehat A\otimes_k k'$ 至少有一个局部分量作为 $k'$-代数同构于形式幂级数
  代数 $K'[[T_1,\ldots,T_n]]$，其中 $K'$ 是 $k'$ 的一个可分扩域。
  \item[c)] $A$ 在 $k$ 上几何正则 (19.6.5)。
\end{enumerate}
\end{proposition}

\begin{proof}
首先注意，若 $k'$ 是 $k$ 的一个纯不可分扩域，则 $A'=A\otimes_k k'$ 中只有
一个位于 $\mathfrak m$ 之上的理想；它由那些 $p^h$ 次幂（其中 $p$ 是 $k$ 的
特征指数）属于 $\mathfrak m$ 的元素组成，这里 $h$ 取适当值
（Bourbaki，\emph{Alg. comm.}，chap.~V，\textsection2，
n\textsuperscript{o}~3，lemme~4）。因此 $A'$ 是局部环，
而
$K\otimes_k k'=(A\otimes_k k')/(\mathfrak m\otimes_k k')$
也是局部环；此外，这两个环的剩余域相同。另一方面，回忆：若 $K$ 是 $k$ 的
一个可分扩域，则对每个有限扩域 $k''$（它是 $k$ 的扩域），
$K\otimes_k k''$ 是若干域的
直积（Bourbaki，\emph{Alg.}，chap.~VIII，\textsection7，
n\textsuperscript{o}~3，th.~1 的 cor.~1）；因此
$\mathfrak m\otimes_k k''$ 是 $A\otimes_k k''$ 的根，而
$K\otimes_k k''$ 的诸域分量就是 $A\otimes_k k''$ 在其诸极大理想处的剩余域；

\oldpage[0\textsubscript{IV-1}]{104}

而且这些域都是 $k''$ 的可分扩域，因为对每个扩域 $k_1$（它是 $k''$ 的扩域），
$K\otimes_k k_1=(K\otimes_k k'')\otimes_{k''}k_1$ 没有根
（\emph{loc. cit.}，th.~1）。把这些说明应用于域 $K$（它在 $k$ 上具有有限
纯不可分重数），便可看出：对\emph{每个}有限扩域 $k''$（它是 $k_0$ 的扩域），
$(K\otimes_k k'')_{\mathrm{red}}$ 在 $k''$ 上可分。

现在证明 (19.6.6)。显然 $b)$ 蕴涵 $b')$。下面证明 $b')$ 蕴涵 $a)$。令
$B'=K'[[T_1,\ldots,T_n]]$。由开头的说明，$K'$ 在 $k'$ 上可分这一假设
蕴涵：对每个有限扩域 $k''$（它是 $k'$ 的扩域），完备半局部环
$B'\otimes_{k'}k''$（它等于形式幂级数环
$(K'\otimes_{k'}k'')[[T_1,\ldots,T_n]]$）的诸分量是
$K_i''[[T_1,\ldots,T_n]]$，其中 $K_i''$ 是
$K'\otimes_{k'}k''$ 的诸域分量。由于 $B'\otimes_{k'}k''$ 的诸局部分量也是
\[
  \widehat A\otimes_k k''
  =(\widehat A\otimes_k k')\otimes_{k'}k'',
\]
的局部分量，可见假设 $b')$ 在以它的任意有限扩域替换 $k'$ 时仍然成立。
特别地，可以假设 $k'$ 是\emph{拟伽罗瓦的}，因而它是一个伽罗瓦扩域，其基域
为纯不可分扩域 $k'_0$（它是 $k$ 的扩域）；于是
$A'=\widehat A\otimes_k k'$ 可写成
$A'_0\otimes_{k'_0}k'$，其中
$A'_0=\widehat A\otimes_k k'_0$ 是一个完备局部环。$k'/k'_0$ 的伽罗瓦群
传递地作用于 $A'$ 的诸局部分量
（Bourbaki，\emph{Alg.}，chap.~VIII，\textsection8，prop.~4），
所以这些分量彼此为伽罗瓦共轭。依假设，其中一个分量作为 $k'$-代数同构于
形式幂级数环 $K'[[T_1,\ldots,T_n]]$，这里 $K'$ 是 $k'$ 的一个可分扩域；
由 (19.6.4)，该分量是一个形式平滑的 $k'$-代数。它的诸共轭分量同样在
$k'$ 上形式平滑，因而由 (19.3.5, (v))，它们的乘积 $A'$ 在 $k'$ 上形式平滑。
由于 $k'$ 作为 $k'_0$-模是忠实平坦、投射且有限型的，由 (19.4.6)，
$A'_0$ 在 $k'_0$ 上形式平滑。最后，$k'_0$ 在 $k$ 上也具有相同性质；再次应用
(19.4.6)，可知 $\widehat A$ 是一个形式平滑的 $k$-代数。
\SourceCorrectionNote{法文印本及外交式转录在此处误称 $A'$ 的诸局部分量作为 $k'_0$-代数同构于 $A'_0$，并以 (19.4.7) 作下降。应为上述伽罗瓦共轭、乘积形式平滑以及两次应用 (19.4.6) 的忠实平坦下降论证。}。

我们已经在 (19.6.5) 中看到 $a)$ 蕴涵 $c)$；还须验证 $c)$ 蕴涵 $b)$。
若 $c)$ 成立，则对每个有限纯不可分扩域 $k'$（它是 $k$ 的扩域并包含 $k_0$），
$A\otimes_k k'$ 是一个正则局部环，其剩余域 $K'$ 是 $k'$ 的一个可分扩域；
因此由 (19.6.4)，它的完备化 $\widehat A\otimes_k k'$ 作为 $k'$-代数同构于
形式幂级数环 $K'[[T_1,\ldots,T_n]]$。
\end{proof}

\begin{remarks}[19.6.7]
\label{0.19.6.7-fr}
\begin{enumerate}
  \item[\rm(i)] 请注意，关于 $K$ 在 $k$ 上具有有限纯不可分重数的假设，只在
  $b')\Rightarrow a)$ 的证明中用到。
  \item[\rm(ii)] 我们以后将在 (22.5.8) 中证明，$a)$ 与 $c)$ 等价，而对
  $K$ 是 $k$ 的扩域\emph{不作任何假设}。
\end{enumerate}
\end{remarks}

\subsection{局部同态的情形；存在性与唯一性定理}
\label{subsection:0.19.7-fr}

在本节中，每当把半局部环视为拓扑环时，总是约定其拓扑为
$\mathfrak r$-准进拓扑，其中 $\mathfrak r$ 是它的根。因此，局部环之间的每个
局部同态都自动\emph{连续}。

\begin{theorem}[19.7.1]
\label{0.19.7.1-fr}
设 $A$、$B$ 是两个 Noether 局部环，$\mathfrak m$、$\mathfrak n$ 分别是它们的
极大理想，$k=A/\mathfrak m$ 是 $A$ 的剩余域；分别赋予 $A$ 与 $B$ 以
$\mathfrak m$-准进拓扑和 $\mathfrak n$-准进拓扑。设
$\varphi:A\to B$ 是一个局部同态，并令 $B_0=B\otimes_A k$。下列性质等价：
\begin{enumerate}
  \item[a)] $B$ 是一个\emph{形式平滑的} $A$-代数。

\oldpage[0\textsubscript{IV-1}]{105}

  \item[b)] $B$ 是一个\emph{平坦的} $A$-模，并且
  $B_0=B/\mathfrak mB$（赋予商拓扑）是一个\emph{形式平滑的} $k$-代数。
\end{enumerate}
\end{theorem}

\begin{proof}
证明分若干步进行。

\medskip
\noindent\phantomsection\label{0.19.7.1.1-fr}\textbf{(19.7.1.1)}\enspace
先证明 $b)$ 蕴涵 $a)$；我们将应用判别法 (19.4.7)，取
$\mathfrak J=\mathfrak m$。由 $b)$ 中的第二个假设，只须证明 $B$ 是一个
\emph{形式投射的} $A$-模。该假设蕴涵：对每个 $h>0$，
$B/\mathfrak m^hB$ 是一个平坦的 $(A/\mathfrak m^h)$-模
（$\mathbf 0_{\mathrm{III}}$，10.2.1）。由于各 $\mathfrak m^h$ 构成 $A$ 中
0 的一组基本邻域，并且
$(B/\mathfrak m^hB)\otimes_{A/\mathfrak m^h}k=B_0$，可以把 $A$ 和 $B$
分别替换为 $A/\mathfrak m^h$ 和 $B/\mathfrak m^hB$，从而可以假设
$A$ 是\emph{Artin 的}（因而是离散的）。由于 $B_0$ 是一个形式平滑的
$k$-代数，所以它是正则环 (19.6.5)；设
$(x_i^0)_{1\leq i\leq n}$ 是 $B_0$ 的一个正则参数系 (17.1.6)，并且对每个
$i$，取 $x_i\in B$，使 $x_i^0$ 是它在 $B_0=B/\mathfrak mB$ 中的像。由于
各 $x_i^0$ 生成极大理想 $\mathfrak n_0=\mathfrak n/\mathfrak mB$（属于 $B_0$），
诸理想
\[
  \mathfrak Q_m=\sum_{i=1}^n(x_i^0)^mB_0 \qquad (m>0)
\]
构成 $B_0$ 中 0 的一组基本邻域，因为 $\mathfrak Q_m$ 显然包含于
$\mathfrak n_0^m$，另一方面又包含 $\mathfrak n_0^{mn}$。令
$\mathfrak J_m=\sum_{i=1}^n x_i^mB$（$m>0$）；显然
$\mathfrak n=\mathfrak J_1+\mathfrak mB$。由于存在 $h>0$ 使
$\mathfrak m^h=0$，故有
$\mathfrak n^m\subset\mathfrak J_1^{m-h}\subset\mathfrak n^{m-h}$（$m>h$）；又因为
已经看到 $\mathfrak J_m\supset\mathfrak J_1^{mn}$，可见各 $\mathfrak J_m$
构成 $B$ 中 0 的一组基本邻域。因此，只须证明各 $B/\mathfrak J_m$ 是
\emph{自由的} $A$-模；而这等价于证明它们是\emph{平坦的} $A$-模
（$\mathbf 0_{\mathrm{III}}$，10.1.3）。现在，$(x_i^0)$ 是 $B_0$ 的极大理想中
诸元素所成的一个 $B_0$-正则序列这一假设蕴涵：序列中的诸元素
$(x_i^0)^m$（$1\leq i\leq n$）对每个 $m>0$ 也具有同一性质 (15.1.20)；结论遂由
(15.1.16, $b)$ 和 $c)$ 得出。
\end{proof}

\begin{lemma}[19.7.1.2]
\label{0.19.7.1.2-fr}
设 $A$ 是一个拓扑环，$B$、$C$ 是两个拓扑 $A$-代数，并且都是 Noether 局部环；
以 $\mathfrak m$、$\mathfrak n'$ 分别表示 $B$、$C$ 的极大理想。再假设 $C$ 是完备的，
而 $B/\mathfrak m$（即 $B$ 的剩余域）是一个有限型 $A$-模。令 $E$ 为完备张量积
$B\widehat\otimes_A C$。则：
\begin{enumerate}
  \item[(i)] $E$ 是一个完备的 Noether 半局部环。
  \item[(ii)] 理想 $\mathfrak mE$ 包含于 $E$ 的根中，并且对每个 $h>0$，
  $E/\mathfrak m^hE$ 同构于 $(B/\mathfrak m^h)\otimes_A C$。
  \item[(iii)] 若 $C$ 是一个平坦 $A$-模，则 $E$ 是一个平坦 $B$-模。
\end{enumerate}
\end{lemma}
\SourceCorrectionNote{法文印本及外交式转录在本引理的陈述中直接使用 $\mathfrak m$ 与 $\mathfrak n'$，却未引入它们。 $\mathfrak m$、$\mathfrak n'$ 分别为 $B$、$C$ 的极大理想。}。

\begin{proof}
依定义，$E$ 是张量积 $B\otimes_A C$ 关于由诸理想
$\operatorname{Im}(\mathfrak m^i\otimes_A C)+
\operatorname{Im}(B\otimes_A\mathfrak n'{}^j)$
所定义的拓扑的分离完备化（$\mathbf 0_{\mathrm I}$，7.7.5）。若置
\[
  \mathfrak r=\operatorname{Im}(\mathfrak m\otimes_A C)+
  \operatorname{Im}(B\otimes_A\mathfrak n'),
\]
则
$\mathfrak r^{2i}\subset
\operatorname{Im}(\mathfrak m^i\otimes_A C)+
\operatorname{Im}(B\otimes_A\mathfrak n'{}^i)\subset\mathfrak r^i$，
所以 $E$ 也是 $B\otimes_A C$ 关于 $\mathfrak r$-准进拓扑的分离完备化。
依假设，
$(B\otimes_A C)/\mathfrak r=(B/\mathfrak m)\otimes_A(C/\mathfrak n')$
是一个有限型 $(C/\mathfrak n')$-模，因而是 Artin 环；此外，
$\mathfrak r/\mathfrak r^2$ 是
\[
  ((\mathfrak m/\mathfrak m^2)\otimes_A C)
  \oplus(B\otimes_A(\mathfrak n'/\mathfrak n'{}^2)),
\]
的一个商，故为有限型 $(B\otimes_A C)$-模。应用
（$\mathbf 0_{\mathrm I}$，7.2.11），可知 $E$ 是一个\emph{Noether 环}；又因
$E/\mathfrak rE$ 同构于 $(B\otimes_A C)/\mathfrak r$，后者是 Artin 的，
所以 $E$ 是\emph{半局部的}。

现在注意到，$E$ 同构于
\[
  \varprojlim_{i,j}((B/\mathfrak m^i)\otimes_A(C/\mathfrak n'{}^j)),
\]
故由双重投射极限定理，它也同构于
\[
  \varprojlim_i\left(\varprojlim_j
  ((B/\mathfrak m^i)\otimes_A(C/\mathfrak n'{}^j))\right)~;
\]

\oldpage[0\textsubscript{IV-1}]{106}

但是，$\varprojlim_j((B/\mathfrak m^i)\otimes_A(C/\mathfrak n'{}^j))$
是 $(B/\mathfrak m^i)\otimes_A C$ 的分离完备化；而 $C$ 是完备的，所以这个
分离完备化就是 $(B/\mathfrak m^i)\otimes_A C$ 本身，因为
$B/\mathfrak m^i$ 是有限型 $A$-模——事实上，$\mathfrak m/\mathfrak m^i$
是有限型 $(B/\mathfrak m)$-模（$\mathbf 0_{\mathrm I}$，7.3.6）。因此
\[
  E=\varprojlim_i((B/\mathfrak m^i)\otimes_A C).
\]
对每个整数 $h>0$，$\mathfrak m^hE$ 是 $E$ 的一个理想，因而在 $E$ 中闭
（$\mathbf 0_{\mathrm I}$，7.3.5），所以是完备的；另一方面，它显然在
\[
  \varprojlim_i
  (\operatorname{Im}(\mathfrak m^h/\mathfrak m^{h+i})\otimes_A C),
\]
中稠密，故等于这个投射极限。此外，这里所考虑的所有投射系统均由
\emph{满同态}定义；因此由（$\mathbf 0_{\mathrm{III}}$，13.2.2），
$E/\mathfrak m^hE$ 同构于
\[
  (B/\mathfrak m^h)\otimes_A C
  =(B\otimes_A C)/\operatorname{Im}(\mathfrak m^h\otimes_A C).
\]
特别地，由于
$\operatorname{Im}(\mathfrak m\otimes_A C)\subset\mathfrak r$，可知
$\mathfrak mE$ 包含于 $\mathfrak rE$，从而包含于 $E$ 的根中。最后，若 $C$
是一个平坦 $A$-模，则
$(B/\mathfrak m^h)\otimes_A C=E/\mathfrak m^hE$
是一个平坦 $(B/\mathfrak m^h)$-模，并且这对每个 $h>0$ 都成立；由于 $B$ 和
$E$ 都是 Noether 的，并且
$\mathfrak mE$ 包含于 $E$ 的根中，由（$\mathbf 0_{\mathrm{III}}$，10.2.2）
可知 $E$ 是一个平坦 $B$-模。
\end{proof}

\begin{lemma}[19.7.1.3]
\label{0.19.7.1.3-fr}
设 $A$ 是一个 Noether 局部环，$\mathfrak m$ 是它的极大理想，$k$ 是它的剩余域，
$B_0$ 是一个 $k$-代数；假设 $B_0$ 是一个完备正则的 Noether 局部环。则存在一个
拓扑 $A$-代数 $B$，它是一个完备 Noether 局部环及平坦 $A$-模，并且使 $B_0$
作为 $k$-代数同构于
$B\otimes_A k=B/\mathfrak mB$。
\end{lemma}

\begin{proof}
由于 $\widehat A$ 在 $A$ 上平坦，并且具有相同的剩余域，我们可以归约到 $A$
\emph{完备}的情形。

令 $K$ 为 $B_0$ 的剩余域，并分两种情形：

\medskip
\noindent I) $K$ 是 $k$ 的一个\emph{可分}扩域。由 (19.6.4)，$B_0$ 作为
$k$-代数同构于形式幂级数环 $K[[T_1,\ldots,T_n]]$。当 $B_0=K$ 时，
本引理已经得到证明（$\mathbf 0_{\mathrm{III}}$，10.3.1）；设 $C$ 是一个完备
Noether 局部环，它是平坦 $A$-模，并且 $C\otimes_A k$ 同构于 $K$。当
$n\geq1$ 时，只须取（沿用前述记号）$B=C[[T_1,\ldots,T_n]]$；事实上，
已知（Bourbaki，\emph{Alg. comm.}，chap.~III，§~3，
n\textsuperscript{o}~4，th.~1 的 cor.~3）$B$ 是一个平坦 $C$-模，因而也是一个
平坦 $A$-模；另一方面，显然 $C[[T_1,\ldots,T_n]]\otimes_A k$ 同构于
$(C/\mathfrak mC)[[T_1,\ldots,T_n]]=B_0$。

\medskip
\noindent II) $K$ 的\emph{特征}为 $p>0$，因而 $k$ 也具有相同特征。以 $P$
记素域 $\mathbf F_p$，以 $W(P)$ 记 $p$-进数的完备局部环 $\mathbf Z_p$；它是一个
离散赋值环（因而是\emph{正则的}），其剩余域为 $P$。先证明存在连续环同态
$W(P)\to A$，从而使 $A$ 成为一个拓扑 $W(P)$-代数。事实上，若
$j:\mathbf Z\to A$ 是典范同态，则依假设有
$j(p\mathbf Z)\subset\mathfrak m$，从而
$j^{-1}(\mathfrak m)=p\mathbf Z$；因此 $j$ 可分解为
$\mathbf Z\to\mathbf Z_{p\mathbf Z}\xrightarrow{f}A$，其中 $f$ 是局部同态，
故为连续的；又因 $A$ 完备，它连续延拓为所求同态 $W(P)\to A$。

由于 $k$ 是 $P$ 的可分扩域，情形 I) 表明存在一个局部同态
$W(P)\to W(k)$，其中 $W(k)$ 是一个完备 Noether 局部环及平坦 $W(P)$-模，
并且 $W(k)\otimes_{W(P)}P$ 同构于 $k$。此外，素元 $p$（属于 $W(P)$）因平坦性而是
$W(k)$-正则元（$\mathbf 0_{\mathrm I}$，6.3.4），又因

\oldpage[0\textsubscript{IV-1}]{107}

$W(k)/pW(k)=k$，所以 $pW(k)$ 是 $W(k)$ 的极大理想；这就蕴涵后一个环是一个
\emph{完备离散赋值环}（Bourbaki，\emph{Alg. comm.}，chap.~VI，§~3，
n\textsuperscript{o}~5，prop.~9）。再由 (19.7.1.1) 可知（由于 $k$ 在 $P$ 上
可分，故为形式平滑的 $P$-代数 (19.6.1)），$W(k)$ 是一个\emph{形式平滑的}
$W(P)$-代数。因此，连续 $W(P)$-同态 $W(k)\to k$ 分解为
$W(k)\to A\to k$ (19.3.11)，从而可以把 $A$ 看作一个拓扑 $W(k)$-代数。
现在把情形 I) 应用于 $B_0$（把它看作 $P$-代数）以及 $W(P)$，可知存在一个
$W(P)$-代数 $B_P$，它是一个完备 Noether 局部环及平坦 $W(P)$-模，并且
$B_P\otimes_{W(P)}P$ 作为 $P$-代数同构于 $B_0$。再次利用 $W(k)$ 是一个
形式平滑的 $W(P)$-代数这一事实，由 (19.3.11)，复合同态
$W(k)\to k\to B_0$ 分解为 $W(k)\to B_P\to B_0$；此外，由于
$k=W(k)/pW(k)$，有
$B_P\otimes_{W(k)}k=B_P/pB_P=B_P\otimes_{W(P)}P=B_0$。现在证明 $B_P$
是一个\emph{平坦的} $W(k)$-模；由于 $W(k)$ 是以 $p$ 为素元的离散赋值环，
只须验证 $B_P$ 是一个\emph{无挠的} $W(k)$-模（$\mathbf 0_{\mathrm I}$，6.3.4），
或者等价地，验证 $p$ 是一个 $B_P$-正则元；而这由 $B_P$ 是一个平坦
$W(P)$-模这一事实得出（$\mathbf 0_{\mathrm I}$，6.3.4）。

现在置
\[
  B=B_P\widehat\otimes_{W(k)}A
\]
并注意到，$A$ 的剩余域等于 $W(k)$ 的剩余域，因而\emph{尤其}是一个有限型
$W(k)$-模。首先由 (19.7.1.2)，$B$ 是一个完备 Noether 半局部环，并且
$\mathfrak mB$ 包含于 $B$ 的根中；此外，$B/\mathfrak mB$ 作为 $k$-代数同构于
$B_P\otimes_{W(k)}k=B_0$，所以 $B$ 实际上是一个局部环。由于 $B_P$ 是一个
平坦 $W(k)$-模，(19.7.1.2) 最后表明 $B$ 是一个平坦 $A$-模。证毕。
\end{proof}

\begin{lemma}[19.7.1.4]
\label{0.19.7.1.4-fr}
设 $A$ 是一个环，$\mathfrak J$ 是 $A$ 的一个理想，$M$、$N$ 是两个 $A$-模，
它们关于 $\mathfrak J$-准进拓扑分离。再假设 $N$ 关于 $\mathfrak J$-准进拓扑
完备，并且 $M$ 是一个平坦 $A$-模。设 $u:N\to M$ 是一个 $A$-同态；若
$u\otimes1:N\otimes_A(A/\mathfrak J)\to M\otimes_A(A/\mathfrak J)$ 是双射，
则 $u$ 是双射。
\end{lemma}

\begin{proof}
相伴分次模均取自 $\mathfrak J$-准进滤过；由关于 $M$、$N$ 的
$\mathfrak J$-准进拓扑的假设，只须证明
$\gr(u):\gr_\bullet(N)\to\gr_\bullet(M)$ 是\emph{双射}的
（Bourbaki，\emph{Alg. comm.}，chap.~III，§~2，n\textsuperscript{o}~8，
th.~1 的 cor.~3）。现在有交换图
\[
\xymatrix{
  \gr_0(N)\otimes_{A_0}\gr_\bullet(A)
    \ar[r]^-{\gr_0(u)\otimes1}\ar[d]_{\varphi_N} &
  \gr_0(M)\otimes_{A_0}\gr_\bullet(A)\ar[d]^{\varphi_M} \\
  \gr_\bullet(N)\ar[r]_-{\gr(u)} & \gr_\bullet(M)
}
\]

\oldpage[0\textsubscript{IV-1}]{108}

其中 $A_0=A/\mathfrak J=\gr_0(A)$，而 $\varphi_M$、$\varphi_N$ 是典范映射
（$\mathbf 0_{\mathrm{III}}$，10.1.1.2）。依假设，$\gr_0(u)$ 与 $\varphi_M$
都是双射（$\mathbf 0_{\mathrm{III}}$，10.2.1），并且 $\varphi_N$ 是满射；
由此先推出 $\gr_0(u)\otimes1$ 是双射，再推出 $\varphi_N$ 是单射，因而是双射，
最后推出 $\gr(u)$ 是双射。
\end{proof}

\begin{lemma}[19.7.1.5]
\label{0.19.7.1.5-fr}
设 $A$ 是一个 Noether 环，$\mathfrak J$ 是 $A$ 的一个理想，$B$、$B'$ 是两个
$A$-代数，并且都是 Noether 局部环；同态 $A\to B$、$A\to B'$ 对于
$\mathfrak J$-准进拓扑（在 $A$ 上）是连续的。假设：
\begin{enumerate}
  \item[$1^{\mathrm o}$] $B$ 和 $B'$ 对于 $\mathfrak J$-准进拓扑是完备的；
  \item[$2^{\mathrm o}$] $B$ 是一个形式平滑的 $A$-代数；
  \item[$3^{\mathrm o}$] $B'$ 是一个平坦 $A$-模。
\end{enumerate}
置 $A_0=A/\mathfrak J$，并设
$u_0:B\otimes_A A_0\to B'\otimes_A A_0$ 是一个 $A_0$-同构；则存在一个
$A$-同构 $u:B\to B'$，使 $u_0=u\otimes1$（这又蕴涵 $B'$ 是一个形式平滑的
$A$-代数，而 $B$ 是一个平坦 $A$-模）。
\end{lemma}

\begin{proof}
置 $B_0=B\otimes_A A_0$、$B'_0=B'\otimes_A A_0$。注意，若
$\mathfrak m$、$\mathfrak m'$ 分别是 $B$、$B'$ 的极大理想，则
$\mathfrak J$-准进拓扑（分别在 $B$ 和 $B'$ 上）是分离的，因为
$\mathfrak JB\subset\mathfrak m$ 且 $\mathfrak JB'\subset\mathfrak m'$；
此外，由于 $\mathfrak JB'$ 在 $B'$ 中关于 $\mathfrak J$-准进拓扑是闭的，
复合同态 $B\to B_0\xrightarrow{u_0}B'_0$ 对于 $\mathfrak J$-准进拓扑是连续的，
因而分解为 $B\xrightarrow{u}B'\to B'_0$，其中 $u$ 是一个连续 $A$-同态
(19.3.10)。显然有 $u_0=u\otimes1$；而 $u_0$ 是双射这一假设由 (19.7.1.4)
蕴涵 $u$ 也是双射。
\end{proof}

\medskip
\noindent\textbf{(19.7.1.6)}\enspace
\emph{证明的结束。}——为完成 (19.7.1) 的证明，须证明 $a)$ 蕴涵 $b)$；
已经知道 $a)$ 蕴涵 $B_0$ 是一个形式平滑的 $k$-代数 (19.3.5, (iii))，
所以一切归结为证明 $B$ 是一个平坦 $A$-模。等价地，只须证明 $\widehat B$
是一个平坦 $\widehat A$-模（Bourbaki，\emph{Alg. comm.}，chap.~III，§~5，
n\textsuperscript{o}~4，prop.~4）；而已知 $\widehat B$ 是一个形式平滑的
$\widehat A$-代数 (19.3.6)，故可以归约到 $A$、$B$ 都完备的情形。由于
$B_0$ 是一个形式平滑的 $k$-代数，它是正则环 (19.6.5)，并且是完备的
（$\mathbf 0_{\mathrm I}$，6.3.5）。应用 (19.7.1.3)，可知存在一个 $A$-代数
$B'$，它是一个完备 Noether 局部环及平坦 $A$-模；还存在局部同态 $A\to B'$
以及一个 $k$-同构 $B_0\simeq B'_0=B'\otimes_A k$。于是，只须在 (19.7.1.5)
中取 $\mathfrak J$ 为 $A$ 的极大理想，便可得到 $B$ 作为 $A$-代数同构于 $B'$，
因而是一个平坦 $A$-模。证毕。

\begin{theorem}[19.7.2]
\label{0.19.7.2-fr}
设 $A$ 是一个 Noether 局部环，$\mathfrak J$ 是包含于 $A$ 的极大理想中的一个
理想，$A_0=A/\mathfrak J$，$B_0$ 是一个完备 Noether 局部环；设
$A_0\to B_0$ 是一个局部同态，使 $B_0$ 成为形式平滑的 $A_0$-代数。则存在一个
完备 Noether 局部环 $B$ 以及一个局部同态 $A\to B$，使 $B$ 成为平坦 $A$-模；
并且存在一个 $A_0$-同构
$u:B\otimes_A A_0\xrightarrow{\sim}B_0$。若 $(B',u')$ 是满足与 $(B,u)$
相同条件的一对对象，则存在一个 $A$-同构
$v:B\xrightarrow{\sim}B'$，使下图交换：
\[
\xymatrix{
  B\otimes_A A_0\ar[rr]^{v\otimes1}_{\sim}\ar[dr]^{\sim}_u &&
  B'\otimes_A A_0\ar[dl]^{u'}_{\sim} \\
  & B_0
}
\]
\end{theorem}

\oldpage[0\textsubscript{IV-1}]{109}

\begin{proof}
令 $\mathfrak m$ 为 $A$ 的极大理想，于是
$\mathfrak m_0=\mathfrak m/\mathfrak J$ 是 $A_0$ 的极大理想；$A$ 与 $A_0$
具有相同的剩余域 $k$。置
$B_{00}=B_0\otimes_{A_0}k$；由于 $B_{00}$ 是一个形式平滑的 $k$-代数
(19.3.5, (iii))，它是一个正则局部环 (19.6.5)。应用 (19.7.1.3)，可知存在
一个拓扑 $A$-代数 $B$，它是一个完备 Noether 局部环及平坦 $A$-模，并且存在
一个 $k$-同构
$u_{00}:B\otimes_A k\xrightarrow{\sim}B_{00}$。注意，由 (19.7.1)，$B$ 是一个
形式平滑的 $A$-代数；因此
$B\otimes_A A_0=B/\mathfrak JB$ 是一个形式平滑的 $A_0$-代数
(19.3.5, (iii))，并且是一个完备 Noether 局部环。此外，依假设及 (19.7.1)，
$B_0$ 是一个平坦 $A_0$-模。由于有 $k$-同构
\[
  u_{00}:B\otimes_A k=(B\otimes_A A_0)\otimes_{A_0}k
  \xrightarrow{\sim}B_0\otimes_{A_0}k=B_{00},
\]
在 (19.7.1.5) 中把 $A$ 替换为 $A_0$、把 $\mathfrak J$ 替换为 $\mathfrak m_0$，
便可推出存在一个 $A_0$-同构
$u:B\otimes_A A_0\xrightarrow{\sim}B_0$，使 $u_{00}=u\otimes1$。

对于唯一性断言，注意诸理想 $\mathfrak J^hB$（相应地，$\mathfrak J^hB'$）
在 $B$（相应地，在 $B'$）中是闭的（$\mathbf 0_{\mathrm I}$，7.3.5），
故 $B$、$B'$ 对于 $\mathfrak J$-准进拓扑是分离且完备的
（Bourbaki，\emph{Top. gén.}，chap.~III，3\textsuperscript{e} éd.，§~3，
n\textsuperscript{o}~5，prop.~9 的 cor.~2）。依假设，有一个 $A_0$-同构
$v_0:B\otimes_A A_0\xrightarrow{\sim}B'\otimes_A A_0$，使
$u'\circ v_0=u$；由于 $B$ 是一个形式平滑的 $A$-代数，而 $B'$ 是一个平坦
$A$-模，可以应用 (19.7.1.5)，从而得到所求 $A$-同构 $v$。
\end{proof}

\begin{remarks}[19.7.3]
\label{0.19.7.3-fr}
\begin{enumerate}
  \item[(i)] 值得注意的是，若只假设 $B$ 和 $B'$ 对于
  $\mathfrak J$-准进拓扑\emph{完备}，(19.7.2) 中的唯一性断言仍然成立。
  我们不知道能否对存在性断言作同样的改进；换言之，能否不再假设局部环
  $B_0$ 对于它的 $\mathfrak n_0$-准进拓扑\emph{完备}（这里
  $\mathfrak n_0$ 表示它的极大理想），而只要求 $B$ 对于
  $\mathfrak J$-准进拓扑完备。当 $A_0$ 对于 $\mathfrak m$-准进拓扑完备时，
  可以看出这个问题等价于下述问题：若 $B_0$ 是一个正则 Noether 局部环
  （不必完备），并且包含素域 $\mathbf F_p=\mathbf Z/p\mathbf Z$，则对每个
  $n>1$，是否存在一个平坦的 $(\mathbf Z/p^n\mathbf Z)$-代数 $B$，使
  $B/pB$ 同构于 $B_0$？
  \item[(ii)] 还应注意，(19.7.2) 中断言存在的同构 $v$ 一般而言
  \emph{并不唯一}（参见 (19.8.7)）。
\end{enumerate}
\end{remarks}

\subsection{Cohen 代数与 Cohen $p$-环；完备局部环结构上的应用}
\label{subsection:0.19.8-fr}

本节的结果是 (19.7) 中诸定理的直接应用；不过鉴于其实用上的重要性，仍值得把它们
明确写出。

\begin{definition}[19.8.1]
\label{0.19.8.1-fr}
设 $A$、$B$ 是两个 Noether 局部环，$\mathfrak m$ 是 $A$ 的极大理想，
$k=A/\mathfrak m$ 是它的剩余域，$\varphi:A\to B$ 是一个局部同态，使 $B$
成为一个 $A$-代数。若 $B$ 满足下列条件，则称它为一个\emph{$A$-Cohen 代数}：
\begin{enumerate}
  \item[(i)] $B$ 是\emph{完备环}。
  \item[(ii)] $B$ 是一个\emph{平坦} $A$-模。
  \item[(iii)] $B\otimes_A k$ 是一个\emph{域}（换言之，$\mathfrak mB$ 是 $B$
  的极大理想），并且是 $k$ 的一个\emph{可分扩域}。
\end{enumerate}
\end{definition}

\oldpage[0\textsubscript{IV-1}]{110}

\begin{theorem}[19.8.2]
\label{0.19.8.2-fr}
设 $A$ 是一个 Noether 局部环，$k$ 是它的剩余域。
\begin{enumerate}
  \item[(i)] 若 $B$ 是一个 $A$-Cohen 代数，则 $B$ 是一个形式平滑的
  $A$-代数。对任意完备 Noether 局部环 $C$、任意局部同态 $A\to C$ 以及
  任意理想 $\mathfrak J\neq C$（属于 $C$），每个 $A$-同态 $B\to C/\mathfrak J$
  因而都分解为 $B\xrightarrow{v}C\to C/\mathfrak J$，其中 $v$ 是一个
  $A$-同态（必为局部同态）。
  \item[(ii)] 对于 $k$ 的任一可分扩域 $K$，存在一个 $A$-Cohen 代数 $B$，
  使 $B\otimes_A k$ 作为 $k$-代数同构于 $K$；而且这样的 $A$-代数在同构意义下
  是唯一的。
\end{enumerate}
\end{theorem}

\begin{proof}
由于 $K$ 是一个形式平滑的 $k$-代数 (19.6.1)，断言 (i) 由 (19.7.1) 得出。
为证明 (ii)，可以归约到 $A$ \emph{完备}的情形，因为说 $B$ 是一个平坦
$A$-模，等价于说它是一个平坦 $\widehat A$-模（$\mathbf 0_{\mathrm{III}}$，
10.2.3）；并且有 $\mathfrak mB=\mathfrak m\widehat AB$，而 $k$ 是
$\widehat A$ 的剩余域。于是只须应用 (19.7.2)，取 $\mathfrak J=\mathfrak m$
以及 $B_0=K$（并应用 (19.6.1)）。
\end{proof}

\begin{definition}[19.8.3]
\label{0.19.8.3-fr}
我们把形如 $\mathbf Z_{p\mathbf Z}$ 的局部环称为\emph{素局部环}，其中
$p\mathbf Z$ 是 $\mathbf Z$ 的一个素理想。我们把素局部环的完备化称为
\emph{完备素局部环}。
\end{definition}

因此，素局部环分为两类：
\begin{enumerate}
  \item[$1^{\mathrm o}$] 对应于极大理想 $p\mathbf Z$ 的情形，其中
  $p\neq0$ 是一个素数；$\mathbf Z_{p\mathbf Z}$ 是一个离散赋值环，其完备化
  是\emph{$p$-进整数环}，通常记作 $\mathbf Z_p$\footnote{这一记号目前已普遍使用，
  但在这里与所采用的记号 $A_f$（$\mathbf 0_{\mathrm I}$，1.2.3）冲突：事实上，
  若 $A=\mathbf Z$ 且 $f=p$，则 $A_f$ 表示由形如 $k/p^n$ 的有理数组成的环
  （$k\in\mathbf Z$，$n$ 是一个 $\geq0$ 的整数）；我们始终避免用记号
  $\mathbf Z_p$ 表示后一环。}。
  \item[$2^{\mathrm o}$] 对于素理想 $p\mathbf Z=(0)$，
  $\mathbf Z_{p\mathbf Z}$ 是有理数域 $\mathbf Q$，并等同于它自身的完备化
  （此处的拓扑自然取为 Noether 局部环的拓扑，故即为离散拓扑）。
\end{enumerate}

（19.8.3）的术语与“素域”这一术语类似，并可按同样方式加以说明：对任意局部环
$A$，考虑典范同态 $\varphi:\mathbf Z\to A$，并令 $p\mathbf Z$ 为极大理想
$\mathfrak m$（属于 $A$）在此同态下的逆像；$p\mathbf Z$ 是 $\mathbf Z$ 的一个素理想，故上述
同态分解为
$\mathbf Z\to\mathbf Z_{p\mathbf Z}\xrightarrow{\psi}A$；此外，由于
$\varphi$ 是从 $\mathbf Z$ 到 $A$ 的\emph{唯一}同态，$p$ 和 $\psi$ 也都由此
\emph{唯一}确定。换言之，对任意局部环 $A$，存在一个\emph{唯一同态}
$\psi:P\to A$，其中 $P$ 是一个素局部环；若 $A$ 还分离且完备，则可通过完备化
延拓这一同态，因而存在一个\emph{唯一同态}
$\widehat\psi:\widehat P\to A$，其中 $\widehat P$ 是一个完备素局部环。此外，
取商后，$\psi$ 给出从剩余域 $\mathbf F_p$（若 $p>0$）或 $\mathbf Q$
（若 $p=0$）到剩余域 $k$（属于 $A$）的同态，因此 $p$ 正是 $k$ 的\emph{特征}。

特别地，若取 $A$ 本身为素局部环（分别地，完备素局部环），便可看出这样的环中
只有\emph{一个自同态}，即恒等同态。

\oldpage[0\textsubscript{IV-1}]{111}

\begin{definition}[19.8.4]
\label{0.19.8.4-fr}
设 $A$ 是一个局部环，$P\to A$ 是从一个素局部环 $P$ 到 $A$ 的唯一同态，
$p$ 是 $P$ 与 $A$ 的剩余域的特征。若 $A$ 是一个\emph{Cohen 环}，也就是说它是
一个 Cohen $P$-代数，则（19.8.1）：
\begin{enumerate}
  \item[$1^{\mathrm o}$] $A$ 是 Noether 环且完备。
  \item[$2^{\mathrm o}$] $A$ 是一个平坦 $P$-模（这也等价于说 $A$ 是一个平坦
  $\widehat P$-模（Bourbaki，\emph{Alg. comm.}，chap.~III，§~5，
  n\textsuperscript{o}~4，prop.~4））。
  \item[$3^{\mathrm o}$] $A/pA$ 是一个域（它必为 $P$ 的剩余域的可分扩域，
  因为后者是一个素域）。
\end{enumerate}
\end{definition}

若 $p=0$，这些条件等价于说 $A$ 是一个\emph{特征为 $0$ 的域}。若 $p>0$，
则必有 $pA\neq0$；条件 $3^{\mathrm o}$ 表示 $pA$ 是极大理想
$\mathfrak m$（属于 $A$）；条件 $2^{\mathrm o}$ 表示 $p$ 是 $A$-正则的，因为 $P$ 是一个
离散赋值环（$\mathbf 0_{\mathrm I}$，6.3.4）。因此，$A$ 是一个正则环
（17.1.1, $d)$），其维数为 $1$，从而是一个离散赋值环，并且由 $1^{\mathrm o}$ 可知它完备；
综上有：

\begin{proposition}[19.8.5]
\label{0.19.8.5-fr}
Cohen 环恰为特征为 $0$ 的域，以及满足下列条件的完备离散赋值环：其剩余域的特征
$p>0$，且其极大理想由 $p\cdot1$ 生成（这里 $1$ 是该环的单位元）。
\end{proposition}

注意，在第二种情形中，$p\cdot1\neq0$，因为 $p$ 是 $A$-正则的，于是可以把
$p\cdot1$ 与整数 $p$ 等同起来；典范同态 $\mathbf Z_p\to A$ 是单射，并且我们把
$p\cdot1$ 与元素 $p$（属于 $\mathbf Z_p$）等同起来；在这种情形下，称 $A$ 为一个
\emph{Cohen $p$-环}。

\begin{theorem}[19.8.6]
\label{0.19.8.6-fr}
\textup{（Cohen）} ---
\begin{enumerate}
  \item[(i)] 设 $W$ 是一个 Cohen 环，$C$ 是一个完备 Noether 局部环，
  $\mathfrak J$ 是 $C$ 的一个不同于 $C$ 的理想。则每个局部同态
  $u:W\to C/\mathfrak J$ 都分解为
  $W\xrightarrow{v}C\to C/\mathfrak J$，其中 $v$ 是一个局部同态。
  \item[(ii)] 设 $K$ 是一个域。存在一个 Cohen 环 $W$，其剩余域同构于 $K$。
  若 $W'$ 是另一个 Cohen 环，$K'$ 是它的剩余域，则每个同构
  $u:K\xrightarrow{\sim}K'$ 都由某个同构 $v:W\xrightarrow{\sim}W'$ 取商得到。
\end{enumerate}
\end{theorem}

\begin{proof}
这不过是把（19.8.2）应用于 $A$ 为素局部环的情形。
\end{proof}

\begin{remarks}[19.8.7]
\label{0.19.8.7-fr}
\begin{enumerate}
  \item[(i)] 当 $K$ 的特征为 $0$ 时，（19.8.6）的第 (ii) 部分成为平凡的。
  \item[(ii)] （19.8.6, (i)）中的同态 $v$ \emph{未必由 $u$ 唯一确定}；
  当 $W$ 是一个特征为 $0$ 的域、$\mathfrak J^2=0$ 且 $u$ 是一个同构时，就已经
  可以看出这一点（参见（21.5.5））。同样，（19.8.6, (ii)）中的同构 $v$
  \emph{未必由 $u$ 唯一确定}（参见（21.5.5））。

  不过，当 $K$ 是一个特征为 $p>0$ 的\emph{完美域}时，我们将在（21.5.5）看到，
  （19.8.6, (ii)）中的同构 $v$ 是唯一的。以后还将看到，在这种情形下，$W$ 可与
  环 $W_\infty(K)$，即 $K$ 上的\emph{无限长度 Witt 向量}环等同。
  \item[(iii)] 在（19.8.6, (i)）中，可以利用（19.3.10）和（19.3.12）减弱
  关于 $C$ 的假设。
\end{enumerate}
\end{remarks}

\begin{theorem}[19.8.8]
\label{0.19.8.8-fr}
设 $A$ 是一个完备 Noether 局部环，$k$ 是它的剩余域。

\oldpage[0\textsubscript{IV-1}]{112}

\begin{enumerate}
  \item[(i)] 存在一个 Cohen 环 $W$，使 $A$ 同构于某个形式幂级数环
  $W[[T_1,\ldots,T_n]]$ 的商环（特别地，$A$ 同构于某个完备正则局部环的商环
  （17.3.8））。若 $A$ 包含一个域，则它同构于
  $k[[T_1,\ldots,T_n]]$ 的一个商环。
  \item[(ii)] 再假设 $A$ 是整环。则存在一个子环 $B$（属于 $A$），使得：
  $1^{\mathrm o}$ $B$ 同构于一个以环 $C$ 为系数的形式幂级数环，而该环是一个域
  或一个 Cohen 环（这蕴含 $B$ 是一个完备正则局部环（17.3.8））；
  $2^{\mathrm o}$ $B$ 与 $A$ 有相同的剩余域，并且单射 $B\to A$ 是一个局部同态；
  $3^{\mathrm o}$ $A$ 是一个有限 $B$-代数。
\end{enumerate}
\end{theorem}

\begin{proof}
令 $\mathfrak m$ 为 $A$ 的极大理想。存在一个 Cohen 环 $W$，其剩余域同构于 $k$
（19.8.6, (ii)）；于是有一个局部同态 $W\to A/\mathfrak m$，因而它分解为
$W\xrightarrow{u}A\to A/\mathfrak m$，其中 $u$ 是一个局部同态。对
任一有限元素族 $(x_i)_{1\leq i\leq n}$（取自 $\mathfrak m$），便存在一个局部同态
$v:W[[T_1,\ldots,T_n]]\to A$，它延拓 $u$，并且有
$v(T_i)=x_i$（对每个 $i$）（Bourbaki，\emph{Alg. comm.}，chap.~III，§~4，
n\textsuperscript{o}~5，prop.~6）。当 $A$ 包含一个域时，它包含一个素域 $P$；
$k$ 是它的一个扩域（必为可分扩域），所以 $A$ 包含一个同构于 $k$ 的域
（19.6.2）；可以把 $W$ 替换为 $k$，用以定义上述 $v$。

\medskip
\noindent(i)\enspace
先取 $x_i$ 为 $\mathfrak m$ 的一组生成元。由于 $W$ 与 $A$ 有相同的剩余域，
而 $x_i$ 在分次环 $\gr_\bullet(A)$ 中的类生成 $\gr_\bullet(A)$（作为一个
$k$-代数），故
$\gr(v):\gr_\bullet(W[[T_1,\ldots,T_n]])\to\gr_\bullet(A)$ 是满射；由此可知
$v$ 本身也是满射（Bourbaki，\emph{Alg. comm.}，chap.~III，§~2，
n\textsuperscript{o}~8，th.~1 的 cor.~2）。回想一下，$A$ 包含一个域的情形
已经证明过，这里只是为了完整起见而再述一次（19.6.3）。

\medskip
\noindent(ii)\enspace
若 $A$ 包含一个域，则由上述论证，它包含一个域 $k'$，后者同构于 $k$；这时取
一个参数系 $(y_j)_{1\leq j\leq m}$（属于 $A$），令 $B=k[[T_1,\ldots,T_m]]$，并考虑
局部同态 $w:B\to A$；它在 $k$ 上与一个同构 $k\to k'$ 一致，并满足
$w(T_j)=y_j$（$1\leq j\leq m$）。若 $A$ 不包含域，则唯一同态
$\mathbf Z_{p\mathbf Z}\to A$（19.8.3）必为单射（否则，由于 $A$ 是整环，
它的核将是 $\mathbf Z_{p\mathbf Z}$ 的极大理想，而它的像将同构于一个域）；
此外，依假设必有 $p>0$，因为 $\mathbf Z_{p\mathbf Z}$ 在 $p=0$ 时是一个域。
$p\cdot1$ 这个 $A$ 的元素（与 $p$ 等同）在 $A$ 中不是零因子，并且属于
$\mathfrak m$，故由（16.3.4）和（16.3.7），存在一个元素族
$(z_i)_{1\leq i\leq m-1}$，它与 $p$ 一同构成 $A$ 的一个参数系。
证明开头所取的 Cohen 环 $W$ 此时是一个剩余域为 $k$ 的离散赋值环，并且 $p$
在其中生成极大理想（19.8.5）；证明开头定义的同态 $u:W\to A$ 把 $p$ 映为自身。
\SourceCorrectionNote{法文印本及外交式转录在此处写作“唯一同态 $u:W\to A$”；但（19.8.7, (ii)）已经明确说明这种提升一般并不唯一。删去“唯一”。}。
于是令 $B=W[[T_1,\ldots,T_{m-1}]]$，并考虑局部同态 $w:B\to A$；它与 $u$
在 $W$ 上一致，并满足 $w(T_i)=z_i$（$1\leq i\leq m-1$）。在两种情形中，
若 $\mathfrak n$ 是 $B$ 的极大理想，则显然 $\mathfrak nA$ 是 $A$ 的一个定义理想；
又因 $B/\mathfrak n$ 与 $A/\mathfrak m$ 同构，$A$ 是一个拟有限 $B$-模
（$\mathbf 0_{\mathrm I}$，7.4.4），故为有限型 $B$-模，因为 $B$ 完备而 $A$ 对于
$\mathfrak n$-准进拓扑分离（$\mathbf 0_{\mathrm I}$，7.4.1）。
\SourceCorrectionNote{法文印本及外交式转录在有限性结论中误作有限型 $A$-模； 为有限型 $B$-模。所引完备性论证作用于拟有限 $B$-模，而定理所需结论也正是 $A$ 对 $B$ 有限。}。
另一方面，

\oldpage[0\textsubscript{IV-1}]{113}

在两种情形中都有 $\dim(B)=\dim(A)=m$；在第一种情形中，这由（17.1.4, (iii)）
得出；在第二种情形中，可以直接看出 $p$ 与诸 $T_i$
（$1\leq i\leq m-1$）构成一个 $B$-正则序列并生成 $\mathfrak n$，或者也可利用
这些元素生成 $\mathfrak n$ 以及（16.3.10）给出的
$\dim(B)\geq\dim(A)$。由于 $A$ 和 $B$ 都是整环，最后由（16.3.10）可知 $w$
是单射，证明完毕。
\end{proof}

\begin{corollary}[19.8.9]
\label{0.19.8.9-fr}
设 $A$ 是含有域 $k_0$ 的完备 Noether 局部整环；设 $k$ 为 $A$ 的剩余域，并假定
$k$ 是 $k_0$ 的有限扩域。则在（19.8.8, (ii)）的结论中，可以把
$1^{\mathrm o}$ 和 $2^{\mathrm o}$ 换成如下条件：$B$ 形如
$k_0[[T_1,\ldots,T_n]]$，并且典范单射 $B\to A$ 是局部
$k_0$-同态（这里的 $k_0$-代数结构，是 $B$ 的通常结构）。
\end{corollary}

\begin{proof}
事实上，沿用（19.8.8, (ii)）的证明，这次把
$w:k_0[[T_1,\ldots,T_n]]\to A$ 定义为在 $k_0$ 上为恒等，并把
$T_j$ 映到 $y_j$（$1\leq j\leq n$）。$k$ 在 $k_0$ 上次数有限这一假设仍推出
$A$ 是拟有限 $B$-模，因而由（$\mathbf 0_{\mathrm I}$，7.4.1）它是有限型的；
于是同（19.8.8）一样得出结论。
\end{proof}

\begin{corollary}[19.8.10]
\label{0.19.8.10-fr}
设 $A$ 是一个 Artin 局部环，其极大理想 $\mathfrak m$ 的平方为零。则存在一个
正则 Noether 局部环 $B$，其极大理想为 $\mathfrak n$，并且
$A$ 同构于 $B/\mathfrak n^2$。
\end{corollary}

\begin{proof}
设 $K$ 为 $A/\mathfrak m$，即 $A$ 的剩余域，并设 $n$ 为
$\mathfrak m/\mathfrak m^2=\mathfrak m$ 作为 $K$-向量空间的维数。若 $A$
含有一个域，则由（19.6.3），$A$ 同构于 $B/\mathfrak b$，其中
$B=K[[T_1,\ldots,T_n]]$，而 $\mathfrak b$ 包含在极大理想
$\mathfrak n$ 的平方中；这是 $B$ 的极大理想。但由于
$\operatorname{long}(B/\mathfrak n^2)=n+1=\operatorname{long}(A)$，必有
$\mathfrak b=\mathfrak n^2$。

现在假定 $A$ 不含域；这推出 $K$ 的特征为 $p>0$，并且
$p\cdot1\neq0$ 于 $A$ 中（19.6.3）。因而 $p\cdot1$ 是 $\mathfrak m$ 的一个元素，
于是另有 $n-1$ 个元素 $x_i$（$2\leq i\leq n$）属于 $\mathfrak m$，它们与
$p\cdot1$ 一同构成 $\mathfrak m$ 的一组 $K$-基。设 $W$ 为剩余域同构于 $K$
的一个 Cohen 环；$W$ 是一个以 $p$ 生成极大理想的离散赋值环。在（19.8.8）的
证明中已经看到，存在一个同态 $u:W\to A$，它把 $p$ 映到自身，并且在取商后
给出 $K$ 上的恒等映射。取 $B=W[[T_2,\ldots,T_n]]$，并考虑局部同态
$w:B\to A$；它与 $u$ 在 $W$ 上一致，并满足
$w(T_i)=x_i$（$2\leq i\leq n$）。显然 $w$ 是满射，并且它的核
$\mathfrak b$ 包含在极大理想
$\mathfrak n=pB+BT_2+\cdots+BT_n$ 的平方中；这是 $B$ 的极大理想。由于
$\operatorname{long}(B/\mathfrak n^2)=n+1=\operatorname{long}(A)$，仍有
$\mathfrak b=\mathfrak n^2$。
\end{proof}

\begin{proposition}[19.8.11]
\label{0.19.8.11-fr}
设 $A$ 是 Artin 局部环，$\mathfrak m$ 是它的极大理想，$k$ 是它的剩余域。
$A$ 同构于某个 Cohen 环的商环，当且仅当 $\mathfrak m$ 由 $p\cdot1$ 生成，
其中 $p$ 是 $k$ 的特征。
\end{proposition}

\begin{proof}
条件的必要性由（19.8.5）显然可知。为证其充分性，同（19.8.8）证明开头一样，
注意到存在一个 Cohen 环 $W$，其剩余域同构于 $k$，以及一个局部同态 $u:W\to A$。
此外，考虑复合同态
$\mathbf Z_{p\mathbf Z}\to W\xrightarrow{u}A$（它必为从
$\mathbf Z_{p\mathbf Z}$ 到 $A$ 的唯一同态），可见 $u$ 把元素
$p\cdot1$（在 $W$ 中）映为元素 $p\cdot1$（在 $A$ 中）；由于元素
$p\cdot1$（在 $W$ 中）生成这个环的极大理想，由假设立即推出

\oldpage[0\textsubscript{IV-1}]{114}

$\gr(u):\gr_\bullet(W)\to\gr_\bullet(A)$ 是满射，因而 $u$ 也是满射
（Bourbaki，\emph{Alg. comm.}，第 III 章，§~2，
n\textsuperscript{o}~8，定理 1 的推论 2）。
\end{proof}

\subsection{相对形式平滑代数}
\label{subsection:0.19.9-fr}

\begin{definition}[19.9.1]
\label{0.19.9.1-fr}
设 $\Lambda$ 是拓扑环，$A$ 是拓扑 $\Lambda$-代数，$B$ 是拓扑 $A$-代数。
若 $B$ 作为 $A$-代数满足如下条件，就称它\emph{相对于 $\Lambda$ 形式平滑}：
对任意离散拓扑 $A$-代数 $C$ 及任意幂零理想 $\mathfrak J$（属于 $C$），每个连续
$A$-同态 $u_0:B\to C/\mathfrak J$，只要它可分解为
$B\xrightarrow{u}C\to C/\mathfrak J$，其中 $u$ 是连续 $\Lambda$-同态，就也可
分解为 $B\xrightarrow{v}C\to C/\mathfrak J$，其中 $v$ 是连续 $A$-同态。
\end{definition}

由此定义可知，若 $B$ 是一个形式平滑的 $A$-代数，则 $B$ 也相对于
$\Lambda$ 形式平滑；这里可以给定任意拓扑 $\Lambda$-代数结构（定义在 $A$ 上）
（换言之，任意连续环同态 $\Lambda\to A$）。

\begin{proposition}[19.9.2]
\label{0.19.9.2-fr}
设 $\Lambda$ 是拓扑环，$A$ 是拓扑 $\Lambda$-代数。
\begin{enumerate}
  \item[(i)] $A$ 是一个相对于 $\Lambda$ 形式平滑的 $\Lambda$-代数。
  \item[(ii)] 若 $B$ 是 $A$-代数且相对于 $\Lambda$ 形式平滑，而 $C$ 是
  $B$-代数且相对于 $\Lambda$ 形式平滑，则 $C$ 是 $A$-代数且相对于
  $\Lambda$ 形式平滑。
  \item[(iii)] 设 $B$ 是 $A$-代数且相对于 $\Lambda$ 形式平滑，$A'$ 是一个
  拓扑 $A$-代数，则拓扑 $A'$-代数 $B\otimes_A A'$ 相对于 $\Lambda$ 形式平滑。
  \item[(iv)] 设 $B$ 是拓扑 $A$-代数，$S$（相应地 $T$）是 $A$（相应地 $B$）
  的乘法子集，并且 $S$ 在 $B$ 中的典范像包含于 $T$。若 $B$ 是一个相对于
  $A$-代数且相对于 $\Lambda$ 形式平滑，则 $T^{-1}B$ 是一个 $S^{-1}A$-代数，
  并且相对于 $\Lambda$ 形式平滑。
  \item[(v)] 设 $B_i$（$1\leq i\leq n$）是拓扑 $A$-代数。乘积
  $\prod_{i=1}^n B_i$ 是一个 $A$-代数且相对于 $\Lambda$ 形式平滑，当且仅当
  每个 $B_i$ 都如此。
\end{enumerate}
\end{proposition}

\begin{proof}
断言 (i) 是平凡的；其余各断言的证明与（19.3.5）的证明完全类似，故留给读者。
\end{proof}

\begin{corollary}[19.9.3]
\label{0.19.9.3-fr}
设 $A$ 是拓扑环，$A'$ 和 $B$ 是两个拓扑 $A$-代数。则拓扑 $A$-代数
$B'=B\otimes_A A'$ 是一个 $B$-代数且相对于 $A$ 形式平滑。
\end{corollary}

\begin{proof}
这由（19.9.2, (i) 和 (iii)）得出。
\end{proof}

\begin{proposition}[19.9.4]
\label{0.19.9.4-fr}
设 $\Lambda$ 是拓扑环，$A$ 是拓扑 $\Lambda$-代数，$B$ 是拓扑 $A$-代数。
下列条件等价：
\begin{enumerate}
  \item[a)] $B$ 是一个 $A$-代数且相对于 $\Lambda$ 形式平滑。
  \item[b)] $\widehat B$ 是一个 $A$-代数且相对于 $\Lambda$ 形式平滑。
  \item[c)] $\widehat B$ 是一个 $\widehat A$-代数且相对于 $\Lambda$ 形式平滑。
  \item[d)] $\widehat B$ 是一个 $\widehat A$-代数且相对于
  $\widehat\Lambda$ 形式平滑。
\end{enumerate}
\end{proposition}

\begin{proof}
证明仍留给读者；它与（19.3.6）的证明相同。
\end{proof}

\oldpage[0\textsubscript{IV-1}]{115}

\begin{env}[19.9.5]
\label{0.19.9.5-fr}
同样，把（19.3.8）中的“形式平滑”换成“相对于 $\Lambda$ 形式平滑”后，该命题
仍然成立（证明也相同）。若把（19.3.10）的命题中的“形式平滑”换成“相对于
$\Lambda$ 形式平滑”，则结论应换成如下结论（证明基本不变）：\emph{每个连续
$A$-同态 $u_0:B\to C/\mathfrak J$，只要它可分解为
$B\xrightarrow{u}C\to C/\mathfrak J$，其中 $u$ 是连续 $\Lambda$-同态，就也可
分解为 $B\xrightarrow{v}C\to C/\mathfrak J$，其中 $v$ 是连续 $A$-同态}。
\end{env}

\begin{env}[19.9.6]
\label{0.19.9.6-fr}
把形式平滑性的判别法（19.4.1）和（19.4.2）中的“形式平滑”换成“相对于
$\Lambda$ 形式平滑”后，它们仍然成立，而证明几乎不变。
\end{env}

\begin{proposition}[19.9.7]
\label{0.19.9.7-fr}
设 $\Lambda$ 是拓扑环，$A$ 是拓扑 $\Lambda$-代数，$B$ 是拓扑 $A$-代数。
假定对任意离散 $A$-代数 $C$ 以及任意理想 $\mathfrak J$（属于 $C$），只要
$\mathfrak J^2=0$，每个连续 $A$-同态
$u_0:B\to C/\mathfrak J$，只要它可分解为
$B\xrightarrow{u}C\to C/\mathfrak J$，其中 $u$ 是连续 $\Lambda$-同态，就也可
分解为 $B\xrightarrow{v}C\to C/\mathfrak J$，其中 $v$ 是连续 $A$-同态。
则 $B$ 是一个 $A$-代数且相对于 $\Lambda$ 形式平滑。
\end{proposition}

\begin{proof}
（19.4.3）的证明可直接照搬。
\end{proof}

\begin{proposition}[19.9.8]
\label{0.19.9.8-fr}
设 $\Lambda$ 是拓扑环，$A$ 是拓扑 $\Lambda$-代数，$B$ 是拓扑 $A$-代数。
$B$ 是一个 $A$-代数且相对于 $\Lambda$ 形式平滑，当且仅当对每个离散拓扑
$B$-模 $L$，若它由 $B$ 的某个开理想零化，就有（参见 18.4.2）
\[
  \operatorname{Exalcotop}_{A/\Lambda}(B,L)=0.
\]
\end{proposition}

\begin{proof}
沿用（19.4.4）证明中的记号，只须注意，在这里可以假定扩张 $E_\Lambda$
（它是 $B_\Lambda$ 的一个扩张）是 $\Lambda$-平凡的；证明的其余部分于是保持不变。
\end{proof}

当 $\Lambda$、$A$ 和 $B$ 都是离散环时，判别法（19.9.8）化为
\begin{equation}
\tag{19.9.8.1}
\label{0.19.9.8.1-fr}
  \operatorname{Exalcom}_{A/\Lambda}(B,L)=0
  \quad\text{对任意 $B$-模 $L$}~;
\end{equation}
换言之，\emph{每个交换 $A$-扩张，若它是 $B$ 由一个 $B$-模给出的，并且是
$\Lambda$-平凡的，则也是 $A$-平凡的}。

\subsection{形式非分歧代数和形式平展代数}
\label{subsection:0.19.10-fr}
\label{0.19.10-fr}

\begin{env}[19.10.1]
\label{0.19.10.1-fr}
我们还要用到两个与形式平滑代数这一概念相近的概念。
\end{env}

\begin{definition}[19.10.2]
\label{0.19.10.2-fr}
设 $A$ 是拓扑环，$B$ 是拓扑 $A$-代数。若 $B$ 作为 $A$-代数满足如下条件，
就称它\emph{形式非分歧}：对任意 $A$-同态 $p:E\to C$（在离散 $A$-代数之间），
只要它的核是幂零的，并对任意连续 $A$-同态 $u:B\to C$，至多存在一个连续
$A$-同态 $v:B\to E$，使 $u$ 可分解为
$B\xrightarrow{v}E\xrightarrow{p}C$。若 $B$ 作为 $A$-代数既形式平滑又形式
非分歧，就称它\emph{形式平展}；换言之，对任意满 $A$-同态 $p:E\to C$
（在离散 $A$-代数之间），只要它的核是幂零的，并对任意连续 $A$-同态
$u:B\to C$，恰有一个连续 $A$-同态
$v:B\to E$，使 $u$ 可分解为 $B\xrightarrow{v}E\xrightarrow{p}C$。
\end{definition}

\oldpage[0\textsubscript{IV-1}]{116}

\begin{examples}[19.10.3]
\label{0.19.10.3-fr}
我们只讨论\emph{离散环}。
\begin{enumerate}
  \item[(i)] 若结构同态 $\rho:A\to B$ 是\emph{满的}，则 $B$ 是一个\emph{形式
  非分歧的 $A$-代数}；事实上，沿用（19.10.2）的记号，复合同态
  $A\xrightarrow{j}B\xrightarrow{v}E$ 必须是结构同态，因此 $v$（若存在）
  是唯一的。反之，$B$ 仅在 $\rho$ 为\emph{双射}时才形式平滑。
  \item[(ii)] 设 $A$ 是环，$S$ 是 $A$ 的一个乘法子集；则
  $B=S^{-1}A$ 是一个\emph{形式平展的 $A$-代数}。事实上，设
  $j:A\to B$ 是典范同态，$\rho:A\to E$ 是一个环同态，$\mathfrak J$ 是 $E$
  的一个幂零理想，$p:E\to C=E/\mathfrak J$ 是典范同态。若 $u:B\to C$ 是
  $A$-同态，则 $u(s/1)$ 在 $C$ 中可逆（$s\in S$），因为 $s/1$ 在 $B$
  中可逆；但 $u(s/1)=p(\rho(s))$，而由于 $\mathfrak J$ 是幂零的，
  $\rho(s)$ 在 $E/\mathfrak J$ 中的像可逆这一事实蕴含 $\rho(s)$ 自身在 $E$
  中可逆（$\mathbf 0_{\mathrm I}$，7.1.12）；由此可知，$\rho$ 能够唯一地分解为
  $A\xrightarrow{j}B\xrightarrow{v}E$，并且显然有
  $u\circ j=(p\circ v)\circ j$，从而 $u=p\circ v$，这就证明了断言。
  \item[(iii)] $k'$ 是域 $k$ 的有限可分扩域，因而是一个形式平滑的 $k$-代数
  （19.6.1）；后面将会看到，$K$ 是域 $k$ 的\emph{有限型}扩域时，它是一个
  形式平展的 $k$-代数，当且仅当它形式非分歧；这又等价于 $K$ 是 $k$ 的有限
  可分扩域
  （21.7.4）。
  \item[(iv)] 多项式代数 $A[X_\alpha]_{\alpha\in I}$ 是环 $A$ 上的代数，并且在
  $A$ 上\emph{形式平滑}（19.3.3），但当 $I$ 非空时，它\emph{并非形式平展}；这由定义
  立即得出。
\end{enumerate}
\end{examples}

\begin{remark}[19.10.4]
\label{0.19.10.4-fr}
与（19.4.3）相同的论证表明，为验证拓扑 $A$-代数 $B$ 形式非分歧，在定义
（19.10.2）中可以只考察 $p:E\to C$ 的核\emph{平方为零}的情形；此外，还可以
只考察 $u(B)\subset p(E)$ 的情形（否则不存在任何 $v$ 分解 $u$），并且通过
把 $C$ 换成 $u(B)$、把 $E$ 换成 $p^{-1}(u(B))$，可以假定 $u$ 和 $p$ 都是满的。
\end{remark}

\section{导射与微分}
\label{section:0.20-fr}
\label{0.20-fr}

本节引入的概念将在第四章 §~16 中以几何形式重新讨论，并将在预概形研究中发挥
重要作用。它们在本章中的重要性，首先来自它们同形式平滑代数概念的关系，特别是
定理（20.4.9）、（20.5.7）和（20.5.12）；这些定理将在第四章论述平滑态射的一节
中改写为几何语言，并在应用中经常使用。另一方面，微分概念将在 §~22 中用来证明
若干重要的正则性判别法；这些判别法对于第四章 §~7 对 Noether 局部环所作的深入
研究至关重要。

\subsection{导射与扩充代数}
\label{subsection:0.20.1-fr}

\oldpage[0\textsubscript{IV-1}]{117}
\begin{proposition}[20.1.1]
\label{0.20.1.1-fr}
设 $A$ 是环（不必交换），$B$、$E$、$C$ 是 $A$-环，$p:E\to C$ 是一个
$A$-同态，其核 $\mathfrak{Q}$ 的平方为零，而 $u:B\to C$ 是一个 $A$-同态。
假定存在一个 $A$-同态 $v_0:B\to E$，使 $u$ 可分解为
$B\xrightarrow{v_0}E\xrightarrow{p}C$。则所有 $A$-同态 $v:B\to E$
中满足 $u=p\circ v$ 的那些同态所成的集合，恰等于所有映射 $v_0+D$ 的集合，其中
$D:B\to\mathfrak{Q}$ 是满足下列两个条件的映射：
\begin{enumerate}
  \item[\rm(i)] $D$ 是 $A$-双模同态。
  \item[\rm(ii)] 对任意 $f$、$g$（它们属于 $B$），都有
\end{enumerate}
\begin{equation}
\tag{20.1.1.1}
\label{0.20.1.1.1-fr}
  D(fg)=f.D(g)+D(f).g.
\end{equation}
\end{proposition}

\begin{proof}
等式 $p(v(f))=p(v_0(f))$（其中 $f\in B$）表示
$D(f)=v(f)-v_0(f)$ 属于 $\mathfrak{Q}$；写出 $v(fg)=v(f)v(g)$，并利用
$\mathfrak{Q}$ 的平方为零，便得到关系 \sref{0.20.1.1.1-fr}；又因 $v$ 和
$v_0$ 都是 $A$-同态，条件 (i) 随之得出。

若 $\rho:A\to B$ 是结构同态，则由 \sref{0.20.1.1.1-fr} 得到
$D(\rho(a)f)=\rho(a)D(f)+D(\rho(a))f$（$a\in A$）；但由 (i)，必须有
$D(\rho(a)f)=\rho(a)D(f)$，故令 $f=1$ 即得
\[
  D(\rho(a))=0 \qquad\text{对}\quad a\in A\,;
\]
反过来，若 $D$ 在 $\rho(A)$ 上为零并满足 \sref{0.20.1.1.1-fr}，则它也满足 (i)。
\end{proof}

\begin{definition}[20.1.2]
\label{0.20.1.2-fr}
给定一个 $A$-环 $B$ 和一个 $B$-双模 $L$，所谓\emph{$A$-导射（从 $B$ 到
$L$）}，是指满足 \sref{0.20.1.1-fr} 中条件 (i) 和 (ii) 的映射 $D:B\to L$。
\end{definition}

由 \sref{0.20.1.1-fr} 可知，一个 $A$-导射 $D:B\to L$ 的\emph{核}是一个
$A$-子环（包含于 $B$）。

有时把 $B$ 到 $L$ 的一个 $\mathbf{Z}$-导射简称为\emph{导射}；也就是说，它是
一个从 $B$ 到 $L$ 的\emph{加性}映射，并满足 \sref{0.20.1.1.1-fr}；因而对这种
映射有 $D(1)=0$。若 $B$ 是一个\emph{素域} $P$ 上的代数，则每个
$\mathbf{Z}$-导射（在 $B$ 上）都是一个 $P$-导射：当 $P$ 的特征 $>0$ 时，这由
前述结论显然可知；在相反情形中，关系 $n.n^{-1}=1$（$n\in\mathbf{Z}$）给出
\[
  nD(n^{-1})=0
\]
故 $D(n^{-1})=0$，因为 $P$ 的特征为 $0$。

由定义 \sref{0.20.1.2-fr} 立即可知，若 $D$、$D'$ 是两个 $A$-导射（从 $B$
到 $L$），则 $D-D'$ 也是如此。换言之，所有 $A$-导射（从 $B$ 到 $L$）构成一个
\emph{加法群}；把这个群记为 $\operatorname{D\acute{e}r}_A(B,L)$。若 $A$ 是
\emph{交换的}，$B$ 是一个交换 $A$-\emph{代数}，而 $L$ 是一个 $B$-模，则对
任意 $A$-导射 $D$（从 $B$ 到 $L$）以及任意 $a\in A$，$aD$ 仍是一个
$A$-导射（从 $B$ 到 $L$）；换言之，$\operatorname{D\acute{e}r}_A(B,L)$ 这时具有
$A$-\emph{模}结构。

\oldpage[0\textsubscript{IV-1}]{118}
命题 \sref{0.20.1.1-fr} 可以解释如下：

\begin{corollary}[20.1.3]
\label{0.20.1.3-fr}
给定 $A$-环之间的两个 $A$-同态 $p:E\to C$、$u:B\to C$，并且前者的核
$\mathfrak{Q}$ 的平方为零，则所有 $A$-同态 $v:B\to E$ 中满足 $u=p\circ v$ 的
那些同态所成的集合或者为空，或者是群 $\operatorname{D\acute{e}r}_A(B,\mathfrak{Q})$
的一个主齐性空间。
\end{corollary}

特别地：

\begin{corollary}[20.1.4]
\label{0.20.1.4-fr}
设 $A$ 是环，$C$ 是 $A$-环，$L$ 是 $C$-双模，$E$ 是一个 $A$-扩充（把
$C$ 以 $L$ 扩充），而 $p:E\to C$ 是增广。同每个导射
$D\in\operatorname{D\acute{e}r}_A(C,L)$ 对应映射 $x\mapsto x+D(p(x))$，所得
映射是从群 $\operatorname{D\acute{e}r}_A(C,L)$ 到 $A$-等价所成群（作用于
$E$ 自身）
的一个同构。
\end{corollary}

\begin{proof}
把 \sref{0.20.1.1-fr} 应用于 $B=E$、$v_0=\mathrm{I}_E$：所有
$A$-同态 $v:E\to E$ 中满足 $p\circ v=p$ 的那些同态所成的集合，恰等于映射
$v=\mathrm{I}_E+D'$ 所成的集合，其中
$D'\in\operatorname{D\acute{e}r}_A(E,L)$。说这样的 $A$-同态 $v$ 是一个
$A$-等价，等价于说 $v$ 在 $L$ 上约化为典范嵌入，也等价于说
$D'(x)=0$ 于 $L$ 中；但这又表示 $D'$ 可分解为
$E\xrightarrow{p}C\xrightarrow{D}L$，其中 $D$ 是一个 $A$-导射。推论得证。
\end{proof}

对于平凡扩充，\sref{0.20.1.3-fr} 给出：

\begin{corollary}[20.1.5]
\label{0.20.1.5-fr}
设 $A$ 是环，$B$、$C$ 是两个 $A$-环，$L$ 是 $C$-双模，而 $u:B\to C$ 是
一个 $A$-同态。对每个导射 $D\in\operatorname{D\acute{e}r}_A(B,L)$，令其对应
映射 $\alpha:x\mapsto(u(x),D(x))$；这样得到一个双射，其值域是所有
$A$-同态 $v:B\to D_C(L)$ 中满足如下条件的同态所成的集合：$u$ 可分解为
$B\xrightarrow{v}D_C(L)\to C$（参见 18.2.3）。
\end{corollary}

更特别地：

\begin{corollary}[20.1.6]
\label{0.20.1.6-fr}
设 $A$ 是环，$B$ 是 $A$-环，$L$ 是 $B$-双模。对每个导射
$D\in\operatorname{D\acute{e}r}_A(B,L)$，令它对应于：$1^\circ$ 一个
$A$-等价 $(x,y)\mapsto(x,y+D(x))$（作用于扩充 $D_B(L)$ 本身）；$2^\circ$ 一个
$A$-同态 $x\mapsto(x,D(x))$（从 $B$ 到 $D_B(L)$），它是增广同态
$D_B(L)\to B$ 的右逆。
由此在下列集合之间定义了典范双射：
\begin{enumerate}
  \item[\rm(i)] 集合 $\operatorname{D\acute{e}r}_A(B,L)$；
  \item[\rm(ii)] $A$-等价所成的集合（作用于 $D_B(L)$ 自身）；
  \item[\rm(iii)] $A$-同态右逆所成的集合（针对增广同态 $D_B(L)\to B$）。
\end{enumerate}
\end{corollary}

应当注意，这样建立的 (i) 与 (ii) 之间的双射保持群结构，而由此导出的 (ii) 与
(iii) 之间的双射，正是（18.3.8）中已经定义的双射。

\begin{corollary}[20.1.7]
\label{0.20.1.7-fr}
设 $A$ 是环，$C$ 是 $A$-环，$L$ 是 $A$-双模，$E$ 是一个 $A$-平凡扩充（把
$C$ 以 $L$ 扩充），而 $p:E\to C$ 是增广。所有 $A$-导射 $D$（从 $E$ 到 $L$）
中满足 $D(x)=x$ 于 $L$ 中的那些导射所成的集合，恰等于映射
$x\mapsto x-w(p(x))$ 所成的集合，其中 $w$ 遍历所有 $A$-同态右逆（针对 $p$）。
\end{corollary}

\begin{proof}
仍只须把 \sref{0.20.1.1-fr} 应用于 $B=E$、$v_0=\mathrm{I}_E$；若
$v=\mathrm{I}_E-D$，则条件 $D(x)=x$（$x\in L$）等价于 $v(x)=0$
（$x\in L$），也就是 $v=w\circ p$，其中 $w:C\to E$ 是一个 $A$-同态；
此外，条件 $p\circ v=p$
等价于 $p\circ w=\mathrm{I}_C$，也就是说 $w$ 是 $p$ 的右逆。
\end{proof}

\subsection{导射的函子性质}
\label{subsection:0.20.2-fr}

\oldpage[0\textsubscript{IV-1}]{119}
\begin{env}[20.2.1]
\label{0.20.2.1-fr}
设 $A$ 是环，$B$ 是 $A$-环，$L$ 是 $B$-双模；若 $L'$ 是另一个 $B$-双模，
$w:L\to L'$ 是一个 $B$-双模同态，则显然映射 $D\mapsto w\circ D$ 是加法群同态
\begin{equation}
\tag{20.2.1.1}
\label{0.20.2.1.1-fr}
  w_0:\operatorname{D\acute{e}r}_A(B,L)\to\operatorname{D\acute{e}r}_A(B,L')
\end{equation}
并且若 $w':L'\to L''$ 是另一个 $B$-双模同态，则
$(w'\circ w)_0=w'_0\circ w_0$。当 $A$ 交换，$B$ 是交换 $A$-代数且 $L$ 是
$B$-模时，\sref{0.20.2.1.1-fr} 是 $A$-模同态。

其次，设 $B'$ 是 $A$-环，$v:B'\to B$ 是一个 $A$-同态，它使 $L$ 成为
$B'$-双模；则映射 $D\mapsto D\circ v$ 是加法群同态
\begin{equation}
\tag{20.2.1.2}
\label{0.20.2.1.2-fr}
  v^0:\operatorname{D\acute{e}r}_A(B,L)\to\operatorname{D\acute{e}r}_A(B',L)
\end{equation}
这是 \sref{0.20.1.1.1-fr} 的直接结果；若 $v':B''\to B'$ 是另一个
$A$-同态，则 $(v\circ v')^0=v'^0\circ v^0$。当 $A$、$B$、$B'$ 交换且
$L$ 是 $B$-模时，\sref{0.20.2.1.2-fr} 是 $A$-模同态。

最后，设 $u:A'\to A$ 是一个使 $B$ 成为 $A'$-环的环同态；每个
$A$-导射 $D\in\operatorname{D\acute{e}r}_A(B,L)$ 也是一个 $A'$-导射，因而得到
交换群的典范单射
\begin{equation}
\tag{20.2.1.3}
\label{0.20.2.1.3-fr}
  u^0:\operatorname{D\acute{e}r}_A(B,L)\to\operatorname{D\acute{e}r}_{A'}(B,L)
\end{equation}
并且若 $u':A''\to A'$ 是另一个环同态，则
$(u\circ u')^0=u'^0\circ u^0$；当 $A$、$A'$、$B$ 交换且 $L$ 是 $B$-模时，
\sref{0.20.2.1.3-fr} 是模的半线性同态（相对于 $u$）。

也可以说
\[
  (A,B,L)\leadsto\operatorname{D\acute{e}r}_A(B,L)
\]
是从（18.3.5）中定义的范畴 $\boldsymbol{K}$ 到交换群范畴
$\boldsymbol{Ab}$ 的一个\emph{协变函子}：它把每个三元组 $(u,v,w)$——它构成
$\boldsymbol{K}$ 的一个态射——对应到同态 $w_0\circ v^0\circ u^0$。函子性的验证
归结为下列图表的交换性：
\[
\xymatrix{
  \operatorname{D\acute{e}r}_A(B,L) \ar[r]^{v^0} \ar[d]_{w_0} &
    \operatorname{D\acute{e}r}_A(B',L) \ar[d]^{w_0} \\
  \operatorname{D\acute{e}r}_A(B,L') \ar[r]^{v^0} &
    \operatorname{D\acute{e}r}_A(B',L')
}
\qquad
\xymatrix{
  \operatorname{D\acute{e}r}_A(B,L) \ar[r]^{u^0} \ar[d]_{w_0} &
    \operatorname{D\acute{e}r}_{A'}(B,L) \ar[d]^{w_0} \\
  \operatorname{D\acute{e}r}_A(B,L') \ar[r]^{u^0} &
    \operatorname{D\acute{e}r}_{A'}(B,L')
}
\]
这里 $w:L\to L'$ 是任意 $B$-双模同态。
\end{env}

\oldpage[0\textsubscript{IV-1}]{120}
\begin{theorem}[20.2.2]
\label{0.20.2.2-fr}
设 $u:A\to B$、$v:B\to C$ 是两个环同态，$L$ 是一个 $C$-双模。存在交换群的
典范正合列
\begin{align}
\tag{20.2.2.1}
\label{0.20.2.2.1-fr}
  0\to\operatorname{D\acute{e}r}_B(C,L)\xrightarrow{u^0}
  \operatorname{D\acute{e}r}_A(C,L)&\xrightarrow{v^0}
  \operatorname{D\acute{e}r}_A(B,L)\xrightarrow{\partial} \\
  &\to\operatorname{Exan}_B(C,L)\xrightarrow{u^1}
  \operatorname{Exan}_A(C,L)\xrightarrow{v^1}
  \operatorname{Exan}_A(B,L) \notag
\end{align}
其中 $u^0$、$v^0$ 分别是（20.2.1.3）和（20.2.1.2）中的同态，$u^1$、$v^1$
分别是（18.3.4.1）和（18.3.3.1）中定义的同态，而 $\partial$ 定义如下：对任意
$A$-导射 $D$（从 $B$ 到 $L$），$\partial(D)$ 是一个 $B$-扩充（把 $C$ 以
$L$ 扩充）的类；这个扩充由环 $D_C(L)$ 上的 $A$-同态
$\alpha:x\mapsto(v(x),D(x))$ 定义（参见 \sref{0.20.1.5-fr}）。此外，正合列
\sref{0.20.2.2.1-fr} 关于 $L$ 具有函子性（分别相对于（20.2.1.1）和
（18.3.1.1）中定义的同态）。
\end{theorem}

\begin{proof}
由于 $D$ 是一个 $A$-导射（并且\emph{更不必说}是一个 $\mathbf{Z}$-导射）
（从 $B$ 到 $L$），$A$-同态 $\alpha:x\mapsto(v(x),D(x))$ 确实在
$D_C(L)$ 上定义了一个 $B$-扩充结构，故 $\partial(D)$ 确有定义
\sref{0.20.1.5-fr}。需要在 $5$ 个位置验证正合性：

\medskip
\noindent 1)\enspace
在 $\operatorname{D\acute{e}r}_B(C,L)$ 处的正合性是平凡的（参见
\sref{0.20.2.1-fr}）。

\medskip
\noindent 2)\enspace
依定义 \sref{0.20.2.1-fr}，$v^0$ 的核是所有 $A$-导射（从 $C$ 到 $L$）中在
$v(B)$ 上为零的那些导射所成的集合，也就是那些 $A$-导射中同时为 $B$-导射的导射
\sref{0.20.1.1-fr}；这就给出在 $\operatorname{D\acute{e}r}_A(C,L)$ 处的正合性。

\medskip
\noindent 3)\enspace
$\partial$ 的核由那些 $D\in\operatorname{D\acute{e}r}_A(B,L)$ 构成；对于这些
导射，$B$-扩充——它由 $\alpha:x\mapsto(v(x),D(x))$ 定义——是 $B$-平凡的。
这表示（18.2.3）存在一个 $B$-同态 $z\mapsto(z,w(z))$（从 $C$ 到 $D_C(L)$；
这里的 $B$-环结构在 $D_C(L)$ 上由 $\alpha$ 定义）。但这样的同态\emph{更不必说}是一个
$A$-同态，因而具有形式 $z\mapsto(z,D'(z))$，其中
$D'\in\operatorname{D\acute{e}r}_A(C,L)$（\sref{0.20.1.6-fr}）；写出它是
$B$-同态这一条件，便得到
\[
  D(x)z+v(x)D'(z)=D'(v(x)z)=v(x)D'(z)+D'(v(x))z
\]
其中 $x\in B$、$z\in C$，从而有 $D'(v(x))=D(x)$；因此 $\partial$ 的核就是
$v^0$ 的像。

\medskip
\noindent 4)\enspace
$u^1$ 的核是所有这样的 $B$-扩充（把 $C$ 以 $L$ 扩充）的类所成的集合：这些
扩充是 $A$-平凡的（18.3.7），因而在等价意义下具有形式 $D_C(L)$，其中
$A$-环结构由同态 $t\mapsto(u(t),0)$ 定义。另一方面，每个 $B$-扩充结构（在
$D_C(L)$ 上）都由同态 $\alpha:x\mapsto(v(x),D(x))$ 定义，其中 $D$ 是一个
$\mathbf{Z}$-导射（从 $B$ 到 $L$）（\sref{0.20.1.5-fr}）；说这个扩充的
$A$-环结构是从该 $B$-扩充的 $B$-环结构借助 $u:A\to B$ 导出的，就表示
$D(u(t))=0$（$t\in A$），因而表示 $D$ 是一个 $A$-导射，也就是说，所考察的
$B$-扩充的类
具有形式 $\partial(D)$；由此得到在 $\operatorname{Exan}_B(C,L)$ 处的正合性。

\medskip
\noindent 5)\enspace
$v^1$ 的核是所有这样的 $A$-扩充 $E$（把 $C$ 以 $L$ 扩充）的类所成的集合：
这些扩充在 $v(B)$ 上平凡，也就是说，存在一个 $A$-同态 $w:B\to E$，使 $v$
可分解为 $B\xrightarrow{w}E\to C$。而这样的 $A$-同态在 $E$ 上定义一个
$B$-扩充结构，这个结构的类在 $u^1$ 下的像就是给定 $A$-扩充的类；反命题显然，
所以已经证明了在 $\operatorname{Exan}_A(C,L)$ 处的正合性。

最后，关于 $L$ 的函子性从定义立即得出。
\end{proof}

\oldpage[0\textsubscript{IV-1}]{121}
\begin{corollary}[20.2.3]
\label{0.20.2.3-fr}
设 $A$、$B$、$C$ 是三个交换环，$u:A\to B$、$v:B\to C$ 是两个环同态，
$L$ 是一个 $C$-模。存在 $A$-模的典范正合列
\begin{align}
\tag{20.2.3.1}
\label{0.20.2.3.1-fr}
  0\to\operatorname{D\acute{e}r}_B(C,L)\xrightarrow{u^0}
  \operatorname{D\acute{e}r}_A(C,L)&\xrightarrow{v^0}
  \operatorname{D\acute{e}r}_A(B,L)\xrightarrow{\partial} \\
  &\to\operatorname{Exalcom}_B(C,L)\xrightarrow{u^1}
  \operatorname{Exalcom}_A(C,L)\xrightarrow{v^1}
  \operatorname{Exalcom}_A(B,L) \notag
\end{align}
它关于 $L$ 具有函子性。
\end{corollary}

\begin{proof}
只要验证下述事实，论证便与 \sref{0.20.2.2-fr} 中相同：对每个导射
$D\in\operatorname{D\acute{e}r}_A(B,L)$，$\partial(D)$ 确实是一个 $B$-扩充
（把 $C$ 以 $L$ 扩充）的类，并且这个扩充是交换 $B$-代数；但这立即由 $C$
的交换性以及 $L$ 是 $C$-模这一事实得出。
\end{proof}

\begin{corollary}[20.2.4]
\label{0.20.2.4-fr}
在 \sref{0.20.2.2-fr}（相应地，\sref{0.20.2.3-fr}）的假设下，存在关于 $L$
具有函子性的典范正合列
\begin{equation}
\tag{20.2.4.1}
\label{0.20.2.4.1-fr}
  0\to\operatorname{D\acute{e}r}_B(C,L)\xrightarrow{u^0}
  \operatorname{D\acute{e}r}_A(C,L)\xrightarrow{v^0}
  \operatorname{D\acute{e}r}_A(B,L)\xrightarrow{\partial}
  \operatorname{Exan}_{B/A}(C,L)\to0
\end{equation}
（相应地，
\begin{equation}
\tag{20.2.4.2}
\label{0.20.2.4.2-fr}
  0\to\operatorname{D\acute{e}r}_B(C,L)\xrightarrow{u^0}
  \operatorname{D\acute{e}r}_A(C,L)\xrightarrow{v^0}
  \operatorname{D\acute{e}r}_A(B,L)\xrightarrow{\partial}
  \operatorname{Exalcom}_{B/A}(C,L)\to0\text{）。}
\end{equation}
这由 $\operatorname{Exan}_{B/A}(C,L)$（相应地，
$\operatorname{Exalcom}_{B/A}(C,L)$）的定义得出（（18.3.7）和（18.4.2））。
\end{corollary}

\begin{remark}[20.2.5]
\label{0.20.2.5-fr}
假定有一个由环同态组成的交换图表
\[
\xymatrix{
  A \ar[r] & B \ar[r] & C \\
  A' \ar[u] \ar[r] & B' \ar[u] \ar[r] & C' \ar[u]
}
\]
则有交换图表
\[
\xymatrix@C=0.5em{
  0 \ar[r] & \operatorname{D\acute{e}r}_B(C,L) \ar[r] \ar[d] &
  \operatorname{D\acute{e}r}_A(C,L) \ar[r] \ar[d] &
  \operatorname{D\acute{e}r}_A(B,L) \ar[r]^-\partial \ar[d] &
  \operatorname{Exan}_B(C,L) \ar[r] \ar[d] &
  \operatorname{Exan}_A(C,L) \ar[r] \ar[d] &
  \operatorname{Exan}_A(B,L) \ar[d] \\
  0 \ar[r] & \operatorname{D\acute{e}r}_{B'}(C',L) \ar[r] &
  \operatorname{D\acute{e}r}_{A'}(C',L) \ar[r] &
  \operatorname{D\acute{e}r}_{A'}(B',L) \ar[r]^-\partial &
  \operatorname{Exan}_{B'}(C',L) \ar[r] &
  \operatorname{Exan}_{A'}(C',L) \ar[r] &
  \operatorname{Exan}_{A'}(B',L)
}
\]
对正合列 \sref{0.20.2.3.1-fr}、\sref{0.20.2.4.1-fr} 和
\sref{0.20.2.4.2-fr} 也有同样的结论。
\end{remark}

\subsection{拓扑环上的连续导射}
\label{subsection:0.20.3-fr}

\begin{env}[20.3.1]
\label{0.20.3.1-fr}
给定两个拓扑环 $A$、$B$（一如既往，都带线性拓扑），我们用
$\operatorname{Hom.\ cont}(A,B)$ 表示从 $A$ 到 $B$ 的连续同态所成的集合。给定
一个拓扑环 $A$，拓扑 $A$-环的范畴定义为 $A$-环的范畴（18.1.4），只须处处以
“拓扑环”代替“环”，以“连续同态”代替“同态”；若 $B$

\oldpage[0\textsubscript{IV-1}]{122}
和 $C$ 是两个拓扑 $A$-环，则用 $\operatorname{Hom.\ cont}_A(B,C)$ 表示连续
$A$-同态（从 $B$ 到 $C$）所成的集合。

设 $A$ 是拓扑环，$B$ 是拓扑 $A$-环，$L$ 是拓扑 $B$-双模；用
$\operatorname{D\acute{e}r.\ cont}_A(B,L)$ 表示\emph{连续} $A$-导射（从 $B$
到 $L$）所成的集合；显然，它是 $\operatorname{D\acute{e}r}_A(B,L)$ 的一个
\emph{子群}（并且它是一个子 $A$-模，当 $A$、$B$ 交换且 $L$ 是 $B$-模时）。
\end{env}

\begin{env}[20.3.2]
\label{0.20.3.2-fr}
显然，把命题 \sref{0.20.1.1-fr} 中的“环”换成“拓扑环”，把“同态”换成“连续
同态”，该命题仍然成立。同样，在 \sref{0.20.1.3-fr} 和
\sref{0.20.1.4-fr} 中把 $\operatorname{D\acute{e}r}$ 换成
$\operatorname{D\acute{e}r.\ cont}$，以“拓扑环”（相应地，“拓扑双模”）代替“环”
（相应地，“双模”），并以“连续同态”代替“同态”，所得结论仍然成立；当然，
在 \sref{0.20.1.4-fr} 中须假定 $p$ 连续，并以“连续 $A$-等价”代替
“$A$-等价”。对 \sref{0.20.1.5-fr} 和 \sref{0.20.1.6-fr} 也有同样的结果，
只须在 $D_C(L)$（相应地，$D_B(L)$）上取\emph{积}拓扑（即在 $C\times L$
上，相应地在 $B\times L$ 上的积拓扑）。
\end{env}

\begin{proposition}[20.3.3]
\label{0.20.3.3-fr}
设 $A$ 是拓扑环，$B$ 是拓扑 $A$-环，$L$ 是一个离散拓扑 $B$-双模，并被
$B$ 的某个开双边理想零化。若在 $B$ 中每个开双边理想的平方仍为开集，则
$\operatorname{D\acute{e}r.\ cont}_A(B,L)=\operatorname{D\acute{e}r}_A(B,L)$。
\end{proposition}

\begin{proof}
事实上，若 $\mathfrak{R}$ 是 $B$ 的一个开双边理想并零化 $L$，而 $D$ 是一个
$A$-导射（从 $B$ 到 $L$），则
$D(\mathfrak{R}^2)\subset\mathfrak{R}.D(\mathfrak{R})+
D(\mathfrak{R}).\mathfrak{R}=0$（\sref{0.20.1.1.1-fr}），所以 $D$ 连续。
\end{proof}

\begin{env}[20.3.4]
\label{0.20.3.4-fr}
把 \sref{0.20.2.1-fr} 中的“环”换成“拓扑环”，把“双模”换成“拓扑双模”，
把“同态”换成“连续同态”，并把 $\operatorname{D\acute{e}r}$ 换成
$\operatorname{D\acute{e}r.\ cont}$，其中的所有结果仍然成立。
\end{env}

\begin{proposition}[20.3.5]
\label{0.20.3.5-fr}
设 $A$ 是拓扑环，$B$ 是拓扑 $A$-环，$L$ 是一个离散拓扑 $B$-双模，并被
$B$ 的某个开双边理想零化。于是存在典范同构
\begin{equation}
\tag{20.3.5.1}
\label{0.20.3.5.1-fr}
  \varinjlim\operatorname{D\acute{e}r}_{A/\mathfrak{J}}
  (B/\mathfrak{R},L)\xrightarrow{\sim}
  \operatorname{D\acute{e}r.\ cont}_A(B,L)
\end{equation}
其中左端的归纳极限取遍双边理想偶 $(\mathfrak{J},\mathfrak{R})$ 所成的滤过有序集，
这些理想偶满足
$\mathfrak{R}.L=L.\mathfrak{R}=0$、
$\mathfrak{J}.B\subset\mathfrak{R}$、
$B.\mathfrak{J}\subset\mathfrak{R}$。
\end{proposition}

\begin{proof}
由于 $A/\mathfrak{J}$ 和 $B/\mathfrak{R}$ 都是离散的，存在典范同态
$w_{\mathfrak{J},\mathfrak{R}}:
\operatorname{D\acute{e}r}_{A/\mathfrak{J}}(B/\mathfrak{R},L)\to
\operatorname{D\acute{e}r.\ cont}_A(B,L)$，它们构成一个归纳系
\sref{0.20.3.4-fr}；取归纳极限即得同态 \sref{0.20.3.5.1-fr}。因为同态
$B/\mathfrak{R}'\to B/\mathfrak{R}$ 在 $\mathfrak{R}\supset\mathfrak{R}'$
时是满的，故由定义立即可知，同态
$\operatorname{D\acute{e}r}_{A/\mathfrak{J}}(B/\mathfrak{R},L)\to
\operatorname{D\acute{e}r}_{A/\mathfrak{J}}(B/\mathfrak{R}',L)$（其中
$\mathfrak{R}.L=L.\mathfrak{R}=0$、
$\mathfrak{J}.B\subset\mathfrak{R}'$、
$B.\mathfrak{J}\subset\mathfrak{R}'$）是单射；显然，同态
$\operatorname{D\acute{e}r}_{A/\mathfrak{J}}(B/\mathfrak{R},L)\to
\operatorname{D\acute{e}r}_{A/\mathfrak{J}'}(B/\mathfrak{R},L)$ 在
$\mathfrak{J}'\subset\mathfrak{J}$ 时也是单射
（\sref{0.20.2.1-fr}）。由此可知，同态 \sref{0.20.3.5.1-fr} 是单射。
另一方面，若 $D$ 是一个连续 $A$-导射（从 $B$ 到 $L$），则它的核包含一个开
双边理想 $\mathfrak{R}_0$（属于 $B$）；若 $\mathfrak{J}_0$ 是 $A$ 的一个开双边理想，
满足 $\mathfrak{J}_0B\subset\mathfrak{R}_0$ 和
$B\mathfrak{J}_0\subset\mathfrak{R}_0$，则显然 $D$ 是一个
$(A/\mathfrak{J}_0)$-导射（从 $B/\mathfrak{R}_0$ 到 $L$）的典范像；因此
\sref{0.20.3.5.1-fr} 是满射。
\end{proof}

\oldpage[0\textsubscript{IV-1}]{123}
\begin{proposition}[20.3.6]
\label{0.20.3.6-fr}
设 $u:A\to B$、$v:B\to C$ 是拓扑环之间的两个连续同态，$L$ 是一个离散
$C$-双模，并被 $C$ 的某个开双边理想零化。存在典范正合列
\begin{align}
\tag{20.3.6.1}
\label{0.20.3.6.1-fr}
  0\to\operatorname{D\acute{e}r.\ cont}_B(C,L)\xrightarrow{u^0}
  \operatorname{D\acute{e}r.\ cont}_A(C,L)&\xrightarrow{v^0}
  \operatorname{D\acute{e}r.\ cont}_A(B,L)\xrightarrow{\partial} \\
  &\to\operatorname{Exantop}_B(C,L)\xrightarrow{u^1}
  \operatorname{Exantop}_A(C,L)\xrightarrow{v^1}
  \operatorname{Exantop}_A(B,L) \notag
\end{align}
其中 $\partial$ 由 \sref{0.20.2.2.1-fr} 中的同态 $\partial$ 取归纳极限而定义；
这个正合列关于 $L$ 具有函子性（在被开双边理想零化的离散 $C$-双模范畴中）。
\end{proposition}

\begin{proof}
从 \sref{0.20.2.2-fr} 出发，这由函子 $\varinjlim$ 的正合性得出。
\end{proof}

\begin{corollary}[20.3.7]
\label{0.20.3.7-fr}
设 $A$、$B$、$C$ 是三个交换拓扑环，$u:A\to B$、$v:B\to C$ 是两个连续同态，
$L$ 是一个离散 $C$-模，并被 $C$ 的某个开理想零化。存在 $A$-模的典范正合列，
它关于 $L$ 具有函子性：
\begin{align}
\tag{20.3.7.1}
\label{0.20.3.7.1-fr}
  0\to\operatorname{D\acute{e}r.\ cont}_B(C,L)\xrightarrow{u^0}
  \operatorname{D\acute{e}r.\ cont}_A(C,L)&\xrightarrow{v^0}
  \operatorname{D\acute{e}r.\ cont}_A(B,L)\xrightarrow{\partial} \\
  &\to\operatorname{Exalcotop}_B(C,L)\xrightarrow{u^1}
  \operatorname{Exalcotop}_A(C,L)\xrightarrow{v^1}
  \operatorname{Exalcotop}_A(B,L). \notag
\end{align}
\end{corollary}

\begin{corollary}[20.3.8]
\label{0.20.3.8-fr}
在 \sref{0.20.3.5-fr}（相应地，\sref{0.20.3.6-fr}）的假设下，存在关于 $L$
具有函子性的典范正合列
\begin{equation}
\tag{20.3.8.1}
\label{0.20.3.8.1-fr}
  0\to\operatorname{D\acute{e}r.\ cont}_B(C,L)\xrightarrow{u^0}
  \operatorname{D\acute{e}r.\ cont}_A(C,L)\xrightarrow{v^0}
  \operatorname{D\acute{e}r.\ cont}_A(B,L)\to
  \operatorname{Exantop}_{B/A}(C,L)\to0
\end{equation}
（相应地，
\begin{equation}
\tag{20.3.8.2}
\label{0.20.3.8.2-fr}
  0\to\operatorname{D\acute{e}r.\ cont}_B(C,L)\xrightarrow{u^0}
  \operatorname{D\acute{e}r.\ cont}_A(C,L)\xrightarrow{v^0}
  \operatorname{D\acute{e}r.\ cont}_A(B,L)\to
  \operatorname{Exalcotop}_{B/A}(C,L)\to0\text{）。}
\end{equation}
\end{corollary}

类似于 \sref{0.20.2.5-fr} 中的图表，留给读者写出。

\subsection{主部与微分}
\label{subsection:0.20.4-fr}

在本节以及以下三节中，所有环都假定为\emph{交换的}。

\begin{env}[20.4.1]
\label{0.20.4.1-fr}
设 $A$ 是拓扑环，$B$ 是拓扑 $A$-代数；在 $A$-代数 $B\otimes_A B$ 上赋予
\emph{张量积拓扑}，它由此成为拓扑 $A$-代数；我们用 $p$（或 $p_{B/A}$）表示
典范满 $A$-同态
\begin{equation}
\tag{20.4.1.1}
\label{0.20.4.1.1-fr}
  p:B\otimes_A B\to B
\end{equation}
它满足 $p(b\otimes b')=bb'$；显然 $p$ 连续。$p$ 的\emph{核}记为
$\mathfrak{J}_{B/A}$（若不致混淆，也只记为 $\mathfrak{J}$）。用
$j_1:B\to B\otimes_A B$ 和 $j_2:B\to B\otimes_A B$ 表示两个典范
$A$-同态，它们满足
\[
  j_1(b)=b\otimes1, \qquad j_2(b)=1\otimes b
\]
并且都是连续的。
\end{env}

\oldpage[0\textsubscript{IV-1}]{124}
\begin{definition}[20.4.2]
\label{0.20.4.2-fr}
所谓\emph{$B$-增广的 $1$ 阶主部代数（$B$ 关于 $A$）}，记为
$\mathrm{P}^1_{B/A}$，是拓扑 $A$-商代数
\begin{equation}
\tag{20.4.2.1}
\label{0.20.4.2.1-fr}
  \mathrm{P}^1_{B/A}=(B\otimes_A B)/\mathfrak{J}^2
\end{equation}
在它上面赋予拓扑 $B$-代数结构，该结构由同态
$\bar{j}_1:B\to\mathrm{P}^1_{B/A}$ 定义（该同态由 $j_1$ 与典范同态
$B\otimes_A B\to\mathrm{P}^1_{B/A}$ 复合而得），并赋予 $B$-代数增广
$\varepsilon:\mathrm{P}^1_{B/A}\to B$（也记为 $\varepsilon_{B/A}$）；后者由
$p$ 取商而得。

因为依定义有 $p(b\otimes1)=b$，显然 $\varepsilon$ 确实是一个 $B$-代数增广。
\end{definition}

\begin{definition}[20.4.3]
\label{0.20.4.3-fr}
增广 $\varepsilon:\mathrm{P}^1_{B/A}\to B$ 的核
\begin{equation}
\tag{20.4.3.1}
\label{0.20.4.3.1-fr}
  \Omega^1_{B/A}=\mathfrak{J}_{B/A}/(\mathfrak{J}_{B/A})^2
\end{equation}
赋予由 $\mathrm{P}^1_{B/A}$ 的拓扑诱导的拓扑，因而成为拓扑 $B$-模；这个
$B$-模称为\emph{$1$-微分}（或简称\emph{微分}）的 $B$-模（关于 $A$）。
\end{definition}

应当注意，$\Omega^1_{B/A}$ 的拓扑也等于如下商拓扑：在
$\mathfrak{J}_{B/A}$ 上由 $B\otimes_A B$ 的拓扑诱导拓扑，再取其商拓扑（Bourbaki，
\emph{Top. gén.}，chap.~III，第 3 版，§~2，n\textsuperscript{o}~7，
prop.~20）。若 $B$ 离散，则 $\Omega^1_{B/A}$ 也离散。用
$\widehat{\Omega}^1_{B/A}$ 表示拓扑 $B$-模 $\Omega^1_{B/A}$ 的\emph{分离完备化}。

每个拓扑环 $B$ 都可以看成拓扑 $\mathbf{Z}$-代数（这里 $\mathbf{Z}$ 赋予离散
拓扑），所以可以定义拓扑 $B$-模 $\Omega^1_{B/\mathbf{Z}}$；它有时也称为 $B$
上的\emph{绝对微分} $B$-模，并记为 $\Omega^1_B$。若 $B$ 是素域 $P$（离散）
上的拓扑代数，则
$B\otimes_P B=B\otimes_{\mathbf{Z}}B$，因而
$\Omega^1_{B/P}=\Omega^1_B$。

\begin{lemma}[20.4.4]
\label{0.20.4.4-fr}
设 $A$ 是环，$B$ 是 $A$-代数。理想 $\mathfrak{J}_{B/A}$（在
$B\otimes_A B$ 中）由元素 $1\otimes s-s\otimes1$ 生成，其中 $s$ 遍历
$A$-代数 $B$ 的一个生成元集。
\end{lemma}

\begin{proof}
显然，对每个 $x\in B$，都有
$x\otimes1-1\otimes x\in\mathfrak{J}$；另一方面，对任意 $x$、$y$（属于
$B$），都有
$x\otimes y=xy\otimes1+(x\otimes1)(1\otimes y-y\otimes1)$。若
$\sum_i(x_i\otimes y_i)\in\mathfrak{J}$，则依定义有
$\sum_i x_i y_i=0$，因而
\begin{equation}
\tag{20.4.4.1}
\label{0.20.4.4.1-fr}
  \sum_i(x_i\otimes y_i)=
  \sum_i(x_i\otimes1)(1\otimes y_i-y_i\otimes1)
\end{equation}
这就证明了 $\mathfrak{J}$ 是由元素 $1\otimes x-x\otimes1$ 生成的理想。此外，
若 $x=st$，则
\begin{equation}
\tag{20.4.4.2}
\label{0.20.4.4.2-fr}
  x\otimes1-1\otimes x=(s\otimes1)(t\otimes1-1\otimes t)
  +(s\otimes1-1\otimes s)(1\otimes t)
\end{equation}
由此通过归纳立即完成证明。
\end{proof}

\begin{proposition}[20.4.5]
\label{0.20.4.5-fr}
设 $A$ 是拓扑环，$B$ 是拓扑 $A$-代数。$\Omega^1_{B/A}$ 的拓扑比由 $B$ 的
拓扑导出的拓扑更粗；若在 $B$ 中每个开理想的平方都是开集，则这两个拓扑相同。
\end{proposition}

\begin{proof}
第一项断言是平凡的，因为 $B\otimes_A B$ 上的张量积拓扑比由 $B$ 的拓扑导出的
拓扑更粗；\emph{更不必说}，$\mathfrak{J}=\mathfrak{J}_{B/A}$ 上由

\oldpage[0\textsubscript{IV-1}]{125}
$B\otimes_A B$ 的拓扑诱导的拓扑，比 $\mathfrak{J}$ 上由 $B$ 的拓扑导出的拓扑
更粗。为证明第二项断言，略微滥用记号，用 $M\otimes N$ 表示子模
$\operatorname{Im}(M\otimes_A N)$；这里考虑的是两个 $A$-子模 $M$、$N$
（都包含于 $B$）。利用关系
\[
\begin{aligned}
  (xy)\otimes z-x\otimes(yz)
  &=(x\otimes1)(1\otimes z)(y\otimes1-1\otimes y)\\
  &=xz.(y\otimes1-1\otimes y)
    -x.(z\otimes1-1\otimes z)(y\otimes1-1\otimes y)
\end{aligned}
\]
在 $B$-模 $B\otimes_A B$ 中（由 $j_1$ 定义），并考虑到
\sref{0.20.4.4-fr}，立即可见：若 $\mathfrak{R}$ 是 $B$ 的一个理想，则
\[
  ((\mathfrak{R}^2\otimes B)+(B\otimes\mathfrak{R}^2))\cap\mathfrak{J}
  \subset(\mathfrak{R}\otimes\mathfrak{R})\cap\mathfrak{J}
  +\mathfrak{R}\mathfrak{J}+\mathfrak{J}^2
\]
另一方面，若 $x_i$、$y_i$ 是 $\mathfrak{R}$ 的元素并且
$\sum_i(x_i\otimes y_i)\in\mathfrak{J}$，则由 \sref{0.20.4.4.1-fr} 得到
$\sum_i(x_i\otimes y_i)\in\mathfrak{R}\mathfrak{J}$，所以最终有
\begin{equation}
\tag{20.4.5.1}
\label{0.20.4.5.1-fr}
  (\mathfrak{R}^2\otimes B+B\otimes\mathfrak{R}^2)\cap\mathfrak{J}
  \subset\mathfrak{R}\mathfrak{J}+\mathfrak{J}^2.
\end{equation}
现在，$0$ 在 $\mathfrak{J}$ 中有一个基本邻域系（就由 $B\otimes_A B$ 的拓扑
诱导的拓扑而言）；取作 $0$ 的邻域的是集合
$(\mathfrak{R}\otimes B+B\otimes\mathfrak{R})\cap\mathfrak{J}$，其中
$\mathfrak{R}$ 遍历所有开理想。又因为 $\Omega^1_{B/A}$ 上由 $B$ 的拓扑导出的
拓扑，也就是对 $\mathfrak{J}^2$ 取商的、在 $\mathfrak{J}$ 上由 $B$ 的拓扑导出的
拓扑，所以关于 $B$ 的开理想的
假设和关系 \sref{0.20.4.5.1-fr} 完成了第二项断言的证明。
\end{proof}

\begin{definition}[20.4.6]
\label{0.20.4.6-fr}
设 $p_1$ 和 $p_2$ 是复合映射
$B\xrightarrow{j_1}B\otimes_A B\to\mathrm{P}^1_{B/A}$、
$B\xrightarrow{j_2}B\otimes_A B\to\mathrm{P}^1_{B/A}$；它们是连续 $A$-同态，
并满足 $\varepsilon\circ p_1=\varepsilon\circ p_2=\mathrm{I}_B$。连续 $A$-同态
（在 $A$-模之间）
\begin{equation}
\tag{20.4.6.1}
\label{0.20.4.6.1-fr}
  d_{B/A}=p_1-p_2:B\to\Omega^1_{B/A}
\end{equation}
称为 $B$ 关于 $A$ 的\emph{外微分}；对每个 $x\in B$，$d_{B/A}(x)$（也记为
$d(x)$ 或 $dx$）称为 $x$ 的\emph{微分}（关于 $A$）。
\end{definition}

当 $A=\mathbf{Z}$ 时，用 $d_B$ 代替 $d_{B/\mathbf{Z}}$；若 $B$ 是素域 $P$ 上的
代数，则 $d_B=d_{B/P}$。

\begin{proposition}[20.4.7]
\label{0.20.4.7-fr}
$B$-模 $\Omega^1_{B/A}$ 由元素 $d_{B/A}(x)$ 生成，其中 $x$ 遍历 $A$-代数
$B$ 的一个生成元系。
\end{proposition}

\begin{proof}
由于 $d_{B/A}(x)$ 是 $x\otimes1-1\otimes x$ 在 $\Omega^1_{B/A}$ 中的典范像，
该命题立即由 \sref{0.20.4.4-fr} 得出。
\end{proof}

\begin{theorem}[20.4.8]
\label{0.20.4.8-fr}
设 $A$ 是拓扑环，$B$ 是拓扑 $A$-代数。
\begin{enumerate}
  \item[\rm(i)] 存在唯一的拓扑增广 $B$-代数同构
\begin{equation}
\tag{20.4.8.1}
\label{0.20.4.8.1-fr}
  \varphi:\mathrm{P}^1_{B/A}\xrightarrow{\sim}D_B(\Omega^1_{B/A})
\end{equation}
  它在 $\Omega^1_{B/A}$ 上约化为恒等映射。
  \item[\rm(ii)] 同态 $d_{B/A}$ 是一个 $A$-导射（从 $B$ 到
  $\Omega^1_{B/A}$），并具有如下泛性质：对每个拓扑 $B$-模 $L$，映射
  $u\mapsto u\circ d_{B/A}$ 是一个 $A$-模同构
\begin{equation}
\tag{20.4.8.2}
\label{0.20.4.8.2-fr}
  \operatorname{Hom.\ cont}_B(\Omega^1_{B/A},L)
  \xrightarrow{\sim}\operatorname{D\acute{e}r.\ cont}_A(B,L).
\end{equation}
\end{enumerate}
\end{theorem}

\begin{proof}
\oldpage[0\textsubscript{IV-1}]{126}
(i) 显然，$\varphi$（采用 \sref{0.20.4.6-fr} 的记号）必为映射
$z\mapsto(\varepsilon(z),z-p_1(\varepsilon(z)))$，而逆同构必为映射
$(b,x)\mapsto p_1(b)+x$；这两个映射都是连续的，所以第一项断言得证。应当注意，
这还蕴含 $B$ 的拓扑（通过 $p_1$）等同于 $B\otimes_A B$ 的拓扑对
$\mathfrak{J}_{B/A}$ 取\emph{商}所得的拓扑。

(ii) $d_{B/A}$ 是一个 $A$-导射（从 $B$ 出发）这一事实，由定义
\sref{0.20.4.6-fr} 和 \sref{0.20.1.1-fr} 得出。

为证明泛性质，回忆 $\operatorname{D\acute{e}r.\ cont}_A(B,L)$ 典范地等同于集合
$G$；该集合由连续 $A$-代数同态 $u:B\to D_B(L)$ 构成，并且复合映射
$B\xrightarrow{u}D_B(L)\xrightarrow{q}B$ 是恒等映射（这里 $q(b,0)=b$ 是投影；
参见 \sref{0.20.1.6-fr} 和 \sref{0.20.3.2-fr}）。另一方面，借助同构 $\varphi$，
$\operatorname{Hom.\ cont}_B(\Omega^1_{B/A},L)$ 典范地等同于连续 $B$-代数同态
$v:\mathrm{P}^1_{B/A}\to D_B(L)$ 所成的集合，并且复合映射
$\mathrm{P}^1_{B/A}\xrightarrow{v}D_B(L)\xrightarrow{q}B$ 等于增广
$\varepsilon$。由于依定义有 $p_2=p_1-d_{B/A}$，问题归结为证明每个 $u\in G$
都可分解为
\[
\xymatrix{
  B \ar[r]^u \ar[d]_{p_2} & D_B(L) \\
  \mathrm{P}^1_{B/A} \ar[ur]_v &
}
\]
其中 $v$ 是连续 $B$-同态。现在，已经有一个连续 $A$-代数同态
$j:b\mapsto(b,0)$（从 $B$ 到 $D_B(L)$），它属于 $G$；根据拓扑代数张量积的
定义（$\mathbf{0_I}$，7.7.6），因此存在连续 $A$-代数同态
$w:B\otimes_A B\to D_B(L)$，使图表
\[
\xymatrix{
  B\otimes_A B \ar[dr]^w & B \ar[l]_{j_1} \ar[d]^j \\
  B \ar[u]^{j_2} \ar[r]_u & D_B(L)
}
\]
交换。于是依定义有
$w(b\otimes1-1\otimes b)=j(b)-u(b)\in L$；根据 \sref{0.20.4.4-fr}，
这蕴含 $w(\mathfrak{J})\subset L$，从而 $w(\mathfrak{J}^2)=0$；因此 $w$ 可分解为
\[
  B\otimes_A B\to\mathrm{P}^1_{B/A}\xrightarrow{v}D_B(L)
\]
其中 $v$ 是连续 $A$-代数同态。此外，因为 $v\circ p_1=j$ 是 $B$-代数同态，
$v$ 也是 $B$-代数同态；这是由 $\mathrm{P}^1_{B/A}$ 上该结构的定义得出的。又因
依定义有 $u=v\circ p_2$，证明完成。
\end{proof}

\begin{theorem}[20.4.9]
\label{0.20.4.9-fr}
假定 $B$ 是形式平滑的拓扑 $A$-代数。则拓扑 $B$-模 $\Omega^1_{B/A}$ 是
\emph{形式投射的}。
\end{theorem}

\begin{proof}
事实上，该假设蕴含 $B\otimes_A B$（赋予由 $B$-代数结构、亦即由 $j_1$ 定义的
结构）是形式平滑的拓扑 $B$-代数（\sref{0.19.3.5-fr}，(iii)），因而也是形式
平滑的拓扑 $A$-代数（\sref{0.19.3.5-fr}，(ii)）；因为 $B$ 拓扑同构于
$A$-商代数 $(B\otimes_A B)/\mathfrak{J}$，并且是形式平滑的 $A$-代数，结论由
\sref{0.19.5.3-fr} 得出。
\end{proof}

\oldpage[0\textsubscript{IV-1}]{127}
\begin{corollary}[20.4.10]
\label{0.20.4.10-fr}
再假定在 $B$ 中每个开理想的平方都是开集；则对每个开理想
$\mathfrak{R}$（属于 $B$），
$\Omega^1_{B/A}\otimes_B(B/\mathfrak{R})
=\Omega^1_{B/A}/\mathfrak{R}.\Omega^1_{B/A}$ 是一个投射
$(B/\mathfrak{R})$-模。
\end{corollary}

\begin{proof}
事实上，此时 $\Omega^1_{B/A}$ 的拓扑由 $B$ 的拓扑导出
（\sref{0.20.4.5-fr}），应用 \sref{0.19.2.4-fr} 即可。
\end{proof}

\begin{corollary}[20.4.11]
\label{0.20.4.11-fr}
设 $A$、$B$ 是两个 Noether 局部环，$\rho:A\to B$ 是一个局部同态，它使 $B$
成为形式平滑的拓扑 $A$-代数（相对于预进拓扑）。则对每个定义理想
$\mathfrak{b}$（属于 $B$），$\Omega^1_{B/A}\otimes_B(B/\mathfrak{b})$ 是自由
$(B/\mathfrak{b})$-模。
\end{corollary}

\begin{proof}
事实上，由 \sref{0.20.4.10-fr} 可知这个模是投射的；又因为 $B/\mathfrak{b}$
是 Artin 环，每个投射 $(B/\mathfrak{b})$-模都是自由的
（$\mathbf{0_{III}}$，10.1.3）。
\end{proof}

\begin{proposition}[20.4.12]
\label{0.20.4.12-fr}
设 $A$ 是拓扑环，$B$ 是拓扑 $A$-代数。若结构同态 $A\to B$ 是满的，则
$\Omega^1_{B/A}=0$。
\end{proposition}

\begin{proof}
事实上，根据 \sref{0.20.1.1-fr}，有
$\operatorname{D\acute{e}r}_A(B,L)=0$；这里所指的是任意 $B$-模 $L$。因此该命题立即由
\sref{0.20.4.8-fr} 得出。
\end{proof}

\begin{examples}[20.4.13]
\label{0.20.4.13-fr}
\begin{enumerate}
  \item[\rm(i)] 设 $A$ 是环，
  $B=A[X_\alpha]_{\alpha\in I}$ 是 $A$ 上的多项式代数。则
  $\Omega^1_{B/A}$ 是自由 $B$-模，而各 $dX_\alpha$ 构成它的一组基。事实上，
  这些 $dX_\alpha$ 生成该 $B$-模（\sref{0.20.4.7-fr}）。另一方面，设 $L$ 是
  自由 $B$-模，并且以 $(z_\alpha)_{\alpha\in I}$ 为一组基，其中 $I$ 是指标集；
  则存在一个 $A$-同态 $u$，它从 $B$ 到 $D_B(L)$，使得
  $u(X_\alpha)=(X_\alpha,z_\alpha)$ 对每个 $\alpha\in I$ 都成立，从而根据
  \sref{0.20.1.5-fr}，存在一个 $A$-导射 $D$，它从 $B$ 到 $L$，使得
  $D(X_\alpha)=z_\alpha$ 对每个 $\alpha$ 都成立；由
  \sref{0.20.4.8.1-fr}，这就证明了各
  $dX_\alpha$ 线性无关。
  \item[\rm(ii)] 设 $A$ 是环，$L$ 是 $A$-模，$B$ 是 $A$-代数
  $D_A(L)$；则典范同态 $x\mapsto d_{B/A}x$ 从 $L$ 到
  $\Omega^1_{B/A}$，并且它是\emph{双射}。事实上，立即可见，各 $B$-导射从
  $B=D_A(L)$ 到一个 $B$-模 $M$ 时，恰为形如
  $(b,x)\mapsto u(x)$ 的映射，其中 $u\in\operatorname{Hom}_A(L,M)$；于是由
  \sref{0.20.4.8-fr}，(ii) 即得结论。
  \item[\rm(iii)] 假定 $B$ 是两个拓扑 $A$-代数 $B_1$、$B_2$ 的乘积
  （把它们视为 $B$ 的理想）。则 $B_1\otimes_A B_2$ 和
  $B_2\otimes_A B_1$ 在同态 $p$（\sref{0.20.4.1.1-fr}）下的像均为零；
  因而立即得到，$\mathrm{P}^1_{B/A}$ 等同于\emph{乘积}
  $\mathrm{P}^1_{B_1/A}\times\mathrm{P}^1_{B_2/A}$，而 $B$-模
  $\Omega^1_{B/A}$ 是 $B$-模 $\Omega^1_{B_1/A}$ 与
  $\Omega^1_{B_2/A}$ 的（拓扑）\emph{直和}（它们分别被 $B_2$ 与 $B_1$
  零化）。
  \item[\rm(iv)] 设 $A$、$B$ 是两个整环，且 $A\subset B$；再假定
  $A$ 整闭，$B$ 整于 $A$，并且 $B$ 的分式域是 $A$ 的分式域的可分扩张。
  则 $B$-模 $\Omega^1_{B/A}$ 是\emph{挠模}。事实上，对每个 $x\in B$，
  $x$ 关于 $K$（即 $A$ 的分式域）的最小多项式 $f(T)$，是属于 $A[T]$ 的
  一个多项式；因为 $x$ 在 $K$ 上可分，故 $f'(x)\neq0$，另一方面由关系
  $f(x)=0$ 推得 $f'(x)d_{B/A}x=0$，于是根据 \sref{0.20.4.7-fr} 得到断言。
\end{enumerate}
\end{examples}

\begin{remarks}[20.4.14]
\label{0.20.4.14-fr}
\begin{enumerate}
  \item[\rm(i)] 应当注意，若不计其拓扑，$B$-模 $\Omega^1_{B/A}$ 与
  $A$、$B$ 的拓扑无关。
  \item[\rm(ii)] 稍后我们将引入“$n$ 阶主部代数”，这是 $B$ 关于 $A$ 的代数
  $\mathrm{P}^n_{B/A}=(B\otimes_A B)/(\mathfrak{J}_{B/A})^{n+1}$；它是
  “$n$ 阶微分演算”的基础。
\end{enumerate}
\end{remarks}

\subsection{
\texorpdfstring{$\Omega^1_{B/A}$}{Omega1(B/A)} 的一些基础性的函子性质}
\label{subsection:0.20.5-fr}

\oldpage[0\textsubscript{IV-1}]{128}
\begin{env}[20.5.1]
\label{0.20.5.1-fr}
在本小节及下一小节中，除非另有明确说明，所考虑的环和模均假定赋予
\emph{离散}拓扑。
\end{env}

\begin{env}[20.5.2]
\label{0.20.5.2-fr}
设 $A$ 是环，$B$、$C$ 是两个 $A$-代数，$u:B\to C$ 是一个
$A$-同态；有交换图
\[
\xymatrix{
  B\otimes_A B \ar[r]^{u\otimes u} \ar[d]_{p_{B/A}} &
  C\otimes_A C \ar[d]^{p_{C/A}} \\
  B \ar[r]^u & C
}
\tag{20.5.2.1}
\label{0.20.5.2.1-fr}
\]
由此取商得到一个代数的 $A$-同态
\begin{equation}
\tag{20.5.2.2}
\label{0.20.5.2.2-fr}
  u':\mathrm{P}^1_{B/A}\to\mathrm{P}^1_{C/A}
\end{equation}
使得图
\[
\xymatrix{
  \mathrm{P}^1_{B/A} \ar[r]^{u'} & \mathrm{P}^1_{C/A} \\
  B \ar[u]^{p_1} \ar[r]^u & C \ar[u]_{p_1}
}
\]
交换；因为 $u\otimes u$ 把 $\mathfrak{J}_{B/A}$ 映入
$\mathfrak{J}_{C/A}$，将 $u'$ 限制于 $\Omega^1_{B/A}$，就得到一个映射
\begin{equation}
\tag{20.5.2.3}
\label{0.20.5.2.3-fr}
  u'':\Omega^1_{B/A}\to\Omega^1_{C/A}
\end{equation}
使得偶对 $(u'',u)$ 对 $B$-模 $\Omega^1_{B/A}$ 的结构和
$C$-模 $\Omega^1_{C/A}$ 的结构而言，是一个\emph{双同态}；由后一事实典范地
导出一个\emph{$C$-模同态}
\begin{equation}
\tag{20.5.2.4}
\label{0.20.5.2.4-fr}
  u_{C/B/A}:\Omega^1_{B/A}\otimes_B C\to\Omega^1_{C/A}.
\end{equation}

\oldpage[0\textsubscript{IV-1}]{129}
此外，因为图
\[
\xymatrix{
  \mathrm{P}^1_{B/A} \ar[r]^{u'} & \mathrm{P}^1_{C/A} \\
  B \ar[u]^{p_2} \ar[r]^u & C \ar[u]_{p_2}
}
\tag{20.5.2.5}
\label{0.20.5.2.5-fr}
\]
也交换，遂知图
\[
\xymatrix{
  \Omega^1_{B/A} \ar[r]^{u''} & \Omega^1_{C/A} \\
  B \ar[u]^{d_{B/A}} \ar[r]^u & C \ar[u]_{d_{C/A}}
}
\tag{20.5.2.6}
\label{0.20.5.2.6-fr}
\]
交换。

最后，若 $w:C\to D$ 是第二个 $A$-代数同态，则有传递性
\begin{equation}
\tag{20.5.2.7}
\label{0.20.5.2.7-fr}
  (w\circ u)_{D/B/A}
  =w_{D/C/A}\circ(u_{C/B/A}\otimes\mathrm{I}_D),
\end{equation}
这是从定义得出的。
\end{env}

\begin{env}[20.5.3]
\label{0.20.5.3-fr}
现在设 $A$、$B$ 是两个环，$v:A\to B$ 是环同态，$C$ 是一个
$B$-代数，并成为一个 $A$-代数，这是借助 $v$ 得到的；则典范映射
$v_0:C\otimes_A C\to C\otimes_B C$ 是代数的满\emph{双同态}
（相对于 $v:A\to B$），且图
\[
\xymatrix{
  C\otimes_A C \ar[r]^{v_0} \ar[d]_{p_{C/A}} &
  C\otimes_B C \ar[d]^{p_{C/B}} \\
  C \ar[r]^{\mathrm{I}_C} & C
}
\tag{20.5.3.1}
\label{0.20.5.3.1-fr}
\]
交换；取商即得到代数的一个双同态
\begin{equation}
\tag{20.5.3.2}
\label{0.20.5.3.2-fr}
  v':\mathrm{P}^1_{C/A}\to\mathrm{P}^1_{C/B}
\end{equation}
使得图
\[
\xymatrix{
  \mathrm{P}^1_{C/A} \ar[r]^{v'} & \mathrm{P}^1_{C/B} \\
  C \ar[u]^{p_1} \ar[r]^{\mathrm{I}_C} & C \ar[u]_{p_1}
}
\]
交换。因为 $v_0$ 把 $\mathfrak{J}_{C/A}$ 映入
$\mathfrak{J}_{C/B}$，将 $v'$ 限制于 $\Omega^1_{C/A}$，就得到映射
\begin{equation}
\tag{20.5.3.3}
\label{0.20.5.3.3-fr}
  v_{C/B/A}:\Omega^1_{C/A}\to\Omega^1_{C/B},
\end{equation}
它是 $C$-模同态。

\oldpage[0\textsubscript{IV-1}]{130}
此外，因为图
\[
\xymatrix{
  \mathrm{P}^1_{C/A} \ar[r]^{v'} & \mathrm{P}^1_{C/B} \\
  C \ar[u]^{p_2} \ar[r]^{\mathrm{I}_C} & C \ar[u]_{p_2}
}
\tag{20.5.3.4}
\label{0.20.5.3.4-fr}
\]
也交换，遂知图
\[
\xymatrix{
  \Omega^1_{C/A} \ar[r]^{v_{C/B/A}} & \Omega^1_{C/B} \\
  C \ar[u]^{d_{C/A}} \ar[r]^{\mathrm{I}_C} & C \ar[u]_{d_{C/B}}
}
\tag{20.5.3.5}
\label{0.20.5.3.5-fr}
\]
交换。

最后，若 $s:A'\to A$ 是第二个环同态，则有传递性
\begin{equation}
\tag{20.5.3.6}
\label{0.20.5.3.6-fr}
  (v\circ s)_{C/B/A'}=v_{C/B/A}\circ s_{C/A/A'}.
\end{equation}
\end{env}

\begin{env}[20.5.4]
\label{0.20.5.4-fr}
现在若有环同态的交换图
\[
\xymatrix{
  B \ar[r]^v & B' \\
  A \ar[u]^f \ar[r]^u & A' \ar[u]_g
}
\]
则由 \sref{0.20.5.2.4-fr} 和 \sref{0.20.5.3.3-fr} 复合得到一个
$B'$-模同态
\begin{equation}
\tag{20.5.4.1}
\label{0.20.5.4.1-fr}
  \Omega^1_{B/A}\otimes_B B'
  \xrightarrow{v_{B'/B/A}}\Omega^1_{B'/A}
  \xrightarrow{u_{B'/A'/A}}\Omega^1_{B'/A'}
\end{equation}
使得 $A'$-同态图
\[
\xymatrix{
  \Omega^1_{B/A}\otimes_B B' \ar[r] & \Omega^1_{B'/A'} \\
  B' \ar[u]^{d_{B/A}\otimes1} \ar[r]^{\mathrm{I}_{B'}} &
  B' \ar[u]_{d_{B'/A'}}
}
\tag{20.5.4.2}
\label{0.20.5.4.2-fr}
\]
交换。

此外，同态 \sref{0.20.5.4.1-fr} 对应于一个 $B$-模双同态
\begin{equation}
\tag{20.5.4.3}
\label{0.20.5.4.3-fr}
  \Omega^1_{B/A}\to\Omega^1_{B'/A'}.
\end{equation}
\end{env}

\begin{proposition}[20.5.5]
\label{0.20.5.5-fr}
若 $A'$、$B$ 是两个 $A$-代数，并且 $B'=B\otimes_A A'$，则典范同态
\sref{0.20.5.4.1-fr}
\begin{equation}
\tag{20.5.5.1}
\label{0.20.5.5.1-fr}
  \Omega^1_{B/A}\otimes_B B'\to\Omega^1_{B'/A'}
\end{equation}
是双射。
\end{proposition}

\begin{proof}
\oldpage[0\textsubscript{IV-1}]{131}
此时，\sref{0.20.5.5.1-fr} 的左端正是
$\Omega^1_{B/A}\otimes_A A'$。在相差一个典范同构的意义下，可以写成
$B'\otimes_{A'}B'=(B\otimes_A B)\otimes_A A'$，而 $p_{B'/A'}$ 此时等同于
$p_{B/A}\otimes\mathrm{I}_{A'}$；因而（由于 $p_{B/A}$ 是满射）
$\mathfrak{J}_{B'/A'}=\operatorname{Im}(\mathfrak{J}_{B/A}\otimes_A A')$，并且
$\mathfrak{J}^2_{B'/A'}
=\operatorname{Im}(\mathfrak{J}^2_{B/A}\otimes_A A')$，由此在相差一个典范同构
的意义下得到 $\mathrm{P}^1_{B'/A'}=\mathrm{P}^1_{B/A}\otimes_A A'$；该同构
把增广理想分别变为其自身。由于 $A$-模 $\mathrm{P}^1_{B/A}$ 典范地等同于
$B$ 与 $\Omega^1_{B/A}$ 的直和，所以在同一个同构下有
\begin{equation}
\tag{20.5.5.2}
\label{0.20.5.5.2-fr}
  \Omega^1_{B/A}\otimes_A A'=\Omega^1_{B'/A'}
\end{equation}
并且立即验证，该同构与典范同构
$\Omega^1_{B/A}\otimes_B B'\xrightarrow{\sim}
\Omega^1_{B/A}\otimes_A A'$ 的复合正是 \sref{0.20.5.5.1-fr}。
\end{proof}

\begin{env}[20.5.6]
\label{0.20.5.6-fr}
由函子性，典范同态 \sref{0.20.5.2.4-fr} 和
\sref{0.20.5.3.3-fr} 对每个 $C$-模 $L$ 给出典范同态
\begin{align}
\tag{20.5.6.1}
\label{0.20.5.6.1-fr}
  &\operatorname{Hom}_C(\Omega^1_{C/A},L)
  \to\operatorname{Hom}_C(\Omega^1_{B/A}\otimes_B C,L)
  =\operatorname{Hom}_B(\Omega^1_{B/A},L) \\
\tag{20.5.6.2}
\label{0.20.5.6.2-fr}
  &\operatorname{Hom}_C(\Omega^1_{C/B},L)
  \to\operatorname{Hom}_C(\Omega^1_{C/A},L).
\end{align}
考虑到 \sref{0.20.4.8.2-fr} 以及交换图 \sref{0.20.5.2.6-fr} 和
\sref{0.20.5.3.5-fr}，这些同态在相差一个典范等同的意义下，分别正是同态
\sref{0.20.2.1.2-fr} 和 \sref{0.20.2.1.3-fr}。
\end{env}

\begin{theorem}[20.5.7]
\label{0.20.5.7-fr}
设 $u:A\to B$、$v:B\to C$ 是两个环同态。
\begin{enumerate}
  \item[\rm(i)] $C$-模序列
\begin{equation}
\tag{20.5.7.1}
\label{0.20.5.7.1-fr}
  \Omega^1_{B/A}\otimes_B C
  \xrightarrow{v_{C/B/A}}\Omega^1_{C/A}
  \xrightarrow{u_{C/B/A}}\Omega^1_{C/B}\to0
\end{equation}
是正合的。
  \item[\rm(ii)] $v_{C/B/A}$ 有左逆，当且仅当 $C$ 作为一个 $B$-代数，
  对于 $A$ 是形式平滑的（对于离散拓扑（参见 \sref{0.19.9.1-fr}））；特别地，
  为此只须 $C$ 是一个形式平滑的 $B$-代数（对于离散拓扑）。
\end{enumerate}
\end{theorem}

\begin{proof}
(i) 考虑到 \sref{0.20.5.6-fr}，序列 \sref{0.20.2.4.2-fr} 的正合性首先表明，
对每个 $C$-模 $L$，序列
\[
  0\to\operatorname{Hom}_C(\Omega^1_{C/B},L)
  \to\operatorname{Hom}_C(\Omega^1_{C/A},L)
  \to\operatorname{Hom}_C(\Omega^1_{B/A}\otimes_B C,L)
\]
是正合的。已知这蕴含序列 \sref{0.20.5.7.1-fr} 的正合性（Bourbaki，
\emph{Alg.}，第 II 章，第 3 版，§~2，n\textsuperscript{o}~1，定理~1）。

(ii) 根据 \sref{0.20.5.7.1-fr} 的正合性，说 $v_{C/B/A}$ 有左逆，意即序列
\begin{equation}
\tag{20.5.7.2}
\label{0.20.5.7.2-fr}
  0\to\Omega^1_{B/A}\otimes_B C
  \xrightarrow{v_{C/B/A}}\Omega^1_{C/A}
  \xrightarrow{u_{C/B/A}}\Omega^1_{C/B}\to0
\end{equation}
是\emph{正合且分裂的}；已知（Bourbaki，\emph{同上}，
n\textsuperscript{o}~1，命题~1），这等价于：对每个 $C$-模 $L$，序列
\[
  0\to\operatorname{Hom}_C(\Omega^1_{C/B},L)
  \to\operatorname{Hom}_C(\Omega^1_{C/A},L)
  \to\operatorname{Hom}_C(\Omega^1_{B/A}\otimes_B C,L)\to0
\]
是正合的；考虑到 \sref{0.20.5.6-fr} 和 \sref{0.20.2.4.2-fr}，该条件等价于
\[
  \operatorname{Exalcom}_{B/A}(C,L)=0
\]
对每个 $C$-模 $L$ 成立，故结论由 \sref{0.19.9.8.1-fr} 得出。
\end{proof}

\oldpage[0\textsubscript{IV-1}]{132}
此外应当注意，若有环同态的交换图
\[
\xymatrix{
  A' \ar[r] & B' \ar[r] & C' \\
  A \ar[u] \ar[r] & B \ar[u] \ar[r] & C \ar[u]
}
\]
则由此得到交换图
\[
\xymatrix{
  \Omega^1_{B/A}\otimes_B C \ar[r] \ar[d] &
  \Omega^1_{C/A} \ar[r] \ar[d] &
  \Omega^1_{C/B} \ar[r] \ar[d] & 0 \\
  \Omega^1_{B'/A'}\otimes_{B'} C' \ar[r] &
  \Omega^1_{C'/A'} \ar[r] &
  \Omega^1_{C'/B'} \ar[r] & 0
}
\tag{20.5.7.3}
\label{0.20.5.7.3-fr}
\]
其中各竖直箭头来自双同态 \sref{0.20.5.4.3-fr}。

\begin{corollary}[20.5.8]
\label{0.20.5.8-fr}
假定同态 $v:B\to C$ 使 $C$ 成为一个形式平展的 $B$-代数（对于离散拓扑，
\sref{0.19.10.2-fr}）；则同态 \sref{0.20.5.3.3-fr}
\[
  v_{C/B/A}:\Omega^1_{B/A}\otimes_B C\to\Omega^1_{C/A}
\]
是双射。
\end{corollary}

\begin{proof}
事实上，若 $C$ 是一个形式非分歧的 $B$-代数（对于离散拓扑），则由
\sref{0.19.10.4-fr}、\sref{0.20.4.8-fr} 和 \sref{0.20.1.1-fr} 得到
$\operatorname{Hom}_C(\Omega^1_{C/B},L)=0$；这里 $C$-模 $L$ 可以任取。从而
$\Omega^1_{C/B}=0$（参见 (20.7.4)）；另一方面，若 $C$ 是一个形式平滑的
$B$-代数（对于离散拓扑），则序列 \sref{0.20.5.7.2-fr} 是正合的；推论即得。
\end{proof}

\begin{corollary}[20.5.9]
\label{0.20.5.9-fr}
设 $A$ 是环，$B$ 是 $A$-代数，$S$ 是 $B$ 的一个乘法子集；则典范同态
\begin{equation}
\tag{20.5.9.1}
\label{0.20.5.9.1-fr}
  S^{-1}\Omega^1_{B/A}\to\Omega^1_{S^{-1}B/A}
\end{equation}
是双射。
\end{corollary}

\begin{proof}
只须对 $C=S^{-1}B$ 应用 \sref{0.20.5.8-fr}；它是一个形式平展的
$B$-代数（对于离散拓扑，\sref{0.19.10.3-fr}，(ii)）。
\end{proof}

考虑到 \sref{0.20.5.5-fr}，因而在相差典范同构的意义下可以写成
\begin{equation}
\tag{20.5.9.2}
\label{0.20.5.9.2-fr}
  \Omega^1_{S^{-1}B/S^{-1}A}
  =S^{-1}\Omega^1_{B/A}
  =\Omega^1_{S^{-1}B/A},
\end{equation}

\begin{corollary}[20.5.10]
\label{0.20.5.10-fr}
若 $k$ 是域，且 $K=k(X_\alpha)_{\alpha\in I}$ 是 $k$ 的纯扩张，则各
$dX_\alpha$ 构成 $K$-向量空间 $\Omega^1_{K/k}$ 的一组基。
\end{corollary}

\begin{proof}
因为 $K$ 是多项式环 $k[X_\alpha]_{\alpha\in I}$ 的分式域，结论由
\sref{0.20.4.13-fr}，(i) 和 \sref{0.20.5.9-fr} 得出。
\end{proof}

\oldpage[0\textsubscript{IV-1}]{133}
\begin{env}[20.5.11]
\label{0.20.5.11-fr}
设 $A$ 是环，$B$ 是 $A$-代数，$\mathfrak{R}$ 是 $B$ 的理想，$C$ 是商
$A$-代数 $B/\mathfrak{R}$；考虑 $A$-模的复合同态
\begin{equation}
\tag{20.5.11.1}
\label{0.20.5.11.1-fr}
  \delta:\mathfrak{R}\longrightarrow B\xrightarrow{d_{B/A}}\Omega^1_{B/A},
\end{equation}
其中第一支箭头是典范单射；由于 $d(xy)=xdy+ydx$，可见
$\delta(\mathfrak{R}^2)\subset\mathfrak{R}.\Omega^1_{B/A}$，故取商得到一个
$A$-模同态
\begin{equation}
\tag{20.5.11.2}
\label{0.20.5.11.2-fr}
  \delta_{C/B/A}:\mathfrak{R}/\mathfrak{R}^2\longrightarrow
  \Omega^1_{B/A}\otimes_B C=
  \Omega^1_{B/A}/(\mathfrak{R}.\Omega^1_{B/A}).
\end{equation}
但事实上，$\delta_{C/B/A}$ 是 $C$-模同态，因为对 $x\in B$ 和
$y\in\mathfrak{R}$，有 $ydx\in\mathfrak{R}.\Omega^1_{B/A}$，从而
$d(xy)\equiv xdy\pmod{\mathfrak{R}.\Omega^1_{B/A}}$；这首先证明
(20.5.11.2) 是 $B$-模同态，而由于 $\mathfrak{R}$ 零化两端，断言得证。

若 $B'$ 是第二个 $A$-代数，$u:B\to B'$ 是一个 $A$-同态，
$\mathfrak{R}'$ 是 $B'$ 的理想并满足
$u(\mathfrak{R})\subset\mathfrak{R}'$，而 $C'=B'/\mathfrak{R}'$ 是商代数，
则有交换图
\begin{equation}
\tag{20.5.11.3}
\label{0.20.5.11.3-fr}
\xymatrix{
  \mathfrak{R}/\mathfrak{R}^2 \ar[r]^-{\delta_{C/B/A}} \ar[d] &
  \Omega^1_{B/A}\otimes_B C \ar[d] \\
  \mathfrak{R}'/\mathfrak{R}'^2 \ar[r]_-{\delta_{C'/B'/A}} &
  \Omega^1_{B'/A}\otimes_{B'} C'
}
\end{equation}
其中各竖直箭头来自 $u$（20.5.2.4）。
\end{env}

\begin{theorem}[20.5.12]
\label{0.20.5.12-fr}
设 $A$ 是环，$B$ 是 $A$-代数，$C$ 是商 $A$-代数
$B/\mathfrak{R}$，$u:B\to C$ 是典范同态。
\begin{enumerate}
  \item[\rm(i)] 有一个 $C$-模正合列
\begin{equation}
\tag{20.5.12.1}
\label{0.20.5.12.1-fr}
  \mathfrak{R}/\mathfrak{R}^2
  \xrightarrow{\delta_{C/B/A}}\Omega^1_{B/A}\otimes_B C
  \xrightarrow{u_{C/B/A}}\Omega^1_{C/A}\longrightarrow 0,
\end{equation}
其中 $u_{C/B/A}$ 和 $\delta_{C/B/A}$ 分别由 (20.5.2.4) 和
(20.5.11.2) 定义。
  \item[\rm(ii)] 若置 $E=B/\mathfrak{R}^2$，则典范同态 (20.5.2.4)
\[
  \Omega^1_{B/A}\otimes_B C\longrightarrow
  \Omega^1_{E/A}\otimes_E C
\]
是双射。
  \item[\rm(iii)] 下列三个条件等价：
  \begin{enumerate}
    \item[\rm a)] $\delta_{C/B/A}$ 有左逆。
    \item[\rm b)] 每个 $A$-扩充，即由 $C$ 的某个 $C$-模 $L$ 对它所作的扩充，若其通过
    $u:B\to C$ 的逆像是 $A$-平凡的，则它自身也是 $A$-平凡的。
    \item[\rm c)] $A$-代数 $E=B/\mathfrak{R}^2$ 是一个 $A$-平凡扩充；这是
    对 $C$ 所作的，其扩充模为 $\mathfrak{R}/\mathfrak{R}^2$。
  \end{enumerate}
  \item[\rm(iv)] $\delta_{C/B/A}$ 的各左逆与典范同态 $E\to C$ 的各右逆之间，
  有一个典范双射对应。
\end{enumerate}
\end{theorem}

\begin{proof}
\oldpage[0\textsubscript{IV-1}]{134}
(i) 因为 $u$ 是满射，所以根据 (20.1.1) 有
$\operatorname{D\acute{e}r}_B(C,L)=0$；这里 $C$-模 $L$ 可以任取。因此正合列
(20.2.3.1) 变为
\begin{equation}
\tag{20.5.12.2}
\label{0.20.5.12.2-fr}
  0\longrightarrow\operatorname{D\acute{e}r}_A(C,L)\xrightarrow{v^0}\operatorname{D\acute{e}r}_A(B,L)
  \xrightarrow{\partial}\operatorname{Exalcom}_B(C,L)
  \xrightarrow{v^1}\operatorname{Exalcom}_A(C,L)
  \xrightarrow{u^1}\operatorname{Exalcom}_A(B,L),
\end{equation}
其中 $v$ 是同态 $A\to B$。另一方面，回忆 (18.3.8)，
$\operatorname{Exalcom}_B(C,L)$ 典范地等同于
$\operatorname{Hom}_C(\mathfrak{R}/\mathfrak{R}^2,L)$；因而由 (20.5.12.2)
和 (20.4.8) 得到正合列
\begin{equation}
\tag{20.5.12.3}
\label{0.20.5.12.3-fr}
\begin{split}
  0\longrightarrow\operatorname{Hom}_C(\Omega^1_{C/A},L)
  \longrightarrow\operatorname{Hom}_C(\Omega^1_{B/A}\otimes_B C,L)
  \xrightarrow{\varphi}\operatorname{Hom}_C(\mathfrak{R}/\mathfrak{R}^2,L)\\
  \xrightarrow{\psi}\operatorname{Ker}(u^1)\longrightarrow 0
\end{split}
\end{equation}
其中 $\varphi=\partial\circ\eta^{-1}$ 且 $\psi=v^1\circ\eta$。回到
$\partial$（20.2.2）和 $\eta$（18.3.8）的定义，立即看出 $\varphi$ 正是同态
$\operatorname{Hom}(\delta_{C/B/A},1_L)$。由 (20.5.12.3) 的前四项组成的正合列
存在，因而表明序列 (20.5.12.1) 是正合的（Bourbaki，\emph{Alg.}，第 II 章，
第 3 版，\S~2，n\textsuperscript{o}~1，定理~1）。

\medskip
(ii) 把正合列 (20.5.12.1) 应用于 $B$ 和理想 $\mathfrak{R}^2$，得到
\begin{equation}
\tag{20.5.12.4}
\label{0.20.5.12.4-fr}
  \mathfrak{R}^2/\mathfrak{R}^4\longrightarrow
  \Omega^1_{B/A}\otimes_B E\longrightarrow\Omega^1_{E/A}\longrightarrow 0,
\end{equation}
由此与 $C$（视为 $E$-代数）张量，得到正合列
\[
  \mathfrak{R}^2/\mathfrak{R}^4\otimes_E C\longrightarrow
  \Omega^1_{B/A}\otimes_B C\longrightarrow
  \Omega^1_{E/A}\otimes_E C\longrightarrow 0.
\]
现在，若 $x,y$ 是 $\mathfrak{R}$ 的两个元素，而 $\zeta$ 是 $xy$ 模
$\mathfrak{R}^4$ 的类，则像 $\delta'(\zeta\otimes1)$ 依定义为
$d_{B/A}(xy)\otimes1=(xd_{B/A}(y)+yd_{B/A}(x))\otimes1$；但因为 $x$、$y$
在 $C$ 中的像都为零，所以也有 $d_{B/A}(xy)\otimes1=0$；这个等式成立于
$\Omega^1_{B/A}\otimes_B C$ 中。这就证明了断言。

\medskip
(iii) 考虑到 (20.5.12.1) 的正合性，说 $\delta_{C/B/A}$ 有左逆，等价于说序列
\begin{equation}
\tag{20.5.12.5}
\label{0.20.5.12.5-fr}
  0\longrightarrow\mathfrak{R}/\mathfrak{R}^2
  \xrightarrow{\delta_{C/B/A}}\Omega^1_{B/A}\otimes_B C
  \xrightarrow{u_{C/B/A}}\Omega^1_{C/A}\longrightarrow0
\end{equation}
是\emph{正合且分裂的}；这又等价于（Bourbaki，\emph{同上}）说：在正合列
(20.5.12.3) 中有 $\operatorname{Ker}(u^1)=0$，且 $L$ 可以任取。这就证明了条件
a) 与 b) 等价（参见 (18.3.6.2)）。

b) 蕴含 c)，源于如下事实：通过 $u:B\to C$ 取逆像时，$A$-扩充 $E$
（这是对 $C$ 所作的、扩充模为 $\mathfrak{R}/\mathfrak{R}^2$ 的扩充）是
$A$-平凡的；此处 $u$ 是典范同态 $B\to E\to C$ 的复合（18.1.6）。反过来，
c) 蕴含 b)，因为每个 $B$-扩充，即对 $C$ 所作且扩充模为 $L$ 的扩充，都在
$B$-意义下等价于 $E\oplus_{\mathfrak{R}/\mathfrak{R}^2}L$（18.3.8）；换言之，
它的类是 $E$ 的类（视为 $B$-扩充）在同态

\oldpage[0\textsubscript{IV-1}]{135}
$w_*:\operatorname{Exan}_B(C,\mathfrak{R}/\mathfrak{R}^2)
\to\operatorname{Exan}_B(C,L)$ 下的像；该同态对应于一个 $B$-同态
$w:\mathfrak{R}/\mathfrak{R}^2\to L$。于是 c) 蕴含 b) 由图 (18.3.6.5) 的
交换性得出：
\[
\xymatrix{
  \operatorname{Exan}_B(C,\mathfrak{R}/\mathfrak{R}^2)
    \ar[r]^{w_*} \ar[d]_{u^*} &
  \operatorname{Exan}_B(C,L) \ar[d]^{u^*} \\
  \operatorname{Exan}_A(C,\mathfrak{R}/\mathfrak{R}^2)
    \ar[r]_{w_*} & \operatorname{Exan}_A(C,L).
}
\]

\medskip
(iv) 我们已经看到（20.1.7），$E\to C$ 的各右逆与满足下述条件的元素所成的集合
$D\in\operatorname{D\acute{e}r}_A(E,\mathfrak{R}/\mathfrak{R}^2)$ 典范一一对应：有
$D(x)=x$，这个等式成立于 $\mathfrak{R}/\mathfrak{R}^2$ 中；因此根据
(20.4.8)，它们也与满足
下述条件的 $E$-同态所成的集合一一对应：
$h:\Omega^1_{E/A}\to\mathfrak{R}/\mathfrak{R}^2$，且复合映射
\[
  \mathfrak{R}/\mathfrak{R}^2\xrightarrow{d_{E/A}}
  \Omega^1_{E/A}\xrightarrow{h}\mathfrak{R}/\mathfrak{R}^2
\]
是恒等映射。与 $C$ 张量后，由此得出（因为
$\mathfrak{R}/\mathfrak{R}^2$ 是 $C$-模）复合映射
\[
  \mathfrak{R}/\mathfrak{R}^2\xrightarrow{\delta_{C/E/A}}
  \Omega^1_{E/A}\otimes_E C\xrightarrow{h\otimes1}
  \mathfrak{R}/\mathfrak{R}^2
\]
是恒等映射；现在，因为 $\mathfrak{R}/\mathfrak{R}^2$ 是 $C$-模，
$h\mapsto h\otimes1$ 是从集合
$\operatorname{Hom}_E(\Omega^1_{E/A},\mathfrak{R}/\mathfrak{R}^2)$ 到
$\operatorname{Hom}_C(\Omega^1_{E/A}\otimes_E C,\mathfrak{R}/\mathfrak{R}^2)$ 的
同构；另一方面，(ii) 证明可以典范地等同
$\Omega^1_{E/A}\otimes_E C$ 与 $\Omega^1_{B/A}\otimes_B C$，而
$\delta_{C/E/A}$ 此时等同于 $\delta_{C/B/A}$。证毕。
\end{proof}

\begin{example}[20.5.13]
\label{0.20.5.13-fr}
设 $B=A[X_\alpha]_{\alpha\in I}$ 是 $A$ 上的多项式代数，
$\mathfrak{R}$ 是 $B$ 的理想，$(P_\lambda)$ 是 $\mathfrak{R}$ 的一组生成元，
并且 $C=B/\mathfrak{R}$；已知 $\Omega^1_{B/A}$ 是自由 $B$-模，其一组基为
各 $dX_\alpha$（20.4.13，(i)），所以各 $dX_\alpha$ 也构成自由 $C$-模
$\Omega^1_{B/A}\otimes_B C$ 的一组基。另一方面，由定义立即可知，
$\mathfrak{R}/\mathfrak{R}^2$ 在 $\delta_{C/B/A}$ 下的像，是由下列各元素生成的
$C$-子模：
\[
  dP_\lambda=\sum_\alpha
  \frac{\partial P_\lambda}{\partial X_\alpha}\,dX_\alpha.
\]
由此得到，$\Omega^1_{C/A}$ 同构于如下商模：自由 $C$-模（其一组基为各
$dX_\alpha$），模掉 $C$-子模（由各 $dP_\lambda$ 生成）。这就给出了
\emph{任意}代数的
微分模的一种描述，因为每个 $A$-代数 $C$ 都可以由上述方式得到。
\end{example}

\begin{corollary}[20.5.14]
\label{0.20.5.14-fr}
若 $C$ 是一个形式平滑的 $A$-代数（对于离散拓扑），则序列
\begin{equation}
\tag{20.5.14.1}
\label{0.20.5.14.1-fr}
  0\longrightarrow\mathfrak{R}/\mathfrak{R}^2
  \xrightarrow{\delta_{C/B/A}}\Omega^1_{B/A}\otimes_B C
  \xrightarrow{u_{C/B/A}}\Omega^1_{C/A}\longrightarrow0
\end{equation}
是正合且分裂的。
\end{corollary}

\begin{proof}
事实上，此时每个 $A$-扩充，即由 $C$ 的一个 $C$-模对它所作的扩充，都是平凡的
（19.4.4.1）。
\end{proof}

\oldpage[0\textsubscript{IV-1}]{136}
\begin{remark}[20.5.15]
\label{0.20.5.15-fr}
设 $u:A\to B$ 是环的\emph{满}同态；那么，对每个环同态
$v:B\to C$，典范同态
\begin{equation}
\tag{20.5.15.1}
\label{0.20.5.15.1-fr}
  u_{C/B/A}:\Omega^1_{C/A}\longrightarrow\Omega^1_{C/B}
\end{equation}
是\emph{双射}；事实上，这由正合列 (20.5.7.1) 得出，因为
$\Omega^1_{B/A}=0$（20.4.12）。
\end{remark}

\subsection{不恰模与特征同态}
\label{subsection:0.20.6-fr}
\label{0.20.6-fr}

\begin{definition}[20.6.1]
\label{0.20.6.1-fr}
给定两个环同态 $u:A\to B$、$v:B\to C$，把\emph{$B$-代数 $C$ 相对于
$A$ 的不恰模}记为 $\Upsilon_{C/B/A}$；它是 $C$-模，即同态
$v_{C/B/A}:\Omega^1_{B/A}\otimes_B C\to\Omega^1_{C/A}$ 的核。因而依定义
（参见 (20.5.7)），有正合列
\begin{equation}
\tag{20.6.1.1}
\label{0.20.6.1.1-fr}
  0\longrightarrow\Upsilon_{C/B/A}\longrightarrow
  \Omega^1_{B/A}\otimes_B C\xrightarrow{v_{C/B/A}}
  \Omega^1_{C/A}\xrightarrow{u_{C/B/A}}\Omega^1_{C/B}
  \longrightarrow0.
\end{equation}
\end{definition}

当 $A=\mathbb{Z}$ 时（因而模 $\Omega^1_{B/A}$ 和 $\Omega^1_{C/A}$ 是“绝对”
微分模 $\Omega^1_B$ 和 $\Omega^1_C$），用 $\Upsilon_{C/B}$ 代替
$\Upsilon_{C/B/\mathbb{Z}}$；当 $B$、$C$ 是素域 $P$ 上的代数时，有
$\Upsilon_{C/B/P}=\Upsilon_{C/B}$。

设 $R$、$S$ 分别是 $B$、$C$ 的乘法子集，并且 $R$ 的像包含于 $S$。于是由
正合列 (20.6.1.1) 和 (20.5.9) 得到
\begin{equation}
\tag{20.6.1.2}
\label{0.20.6.1.2-fr}
  \Upsilon_{S^{-1}C/R^{-1}B/A}=S^{-1}\Upsilon_{C/B/A}.
\end{equation}

\begin{proposition}[20.6.2]
\label{0.20.6.2-fr}
若 $C$ 作为 $B$-代数对于 $A$ 是形式平滑的（特别地，若 $C$ 是形式平滑的
$B$-代数），则有 $\Upsilon_{C/B/A}=0$。
\end{proposition}

\begin{proof}
这由 (20.5.7，(ii)) 得出。
\end{proof}

\begin{proposition}[20.6.3]
\label{0.20.6.3-fr}
设 $k$ 是域，$K$ 是 $k$ 的扩张。$K$ 是 $k$ 的可分扩张，当且仅当
$\Upsilon_{K/k}=0$（换言之，典范同态
$\Omega^1_k\otimes_k K\to\Omega^1_K$ 是单射；或者仍根据 (20.4.8)，从 $k$
到 $K$ 的每个导射都能延拓为从 $K$ 到自身的导射）。
\end{proposition}

\begin{proof}
事实上，说 $K$ 在 $k$ 上可分，等价于说它是形式平滑的 $k$-代数
（19.6.1）；另一方面，若 $P$ 是 $k$ 的素域，则 $K$ 在 $P$ 上可分，因而是
形式平滑的 $P$-代数，所以，说 $K$ 是形式平滑的 $k$-代数，等价于说它是
一个形式平滑的 $k$-代数且相对于 $P$。最后，根据 (20.5.7，(ii))，说 $K$
是一个形式平滑的 $k$-代数且相对于 $P$，等价于说同态
$v_{K/k/P}:\Omega^1_k\otimes_k K\to\Omega^1_K$ 有左逆；但因为 $K$ 是域，
后一条件等价于说 $v_{K/k/P}$ 的核，即 $\Upsilon_{K/k}$，为零。
\end{proof}

\begin{env}[20.6.4]
\label{0.20.6.4-fr}
考虑交换图
\begin{equation}
\tag{20.6.4.1}
\label{0.20.6.4.1-fr}
\xymatrix{
  A' \ar[r]^{u'} & B' \ar[r]^{v'} & C' \\
  A \ar[u]^f \ar[r]_u & B \ar[u]^g \ar[r]_v & C \ar[u]^h
}
\end{equation}

\oldpage[0\textsubscript{IV-1}]{137}
其中各环均交换。相应图 (20.5.7.3) 的交换性蕴含存在唯一的 $C$-同态
\begin{equation}
\tag{20.6.4.2}
\label{0.20.6.4.2-fr}
  \Upsilon_{C/B/A}\longrightarrow\Upsilon_{C'/B'/A'}
\end{equation}
它从 (20.6.4.1) 典范地导出，并使图
\begin{equation}
\tag{20.6.4.3}
\label{0.20.6.4.3-fr}
\xymatrix{
  0 \ar[r] & \Upsilon_{C/B/A} \ar[r] \ar[d] &
  \Omega^1_{B/A}\otimes_B C \ar[r]^{v_{C/B/A}} \ar[d] &
  \Omega^1_{C/A} \ar[r]^{u_{C/B/A}} \ar[d] &
  \Omega^1_{C/B} \ar[r] \ar[d] & 0 \\
  0 \ar[r] & \Upsilon_{C'/B'/A'} \ar[r] &
  \Omega^1_{B'/A'}\otimes_{B'} C' \ar[r]_{v_{C'/B'/A'}} &
  \Omega^1_{C'/A'} \ar[r]_{u_{C'/B'/A'}} &
  \Omega^1_{C'/B'} \ar[r] & 0
}
\end{equation}
交换。

此外，给定同态 (20.6.4.2)，等价于给定一个 $C'$-同态
\begin{equation}
\tag{20.6.4.4}
\label{0.20.6.4.4-fr}
  \Upsilon_{C/B/A}\otimes_C C'\longrightarrow\Upsilon_{C'/B'/A'}
\end{equation}
它与典范同态 $\Upsilon_{C/B/A}\to\Upsilon_{C/B/A}\otimes_C C'$ 复合后，
重新给出 (20.6.4.2)。显然，(20.6.4.2) 具有自然的传递性；由此可以说
$\Upsilon_{C/B/A}$ 是三元组 $(A,B,C)$ 的一个\emph{函子}。
\end{env}

\begin{env}[20.6.5]
\label{0.20.6.5-fr}
在 (20.6.1) 的条件下，为便于后文，我们引入一个 $C$-模的\emph{（链）复形}
$\mathrm{K}_\bullet(C/B/A)$；它除次数 0、1 外各项均为零，而在这两个次数取
\begin{equation}
\tag{20.6.5.1}
\label{0.20.6.5.1-fr}
\left\{
\begin{aligned}
  \mathrm{K}_0(C/B/A)&=\Omega^1_{C/A},\\
  \mathrm{K}_1(C/B/A)&=\Omega^1_{B/A}\otimes_B C,
\end{aligned}
\right.
\end{equation}
而微分算子 $\mathrm{K}_1\to\mathrm{K}_0$ 为 $v_{C/B/A}$。这样就可以在相差
典范同构的意义下，把 $\Omega^1_{C/B}$ 和 $\Upsilon_{C/B/A}$ 写成该复形的
\emph{同调模}：
\begin{equation}
\tag{20.6.5.2}
\label{0.20.6.5.2-fr}
  \mathrm{H}_0(\mathrm{K}_\bullet(C/B/A))=\Omega^1_{C/B},\qquad
  \mathrm{H}_1(\mathrm{K}_\bullet(C/B/A))=\Upsilon_{C/B/A}.
\end{equation}
\end{env}

同样有：
\begin{proposition}[20.6.6]
\label{0.20.6.6-fr}
在 (20.6.1) 的假设下，对每个 $C$-模 $L$，有典范 $C$-同构
\begin{align}
\tag{20.6.6.1}
\label{0.20.6.6.1-fr}
  \mathrm{H}^0(\mathrm{K}_\bullet(C/B/A),L)&\simeq\operatorname{D\acute{e}r}_B(C,L),\\
\tag{20.6.6.2}
\label{0.20.6.6.2-fr}
  \mathrm{H}^1(\mathrm{K}_\bullet(C/B/A),L)&\simeq
  \operatorname{Exalcom}_{B/A}(C,L).
\end{align}
\end{proposition}

\begin{proof}
事实上，根据 (20.4.8) 和 (20.5.6)，上链复形
$\operatorname{Hom}_C(\mathrm{K}_\bullet(C/B/A),L)$ 正是复形
\[
  \cdots\longrightarrow0\longrightarrow\operatorname{D\acute{e}r}_A(C,L)
  \longrightarrow\operatorname{D\acute{e}r}_A(B,L)\longrightarrow0\longrightarrow\cdots
\]
\oldpage[0\textsubscript{IV-1}]{138}
其中微分算子为 $v^0$（采用 (20.2.1) 的记号）。于是命题由正合列
(20.2.4.2) 和上同调模的定义得出：
\begin{equation}
\tag{20.6.6.3}
\label{0.20.6.6.3-fr}
  \mathrm{H}^i(\mathrm{K}_\bullet,L)=
  \mathrm{H}^i(\operatorname{Hom}_C(\mathrm{K}_\bullet,L)).
\end{equation}
\end{proof}

若有环同态的交换图
\[
\xymatrix{
  A' \ar[r] & B' \ar[r] & C' \\
  A \ar[u] \ar[r] & B \ar[u] \ar[r] & C \ar[u]
}
\]
则双同态 (20.5.4.3) 定义一个模复形的双同态
\begin{equation}
\tag{20.6.6.4}
\label{0.20.6.6.4-fr}
  \mathrm{K}_\bullet(C/B/A)\longrightarrow
  \mathrm{K}_\bullet(C'/B'/A')
\end{equation}
而由它在同调或上同调上导出的双同态，通过公式 (20.6.5.2)、(20.6.6.1) 和
(20.6.6.2)，分别等同于已经在 (20.5.4.3)、(20.6.4.2)、(20.2.1) 和
(18.3.7.2) 中定义的双同态（最后一个还要考虑正合列 (20.2.4.2) 中算子
$\partial$ 的定义）。

\begin{env}[20.6.7]
\label{0.20.6.7-fr}
已知，对于一个复形 $\mathrm{K}_\bullet$（其各项为 $C$-模）以及一个
$C$-模 $L$，有典范同态
\[
  \alpha_i:\mathrm{H}^i(\mathrm{K}_\bullet,L)\longrightarrow
  \operatorname{Hom}_C(\mathrm{H}_i(\mathrm{K}_\bullet),L)
\]
（M，IV，6）。在这里，典范同态
\[
  \alpha_1:\mathrm{H}^1(\mathrm{K}_\bullet(C/B/A),L)
  \longrightarrow
  \operatorname{Hom}_C(\mathrm{H}_1(\mathrm{K}_\bullet(C/B/A)),L)
\]
立即定义为如下所得的同态：对限制同态
\[
  \operatorname{Hom}_C(\mathrm{K}_1(C/B/A),L)\longrightarrow
  \operatorname{Hom}_C(\mathrm{H}_1(\mathrm{K}_\bullet(C/B/A)),L)
\]
关于 $\operatorname{Hom}_C(\mathrm{K}_0,L)$ 的像取商；这是因为
$\mathrm{H}_1(\mathrm{K}_\bullet(C/B/A))$ 正是 $\mathrm{K}_1\to\mathrm{K}_0$
的核。考虑 (20.6.6.2) 和 (20.6.5.2)，因而得到一个典范 $C$-同态
\begin{equation}
\tag{20.6.7.1}
\label{0.20.6.7.1-fr}
  \operatorname{Exalcom}_{B/A}(C,L)\longrightarrow
  \operatorname{Hom}_C(\Upsilon_{C/B/A},L)
\end{equation}
它可以具体说明如下：根据 (20.2.4.2)，每个 $B$-扩充（它是对 $C$ 由
$L$ 作成且为 $A$-平凡的）都来自某个 $A$-导射 $D$，该导射把 $B$ 映入
$L$；因而根据 (20.4.8)，它来自某个 $C$-同态 $f$，该同态把
$\Omega^1_{B/A}\otimes_B C$ 映入 $L$。把这一扩充的类对应于 $f$ 在
$\Upsilon_{C/B/A}$ 上的
\emph{限制}；它事实上只依赖于这个类，而不依赖于 $D$ 的选择。
\end{env}

\begin{definition}[20.6.8]
\label{0.20.6.8-fr}
设 $u:A\to B$、$v:B\to C$ 是两个环同态。对于任一 $B$-扩充 $E$，它以
$C$ 为商、以某个 $C$-模 $L$ 为核，并且是 $A$-平凡的 (18.3.7)，称由
典范同态 (20.6.7.1) 所给出的、该扩充类 $E$ 的像，也就是记为 $\chi_E$ 的
$C$-同态 $\Upsilon_{C/B/A}\to L$，为 $E$ 的\emph{特征同态}。
\end{definition}

\oldpage[0\textsubscript{IV-1}]{139}
还可以用另一种方式定义同态 $\chi_E$：
\begin{proposition}[20.6.9]
\label{0.20.6.9-fr}
设 $E$ 是一个 $B$-扩充，它扩充 $C$，其核为一个 $C$-模 $L$，并且是
$A$-平凡的 (18.3.7)；那么图
\begin{equation}
\tag{20.6.9.1}
\label{0.20.6.9.1-fr}
\xymatrix{
  0 \ar[r] & \Upsilon_{C/B/A} \ar[r] \ar[d]^{\chi_E} &
  \Omega^1_{B/A}\otimes_B C \ar[r]^{v_{C/B/A}}
    \ar[d]^{q_{E/B/A}\otimes1_C} &
  \Omega^1_{C/A} \ar[d]^{1} \\
  0 \ar[r] & L \ar[r]_{\delta_{C/E/A}} &
  \Omega^1_{E/A}\otimes_E C \ar[r]_{p_{C/E/A}} &
  \Omega^1_{C/A} \ar[r] & 0
}
\end{equation}
是交换的，且两行均正合；这里 $q:B\to E$ 定义 $B$-扩充结构于 $E$ 上，
而 $p:E\to C$ 是增广同态。
\end{proposition}

\begin{proof}
图的下行是与两个同态 $A\xrightarrow{q\circ u}E$ 和
$p:E\to C=E/L$ 有关的序列 (20.5.12.5)；由于 $p$ 是满射且 $E$ 是
$A$-平凡扩充，该序列根据 (20.5.12, (iii)) 正合并且分裂。由关系
$v=p\circ q$ (20.5.2.7)，(20.6.9.1) 的右方格交换；于是
$q_{E/B/A}\otimes1_C$ 把核 $\Upsilon_{C/B/A}$（即 $v_{C/B/A}$ 的核）映入
核 $L$（即 $p_{C/E/A}$ 的核）。另一方面，设 $h:C\to E$ 是一个
$A$-同态，并且是 $p$ 的右逆；设 $j:L\to E$ 为典范单射，于是有
$q(b)=h(v(b))+j(D(b))$，对所有 $b\in B$ 成立；这里 $D$ 是一个
$A$-导射，它把 $B$ 映入 $L$，并定义 $B$-扩充 $E$。可以写成
$D=f\circ d_{B/A}$，其中
$f:\Omega^1_{B/A}\to L$ 是一个 $B$-同态。根据 (20.5.2.6)，有
\[
  q_{E/B/A}(d_{B/A}(b))=d_{E/A}(q(b))
  =d_{E/A}(h(v(b)))+d_{E/A}(j(f(d_{B/A}(b))))
\]
对所有 $b\in B$ 成立。现在令
$z=\sum_i d_{B/A}(b_i)\otimes c_i$，其中 $b_i\in B$ 且 $c_i\in C$，为
$\Omega^1_{B/A}\otimes_B C$ 的一个元素；则
\[
  (q_{E/B/A}\otimes1_C)(z)=
  \sum_i d_{E/A}(h(v(b_i)))\otimes c_i+
  \sum_i d_{E/A}(j(f(d_{B/A}(b_i))))\otimes c_i.
\]
在第一项和中，
$d_{E/A}(h(v(b_i)))=h_{E/C/A}(d_{C/A}(v(b_i)))$，故根据
(20.5.2.6)，该项和为
$(h_{E/C/A}\otimes1_C)(v_{C/B/A}(z))$。若取
$z\in\Upsilon_{C/B/A}$，该项和便为零，余下的项依照
$\delta_{C/E/A}$ 的定义给出
\[
  (q_{E/B/A}\otimes1_C)(z)=
  \sum_i\delta_{C/E/A}(c_i f(d_{B/A}(b_i)))
  =\delta_{C/E/A}(f(z)),
\]
这就证明了 (20.6.9.1) 的左方格交换。
\end{proof}

请注意，若只假设 $\delta_{C/E/A}$ 是\emph{单射}（未必有左逆），上述解释
仍然可以把 $\chi_E$ 定义为 $q_{E/B/A}\otimes1_C$ 在
$\Upsilon_{C/B/A}$ 上的限制。

\begin{corollary}[20.6.10]
\label{0.20.6.10-fr}
设 $\mathfrak{R}$ 是 $B$ 的一个理想，$C=B/\mathfrak{R}$，
$E=B/\mathfrak{R}^2$；若 $E$ 是一个 $B$-扩充，而且是 $A$-平凡的
(18.3.7)，它扩充 $C$，核为 $\mathfrak{R}/\mathfrak{R}^2$，那么特征同态
$\chi_E:\Upsilon_{C/B/A}\to\mathfrak{R}/\mathfrak{R}^2$ 是双射。
\end{corollary}

\begin{proof}
\oldpage[0\textsubscript{IV-1}]{140}
事实上，在图 (20.6.9.1) 中，右边的两个竖直箭头都是双射同态
(20.5.12, (ii))。
\end{proof}

\begin{theorem}[20.6.11]
\label{0.20.6.11-fr}
设 $u:A\to B$、$v:B\to C$ 是两个环同态，$L$ 是一个 $C$-模。假设下列
条件之一成立：
\begin{enumerate}
  \item[\rm(i)] $L$ 是一个内射 $C$-模。
  \item[\rm(ii)] $\Upsilon_{C/B/A}$ 是 $C$-模
$\Omega^1_{B/A}\otimes_B C$ 的一个直和因子，并且
$u_{C/B/A}:\Omega^1_{C/A}\to\Omega^1_{C/B}$ 有右逆。
\end{enumerate}
那么典范同态 (20.6.7.1)
\[
  \operatorname{Exalcom}_{B/A}(C,L)\longrightarrow
  \operatorname{Hom}_C(\Upsilon_{C/B/A},L)
\]
是双射。

特别地，如果 $\Omega^1_{C/B}$ 和 $\Omega^1_{C/A}$ 都是投射 $C$-模，
那么典范同态 (20.6.7.1) 是双射。
\end{theorem}

\begin{proof}
条件 (i)、(ii) 中的每一个都蕴含 (20.6.7.1) 为双射；在两种情形下，这都
直接来自 $\alpha_1$ 的定义。还应指出，条件 (ii) 是 (20.6.7.1) 对
\emph{每个} $C$-模 $L$ 都为双射的\emph{充分必要}条件（Bourbaki，
\emph{Alg.}，第 II 章，第 3 版，\S~2，n\textsuperscript{o}~1，命题 1）。
若假设 $\Omega^1_{C/B}$ 和 $\Omega^1_{C/A}$ 是投射 $C$-模，则在正合列
(20.6.1.1) 中，$\operatorname{Ker}(u_{C/B/A})$ 是投射 $C$-模，因为正合列
\[
  0\longrightarrow\operatorname{Ker}(u_{C/B/A})\longrightarrow\Omega^1_{C/A}
  \longrightarrow\Omega^1_{C/B}\longrightarrow0
\]
分裂，这是由于 $\Omega^1_{C/B}$ 为投射模；又因为
$\operatorname{Ker}(u_{C/B/A})=\operatorname{Im}(v_{C/B/A})$，正合列
\[
  0\longrightarrow\Upsilon_{C/B/A}\longrightarrow
  \Omega^1_{B/A}\otimes_B C\longrightarrow
  \operatorname{Im}(v_{C/B/A})\longrightarrow0
\]
也分裂。
\end{proof}

\begin{corollary}[20.6.12]
\label{0.20.6.12-fr}
假设 $C$ 是一个形式光滑 $A$-代数。于是存在典范同态
\begin{equation}
\tag{20.6.12.1}
\label{0.20.6.12.1-fr}
  \operatorname{Exalcom}_B(C,L)\longrightarrow
  \operatorname{Hom}_C(\Upsilon_{C/B/A},L).
\end{equation}
此外，如果 (20.6.11) 的条件 (i)、(ii) 之一成立，该同态就是双射。
\end{corollary}

\begin{proof}
事实上，由关于 $C$ 的假设可得
$\operatorname{Exalcom}_B(C,L)=\operatorname{Exalcom}_{B/A}(C,L)$，而同态
(20.6.12.1) 正是 (20.6.7.1)。
\end{proof}

\begin{corollary}[20.6.13]
\label{0.20.6.13-fr}
如果 $C$ 是一个形式光滑 $A$-代数，并且 $\Omega^1_{C/B}$ 是一个投射
$C$-模，那么同态 (20.6.12.1) 是双射。
\end{corollary}

\begin{proof}
事实上，此时知道 $\Omega^1_{C/A}$ 是一个投射 $C$-模 (20.4.9)，这里所取
拓扑都是离散的。
\end{proof}

\begin{env}[20.6.14]
\label{0.20.6.14-fr}
沿用上述记号，现在进一步假设 $A$、$B$、$C$ 都是 $\Lambda$-代数，并且
$u$、$v$ 是 $\Lambda$-同态；这等价于给定三个环同态
\[
  \Lambda\xrightarrow{s}A\xrightarrow{u}B\xrightarrow{v}C.
\]

\oldpage[0\textsubscript{IV-1}]{141}
因此，除了不恰模 $\Upsilon_{C/B/A}$ 之外，还有不恰模
$\Upsilon_{B/A/\Lambda}$、$\Upsilon_{C/A/\Lambda}$ 和
$\Upsilon_{C/B/\Lambda}$；而且已经定义了下列典范 $C$-模同态 (20.6.4.2)：
\begin{align}
\tag{20.6.14.1}
\label{0.20.6.14.1-fr}
  u':\Upsilon_{C/A/\Lambda}&\longrightarrow\Upsilon_{C/B/\Lambda},\\
\tag{20.6.14.2}
\label{0.20.6.14.2-fr}
  s':\Upsilon_{C/B/\Lambda}&\longrightarrow\Upsilon_{C/B/A}.
\end{align}

在交换图 (20.5.7.3)
\begin{equation}
\tag{20.6.14.3}
\label{0.20.6.14.3-fr}
\xymatrix{
  \Omega^1_{A/\Lambda}\otimes_A B \ar[r] \ar[d] &
  \Omega^1_{B/\Lambda} \ar[r] \ar[d] &
  \Omega^1_{B/A} \ar[r] \ar[d] & 0 \\
  \Omega^1_{A/\Lambda}\otimes_A C \ar[r] &
  \Omega^1_{C/\Lambda} \ar[r] & \Omega^1_{C/A} \ar[r] & 0
}
\end{equation}
中，下行由 $C$-模组成，因此张量化后得到交换图
\begin{equation}
\tag{20.6.14.4}
\label{0.20.6.14.4-fr}
\xymatrix{
  \Omega^1_{A/\Lambda}\otimes_A C \ar[r] \ar@{=}[d] &
  \Omega^1_{B/\Lambda}\otimes_B C \ar[r] \ar[d] &
  \Omega^1_{B/A}\otimes_B C \ar[r] \ar[d] & 0 \\
  \Omega^1_{A/\Lambda}\otimes_A C \ar[r] &
  \Omega^1_{C/\Lambda} \ar[r] & \Omega^1_{C/A} \ar[r] & 0
}
\end{equation}
其中第一行仍然正合，左边的竖直箭头是恒等同态。若置
\begin{equation}
\tag{20.6.14.5}
\label{0.20.6.14.5-fr}
  \Upsilon^C_{B/A/\Lambda}=
  \operatorname{Ker}(u_{B/A/\Lambda}\otimes1_C)=
  \operatorname{Ker}(\Omega^1_{A/\Lambda}\otimes_A C\longrightarrow
       \Omega^1_{B/\Lambda}\otimes_B C),
\end{equation}
考虑到 $\Upsilon_{C/A/\Lambda}$ 的定义，可见存在唯一的 $C$-同态
\begin{equation}
\tag{20.6.14.6}
\label{0.20.6.14.6-fr}
  v':\Upsilon^C_{B/A/\Lambda}\longrightarrow\Upsilon_{C/A/\Lambda}
\end{equation}
使图
\begin{equation}
\tag{20.6.14.7}
\label{0.20.6.14.7-fr}
\xymatrix{
  0 \ar[r] & \Upsilon^C_{B/A/\Lambda} \ar[r] \ar[d]^{v'} &
  \Omega^1_{A/\Lambda}\otimes_A C \ar[r] \ar@{=}[d] &
  \Omega^1_{B/\Lambda}\otimes_B C \ar[r] \ar[d] &
  \Omega^1_{B/A}\otimes_B C \ar[r] \ar[d] & 0 \\
  0 \ar[r] & \Upsilon_{C/A/\Lambda} \ar[r] &
  \Omega^1_{A/\Lambda}\otimes_A C \ar[r] &
  \Omega^1_{C/\Lambda} \ar[r] & \Omega^1_{C/A} \ar[r] & 0
}
\end{equation}
交换；它的两行均正合。

当 $\Lambda=\mathbb{Z}$ 时，用 $\Upsilon^C_{B/A}$ 代替
$\Upsilon^C_{B/A/\mathbb{Z}}$。如果 $\Lambda$ 是素域，则
$\Upsilon^C_{B/A/\Lambda}=\Upsilon^C_{B/A}$。
\end{env}

\begin{env}[20.6.15]
\label{0.20.6.15-fr}
为了研究上述各模之间的关系，一方面引入 $B$-模复形
$\mathrm{K}_\bullet(B/A/\Lambda)$，另一方面引入 $C$-模复形
\oldpage[0\textsubscript{IV-1}]{142}
$\mathrm{K}_\bullet(C/A/\Lambda)$ 和 $\mathrm{K}_\bullet(C/B/\Lambda)$
(20.6.5)，此外再引入下列 $C$-模复形。首先置
\begin{equation}
\tag{20.6.15.1}
\label{0.20.6.15.1-fr}
  \mathrm{K}^C_\bullet(B/A/\Lambda)=
  \mathrm{K}_\bullet(B/A/\Lambda)\otimes_B C.
\end{equation}
其次，用 $\mathrm{T}_\bullet(C/B/\Lambda)$ 表示如下 $C$-模复形：除次数 0
和 1 外各项均为零，并且
\begin{equation}
\tag{20.6.15.2}
\label{0.20.6.15.2-fr}
  \mathrm{T}_0(C/B/\Lambda)=\mathrm{T}_1(C/B/\Lambda)=
  \Omega^1_{B/\Lambda}\otimes_B C,
\end{equation}
微分算子取恒等同态，因而这一复形零伦；最后置
\begin{equation}
\tag{20.6.15.3}
\label{0.20.6.15.3-fr}
  \mathrm{K}'_\bullet(C/A/\Lambda)=
  \mathrm{K}_\bullet(C/A/\Lambda)\oplus
  \mathrm{T}_\bullet(C/B/\Lambda).
\end{equation}
由 $\mathrm{T}_\bullet$ 的平凡性，显然对每个 $C$-模 $L$ 和每个 $i$ 都有
\begin{equation}
\tag{20.6.15.4}
\label{0.20.6.15.4-fr}
  \mathrm{H}^i(\mathrm{K}'_\bullet,L)=
  \mathrm{H}^i(\mathrm{K}_\bullet,L)
  \quad\text{et}\quad
  \mathrm{H}_i(\mathrm{K}'_\bullet,L)=
  \mathrm{H}_i(\mathrm{K}_\bullet,L).
\end{equation}
\end{env}
\begin{env}[20.6.16]
\label{0.20.6.16-fr}
现在按如下方式定义一个在每个次数都\emph{分裂}的复形正合列
\begin{equation}
\tag{20.6.16.1}
\label{0.20.6.16.1-fr}
  0\to\mathrm{K}^C_\bullet(B/A/\Lambda)
  \xrightarrow{j}\mathrm{K}'_\bullet(C/A/\Lambda)
  \xrightarrow{p}\mathrm{K}_\bullet(C/B/\Lambda)\to0.
\end{equation}
暂且用
\[
  f:\Omega^1_{A/\Lambda}\otimes_A C\to
    \Omega^1_{B/\Lambda}\otimes_B C,
  \qquad
  g:\Omega^1_{B/\Lambda}\otimes_B C\to\Omega^1_{C/\Lambda}
\]
分别表示典范同态 $u_{B/A/\Lambda}\otimes1_C$ 和
$v_{C/B/\Lambda}$；它们的复合为
$g\circ f=(v\circ u)_{C/A/\Lambda}$（参见 \sref{0.20.6.14.4-fr}）。取
$j_1(x)=(x,f(x))$，$p_1(y,z)=z-f(y)$，
$j_0(x)=(g(x),x)$，$p_0(y,z)=g(z)-y$。于是
$\operatorname{Im}(j_1)=\operatorname{Ker}(p_1)$ 是 $f$ 的图形，它与
$\{0\}\oplus\mathrm{T}_1$ 互补；而
$\operatorname{Im}(j_0)=\operatorname{Ker}(p_0)$ 是 $g$ 的图形，它与
$\mathrm{K}_1(C/A/\Lambda)\oplus\{0\}$ 互补。图
\[
\xymatrix{
  0 \ar[r] & \mathrm{K}^C_1(B/A/\Lambda) \ar[r]^{j_1} \ar[d] &
  \mathrm{K}'_1(C/A/\Lambda) \ar[r]^{p_1} \ar[d] &
  \mathrm{K}_1(C/B/\Lambda) \ar[r] \ar[d] & 0 \\
  0 \ar[r] & \mathrm{K}^C_0(B/A/\Lambda) \ar[r]_{j_0} &
  \mathrm{K}'_0(C/A/\Lambda) \ar[r]_{p_0} &
  \mathrm{K}_0(C/B/\Lambda) \ar[r] & 0
}
\]
中竖直箭头都是微分算子；立即可以验证此图交换。
\end{env}

\begin{theorem}[20.6.17]
\label{0.20.6.17-fr}
有如下 $C$-模正合列：
\begin{equation}
\tag{20.6.17.1}
\label{0.20.6.17.1-fr}
  0\to\Upsilon^C_{B/A/\Lambda}\xrightarrow{v'}
  \Upsilon_{C/A/\Lambda}\xrightarrow{u'}
  \Upsilon_{C/B/\Lambda}\xrightarrow{\partial}
  \Omega^1_{B/A}\otimes_B C\xrightarrow{v_{C/B/A}}
  \Omega^1_{C/A}\xrightarrow{u_{C/B/A}}\Omega^1_{C/B}\to0,
\end{equation}
\oldpage[0\textsubscript{IV-1}]{143}
其中边缘算子 $\partial$ 是复合
\begin{equation}
\tag{20.6.17.2}
\label{0.20.6.17.2-fr}
  \Upsilon_{C/B/\Lambda}\xrightarrow{s'}\Upsilon_{C/B/A}
  \longrightarrow\Omega^1_{B/A}\otimes_B C,
\end{equation}
而第二个箭头为典范单射。
\end{theorem}

\begin{proof}
对复形正合列 \sref{0.20.6.16.1-fr} 写出同调正合列，就得到
\sref{0.20.6.17.1-fr}，因为次数 0、1 以外的同调均为零；而该正合列中由
$j$ 和 $p$ 的函子性所给出的同态确实就是命题中的同态，这是立即可见的。
只须验证 $\partial$ 等于 \sref{0.20.6.17.2-fr}。事实上，一个元素
$z\in\Upsilon_{C/B/\Lambda}$ 是经 $p_1$ 由 $(0,z)$ 所给出的像，由此立即推出
$\partial(z)$ 是 $z$ 在典范同态
\[
  s_{B/A/\Lambda}\otimes1_C:
  \Omega^1_{B/\Lambda}\otimes_B C\to\Omega^1_{B/A}\otimes_B C
\]
下的像。我们的断言由交换图
\begin{equation}
\tag{20.6.17.3}
\label{0.20.6.17.3-fr}
\xymatrix{
  \Upsilon_{C/B/\Lambda} \ar[r]^{s'} \ar[dd] & \Upsilon_{C/B/A} \ar[dd] \\
  & \\
  \Omega^1_{B/\Lambda}\otimes_B C \ar[r] &
  \Omega^1_{B/A}\otimes_B C
}
\end{equation}
得出（参见 \sref{0.20.6.4.3-fr}）。
\end{proof}

\begin{corollary}[20.6.18]
\label{0.20.6.18-fr}
\begin{enumerate}
  \item[\rm(i)] $C$-模序列
\begin{equation}
\tag{20.6.18.1}
\label{0.20.6.18.1-fr}
  0\to\Upsilon^C_{B/A/\Lambda}\xrightarrow{v'}
  \Upsilon_{C/A/\Lambda}\xrightarrow{u'}
  \Upsilon_{C/B/\Lambda}\xrightarrow{s'}\Upsilon_{C/B/A}\to0
\end{equation}
是正合的。
  \item[\rm(ii)] 如果 $B$ 是一个 $A$-代数，并且相对于 $\Lambda$ 形式光滑，
则有 $\Upsilon^C_{B/A/\Lambda}=0$。
\end{enumerate}
\end{corollary}

\begin{proof}
(i) 在 \sref{0.20.6.17.1-fr} 中，$\Upsilon_{C/B/\Lambda}$ 在
$\partial$ 下的像是 $v_{C/B/A}$ 的核，因而按定义就是
$\Upsilon_{C/B/A}$。

(ii) 该假设蕴含序列
\[
  0\to\Omega^1_{A/\Lambda}\otimes_A B\to
  \Omega^1_{B/\Lambda}\to\Omega^1_{B/A}\to0
\]
正合且分裂 \sref{0.20.5.7-fr}；与 $C$ 张量后，序列
\[
  0\to\Omega^1_{A/\Lambda}\otimes_A C\to
  \Omega^1_{B/\Lambda}\otimes_B C\to
  \Omega^1_{B/A}\otimes_B C\to0
\]
仍然正合，断言由此得出。
\end{proof}

\begin{corollary}[20.6.19]
\label{0.20.6.19-fr}
设 $K$ 是一个域，$E$、$F$ 是 $K$ 的两个扩张，并且
$K\subset E\subset F$。
\begin{enumerate}
  \item[\rm(i)] 如果 $F$ 是 $E$ 的可分扩张，则
  $\Upsilon_{F/E/K}=0$（换言之，典范同态
  $\Omega^1_{E/K}\otimes_E F\to\Omega^1_{F/K}$ 是单射）。
  \item[\rm(ii)] 反过来，如果 $F$ 是 $K$ 的可分扩张，并且
  $\Upsilon_{F/E/K}=0$，那么 $F$ 是 $E$ 的可分扩张。
\end{enumerate}
\end{corollary}

\begin{proof}
若 $P$ 是 $K$ 的素域，则确有正合列 \sref{0.20.6.18-fr}
\[
  \Upsilon_{F/K/P}\to\Upsilon_{F/E/P}\to\Upsilon_{F/E/K}\to0.
\]
如果 $\Upsilon_{F/E/P}=\Upsilon_{F/E}=0$，那么
$\Upsilon_{F/E/K}=0$，从而由 \sref{0.20.6.3-fr} 得到 (i)；反过来，如果
$\Upsilon_{F/E/K}=0$ 且 $\Upsilon_{F/K/P}=\Upsilon_{F/K}=0$，则
$\Upsilon_{F/E/P}=0$，从而由 \sref{0.20.6.3-fr} 得到 (ii)。
\end{proof}

\oldpage[0\textsubscript{IV-1}]{144}
\begin{corollary}[20.6.20]
\label{0.20.6.20-fr}
\begin{enumerate}
  \item[\rm(i)] 如果 $K$ 是 $k$ 的代数可分扩张，则
  $\Omega^1_{K/k}=0$。
  \item[\rm(ii)] 设 $k$ 是特征为 0 的域，$K$ 是 $k$ 的一个扩张。
  $\Omega^1_{K/k}=0$ 的充分必要条件是 $K$ 为 $k$ 的代数扩张。特别地，
  对于特征为 0 的域 $K$，$\Omega^1_K=0$ 的充分必要条件是 $K$ 为
  $\mathbb{Q}$ 的代数扩张（参见 (21.4.4) 和 (21.7.5)）。
\end{enumerate}
\end{corollary}

\begin{proof}
(i) 对每个 $x\in K$，设 $f$ 是 $x$ 在 $k$ 上的极小多项式；由于
$f'(x)\ne0$ 且 $d_{K/k}(f(x))=f'(x)d_{K/k}(x)=0$，故
$d_{K/k}(x)=0$，断言由 \sref{0.20.4.7-fr} 得出。

(ii) 存在一个纯扩张 $L$，它是 $k$ 的扩张，使得 $k\subset L\subset K$，并且
$K$ 是 $L$ 的代数扩张。由于 $K$ 在 $L$ 上可分，由
\sref{0.20.6.19-fr}, (i) 可知序列 \sref{0.20.5.7.2-fr}
\[
  0\to\Omega^1_{L/k}\otimes_L K\to\Omega^1_{K/k}
  \to\Omega^1_{K/L}\to0
\]
正合；又由 (i) 有 $\Omega^1_{K/L}=0$。于是关系
$\Omega^1_{K/k}=0$ 等价于 $\Omega^1_{L/k}=0$；因为 $L$ 是 $k$ 的纯扩张，
由 \sref{0.20.5.10-fr} 可知关系 $\Omega^1_{L/k}=0$ 等价于 $L=k$。
\end{proof}

\begin{remarks}[20.6.21]
\label{0.20.6.21-fr}
\begin{enumerate}
  \item[\rm(i)] 由于 $\Upsilon_{B/A/\Lambda}$ 是同态
  \[
    u_{B/A/\Lambda}:\Omega^1_{A/\Lambda}\otimes_A B
    \to\Omega^1_{B/\Lambda}
  \]
  的核，而 $\Upsilon^C_{B/A/\Lambda}$ 是
  $u_{B/A/\Lambda}\otimes1_C$ 的核，故有典范同态
\begin{equation}
\tag{20.6.21.1}
\label{0.20.6.21.1-fr}
  \Upsilon_{B/A/\Lambda}\otimes_B C\to\Upsilon^C_{B/A/\Lambda}.
\end{equation}
当序列
\begin{equation}
\tag{20.6.21.2}
\label{0.20.6.21.2-fr}
  0\to\Upsilon_{B/A/\Lambda}\otimes_B C\to
  \Omega^1_{A/\Lambda}\otimes_A C\to
  \Omega^1_{B/\Lambda}\otimes_B C\to
  \Omega^1_{B/A}\otimes_B C\to0
\end{equation}
正合时，此同态为双射；下列情形中确实如此：
\begin{enumerate}
  \item[$1^\circ$] $C$ 是一个平坦 $B$-模。
  \item[$2^\circ$] $B$-模 $\Omega^1_{B/\Lambda}$ 和
  $\Omega^1_{B/A}$ 都是\emph{平坦的}：事实上，此时
  $\operatorname{Ker}(\Omega^1_{B/\Lambda}\to\Omega^1_{B/A})$ 也平坦
  （$\mathbf{0_I}$, 6.1.2），于是由（$\mathbf{0_I}$, 6.1.2），序列
  \sref{0.20.6.21.2-fr} 正合。
\end{enumerate}
  \item[\rm(ii)] 考虑环同态的交换图
\[
\xymatrix{
  \Lambda' \ar[r] & A' \ar[r] & B' \ar[r] & C' \\
  \Lambda \ar[u] \ar[r] & A \ar[u] \ar[r] & B \ar[u] \ar[r] & C \ar[u]
}
\]
那么 \sref{0.20.6.16-fr} 中的定义表明有如下复形交换图（其中竖直箭头来自
\sref{0.20.6.6.4-fr}）：
\[
\xymatrix{
  0 \ar[r] & \mathrm{K}^C_\bullet(B/A/\Lambda) \ar[r] \ar[d] &
  \mathrm{K}'_\bullet(C/A/\Lambda) \ar[r] \ar[d] &
  \mathrm{K}_\bullet(C/B/\Lambda) \ar[r] \ar[d] & 0 \\
  0 \ar[r] & \mathrm{K}^{C'}_\bullet(B'/A'/\Lambda') \ar[r] &
  \mathrm{K}'_\bullet(C'/A'/\Lambda') \ar[r] &
  \mathrm{K}_\bullet(C'/B'/\Lambda') \ar[r] & 0
}
\]
\oldpage[0\textsubscript{IV-1}]{145}
在同调上取诱导映射，便得到\emph{交换图}
\begin{equation}
\tag{20.6.21.3}
\label{0.20.6.21.3-fr}
\xymatrix@C=0.7em{
  0 \ar[r] & \Upsilon^C_{B/A/\Lambda} \ar[r] \ar[d] &
  \Upsilon_{C/A/\Lambda} \ar[r] \ar[d] &
  \Upsilon_{C/B/\Lambda} \ar[r] \ar[d] &
  \Omega^1_{B/A}\otimes_B C \ar[r] \ar[d] &
  \Omega^1_{C/A} \ar[r] \ar[d] &
  \Omega^1_{C/B} \ar[r] \ar[d] & 0 \\
  0 \ar[r] & \Upsilon^{C'}_{B'/A'/\Lambda'} \ar[r] &
  \Upsilon_{C'/A'/\Lambda'} \ar[r] &
  \Upsilon_{C'/B'/\Lambda'} \ar[r] &
  \Omega^1_{B'/A'}\otimes_{B'} C' \ar[r] &
  \Omega^1_{C'/A'} \ar[r] & \Omega^1_{C'/B'} \ar[r] & 0
}
\end{equation}
对于 \sref{0.20.6.18.1-fr} 也有类似的交换图。
\end{enumerate}
\end{remarks}

\begin{proposition}[20.6.22]
\label{0.20.6.22-fr}
设 $s:\Lambda\to A$、$u:A\to B$ 是两个环同态，$\mathfrak{R}$ 是 $B$
的一个理想，$C$ 是商环 $B/\mathfrak{R}$。假设
$E=B/\mathfrak{R}^2$ 是一个 $B$-扩充，它是 $\Lambda$-平凡的，并且扩充
$C$，核为 $\mathfrak{R}/\mathfrak{R}^2$。于是有正合列
\begin{equation}
\tag{20.6.22.1}
\label{0.20.6.22.1-fr}
  0\to\Upsilon^C_{B/A/\Lambda}\xrightarrow{v'}
  \Upsilon_{C/A/\Lambda}\xrightarrow{\chi_E}
  \mathfrak{R}/\mathfrak{R}^2\xrightarrow{\delta_{C/B/A}}
  \Omega^1_{B/A}\otimes_B C\xrightarrow{v_{C/B/A}}
  \Omega^1_{C/A}\to0.
\end{equation}
\end{proposition}

\begin{proof}
事实上，由于 $v:B\to C$ 是满射，故 $\Omega^1_{C/B}=0$
\sref{0.20.4.12-fr}。此外，由 \sref{0.20.6.10-fr} 可知
$\Upsilon_{C/B/\Lambda}$ 与 $\mathfrak{R}/\mathfrak{R}^2$ 典范同构。于是只须
应用正合列 \sref{0.20.6.17.1-fr} 和 \sref{0.20.6.18.1-fr}，并注意有交换图
\[
\xymatrix{
  \Upsilon_{C/A/\Lambda} \ar[r] \ar[d]_{\chi_E} &
  \Omega^1_{A/\Lambda}\otimes_A C \ar[d] \\
  \Upsilon_{C/E/\Lambda}=\Upsilon_{C/B/\Lambda}=\mathfrak{R}/\mathfrak{R}^2
    \ar[r]^-{\delta_{C/B/A}} \ar[d]_{\ell} &
  \Omega^1_{B/A}\otimes_B C=\Omega^1_{E/A}\otimes_E C \ar[d] \\
  \mathfrak{R}/\mathfrak{R}^2 \ar[r]^{\delta_{C/B/A}} &
  \Omega^1_{B/A}\otimes_B C
}
\]
并使用图 \sref{0.20.6.17.3-fr} 的交换性即可。
\end{proof}

这样便明确了 $\delta_{C/B/A}$ 的\emph{核}；这里所说的情形是：存在环
$\Lambda$，使得 $A$ 是一个 $\Lambda$-代数且 $B/\mathfrak{R}^2$ 是
$\Lambda$-平凡的。应当注意，当 $C$ 是一个\emph{形式光滑
$\Lambda$-代数}（取离散拓扑）时尤其如此。

再假设有环同态的交换图
\[
\xymatrix{
  \Lambda' \ar[r] & A' \ar[r] & B' \\
  \Lambda \ar[u] \ar[r] & A \ar[u] \ar[r] & B \ar[u]^f
}
\]
\oldpage[0\textsubscript{IV-1}]{146}
并且 $\mathfrak{R}'$ 是 $B'$ 的一个理想，满足
$f(\mathfrak{R})\subset\mathfrak{R}'$；还假设
$E'=B'/\mathfrak{R}'^2$ 是一个 $B'$-扩充，它是 $\Lambda'$-平凡的，并且
扩充 $C'=B'/\mathfrak{R}'$，核为 $\mathfrak{R}'/\mathfrak{R}'^2$。于是有
交换图
\begin{equation}
\tag{20.6.22.2}
\label{0.20.6.22.2-fr}
\xymatrix@C=0.7em{
  0 \ar[r] & \Upsilon^{C'}_{B'/A'/\Lambda'} \ar[r] &
  \Upsilon_{C'/A'/\Lambda'} \ar[r] &
  \mathfrak{R}'/\mathfrak{R}'^2 \ar[r] &
  \Omega^1_{B'/A'}\otimes_{B'} C' \ar[r] &
  \Omega^1_{C'/A'} \ar[r] & 0 \\
  0 \ar[r] & \Upsilon^C_{B/A/\Lambda} \ar[r] \ar[u] &
  \Upsilon_{C/A/\Lambda} \ar[r] \ar[u] &
  \mathfrak{R}/\mathfrak{R}^2 \ar[r] \ar[u] &
  \Omega^1_{B/A}\otimes_B C \ar[r] \ar[u] &
  \Omega^1_{C/A} \ar[r] \ar[u] & 0
}
\end{equation}
这是由 \sref{0.20.6.21.3-fr} 和 \sref{0.20.5.11.3-fr} 得出的。

\begin{corollary}[20.6.23]
\label{0.20.6.23-fr}
在 \sref{0.20.6.22-fr} 的假设下，再假设 $B$ 是一个形式光滑
$A$-代数。于是有正合列
\begin{equation}
\tag{20.6.23.1}
\label{0.20.6.23.1-fr}
  0\to\Upsilon_{C/A/\Lambda}\xrightarrow{\chi_E}
  \mathfrak{R}/\mathfrak{R}^2\xrightarrow{\delta_{C/B/A}}
  \Omega^1_{B/A}\otimes_B C\to\Omega^1_{C/A}\to0.
\end{equation}
\end{corollary}

\begin{proof}
这确实由 \sref{0.20.6.18-fr}, (ii) 得出。
\end{proof}

\begin{env}[20.6.24]
\label{0.20.6.24-fr}
当 \sref{0.20.6.22-fr} 的假设成立时，还称特征同态 $\chi_E$ 为
\emph{$A$-代数 $B$ 关于 $\Lambda$ 和理想 $\mathfrak{R}$ 的特征同态}
（在不会混淆时省略后面这些说明）；有时把它记为 $\chi_B$ 或
$\chi_{B/A}$。
\end{env}

\begin{proposition}[20.6.25]
\label{0.20.6.25-fr}
设 $s:\Lambda\to A$、$u:A\to B$、$v:B\to C$ 是三个环同态，$L$ 是一个
$C$-模。于是有正合列
\begin{equation}
\tag{20.6.25.1}
\label{0.20.6.25.1-fr}
\begin{split}
  0\to\operatorname{D\acute{e}r}_B(C,L)\xrightarrow{u^1}
  \operatorname{D\acute{e}r}_A(C,L)\xrightarrow{v^1}
  \operatorname{D\acute{e}r}_A(B,L)\xrightarrow{\partial} \\
  \to\operatorname{Exalcom}_{B/\Lambda}(C,L)\xrightarrow{u^1}
  \operatorname{Exalcom}_{A/\Lambda}(C,L)\xrightarrow{v^1}
  \operatorname{Exalcom}_{A/\Lambda}(B,L)\to0,
\end{split}
\end{equation}
其中 $u^1$、$v^1$ 是 (18.3.6.4) 和 (18.3.6.2) 中定义的同态，而
$\partial$ 按 \sref{0.20.2.2-fr} 中的方式定义。
\end{proposition}

\begin{proof}
事实上，由于正合列 \sref{0.20.6.16.1-fr} 分裂，便导出正合列
\[
  0\to\operatorname{Hom}_C(\mathrm{K}_\bullet(C/B/\Lambda),L)\to
  \operatorname{Hom}_C(\mathrm{K}'_\bullet(C/A/\Lambda),L)\to
  \operatorname{Hom}_C(\mathrm{K}^C_\bullet(B/A/\Lambda),L)\to0.
\]
对这个复形应用上同调正合列，并考虑 \sref{0.20.6.15.4-fr} 和
\sref{0.20.6.6-fr}，就得到 \sref{0.20.6.25.1-fr}，因为按定义有
\[
  \operatorname{Hom}_C(\mathrm{K}^C_\bullet(B/A/\Lambda),L)
  =\operatorname{Hom}_B(\mathrm{K}_\bullet(B/A/\Lambda),L).
\]
$u^1$ 和 $v^1$ 与 (18.3.4.2) 中各同态的等同性，由
\sref{0.20.6.6.4-fr} 得出。
\end{proof}

\begin{remark}[20.6.26]
\label{0.20.6.26-fr}
在本节中，复形 $\mathrm{K}_\bullet(C/B/A)$ 作为一种技术手段出现，用来简化
某些函子性行为的叙述。事实上，把这些复形视为
\oldpage[0\textsubscript{IV-1}]{147}
$C$-模复形的“差一个伦”的
范畴中的对象（也就是说，其态射是相互同伦的同态之类），它们是十分重要的
不变量，并且比由 $\Omega^1_{C/B}$ 和 $\Upsilon_{C/B/A}$ 组成的二元组更精细。
当 $k$ 是一个素域，并且 $C$ 是一个形式光滑 $k$-代数（例如，在 $k$ 的
某个扩张上有限型的正则环（参见（$\mathbf{IV}$, 6.8.6）））时，可以证明
复形 $\mathrm{K}_\bullet(C/A/k)$ 只用 $C$ 和 $A$ 就能描述（不涉及 $k$）：
把 $C$ 表示为一个多项式代数 $B$ 的商；该代数定义在 $A$ 上，所取的理想
记为 $\mathfrak{Q}$。再考虑只有两个非零项的复形
$\mathrm{F}_\bullet(C/A)$
\[
  \cdots\to0\to\mathfrak{Q}/\mathfrak{Q}^2\to
  \Omega^1_{B/A}\otimes_B C\to0\to\cdots
\]
（由 \sref{0.20.6.23.1-fr}，它的同调确实与
$\mathrm{K}_\bullet(C/A/k)$ 的同调相同）。从同调代数的观点看，这些复形
$\mathrm{F}_\bullet(C/A)$ 对于 $\operatorname{Spec}(C)$ 在
$\operatorname{Spec}(A)$ 上所起的作用就是余法丛；它们将在本书论述凝聚层
对偶性和 Riemann--Roch 定理的各章中占据重要地位。
\end{remark}

\subsection{推广到拓扑环上}
\label{subsection:0.20.7-fr}

\begin{env}[20.7.1]
\label{0.20.7.1-fr}
由定义立即可知，如果在 \sref{0.20.5.2-fr} 和 \sref{0.20.5.3-fr} 中假设环
$A$、$B$、$C$ 都是拓扑环，而环同态 $u$、$v$ 都连续，那么同态
$u_{C/B/A}$ 和 $v_{C/B/A}$ 都是\emph{连续的}（当然，必须在
$\Omega^1_{B/A}\otimes_B C$ 上取张量积拓扑）。

此外，$u_{C/B/A}:\Omega^1_{C/A}\to\Omega^1_{C/B}$ 是拓扑 $C$-模的一个
\emph{严格满态射}；事实上，根据张量积拓扑的定义，典范同态
$C\otimes_A C\to C\otimes_B C$ 具有这一性质；取商后，由此推出相应的同态
$\mathrm{P}^1_{C/A}\to\mathrm{P}^1_{C/B}$ 是严格满态射，而
$u_{C/B/A}$ 是后一个同态在 $\Omega^1_{C/A}$ 上的限制。

在 \sref{0.20.5.5-fr} 中，如果 $A'$ 和 $B$ 是拓扑 $A$-代数，并且在
$B'$ 与 $\Omega^1_{B/A}\otimes_B B'$ 上都赋予张量积拓扑，则考虑到
$\mathrm{P}^1_{B/A}$ 是 $B$ 与 $\Omega^1_{B/A}$ 的拓扑直和
\sref{0.20.4.8-fr}，典范同态 \sref{0.20.5.5.1-fr} 是拓扑同构。
\end{env}

\begin{proposition}[20.7.2]
\label{0.20.7.2-fr}
设 $u:A\to B$、$v:B\to C$ 是拓扑环的两个连续同态。连续同态
$v_{C/B/A}:\Omega^1_{B/A}\otimes_B C\to\Omega^1_{C/A}$ 形式左可逆
（参见 \sref{0.19.1.5-fr}）的充分必要条件，是 $C$ 为一个 $B$-代数且相对于
$A$ 形式平滑 \sref{0.19.9.1-fr}（特别地，只要 $C$ 是一个形式平滑
$B$-代数就已足够）。
\end{proposition}

\begin{proof}
说 $v_{C/B/A}$ 形式左可逆，根据 $\Omega^1_{C/A}$ 和
$\Omega^1_{B/A}\otimes_B C$ 的拓扑都不细于由 $C$ 的拓扑导出的拓扑
\sref{0.20.4.5-fr}，就是指：对每个离散 $C$-模 $L$，若它被 $C$ 的一个
开理想零化，则典范同态
\[
  \operatorname{Hom.\ cont}_C(\Omega^1_{C/A},L)\to
  \operatorname{Hom.\ cont}_C(\Omega^1_{B/A}\otimes_B C,L)
\]
\oldpage[0\textsubscript{IV-1}]{148}
都是\emph{满射} \sref{0.19.1.5-fr}。由张量积拓扑的定义，
$\operatorname{Hom.\ cont}_C(\Omega^1_{B/A}\otimes_B C,L)=
\operatorname{Hom.\ cont}_B(\Omega^1_{B/A},L)$；因此根据
\sref{0.20.4.8-fr}，这等价于说典范同态
\[
  \operatorname{D\acute{e}r.\ cont}_A(C,L)\to
  \operatorname{D\acute{e}r.\ cont}_A(B,L)
\]
为满射。但是正合列 \sref{0.20.3.8.2-fr} 表明，这一条件等价于
$\operatorname{Exalcotop}_{B/A}(C,L)=0$，也就是说，恰好等价于 $C$ 相对于
$A$ 形式平滑 \sref{0.19.9.8-fr}。
\end{proof}

\begin{corollary}[20.7.3]
\label{0.20.7.3-fr}
假设在 $B$ 和 $C$ 中，每个开理想的平方仍为开理想。$C$ 是一个 $B$-代数且
相对于 $A$ 形式平滑的充分必要条件是：若用 $(\mathfrak{R}_\lambda)$ 表示一个
$0$ 的邻域基，它由 $C$ 的理想组成，则对每个 $\lambda$，同态
\begin{equation}
\tag{20.7.3.1}
\label{0.20.7.3.1-fr}
  v_{C/B/A}\otimes 1_{C/\mathfrak{R}_\lambda}:
  \Omega^1_{B/A}\otimes_B(C/\mathfrak{R}_\lambda)\to
  \Omega^1_{C/A}\otimes_C(C/\mathfrak{R}_\lambda)
\end{equation}
都有左逆。
\end{corollary}

\begin{proof}
事实上，此时知道 $\Omega^1_{B/A}$（相应地，$\Omega^1_{C/A}$）的拓扑由
$B$（相应地，$C$）的拓扑导出 \sref{0.20.4.5-fr}；由此立即推出
$\Omega^1_{B/A}\otimes_B C$ 的拓扑也由 $C$ 的拓扑导出；于是本推论由
\sref{0.20.7.2-fr} 和 \sref{0.19.1.7-fr} 得出。
\end{proof}

\begin{proposition}[20.7.4]
\label{0.20.7.4-fr}
设 $A$ 是一个拓扑环，$B$ 是一个拓扑 $A$-代数。$B$ 为形式非分歧的
\sref{0.19.10.2-fr} 充分必要条件，是分离完备化
$\widehat{\Omega}^1_{B/A}$ 为零。
\end{proposition}

\begin{proof}
事实上，由 \sref{0.19.10.4-fr} 和 \sref{0.20.1.1-fr} 立即可知，$B$ 为形式
非分歧的充分必要条件是：对于每个开理想 $\mathfrak{R}$（属于 $B$）、每个
开理想 $\mathfrak{H}$（属于 $A$）且满足 $\mathfrak{H}B\subset\mathfrak{R}$，以及每个
$(B/\mathfrak{R})$-模 $L$，都有
$\operatorname{D\acute{e}r}_{A/\mathfrak{H}}(B/\mathfrak{R},L)=0$；这就是说，
$\operatorname{D\acute{e}r.\ cont}_A(B,L)=0$ 对任意离散 $B$-模 $L$ 均成立，
只要它被 $B$ 的一个开理想零化。根据 \sref{0.20.4.8-fr}，这等价于
$\operatorname{Hom.\ cont}_B(\Omega^1_{B/A},L)=0$ 对这样的 $B$-模成立，
命题立即得证。

当 $B$ 为\emph{离散的}时，\sref{0.20.7.4-fr} 命题中的条件因而等价于
$\Omega^1_{B/A}=0$。
\end{proof}

\begin{corollary}[20.7.5]
\label{0.20.7.5-fr}
设 $A$ 是一个环，$\mathfrak{m}$ 是 $A$ 的一个理想，$B$ 是一个
$A$-代数；在 $A$ 上赋予 $\mathfrak{m}$-预进拓扑，在 $B$ 上赋予由 $A$
的拓扑导出的拓扑，并置 $A_0=A/\mathfrak{m}$、
$B_0=B/\mathfrak{m}B=B\otimes_A A_0$。那么，$B$ 为形式非分歧
$A$-代数的充分必要条件是 $\Omega^1_{B_0/A_0}=0$（也就是说，$B_0$ 为形式
非分歧 $A_0$-代数）。
\end{corollary}

\begin{proof}
事实上，由 \sref{0.20.5.5-fr} 有
$\Omega^1_{B/A}\otimes_B B_0=\Omega^1_{B/A}/\mathfrak{m}.\Omega^1_{B/A}
=\Omega^1_{B_0/A_0}$；说 $\Omega^1_{B/A}$ 的每个开子模都等于
$\Omega^1_{B/A}$，根据 \sref{0.20.4.5-fr}，等价于说
$\mathfrak{m}.\Omega^1_{B/A}=\Omega^1_{B/A}$，结论由此得出。
\end{proof}

\begin{proposition}[20.7.6]
\label{0.20.7.6-fr}
设 $u:A\to B$、$v:B\to C$ 是拓扑环的两个连续同态，并假设 $v$ 使 $C$
成为一个形式平展 $B$-代数；那么
$v_{C/B/A}:\Omega^1_{B/A}\otimes_B C\to\Omega^1_{C/A}$ 是一个形式双态射
\sref{0.19.1.2-fr}。
\end{proposition}

\begin{proof}
事实上，对每个离散 $C$-模 $L$，只要它被 $C$ 的一个开理想零化，此时就有
\sref{0.20.7.4-fr}
$\operatorname{D\acute{e}r.\ cont}_B(C,L)=0$，因而由
\sref{0.20.3.6.1-fr}，典范同态
\oldpage[0\textsubscript{IV-1}]{149}
$\operatorname{D\acute{e}r.\ cont}_A(C,L)\to
\operatorname{D\acute{e}r.\ cont}_A(B,L)$ 为单射；另一方面，由
\sref{0.20.7.2-fr} 它也是满射，故为双射。由此可知 $v_{C/B/A}$ 的像必在
$\Omega^1_{C/A}$ 中稠密；否则，$\Omega^1_{C/A}$ 对
$\operatorname{Im}(v_{C/B/A})$ 闭包的商将是分离且 $\neq0$ 的，因而会有一个
离散商 $L\neq0$，这与刚刚证明的事实矛盾（还要考虑
\sref{0.20.4.8-fr}）。因此，由 \sref{0.20.7.2-fr}，$v_{C/B/A}$ 已是一个
形式单态射；它同时也是形式满态射 \sref{0.19.1.2-fr}，故为形式双态射。
\end{proof}

\begin{corollary}[20.7.7]
\label{0.20.7.7-fr}
假设在 $B$ 和 $C$ 中，每个开理想的平方仍为开理想。如果 $C$ 是一个形式平展
$B$-代数，那么对每个开理想 $\mathfrak{R}_\lambda$（属于 $C$），同态
\sref{0.20.7.3.1-fr} 都是双射。
\end{corollary}

\begin{proof}
事实上，由 \sref{0.19.1.1-fr}，该同态为满射；由
\sref{0.20.7.3-fr}，它又为单射。
\end{proof}

\begin{proposition}[20.7.8]
\label{0.20.7.8-fr}
设 $u:A\to B$ 是拓扑环的一个连续同态，$\mathfrak{R}$ 是 $B$ 的一个理想，
$C$ 是商拓扑环 $B/\mathfrak{R}$，$v:B\to C$ 是典范同态。那么：
\begin{enumerate}
  \item[\rm(i)] 在正合列 \sref{0.20.5.12.1-fr}
\[
  \mathfrak{R}/\mathfrak{R}^2\xrightarrow{\delta_{C/B/A}}
  \Omega^1_{B/A}\otimes_B C\xrightarrow{v_{C/B/A}}\Omega^1_{C/A}\to0
\]
中，同态 $\delta_{C/B/A}$ 连续，而同态 $v_{C/B/A}$ 是拓扑 $C$-模的严格
态射。
  \item[\rm(ii)] $\delta_{C/B/A}$ 形式左可逆 \sref{0.19.1.5-fr} 的充分必要
条件是：对于每个离散 $C$-模 $L$，若它被 $C$ 的一个开理想零化，则典范同态
\begin{equation}
\tag{20.7.8.1}
\label{0.20.7.8.1-fr}
  \operatorname{Exalcotop}_A(C,L)\to\operatorname{Exalcotop}_A(B,L)
\end{equation}
都是单射。
\end{enumerate}
\end{proposition}

\begin{proof}
(i) 第一个断言是显然的。另一方面，由张量积拓扑的定义，典范同态
$v\otimes v:B\otimes_A B\to C\otimes_A C$ 是严格态射；立即可知（参见
\sref{0.20.5.2-fr}），$v_{C/B/A}$ 也具有这一性质。

(ii) 说 $\delta_{C/B/A}$ 形式左可逆，就是说对每个离散 $C$-模 $L$，若它被
$C$ 的一个开理想零化，则典范同态
\[
  \operatorname{Hom.\ cont}_C(\Omega^1_{B/A}\otimes_B C,L)\to
  \operatorname{Hom.\ cont}_C(\mathfrak{R}/\mathfrak{R}^2,L)
\]
为满射。考虑到 \sref{0.18.4.3-fr} 和 \sref{0.20.4.8-fr}，这等价于说典范
同态
\[
  \operatorname{D\acute{e}r.\ cont}_A(B,L)\to
  \operatorname{Exalcotop}_B(C,L)
\]
为满射，因而结论由正合列 \sref{0.20.3.6.1-fr} 得出。

考虑两边的定义 \sref{0.18.4.1-fr}，\sref{0.20.7.8.1-fr} 为单射这一事实还可
表述如下：对于每个开理想 $\mathfrak{M}$（属于 $A$）、每个开理想
$\mathfrak{N}$（属于 $B$）且满足 $\mathfrak{M}B\subset\mathfrak{N}$，以及
每个 $(A/\mathfrak{M})$-扩充 $E$，它扩充
$B/(\mathfrak{R}+\mathfrak{N})$，核为一个
$(B/(\mathfrak{R}+\mathfrak{N}))$-模 $L$；若 $E$ 关于典范同态
$B/\mathfrak{N}\to B/(\mathfrak{R}+\mathfrak{N})$ 的逆像是
$(A/\mathfrak{M})$-平凡的，则存在一个开理想
$\mathfrak{M}'\subset\mathfrak{M}$（属于 $A$），
\oldpage[0\textsubscript{IV-1}]{150}
以及一个开理想 $\mathfrak{N}'\subset\mathfrak{N}$（属于 $B$），使得
$\mathfrak{M}'B\subset\mathfrak{N}'$，并且 $E$ 关于典范同态
$B/(\mathfrak{R}+\mathfrak{N}')\to B/(\mathfrak{R}+\mathfrak{N})$ 的逆像是
$(A/\mathfrak{M}')$-平凡的。
\end{proof}

\begin{corollary}[20.7.9]
\label{0.20.7.9-fr}
如果拓扑 $A$-代数 $C=B/\mathfrak{R}$ 是形式平滑的，那么典范同态
$\delta_{C/B/A}$ 形式左可逆。
\end{corollary}

\begin{env}[20.7.10]
\label{0.20.7.10-fr}
在 \sref{0.20.6.1-fr} 中，当 $A$、$B$、$C$ 是拓扑环，$u$、$v$ 是连续
同态时，在 $\Upsilon_{C/B/A}$ 上赋予由 $\Omega^1_{B/A}\otimes_B C$ 的拓扑
诱导的拓扑；此时，只要图 \sref{0.20.6.4.1-fr} 中的同态连续，同态
\sref{0.20.6.4.2-fr} 和 \sref{0.20.6.4.4-fr} 也连续。此外，如果在
\sref{0.20.6.7-fr} 中假设 $L$ 是一个离散 $C$-模，并且被 $C$ 的一个开理想
零化，那么取归纳极限便导出典范 $C$-同态
\begin{equation}
\tag{20.7.10.1}
\label{0.20.7.10.1-fr}
  \operatorname{Exalcotop}_{B/A}(C,L)\to
  \operatorname{Hom.\ cont}_C(\Upsilon_{C/B/A},L)
\end{equation}
（考虑 \sref{0.18.5.3.1-fr}）：对于每个开理想 $\mathfrak{M}$（属于
$A$）、每个开理想 $\mathfrak{N}$（属于 $B$）且
$\mathfrak{M}B\subset\mathfrak{N}$、每个开理想 $\mathfrak{P}$（属于 $C$）且
$\mathfrak{N}C\subset\mathfrak{P}$、$\mathfrak{P}.L=0$，每个
$(B/\mathfrak{N})$-扩充 $E$，它扩充 $C/\mathfrak{P}$、核为 $L$，并且是
$(A/\mathfrak{M})$-平凡的，都来自某个连续 $C$-同态 $f$；该同态把
$\Omega^1_{B/A}\otimes_B C$ 映入 $L$。同态 \sref{0.20.7.10.1-fr} 把
$\operatorname{Exalcotop}_{B/A}(C,L)$ 中 $E$ 的类的像，对应到 $f$ 在
$\Upsilon_{C/B/A}$ 上的限制；后者称为 $E$ 的\emph{特征同态}，仍记为
$\chi_E$。
\end{env}

\begin{proposition}[20.7.11]
\label{0.20.7.11-fr}
假设 $C$ 的拓扑满足开理想的平方仍为开理想。如果 $C$ 是形式光滑拓扑
$A$-代数，并且 $\Omega^1_{C/B}$ 是一个形式投射 $C$-模，那么对于每个离散
$C$-模 $L$，若它被 $C$ 的一个开理想零化，便有典范同构
\begin{equation}
\tag{20.7.11.1}
\label{0.20.7.11.1-fr}
  \operatorname{Exalcotop}_B(C,L)\xrightarrow{\sim}
  \operatorname{Hom.\ cont}_C(\Upsilon_{C/B/A},L).
\end{equation}
\end{proposition}

\begin{proof}
事实上，在正合列 \sref{0.20.3.7.1-fr} 中有
$\operatorname{Exalcotop}_A(C,L)=0$ \sref{0.19.4.4-fr}，所以
$\operatorname{Exalcotop}_B(C,L)=\operatorname{Exalcotop}_{B/A}(C,L)$，而同态
\sref{0.20.7.8.1-fr} 正是 \sref{0.20.7.10.1-fr}。考虑到
$\Omega^1_{C/B}$ 的拓扑此时由 $C$ 的拓扑导出 \sref{0.20.4.5-fr}，并利用
\sref{0.19.2.4-fr}，由 \sref{0.20.6.13-fr} 取归纳极限可知该同态为双射。
\end{proof}

\begin{env}[20.7.12]
\label{0.20.7.12-fr}
在 \sref{0.20.6.14-fr} 中，仍在 $\Upsilon^C_{B/A/\Lambda}$ 上赋予由
$\Omega^1_{A/\Lambda}\otimes_A C$ 的拓扑诱导的拓扑；当所考虑的环都是
拓扑环、环同态都连续时，同态 \sref{0.20.6.14.6-fr} 是连续的。
\end{env}

\begin{env}[20.7.13]
\label{0.20.7.13-fr}
如果在 \sref{0.20.6.23-fr} 中假设环 $\Lambda$、$A$、$B$、$C$ 都是拓扑环，
同态 $s$、$u$、$v$ 都连续，而 $L$ 是一个离散 $C$-模且被 $C$ 的一个开理想
零化，那么可以像 \sref{0.20.3.6-fr} 那样取归纳极限，得到正合列
\begin{equation}
\tag{20.7.13.1}
\label{0.20.7.13.1-fr}
\begin{split}
  0\to\operatorname{D\acute{e}r.\ cont}_B(C,L)\to
  \operatorname{D\acute{e}r.\ cont}_A(C,L)\to
  \operatorname{D\acute{e}r.\ cont}_A(B,L)\to \\
  \to\operatorname{Exalcotop}_{B/\Lambda}(C,L)\to
  \operatorname{Exalcotop}_{A/\Lambda}(C,L)\to
  \operatorname{Exalcotop}_{A/\Lambda}(B,L)\to0.
\end{split}
\end{equation}
\end{env}

\begin{env}[20.7.14]
\label{0.20.7.14-fr}
设 $A$ 是一个拓扑环，$B$ 是一个拓扑 $A$-代数，$\mathfrak{M}'$ 是 $A$ 的
一个开理想，$\mathfrak{N}'$ 是 $B$ 的一个开理想，并且
$\mathfrak{M}'B\subset\mathfrak{N}'$；置 $A'=A/\mathfrak{M}'$，
\oldpage[0\textsubscript{IV-1}]{151}
$B'=B/\mathfrak{N}'$。同态 $B\otimes_A B\to B'\otimes_{A'}B'$ 的核是
$\mathfrak{U}'=\operatorname{Im}(\mathfrak{N}'\otimes B+B\otimes\mathfrak{N}')$，
由此立即可知同态
\begin{equation}
\tag{20.7.14.1}
\label{0.20.7.14.1-fr}
  \varphi_{(\mathfrak{M}',\mathfrak{N}')}:
  \Omega^1_{B/A}\to\Omega^1_{B'/A'}
\end{equation}
的核为 $((\mathfrak{J}\cap\mathfrak{U}')+\mathfrak{J}^2)/\mathfrak{J}^2$；
另一方面，由 \sref{0.20.4.7-fr}，同态 \sref{0.20.7.14.1-fr} 是满射。若
$\mathfrak{M}''$（相应地，$\mathfrak{N}''$）是 $A$（相应地，$B$）的另一个
开理想，满足 $\mathfrak{M}''\subset\mathfrak{M}'$、
$\mathfrak{N}''\subset\mathfrak{N}'$ 且
$\mathfrak{M}''B\subset\mathfrak{N}''$，并置
$A''=A/\mathfrak{M}''$、$B''=B/\mathfrak{N}''$，则同样有满射同态
\[
  \varphi_{(\mathfrak{M}'',\mathfrak{N}''),(\mathfrak{M}',\mathfrak{N}')}:
  \Omega^1_{B''/A''}\to\Omega^1_{B'/A'}
\]
而这些同态显然组成一个\emph{投射系}。注意到各
$((\mathfrak{J}\cap\mathfrak{U}')+\mathfrak{J}^2)/\mathfrak{J}^2$ 组成
$0$ 在 $\Omega^1_{B/A}$ 中的一个基本邻域系，便可知\emph{分离完备化}
$\widehat{\Omega}^1_{B/A}$，也就是拓扑 $B$-模 $\Omega^1_{B/A}$ 的分离完备化，在相差典范
同构的意义下由下式给出：
\begin{equation}
\tag{20.7.14.2}
\label{0.20.7.14.2-fr}
  \widehat{\Omega}^1_{B/A}=\varprojlim(\Omega^1_{B'/A'}).
\end{equation}

此外，典范同态 $j:\Omega^1_{B/A}\to\widehat{\Omega}^1_{B/A}$ 是诸
$\varphi_{(\mathfrak{M}',\mathfrak{N}')}$ 所成投射系的投射极限，故
$\varphi_{(\mathfrak{M}',\mathfrak{N}')}$ 可分解为
$\Omega^1_{B/A}\xrightarrow{j}\widehat{\Omega}^1_{B/A}\to
\Omega^1_{B'/A'}$；由于它是满射，可知对每一对
$(\mathfrak{M}',\mathfrak{N}')$，典范同态
\begin{equation}
\tag{20.7.14.3}
\label{0.20.7.14.3-fr}
  \widehat{\Omega}^1_{B/A}\to\Omega^1_{B'/A'}
\end{equation}
都是满射。

在上述论述中，还可以处处把 $A'$ 换成 $A$。

最后，如果 $L$ 是一个\emph{分离且完备的}拓扑 $B$-模，则每个连续
$B$-同态，只要从 $\Omega^1_{B/A}$ 映入 $L$，都唯一延拓成一个连续
$\widehat B$-同态，它从 $\widehat{\Omega}^1_{B/A}$ 映入 $L$；反过来，把这样的
同态与 $\Omega^1_{B/A}\to\widehat{\Omega}^1_{B/A}$ 复合，又得到连续
$B$-同态。因此有典范同构
\[
  \operatorname{Hom.\ cont}_B(\Omega^1_{B/A},L)\xrightarrow{\sim}
  \operatorname{Hom.\ cont}_{\widehat B}(\widehat{\Omega}^1_{B/A},L).
\]
特别地，如果 $L$ 是一个被 $B$ 的某个开理想零化的离散 $B$-模，则典范
同构 \sref{0.20.4.8.2-fr} 还可以写成
\begin{equation}
\tag{20.7.14.4}
\label{0.20.7.14.4-fr}
  \operatorname{Hom.\ cont}_{\widehat B}(\widehat{\Omega}^1_{B/A},L)
  \xrightarrow{\sim}\operatorname{D\acute{e}r.\ cont}_A(B,L).
\end{equation}
\end{env}

\begin{proposition}[20.7.15]
\label{0.20.7.15-fr}
设 $A$ 是一个离散环，$B$ 是一个赋有进拓扑的拓扑 $A$-代数
（$\mathbf{0}_{\mathrm I}$, 7.1.9），$\mathfrak{n}$ 是 $B$ 的一个定义理想，并置
$B_0=B/\mathfrak{n}$。假设 $\Omega^1_{B_0/A}$ 和
$\mathfrak{n}/\mathfrak{n}^2$ 都是有限型 $B_0$-模。那么
$\widehat{\Omega}^1_{B/A}$ 是一个有限型 $\widehat B$-模。
\end{proposition}

\begin{proof}
由于 $B$ 的每个开理想的平方仍为开理想，$\Omega^1_{B/A}$ 的拓扑就是
$\mathfrak{n}$-预进拓扑 \sref{0.20.4.5-fr}。考虑到 $B$ 是一个进环，依据
（$\mathbf{0}_{\mathrm I}$, 7.2.7 和 7.2.9），只须证明
\oldpage[0\textsubscript{IV-1}]{152}
$\Omega^1_{B/A}/\mathfrak{n}.\Omega^1_{B/A}=
\Omega^1_{B/A}\otimes_B B_0$ 是一个有限型 $B_0$-模。然而，这由假设和正合列
\sref{0.20.5.12.1-fr}
\[
  \mathfrak{n}/\mathfrak{n}^2\to
  \Omega^1_{B/A}\otimes_B B_0\to\Omega^1_{B_0/A}\to0
\]
得出。
\end{proof}

\begin{env}[20.7.16]
\label{0.20.7.16-fr}
例如，当 $A$ 是一个域 $k$，$B=k'[[T_1,\ldots,T_n]]$ 是赋有通常拓扑的
形式幂级数代数，而 $k'$ 是 $k$ 的一个\emph{有限}扩张时，命题
\sref{0.20.7.15-fr} 适用（参见 \sref{0.21.9.2-fr}）。还应注意，分式域
$K$（即 $B$ 的分式域）在 $k$ 上有无限超越次数；当 $k$ 的特征为 $0$ 时，
由此立即推出（特别是利用 \sref{0.20.4.13-fr}, (i) 和
\sref{0.20.5.9-fr}，并且还利用特征为 $0$ 的域上的每个导子都能延拓到任一
扩张这一事实）：$\Omega^1_{B/k}$ \emph{不是有限型 $B$-模}。
\end{env}

\begin{env}[20.7.17]
\label{0.20.7.17-fr}
设 $A$、$B$、$C$ 是三个拓扑环，$u:A\to B$、$v:B\to C$ 是两个连续同态；
把 $A$、$B$、$C$ 分别换成关于开理想的商
$A'=A/\mathfrak{M}'$、$B'=B/\mathfrak{N}'$、$C'=C/\mathfrak{P}'$，其中
$u(\mathfrak{M}')\subset\mathfrak{N}'$、
$v(\mathfrak{N}')\subset\mathfrak{P}'$，于是有同态
$u':A'\to B'$、$v':B'\to C'$。由此得到典范同态
$u'_{C'/B'/A'}$、$v'_{C'/B'/A'}$；依据 \sref{0.20.5.4-fr}，它们组成
投射系，故取极限便得到以下典范同态；它们是正合列
\sref{0.20.5.7.1-fr}
\[
  \Omega^1_{B/A}\otimes_B C\xrightarrow{v_{C/B/A}}\Omega^1_{C/A}
  \xrightarrow{u_{C/B/A}}\Omega^1_{C/B}\to0:
\]
中各同态在分离完备化上的延拓：
\begin{align}
  \widehat v_{C/B/A}:\widehat{\Omega}^1_{B/A}
  \widehat\otimes_{\widehat B}\widehat C&\to\widehat{\Omega}^1_{C/A},
  \tag{20.7.17.1}\label{0.20.7.17.1-fr}\\
  \widehat u_{C/B/A}:\widehat{\Omega}^1_{C/A}&\to
  \widehat{\Omega}^1_{C/B}. \tag{20.7.17.2}\label{0.20.7.17.2-fr}
\end{align}
而在序列
\begin{equation}
\tag{20.7.17.3}
\label{0.20.7.17.3-fr}
  \widehat{\Omega}^1_{B/A}\widehat\otimes_{\widehat B}\widehat C
  \xrightarrow{\widehat v_{C/B/A}}\widehat{\Omega}^1_{C/A}
  \xrightarrow{\widehat u_{C/B/A}}\widehat{\Omega}^1_{C/B}\to0
\end{equation}
中，相邻两个同态的复合为 $0$，但该序列\emph{不一定正合}。不过，如果 $B$
和 $C$ 都\emph{可度量}，那么同态 $\widehat u_{C/B/A}$ \emph{是满射}，而
$\operatorname{Im}(\widehat v_{C/B/A})$ 在
$\operatorname{Ker}(\widehat u_{C/B/A})$ 中\emph{稠密}。这立即由如下事实得出
（参见（$\mathbf{0}_{\mathrm I}$, 7.3.1））：若
$(\mathfrak{M}_k)$（相应地，$(\mathfrak{N}_k)$）是 $B$（相应地，$C$）中
各理想组成的一个递减序列，并且构成 $0$ 的一个基本邻域系，且置
$B_k=B/\mathfrak{M}_k$、$C_k=C/\mathfrak{N}_k$，则过渡同态
$\Omega^1_{C_{k+1}/B_{k+1}}\to\Omega^1_{C_k/B_k}$、
$\Omega^1_{C_{k+1}/A}\to\Omega^1_{C_k/A}$ 和
$\Omega^1_{B_{k+1}/A}\to\Omega^1_{B_k/A}$ 都是满射
\sref{0.20.7.14-fr}。
\end{env}

\begin{proposition}[20.7.18]
\label{0.20.7.18-fr}
设 $A$、$B$、$C$ 是三个拓扑环，$u:A\to B$、$v:B\to C$ 是两个连续同态。
假设 $B$ 和 $C$ 都是可容且可度量的（$\mathbf{0}_{\mathrm I}$, 7.1.2）。
典范同态
\[
  \widehat v_{C/B/A}:\widehat{\Omega}^1_{B/A}
  \widehat\otimes_{\widehat B}\widehat C\to\widehat{\Omega}^1_{C/A}
\]
有一个连续 $C$-同态作为左逆的充分必要条件，是 $C$ 为一个 $B$-代数且
相对于 $A$ 形式光滑。
\end{proposition}

\begin{proof}
由 \sref{0.20.7.2-fr} 和 \sref{0.19.1.6-fr} 可知，该条件是必要的。为证明
充分性，注意拓扑 $C$-模
$L=\widehat{\Omega}^1_{B/A}\widehat\otimes_{\widehat B}\widehat C$ 由假设是
可度量且
\oldpage[0\textsubscript{IV-1}]{153}
完备的；因此，赋有乘积拓扑的 $E=D_C(L)$ 也是可度量且完备的。此外，它还是
可容的，因为若 $\mathfrak{R}$ 是 $C$ 的一个定义理想，则序列
$(\mathfrak{R}\oplus L)^n=\mathfrak{R}^n\oplus\mathfrak{R}^{n-1}L$ 趋于 $0$。
由于复合映射
\[
  D:B\xrightarrow{d_{B/A}}\Omega^1_{B/A}\to
  \widehat{\Omega}^1_{B/A}\widehat\otimes_{\widehat B}\widehat C=L
\]
是一个连续 $A$-导子，把 $B$ 映入 $L$，故连续 $A$-同态
\[
  f:x\longmapsto(v(x),D(x))
\]
从 $B$ 到 $E$，并在 $E$ 上定义了一个 $B$-扩充结构。由于 $L$ 是 $E$ 中的
闭理想，由 \sref{0.19.9.5-fr} 和假设可知，恒等映射 $C\to E/L$（它是一个
$B$-同态）可分解为 $C\xrightarrow{g}E\to E/L$，其中 $g$ 是满足
$g\circ v=f$ 的连续同态；于是由 \sref{0.20.1.3-fr}，$g$ 形如
$y\mapsto(y,D'(y))$，其中 $D'$ 是一个连续 $B$-导子，把 $C$ 映入 $L$，也就是说
$D'\circ v=D$。考虑到 \sref{0.20.7.14.4-fr}，有
$D'=h\circ\widehat d_{C/A}$，其中
$h:\widehat{\Omega}^1_{C/A}\to L$ 是连续 $C$-同态；但根据定义，有
$\widehat d_{C/A}\circ v=\widehat v_{C/B/A}\circ D$，而 $B$ 在 $D$ 下的像
（在拓扑意义下）生成 $C$-模 $L$ \sref{0.20.4.7-fr}，所以关系
$D'\circ v=D$ 确实给出 $h\circ\widehat v_{C/B/A}=1_L$。
\end{proof}

\begin{corollary}[20.7.19]
\label{0.20.7.19-fr}
在 \sref{0.20.7.18-fr} 的假设下，如果还假设 $B$ 和 $C$ 中每个开理想的
平方仍为开理想，那么 $C$ 为一个 $B$-代数且相对于 $A$ 形式光滑的充分必要条件，
是 $\widehat v_{C/B/A}$ 左可逆。
\end{corollary}

\begin{proof}
事实上，
$\widehat{\Omega}^1_{B/A}\widehat\otimes_{\widehat B}\widehat C$ 和
$\widehat{\Omega}^1_{C/A}$ 的拓扑此时都由 $C$ 的拓扑导出
\sref{0.20.4.5-fr}，而从其中一个到另一个的每个 $C$-同态都必定连续。
\end{proof}

\begin{env}[20.7.20]
\label{0.20.7.20-fr}
设 $A$ 是一个拓扑环，$B$ 是一个\emph{可度量且完备的}拓扑 $A$-代数，
$\mathfrak{R}$ 是 $B$ 的一个\emph{闭}理想，而 $C=B/\mathfrak{R}$ 是商拓扑环；
它也是可度量且完备的。设 $(\mathfrak{M}_k)$ 是由各理想组成的、$0$ 在 $B$ 中的
一个递减基本邻域系，并置
$B_k=B/\mathfrak{M}_k$、
$\mathfrak{R}_k=(\mathfrak{R}+\mathfrak{M}_k)/\mathfrak{M}_k$、
$C_k=B_k/\mathfrak{R}_k$。我们有一个同态投射系
\[
  \delta_{C_k/B_k/A}:\mathfrak{R}_k/\mathfrak{R}_k^2\to
  \Omega^1_{B_k/A}\otimes_{B_k}C_k
\]
\sref{0.20.5.11.3-fr}；取极限便由此得到典范同态
\[
  \widehat\delta_{C/B/A}:\mathfrak{R}/\mathfrak{R}^2\to
  \widehat{\Omega}^1_{B/A}\widehat\otimes_{\widehat B}\widehat C
\]
而且，像 \sref{0.20.7.17-fr} 那样论证，可知典范同态
$\widehat u_{C/B/A}:\widehat{\Omega}^1_{B/A}
\widehat\otimes_{\widehat B}\widehat C\to\widehat{\Omega}^1_{C/A}$
\emph{是满射}，并且 $\operatorname{Im}(\widehat\delta_{C/B/A})$ 在
$\operatorname{Ker}(\widehat u_{C/B/A})$ 中\emph{稠密}。
\end{env}

\section{特征 p 的环里的微分}
本节的结果比第 19、20 和 22 节的结果更为专门且更具技术性，在第 IV 章的论述中
只会偶尔使用。它们在这里的主要作用，是证明第 22 节的三个定理
（22.3.3、22.5.8 和 22.7.3）；其中第一个和最后一个定理在第 IV 章第 7 节关于
Noether 局部环的“精细”理论中起着实质性作用。
\oldpage[0\textsubscript{IV-1}]{154}
\subsection{$p$-生成元组和 $p$-基底}
\label{0.21.1-fr}

\begin{env}[21.1.1]
\label{0.21.1.1-fr}
给定一个数 $p$，它或者是 $0$，或者是一个素数。如果环 $A$ 的特征为 $p$，
我们的意思是存在一个环同态 $P\to A$，其中 $P$ 是特征为 $p$ 的素域；应注意，
这个同态必定唯一，因为复合
$\mathbf Z\to P\to A$ 是从 $\mathbf Z$ 到 $A$ 的唯一同态。
如果 $A\neq0$，这也等价于说 $A$ 包含一个特征为 $p$ 的域；$P$ 的像必定是
一个与 $P$ 同构的域（并且还是 $A$ 中唯一与 $P$ 同构的子域）。
\end{env}

\begin{env}[21.1.2]
\label{0.21.1.2-fr}
若 $p>0$，则说 $A$ 的特征为 $p$，等价于在 $A$ 中有 $p\cdot1=0$，也等价于
$pA=0$。若 $p=0$，则说 $A$ 的特征为 $p$，等价于对每个整数 $n\neq0$，
$n\cdot1$ 在 $A$ 中可逆。若 $A\neq0$，则数 $p$（素数或 $0$）至多有一个
能使 $A$ 的特征为 $p$；这由 Bézout 恒等式 $ap+bq=1$ 得出，其中 $p,q$
是两个不同素数。相反，零环 $0$ 对每个 $p$ 而言都有特征 $p$。
\end{env}

\begin{env}[21.1.3]
\label{0.21.1.3-fr}
若 $A$ 的特征为 $p$，则 $A$ 上的每个代数也具有同一特征。特别地，对每个
素理想 $\mathfrak p$（属于 $A$），$A$ 在 $\mathfrak p$ 处的剩余域都有特征 $p$。
反过来，若 $p=0$，并且对每个极大理想 $\mathfrak m$（属于 $A$），$A$ 在
$\mathfrak m$ 处的剩余域都有特征 $0$，那么 $A$ 也具有同一特征；因为对每个整数
$n\neq0$，$n\cdot1$ 此时在所有 $A_{\mathfrak m}$ 中都可逆，因而在 $A$ 中也
可逆。相反，若 $p>0$，一个局部环的剩余域可以有特征 $p$，而该局部环本身
并无特征 $p$；素环 $\mathbf Z_{p\mathbf Z}$（19.8.3）或 Artin 局部环
$\mathbf Z/p^n\mathbf Z$（其中 $n\geq2$）便给出例子，它们都不包含域。

最后应注意，一个环（即使是既约环）在各素理想处的剩余域也可能有不同的特征，
$\mathbf Z$ 就给出了这样的例子。
\end{env}

\begin{env}[21.1.4]
\label{0.21.1.4-fr}
在本节以下各处，我们固定一个素数 $p$；除非另有明言，所有环均假设有特征 $p$。
对这样的环 $A$，映射 $x\mapsto x^p$ 是 $A$ 的一个自同态，记为 $F_A$；若
$A$ 既约，则 $F_A$ 为单射。置 $A^p=F_A(A)$（这是 $A$ 的子环，由各
$x^p$（$x\in A$）组成）；自然可以把 $A$ 看成一个 $A^p$-代数。

还可以通过从 $A$ 到 $A$ 的同态 $F_A:x\mapsto x^p$，把 $A$ 看成一个
$A$-代数；换句话说，这就是代数 $A^{(p)}$，其中积 $\lambda\cdot x$，即
$x\in A$ 与标量 $\lambda\in A$ 的积，等于积 $\lambda^p x$，此积取在环
$A$ 中。显然，
对每个环同态 $u:A\to B$，偶对 $(u,u)$ 都是代数的双同态
$A^{(p)}\to B^{(p)}$。对每个 $A$-模 $E$，置
$E^{(p)}=E\otimes_A A^{(p)}$；这里把 $A^{(p)}$ 看成一个 $(A,A)$-双模，
其左 $A$-模结构是刚才定义的结构，而其右 $A$-模结构由 $A$ 中的乘法定义。
因此，对 $x\in E$、$\alpha,\beta\in A$，有
\[
  \alpha(x\otimes\beta)=x\otimes(\alpha\beta),\qquad
  (\alpha x)\otimes\beta=x\otimes\alpha^p\beta.
\]
\end{env}

\oldpage[0\textsubscript{IV-1}]{155}
置 $x^{(p)}=x\otimes1$，于是 $(\alpha x)^{(p)}=\alpha^p x^{(p)}$。当 $F_A$
为单射时，可以把 $A^{(p)}$ 等同于环 $A$，后者被看成其子环 $A^p$ 上的代数。

\begin{proposition}[21.1.5]
\label{0.21.1.5-fr}
设 $A$、$B$ 是两个环，$u:A\to B$ 是一个同态。
\begin{enumerate}
  \item[(i)] 每个 $A$-导子，若它把 $B$ 映入一个 $B$-模 $L$，就在 $A[B^p]$
  上为零（因此它是一个 $A[B^p]$-导子）。
  \item[(ii)] 对每个 $A$-子代数 $A'$，若它包含于 $A[B^p]$，且
  $j:A'\to B$ 是典范单射，则典范同态
  \[
  j_{B/A'/A}:\Omega^1_{B/A}\longrightarrow\Omega^1_{B/A'}
  \]
  为双射。
  \item[(iii)] 假设存在一个整数 $s\geq0$，使得
  $u(A)\supset B^{p^s}$。那么在环 $B\otimes_A B$ 中，
  $\mathfrak J_{B/A}$ 是一个幂零理想。
\end{enumerate}
\end{proposition}

\begin{proof}
(i) 从（20.1.1.1）出发归纳可知，对每个 $A$-导子 $D:B\to L$，都有
$D(x^n)=nx^{n-1}D(x)$；特别地，$D(x^p)=0$。

(ii) 考虑（20.4.8），断言 (ii) 只是 (i) 的另一种表述；后者写成
$\operatorname{Der}_{A[B^p]}(B,L)=\operatorname{Der}_A(B,L)$，对每个
$B$-模 $L$ 都成立。

(iii) 对每个 $x\in B$，有
$(x\otimes1-1\otimes x)^{p^s}=x^{p^s}\otimes1-1\otimes x^{p^s}=0$，
因为 $x^{p^s}\in u(A)$。结论由各元素
$1\otimes x-x\otimes1$ 生成 $\mathfrak J_{B/A}$（20.4.4）这一事实得出。
\end{proof}

由 (21.1.5) 立即可知，对每一对环同态 $A\to B\to C$，都有
\begin{equation}
\tag{21.1.5.1}
\label{0.21.1.5.1-fr}
\Upsilon_{C/B/A}=\Upsilon_{C/B/A'}
\end{equation}
其中对任意 $A$-子代数 $A'\subset A[B^p]$（它包含于 $B$），等式都成立。

另一方面，(21.1.5) 特别表明，对“绝对”微分模有
\begin{equation}
\tag{21.1.5.2}
\label{0.21.1.5.2-fr}
\Omega^1_A=\Omega^1_{A/A^p}.
\end{equation}

\begin{corollary}[21.1.6]
\label{0.21.1.6-fr}
假设 $B$ 是有限型 $A$-代数，并且存在一个整数 $s\geq0$，使得
$u(A)\supset B^{p^s}$。如果 $\Omega^1_{B/A}=0$，那么 $u$ 是满射。
\end{corollary}

\begin{proof}
\oldpage[0\textsubscript{IV-1}]{156}
根据 (21.1.5, (iii))，$\mathfrak J_{B/A}$ 是幂零理想；此外，由 (20.4.4)
和假设，$\mathfrak J_{B/A}$ 是有限型理想，因而是幂零的；关系
$\Omega^1_{B/A}=0$ 意味着
$\mathfrak J_{B/A}=\mathfrak J_{B/A}^2$，所以可得
$\mathfrak J_{B/A}=0$，也就是说，$\pi:B\otimes_A B\to B$ 为双射。
另一方面，$B$ 的每个元素的 $p^s$ 次幂都属于 $u(A)$，所以 $B$ 整于 $A$；
又因为它是有限型的，故它是有限 $A$-代数。因此，问题归结为证明下述引理
（其中并不假设所考虑的环有特征 $p$）：
\end{proof}

\begin{lemma}[21.1.6.1]
\label{0.21.1.6.1-fr}
设 $R$ 是一个环，$S$ 是一个有限 $R$-代数；若典范同态
$\pi:S\otimes_R S\to S$ 为双射，则结构同态 $u:R\to S$ 为满射。
\end{lemma}

\begin{proof}
只须证明，对每个极大理想 $\mathfrak m$（属于 $R$），若置
$T=R-\mathfrak m$，则同态
$u_{\mathfrak m}:R_{\mathfrak m}\to T^{-1}S$ 为满射（Bourbaki，
\emph{Alg. comm.}，第 II 章，§~3，n$^\circ$~3，定理 1）。事实上，假设蕴涵同态
$(T^{-1}S)\otimes_{R_{\mathfrak m}}(S/\mathfrak mS)\to T^{-1}S$
为双射（0$_I$, 1.3.4）；又因 $T^{-1}S$ 是有限
$R_{\mathfrak m}$-代数，可知只须考虑 $R$ 为局部环的情形。

仍用 $\mathfrak m$ 表示它的极大理想。只须证明
$u\otimes1:R/\mathfrak m\to S/\mathfrak mS$ 为满射；这由 Nakayama 引理得出，
因为 $S$ 是有限型 $R$-模。由于典范同态
$(S/\mathfrak mS)\otimes_{R/\mathfrak m}(S/\mathfrak mS)\to S/\mathfrak mS$
为双射，并且 $S/\mathfrak mS$ 是有限 $(R/\mathfrak m)$-代数，最终只须考虑
$R$ 为域的情形；但此时，$S\otimes_R S$ 与 $S$ 作为 $R$-模的秩只有在
$S$ 的秩为 $0$ 或 $1$ 时才可能相等，因而 $u:R\to S$ 为满射。
\end{proof}

\begin{proposition}[21.1.7]
\label{0.21.1.7-fr}
设 $A$ 是一个环，$B$ 是一个 $A$-代数，$(x_\alpha)_{\alpha\in I}$ 是 $B$
的一族元素。考虑以下性质：
\begin{enumerate}
  \item[a)] $B=A[B^p,(x_\alpha)]$；换句话说，$A[B^p]$-代数 $B$ 由族
  $(x_\alpha)_{\alpha\in I}$ 生成。
  \item[b)] $A[B^p]$-模 $B$ 由单项式
  $\prod_\alpha x_\alpha^{n(\alpha)}$ 生成；这里
  $(n(\alpha))_{\alpha\in I}$ 是一族具有有限支集的整数，并且
  $0\leq n(\alpha)<p$ 对每个 $\alpha\in I$ 都成立。
  \item[c)] $B$-模 $\Omega^1_{B/A}$ 由各
  $d_{B/A}(x_\alpha)$ $(\alpha\in I)$ 生成。
\end{enumerate}
那么性质 a) 与 b) 等价，并且它们蕴涵 c)；此外，如果 $B$ 是有限型
$A[B^p]$-代数，则 c) 与 a)、b) 等价。
\end{proposition}

\begin{proof}
显然，b) 蕴涵 a)。反过来，a) 也蕴涵 b)，因为每个单项式
$\prod_\alpha x_\alpha^{m(\alpha)}$（其中各 $m(\alpha)$ 是
$\geq0$ 的整数，并组成一族具有有限支集的族），都可以写成
$(\prod_\alpha x_\alpha^{q(\alpha)})^p\prod_\alpha x_\alpha^{r(\alpha)}$；
只须把每个 $m(\alpha)$ 除以 $p$，便得到
$m(\alpha)=p q(\alpha)+r(\alpha)$，其中 $0\leq r(\alpha)<p$。
a) 蕴涵 c) 由 (21.1.5, (ii)) 和 (20.4.7) 得出。

反过来，假设 c) 成立，并且 $B$ 是有限型 $A[B^p]$-代数；设 $B'$ 是一个
$A[B^p]$-子代数，它包含于 $B$ 且由各 $x_\alpha$ 生成。在正合列 (20.5.7.1)
\[
\Omega^1_{B'/A[B^p]}\otimes_{B'}B\longrightarrow
\Omega^1_{B/A[B^p]}\longrightarrow\Omega^1_{B/B'}\longrightarrow0
\]
中，考虑到 (21.1.5, (ii))，假设蕴涵左边的箭头为满射；因此
$\Omega^1_{B/B'}=0$。又因为 $B^p\subset B'\subset B$，并且 $B$ 是有限型
$B'$-代数，故由 (21.1.6) 必有 $B'=B$。
\end{proof}

\begin{remark}[21.1.8]
\label{0.21.1.8-fr}
当 $B$ 是域时，我们将在没有有限性假设的情形下证明性质 a)、b)、c) 等价
（21.4.5）。
\end{remark}

\begin{definition}[21.1.9]
\label{0.21.1.9-fr}
设 $A$ 是一个环（特征为 $p$），$B$ 是一个 $A$-代数。称一族
$(x_\alpha)_{\alpha\in I}$（由 $B$ 的元素组成）是 $p$-自由的，相对于 $A$
（相应地，是一个 $p$-生成元组，生成 $B$ 而相对于 $A$；相应地，是一个
$p$-基底，属于 $B$ 而相对于 $A$），其含义是
单项式族 $\prod_\alpha x_\alpha^{n(\alpha)}$（其中
$0\leq n(\alpha)<p$，且 $(n(\alpha))_{\alpha\in I}$ 具有有限支集）在
$A[B^p]$-模 $B$ 中是一个自由族（相应地，一个生成元系；相应地，一组基）。
\end{definition}

对一个集合 $M$，若其元素属于 $B$，也有相应的定义；只须考虑由典范单射 $M\to B$ 定义的
族。当取 $A$ 为素域 $\mathbf F_p$ 时（此时 $A[B^p]=B^p$），在上述定义中
省略 $A$（也可以说某族是“绝对”$p$-自由的，相应地，是“绝对”$p$-生成元组，
相应地，是“绝对”$p$-基底）。

\oldpage[0\textsubscript{IV-1}]{157}
\begin{lemma}[21.1.10]
\label{0.21.1.10-fr}
设 $C$ 是一个 $A$-子代数，它包含于 $B$ 且满足 $B^p\subset C$；$M$ 是 $C$ 的一个
子集，$N$ 是 $B$ 的一个子集。
\begin{enumerate}
  \item[(i)] 如果 $M$ 是一个 $p$-生成元组，生成 $C$ 而底环为 $A$，而 $N$
  是一个 $p$-生成元组，生成 $B$ 而底环为 $C$，那么 $M\cup N$ 是一个
  $p$-生成元组，生成 $B$ 而底环为 $A$。
  \item[(ii)] 假设 $M$ 是一个 $p$-基底，属于 $C$ 而相对于 $A$；那么 $N$
  是 $p$-自由的且相对于 $C$，其充分必要条件是 $M\cup N$ 是 $p$-自由的且
  相对于 $A$。
\end{enumerate}
\end{lemma}

\begin{proof}
只须考虑 $B^p\subset A\subset C\subset B$ 的情形。于是，(i) 是下述事实的
一个特例：如果 $P$（相应地，$Q$）是 $A$-模 $C$（相应地，$C$-模 $B$）的
一个生成元系，则所有积 $xy$（其中 $x\in P$、$y\in Q$）组成 $A$-模 $B$
的一个生成元系。

保留同样的记号。若 $P$ 是 $A$-模 $C$ 的一组基，则说 $Q$ 在 $C$ 上是
一个自由族，就是说关系
$\sum_{\lambda,\mu}a_{\lambda\mu}x_\lambda y_\mu=0$（其中
$a_{\lambda\mu}\in A$、$x_\lambda\in P$、$y_\mu\in Q$，各
$x_\lambda$（相应地，各 $y_\mu$）两两不同），等价于
$a_{\lambda\mu}=0$ 对每一对 $(\lambda,\mu)$ 都成立，也等价于
$\sum_\lambda a_{\lambda\mu}x_\lambda=0$ 对每个 $\mu$ 都成立；由此得到断言 (ii)。
\end{proof}

\subsection{$p$-基底与形式平滑性}
\label{0.21.2-fr}

\begin{theorem}[21.2.1]
\label{0.21.2.1-fr}
设 $B$ 是一个环，$A$ 是 $B$ 的一个子环，满足 $B^p\subset A\subset B$，
$(x_\alpha)_{\alpha\in I}$ 是一个 $p$-基底，它属于 $B$ 且相对于 $A$。
设 $E$ 是一个 $A$-代数，$u:A\to E$ 是结构同态，$q:E\to B$ 是一个满的
$A$-同态，$\mathfrak R$ 是它的核，并假设 $\mathfrak R^p=0$。那么：
\begin{enumerate}
  \item[(i)] 存在一个 $A$-同态 $v:B\to E$，它是同态 $q:E\to B$ 的右逆，
  其充分必要条件是，$u(q(z)^p)=z^p$ 对每个 $z\in E$ 都成立。
  \item[(ii)] 当 (i) 的条件满足时，对任意一族
  $(z_\alpha)_{\alpha\in I}$ 的 $E$ 中元素，若 $q(z_\alpha)=x_\alpha$ 对每个
  $\alpha\in I$ 都成立，则存在唯一的 $A$-同态 $v:B\to E$，使得
  $v(x_\alpha)=z_\alpha$ 对每个 $\alpha\in I$ 都成立；并且 $v$ 是 $q$ 的右逆。
\end{enumerate}
\end{theorem}

\begin{proof}
如果存在上述 $A$-同态 $v:B\to E$，那么必有 $v(a)=u(a)$，对每个
$a\in A\subset B$ 都成立；所以 $v(q(z)^p)=u(q(z)^p)$ 对每个 $z\in E$
都成立，因为 $B^p\subset A$。另一方面，有
$v(q(z)^p)=(v(q(z)))^p$；并且依定义，有
$v(q(z))\equiv z\pmod{\mathfrak R}$，只要 $q\circ v=1_B$；所以有
$(v(q(z)))^p=z^p$，因为 $\mathfrak R^p=0$。这就证明了 (i) 中条件的必要性；
其充分性将由 (ii) 得出。

现在，在 (ii) 的假设下，$v$ 的唯一性是显然的，因为各单项式
$\prod_\alpha x_\alpha^{n(\alpha)}$ 生成 $A$-模 $B$。此外，这些单项式组成
$A$-模 $B$ 的一组基，故存在一个 $A$-\emph{线性}映射 $v$，它从 $B$ 到 $E$，使得
$v(\prod_\alpha x_\alpha^{n(\alpha)})=\prod_\alpha z_\alpha^{n(\alpha)}$；
这里 $(n(\alpha))_{\alpha\in I}$ 是任意具有有限支集的族，并且
$0\leq n(\alpha)<p$ 对每个 $\alpha$ 都成立。还须证明 $v$ 是环同态。可以写成
$(\prod_\alpha x_\alpha^{m(\alpha)})(\prod_\alpha x_\alpha^{n(\alpha)})
=a\prod_\alpha x_\alpha^{r(\alpha)}$，其中
$r(\alpha)=m(\alpha)+n(\alpha)$，若 $m(\alpha)+n(\alpha)<p$；在相反情形下，
$r(\alpha)=m(\alpha)+n(\alpha)-p$；而 $a\in A$
\oldpage[0\textsubscript{IV-1}]{158}
是各 $x_\alpha^p$ 的积，其中 $\alpha\in I$ 满足
$m(\alpha)+n(\alpha)\geq p$。因此只须证明 $u(a)=\prod_\alpha z_\alpha^p$；但由于
$x_\alpha=q(z_\alpha)$，这正由假设得出。
\end{proof}

\begin{corollary}[21.2.2]
\label{0.21.2.2-fr}
设 $B$ 是一个环，$A$ 是 $B$ 的一个子环，满足 $B^p\subset A\subset B$。
如果存在一个 $p$-基底，它属于 $B$ 且相对于 $A$，那么 $B$ 是一个
$A$-代数，并且相对于 $B^p$ 形式光滑（各环均取离散拓扑）。
\end{corollary}

\begin{proof}
设 $E$ 是一个 $A$-扩充，它由 $B$ 被一个 $B$-模 $L$ 扩充而得；$q:E\to B$ 是增广，
$u:A\to E$ 是结构同态。说 $E$ 是 $B^p$-平凡的，意味着存在一个环同态
$v:B\to E$，使得 $q(v(b))=b$ 和 $v(b^p)=u(b^p)$ 对每个 $b\in B$ 都成立。
由此，对 $z\in E$，有 $u(q(z)^p)=v(q(z)^p)=(v(q(z)))^p$；另一方面，
由关系 $L^2=0$ 也有 $L^p=0$，而 $v(q(z))-z\in L$，故
$(v(q(z)))^p=z^p$。因此 (21.2.1, (i)) 的条件得到满足，且 $E$ 是
$B^p$-平凡的。
\end{proof}

\begin{corollary}[21.2.3]
\label{0.21.2.3-fr}
设 $B$ 是一个环，$A$ 是 $B$ 的一个子环，满足 $B^p\subset A\subset B$，
而 $\left(x_\alpha\right)_{\alpha\in I}$ 是一个 $p$-基底，它属于 $B$
且相对于 $A$。设 $L$ 是一个 $B$-模。那么：
\begin{enumerate}
  \item[(i)] 一个导子 $D$，它从 $A$ 到 $L$，能延拓为从 $B$ 到 $L$ 的导子，
  其充分必要条件是 $D$ 在 $B^p$ 上为零。
  \item[(ii)] 如果 $D$ 在 $B^p$ 上为零，那么对任意一族
  $(y_\alpha)_{\alpha\in I}$ 的 $L$ 中元素，存在唯一一个导子 $D'$，它从
  $B$ 到 $L$，延拓 $D$，并且 $D'(x_\alpha)=y_\alpha$ 对每个
  $\alpha\in I$ 都成立。
\end{enumerate}
\end{corollary}

\begin{proof}
给定一个导子 $D$，它从 $A$ 到 $L$，考虑环 $E=D_B(L)$，以及同态
$u:A\to E$，其定义为 $u(a)=(a,D(a))$。一个 $A$-同态 $v:B\to E$，
它是同态 $q:E\to B$ 的右逆，此时形如 $x\mapsto(x,D'(x))$，其中 $D'$ 是从
$B$ 到 $L$ 的一个导子，它延拓 $D$（20.1.5）。由于
$u(q(z)^p)=(q(z)^p,D(q(z)^p))$ 对 $z\in E$ 成立，该推论立即由把
(21.2.1) 应用于 $E$ 得出。
\end{proof}

由 (21.2.3) 可知序列
\begin{equation}
\tag{21.2.3.1}
\label{0.21.2.3.1-fr}
0\to\operatorname{Der}_A(B,L)\to\operatorname{Der}_{B^p}(B,L)
\to\operatorname{Der}_{B^p}(A,L)\to0
\end{equation}
（参见 (20.2.2.1)）是正合的，并且映射
\begin{equation}
\tag{21.2.3.2}
\label{0.21.2.3.2-fr}
D'\longmapsto(D'(x_\alpha))_{\alpha\in I}
\end{equation}
是 $B$-模 $\operatorname{Der}_A(B,L)$ 到积 $B$-模 $L^I$ 的一个同构
（在 (21.2.3) 中令 $D=0$）。

\begin{corollary}[21.2.4]
\label{0.21.2.4-fr}
在 (21.2.3) 的假设下，序列
\begin{equation}
\tag{21.2.4.1}
\label{0.21.2.4.1-fr}
0\to\Omega^1_{A/B^p}\otimes_A B\to\Omega^1_{B/B^p}
\to\Omega^1_{B/A}\to0
\end{equation}
是正合且可裂的，并且族 $(d_{B/A}(x_\alpha))_{\alpha\in I}$ 是 $B$-模
$\Omega^1_{B/A}$ 的一组基。
\end{corollary}

\begin{proof}
这立即由 (21.2.3) 以及公式
$\operatorname{Der}_A(B,L)=\operatorname{Hom}_B(\Omega^1_{B/A},L)$（20.4.8）得出。
\end{proof}

\begin{corollary}[21.2.5]
\label{0.21.2.5-fr}
设 $A$ 是一个环，$B$ 是一个 $A$-代数，$(x_\alpha)_{\alpha\in I}$ 是
一个 $p$-基底，它属于 $B$ 且相对于 $A$。那么族 $(d_{B/A}(x_\alpha))_{\alpha\in I}$
是 $B$-模 $\Omega^1_{B/A}$ 的一组基。
\end{corollary}

\begin{proof}
利用 (21.1.5, (ii))，可归结到 $A=A[B^p]$ 的情形；这时只须应用 (21.2.4)。
\end{proof}

\oldpage[0\textsubscript{IV-1}]{159}
\begin{env}[21.2.6]
\label{0.21.2.6-fr}
设 $A$、$B$ 是两个环，$u:A\to B$ 是一个环同态；采用 (21.1.4) 的记号，
有交换图
\begin{equation}
\tag{21.2.6.1}
\label{0.21.2.6.1-fr}
\xymatrix{
 A^{(p)} \ar[r]^{u} & B^{(p)}\\
 A \ar[u]^{F_A} \ar[r]^{u} & B \ar[u]_{F_B}
}
\end{equation}
由此典范地得到一个 $B$-代数同态
\begin{equation}
\tag{21.2.6.2}
\label{0.21.2.6.2-fr}
u\otimes1:A^{(p)}\otimes_A B\longrightarrow B^{(p)}
\end{equation}
它的像是环 $A[B^p]$（这是一个 $B$-代数，其结构来自同态 $F_B:B\to B^p$）。
\end{env}

\begin{theorem}[21.2.7]
\label{0.21.2.7-fr}
设 $A$ 是一个环，$B$ 是一个 $A$-代数，$u:A\to B$ 是结构同态。假设满足
以下条件：
\begin{enumerate}
  \item[(i)] 由 $u$ 得出的同态 (21.2.6.2) 是单射。
  \item[(ii)] $B$ 有一个 $p$-基底 $(x_\alpha)_{\alpha\in I}$，它相对于 $A$。
\end{enumerate}
那么 $B$ 是一个形式光滑的 $A$-代数（各环均取离散拓扑）。更确切地说，
给 $A$、$B$ 赋予离散拓扑；设 $E$ 是一个容许的拓扑 $A$-代数
（0$_I$, 7.1.2），$\mathfrak R$ 是 $E$ 的一个定义理想，
$C=E/\mathfrak R$，$v:B\to C$ 是一个 $A$-同态，而 $q:E\to C$ 是增广。
那么，对任意一族 $(z_\alpha)_{\alpha\in I}$ 的 $E$ 中元素，若
$q(z_\alpha)=v(x_\alpha)$ 对每个 $\alpha\in I$ 都成立，则存在唯一一个
$A$-同态 $w:B\to E$，使得 $v=q\circ w$，并且
$w(x_\alpha)=z_\alpha$ 对每个 $\alpha\in I$ 都成立。
\end{theorem}

\begin{proof}
先考虑 $E$ 为离散环的情形；于是 $\mathfrak R$ 是幂零的。可以只考虑 $v$
为满射的情形；此外，(19.4.3) 的论证使我们可以假设 $\mathfrak R^2=0$，
从而 $\mathfrak R^p=0$。最后，考虑由 $\mathfrak R$ 扩张 $C$ 所得的扩张 $E$
在 $v$ 下的逆像，便可以只考虑 $C=B$ 且 $v$ 为恒等映射的情形 (19.4.4)。
由于环同态 $F_E:z\rightsquigarrow z^p$ 在 $\mathfrak R$ 上为零，它分解为
$E\xrightarrow q B\xrightarrow{F'}E$；而将 $F'$ 看作从 $B$ 到 $E^{(p)}$ 的
同态时，按照借助复合同态 $A\xrightarrow r E\xrightarrow{F_E}E^{(p)}$
定义的 $E^{(p)}$ 的 $A$-代数结构，$F'$ 是一个\emph{$A$-同态}，其中
$r:A\to E$ 是 $A$-代数 $E$ 的结构同态；此外，$r:A^{(p)}\to E^{(p)}$
也同样是一个 $A$-同态。因此存在唯一一个 $A$-同态
$f:A^{(p)}\otimes_A B\to E^{(p)}$，使得 $f$ 与典范同态
$A^{(p)}\to A^{(p)}\otimes_A B$ 和 $B\to A^{(p)}\otimes_A B$ 的复合
分别为 $r$ 和 $F'$。而由假设，可以把 $A^{(p)}\otimes_A B$ 等同于
$A[B^p]$，此时典范同态 $A^{(p)}\to A^{(p)}\otimes_A B$ 和
$B\to A^{(p)}\otimes_A B$ 分别等同于 $u$ 和 $F_B$。因此，现在可以通过
同态 $f$ 把 $E$ 看成一个 $A[B^p]$-代数；而且依构造，对每个 $z\in E$
都有 $f(q(z)^p)=z^p$。于是 (21.2.1) 的适用条件得到满足，从而在这一情形下
定理得证。

为过渡到一般情形，考虑 $E$ 的开理想的一个基本系 $(\mathfrak J_\lambda)$，
并置 $E_\lambda=E/\mathfrak J_\lambda$、
$\mathfrak R_\lambda=(\mathfrak R+\mathfrak J_\lambda)/\mathfrak J_\lambda$、
$C_\lambda=E_\lambda/\mathfrak R_\lambda=E/(\mathfrak R+\mathfrak J_\lambda)$；
由于 $E$ 是容许的，有 $E=\varprojlim E_\lambda$（0$_I$, 7.2.2）。
对每一对 $(\alpha,\lambda)$，以 $z_{\alpha\lambda}$ 表示 $z_\alpha$ 在
$E_\lambda$ 中的典范像，以 $p_\lambda:C\to C_\lambda$ 表示典范同态，
\oldpage[0\textsubscript{IV-1}]{160}
以 $q_\lambda$ 表示典范同态 $E_\lambda\to C_\lambda$。
由于依假设 $\mathfrak R_\lambda$ 在 $E_\lambda$ 中是幂零的，证明的第一部分
证明了满足 $p_\lambda\circ v=q_\lambda\circ w_\lambda$ 以及对每个
$\alpha$ 都有 $w_\lambda(x_\alpha)=z_{\alpha\lambda}$ 的 $A$-同态
$w_\lambda:B\to E_\lambda$ 的存在性和唯一性。
$w_\lambda$ 的唯一性随即表明 $(w_\lambda)$ 是一个同态投射系，而
$w=\varprojlim w_\lambda$ 即为所求。
\end{proof}

\begin{remark}[21.2.8]
\label{0.21.2.8-fr}
同态 $w:B\to E$ 的存在性与唯一性（它满足
$w(x_\alpha)=z_\alpha$ 对每个 $\alpha$ 都成立），也可以直接从下列事实得出：
$B$ 是一个形式光滑的 $A$-代数，并且各 $d_{B/A}(x_\alpha)$ 组成 $B$-模
$\Omega^1_{B/A}$ 的一组基 (21.2.4)；无须利用它们构成 $p$-基底这一事实
（因此，即使不假设各环具有特征 $p$，该结果也成立）。事实上，只须考虑
$\mathfrak R^2=0$ 的情形；由于
$\operatorname{Der}_A(B,\mathfrak R)=\operatorname{Hom}_B(\Omega^1_{B/A},\mathfrak R)$
（20.4.8），所以存在唯一一个 $A$-导子 $D$，它从 $B$ 到 $\mathfrak R$，
并满足 $D(x_\alpha)=z'_\alpha$，其中 $(z'_\alpha)$ 是 $\mathfrak R$ 中
任意一族元素；结论由 (20.1.1) 得出。
\end{remark}

\subsection{$p$-基底与不恰模}
\label{0.21.3-fr}

\begin{env}[21.3.1]
\label{0.21.3.1-fr}
设 $A$、$B$ 是两个环（特征为 $p$），$i:A\to B$、$j:B\to A$ 是两个
环同态，满足
\begin{align}
\tag{21.3.1.1}\label{0.21.3.1.1-fr}
j(i(a))&=a^p\quad\text{对每个 }a\in A,\\
\tag{21.3.1.2}\label{0.21.3.1.2-fr}
i(j(b))&=b^p\quad\text{对每个 }b\in B.
\end{align}
通常 $i$ 是单射，因而 $A$ 通过 $i:A\to B$ 等同于 $B$ 的一个子环；作此
等同后，$j$ 的存在蕴涵 $B^p\subset A$，而 $j$ 此时等同于 $F_B$。
\end{env}

\begin{env}[21.3.2]
\label{0.21.3.2-fr}
如果 $h:A^p\to A$ 是典范单射，则由 (21.3.1.1) 和 (21.3.1.2) 有交换图
\begin{equation}
\tag{21.3.2.1}\label{0.21.3.2.1-fr}
\xymatrix{
 B \ar[r]^{j} & A\\
 A \ar[u]^{i} \ar[r]^{F_A} & A^p \ar[u]_{h}
}
\end{equation}
所以偶对 $(j,F_A)$ 可看成从 $A$-代数 $B$（结构由 $i$ 给出）到
$A^p$-代数 $A$（结构由 $h$ 给出）的一个双同态。由此得到 $A$-模的典范同态
\begin{equation}
\tag{21.3.2.2}\label{0.21.3.2.2-fr}
\pi_{B/A}:\Omega^1_{B/A}\otimes_B A_{[j]}\longrightarrow
\Omega^1_{A/A^p}=\Omega^1_A
\end{equation}
（参见 (20.5.4)；$\Omega^1_{A/A^p}$ 与绝对微分模 $\Omega^1_A$ 的等同来自
(21.1.5, (ii)））。置
\begin{align}
\tag{21.3.2.3}\label{0.21.3.2.3-fr}
\Theta_{B/A}&=\Omega^1_{B/A}\otimes_B A_{[j]},\\
\tag{21.3.2.4}\label{0.21.3.2.4-fr}
\Xi_{B/A}&=\operatorname{Ker}(\pi_{B/A}).
\end{align}
\oldpage[0\textsubscript{IV-1}]{161}
于是有正合列
\begin{equation}
\tag{21.3.2.5}\label{0.21.3.2.5-fr}
0\to\Xi_{B/A}\to\Theta_{B/A}\xrightarrow{\pi_{B/A}}\Omega^1_A
\end{equation}
（应注意，这些记号可能引起混淆，因为 $\Theta_{B/A}$ 和 $\Xi_{B/A}$ 不仅依赖于
$A$ 和 $B$，也依赖于 $i$ 和 $j$）。
\end{env}

\begin{env}[21.3.3]
\label{0.21.3.3-fr}
由于 (21.3.1.2) 给出 $F_B=i\circ j$，所以对每个 $B$-模 $M$ 都可以写成
（参见 (21.1.4)）
\begin{equation}
\tag{21.3.3.1}\label{0.21.3.3.1-fr}
M^{(p)}=M\otimes_B B^{(p)}=(M\otimes_B A_{[j]})\otimes_A B_{[i]},
\end{equation}
特别地，
\begin{equation}
\tag{21.3.3.2}\label{0.21.3.3.2-fr}
(\Omega^1_{B/A})^{(p)}=\Theta_{B/A}\otimes_A B_{[i]},
\end{equation}
因而从 (21.3.2.2) 得到 $B$-模的典范同态
\begin{equation}
\tag{21.3.3.3}\label{0.21.3.3.3-fr}
\pi_{B/A}\otimes1_B:(\Omega^1_{B/A})^{(p)}\longrightarrow
\Omega^1_A\otimes_A B_{[i]}.
\end{equation}
考虑到 (20.5.4.2)，这个同态的像是 $B$-模，它由各元素
$d_A(j(b))\otimes1$（其中 $b\in B$）生成。这些元素在典范同态
$i_{B/A}:\Omega^1_A\otimes_A B\to\Omega^1_B$（它由 $i$ 导出，见 (20.5.2)）
下的像，因而属于一个子-$B$-模，而这个模是 $\Omega^1_B$ 的子模，并由各
$d_B(i(j(b)))$（其中 $b\in B$）生成；
由 (21.3.1.2)，这个像为零。换言之，有同态序列
\begin{equation}
\tag{21.3.3.4}\label{0.21.3.3.4-fr}
0\to\Xi_{B/A}\otimes_A B_{[i]}\to(\Omega^1_{B/A})^{(p)}
\to\Upsilon_{B/A}
\end{equation}
它不一定正合，但任意两个相继同态的复合为零。
\end{env}

\begin{proposition}[21.3.4]
\label{0.21.3.4-fr}
如果 $B$ 是一个平坦 $A$-模（其结构由 $i$ 给出），则序列 (21.3.3.4) 正合。
\end{proposition}
\begin{proof}
这由平坦性的定义得出。特别地，当 $B^p\subset A\subset B$，$i$ 和 $j$
分别是典范单射和 $F_B$，并且 $B$ 有一个 $p$-基底且相对于 $A$ 时，该命题适用
（此时 $B$ 是一个自由 $A$-模）。但即使在这种情形下，核 $\Xi_{B/A}$ 也不一定
为零。不过有：
\end{proof}

\begin{proposition}[21.3.5]
\label{0.21.3.5-fr}
设 $B$ 是一个既约环，$A$ 是 $B$ 的一个子环，满足 $B^p\subset A\subset B$。
假设存在一个 $p$-基底 $(x_\lambda)_{\lambda\in L}$，它属于 $B$ 且相对于 $A$；
又存在一个 $p$-基底 $(y_\mu)_{\mu\in M}$，它属于 $A$ 且相对于 $B^p$。
那么 (21.3.3.4) 中的典范同态
\begin{equation}
\tag{21.3.5.1}\label{0.21.3.5.1-fr}
(\Omega^1_{B/A})^{(p)}\longrightarrow\Upsilon_{B/A}
\end{equation}
是双射，并且各元素 $d_A(x_\lambda^p)\otimes1$ 组成 $B$-模
$\Upsilon_{B/A}$ 的一组基。
\end{proposition}

\begin{proof}
\oldpage[0\textsubscript{IV-1}]{162}
由于 $B$ 是既约的，$x\mapsto x^p$ 是从 $B$ 到 $B^p$ 的一个结构同构；
借助这个同构传递结构，可知 $(x_\lambda^p)_{\lambda\in L}$ 与各 $y_\mu$
组成一个 $p$-基底，它属于 $B^p$ 且相对于 $A^p$ (21.1.10)。由 (21.1.5, (ii))
和 (21.2.5) 可知，各 $d_A(x_\lambda^p)$ 与各 $d_A(y_\mu)$ 组成 $A$-模
$\Omega^1_A$ 的一组基。现在，$d_A(x_\lambda^p)\otimes1$ 在
$i_{B/A}$ 下的像是 $d_B(x_\lambda^p)=0$；另一方面，$d_A(y_\mu)$ 在
$i_{B/A}$ 下的像是 $d_B(i(y_\mu))$。由于各 $y_\mu$ 与各 $x_\lambda$
组成一个 $p$-基底，它属于 $B$ 且相对于 $B^p$ (21.1.10)，各 $d_B(y_\mu)$ 是
$\Omega^1_B$ 的一组基的一部分 (21.2.5)，所以它们在 $B$ 上线性无关。
由此可知，$i_{B/A}$ 的核以各 $d_A(x_\lambda^p)\otimes1$ 为一组基；这些元素
是 $(\Omega^1_{B/A})^{(p)}$ 中相应元素在 (21.3.5.1) 下的像，而后者组成
$B$-模 $(\Omega^1_{B/A})^{(p)}$ 的一组基 (21.2.5)，所以 $B$-模同态
(21.3.5.1) 是双射。
\end{proof}

\begin{remarks}[21.3.6]
\label{0.21.3.6-fr}
\begin{enumerate}
  \item[(i)] 设 $A'$、$B'$ 是两个特征为 $p$ 的环，并给定两个同态
  $i':A'\to B'$、$j':B'\to A'$，它们满足 (21.3.1.1) 和 (21.3.1.2)；
  又给定环同态 $f:A\to A'$、$g:B\to B'$，使图表
  \[
  \xymatrix{
  A'\ar[r]^{i'}&B'\ar[r]^{j'}&A'\\
  A\ar[u]^f\ar[r]^i&B\ar[u]^g\ar[r]^j&A\ar[u]^f
  }
  \]
  交换。那么典范同态 $\Omega^1_{B/A}\to\Omega^1_{B'/A'}$（20.5.4.3）
  给出典范双同态 $\Theta_{B/A}\to\Theta_{B'/A'}$，使图表
  \[
  \xymatrix{
  0\ar[r]&\Xi_{B/A}\ar[r]\ar[d]&\Theta_{B/A}\ar[r]^{\pi_{B/A}}\ar[d]&\Omega^1_A\ar[d]\\
  0\ar[r]&\Xi_{B'/A'}\ar[r]&\Theta_{B'/A'}\ar[r]^{\pi_{B'/A'}}&\Omega^1_{A'}
  }
  \]
  交换。
  \item[(ii)] 设 $A'$ 是任意一个 $A$-代数，并置 $B'=B\otimes_A A'$；
  那么 $i'=i\otimes1:A'\to B'$ 与 $j'=j\otimes1:B'\to A'$ 满足
  (21.3.1.1) 和 (21.3.1.2)，并且
  $\Omega^1_{B'/A'}=\Omega^1_{B/A}\otimes_A A'=\Omega^1_{B/A}\otimes_B B'$
  (20.5.5)。由此得到典范的 $A'$-同构
  \begin{equation}
  \tag{21.3.6.1}\label{0.21.3.6.1-fr}
  \Theta_{B/A}\otimes_A A'\xrightarrow{\sim}\Theta_{B'/A'}
  \end{equation}
  以及典范的 $B'$-同态
  \begin{equation}
  \tag{21.3.6.2}\label{0.21.3.6.2-fr}
  (\Omega^1_{B/A})^{(p)}\otimes_B B'\xrightarrow{\sim}
  (\Omega^1_{B'/A'})^{(p)}.
  \end{equation}
\end{enumerate}
\end{remarks}

\subsection{域扩张的情形}
\label{0.21.4-fr}

\begin{env}[21.4.0]
\label{0.21.4.0-fr}
设 $K$ 是一个特征为 $p>0$ 的域，$k$ 是 $K$ 的一个子域；那么环
$k[K^p]$ 等于域 $k(K^p)$，因为 $k$ 在 $K^p$ 上是代数的。因此，可以应用
前面各号的结果，只须处处把 $A$、$B$ 和 $A[B^p]$ 分别替换为 $k$、$K$ 和
$k(K^p)$。
\end{env}

\begin{lemma}[21.4.1]
\label{0.21.4.1-fr}
设 $k$ 是一个特征为 $p>0$ 的域，$K$ 是 $k$ 的一个扩张。一个元素
$x\in K$ 为 $p$-自由且相对于 $k$，其充分必要条件是 $x\notin k(K^p)$。
\end{lemma}

\begin{proof}
事实上，$x$ 是多项式 $X^p-x^p$（属于 $k(K^p)[X]$）的一个根；已知
（Bourbaki，\emph{Alg.}，第 V 章，§~8，n$^\circ$~1，命题 1），如果
$x\notin k(K^p)$，则这个多项式不可约，因而各元素
$1,x,\ldots,x^{p-1}$ 组成 $k(K^p)$-模 $k(K^p)(x)$ 的一组基。
\end{proof}

\begin{theorem}[21.4.2]
\label{0.21.4.2-fr}
设 $k$ 是一个特征为 $p>0$ 的域，$K$ 是 $k$ 的一个扩张，$S$ 是一个
$p$-生成元组，它生成 $K$ 而相对于 $k$；$L\subset S$ 是一个 $p$-自由且
相对于 $k$ 的子集。那么存在一个 $p$-基底 $B$，它属于 $K$ 且相对于 $k$，
并满足 $L\subset B\subset S$。特别地，$k$ 的每个扩张都有一个
$p$-基底且相对于 $k$。
\end{theorem}

\begin{proof}
只须考虑 $K^p\subset k$ 的情形。由 Zorn 定理，$K$ 中存在一个子集 $B$，
使得 $L\subset B\subset S$；它为 $p$-自由且相对于 $k$，并且在 $S$ 中
所有具有这些性质的子集中为极大。只须证明子域 $K'$（属于 $K$），它由
$k$ 与 $B$ 生成，等于 $K$。否则会有 $x\in S$ 不属于 $K'$；由于
$K^p\subset k$，有 $K'=K'(K^p)$，所以 $x$ 为 $p$-自由且相对于 $K'$
(21.4.1)，从而 $B\cup\{x\}$ 为 $p$-自由且相对于 $k$ (21.1.10)，这与
$B$ 的定义矛盾。
(21.4.2) 的最后一个断言由取 $L=\varnothing$、$S=K$ 得出。
\end{proof}

\begin{corollary}[21.4.3]
\label{0.21.4.3-fr}
设 $k$ 是一个特征为 $p>0$ 的域，$K$ 是 $k$ 的一个扩张。一族
$(x_\alpha)$ 的 $K$ 中元素为 $p$-自由且相对于 $k$，其充分必要条件是，
对每个 $\alpha$，$x_\alpha$ 不属于域 $K_\alpha$；该域由 $k(K^p)$ 以及各
$x_\beta$（其指标满足 $\beta\neq\alpha$）生成。
\end{corollary}

\begin{proof}
由 (21.1.10) 可知该条件是必要的。反过来，假设该条件满足；只须考虑
$(x_\alpha)$ 是一个 $p$-生成元组且生成 $K$ 的情形。此时，$(x_\alpha)$
有一个子族，它是一个 $p$-基底且属于 $K$ (21.4.2)；但这个子族不可能不同于
$(x_\alpha)$，否则由假设，它便不是一个 $p$-生成元组且不生成 $K$。
\end{proof}

\oldpage[0\textsubscript{IV-1}]{163}
\begin{corollary}[21.4.4]
\label{0.21.4.4-fr}
设 $k$ 是一个特征为 $p>0$ 的域，$K$ 是 $k$ 的一个扩张。有
$\Omega^1_{K/k}=0$ 的充分必要条件是 $K=k(K^p)$。特别地，有
$\Omega^1_K=0$ 的充分必要条件是 $K$ 为完美域。
\end{corollary}

\begin{proof}
事实上，如果 $(x_\alpha)$ 是一个 $p$-基底，它属于 $K$ 且相对于 $k$，则各
$d_{K/k}(x_\alpha)$ 组成向量空间 $\Omega^1_{K/k}$ 的一组基 (21.2.5)。
\end{proof}

\begin{theorem}[21.4.5]
\label{0.21.4.5-fr}
设 $k$ 是一个特征为 $p>0$ 的域，$K$ 是 $k$ 的一个扩张，$(x_\alpha)$ 是
$K$ 中的一族元素。$(x_\alpha)$ 为 $p$-自由且相对于 $k$（相应地，是一个
$p$-生成元组，生成 $K$ 而相对于 $k$；相应地，是一个 $p$-基底，属于 $K$
且相对于 $k$），其充分必要条件是族 $(d_{K/k}(x_\alpha))$ 在向量空间
$\Omega^1_{K/k}$ 中是一个自由族（相应地，一个生成元系；相应地，一组基）。
\end{theorem}

\begin{proof}
\oldpage[0\textsubscript{IV-1}]{164}
设 $K'$ 是 $K$ 的子域（也就是子环），它由 $k(K^p)$ 与各 $x_\alpha$ 生成；
考虑到 (21.4.3)，只须证明对每个 $\alpha$，$x_\alpha$ 不属于 $K_\alpha$。
在正合列
\[
\Omega^1_{K_\alpha/k(K^p)}\otimes_{K_\alpha}K\longrightarrow
\Omega^1_{K/k(K^p)}\longrightarrow\Omega^1_{K/K_\alpha}\longrightarrow0,
\]
中，左边箭头把各 $d_{K_\alpha/k}(x_\beta)$（其中 $\beta\neq\alpha$）映为
$d_{K/k}(x_\beta)$；所以这些元素生成向量空间 $\Omega^1_{K/k(K^p)}$ 的一个
子空间，而它不包含 $d_{K/k}(x_\alpha)$。由于这个子空间是右边箭头的核，
可知 $d_{K/K_\alpha}(x_\alpha)\neq0$，从而 $x_\alpha\notin K_\alpha$。

现在，如果 $(x_\alpha)$ 是一个 $p$-自由族，则它是某个 $p$-基底的一个子族，
而该基底属于 $K$ 且相对于 $k$ (21.4.2)；所以各 $d_{K/k}(x_\alpha)$ 是向量空间
$\Omega^1_{K/k}$ 的一组基的一部分 (21.2.5)，从而在 $K$ 上线性无关。
反过来，如果各 $d_{K/k}(x_\alpha)$ 在 $K$ 上线性无关，那么前两段表明
$(x_\alpha)$ 为 $p$-自由且相对于 $k$。
\end{proof}

\begin{corollary}[21.4.6]
\label{0.21.4.6-fr}
设 $k$ 是一个特征为 $p>0$ 的域，$K$ 是 $k$ 的一个扩张，$x$ 是 $K$ 的一个
元素。以下三个条件等价：
\begin{enumerate}
  \item[a)] $x\notin k(K^p)$。
  \item[b)] $d_{K/k}(x)\neq0$。
  \item[c)] 元素 $x$ 为 $p$-自由且相对于 $k$。
\end{enumerate}
\end{corollary}

\begin{proposition}[21.4.7]
\label{0.21.4.7-fr}
设 $L$ 是一个特征为 $p>0$ 的域，$K$ 是 $L$ 的一个子域，满足
$L^p\subset K\subset L$。那么下列 $L$ 上向量空间的序列
\begin{equation}
\tag{21.4.7.1}\label{0.21.4.7.1-fr}
0\to\Omega^1_{K/L^p}\otimes_KL\to\Omega^1_{L/L^p}
\to\Omega^1_{L/K}\to0
\end{equation}
是正合的；换言之，有 $\Upsilon_{L/K/L^p}=0$。
\end{proposition}
\begin{proof}
考虑到 (21.4.2)，这是 (21.2.4) 的一个特例。
\end{proof}

\begin{proposition}[21.4.8]
\label{0.21.4.8-fr}
在 (21.4.7) 的假设下，典范同态
\begin{equation}
\tag{21.4.8.1}\label{0.21.4.8.1-fr}
(\Omega^1_{L/K})^{(p)}\longrightarrow\Upsilon_{L/K}
\end{equation}
是双射。如果 $(x_\lambda)$ 是一个 $p$-基底，它属于 $L$ 且相对于 $K$，则各元素
$d_K(x_\lambda^p)\otimes1$ 组成向量空间 $\Upsilon_{L/K}$ 的一组基。
\end{proposition}
\begin{proof}
这是 (21.3.5) 的一个特例。
\end{proof}

\subsection{应用：可分性判别法}
\label{0.21.5-fr}

在本号以及以下两号中，不再假设所考虑的环具有特征 $p>0$。

\begin{env}[21.5.1]
\label{0.21.5.1-fr}
先注意，判别法 (21.2.7) 能够证明 Cohen 关于可分扩张的定理 (19.6.1) 的一部分，
即：如果 $k$ 的特征为 $p>0$，而 $K$ 是 $k$ 的一个可分扩张，那么 $K$
是一个形式光滑的 $k$-代数。事实上，$K$ 有一个 $p$-基底且相对于 $k$
(21.4.2)；另一方面，由 MacLane 判别法（Bourbaki，\emph{Alg.}，第 V 章，
§~8，n$^\circ$~2，命题 3）可知，在 $K$ 的一个代数闭包中，$k^{p^{-1}}$
与 $K$ 在 $k$ 上线性无交；因而典范同态
$k^{p^{-1}}\otimes_k K\to k^{p^{-1}}(K)$ 是双射。这恰好是 (21.2.7) 的
条件 (i)，其中通过同构 $x\mapsto x^p$ 传递了结构。
\end{env}
\begin{proposition}[21.5.2]
\label{0.21.5.2-fr}
\textnormal{(MacLane).}
设 $A$、$B$ 是两个完备离散赋值环，$\mathfrak m$、$\mathfrak n$ 分别为其
极大理想，$K=A/\mathfrak m$ 和 $L=B/\mathfrak n$ 为其剩余域；设
$u:A\to B$ 是一个同态，满足 $Bu(\mathfrak m)=\mathfrak n$，并设
$u_0=u\otimes 1:K\to L$ 是相应的剩余域同态。考虑以下条件：
\begin{enumerate}
  \item[\rm a)] $L$ 是 $K$ 的一个可分扩张（相对于 $u_0$）。
  \item[\rm b)] 对任意完备离散赋值环 $B'$，若其极大理想为 $\mathfrak n'$，
  剩余域为 $L'$，
\oldpage[0\textsubscript{IV-1}]{165}
  则对任意同态 $u':A\to B'$，若 $B'u'(\mathfrak m)=\mathfrak n'$，以及
  对任意 $K$-同构 $\sigma:L\to L'$（相对于 $u_0$ 和
  $u'_0=u'\otimes1:K\to L'$），都存在一个同构 $w:B\to B'$，使得
  $u'=w\circ u$，并且 $w_0=w\otimes1:L\to L'$ 等于 $\sigma$。
  \item[\rm b$'$)] 对任意同态 $u':A\to B$，若
  $Bu'(\mathfrak m)=\mathfrak n$，且 $u'_0=u'\otimes1:K\to L$ 等于 $u_0$，
  则存在一个自同构 $w$，它作用在 $B$ 上，使得 $u'=w\circ u$，并且
  $w_0=w\otimes1:L\to L$ 为恒等映射。
  \item[\rm c)] 以 $u_1:A/\mathfrak m^2\to B/\mathfrak n^2$ 表示由 $u$
  取商所得的同态；那么，对任意局部同态
  $u'_1:A/\mathfrak m^2\to B/\mathfrak n^2$，若
  $\operatorname{gr}_0(u'_1)=\operatorname{gr}_0(u_1)(=u_0)$，则存在
  一个自同构 $w_1$，它作用在 $B/\mathfrak n^2$ 上，使 $\operatorname{gr}_0(w_1)$
  为恒等映射，并且 $u'_1=w_1\circ u_1$。
  \item[\rm c$'$)] 对任意局部同态 $u'_1:A/\mathfrak m^2\to
  B/\mathfrak n^2$，若 $\operatorname{gr}_0(u'_1)=\operatorname{gr}_0(u_1)$
  且 $\operatorname{gr}_1(u'_1)=\operatorname{gr}_1(u_1)$
  （这是从 $\mathfrak m/\mathfrak m^2$ 到 $\mathfrak n/\mathfrak n^2$ 的同态），
  则存在一个自同构 $w_1$，它作用在 $B/\mathfrak n^2$ 上，使
  $\operatorname{gr}_0(w_1)$ 为恒等映射，并且 $u'_1=w_1\circ u_1$。
\end{enumerate}
于是有蕴涵 $c)\Leftrightarrow c')\Leftrightarrow a)\Rightarrow b) \Rightarrow b')$。
如果还假设 $A$ 是一个 Cohen 环，则上述五个条件等价。

\end{proposition}
\begin{proof}
蕴涵 $b)\Rightarrow b')$ 和 $c)\Rightarrow c')$ 是显然的。来证明 $a)$ 蕴涵
$b)$。同态 $u$ 使 $B$ 成为一个无挠 $A$-模，因为依假设，它把 $A$ 的一个
素元变为 $B$ 的一个素元；由此可知 $B$ 是一个平坦 $A$-模
（0$_{\rm I}$, 6.3.4），因而是一个 Cohen $A$-代数 (19.8.1)。于是，
$a)$ 蕴涵 $b)$ 由把 (19.8.2, (i)) 应用于 $G=B'$、
$\mathfrak J=\mathfrak n'$ 得出；这里把 $B'$ 看成一个 $A$-代数（其结构同态为 $u'$），
而同态 $B\to B'/\mathfrak n'$ 是复合
$B\to B/\mathfrak n\to B'/\mathfrak n'$；这是一个 $A$-同态，因为假设给出
$u'_0=\sigma\circ u_0$。同一论证证明 $a)$ 蕴涵 $c)$；这次取
$G=B/\mathfrak n^2$、$\mathfrak J=\mathfrak n/\mathfrak n^2$，并把
$G$ 看成一个 $A$-代数，其结构同态为 $u'_1$。

其次证明 $c')$ 蕴涵 $a)$。两个同态 $u_1$ 和 $u'_1$ 满足
$u'_1=u_1+D$，其中 $D$ 是从 $A/\mathfrak m^2$ 到
$\mathfrak n/\mathfrak n^2$ 的一个导子 (20.1.1)。此外，假设
$\operatorname{gr}_1(u'_1)=\operatorname{gr}_1(u_1)$ 意味着 $D$ 在
$\mathfrak m/\mathfrak m^2$ 上为零，因而可以把它看成从
$K=A/\mathfrak m$ 到 $\mathfrak n/\mathfrak n^2$ 的一个导子；而
$\mathfrak n/\mathfrak n^2$（选择 $B$ 的一个素元后）等同于 $K$-模 $L$
（相对于 $u_0$）。另一方面，对 $w_1$ 施加的条件蕴涵
$\operatorname{gr}_1(w_1)$ 也为恒等映射（因为
$u_1(\mathfrak m/\mathfrak m^2)$ 生成 $\mathfrak n/\mathfrak n^2$）；
所以 $w_1$ 形如 $x\mapsto x+D'(x)$，其中 $D'$ 这一次是从
$B/\mathfrak n^2$ 到 $(B/\mathfrak n^2)$-模 $\mathfrak n/\mathfrak n^2$
的一个导子。由于 $w_1$ 在 $\mathfrak n/\mathfrak n^2$ 上为恒等映射，
$D'$ 仍可（通过此前 $\mathfrak n/\mathfrak n^2$ 与 $L$ 的等同）看成从
$L$ 到 $L$ 的一个导子。最后，关系 $u'_1=w_1\circ u_1$ 意味着 $D'$ 延拓
$D$。另一方面应注意，每个导子 $D$——从 $K$ 到 $L$——都对应于一个同态
$u'_1$，且该同态满足 $c')$ 的条件 (20.1.1)；所以条件 $c')$ 意味着，从 $K$ 到 $L$
的每个导子都可以延拓为从 $L$ 到自身的一个导子，也就是说 (20.6.5)，
$L$ 在 $K$ 上是可分的。

最后证明：当 $A$ 是 Cohen 环时，$b')$ 蕴涵 $c')$。先证明在 $c')$ 的假设下，
存在一个同态 $u':A\to B$，它满足 $b')$ 的假设，并且取商后给出
$u'_1:A/\mathfrak m^2\to B/\mathfrak n^2$。事实上，这由把 (19.8.6, (i))
应用于复合同态
$v':A\to A/\mathfrak m^2\xrightarrow{u'_1}B/\mathfrak n^2$
得出；
\oldpage[0\textsubscript{IV-1}]{166}
后者因而分解为 $A\xrightarrow{u'}B\to B/\mathfrak n^2$，而关于
$\operatorname{gr}_0(u'_1)$ 和 $\operatorname{gr}_1(u'_1)$ 的假设蕴涵
$\operatorname{gr}_0(u')=\operatorname{gr}_0(u)=u_0$ 以及
$\operatorname{gr}_1(u')=\operatorname{gr}_1(u)$；所以 $u'$ 把 $A$ 的一个
素元映为 $B$ 的一个素元，因而有 $Bu'(\mathfrak m)=\mathfrak n$。
于是，为得到符合要求的自同构 $w_1$，只须取一个自同构 $w$（它作用在 $B$
上）所诱导的商自同构；后者由把 $b')$ 应用于同态 $u'$ 得到。
\end{proof}
\begin{remarks}[21.5.3]
\label{0.21.5.3-fr}
\begin{enumerate}
  \item[(i)] 域的微分性质能够解决完备 Noether 局部环 $C$ 中代表域的唯一性问题
  (19.8.7, (ii))。设 $\mathfrak J$ 是 $C$ 的极大理想，$K=C/\mathfrak J$
  是 $C$ 的剩余域；可以只考虑 $\mathfrak J\ne0$ 的情形。假设存在一个同态
  $u:K\to C$，它与增广 $C\to K$ 的复合为恒等映射；那么，这个同态唯一的
  充分必要条件是 $\Omega^1_K=0$。事实上，条件 $\Omega^1_K=0$ 蕴涵 $K$
  在其素域上形式非分歧 (20.7.4)。如果还存在第二个符合要求的同态
  $v\ne u$，则存在一个最大的整数 $n$，使 $\mathfrak J^n$ 包含所有
  $u(x)-v(x)$（其中 $x\in K$）。取商后，$u$ 和 $v$ 将给出两个不同同态
  $u_1,v_1$，其定义域为 $K$，值域为 $C/\mathfrak J^{n+1}$，而它们与
  $C/\mathfrak J^{n+1}\to C/\mathfrak J^n$ 的复合相等，这与定义 (19.10.2)
  矛盾。反过来，假设 $\Omega^1_K\ne0$；此时存在一个非零导子 $D\ne0$，
  它从 $K$ 到 $\mathfrak J/\mathfrak J^2$ (20.4.8)，因而存在同态
  $v_1:K\to C/\mathfrak J^2$；若 $u_1:K\to C/\mathfrak J^2$ 是由 $u$
  取商所得的同态，则 $v_1=u_1+D$ (20.1.1)。如果 $k$ 是 $K$ 的素域，
  则 $C$ 是一个 $k$-代数，而 $K$ 是一个形式光滑的 $k$-代数 (19.6.1)；
  $u_1$ 和 $v_1$ 在 $k$ 上的限制相同，所以 $v_1$ 分解为
  $K\xrightarrow{v}C\to C/\mathfrak J^2$，并且 $v\ne u$。
  回忆（(20.6.20) 和 (21.4.4)），条件 $\Omega^1_K=0$ 意味着：如果 $K$
  的特征 $\ne0$，则它是完美的；在其余情形中，它是 $\mathbf Q$ 的代数扩张，
  并且特征为 $0$。
  \item[(ii)] 同样，设 $W$ 是一个 Cohen 环，$C$ 是一个完备 Noether 局部环，
  $\mathfrak J\ne0$ 是 $C$ 的一个包含于极大理想的理想。分解
  $W\to C\to C/\mathfrak J$（见 (19.8.6, (i)））唯一的充分必要条件是，
  剩余域 $K$（属于 Cohen 环 $W$）满足 $\Omega^1_K=0$。事实上，如果
  $\Omega^1_K\ne0$，而 $\mathfrak m$ 表示 $W$ 的极大理想，则只须把一个
  非零导子 $D\ne0$——它从 $K$ 到 $\mathfrak J/\mathfrak m\mathfrak J$——与
  增广 $W\to K$ 复合，便得到一个非零导子 $D'\ne0$，它从 $W$ 到
  $\mathfrak J/\mathfrak m\mathfrak J$；然后像 (i) 中
  那样完成论证：借助 $D'$ 构造同态
  $v_1:W\to C/\mathfrak m\mathfrak J$，它不同于同态
  $u_1:W\to C/\mathfrak m\mathfrak J$；后者由 $u$ 取商所得；最后调用 (19.8.6, (i))。
  反之，如果 $\Omega^1_K=0$，那么 $v$ 的唯一性在下列情形已经由 (i) 得出：
  $W=K$ 是特征为 $0$ 的域。在相反情形下，有 $\widehat{\Omega}^1_W=0$；
  事实上，此时 $\mathfrak m=pW$，其中 $p$ 是 $K$ 的特征 (19.8.5)，所以
  典范同态 $\mathfrak m/\mathfrak m^2\to\Omega^1_W/\mathfrak m\Omega^1_W$
  (20.5.11.2) 为零。把正合列
  (20.5.12.1) 应用于 $W$ 和 $K=W/\mathfrak m$，便得到
  $\Omega^1_W=\mathfrak m\Omega^1_W$，从而得到上述断言。但这时，由
  (20.7.4)，$W$ 在 $\mathbf Z$ 上形式非分歧（相对于它的 adic 拓扑），
  而 $v$ 的唯一性可像 (i) 中那样证明。
\end{enumerate}
\end{remarks}

\oldpage[0\textsubscript{IV-1}]{167}

\subsection{扩张的可容域}
\label{section:0.21.6-fr}

\begin{env}[21.6.1]
\label{0.21.6.1-fr}
给定四个域 $k_0\subset k\subset K\subset L$，由 (20.6.16) 和 (20.6.17)
可得正合列
\begin{equation}
\tag{21.6.1.1}
\label{0.21.6.1.1-fr}
0\to\Upsilon_{K/k/k_0}\otimes_KL\xrightarrow{v}\Upsilon_{L/k/k_0}
\xrightarrow{u}\Upsilon_{L/K/k_0}\xrightarrow{s}\Upsilon_{L/K/k}\to0.
\end{equation}

当保持 $k_0$、$K$ 和 $L$ 固定，而让中间域 $k$ 在 $k_0$ 与 $K$ 之间
“变动”时，显然有 $\Upsilon_{L/K/k_0}=\Upsilon_{L/K/k}$，只要
$k=k_0$。如果 (21.6.1.1) 中的典范同态 $s$
仍为双射，就称 $k$ 为一个 $k_0$-可容域，这里所指的是 $L$ 关于 $K$
的扩张。在某些条件下，存在这样的域 $k$，而它又“充分接近” $K$（例如
$[K:k]$ 有限）；这种存在性之所以有用，是因为在某些问题中，可以用较容易
处理微分模 $\Omega^1_{K/k_0}$ 和 $\Omega^1_{L/k_0}$（它们可能“过大”，
例如当 $k_0$ 是素域时），而改用 $\Omega^1_{K/k}$ 和 $\Omega^1_{L/k}$。

当 $k_0$ 是\emph{素域}时，将简称“可容域”，不再说“$k_0$-可容域”。

作如下记号：
\end{env}

\begin{env}[21.6.1.2]
\label{0.21.6.1.2-fr}
\[
\Delta(L/K,k/k_0)=\operatorname{Coker}
(\Upsilon_{K/k/k_0}\otimes_KL\to\Upsilon_{L/k/k_0})
\simeq\operatorname{Ker}(\Upsilon_{L/K/k_0}\to\Upsilon_{L/K/k})
\]
（这是一个 $L$-向量空间）；其秩记为 $d(L/K,k/k_0)$，称为相对于
$k_0$，$k$ 对扩张 $L$（相对于 $K$）的\emph{可容性缺陷}（显然，它为零
当且仅当 $k$ 对这一扩张是 $k_0$-可容的）。当 $k_0$ 是素域时，分别以
$\Delta(L/K,k)$ 和 $d(L/K,k)$ 代替 $\Delta(L/K,k/k_0)$ 和
$d(L/K,k/k_0)$。
\end{env}

\begin{proposition}[21.6.2]
\label{0.21.6.2-fr}
设 $k_0\subset k\subset K\subset L$ 是四个域。
\begin{enumerate}
  \item[\rm (i)] 下列条件等价：
  \begin{enumerate}
    \item[\rm a)] 域 $k$ 是 $k_0$-可容的，这里所指的是 $L$ 关于 $K$ 的扩张
    （换言之，同态 $\Upsilon_{L/K/k_0}\to
    \Upsilon_{L/K/k}$ 是单射，从而是双射）。
    \item[\rm b)] 典范同态
    $u:\Upsilon_{K/k/k_0}\to\Upsilon_{L/k/k_0}$ 为零。
    \item[\rm c)] 典范同态
    $v:\Upsilon_{K/k/k_0}\otimes_KL\to\Upsilon_{L/k/k_0}$ 是满射
    （从而是双射）。
    \item[\rm d)] 有 $d(L/K,k/k_0)=0$（或
    $\Delta(L/K,k/k_0)=0$）。
  \end{enumerate}
  \item[\rm (ii)] 在以下任一种情形中，(i) 的等价条件均成立：
  $\alpha)$ $L$ 在 $k$ 上可分；$\beta)$ $L$ 在 $K$ 上可分；
  $\gamma)$ 有 $k\subset k_0(K^p)$，其中 $p$ 表示 $k_0$ 的特征指数。
\end{enumerate}
\end{proposition}

\begin{proof}
\begin{enumerate}
  \item[(i)] 这些断言直接由序列 (21.6.1.1) 的正合性得出。
  \item[(ii)] 如果 $L$ 在 $k$ 上可分，则
  $\Upsilon_{L/K/k_0}=0$ (20.6.19)，故 (i) 的条件 b) 成立；如果 $L$
  在 $K$ 上可分，则 $\Upsilon_{L/K/k_0}=0$ (20.6.19)，故 (i) 的条件 a)
  成立；最后，如果 $k\subset k_0(K^p)$，那么由 (21.1.5.1) 可知
  $\Upsilon_{L/K/k_0}=\Upsilon_{L/K/k}$，因而 (i) 的条件 a) 成立。
\end{enumerate}
\end{proof}
\oldpage[0\textsubscript{IV-1}]{168}
\begin{env}[21.6.3]
\label{0.21.6.3-fr}
假设有如下由域的单同态组成的交换图：
\[
\xymatrix{
 k'_0\ar[r]&k'\ar[r]&K'\ar[r]&L'\\
 k_0\ar[u]\ar[r]&k\ar[u]\ar[r]&K\ar[u]\ar[r]&L\ar[u]
}
\]
由 (20.6.17, (ii)) 可得一个典范同态
\begin{equation}
\tag{21.6.3.1}
\label{0.21.6.3.1-fr}
\Delta(L/K,k/k_0)\longrightarrow\Delta(L'/K',k'/k'_0)
\end{equation}
它具有显然的传递性；因此，可以说 $\Delta(L/K,k/k_0)$ 是四元组
$(k_0,k,K,L)$ 的一个\emph{函子}。
\end{env}

\begin{proposition}[21.6.4]
\label{0.21.6.4-fr}
\begin{enumerate}
  \item[\rm (i)] 设 $k\subset k'\subset k''\subset K\subset L$ 是五个域。
  有典范同态的正合列
\begin{equation}
\tag{21.6.4.1}
\label{0.21.6.4.1-fr}
0\to\Delta(L/K,k'/k)\to\Delta(L/K,k''/k)\to\Delta(L/K,k''/k')\to0
\end{equation}
从而有等式
\begin{equation}
\tag{21.6.4.2}
\label{0.21.6.4.2-fr}
d(L/K,k''/k)=d(L/K,k''/k')+d(L/K,k'/k).
\end{equation}
  \item[\rm (ii)] 设 $k_0\subset k\subset K\subset L\subset M$ 是五个域。
  有典范同态的正合列
\begin{equation}
\tag{21.6.4.3}
\label{0.21.6.4.3-fr}
0\to\Delta(L/K,k/k_0)\otimes_LM\to\Delta(M/K,k/k_0)
\to\Delta(M/L,k/k_0)\to0
\end{equation}
从而有等式
\begin{equation}
\tag{21.6.4.4}
\label{0.21.6.4.4-fr}
d(M/K,k/k_0)=d(M/L,k/k_0)+d(L/K,k/k_0).
\end{equation}
\end{enumerate}
\end{proposition}

\begin{proof}
\begin{enumerate}
  \item[(i)] 考虑交换图
\[
\xymatrix{
0\ar[r]&\Upsilon_{K/k'/k}\otimes_KL\ar[r]\ar[d]&
\Upsilon_{K/k''/k}\otimes_KL\ar[r]\ar[d]&
\Upsilon_{K/k''/k'}\otimes_KL\ar[r]\ar[d]&0\\
0\ar[r]&\Upsilon_{L/k'/k}\ar[r]\ar[d]&
\Upsilon_{L/k''/k}\ar[r]\ar[d]&
\Upsilon_{L/k''/k'}\ar[r]\ar[d]&0\\
0\ar[r]&\Delta(L/K,k'/k)\ar[r]&\Delta(L/K,k''/k)\ar[r]&
\Delta(L/K,k''/k')\ar[r]&0
}
\]
可以分别把三行看成\emph{复形} $T_1,T_2,T_3$；正合列 (21.6.1.1)
和 $\Delta$ 的定义 (21.6.1.2) 表明，有复形的正合列
$0\to T_1\to T_2\to T_3\to0$；对它应用上同调正合列，并注意到，由
\oldpage[0\textsubscript{IV-1}]{169}
(21.6.1.1) 的正合性，$T_1$ 和 $T_2$ 的上同调除在同一个次数外均为零，
而在该次数，其上同调模都等于 $\Upsilon_{k''/k'/k}\otimes_{k''}L$；由于
$T_1$ 和 $T_2$ 具有相同的上同调，$T_3$ 的上同调必为零，这就证明了 (i)。
  \item[(ii)] 同样考虑交换图
\[
\xymatrix{
0\ar[r]&\Delta(L/K,k/k_0)\otimes_LM\ar[r]\ar[d]&
\Delta(M/K,k/k_0)\ar[r]\ar[d]&\Delta(M/L,k/k_0)\ar[r]&0\\
0\ar[r]&\Upsilon_{L/K/k_0}\otimes_LM\ar[r]\ar[d]&
\Upsilon_{M/K/k_0}\ar[r]\ar[d]&\Upsilon_{M/L/k_0}\ar[r]&0\\
0\ar[r]&\Upsilon_{L/K/k}\otimes_LM\ar[r]&
\Upsilon_{M/K/k}\ar[r]&\Upsilon_{M/L/k}\ar[r]&0
}
\]
这里仍把三行看成\emph{复形} $T'_1,T'_2,T'_3$；正合列 (21.6.1.1)
和 $\Delta$ 的定义 (21.6.1.2) 在这里给出复形的正合列
$0\to T'_1\to T'_2\to T'_3\to0$，再对它应用上同调正合列；这一次，
由 (21.6.1.1) 的正合性，$T'_2$ 和 $T'_3$ 的上同调除在同一个次数外均为零，
而在该次数，其上同调模都等于 $\Upsilon_{M/L/k_0}$；因此，$T'_1$ 的上同调
为零，这就证明了 (ii)。
\end{enumerate}
\end{proof}

\begin{corollary}[21.6.5]
\label{0.21.6.5-fr}
\begin{enumerate}
  \item[(i)] 给定五个域 $k\subset k'\subset k''\subset K\subset L$，
  $k''$ 为 $k$-可容域（相对于 $L$ 关于 $K$ 的扩张）的充分必要条件是：
  $k'$ 为 $k$-可容域，并且 $k''$ 为 $k'$-可容域（仍相对于 $L$ 关于 $K$
  的扩张）。
  \item[(ii)] 给定五个域
  $k_0\subset k\subset K\subset L\subset M$，$k$ 为一个 $k_0$-可容域
  （相对于 $M$ 关于 $K$ 的扩张）的充分必要条件是：它对 $L$ 关于 $K$ 的扩张
  以及 $M$ 关于 $L$ 的扩张都是如此。
\end{enumerate}

这立即由关系式 (21.6.4.2) 和 (21.6.4.4) 得出，因为 $d$ 的值均为
$\geqslant0$。
\end{corollary}

\begin{corollary}[21.6.6]
\label{0.21.6.6-fr}
设 $k_0\subset k\subset K\subset L$ 是四个域，并假设 $k$ 是一个
$k_0$-可容域，相对于 $L$ 关于 $K$ 的扩张。于是，如果 $k'_0$、$k'$、$K'$、$L'$
是四个域，并且
$k_0\subset k'_0\subset k'\subset k\subset K\subset K'\subset L'\subset L$，
那么 $k'$ 是一个 $k'_0$-可容域，相对于 $L'$ 关于 $K'$ 的扩张。
\end{corollary}

\subsection{Cartier 恒等式}
\label{subsection:0.21.7-fr}

以下结果用微分的语言表述了 MacLane 关于导子的一个定理：

\oldpage[0\textsubscript{IV-1}]{170}
\begin{theorem}[21.7.1]
\label{0.21.7.1-fr}
\textnormal{(Cartier).}
设 $K$ 是一个域，$L$ 是 $K$ 的一个有限型扩张。于是
$\Omega^1_{L/K}$ 和 $\Upsilon_{L/K}$ 都是有限秩的 $L$-向量空间，并且有
\begin{equation}
\tag{21.7.1.1}
\label{0.21.7.1.1-fr}
\operatorname{rg}\Omega^1_{L/K}-\operatorname{rg}\Upsilon_{L/K}
=\deg.\operatorname{tr}_K L.
\end{equation}
\end{theorem}

\begin{proof}
如果 $L'$ 是满足 $K\subset L'\subset L$ 的一个域，则有正合列
(20.6.15.1)（其中取 $\Lambda=P$，即 $K$ 的素域，$A=K$、
$B=L'$、$C=L$）：
\[
0\to\Upsilon_{L'/K}\otimes_{L'}L\to\Upsilon_{L/K}
\to\Omega^1_{L'/K}\otimes_{L'}L\to\Omega^1_{L/K}
\to\Omega^1_{L/L'}\to0
\]
由此得到（0$_{\rm III}$, 11.10.2）
\[
\operatorname{rg}_L\Omega^1_{L/K}-\operatorname{rg}_L\Upsilon_{L/K}
=(\operatorname{rg}_{L'}\Omega^1_{L/L'}-\operatorname{rg}_{L'}\Upsilon_{L/L'})
+(\operatorname{rg}_{L'}\Omega^1_{L'/K}-\operatorname{rg}_{L'}\Upsilon_{L'/K}).
\]
另一方面，
$\deg.\operatorname{tr}_K L=\deg.\operatorname{tr}_K L'
+\deg.\operatorname{tr}_{L'}L$；因此，对扩张 $L$ 的生成元数作归纳，可将
(21.7.1.1) 的证明归结为 $L=K(x)$ 的情形。此时分三种情形：
\begin{enumerate}
  \item[\rm a)] $x$ 在 $K$ 上超越；由于 $L$ 在 $K$ 上可分，有
  $\Upsilon_{L/K}=0$ (20.6.3)；另一方面，(20.5.9) 和 (20.4.13) 表明
  $\Omega^1_{L/K}$ 的秩为 $1$，从而此时 (21.7.1.1) 成立。
  \item[\rm b)] $L$ 是 $K$ 的一个可分代数扩张，因而仍有
  $\Upsilon_{L/K}=0$ (20.6.3)。另一方面，由 (20.6.20) 有
  $\Omega^1_{L/K}=0$，所以 (21.7.1.1) 仍成立。
  \item[\rm c)] $L$ 是 $K$ 的一个不可分代数扩张；开头的论证表明，可以只考虑
  $L^p\subset K$ 的情形；于是由 (21.4.8) 得到
  $\operatorname{rg}_L\Upsilon_{L/K}=\operatorname{rg}_L\Omega^1_{L/K}$，
  所以 (21.7.1.1) 仍成立。
\end{enumerate}
\end{proof}

\begin{corollary}[21.7.2]
\label{0.21.7.2-fr}
设 $K$ 是一个域，$L$ 是 $K$ 的一个有限型扩张，$k$ 是 $K$ 的一个子域。
于是 $\Omega^1_{L/K}$ 和 $\Upsilon_{L/K/k}$ 都是有限秩的 $L$-向量空间，
并且有
\begin{equation}
\tag{21.7.2.1}
\label{0.21.7.2.1-fr}
\operatorname{rg}\Omega^1_{L/K}-\operatorname{rg}\Upsilon_{L/K/k}
=\deg.\operatorname{tr}_K L+d(L/K,k).
\end{equation}
因而有不等式
\begin{equation}
\tag{21.7.2.2}
\label{0.21.7.2.2-fr}
\operatorname{rg}\Omega^1_{L/K}-\operatorname{rg}\Upsilon_{L/K/k}
\geqslant\deg.\operatorname{tr}_K L.
\end{equation}
此外，(21.7.2.2) 两边相等的充分必要条件是，$k$ 是一个可容域，相对于
$L$ 关于 $K$ 的扩张。
\end{corollary}

\begin{proof}
事实上，由于同态
$\Upsilon_{L/K}\to\Upsilon_{L/K/k}$ 是满射，按定义 (21.6.1.2) 有
$\operatorname{rg}\Upsilon_{L/K}-\operatorname{rg}\Upsilon_{L/K/k}
=d(L/K,k)$，所以该推论立即由 (21.7.1) 和 (21.6.2) 得出。
\end{proof}

\begin{corollary}[21.7.3]
\label{0.21.7.3-fr}
设 $k\subset K\subset L$ 是三个域，并且 $L$ 是 $k$ 的一个有限型扩张。
于是有
\begin{equation}
\tag{21.7.3.1}
\label{0.21.7.3.1-fr}
\operatorname{rg}_L\Omega^1_{L/k}-\operatorname{rg}_K\Omega^1_{K/k}
=\deg.\operatorname{tr}_K L+d(L/K,k).
\end{equation}
\oldpage[0\textsubscript{IV-1}]{171}
因而有不等式
\begin{equation}
\tag{21.7.3.2}
\label{0.21.7.3.2-fr}
\operatorname{rg}_L\Omega^1_{L/k}-\operatorname{rg}_K\Omega^1_{K/k}
\geqslant\deg.\operatorname{tr}_K L.
\end{equation}
并且，等号成立的充分必要条件是，$k$ 是一个可容域，相对于 $L$ 关于 $K$
的扩张。
\end{corollary}

\begin{proof}
事实上，有正合列 (20.6.1.1)
\[
0\to\Upsilon_{L/K/k}\to\Omega^1_{K/k}\otimes_KL\to\Omega^1_{L/k}
\to\Omega^1_{L/K}\to0
\]
所以 (21.7.3.1) 的左端等于
\[
\operatorname{rg}_L\Omega^1_{L/k}-\operatorname{rg}_L\Upsilon_{L/K/k}
\]
只须应用 (21.7.2) 即得结论。
\end{proof}

最后这个推论的意义在于，它只涉及微分模，而完全不涉及不恰模。下文还将看到
(21.8.6)，对于每一个有限型扩张 $L$（它是 $K$ 的扩张），都存在一个子域
$k$（它是 $K$ 的子域），
使 $[K:k]$ 有限，并且该子域是可容的，相对于 $L$ 关于 $K$ 的扩张；此时
(21.7.3.2) 的两端相等。

\begin{corollary}[21.7.4]
\label{0.21.7.4-fr}
设 $K$ 是域 $k$ 的一个有限型扩张。下列条件等价：
\begin{enumerate}
  \item[\rm a)] $K$ 是 $k$ 的有限可分扩张。
  \item[\rm b)] $K$ 是一个形式非分歧的 $k$-代数 (19.10.2)。
  \item[\rm c)] $K$ 是一个形式 étale 的 $k$-代数 (19.10.2)。
  \item[\rm d)] 有 $\Omega^1_{K/k}=0$。
\end{enumerate}
\end{corollary}

\begin{proof}
事实上，已知（这里拓扑均为离散拓扑）条件 b) 与 d) 等价 (20.7.4)，而 a)
蕴涵 $K$ 是一个形式光滑的 $k$-代数 (19.6.1)，所以 a) 与 b) 的合取等价于
c)。因此只需证明 d) 蕴涵 a)。但是，由 (21.7.1.1)，关系
$\Omega^1_{K/k}=0$ 蕴涵 $K$ 在 $k$ 上是代数的（因而是有限的），并且
$\Upsilon_{K/k}=0$；换言之，由 (20.6.3)，$K$ 在 $k$ 上可分。
\end{proof}

\begin{remark}[21.7.5]
\label{0.21.7.5-fr}
由 (20.6.3.2) 和（0$_{\rm III}$, 11.10.2），(21.7.1.1) 的左端正是复形
$K_\bullet(C/B/A)$ 的 \emph{Euler--Poincaré 示性数}；这个复形在 (20.6.3)
中引入，这里取 $C=L$、$B=K$ 和 $A=P$（即 $K$ 的素域）。在讨论交叉理论与
Riemann--Roch 定理的章节中，取值于 $\mathcal O_X$-模类群的推广
Euler--Poincaré 示性数也将起重要作用；它们属于推广复形，而这些复形推广了
(20.6.22) 中考虑的 $F_\bullet(C/A)$。
\end{remark}

\subsection{可容性判别法}
\label{subsection:0.21.8-fr}

现在回到先前的约定；因此，本号中所考虑的所有域都假设具有特征 $p>0$。

\begin{lemma}[21.8.1]
\label{0.21.8.1-fr}
设 $K$ 是一个域，$k$ 是 $K$ 的一个子域，$(k_\alpha)_{\alpha\in I}$ 是
$K$ 的一个滤过递减子域族，并且
$k=\bigcap_{\alpha\in I}k_\alpha$。设 $V$ 是一个 $K$-向量空间，
$(a_i)_{1\leqslant i\leqslant n}$ 是 $V$ 中的一个有限向量族；如果族
$(a_i)$ 在 $k$ 上线性无关，则存在一个指标 $\gamma$，使它在 $k_\gamma$
上也线性无关。
\end{lemma}

\oldpage[0\textsubscript{IV-1}]{172}
\begin{proof}
设 $r$ 是族 $(a_i)_{1\leqslant i\leqslant n}$ 在 $K$ 上的秩，对 $n-r$
作归纳；当 $n=r$ 时断言显然成立，因为此时族
$(a_i)_{1\leqslant i\leqslant r}$ 在 $K$ 上线性无关，因而在 $K$ 的每个子域
上都线性无关。例如，假设 $(a_i)_{1\leqslant i\leqslant r}$ 在 $K$ 上
线性无关，并写成
$a_{r+1}=\sum_{i=1}^r\lambda_i a_i$，其中 $\lambda_i\in K$；由于族
$(a_i)_{1\leqslant i\leqslant r+1}$ 在 $k$ 上线性无关，诸 $\lambda_i$
不可能全属于 $k$；例如假设
$\lambda_{i_0}\notin k$。于是存在一个指标 $\beta$，使
$\lambda_{i_0}\notin k_\beta$；由此可知族
$(a_i)_{1\leqslant i\leqslant r+1}$ 在 $k_\beta$ 上线性无关；事实上，族
$(a_i)_{1\leqslant i\leqslant r}$ 在 $K$ 的每个子域上线性无关，若族
$(a_i)_{1\leqslant i\leqslant r+1}$ 在 $k_\beta$ 上不是线性无关的，
$a_{r+1}$ 就等于诸 $a_i$ 的一个线性组合，其系数属于 $k_\beta$；这些系数
必为诸 $\lambda_i$，于是得到矛盾。现在只须应用归纳假设：把 $K$ 换成
$k_\beta$，并把族 $(k_\alpha)_{\alpha\in I}$ 换成由其中的 $k_\alpha$
（它们包含于 $k_\beta$）组成的子族。
\end{proof}

\begin{lemma}[21.8.2]
\label{0.21.8.2-fr}
设 $K$ 是一个域，$p$ 是它的特征指数，$k_0$ 是 $K$ 的一个子域，
$(k_\alpha)_{\alpha\in I}$ 是 $K$ 的一个滤过递减子域族，并且
$\bigcap_{\alpha\in I}k_\alpha(K^p)=k_0(K^p)$。如果
$(x_i)_{1\leqslant i\leqslant n}$ 是 $K$ 的一个有限元素族，并且是
$p$-自由的（相对于 $k_0$），则存在一个指标 $\alpha$，使 $(x_i)$ 仍为
$p$-自由的（相对于 $k_\alpha$）。

事实上，族 $(x_i)$ 为 $p$-自由的（相对于 $k$，它是 $K$ 的一个子域），
是指有限单项式族
$\prod_{i=0}^n x_i^{m(i)}$（其中 $0\leqslant m(i)<p$）在 $k(K^p)$ 上
线性无关；因此，只须把引理 (21.8.1) 应用于向量空间 $V=K$ 中的这个单项式族，
以及两个子域 $k_\alpha(K^p)$ 和 $k_0(K^p)$；它们都属于 $K$。
\end{lemma}

\begin{theorem}[21.8.3]
\label{0.21.8.3-fr}
设 $K$ 是一个特征为 $p>0$ 的域，$k_0$ 是 $K$ 的一个子域，
$(k_\alpha)_{\alpha\in I}$ 是 $K$ 的一个滤过递减子域族，并且各成员都包含
$k_0$。下列条件等价：
\begin{enumerate}
  \item[\rm a)] $\bigcap_{\alpha\in I}k_\alpha(K^p)=k_0(K^p)$。
  \item[\rm b)] 对每个扩张 $L$（它是 $K$ 的扩张），如果
  $\Upsilon_{L/K/k_\alpha}$ 是一个有限秩的 $L$-向量空间（特别地，由
  (21.7.2)，当 $L$ 是有限型扩张时就是如此），则存在 $\alpha\in I$，
  使 $k_\alpha$ 是一个 $k_0$-可容域，相对于 $L$ 关于 $K$ 的扩张。
  \item[\rm b$'$)] 对每个扩张 $L=K(x)$（它是 $K$ 的扩张），如果 $x^p\in K$，
  则存在 $\alpha\in I$，使 $k_\alpha$ 是一个 $k_0$-可容域，相对于
  $L$ 关于 $K$ 的扩张。
  \item[\rm c)] 典范映射
\begin{equation}
\tag{21.8.3.1}
\label{0.21.8.3.1-fr}
\Omega^1_{K/k_0}\longrightarrow\varprojlim_{\alpha}\Omega^1_{K/k_\alpha}
\end{equation}
是单射。
\end{enumerate}
\end{theorem}

\begin{proof}
当然，典范映射 (21.8.3.1) 是在同态的投射系
$\Omega^1_{K/k_0}\to\Omega^1_{K/k_\alpha}$ (20.5.3.3) 中取投射极限所得的。
我们将按如下逻辑次序证明定理：$c)\Rightarrow b)\Rightarrow b')
\Rightarrow a)\Rightarrow c)$。

称 $k_\alpha$ 是一个 $k_0$-可容域（相对于 $L$ 关于 $K$ 的扩张），是指
典范同态
\oldpage[0\textsubscript{IV-1}]{173}
$\Upsilon_{L/K/k_0}\to\Upsilon_{L/K/k_\alpha}$ 是单射；设 $N_\alpha$ 是这个
同态的核，则诸 $N_\alpha$ 构成 $\Upsilon_{L/K/k_0}$ 的一个滤过递减
向量子空间系；由于 $\Upsilon_{L/K/k_0}$ 按假设秩有限，诸 $N_\alpha$ 中有一个
为零，与它们的交为零，是等价的。但是，这个交恰是投射极限同态
$\Upsilon_{L/K/k_0}\to\varprojlim_\alpha\Upsilon_{L/K/k_\alpha}$ 的核。
另一方面，有交换图
\[
\xymatrix{
\Upsilon_{L/K/k_0}\ar[r]\ar[d]&\varprojlim\Upsilon_{L/K/k_\alpha}\ar[d]\\
\Omega^1_{K/k_0}\otimes_KL\ar[r]&
\varprojlim(\Omega^1_{K/k_\alpha}\otimes_KL)
}
\]
其中左边的竖箭头按定义为单射。因此，为证明 $c)$ 蕴涵 $b)$，只须证明
典范映射
\begin{equation}
\tag{21.8.3.2}
\label{0.21.8.3.2-fr}
\Omega^1_{K/k_0}\otimes_KV\longrightarrow
\varprojlim_{\alpha}(\Omega^1_{K/k_\alpha}\otimes_KV)
\end{equation}
对每个向量空间 $V$（它在 $K$ 上）都是单射（特别地，对 $V=L$ 如此）。当
$V=K^n$ 时，这是显然的，因为等式 $W\otimes_KV=W^n$ 对任意向量空间
$W$（它在 $K$ 上）都成立，而乘积与投射极限可交换。另一方面，对 (21.8.3.2)
左端的任意元素 $z$，都存在一个有限秩子空间 $V'$，它包含于 $V$，使
$z\in\Omega^1_{K/k_0}\otimes_KV'$（将后者典范地视为
$\Omega^1_{K/k_0}\otimes_KV$ 的子空间）；如果 $z\ne0$，那么它在
$\varprojlim_\alpha(\Omega^1_{K/k_\alpha}\otimes_KV')$ 中的像不为零；由于
函子 $\varprojlim$ 左正合，$z$ 在 (21.8.3.2) 右端中的像也
$\ne0$，这就证明了 $c)$ 蕴涵 $b)$。

$b)$ 蕴涵 $b')$ 是显然的；来证明 $b')$ 蕴涵 $a)$。沿用 $b')$ 的记号，
可以假设 $L\ne K$。由 (21.4.7) 可知
$\operatorname{rg}\Upsilon_{L/K}\leqslant1$。如果
$\Upsilon_{L/K/k_\alpha}=0$，则由 (21.6.1.1) 无须再证。否则，
$\Upsilon_{L/K/k_\alpha}$ 典范地等同于 $\Upsilon_{L/K}$；若令
$a=x^p$，则 $\Upsilon_{L/K}$ 有一个仅由元素 $d_K(a)\otimes1$ 构成的基
(21.4.7)；所以 $\Upsilon_{L/K/k_\alpha}$ 有一个仅由元素
$d_{K/k_\alpha}(a)\otimes1$ 构成的基，而说一个子域 $k_\alpha$ 是
$k_0$-可容的（相对于 $L$ 关于 $K$ 的扩张），就是要求
$d_{K/k_\alpha}(a)\ne0$，或者等价地，由 (21.4.5)，要求
$a\notin k_\alpha(K^p)$。但是，对任意 $a\notin k_0(K^p)$，都有
$d_{K/k_0}(a)\ne0$ 以及 $K(a^{1/p})\ne K$；因而可以应用 $b')$，它证明
存在一个 $\alpha$，使 $a\notin k_\alpha(K^p)$；换言之，$b')$ 蕴涵 $a)$。

还须证明 $a)$ 蕴涵 $c)$。设 $(x_\mu)_{\mu\in M}$ 是一个 $p$-基底；这里
它是 $K$ 在 $k_0$ 上的基底。于是，诸 $d_{K/k_0}(x_\mu)$ 构成 $K$-向量空间
$\Omega^1_{K/k_0}$ 的一个基 (21.4.5)，而条件 $c)$ 是指：对每个有限子集
$J$（它属于指标集 $M$），都存在一个 $\alpha\in I$，使诸
$d_{K/k_\alpha}(x_\mu)$（其中 $\mu\in J$）在 $\Omega^1_{K/k_\alpha}$ 中
线性无关；但是，由 (21.4.5)，这也意味着诸 $x_\mu$（其中 $\mu\in J$）
构成一个 $p$-自由族（相对于 $k_\alpha$）；具有这一性质的 $\alpha$ 的存在性
由 (21.8.2) 和假设 $a)$ 得出。
\end{proof}
\oldpage[0\textsubscript{IV-1}]{174}

\begin{corollary}[21.8.4]
\label{0.21.8.4-fr}
设 $K$ 是一个特征为 $p>0$ 的域，$k_0$ 是 $K$ 的一个子域，
$(k_\alpha)_{\alpha\in I}$ 是 $K$ 的一个滤过递减子域族，各成员都包含
$k_0$，并且
\[
  \bigcap_{\alpha\in I}k_\alpha(K^p)=k_0(K^p).
\]
设 $L$ 是 $K$ 的一个扩张，并且 $\Upsilon_{L/K/k_0}$ 是一个有限秩的
$L$-向量空间；于是，对每个域 $K'$，若 $K\subset K'\subset L$，
都存在 $\alpha\in I$，使 $k_\alpha$ 是一个 $k_0$-可容域，相对于 $L$
关于 $K'$ 的扩张。
\end{corollary}

\begin{proof}
只须应用 (21.8.3) 和 (21.6.6)。
\end{proof}

\begin{corollary}[21.8.5]
\label{0.21.8.5-fr}
设 $K$ 是一个特征为 $p>0$ 的域，$k_0$ 是 $K$ 的一个子域，
$(k_\alpha)_{\alpha\in I}$ 是 $K$ 的一个滤过递减子域族，各成员都包含
$k_0$，并且
\[
  \bigcap_{\alpha\in I}k_\alpha(K^p)=k_0(K^p).
\]
于是，对每个扩张 $L$（它是 $K$ 的扩张），如果 $\Upsilon_{L/K/k_0}$ 是一个有限秩的
$L$-向量空间，就有
\[
  \bigcap_{\alpha\in I}k_\alpha(L^p)=k_0(L^p).
\]
\end{corollary}

\begin{proof}
事实上，假设 $M$ 是 $L$ 的一个扩张，并且 $\Upsilon_{M/L/k_0}$ 是一个有限秩的
$M$-向量空间。正合列 (21.6.1.1)
\[
  0\longrightarrow \Upsilon_{L/K/k_0}\otimes_L M
  \longrightarrow \Upsilon_{M/K/k_0}
  \longrightarrow \Upsilon_{M/L/k_0}
\]
以及假设表明，$\Upsilon_{M/K/k_0}$ 也是一个有限秩的 $M$-向量空间。因此，
存在一个指标 $\alpha$，使 $k_\alpha$ 是一个 $k_0$-可容域，相对于 $M$
关于 $K$ 的扩张 (21.8.3)，因而也相对于 $M$ 关于 $L$ 的扩张 (21.6.6)；
这对每个扩张 $M$（它是 $L$ 的扩张）都成立，只要 $\Upsilon_{M/L/k_0}$ 秩有限，故该推论
由 (21.8.3) 中 $a)$ 与 $b)$ 的等价性得出。
\end{proof}

\begin{corollary}[21.8.6]
\label{0.21.8.6-fr}
设 $K$ 是一个域，$p$ 是它的特征指数，$k_0$ 是 $K$ 的一个子域。如果 $L$
是 $K$ 的一个扩张，并且 $\Upsilon_{L/K/k_0}$ 是一个有限秩的 $L$-向量空间，
则存在一个子域 $k$，它属于 $K$，包含 $k_0(K^p)$，满足 $[K:k]$ 有限，
并且是一个 $k_0$-可容域，相对于 $L$ 关于 $K$ 的扩张。
\end{corollary}

\begin{proof}
事实上，由 (21.8.3)，只须构造一个滤过递减子域族
$(k_\alpha)_{\alpha\in I}$，其成员是 $K$ 的子域，并且包含 $K^p$ 和 $k_0$，满足
$[K:k_\alpha]<+\infty$，并且
$\bigcap_\alpha k_\alpha=k_0(K^p)$。为此，取一个 $p$-基底
$(x_\lambda)_{\lambda\in J}$；这是 $K$ 在 $k_0$ 上的基底。对每个有限子集
$H$（它属于 $J$），考虑子域 $k_H$，它属于 $K$，并由 $k_0(K^p)$ 与诸
$x_\lambda$（其指标满足 $\lambda\in J-H$）生成；由此定义可知，
$(x_\lambda)_{\lambda\in H}$ 是一个 $p$-基底；这是 $K$ 在 $k_H$ 上的基底，
因而立刻得出诸 $k_H$ 满足所需条件。
\end{proof}

\begin{remarks}[21.8.7]
\label{0.21.8.7-fr}
\begin{enumerate}
  \item[(i)] 已经看到 (21.7.2)，如果 $L$ 是 $K$ 的一个\emph{有限型}扩张，
  那么 $\Upsilon_{L/K/k_0}$ 对每个子域 $k_0$（它属于 $K$）都秩有限。如果 $L$
  是 $K$ 的一个\emph{可分}扩张，情形也一样，因为由 (20.6.19) 有
  $\Upsilon_{L/K/k_0}=0$。最后，如果 $L$ 是一个可分扩张 $L_0$ 的
  \emph{有限型扩张}，而后者又在 $K$ 上，则与 (21.8.5) 中相同的论证表明，
  $\Upsilon_{L/K/k_0}$ 仍秩有限（事实上，它同构于
  $\Upsilon_{L/L_0/k_0}$ 的一个子空间）。

  \oldpage[0\textsubscript{IV-1}]{175}
  \item[(ii)] 在 (21.8.5) 的陈述中，从而也在 (21.8.3, $b)$ 的陈述中，
  不能省略 $\Upsilon_{L/K/k_0}$ 在 $L$ 上秩有限这一假设。取 $k_0$ 为素域
  $\mathbf F_p$，取 $K$ 为满足 $[K:K^p]$ 可数无穷的一个域（例如无穷多个
  不定元上的有理函数域 $\mathbf F_p(X_1,\ldots,X_n,\ldots)$）；(21.8.6)
  的构造方法立即表明，存在一个严格递减无穷子域序列 $(K_n)$，其成员属于 $K$，
  使 $K_0=K$ 且 $\bigcap_nK_n=K^p$。我们将构造有限扩张 $M_n$ 的一个递增序列；
  它们都是 $K$ 的扩张。若 $M$ 是诸 $M_n$ 的并，则扩张 $L=M^{1/p}$（它是
  $K$ 的扩张）将使 (21.8.5)
  失效。为此，取一个元素 $x$，它属于 $K-K^p$；令 $M_0=K^p$，并令
  $M_n=K^p(a_1x,a_2x,\ldots,a_nx)$（$n\geqslant1$）；诸 $a_n$ 递归构造，
  使得 $a_n\in K_n$ 以及 $a_{n+1}\notin M_n(x)$ 对每个 $n$ 都成立：这是可能的，
  因为 $M_n(x)$ 在 $K^p$ 上次数有限，而 $K_n$ 并非如此。于是立即得到
  $x\notin M=L^p$；但是，由 $x=(a_nx)a_n^{-1}$，有
  $x\in K_n(M)=K_n(L^p)$ 对每个 $n$ 都成立。
\end{enumerate}
\end{remarks}

\begin{proposition}[21.8.8]
\label{0.21.8.8-fr}
设 $k$ 是一个特征为 $p>0$ 的域，并设
$A=k[[T_1,\ldots,T_r]]$ 是一个形式幂级数环，其中不定元数为 $r$，底域为
$k$；$K=k((T_1,\ldots,T_r))$ 是它的分式域。于是存在一个滤过递减
Noether 子环族 $(A_\alpha)$，其成员属于 $A$，使得 $A$ 是一个有限型自由
$A_\alpha$-模，对每个 $\alpha$ 都如此；并且，如果 $K_\alpha$ 是 $A_\alpha$ 的分式域，则有
$\bigcap_\alpha K_\alpha=K^p$。
\end{proposition}

\begin{proof}
可以写成 $K^p=k^p((T_1^p,\ldots,T_r^p))$；在 (21.8.6) 的证明中已经看到，
存在一个递减子域族 $(k_\alpha)$，其成员属于 $k$，使得
$[k:k_\alpha]$ 对每个 $\alpha$ 都有限，并且 $\bigcap_\alpha k_\alpha=k^p$；显然，如果令
$A_\alpha=k_\alpha[[T_1^p,\ldots,T_r^p]]$，则 $A$ 是一个有限型自由
$A_\alpha$-模；因此，只须证明关系式
\begin{equation}
\tag{21.8.8.1}
\label{0.21.8.8.1-fr}
  \bigcap_\alpha k_\alpha((T_1^p,\ldots,T_r^p))
  =k^p((T_1^p,\ldots,T_r^p)).
\end{equation}
由于 $\bigcap_\alpha k_\alpha=k^p$，这将由以下两个引理得出。
\end{proof}

\begin{lemma}[21.8.8.2]
\label{0.21.8.8.2-fr}
如果 $k'$ 是域 $k$ 的一个扩张，则有
\[
  k'[[T_1,\ldots,T_r]]\cap k((T_1,\ldots,T_r))=k[[T_1,\ldots,T_r]].
\]
\end{lemma}

\begin{proof}
事实上，令 $C=k[[T_1,\ldots,T_r]]$、$D=k'[[T_1,\ldots,T_r]]$；由于
$k((T_1,\ldots,T_r))$ 是 $C$ 的分式域，只须证明 $D$ 是一个\emph{忠实平坦}
$C$-模（Bourbaki，\emph{Alg. comm.}，第 I 章，\S~3，
n\textsuperscript{o}~5，命题 10）。但是，$C$ 和 $D$ 都是 Noether 局部环；
如果 $\mathfrak m$ 是 $C$ 的极大理想，则
$D/\mathfrak m^jD=(C/\mathfrak m^j)\otimes_k k'$，所以
$D/\mathfrak m^jD$ 是一个平坦的 $(C/\mathfrak m^j)$-模；只须应用
($\mathbf 0_{\mathrm{III}}$, 10.2.1) 和
($\mathbf 0_{\mathrm I}$, 6.6.2)。
\end{proof}

\begin{lemma}[21.8.8.3]
\label{0.21.8.8.3-fr}
设 $k$ 是一个域，$(k_\alpha)$ 是 $k$ 的一个滤过递减子域族，并令
$k_0=\bigcap_\alpha k_\alpha$。假设存在一个幂 $q$，它是 $k$ 的特征指数的幂，使
$k^q\subset k_0$。于是有
\begin{equation}
\tag{21.8.8.4}
\label{0.21.8.8.4-fr}
  \bigcap_\alpha k_\alpha((T_1,\ldots,T_r))=k_0((T_1,\ldots,T_r)).
\end{equation}
\end{lemma}

\begin{proof}
\oldpage[0\textsubscript{IV-1}]{176}
只须证明 (21.8.8.4) 左端的一个元素 $f\ne0$ 属于其右端。取一个指标
$\gamma$，并取 $g\in k_\gamma[[T_1,\ldots,T_r]]$，使
$gf\in k_\gamma[[T_1,\ldots,T_r]]$；把 $g$ 换成 $g^q$ 后，可以假设
$g\in k_0[[T_1,\ldots,T_r]]$。于是，对每个 $\alpha\geqslant\gamma$，有
\[
  gf\in k_\gamma[[T_1,\ldots,T_r]]\cap k_\alpha((T_1,\ldots,T_r))
  =k_\alpha[[T_1,\ldots,T_r]]
\]
这是引理 (21.8.8.2) 的结果。但是显然，诸环
$k_\alpha[[T_1,\ldots,T_r]]$ 的交正是 $k_0[[T_1,\ldots,T_r]]$，因而确有
$f\in k_0((T_1,\ldots,T_r))$。
\end{proof}

\subsection{形式幂级数环的完备微分模}
\label{subsection:0.21.9-fr}

在本款中，不再假设域必具有特征 $>0$。

\begin{lemma}[21.9.1]
\label{0.21.9.1-fr}
设 $k$ 是一个域，$A$ 是一个完备 Noether 局部环且为一个 $k$-代数，
$K$ 是 $A$ 的剩余域。如果 $K$ 是 $k$ 的有限型扩张，则 $A$-模
$\widehat\Omega^1_{A/k}$ 是有限型的。
\end{lemma}

\begin{proof}
由 (20.7.15)，只须证明 $\Omega^1_{K/k}$ 是有限秩的 $K$-向量空间。
而依假设，$K$ 是一个有限型 $k$-代数 $B$ 的分式域；由 (20.4.7)，
$\Omega^1_{B/k}$ 是有限型 $B$-模，故结论由 (20.5.9) 得出。
\end{proof}

\begin{proposition}[21.9.2]
\label{0.21.9.2-fr}
设 $k_0$ 是一个域，$A$ 是一个完备 Noether 局部环且为一个形式光滑的
$k_0$-代数，$\mathfrak m$ 是它的极大理想，$K=A/\mathfrak m$ 是它的剩余域
（因而作为环，$A$ 同构于一个形式幂级数环 $K[[T_1,\ldots,T_r]]$
（19.6.5））。如果 $K$ 是 $k_0$ 的有限型扩张，则
$\widehat\Omega^1_{A/k_0}$ 是一个自由 $A$-模，其秩等于
$\dim(A)+\deg.\operatorname{tr}_{k_0}K$。此外，对每个子域 $k$（属于
$k_0$），只要 $\Omega^1_{k_0/k}$ 的秩有限，
$\widehat\Omega^1_{A/k}$ 就是秩等于
\[
  \dim(A)+\deg.\operatorname{tr}_{k_0}K
  +\operatorname{rg}_{k_0}(\Omega^1_{k_0/k}).
\]
的自由 $A$-模。
\end{proposition}

\begin{proof}
显然，如果 $P$ 是 $k_0$ 的素子域，则 $K$ 是一个形式光滑的 $P$-代数
（19.6.1）。另一方面，注意到 $\Upsilon^K_{A/k_0/P}=0$：事实上，这等价于
同态
\[
  u_{A/k_0/P}\otimes1_K:
  \Omega^1_{k_0/P}\otimes_{k_0}K
  \longrightarrow \Omega^1_{A/P}\otimes_A K
\]
为单射（20.6.14.5），其中 $u$ 表示结构同态 $k_0\to A$。但依假设，
$A$ 是形式光滑的 $k_0$-代数，故由 (20.7.2)，$u_{A/k_0/P}$ 形式左可逆，
所以由 (19.1.7)，$u_{A/k_0/P}\otimes1_K$ 左可逆，因而更是单射；这就证明了
上述断言。在正合序列 (20.6.22.1) 中，把 $\Lambda,A,B,C$ 分别换成
$P,k_0,A,K$，得到正合序列
\[
  0\longrightarrow \Upsilon^K_{A/k_0}
  \longrightarrow \mathfrak m/\mathfrak m^2
  \longrightarrow \Omega^1_{A/k_0}\otimes_A K
  \longrightarrow \Omega^1_{K/k_0}\longrightarrow0,
\]
从而
\[
\begin{split}
  \operatorname{rg}_K(\Omega^1_{A/k_0}\otimes_AK)
  &={\operatorname{rg}}_K(\mathfrak m/\mathfrak m^2)
    +\bigl(\operatorname{rg}_K(\Omega^1_{K/k_0})
    -\operatorname{rg}_K(\Upsilon_{K/k_0})\bigr)\\
  &={\operatorname{rg}}_K(\mathfrak m/\mathfrak m^2)
    +\deg.\operatorname{tr}_{k_0}K
\end{split}
\]
其中用了 (21.7.1)。另一方面，$\widehat\Omega^1_{A/k_0}$ 是形式投射
$A$-模（20.4.9）；由于其拓扑是 $\mathfrak m$-预进拓扑（20.4.5），
$\widehat\Omega^1_{A/k_0}/\mathfrak m^{j+1}\widehat\Omega^1_{A/k_0}$ 是投射
$(A/\mathfrak m^{j+1})$-模（19.2.4），对每个 $j$ 都如此，因而是自由模
（$\mathbf 0_{\mathrm{III}}$, 10.1.2），且其秩为

\oldpage[0\textsubscript{IV-1}]{177}
$n=\operatorname{rg}_K(\Omega^1_{A/k_0}\otimes_AK)$。由 (21.9.1)，
$\widehat\Omega^1_{A/k_0}$ 是有限型 $A$-模；此外，这个 $A$-模由
（$\mathbf 0_{\mathrm{III}}$, 10.2.1）是平坦的，故由
（$\mathbf 0_{\mathrm{III}}$, 10.1.3）它是秩为 $n$ 的自由模，第一项断言
由此得证。

依假设，$\Omega^1_{k_0/k}$ 在 $k_0$ 上秩有限，所以完备张量积
$\Omega^1_{k_0/k}\widehat\otimes_{k_0}A$ 等同于
$\Omega^1_{k_0/k}\otimes_{k_0}A$；由于 $A$ 是形式光滑的 $k_0$-代数，
由 (20.7.18)，同态
\[
  \widehat u_{A/k_0/k}:
  \Omega^1_{k_0/k}\otimes_{k_0}A
  \longrightarrow\widehat\Omega^1_{A/k}
\]
有一个连续 $A$-同态作为左逆，特别地，这蕴含其像在
$\widehat\Omega^1_{A/k}$ 中是闭的；因此，由 (20.7.17) 及上述结论，正合序列
\[
  0\longrightarrow \Omega^1_{k_0/k}\otimes_{k_0}A
  \xrightarrow{\widehat u_{A/k_0/k}}\widehat\Omega^1_{A/k}
  \longrightarrow\widehat\Omega^1_{A/k_0}\longrightarrow0
\]
是分裂的；这就完成了命题的证明。
\end{proof}

\begin{corollary}[21.9.3]
\label{0.21.9.3-fr}
设 $k$ 是一个域，$A$ 是形式幂级数环 $k[[T_1,\ldots,T_r]]$，并赋予其通常的
$k$-代数结构。于是 $\widehat\Omega^1_{A/k}$ 是一个自由 $A$-模，其秩为
$r=\dim(A)$。
\end{corollary}

\begin{proof}
只须注意 $A$ 是形式光滑的 $k$-代数（19.3.4），并在 (21.9.2) 中取
$k_0=k=K$。
\end{proof}

\begin{lemma}[21.9.4]
\label{0.21.9.4-fr}
设 $k$ 是一个特征为 $p>0$ 的域，$A$ 是一个 $k$-代数且为一个完备
Noether 局部环，$B$ 是一个 $k$-子代数，包含在 $A$ 中；它同构于拓扑
$k$-代数 $k[[T_1,\ldots,T_r]]$，并且 $A$ 是有限型 $B$-代数。
如果 $B_1$ 是子代数 $k[[T_1^p,\ldots,T_r^p]]$（属于 $B$），则
$\widehat\Omega^1_{A/k}$ 典范地等同于 $\Omega^1_{A/B_1}$。
\end{lemma}

\begin{proof}
$B_1$ 到一个拓扑 $A$-模 $P$ 的任何连续导子，只要它是一个连续
$k$-导子的限制，而该导子从 $A$ 到 $P$，就必为\emph{零}：因为它在多项式子环
$k[T_1^p,\ldots,T_r^p]$ 上为零，而该子环在 $B_1$ 中稠密。因此，正合序列
(20.3.7.1) 表明典范同态
\[
  \operatorname{D\acute{e}r.\,cont}_{B_1}(A,P)
  \longrightarrow \operatorname{D\acute{e}r.\,cont}_k(A,P)
\]
是双射；考虑到 $\Omega^1_{A/B_1}$ 是有限型 $A$-模（20.4.7），而且由于
$A$ 是完备 Noether 局部环，它是分离且完备的，结论由 (20.7.14.4) 得出。
\end{proof}

\begin{proposition}[21.9.5]
\label{0.21.9.5-fr}
设 $A$ 是一个完备 Noether 局部整环，$A_0$ 是 $A$ 的一个子环，它同构于
形式幂级数环 $k[[T_1,\ldots,T_r]]$，其中 $k$ 是一个域；$A$ 是有限
$A_0$-代数，而分式域 $E$（属于 $A$）是分式域 $L_0$（属于 $A_0$）的
可分扩张。
于是有
\begin{equation}
\tag{21.9.5.1}
\label{0.21.9.5.1-fr}
  \operatorname{rg}_A\widehat\Omega^1_{A/k}
  =\operatorname{rg}_E(\widehat\Omega^1_{A/k}\otimes_AE)=\dim(A).
\end{equation}
\end{proposition}

\begin{proof}
注意，如果 $\mathfrak m_0$ 是 $A_0$ 的极大理想，则 $A$ 的拓扑就是
$\mathfrak m_0$-进拓扑；这是因为 $A$ 是有限 $A_0$-代数（Bourbaki，
\emph{Alg. comm.}，第 IV 章，\S~2，n\textsuperscript{o}~5，命题 9 的
推论 3），并且它在 $A_0$ 上诱导 $\mathfrak m_0$-进拓扑（Bourbaki，
\emph{Alg. comm.}，第 III 章，\S~3，n\textsuperscript{o}~4，定理 3）。
由 (21.9.1)，$\widehat\Omega^1_{A/k}$ 是有限型 $A$-模；并且，

\oldpage[0\textsubscript{IV-1}]{178}
由 (21.9.3)，$\widehat\Omega^1_{A_0/k}$ 是一个自由 $A_0$-模，其秩为
$r=\dim(A_0)=\dim(A)$（16.1.5）；所以张量积
$\widehat\Omega^1_{A_0/k}\otimes_{A_0}A$ 是一个自由 $A$-模，其秩为 $r$，且与
它自身的分离完备化相同；最后，$\widehat\Omega^1_{A/A_0}$ 是有限型
$A$-模（20.4.7），故与它自身的分离完备化相同
（$\mathbf 0_{\mathrm I}$, 7.3.6）。在这些条件下，在同态序列
\begin{equation}
\tag{21.9.5.2}
\label{0.21.9.5.2-fr}
  \widehat\Omega^1_{A_0/k}\otimes_{A_0}A
  \xrightarrow{u}\widehat\Omega^1_{A/k}
  \xrightarrow{v}\widehat\Omega^1_{A/A_0}\longrightarrow0
\end{equation}
中，已知 $v$ 是满射，并且 $\operatorname{Im}(u)$ 在
$\operatorname{Ker}(v)$ 中稠密（20.7.17）；但是
每个 $A_0$-子模，只要包含于 $\widehat\Omega^1_{A/k}$ 中，对
$\mathfrak m_0$-进拓扑都是闭的（$\mathbf 0_{\mathrm I}$, 7.3.5），所以
序列 (21.9.5.2) 正合。
另一方面，注意到 $A$ 在 $A_0$ 上整，$A_0$ 整闭，并且 $E$ 在 $L_0$ 上
可分，故 $\widehat\Omega^1_{A/A_0}$ 是一个挠 $A$-模（20.4.13, (iv)）；
将正合序列 (21.9.5.2) 与 $E$ 作张量积，得到正合序列
\[
  \widehat\Omega^1_{A_0/k}\otimes_{A_0}E
  \longrightarrow \widehat\Omega^1_{A/k}\otimes_AE\longrightarrow0
\]
因而 $\operatorname{rg}_E(\widehat\Omega^1_{A/k}\otimes_AE)\leqslant r$。
所以还须证明 $\widehat\Omega^1_{A/k}\otimes_AE$ 中存在 $r$ 个在 $E$ 上
线性无关的元素。令 $D_i=\partial/\partial T_i$ 为 $A_0$ 到自身的典范导子
$(1\leqslant i\leqslant r)$；它们唯一地延拓为导子（仍记作 $D_i$），
从 $E$ 到自身；这是因为 $E$ 是 $L_0$ 的有限可分扩张。限制到 $A$ 后，
这些导子给出 $A$ 到 $E$ 的导子，并且显然，这些导子的值都落在同一个有限型
$A$-子模中，而该子模包含于 $E$；由于它们在 $A_0$ 上连续，而 $A$ 的拓扑是
$\mathfrak m_0$-进拓扑，这样就构造了 $r$ 个连续导子，它们从 $A$ 到一个
拓扑 $A$-模；由于 $D_i(T_j)=\delta_{ij}$，它们显然线性无关。这就完成了命题的
证明。
\end{proof}

下面的命题是 (21.7.3) 的一个“形式”对应物：

\begin{proposition}[21.9.6]
\label{0.21.9.6-fr}
设 $k_0\subset k\subset K_0$ 是三个特征为 $p>0$ 的域，并且
$[K_0:k]<+\infty$；置
\[
  L_0=K_0((T_1,\ldots,T_r)),\qquad
  M_0=K_0((T_1^p,\ldots,T_r^p)),
\]
\[
  L=k((T_1,\ldots,T_r)),\qquad
  M=k((T_1^p,\ldots,T_r^p)),\qquad
  N=k((T_1^{p^2},\ldots,T_r^{p^2})).
\]
设 $E$ 是 $L_0$ 的一个有限型扩张。
\begin{enumerate}
  \item[(i)] 如果 $\Upsilon_{K_0/k_0}$ 的秩有限，则有
  \begin{equation}
  \tag{21.9.6.1}
  \label{0.21.9.6.1-fr}
    \operatorname{rg}_E\Omega^1_{E/M}
    -\operatorname{rg}_{K_0}\Omega^1_{K_0/k}
    =(r+\deg.\operatorname{tr}_{L_0}E)
      +d(E/K_0,k_0)+d(E/M_0,N/k_0)
  \end{equation}
  （记号同 (21.6.1)），特别地，有
  \begin{equation}
  \tag{21.9.6.2}
  \label{0.21.9.6.2-fr}
    \operatorname{rg}_E\Omega^1_{E/M}
    -\operatorname{rg}_{K_0}\Omega^1_{K_0/k}
    \geqslant r+\deg.\operatorname{tr}_{L_0}E+d(E/K_0,k_0).
  \end{equation}
  此外，如果 (21.9.6.2) 两边相等，则把 $k$ 换成一个域 $k'$，其中
  $k_0\subset k'\subset k$，两边仍然相等。

  \item[(ii)] 再假设 $[k_0:k_0^p]<+\infty$。设 $(k_\alpha)$ 是 $K_0$ 的
  一个滤过递减子域族，各子域都包含 $k_0$，且
  $[K_0:k_\alpha]<+\infty$ 对每个 $\alpha$ 都成立，并且
  $\bigcap_\alpha k_\alpha(K_0^p)=k_0(K_0^p)$；于是存在一个 $\alpha$，
  使得在取 $k=k_\alpha$ 时，(21.9.6.2) 两边相等。
\end{enumerate}
\end{proposition}

\begin{proof}
\begin{enumerate}
  \oldpage[0\textsubscript{IV-1}]{179}
  \item[(i)] 已知（21.1.5, (ii)），
  $\operatorname{rg}_E\Omega^1_{E/M}$ 在把 $M$ 换成 $M(L_0^p)$ 时不变；由于
  $L_0^p=K_0^p((T_1^p,\ldots,T_r^p))$ 且
  $[k(K_0^p):k]<+\infty$，有
  \[
    M(L_0^p)=k(K_0^p)((T_1^p,\ldots,T_r^p))~;
  \]
  同样，在 $\Omega^1_{K_0/k}$ 中把 $k$ 换成 $k(K_0^p)$ 时，它也不变，
  所以可以\emph{假设 $K_0^p\subset k$}；以下始终采用这一假设。我们还引入域
  \[
    k'_0=k_0(K_0^p),\qquad
    N'=k'_0((T_1^{p^2},\ldots,T_r^{p^2}))
  \]
  从而得到域的交换图
  \[
    \xymatrix{
      K_0\ar[r] & M_0\ar[r] & L_0\ar[r] & E\\
      k\ar[u]\ar[r] & N\ar[u]\ar[r] & M\ar[u]\ar[r] & L\ar[u]\\
      k_0\ar[u]\ar[r] & k'_0\ar[u]\ar[r] & N'\ar[u] &
    }
  \]

  我们来计算 (21.9.6.2) 左边与右边之差 $\delta$。

  由 (21.7.3)，有
  \[
    \operatorname{rg}_E\Omega^1_{E/M}
    -\operatorname{rg}_{M_0}\Omega^1_{M_0/M}
    =\deg.\operatorname{tr}_{M_0}E+d(E/M_0,M)
  \]
  并且由于 $L_0$ 是 $M_0$ 的有限扩张，有
  $\deg.\operatorname{tr}_{L_0}E=\deg.\operatorname{tr}_{M_0}E$。另一方面，
  任取一个 $p$-基底：它是 $K_0$ 在 $k=k(K_0^p)$ 上的；它也同样是一个
  $p$-基底，后者是 $M_0$ 在 $M$ 上的。所以由 (21.4.5)，有
  \[
    \operatorname{rg}_{M_0}\Omega^1_{M_0/M}
    =\operatorname{rg}_{K_0}\Omega^1_{K_0/k}
  \]
  从而
  \begin{equation}
  \tag{21.9.6.3}
  \label{0.21.9.6.3-fr}
    \delta=d(E/M_0,M)-d(E/K_0,k_0)-r.
  \end{equation}

  现在记下如下经典引理：

  \begin{lemma}[21.9.6.4]
  \label{0.21.9.6.4-fr}
  对任一域 $k$，形式幂级数域 $K=k((T_1,\ldots,T_r))$ 是 $k$ 的可分扩张。
  \end{lemma}

  \begin{proof}
  为求完整，简要回顾这个引理的证明。只须证明（Bourbaki，\emph{Alg.}，
  第 VIII 章，\S~7，n\textsuperscript{o}~3，定理 1 的证明）：对每个有限扩张
  $k'$（在 $k$ 上），$K\otimes_k k'$ 都没有幂零元。但是，如果置
  $A=k[[T_1,\ldots,T_r]]$ 和 $S=A-\{0\}$，则 $K\otimes_k k'$ 等于
  $S^{-1}(A\otimes_k k')$；而 $A\otimes_k k'$ 典范地等同于整环
  $A'=k'[[T_1,\ldots,T_r]]$，且 $A$ 等同于 $A'$ 的一个子环，结论即得。
  \end{proof}

  利用这个引理和（21.6.2, (ii)），有 $d(L_0/K_0,k_0)=0$；因此公式
  (21.6.4.4) 给出
  \[
    d(E/K_0,k_0)=d(E/L_0,k_0).
  \]

  \oldpage[0\textsubscript{IV-1}]{180}
  另一方面，公式 (21.6.4.4) 和 (21.6.4.2) 给出
  \[
    d(E/M_0,M)=d(E/L_0,M)+d(L_0/M_0,M),
  \]
  \[
    d(L_0/M_0,M)=d(L_0/M_0,M/N)+d(L_0/M_0,N),
  \]
  从而
  \[
    \delta=d(E/L_0,M)+d(L_0/M_0,N)-d(E/L_0,k_0)
    +d(L_0/M_0,M/N)-r.
  \]

  我们将证明关系式
  \begin{equation}
  \tag{21.9.6.5}
  \label{0.21.9.6.5-fr}
    d(L_0/M_0,M/N)=r.
  \end{equation}
  为此，注意到依定义有
  \[
    d(L_0/M_0,M/N)
    =\operatorname{rg}_{L_0}\Upsilon_{L_0/M/N}
    -\operatorname{rg}_{M_0}\Upsilon_{M_0/M/N}.
  \]
  考虑到两个正合序列
  \[
  \begin{aligned}
    0&\longrightarrow\Upsilon_{L_0/M/N}
      \longrightarrow\Omega^1_{M/N}\otimes_ML_0
      \longrightarrow\Omega^1_{L_0/N}
      \longrightarrow\Omega^1_{L_0/M}\longrightarrow0,\\
    0&\longrightarrow\Upsilon_{M_0/M/N}
      \longrightarrow\Omega^1_{M/N}\otimes_MM_0
      \longrightarrow\Omega^1_{M_0/N}
      \longrightarrow\Omega^1_{M_0/M}\longrightarrow0,
  \end{aligned}
  \]
  得到
  \[
  \begin{split}
    d(L_0/M_0,M/N)={}&\operatorname{rg}_{L_0}\Omega^1_{L_0/M}
      -\operatorname{rg}_{L_0}\Omega^1_{L_0/N}\\
      &-\operatorname{rg}_{M_0}\Omega^1_{M_0/M}
      +\operatorname{rg}_{M_0}\Omega^1_{M_0/N}.
  \end{split}
  \]
  由于 $N(L_0^p)\supset M$，有
  $\Omega^1_{L_0/N}=\Omega^1_{L_0/M}$（21.1.5, (ii)）。另一方面，
  $T_i^p$ $(1\leqslant i\leqslant r)$ 构成一个 $p$-基底，即 $M$ 在 $N$
  上的；而一个 $p$-基底，如果它是 $K_0$ 在 $k$ 上的，也同样是一个
  $p$-基底，即 $M_0$ 在 $M$ 上的。所以，由于 $M_0^p\subset N$，根据
  (21.4.5) 和 (21.1.10)，
  \[
    \operatorname{rg}_{M_0}\Omega^1_{M_0/N}
    =r+\operatorname{rg}_{K_0}\Omega^1_{K_0/k}
    =r+\operatorname{rg}_{M_0}\Omega^1_{M_0/M}
  \]
  这就证明了 (21.9.6.5)。

  再由 (21.6.4.2)，有
  \[
    d(E/L_0,M)=d(E/L_0,k_0)+d(E/L_0,M/k_0)
  \]
  以及
  \[
    d(L_0/M_0,N)=d(L_0/M_0,N/k)+d(L_0/M_0,k).
  \]
  但是，由于 $L_0$ 在 $K_0$ 上可分（21.9.6.4），有
  $d(L_0/K_0,k)=0$，因而也有 $d(L_0/M_0,k)=0$（(21.6.2) 和
  (21.6.4)）。所以得到
  \[
    \delta=d(E/L_0,M/k_0)+d(L_0/M_0,N/k).
  \]
  对此公式右边的第一项再应用两次 (21.6.4.2)，得到
  \[
    d(E/L_0,M/k_0)=d(E/L_0,M/N)+d(E/L_0,N/k)+d(E/L_0,k/k_0).
  \]
  由于 $M\subset N(L_0^p)$，有 $d(E/L_0,M/N)=0$（21.6.2），故
  （由 (21.6.4.4)）
  \[
    \delta=d(E/L_0,k/k_0)+d(E/M_0,N/k).
  \]
  但是，$M_0$ 和 $L_0$ 都是 $K_0$ 的可分扩张（21.9.6.4），所以
  \[
    d(L_0/K_0,k/k_0)=d(M_0/K_0,k/k_0)=0
  \]
  （21.6.2）；再应用两次 (21.6.4.4)，

  \oldpage[0\textsubscript{IV-1}]{181}
  得到
  \[
    d(E/L_0,k/k_0)=d(E/K_0,k/k_0)=d(E/M_0,k/k_0),
  \]
  最后应用一次 (21.6.4.2)，便有
  \[
    \delta=d(E/M_0,N/k_0),
  \]
  这就证明了 (21.9.6.1)。此外，当把 $k$ 换成一个子扩张 $k'$（属于
  $k_0$）时，$N$ 被换成它的一个子扩张 $N'$，因而
  $d(E/M_0,N'/k_0)\leqslant d(E/M_0,N/k_0)$（21.6.4.2）；这证明了 (i) 的
  最后一项断言。

  \item[(ii)] 对每个 $\alpha$，置
  $N_\alpha=k_\alpha((T_1^{p^2},\ldots,T_r^{p^2}))$。由 (21.9.6.1) 和
  (21.8.3)，考虑到 $K_0^p\subset k_\alpha$ 且 $E$ 是 $M_0$ 的有限型扩张，
  所要证明的是
  \[
    \bigcap_\alpha N_\alpha=k_0(M_0^p).
  \]
  由 (21.8.8.3) 及关系式
  $\bigcap_\alpha k_\alpha=k_0(K_0^p)$（21.8.6），有
  \begin{equation}
  \tag{21.9.6.6}
  \label{0.21.9.6.6-fr}
    \bigcap_\alpha N_\alpha=N'
    =(k_0(K_0^p))((T_1^{p^2},\ldots,T_r^{p^2}))
  \end{equation}
  另一方面，有
  $k_0(M_0^p)=k_0(K_0^p((T_1^{p^2},\ldots,T_r^{p^2})))$。因此，为完成
  (21.9.6, (ii)) 的证明，只须证明如下引理。

  \begin{lemma}[21.9.6.7]
  \label{0.21.9.6.7-fr}
  设 $k_0$ 是一个特征为 $p>0$ 的域，并且
  $[k_0:k_0^p]<+\infty$；设 $K_0$ 是 $k_0$ 的一个扩张。于是有
  \begin{equation}
  \tag{21.9.6.8}
  \label{0.21.9.6.8-fr}
    (k_0(K_0^p))((T_1,\ldots,T_r))
    =k_0\bigl(K_0^p((T_1,\ldots,T_r))\bigr).
  \end{equation}
  \end{lemma}

  \begin{proof}
  由于 $k_0^p\subset K_0^p$，有
  \[
    k_0^p\bigl(K_0^p((T_1,\ldots,T_r))\bigr)
    =K_0^p((T_1,\ldots,T_r)).
  \]
  如果 $(c_i)_{1\leqslant i\leqslant n}$ 是 $k_0$ 在 $k_0^p$ 上的一个基底，
  则 $(c_i)$ 显然也是 (21.9.6.8) 两边各自在
  $K_0^p((T_1,\ldots,T_r))$ 上的一组生成元。
  \end{proof}
\end{enumerate}
\end{proof}

\begin{remark}[21.9.7]
\label{0.21.9.7-fr}
(21.9.6, (ii)) 的断言\emph{不能推广}到 $[k_0:k_0^p]$ 为无限的情形。
事实上，取 $k_0$ 为有无穷多个不定元的有理函数域
$\mathbf F_p(X_n)_{n\geqslant0}$，于是诸 $X_n$ 构成一个 $p$-基底，即
$k_0$ 在 $\mathbf F_p$ 上的，因而 $[k_0:k_0^p]=+\infty$。在
(21.9.6) 的陈述中，取 $K_0=k_0$（故必有 $k=k_0$）、
$L_0=k_0((T))$，并取 $E$ 为扩张 $L_0(y)$，其中 $y$ 是 $L_0[Y]$ 中
多项式
\[
  Y^p-\sum_{n=0}^{\infty}X_nT^{np}
\]
的一个根。于是 (21.9.6.2) 两边之差 $\delta$ \emph{不为零}。事实上，
$E$ 在 $k_0$ 上\emph{可分}；否则 $y^p$ 将是 $k_0(L_0^p)$ 的元素，而立即可见
这是不可能的（因为有无穷多个代数无关的 $X_n$）。因此有
$d(E/K_0,k_0)=0$，公式 (21.9.6.3) 化为
$\delta=d(E/M_0,M_0)-1$，并且显然
$d(E/M_0,M_0)=\operatorname{rg}_E\Upsilon_{E/M_0}$。一切归结为验证这个秩
等于 $2$。而 $E^p\subset M_0$；一个 $p$-基底——$E$ 在 $L_0$ 上的——
只由元素 $y$ 构成；而 $L_0$ 显然有一个 $p$-基底（在 $M_0$ 上），只由
$T$ 构成；
所以，由于 $E$ 是 $M_0$ 的代数扩张，上述断言由 (21.7.1) 和 (21.4.5)
得出。
\end{remark}

\begin{proposition}[21.9.8]
\label{0.21.9.8-fr}
设 $k_0$ 是一个域（特征任意），$K_0$ 是 $k_0$ 的一个可分扩张，$A$ 是
一个完备 Noether 局部环，它是一个 $k_0$-代数，并且其剩余域

\oldpage[0\textsubscript{IV-1}]{182}
是 $K_0$ 的一个有限扩张。对任一素理想 $\mathfrak p$（属于 $A$），令
$k(\mathfrak p)$ 为
$A$ 在 $\mathfrak p$ 处的剩余域。于是，对任一域 $k$，只要
$k_0\subset k\subset K_0$ 且 $[K_0:k]<+\infty$，就有
\begin{equation}
\tag{21.9.8.1}
\label{0.21.9.8.1-fr}
\begin{split}
  &\operatorname{rg}_{k(\mathfrak p)}
    \bigl(\widehat\Omega^1_{(A/\mathfrak p)/k}
      \otimes_{A/\mathfrak p}k(\mathfrak p)\bigr)
    -\operatorname{rg}_{K_0}\Omega^1_{K_0/k}\\
  &\hspace{35mm}\geqslant
    \dim(A/\mathfrak p)
    +\operatorname{rg}_{k(\mathfrak p)}\Upsilon_{k(\mathfrak p)/k_0}.
\end{split}
\end{equation}
再假设 $[k_0:k_0^p]<+\infty$（其中 $p$ 是 $k_0$ 的特征指数）。于是可以
找到一个域 $k$，满足 $k_0(K_0^p)\subset k\subset K_0$ 和
$[K_0:k]<+\infty$，使 (21.9.8.1) 两边相等。
\end{proposition}

\begin{proof}
由于 $A/\mathfrak p$ 的剩余域 $K$ 也就是 $A$ 的剩余域，并且它是完备的，可以化到
$A$ 为整环且 $\mathfrak p=0$ 的情形；此时令 $E=k(\mathfrak p)$，即 $A$ 的
分式域。由于 $K_0$ 在 $k_0$ 上形式光滑（19.6.1），存在一个
$k_0$-单同态 $K_0\to A$，使 $A$ 成为一个 $K_0$-代数；此外，已知存在
一个 $K_0$-子代数 $A_0$，它包含于 $A$，并且 $K_0$-同构于形式幂级数代数
$K_0[[T_1,\ldots,T_r]]$，并且 $A$ 是\emph{有限} $A_0$-代数（19.8.9）；
这蕴含 $E$ 是分式域 $L_0=K_0((T_1,\ldots,T_r))$（属于 $A_0$）的有限扩张。
于是 $\Upsilon_{K_0/k_0}=0$，因为 $K_0$ 在 $k_0$ 上可分，并且
$\deg.\operatorname{tr}_{L_0}E=0$；此外，因为 $A$ 是有限 $A_0$-代数，
有 $\dim(A)=\dim(A_0)$（16.1.5）。

如果 $p=1$，则由 (21.9.5) 和（17.1.4, (iii)），有
$\operatorname{rg}_E(\widehat\Omega^1_{A/k}\otimes_AE)=r$；又有
$\Omega^1_{K_0/k}=0$（21.7.4），因为 $K_0$ 是 $k$ 的有限可分扩张。
另一方面，$\Upsilon_{E/k_0}=0$，因为 $E$ 在 $k_0$ 上可分；所以在这种情形下，
(21.9.8.1) 两边相等。

现在假设 $p>1$；如果 $B=k[[T_1,\ldots,T_r]]$，则 $A$ 是有限 $B$-代数，
所以置 $B_1=k[[T_1^p,\ldots,T_r^p]]$ 后，
$\widehat\Omega^1_{A/k}$ 等同于 $\Omega^1_{A/B_1}$（21.9.4）；再以 $M$
表示分式域 $k((T_1^p,\ldots,T_r^p))$，由 (20.5.9)，
$\widehat\Omega^1_{A/k}\otimes_AE$ 等同于 $\Omega^1_{E/M}$。
于是，不等式 (21.9.8.1) 正是 (21.9.6.2)，而该推论的最后断言由
（21.9.6, (ii)）得出。
\end{proof}

\section{形式平滑性与正则性的微分判别法}
\label{section:0.22-fr}

本节的主要结果如下：
\begin{enumerate}
  \item[a)] 判别法 (22.2.2) 和 (22.5.8)，它们给出完备 Noether 局部环
  $A$（含有一个域 $k$）在 $k$ 上形式平滑的充分必要条件；判别法
  (22.5.8) 的陈述不涉及微分概念，并补充了 (19.6.6)；陈述 (22.2.2)
  略具技术性，它使我们能够在 $k$ 固定时，对所考虑的局部环 $A$ 给出一种
  \emph{分类}（22.2.5）：这些环由其维数 $n$ 和剩余域 $K$ 确定，唯一的
  条件是
  $n\geqslant\operatorname{rg}_K\Upsilon_{K/k}$。

  \item[b)] 定理 (22.3.3) 以及 Nagata 的 Jacobian 判别法 (22.7.3) 的推论
  (22.7.6)，它们给出一些条件，借此可以断言一个完备 Noether 局部环的
  \emph{局部化}环 $A$（含有一个域 $k$）在 $k$ 上几何正则；它们将在
  第 IV 章第 7 节对局部环的“精细”性质的研究中起重要作用。目前尚无
  不依赖 Nagata 的 Jacobian 判别法的 (22.7.6) 证明。

  \item[c)] Zariski 的 Jacobian 判别法 (22.6.7) 及其变体；它们都是
  (20.7.8) 的容易推论。
\end{enumerate}

\oldpage[0\textsubscript{IV-1}]{183}

\subsection{形式平滑性的提升}
\label{subsection:0.22.1-fr}

\begin{theorem}[22.1.1]
\label{0.22.1.1-fr}
设 $s:\Lambda\to A$、$u:A\to B$ 是两个环同态，$\mathfrak J$ 是 $B$ 的
一个理想，$C$ 是商环 $B/\mathfrak J$。假设 $C$ 对离散拓扑是一个形式平滑的
$\Lambda$-代数。
\begin{enumerate}
  \item[(i)] 假设下列条件成立：
  \begin{enumerate}
    \item[$\alpha$)] $\mathfrak J/\mathfrak J^2$ 是投射 $C$-模。
    \item[$\beta$)] 典范同态
    $\mathbf S_C^*(\mathfrak J/\mathfrak J^2)\to
    \operatorname{gr}_{\mathfrak J}^*(B)$（19.5.2）是双射。
    \item[$\gamma$)] 特征同态
    $\chi:\Upsilon_{C/B/\Lambda}\to\mathfrak J/\mathfrak J^2$（20.6.20）
    是单射。
    \item[$\delta$)] $C$-模
    $\Omega^1_{B/A}\otimes_BC$ 是投射模。
  \end{enumerate}
  于是，赋予 $B$ 以 $\mathfrak J$-预进拓扑后，它是一个形式平滑的
  $\Lambda$-代数。

  \item[(ii)] 反过来，假设 $B$ 是一个形式平滑的 $\Lambda$-代数（对
  $\mathfrak J$-预进拓扑），并且 $A$ 对离散拓扑是一个形式平滑的
  $\Lambda$-代数。
  于是 (i) 的条件 $\alpha$)、$\beta$)、$\gamma$)、$\delta$) 都成立。
\end{enumerate}
\end{theorem}

\begin{proof}
由正合序列 (20.6.22.1)（它可以应用，因为 $B/\mathfrak J^2$ 是一个
$\Lambda$-平凡扩张，即 $C$ 由 $\mathfrak J/\mathfrak J^2$ 构成的，而依假设
$C$ 是形式平滑的 $\Lambda$-代数（19.4.4）），条件 $\gamma$) 等价于
$\Upsilon^C_{B/A/\Lambda}=0$，也等价于下列条件：
\begin{enumerate}
  \item[$\gamma'$)] 同态
  \[
    u_{B/A/\Lambda}\otimes 1_C:
    \Omega^1_{A/\Lambda}\otimes_A C
    \longrightarrow \Omega^1_{B/\Lambda}\otimes_B C
  \]
  是单射。
\end{enumerate}

(i) 条件 $\alpha$) 和 $\beta$) 蕴含 $B$ 是一个形式平滑的
$\Lambda$-代数（对 $\mathfrak J$-预进拓扑）（19.5.4）。因此，断言 $B$（对
$\mathfrak J$-预进拓扑）是形式平滑的 $A$-代数，与断言它是一个
形式平滑的 $A$-代数（\emph{相对于 $\Lambda$}），是一回事。为证明 $B$ 是
形式平滑的 $A$-代数，由 (20.7.2)，只须证明拓扑 $B$-模的连续同态
\[
  u_{B/A/\Lambda}:\Omega^1_{A/\Lambda}\otimes_A B
  \longrightarrow\Omega^1_{B/\Lambda}
\]
\emph{形式左可逆}。由于 $B$ 是形式平滑的 $\Lambda$-代数，拓扑 $B$-模
$\Omega^1_{B/\Lambda}$ 是形式投射的（20.4.9），而其拓扑由 $B$ 的拓扑
导出（20.4.5）。所以由 (19.1.9)，要使 $u_{B/A/\Lambda}$ 形式左可逆，
只须 $u_{B/A/\Lambda}\otimes 1_C$ 左可逆。但由条件 $\gamma'$)，后一映射
是单射，因而有正合序列 (20.6.14.7)
\[
  0\longrightarrow\Omega^1_{A/\Lambda}\otimes_A C
  \longrightarrow\Omega^1_{B/\Lambda}\otimes_B C
  \longrightarrow\Omega^1_{B/A}\otimes_B C\longrightarrow0.
\]
最后，由条件 $\delta$)，$C$-模
$\Omega^1_{B/A}\otimes_BC$ 是投射的，所以前一正合序列分裂；这就完成了
(i) 的证明。

\oldpage[0\textsubscript{IV-1}]{184}

(ii) 如果 $B$ 是一个形式平滑的 $\Lambda$-代数（对 $\mathfrak J$-预进拓扑），
并且 $A$ 对离散拓扑是一个形式平滑的 $\Lambda$-代数，则 $B$ 是一个
形式平滑的 $A$-代数（对 $\mathfrak J$-预进拓扑）（19.3.5, (ii)）；所以条件
$\alpha$) 和 $\beta$) 由 (19.5.6) 得出。此外，(20.7.2) 表明
$u_{B/A/\Lambda}$ 形式左可逆，因而
$u_{B/A/\Lambda}\otimes 1_C$ 左可逆（19.1.7），故\emph{更是}单射。
最后，$\Omega^1_{B/A}$ 是形式投射 $B$-模（20.4.9），所以
$\Omega^1_{B/A}\otimes_BC$ 是投射 $C$-模（19.2.4）。
\end{proof}

\begin{corollary}[22.1.2]
\label{0.22.1.2-fr}
设 $s:\Lambda\to A$、$u:A\to B$ 是两个环同态，$\mathfrak m$ 是 $B$ 的
一个极大理想，$K=B/\mathfrak m$ 是商域。假设 $\Lambda$-代数 $A$ 和 $K$
对离散拓扑都是形式平滑的。于是下列条件等价：
\begin{enumerate}
  \item[a)] $B$ 是一个形式平滑的 $A$-代数（对 $\mathfrak m$-预进拓扑）。
  \item[b)] 典范同态
  \begin{equation}
  \tag{22.1.2.1}
  \label{0.22.1.2.1-fr}
    \mathbf S_K^*(\mathfrak m/\mathfrak m^2)
    \longrightarrow\operatorname{gr}_{\mathfrak m}^*(B)
  \end{equation}
  是双射，并且特征同态
  \begin{equation}
  \tag{22.1.2.2}
  \label{0.22.1.2.2-fr}
    \chi:\Upsilon_{K/B/\Lambda}\longrightarrow\mathfrak m/\mathfrak m^2
  \end{equation}
  是单射。
\end{enumerate}
\end{corollary}

\begin{proof}
这立即由 (22.1.1) 得出，因为在此条件 $\alpha$) 和 $\delta$) 显然成立，
这是由于 $K$ 是一个域。
\end{proof}

\begin{remark}[22.1.3]
\label{0.22.1.3-fr}
假设 (22.1.2) 的假设成立，并再假设 $\mathfrak m/\mathfrak m^2$ 是有限秩
$K$-向量空间。于是，要使 $B$ 成为一个形式平滑的 $A$-代数（对
$\mathfrak m$-预进拓扑），只须下列条件成立：

\begin{env}[22.1.3.1]
\label{0.22.1.3.1-fr}
给定一个多项式代数 $E=K[T_1,\ldots,T_n]$、极大理想 $\mathfrak n$
（属于 $E$），它由诸 $T_i$ $(1\leqslant i\leqslant n)$ 生成，以及一个理想
$\mathfrak b\supset\mathfrak n^k$（属于 $E$），每个 $\Lambda$-同态
$v:B\to E/\mathfrak b$，只要它在取商后给出恒等映射 $K\to K$，就分解为
\[
  B\xrightarrow{w}E/\mathfrak n^k\longrightarrow E/\mathfrak b,
\]
其中 $w$ 是一个 $\Lambda$-同态。
\end{env}

\begin{env}[22.1.3.2]
\label{0.22.1.3.2-fr}
如果 $F$ 是平凡扩张 $D_K(K)$（“$K$ 上的对偶数代数”），
$\rho:A\to F$ 是一个环同态，使复合
$A\xrightarrow{\rho}F\to K$ 为典范同态，那么每个 $\Lambda$-同态
$v:B\to K$，只要它在取商后给出恒等映射 $K\to K$，就分解为
\[
  B\xrightarrow{w}F\longrightarrow K,
\]
其中 $w$ 是一个 $\Lambda$-同态（这里的 $A$-代数结构在 $F$ 上由 $\rho$
定义）。
\end{env}

事实上，已在 (19.5.5) 中看到，条件 (22.1.3.1) 蕴含典范同态
(22.1.2.1) 是双射。由 (22.1.2)，接下来只须证明典范同态
\[
  u_{B/A/\Lambda}\otimes1_K:
  \Omega^1_{A/\Lambda}\otimes_AK
  \longrightarrow\Omega^1_{B/\Lambda}\otimes_BK
\]
是单射；这也等价于：由于这里涉及的是 $K$ 上的向量空间，对每个
$K$-向量空间 $L$，典范同态
$\operatorname{D\acute{e}r}_\Lambda(B,L)\to\operatorname{D\acute{e}r}_\Lambda(A,L)$
是\emph{满射}（由 (20.4.8)）。为此，显然只须且只需典范同态
$\operatorname{D\acute{e}r}_\Lambda(B,K)\to\operatorname{D\acute{e}r}_\Lambda(A,K)$
是满射；由 (20.1.5)，上述断言得证。

注意，条件 (22.1.3.1) 和 (22.1.3.2) 都由下列条件推出：

\begin{env}[22.1.3.3]
\label{0.22.1.3.3-fr}
对每个局部 Artin $\Lambda$-代数 $E$，其剩余域等于 $K$，以及每个幂零理想
$\mathfrak r$（属于 $E$），每个 $\Lambda$-同态 $B\to E/\mathfrak r$，只要
它在取商后给出恒等映射 $K\to K$，就分解为
$B\xrightarrow{u}E\to E/\mathfrak r$，其中 $u$ 是一个
$\Lambda$-同态。
\end{env}
\end{remark}

\oldpage[0\textsubscript{IV-1}]{185}

上述结果可推广如下：

\begin{proposition}[22.1.4]
\label{0.22.1.4-fr}
设 $\varphi:A\to B$ 是 Noether 局部环之间的局部同态，$k$（相应地，$K$）
是 $A$（相应地，$B$）的剩余域。要使 $B$ 成为形式平滑的 $A$-代数
（对于预进拓扑），当且仅当：对于每个局部 Artin 环 $C$，其剩余域等于
$K$；对于每个局部同态 $A\to C$（它使 $C$ 成为拓扑 $A$-代数）；并对于
每个幂零理想 $\mathfrak J$（属于 $C$），凡是局部 $A$-同态
$B\to C/\mathfrak J$，若其取商所得同态 $K\to K$ 是恒等同态，就可分解为
\[
  B\xrightarrow{f}C\longrightarrow C/\mathfrak J,
\]
其中 $f$ 是局部 $A$-同态。
\end{proposition}

\begin{proof}
此条件按定义 (19.3.1) 是必要的；我们来证明它是充分的。局部
$\widehat A$-同态从 $\widehat B$ 到离散局部 $\widehat A$-代数，都由
局部 $A$-同态（从 $B$ 到该代数）延拓而来；故由 (19.3.6)，可只考虑 $A$ 和
$B$ 都\emph{完备}的情形。当 $A$ 是一个\emph{域} $k$ 时，命
$\Lambda$ 为 $k$ 的素子域，本命题由 (22.1.3.3) 和 (19.6.1) 得出。

在一般情形，置 $B_0=B\otimes_Ak=B/\mathfrak mB$（$\mathfrak m$ 是
$A$ 的极大理想）；它是完备的。设 $C_0$ 是一个局部 Artin $k$-代数，其
剩余域等于 $K$，$\mathfrak J_0$ 是 $C_0$ 的幂零理想；每个局部
$k$-同态 $g_0:B_0\to C_0/\mathfrak J_0$，若其取商后在 $K$ 上诱导
恒等同态，则经复合给出
局部 $A$-同态
\[
  g:B\longrightarrow B_0\xrightarrow{g_0}C_0/\mathfrak J_0 .
\]
因此按假设，它分解为
$B\xrightarrow{f}C_0\to C_0/\mathfrak J_0$，其中 $f$ 是局部
$A$-同态；但因 $C_0$ 是 $k$-代数，$f$ 又分解为
$B\to B_0\xrightarrow{f_0}C_0$，其中 $f_0$ 是局部 $k$-同态。由此可见，
把陈述中的 $B_0$ 和 $k$ 分别代替 $B$ 和 $A$ 后，其假设仍然成立；故由
刚才证明的情形，$B_0$ 是形式平滑的 $k$-代数。

应用 (19.7.2) 和 (19.7.1)，可知存在一个完备 Noether 局部环 $B'$、一个
局部同态 $A\to B'$，使 $B'$ 成为\emph{平坦} $A$-模和\emph{形式平滑}
$A$-代数，并存在 $k$-同构
\[
  u'_0:B'_0=B'\otimes_Ak\xrightarrow{\sim}B_0.
\]
令 $v'_0$ 为 $u'_0$ 的逆同构；我们要证明，存在一个局部 $A$-同态
$v:B\to B'$，使 $v'_0$ 由 $v$ 取商而得。为此，以 $\mathfrak n$ 和
$\mathfrak n'$ 分别表示 $B$ 和 $B'$ 的极大理想；对每个 $j$，令
\[
  w_j:B/(\mathfrak n^{j+1}+\mathfrak mB)
  \longrightarrow B'/(\mathfrak n'^{j+1}+\mathfrak mB')
\]
为 $v'_0$ 取商所得的同态。只须对每个 $j$ 构造局部 $A$-同态
$v_j:B/\mathfrak n^{j+1}\to B'/\mathfrak n'^{j+1}$，使下列图表交换：
\[
\xymatrix@C=4em@R=3em{
  B/\mathfrak n^{j+1} \ar[r]^{v_j} \ar[d] &
  B'/\mathfrak n'^{j+1} \ar[d] \\
  B/\mathfrak n^j \ar[r]_{v_{j-1}} & B'/\mathfrak n'^j
}
\qquad
\xymatrix@C=4em@R=3em{
  B/\mathfrak n^{j+1} \ar[r]^{v_j} \ar[d] &
  B'/\mathfrak n'^{j+1} \ar[d] \\
  B/(\mathfrak n^{j+1}+\mathfrak mB) \ar[r]_{w_j} &
  B'/(\mathfrak n'^{j+1}+\mathfrak mB')
}
\]
由于 $B$ 和 $B'$ 完备，同态 $v=\varprojlim v_j$ 就满足要求。

\oldpage[0\textsubscript{IV-1}]{186}

而 $v_j$ 的递归构造，由关于 $B$ 的假设以及 (19.3.10.1) 中的同一论证
得出；这里使用引理 (19.3.10.2) 和等式
\[
  \mathfrak n'^j+(\mathfrak n'^{j+1}+\mathfrak mB')
  =\mathfrak n'^j+\mathfrak mB'.
\]
于是由 (19.7.1.4)，$v$ 是一个 $A$-\emph{同构}：这是因为 $B'$ 是平坦
$A$-模，且 $B$ 对 $\mathfrak m$-预进拓扑完备，而诸理想
$\mathfrak m^iB$ 在 $B$ 的 $\mathfrak n$-进拓扑中是闭的
（$\mathbf 0_{\mathrm I}$，7.3.5）。证毕。
\end{proof}

\subsection{域上的形式平滑局部代数的微分刻画}
\label{subsection:0.22.2-fr}

\begin{env}[22.2.1]
\label{0.22.2.1-fr}
设 $k$ 是一个域，$P$ 是它的素子域，$A$ 是一个同时为\emph{局部环}的
$k$-代数，$\mathfrak m$ 是其极大理想，$K=A/\mathfrak m$ 是其剩余域。
因 $K$ 在 $P$ 上可分，$K$ 对离散拓扑是形式平滑的 $P$-代数
（19.6.1）；所以这个 $P$-扩充 $A/\mathfrak m^2$，即 $K$ 被
$\mathfrak m/\mathfrak m^2$ 扩充所得者，是 $P$-\emph{平凡的}，因而特征同态
\begin{equation}
\tag{22.2.1.1}
\label{0.22.2.1.1-fr}
  \chi_{A/k}:\Upsilon_{K/k}\longrightarrow\mathfrak m/\mathfrak m^2
\end{equation}
得到定义（20.6.23）。
\end{env}

\begin{theorem}[22.2.2]
\label{0.22.2.2-fr}
在 (22.2.1) 的条件下，要使 $A$ 作为 $k$-代数对
$\mathfrak m$-预进拓扑形式平滑，当且仅当下列条件成立：
\begin{enumerate}
  \item[(i)] 典范同态
  $\mathbf S_K^*(\mathfrak m/\mathfrak m^2)\to
  \operatorname{gr}_{\mathfrak m}^*(A)$ 是双射。
  \item[(ii)] 特征同态
  $\chi_{A/k}:\Upsilon_{K/k}\to\mathfrak m/\mathfrak m^2$ 是单射。
\end{enumerate}
\end{theorem}

\begin{proof}
这是 (22.1.2) 的特殊情形，其中把 $\Lambda$ 换成 $P$、把 $A$ 换成 $k$、
把 $B$ 换成 $A$；事实上，$K$ 和 $k$ 在 $P$ 上可分，所以都是形式平滑的
$P$-代数（19.6.1）。
\end{proof}

\begin{remark}[22.2.3]
\label{0.22.2.3-fr}
(22.2.2) 的条件 (ii) 可由下列任一条件代替：
\begin{enumerate}
  \item[a)] 典范同态
  $\Omega^1_k\otimes_kK\to\Omega^1_A\otimes_AK$ 是单射。
  \item[b)] 对每个 $n$，典范同态
  \[
    \Omega^1_k\otimes_k(A/\mathfrak m^{n+1})
    \longrightarrow
    \Omega^1_A\otimes_A(A/\mathfrak m^{n+1})
  \]
  左可逆。
  \item[c)] 典范同态
  $\Omega^1_k\widehat\otimes_kA\to\widehat\Omega^1_A$ 左可逆。
  \item[d)] $A$ 是一个 $k$-代数；它对 $\mathfrak m$-预进拓扑
  \emph{相对于素域形式平滑}。
\end{enumerate}
事实上，在 (22.1.1) 中已经指出，$\chi_{A/k}$ 的单射性等价于 $a)$，
而 $a)$ 与 (22.2.2) 的条件 (i) 合取，等价于 $d)$ 与 (i) 合取。最后，
在 (22.1.1) 中还已看到 $b)$ 与 $d)$ 等价，而 $b)$ 与 $c)$ 的等价性
由 (19.1.8) 得出。
\end{remark}

\begin{env}[22.2.4]
\label{0.22.2.4-fr}
(22.2.2) 的主要应用涉及\emph{完备} Noether 局部 $k$-代数（在这种情形，
$\mathfrak m/\mathfrak m^2$ 是有限秩 $K$-向量空间）。具体而言，设 $k$
是一个域，$K$ 是 $k$ 的一个扩域，$V$ 是有限秩 $K$-向量空间；

\oldpage[0\textsubscript{IV-1}]{187}

我们考虑三元组 $(A,\alpha,\beta)$，其中 $A$ 是一个形式平滑的完备
Noether 局部 $k$-代数，极大理想为 $\mathfrak m$；
$\alpha:A/\mathfrak m\xrightarrow{\sim}K$ 是 $k$-代数同构，
$\beta:\mathfrak m/\mathfrak m^2\xrightarrow{\sim}V$ 是向量空间的
双同构（相对于同构 $\alpha$）。给定两个这样的三元组
$(A,\alpha,\beta)$、$(A',\alpha',\beta')$，从第一个到第二个的一个
\emph{等价}，是一个 $k$-同构 $u:A\xrightarrow{\sim}A'$，使得（以
$\mathfrak m'$ 表示 $A'$ 的极大理想）下列图表交换：
\[
\xymatrix@C=4em@R=3em{
  A/\mathfrak m \ar[r]^{\operatorname{gr}^0(u)} \ar[d]_{\alpha} &
  A'/\mathfrak m' \ar[d]^{\alpha'} \\
  K \ar[r]_{1_K} & K
}
\qquad
\xymatrix@C=4em@R=3em{
  \mathfrak m/\mathfrak m^2 \ar[r]^{\operatorname{gr}^1(u)}
    \ar[d]_{\beta} &
  \mathfrak m'/\mathfrak m'^2 \ar[d]^{\beta'} \\
  V \ar[r]_{1_V} & V
}
\]
显然，三元组 $(A,\alpha,\beta)$ 的诸\emph{等价类}组成一个\emph{集合}，
记作 $\operatorname{FL}(K/k,V)$。现在注意，每个三元组
$(A,\alpha,\beta)$ 都对应于 $K$-向量空间的复合同态
\[
  \Upsilon_{K/k}\xrightarrow{\chi_{A/k}}
  \mathfrak m/\mathfrak m^2\xrightarrow{\beta}V
\]
（这里把 $\chi_{A/k}$ 看成相对于 $\alpha^{-1}$ 的双同态）；上述定义
证明（由 (20.6.22.2)），这个 $K$-同态\emph{只依赖于}
$(A,\alpha,\beta)$ 在 $\operatorname{FL}(K/k,V)$ 中的等价类；换言之，
我们定义了一个典范映射
\begin{equation}
\tag{22.2.4.1}
\label{0.22.2.4.1-fr}
  \operatorname{FL}(K/k,V)\longrightarrow
  \operatorname{Hom}_K(\Upsilon_{K/k},V).
\end{equation}
\end{env}

\begin{proposition}[22.2.5]
\label{0.22.2.5-fr}
设 $k$ 是一个域，$K$ 是 $k$ 的一个扩域，$V$ 是有限秩 $K$-向量空间。
典范映射 (22.2.4.1) 把 $\operatorname{FL}(K/k,V)$ 双射到
所有单射 $K$-同态（从 $\Upsilon_{K/k}$ 到 $V$）所成的集合。
\end{proposition}

\begin{proof}
(i) 为证明 (22.2.4.1) 是单射，需要证明如下断言。设 $A$、$A'$ 是两个
形式平滑的完备 Noether 局部 $k$-代数，$\mathfrak m$、$\mathfrak m'$
分别是其极大理想，
$\alpha:A/\mathfrak m\xrightarrow{\sim}K$、
$\alpha':A'/\mathfrak m'\xrightarrow{\sim}K$ 是两个 $k$-同构；再设给定
两个 $k$-同构
\[
  v^0:A/\mathfrak m\xrightarrow{\sim}A'/\mathfrak m',\qquad
  v^1:\mathfrak m/\mathfrak m^2\xrightarrow{\sim}
      \mathfrak m'/\mathfrak m'^2
\]
并使下列图表
\begin{equation}
\tag{22.2.5.1}
\label{0.22.2.5.1-fr}
\xymatrix@C=4em@R=3em{
  A/\mathfrak m \ar[r]^{v^0} \ar[dr]_{\alpha} &
  A'/\mathfrak m' \ar[dl]^{\alpha'} \\
  & K
}
\qquad
\xymatrix@C=4em@R=3em{
  \mathfrak m/\mathfrak m^2 \ar[r]^{v^1} &
  \mathfrak m'/\mathfrak m'^2 \\
  \Upsilon_{K/k} \ar[u]^{\chi_A} \ar[ur]_{\chi_{A'}} &
}
\end{equation}
交换。所要证明的是，存在一个 $k$-同构
$u:A\xrightarrow{\sim}A'$，使 $\operatorname{gr}^0(u)=v^0$ 且
$\operatorname{gr}^1(u)=v^1$。首先注意，因 $K$ 是一个域，典范同态
(20.6.7.1) 是\emph{双射}（20.6.11），所以已经存在一个 $k$-同构
$v:A/\mathfrak m^2\xrightarrow{\sim}A'/\mathfrak m'^2$，使
$\operatorname{gr}^0(v)=v^0$ 且 $\operatorname{gr}^1(v)=v^1$。现在注意，
因 $A'$ 可度量且完备，而 $A$ 是形式平滑的 $k$-代数，复合同态
$A\to A/\mathfrak m^2\xrightarrow{v}A'/\mathfrak m'^2$ 分解为
$A\xrightarrow{u}A'\to A'/\mathfrak m'^2$，其中 $u$ 是

\oldpage[0\textsubscript{IV-1}]{188}

一个 $k$-同态（19.3.11）；所以已有
$\operatorname{gr}^0(u)=v^0$ 且 $\operatorname{gr}^1(u)=v^1$。另一方面，
有交换图
\[
\xymatrix@C=5em@R=3em{
  \mathbf S_K^*(\mathfrak m/\mathfrak m^2) \ar[r] \ar[d]_{w} &
  \operatorname{gr}_{\mathfrak m}^*(A) \ar[d]^{\operatorname{gr}^*(u)} \\
  \mathbf S_K^*(\mathfrak m'/\mathfrak m'^2) \ar[r] &
  \operatorname{gr}_{\mathfrak m'}^*(A')
}
\]
其中横箭头是典范同态，而 $w$ 由 $v^0$ 和 $v^1$ 得到；由于 $v^0$ 和
$v^1$ 都是双射，$w$ 也如此。又因 $A$ 和 $A'$ 是形式平滑的 $k$-代数，
由 (22.2.2, (i)) 可知 $\operatorname{gr}(u)$ 是双射；但 $A$ 和 $A'$
分离且完备，因此 $u$ 本身是\emph{双射}
（Bourbaki，\emph{Alg. comm.}，第 III 章，§~2，
n\textsuperscript{o}~8，定理 1 的推论 3）。

(ii) 还须证明 (22.2.4.1) 的像正是从 $\Upsilon_{K/k}$ 到 $V$ 的
\emph{单射}同态所成的集合。由 (22.2.2, (ii)) 已知其像包含于这个集合；
所以只须证明下列引理。
\end{proof}

\begin{lemma}[22.2.5.2]
\label{0.22.2.5.2-fr}
对于每个 $K$-同态 $h:\Upsilon_{K/k}\to V$，存在一个 $k$-代数 $A$，
它是正则且完备的 Noether 局部环，其极大理想为 $\mathfrak m$、剩余域为
$K$；并存在一个 $K$-同构
$\beta:\mathfrak m/\mathfrak m^2\xrightarrow{\sim}V$，使得
$h=\beta\circ\chi_{A/k}$。
\end{lemma}

\begin{proof}
事实上，若此引理已证，则 $A$ 的正则性蕴含 (22.2.2) 的条件 (i) 成立
（17.1.1）；若再有 $h$ 是单射，则 (22.2.2) 表明 $A$ 是形式平滑的
$k$-代数，其等价类以 $h$ 为像（反之，若 $h$ 不是单射，则 (22.2.5.2)
所保证存在的 $k$-代数 $A$ 并非形式平滑）。

为证明 (22.2.5.2)，令 $B$ 为 $K$-代数 $\mathbf S_K^*(V)$，
$\mathfrak n$ 为 $B$ 的增广理想，并令 $A=\widehat B$ 为 $K$-代数
$B$ 对 $\mathfrak n$-预进拓扑的完备化（于是 $A$ 同构于形式幂级数代数
$K[[T_1,\ldots,T_n]]$，若 $n=\operatorname{rg}_KV$）。由 (20.6.11)，
在 $A/\mathfrak m^2$（其中 $\mathfrak m=\mathfrak nA$）上存在一个
$k$-扩充结构，是 $K$ 被 $V$ 扩充所得者（它一般\emph{不同于}
$k$-平凡结构，后者由典范嵌入 $k\to K\to A$ 定义），使 $h$ 是这个扩充的
特征同态。
若 $f:k\to A/\mathfrak m^2$ 是结构同态，则 $k$ 在素域上可分这一事实
容许由 (19.6.1) 把 $f$ 分解为
$k\xrightarrow{g}A\to A/\mathfrak m^2$；如此定义的 $k$-代数 $A$
满足要求（20.6.22.2）。
\end{proof}

\begin{theorem}[22.2.6]
\label{0.22.2.6-fr}
设 $k$ 是一个域，$K$ 是 $k$ 的一个扩域。存在一个局部 Noether
$k$-代数 $A$，其极大理想为 $\mathfrak m$，使得
$A/\mathfrak m=K$ 且 $A$ 对其 $\mathfrak m$-预进拓扑形式平滑，
当且仅当 $\Upsilon_{K/k}$ 是有限秩 $K$-向量空间。这时，
$k$-代数 $A$ 由其维数确定（精确到一个 $k$-同构，该同构取商后给出
恒等同态 $K\to K$）；这个维数只须满足
$\geqslant\operatorname{rg}_K\Upsilon_{K/k}$。
\end{theorem}

\begin{proof}
由于 $\mathfrak m/\mathfrak m^2$ 在 $K$ 上秩有限，条件的必要性由
(22.2.2, (ii)) 得出；反过来，(22.2.5) 表明这个条件是充分的，并且可把
$\mathfrak m/\mathfrak m^2$ 取为一个 $K$-向量空间，其秩可任意取为
$\geqslant\operatorname{rg}_K(\Upsilon_{K/k})$。
\end{proof}

\oldpage[0\textsubscript{IV-1}]{189}

特别地，$\Upsilon_{K/k}$ 的恒等同态\emph{典范地}给 $K$ 配上一个形式平滑
的完备 Noether $k$-代数，对这个代数，$\Upsilon_{K/k}$ 同构于
$\mathfrak m/\mathfrak m^2$。

\begin{remarks}[22.2.7]
\label{0.22.2.7-fr}
\begin{enumerate}
  \item[(i)] 设 $\Lambda$ 是剩余域为 $k$ 的 Noether 局部环；由
  (19.7.1) 和 (19.7.2)，确定那些既为形式平滑 $\Lambda$-代数、又是完备
  Noether 局部环的 $A$，\emph{等价于}把 $\Lambda$ 换成 $k$ 后的同一问题。
  因而原则上，(22.2.5) 已把它完全解决：问题归结为所求 $k$-代数的
  \emph{剩余域}作为 $k$ 的扩域所具有的结构。

  \item[(ii)] $\Upsilon_{K/k}$ 在 $K$ 上秩有限这一事实，
  \emph{并不蕴含} $K$ 在 $k$ 上具有“有限纯不可分重数”（参见 (19.6.6)）。
  例如，设 $k_0$ 是特征 $p>0$ 的完美域，$k=k_0(T)$ 是 $k_0$ 上一个
  不定元的有理函数域，而 $K=k^{p^{-\infty}}$ 是 $k$ 的最小完美扩域。
  这时 $\Omega^1_K=\Omega^1_{K/k}=0$（21.4.4），故按定义 (20.6.1.1)，
  $\Upsilon_{K/k}=\Omega^1_k\otimes_kK$；又因 $k^p=k_0(T^p)$，$T$ 是
  一个 $p$-基底，即 $k$ 在 $k^p$ 上的一个基底，所以由 (21.4.5)，
  $\Omega^1_k$ 的秩为 $1$（作为 $k$-模），从而 $\Upsilon_{K/k}$ 的秩为
  $1$（作为 $K$-模）。但是对于每个
  $n$，局部环 $K\otimes_kk^{p^{-n}}$ 的剩余域同构于 $K$，故它在
  $k^{p^{-n}}$ 上不可分。
\end{enumerate}
\end{remarks}

定理 (22.2.2) 可推广如下：

\begin{proposition}[22.2.8]
\label{0.22.2.8-fr}
在 (22.2.1) 的条件下，假设 $p$ 是 $k$ 的特征并满足 $>0$；设
$(k_\alpha)$ 是 $k$ 的子域组成的一个递减有向族，并满足
$\bigcap_\alpha k_\alpha(k^p)=k^p$。则 (22.2.2) 的条件 (ii) 可由
下列任一条件代替：
\begin{enumerate}
  \item[a)] 存在 $\alpha_0$，使对每个
  $\alpha\geqslant\alpha_0$，典范同态
  \[
    \Omega^1_{k/k_\alpha}\otimes_kK
    \longrightarrow\Omega^1_{A/k_\alpha}\otimes_AK
  \]
  都是单射。
  \item[b)] 存在 $\alpha_0$，使 $A$ 是一个 $k$-代数，并且对
  $\mathfrak m$-预进拓扑相对于 $k_\alpha$ 形式平滑；这对每个
  $\alpha\geqslant\alpha_0$ 都成立。
  \item[c)] 对每个 $\alpha$ 以及每个整数 $n\geqslant0$，典范同态
  \[
    \Omega^1_{k/k_\alpha}\otimes_k(A/\mathfrak m^{n+1})
    \longrightarrow
    \Omega^1_{A/k_\alpha}\otimes_A(A/\mathfrak m^{n+1})
  \]
  都左可逆。
\end{enumerate}
\end{proposition}

\begin{proof}
事实上，若 $A$ 是形式平滑的 $k$-代数，则它\emph{更是}相对于 $k$ 的
每个子域形式平滑；由 (20.7.2) 和 (19.1.7)，这蕴含 $c)$。由
(20.7.2)，显然 $c)$ 蕴含 $b)$，而 $b)$ 蕴含 $a)$。最后，考虑
(22.2.3)，只须证明 $a)$ 蕴含典范同态
\[
  u:\Omega^1_k\otimes_kK\longrightarrow\Omega^1_A\otimes_AK
\]
是单射。对每个 $\alpha$，有交换图
\[
\xymatrix@C=5em@R=3em{
  \Omega^1_k\otimes_kK \ar[r]^{u} \ar[d]_{v_\alpha} &
  \Omega^1_A\otimes_AK \ar[d] \\
  \Omega^1_{k/k_\alpha}\otimes_kK \ar[r]_{u_\alpha} &
  \Omega^1_{A/k_\alpha}\otimes_AK
}
\]

\oldpage[0\textsubscript{IV-1}]{190}

并且关于 $u_\alpha$ 对 $\alpha\geqslant\alpha_0$ 为单射的假设，蕴含
$\operatorname{Ker}(u)$ 包含于诸 $\operatorname{Ker}(v_\alpha)$ 的
\emph{交}。但由假设
$\bigcap_\alpha k_\alpha(k^p)=k^p$，已在 (21.8.3) 中看到这个交为零，
从而证明完毕。
\end{proof}

\begin{proposition}[22.2.9]
\label{0.22.2.9-fr}
在 (22.2.1) 的假设下，下列条件等价：
\begin{enumerate}
  \item[a)] $A/\mathfrak m^2$ 是一个 $k$-扩充，而且是 $k$-平凡的；
  它是 $K$ 被 $\mathfrak m/\mathfrak m^2$ 扩充所得者。
  \item[b)] $\chi_{A/k}=0$。
  \item[c)] 典范同态 (20.5.11.2)
  \[
    \delta_{K/A/k}:\mathfrak m/\mathfrak m^2
    \longrightarrow\Omega^1_{A/k}\otimes_AK
  \]
  是单射（换言之，序列
  \begin{equation}
  \tag{22.2.9.1}
  \label{0.22.2.9.1-fr}
    0\longrightarrow\mathfrak m/\mathfrak m^2
    \longrightarrow\Omega^1_{A/k}\otimes_AK
    \longrightarrow\Omega^1_{K/k}\longrightarrow0
  \end{equation}
  正合）。
\end{enumerate}
\end{proposition}

\begin{proof}
$b)$ 与 $c)$ 的等价性由正合列 (20.6.22.1) 得出；$c)$ 与 $a)$ 的等价性
由 (20.5.12) 得出，因为 (22.2.9.1) 若正合，则作为向量空间序列必分裂。

注意，如果 $K$ 在 $k$ 上\emph{可分}，则由 $\chi_{A/k}$ 的定义
(22.2.1) 和 (20.6.3)，(22.2.9) 的这些等价条件成立。
\end{proof}

\begin{corollary}[22.2.10]
\label{0.22.2.10-fr}
在 (22.2.1) 的假设下，考虑一个子域 $k_0$（属于 $k$）。下列条件等价：
\begin{enumerate}
  \item[a)] $A/\mathfrak m^2$ 是一个 $k_0$-扩充，而且是
  $k_0$-平凡的；它是 $K$ 被 $\mathfrak m/\mathfrak m^2$ 扩充所得者。
  \item[b)] 复合同态
  \[
    \Upsilon_{K/k_0}\longrightarrow\Upsilon_{K/k}
    \xrightarrow{\chi_{A/k}}\mathfrak m/\mathfrak m^2
  \]
  为零。
  \item[c)] 特征同态
  $\chi_{A/k}:\Upsilon_{K/k}\to\mathfrak m/\mathfrak m^2$ 分解为
  \[
    \Upsilon_{K/k}\longrightarrow\Upsilon_{K/k/k_0}
    \xrightarrow{\chi'}\mathfrak m/\mathfrak m^2.
  \]
\end{enumerate}
此外，当这些条件成立时，同态 $\chi'$ 被唯一确定，并且等于
$k_0$-平凡 $k_0$-扩充 $A/\mathfrak m^2$（即 $K$ 被
$\mathfrak m/\mathfrak m^2$ 扩充所得者）的特征同态。
\end{corollary}

\begin{proof}
因 $b)$ 中的复合正是 $\chi_{A/k_0}$，$a)$ 与 $b)$ 的等价性由
(22.2.9) 得出；另一方面，由正合列 (20.6.18.1)
\[
  \Upsilon_{K/k_0}\longrightarrow\Upsilon_{K/k}
  \longrightarrow\Upsilon_{K/k/k_0}\longrightarrow0,
\]
$b)$ 与 $c)$ 等价。又因典范同态
$\Upsilon_{K/k}\to\Upsilon_{K/k/k_0}$ 是满射，这还蕴含 $\chi'$ 的
唯一性；而 $\chi'$ 是 $k_0$-扩充 $A/\mathfrak m^2$ 的特征同态这一事实，
由 (20.6.22.2) 得出。
\end{proof}

\begin{corollary}[22.2.11]
\label{0.22.2.11-fr}
在 (22.2.1) 的假设下，设 $p$ 是 $k$ 的特征指数，设 $(k_\alpha)$ 是
$k$ 的子域组成的一个递减有向族，并满足
$\bigcap_\alpha k_\alpha(k^p)=k^p$。如果 $\Upsilon_{K/k}$ 在 $K$ 上秩
有限，则存在一个指标 $\alpha$，使 $A/\mathfrak m^2$ 是一个
$k_\alpha$-扩充，而且是 $k_\alpha$-平凡的；它是 $K$ 被
$\mathfrak m/\mathfrak m^2$ 扩充所得者。
\end{corollary}

\begin{proof}
事实上，由 (21.8.3) 已知，此时存在一个指标 $\alpha$，使典范同态
$\Upsilon_{K/k}\to\Upsilon_{K/k/k_\alpha}$ 是双射；所以 (22.2.10) 的
条件 $b)$（其中以 $k_\alpha$ 代替 $k_0$）成立。
\end{proof}

\oldpage[0\textsubscript{IV-1}]{191}

\subsection{应用于某些局部环与其完备化之间的关系}
\label{subsection:0.22.3-fr}

\begin{lemma}[22.3.1]
\label{0.22.3.1-fr}
设 $A_0$ 是 Noether 半局部环，$A$ 是有限 $A_0$-代数（因而也是 Noether
半局部环；参见 Bourbaki，\emph{Alg. comm.}，第 IV 章，§~2，
n\textsuperscript{o}~5，命题 9 的推论 3）。如果 $\widehat A_0$ 和
$\widehat A$ 是 $A_0$ 和 $A$ 对各自预进拓扑的完备化，那么，当赋予
$A_0$、$A$ 和 $\widehat A$ 离散拓扑时，$\widehat A$ 是一个
$A$-代数，并且相对于 $A_0$ 形式平滑（19.9.1）。
\end{lemma}

\begin{proof}
事实上，已知
$\widehat A=A\otimes_{A_0}\widehat A_0$
（Bourbaki，\emph{Alg. comm.}，第 IV 章，§~2，
n\textsuperscript{o}~5，命题 9 的推论 3，以及第 III 章，§~3，
n\textsuperscript{o}~4，定理 1），所以结论由 (19.9.3) 得出。
\end{proof}

\begin{proposition}[22.3.2]
\label{0.22.3.2-fr}
设 $A$ 是正则 Noether 局部环，$K$ 是它的分式域，$p$ 是 $K$ 的特征。
假设下列假设之一成立：
\begin{enumerate}
  \item[(i)] $p=0$。
  \item[(ii)] $p>0$，并存在一个递减有向族 $(A_\alpha)$，它由
  Noether 子环（属于 $A$）组成，使得 $A$ 是有限 $A_\alpha$-代数
  （对每个 $\alpha$）；有 $A^{p^{h(\alpha)}}\subset A_\alpha$，其中
  $h(\alpha)$ 是某个适当的整数；并且，若以 $K_\alpha$ 表示
  $A_\alpha$ 的分式域，则
  \[
    \bigcap_\alpha K_\alpha(K^p)=K^p.
  \]
\end{enumerate}
于是，设 $B$ 是一个整的有限 $A$-代数并且包含 $A$，$\mathfrak n$ 是
$B$ 的一个素理想，$C$ 是局部环 $B_\mathfrak n$，$\mathfrak q$ 是完备化
$\widehat C$ 的一个素理想，并满足 $\mathfrak q\cap C=0$；于是局部环
$\widehat C_\mathfrak q$ 是分式域 $L$（属于 $B$）上的代数。那么
$\widehat C_\mathfrak q$ 是形式平滑的 $L$-代数（对于其
$\mathfrak q$-进拓扑），因而是在 $L$ 上几何正则的环（19.6.5）。
\end{proposition}

\begin{proof}
分两种情形。

I) 假设 $\mathfrak n$ 包含极大理想 $\mathfrak m$（属于 $A$），因而
$\mathfrak n\cap A=\mathfrak m$。于是 $\mathfrak n$ 是极大理想
（Bourbaki，\emph{Alg. comm.}，第 V 章，§~2，
n\textsuperscript{o}~1，命题 1）；若 $\mathfrak r$ 是半局部环 $B$ 的
根，则 $\widehat C$ 是 $B$ 对 $\mathfrak n$-预进拓扑的分离完备化，
并且是半局部环 $\widehat B$ 的一个分量；后者是 $B$ 对
$\mathfrak r$-预进拓扑的完备化（Bourbaki，\emph{Alg. comm.}，
第 III 章，§~3，n\textsuperscript{o}~4，命题 8）。因此可写成
$\widehat C_\mathfrak q=\widehat B_{\mathfrak q'}$，其中
$\mathfrak q'$ 是 $\widehat B$ 中取逆像为 $\mathfrak q$ 的素理想。
现在注意，$B$ 是 $\widehat B$ 的子环，并且假设
$\mathfrak q\cap C=0$ 蕴含 $\mathfrak q'\cap B=0$；故
$\widehat B_{\mathfrak q'}$ 也是环 $\widehat B\otimes_B L$ 在它的一个
素理想 $\mathfrak q''$ 处的局部化，而 $\mathfrak q'$ 是该素理想在
$\widehat B$ 中的逆像。此外，有
$\widehat B=\widehat A\otimes_A B$
（Bourbaki，\emph{Alg. comm.}，第 IV 章，§~2，
n\textsuperscript{o}~5，命题 9 的推论，以及第 III 章，§~3，
n\textsuperscript{o}~4，定理 1），所以
$\widehat B\otimes_B L=\widehat A\otimes_A L$。最后，若
$\mathfrak p$ 是 $\widehat A$ 中取逆像为 $\mathfrak q''$ 的素理想，
则 $A\to B$ 的单射性以及交换图
\[
  \xymatrix{
    B \ar[r] & \widehat B\\
    A \ar[u] \ar[r] & \widehat A \ar[u]
  }
\]
蕴含 $\mathfrak p\cap A=0$；故
$(\widehat A\otimes_A L)_{\mathfrak q''}$ 也是环
\[
  \widehat A_\mathfrak p\otimes_A L
  =(\widehat A_\mathfrak p\otimes_A K)\otimes_K L
\]
在它的一个素理想处的局部化。

\oldpage[0\textsubscript{IV-1}]{192}

因此，为证明这个局部环是形式平滑的 $L$-代数，只须证明
$\widehat A_\mathfrak p\otimes_A K$ 是形式平滑的 $K$-代数
（19.3.5, (iii) 和 (iv)）；而且
$\widehat A_\mathfrak p\otimes_A K=\widehat A_\mathfrak p$，因为
$\mathfrak p\cap A=0$。我们应用判别法 (22.2.2) 和 (22.2.8)。
首先，因 $A$ 正则，$\widehat A$ 也正则（17.1.5），从而
$\widehat A_\mathfrak p$ 也正则（17.3.2）；所以 (22.2.2) 的条件 (i)
成立（17.1.1）。于是只须（当 $p>0$ 时）验证 (22.2.8) 的条件 b)。
而 (22.3.1) 表明，对每个 $\alpha$，$\widehat A$ 是一个 $A$-代数，
并且相对于 $A_\alpha$ 形式平滑（对于离散拓扑）；因此
\[
  \widehat A\otimes_{A_\alpha}K_\alpha
  =\widehat A\otimes_A(A\otimes_{A_\alpha}K_\alpha)
\]
是一个形式平滑的
$(A\otimes_{A_\alpha}K_\alpha)$-代数，并且相对于 $K_\alpha$ 形式平滑
（对于离散拓扑）（19.9.2, (iii)）。但因 $A$ 是有限
$A_\alpha$-代数，有 $A\otimes_{A_\alpha}K_\alpha=K$；又因
$\widehat A_\mathfrak p$ 是 $\widehat A\otimes_A K$ 的一个分式环，
由 (19.9.2, (iv)) 可知，$\widehat A_\mathfrak p$ 是一个 $K$-代数，
并且相对于 $K_\alpha$ 形式平滑（对于离散拓扑），因而\emph{更是}对其
$\mathfrak p$-进拓扑如此（(19.3.8) 和 (19.9.5)）。这正是 (22.2.8)
的条件 b)，故在 $p>0$ 的这种情形命题得证。当 $p=0$ 时，
$\widehat A_\mathfrak p$ 的剩余域在 $K$ 上可分；又因
$\widehat A_\mathfrak p$ 是正则环，由 (19.6.4)，它是形式平滑的
$K$-代数。

II) 一般情形。令 $\mathfrak s$ 为素理想 $\mathfrak n\cap A$
（属于 $A$），并置 $S=A-\mathfrak s$；于是
$S^{-1}A=A_\mathfrak s$，且
$C=(S^{-1}B)_{S^{-1}\mathfrak n}$。这一次，$S^{-1}\mathfrak n$
包含极大理想 $\mathfrak sA_\mathfrak s$（属于 $A_\mathfrak s$）。由于
$A_\mathfrak s$ 正则（17.3.2），且 $S^{-1}B$ 是一个整的有限
$A_\mathfrak s$-代数并包含 $A_\mathfrak s$，一切归结为对
$A_\mathfrak s$ 验证条件 (ii)，这里是在 $p>0$ 时。为此，置
$\mathfrak s_\alpha=\mathfrak s\cap A_\alpha$，并考虑局部环
$(A_\alpha)_{\mathfrak s_\alpha}$。对每个
$t\notin\mathfrak s$，按假设，置 $q=p^{h(\alpha)}$，则
$t^q\in A_\alpha$ 且 $t^q\notin\mathfrak s$，故
$t^q\notin\mathfrak s_\alpha$。若置
$S_\alpha=A_\alpha-\mathfrak s_\alpha$，由此可见
$A_\mathfrak s=S_\alpha^{-1}A$；并且假设蕴含 $A_\mathfrak s$ 是有限
$(A_\alpha)_{\mathfrak s_\alpha}$-代数。此外，$K_\alpha$ 也是
$(A_\alpha)_{\mathfrak s_\alpha}$ 的分式域，这就完成了证明。
\end{proof}

\begin{theorem}[22.3.3]
\label{0.22.3.3-fr}
设 $A$ 是完备 Noether 局部环，$\mathfrak p$ 是 $A$ 的一个素理想，
$B$ 是局部化环 $A_\mathfrak p$，$\mathfrak q$ 是完备化
$\widehat B$ 的一个素理想，且 $\mathfrak r=\mathfrak q\cap B$。那么，
局部环 $\widehat B_\mathfrak q$（即 $\widehat B$ 的局部化环）对预进拓扑是形式平滑的
$B_\mathfrak r$-代数。
\end{theorem}

\begin{proof}
素理想 $\mathfrak r$（属于 $B$）具有形式 $\mathfrak n_\mathfrak p$，其中
$\mathfrak n\subset\mathfrak p$ 是 $A$ 的一个素理想，并且
$B_\mathfrak r=A_\mathfrak n$。令 $k$ 为 $B_\mathfrak r$ 的剩余域；
它也是整完备局部环 $A'=A/\mathfrak n$ 的分式域。由于 $\widehat B$ 是
平坦 $B$-模 $(\mathbf 0_\mathrm I,~7.3.5)$，
$\widehat B_\mathfrak q$ 是平坦 $B_\mathfrak r$-模
$(\mathbf 0_\mathrm I,~6.3.2)$；为证明 $\widehat B_\mathfrak q$ 是
形式平滑的 $B_\mathfrak r$-代数，只须证明
\[
  \widehat B_\mathfrak q\otimes_{B_\mathfrak r}k
  =\widehat B_\mathfrak q/\mathfrak r\widehat B_\mathfrak q
\]
是形式平滑的 $k$-代数（19.7.1）。而
$\widehat B_\mathfrak q/\mathfrak r\widehat B_\mathfrak q
=(\widehat B/\mathfrak r\widehat B)_{\mathfrak q'}$，其中
$\mathfrak q'$ 是 $\mathfrak q$ 在
$\widehat B/\mathfrak r\widehat B$ 中的典范像；按定义有
$\mathfrak q'\cap(B/\mathfrak r)=0$，且
$B/\mathfrak r=A_{\mathfrak p'}$，其中
$\mathfrak p'=\mathfrak p/\mathfrak n$。此外，
$\widehat B/\mathfrak r\widehat B$ 是 $B/\mathfrak r$ 的完备化。

因此问题归结为证明下列推论：
\end{proof}

\begin{corollary}[22.3.4]
\label{0.22.3.4-fr}
设 $A$ 是整的完备 Noether 局部环，$K$ 是其分式域，$\mathfrak p$ 是
$A$ 的一个素理想，$B$ 是局部化环 $A_\mathfrak p$。那么，对于每个素理想
$\mathfrak q$（属于完备化 $\widehat B$），

\oldpage[0\textsubscript{IV-1}]{193}

若 $\mathfrak q\cap B=0$，则局部环 $\widehat B_\mathfrak q$ 是形式平滑的
$K$-代数（对于其 $\mathfrak q$-进拓扑），因而是在 $K$ 上几何正则的环。
\end{corollary}

\begin{proof}
事实上，由 (19.8.9)，存在一个子环 $A_0$（属于 $A$）；它是特征 $p>0$ 的
某个域上或某个 Cohen 环上的形式幂级数环（回想 Cohen 环的分式域特征为
$0$），并且 $A$ 是有限 $A_0$-代数。已知环 $A_0$ 正则（17.3.8），且由
(21.8.8)，它满足 (22.3.2) 的假设 (i)、(ii) 之一。因此只须应用
(22.3.2)，把其中的 $B$ 换成 $A$，并把其中的 $A$ 换成 $A_0$。
\end{proof}

\subsection{极大理想平方为零的局部环之有限扩充的预备结果}
\label{subsection:0.22.4-fr}
\label{0.22.4-fr}

\begin{proposition}[22.4.1]
\label{0.22.4.1-fr}
设 $A$ 是一个局部环，其极大理想 $\mathfrak m$ 平方为零，
$K=A/\mathfrak m$ 是它的剩余域；用 $V$ 表示理想 $\mathfrak m$，
把它看成 $K$-向量空间。设 $T_i$ $(1\leqslant i\leqslant r)$ 是不定元，
并且对每个 $i$，$F_i$ 是 $A[T_i]$ 中次数为 $d_i$ 的首一多项式；用
$A'$ 表示多项式环 $A[T_1,\ldots,T_r]$ 对理想 $\mathfrak b$ 的商，
其中该理想由这 $r$ 个多项式 $F_i$ 生成。假设 $A'$ 是局部环，并设
$\mathfrak m'$ 是其
极大理想、$K'=A'/\mathfrak m'$ 是其剩余域。于是：
\begin{enumerate}
  \item[(i)] 若置 $d=\prod_{i=1}^r d_i$，则 $A'$ 是自由 $A$-模，
  其秩为 $d$，并且 $[K':K]\leqslant d$。
  \item[(ii)] 要使 $[K':K]=d$，当且仅当
  $A'\otimes_AK=A'/\mathfrak mA'$ 是一个域（必同构于 $K'$）；换言之，
  当且仅当 $\mathfrak m'=\mathfrak mA'$。在这种情形，
  $\mathfrak m'^2=0$；并且若以 $V'$ 表示理想 $\mathfrak m'$，把它看成
  $K'$-向量空间，则典范同态 $V\otimes_KK'\to V'$ 是双射（因而
  $\operatorname{rg}_{K'}V'=\operatorname{rg}_K V$）。
\end{enumerate}
\end{proposition}

\begin{proof}
即使不假设 $A'$ 局部，由欧几里得除法并对 $r$ 归纳，也显然可知：若
$t_i$ 是 $T_i$ 模 $\mathfrak b$ 的类，则单项式
$M_\nu=t_1^{\nu_1}t_2^{\nu_2}\cdots t_r^{\nu_r}$，其中
$0\leqslant\nu_i\leqslant d_i-1$，构成 $A$-模 $A'$ 的一个基。若再假设
$A'$ 局部，则 $K'$ 是
$A'/\mathfrak mA'=A'\otimes_AK$ 的一个商，故
$[K':K]\leqslant[A'\otimes_AK:K]=d$；并且等号只有在
$\mathfrak m'=\mathfrak mA'$ 时才可能成立。这证明了 (i) 和 (ii) 的
第一个断言。另一方面，显然当 $\mathfrak m'=\mathfrak mA'$ 时，
$\mathfrak m'^2=0$；而且
$V\otimes_KK'=\mathfrak m\otimes_AA'=\mathfrak mA'=V'$，因为 $A'$ 是
自由 $A$-模。
\end{proof}

\begin{remark}[22.4.2]
\label{0.22.4.2-fr}
当 $A$ 是 Artin 环时（即 $\operatorname{rg}_K(V)$ 有限），$A'$ 为
局部环这一假设总蕴含
$\operatorname{rg}_{K'}(\mathfrak m'/\mathfrak m'^2)
\geqslant\operatorname{rg}_K(V)$，如稍后 (22.5.2) 所见。注意，这两个秩
可以相等，而不必有 $\mathfrak m'=\mathfrak mA'$。
\end{remark}

\begin{proposition}[22.4.3]
\label{0.22.4.3-fr}
关于 $A$ 的假设和记号同 (22.4.1)；假设诸 $F_i$ 具有形式
\begin{equation}
\label{0.22.4.3.1-fr}
  F_i(T_i)=T_i^{d_i}+\sum_{1\leqslant k\leqslant d_i}
  c_{ik}T_i^{d_i-k},
  \tag{22.4.3.1}
\end{equation}
其中 $d_i\geqslant2$ 对每个 $i$ 都成立，且
$c_{ik}\in\mathfrak m$（“Eisenstein 多项式”）。于是：
\begin{enumerate}
  \item[(i)] $A'$ 是局部环，其剩余域
  $K'=A'/\mathfrak m'$ 同构于 $K$。

  \oldpage[0\textsubscript{IV-1}]{194}

  \item[(ii)] 典范同态 $V\to V'=\mathfrak m'/\mathfrak m'^2$ 的核，是
  理想 $\mathfrak n$（属于 $A$）；该理想由
  $\xi_i=c_{i,d_i}$（诸 $F_i$ 的“常数项”）生成，而这个同态的余核在
  $K'$ 上以诸 $T_i$ 的像为一个基。
  特别地，
  \begin{equation}
  \label{0.22.4.3.2-fr}
    \operatorname{rg}_{K'}(V')=\operatorname{rg}_K(V)+r-r',
    \tag{22.4.3.2}
  \end{equation}
  其中 $r'=\operatorname{rg}_K(\mathfrak n)$。当
  $\operatorname{rg}_K(V)$ 有限时（即当 $A$ 是 Artin 环时），要使
  $\operatorname{rg}_{K'}(V')=\operatorname{rg}_K(V)$，当且仅当诸
  $\xi_i$ 在 $K$ 上线性无关。
\end{enumerate}
\end{proposition}

\begin{proof}
(i) 显然，$A'\otimes_AK=A'/\mathfrak mA'$ 同构于
$K[T_1,\ldots,T_r]/\mathfrak b'$，其中 $\mathfrak b'$ 是由这 $r$ 个
多项式 $T_i^{d_i}$ 生成的理想；因此它是剩余域为 $K$ 的局部环，其极大
理想由诸 $T_i$ 模 $\mathfrak b'$ 的类生成，由此得到 (i)。

(ii) 令 $t_i$ $(1\leqslant i\leqslant r)$ 为 $T_i$ 模
$\mathfrak b$ 的类；由 (i)，极大理想 $\mathfrak m'$（属于 $A'$）为
\[
  \mathfrak m'=\mathfrak mA'+\sum_{i=1}^r t_iA',
\]
于是考虑 $\mathfrak m^2=0$，有
\[
  \mathfrak m'^2=\sum_{i,j}t_it_jA'+\sum_i t_i\mathfrak mA'.
\]
由此可知，$A'/\mathfrak m'^2$ 同构于环
$A[T_1,\ldots,T_r]/\mathfrak c$，其中 $\mathfrak c$ 是由下列元素生成的
理想：
\[
  F_i\ (1\leqslant i\leqslant r),\qquad
  T_iT_j\ (1\leqslant i\leqslant r,\ 1\leqslant j\leqslant r),\qquad
  xT_i\ (1\leqslant i\leqslant r,\ x\in\mathfrak m).
\]
由于 $d_i\geqslant2$ 且诸 $T_i^2$ 属于 $\mathfrak c$，上述假设蕴含
$\mathfrak c$ 也由下列元素生成：
\[
  \xi_i\ (1\leqslant i\leqslant r),\qquad
  T_iT_j\ (1\leqslant i\leqslant r,\ 1\leqslant j\leqslant r),\qquad
  T_ix\ (1\leqslant i\leqslant r,\ x\in\mathfrak m).
\]
令
\[
  \mathfrak c_1=\sum_{i,j}T_iT_jA[T_1,\ldots,T_r],\qquad
  \mathfrak c_2=\mathfrak c_1+\sum_i\mathfrak mT_iA[T_1,\ldots,T_r],
\]
于是 $\mathfrak c=\mathfrak c_2+\sum_i\xi_iA[T_1,\ldots,T_r]$。显然，
$A[T_1,\ldots,T_r]/\mathfrak c_1$ 作为 $A$-模，是 $A$ 与诸
$At_i^{(1)}$ 的直和，其中 $t_i^{(1)}$ 是 $T_i$ 模
$\mathfrak c_1$ 的类。由此推出，$A[T_1,\ldots,T_r]/\mathfrak c_2$
是 $A$ 与诸 $Kt_i^{(2)}$ 的直和，其中 $t_i^{(2)}$ 是 $T_i$ 模
$\mathfrak c_2$ 的类。最后，$A[T_1,\ldots,T_r]/\mathfrak c$ 是
$A/\mathfrak n$ 与诸 $Kt_i'$ 的直和，其中 $t_i'$ 是 $T_i$ 模
$\mathfrak c$ 的类，而 $\mathfrak n$ 是 $A$ 中由诸 $\xi_i$ 生成的理想。
在把 $A'/\mathfrak m'^2$ 与 $A[T_1,\ldots,T_r]/\mathfrak c$ 认同之后，
$V$ 的像与 $\mathfrak m/\mathfrak n$ 认同；所以上述论证证明了 (ii)。
\end{proof}

\begin{proposition}[22.4.4]
\label{0.22.4.4-fr}
关于 $A$ 的假设和记号同 (22.4.1)；假设 $K$ 的特征为 $p>0$（所以有
$p\cdot1\in\mathfrak m$，这是在 $A$ 中），并且诸 $F_i$ 具有形式
\begin{equation}
\label{0.22.4.4.1-fr}
  F_i(T_i)=T_i^p+p\sum_{k=1}^{p-1}c_{ik}T_i^{p-k}-f_i
  \tag{22.4.4.1}
\end{equation}
其中 $c_{ik}\in A$ 且 $f_i\in A$；这对每个 $i$ 和每个 $k$ 都成立，而
后者满足 $1\leqslant k\leqslant p-1$。此外，若 $a_i$ 是 $f_i$ 模
$\mathfrak m$ 的类，则假设族 $(a_i)_{1\leqslant i\leqslant r}$ 在
$p$-自由（在 $K^p$ 上，即“绝对 $p$-自由”）。于是：
\begin{enumerate}
  \item[(i)] $A'$ 是一个极大理想为 $\mathfrak mA'$ 的局部环，其剩余域
  $K'$ 同构于由 $K$ 添加诸 $a_i^{1/p}$
  $(1\leqslant i\leqslant r)$ 所得的扩域。

  \oldpage[0\textsubscript{IV-1}]{195}

  \item[(ii)] 诸元素 $d_Af_i\otimes1_{K'}$ 构成典范同态
  \[
    \Omega^1_A\otimes_AK'\longrightarrow
    \Omega^1_{A'}\otimes_{A'}K'
  \]
  的核 $N=\Upsilon^{K'}_{A'/A}$ 的一个基。
\end{enumerate}
\end{proposition}

\begin{proof}
(i) 这一次，$A'\otimes_AK=A'/\mathfrak mA'$ 同构于
$K[T_1,\ldots,T_r]/\mathfrak b'$，其中 $\mathfrak b'$ 是由诸多项式
$T_i^p-a_i$ 生成的理想；关于诸 $a_i$ 的假设蕴含这个商环是一个域
（21.1.8），由此得到 (i) 的各断言。

在证明 (ii) 之前，我们先证明下列

\begin{lemma}[22.4.4.2]
\label{0.22.4.4.2-fr}
设 $\Lambda$ 是一个环，$A$ 是一个 $\Lambda$-代数，$\mathfrak J$ 是
$A$ 的一个理想，使得环 $C=A/\mathfrak J$ 的特征为 $p>0$（这等价于说
$p\cdot1\in\mathfrak J$；此关系是在 $A$ 中）。令 $\pi$ 为 $p\cdot1$ 在
$\mathfrak J/\mathfrak J^2$ 中的典范像。如果 $C$ 对离散拓扑是形式平滑的
$(\Lambda/p\Lambda)$-代数，则有典范正合列
\begin{equation}
\label{0.22.4.4.3-fr}
  0\longrightarrow(\mathfrak J/\mathfrak J^2)/C\cdot\pi
  \longrightarrow\Omega^1_{A/\Lambda}\otimes_AC
  \longrightarrow\Omega^1_{C/\Lambda}\longrightarrow0.
  \tag{22.4.4.3}
\end{equation}
\end{lemma}

\begin{proof}[引理的证明]
事实上，置 $A'=A/pA$、$\mathfrak J'=\mathfrak J/pA$，于是
$C=A'/\mathfrak J'$。由于
$A'=A\otimes_\Lambda(\Lambda/p\Lambda)$，有典范同构
$\Omega^1_{A/\Lambda}\otimes_AA'=\Omega^1_{A'/\Lambda'}$，其中
$\Lambda'$ 表示环 $\Lambda/p\Lambda$（20.5.5）。现在把正合列
(20.5.14.1) 应用于形式平滑的 $\Lambda'$-代数
$C=A'/\mathfrak J'$，便得到正合列
\[
  0\longrightarrow\mathfrak J'/\mathfrak J'^2
  \longrightarrow\Omega^1_{A'/\Lambda'}\otimes_{A'}C
  \longrightarrow\Omega^1_{C/\Lambda'}\longrightarrow0.
\]
又因同态 $\Lambda\to\Lambda'$ 是满射，由 (20.5.15) 有典范同构
$\Omega^1_{C/\Lambda'}=\Omega^1_{C/\Lambda}$；结合以上论证，便得到正合列
(22.4.4.3)。
\end{proof}

(ii) 由于 $\mathfrak m'^2=0$，仍用 $V'$ 表示理想 $\mathfrak m'$，
把它看成 $K'$-向量空间，并置
\begin{equation}
\label{0.22.4.4.4-fr}
  V_0=V/K\cdot p=\mathfrak m/(\mathfrak m^2+pA),\qquad
  V_0'=V'/K'\cdot p=\mathfrak m'/(\mathfrak m'^2+pA').
  \tag{22.4.4.4}
\end{equation}
把正合列 (22.4.4.3) 应用于如下情形：$\Lambda=\mathbf Z$；取
$\mathbf Z$-代数 $A$（相应地，$A'$）及其理想 $\mathfrak m$
（相应地，$\mathfrak m'$）。由于 $K$（相应地，$K'$）在素域
$\mathbf Z/p\mathbf Z$ 上可分，因而是形式平滑的
$(\mathbf Z/p\mathbf Z)$-代数（19.6.1），便得到正合列
\begin{equation}
\label{0.22.4.4.5-fr}
  0\longrightarrow V_0\longrightarrow\Omega^1_A\otimes_AK
  \longrightarrow\Omega^1_K\longrightarrow0,\qquad
  0\longrightarrow V_0'\longrightarrow\Omega^1_{A'}\otimes_{A'}K'
  \longrightarrow\Omega^1_{K'}\longrightarrow0.
  \tag{22.4.4.5}
\end{equation}
现在考虑图 (22.4.4.6)，其中两个中间列分别是 (22.4.4.5) 的第二个序列，
以及第一个序列与 $K'$ 的张量积；由于 (22.4.4.5) 的第一个正合列由
$K$ 上的向量空间组成，(22.4.4.6) 的两个中间列正合。此外，由
(22.4.4.2) 的证明以及图 (20.5.11.3) 的交换性，这两列与连接它们的横箭头
所成的图交换。(22.4.4.6) 的最后一行，是在 (20.6.1.1) 中把
$A$、$B$、$C$ 分别换成 $\mathbf Z$、$K$ 和 $K'$ 后所得的正合列；
中间一行是以 $K'$ 与如下 $K$-模正合列作张量所得：该正合列来自
(20.6.1.1)，其中把 $A$、$B$、$C$ 分别换成 $\mathbf Z$、$A$、$K$。由图
(20.6.4.3) 的交换性，这两行与连接它们的竖箭头所成的图交换。另一方面，
注意这里满足 (22.4.1, (ii)) 的应用条件，所以图 (22.4.4.6) 的最上一行
正合，而图的两个端列

\oldpage[0\textsubscript{IV-1}]{196}

\begin{equation}
\label{0.22.4.4.6-fr}
  \xymatrix{
    && 0 \ar[d] & 0 \ar[d] && \\
    & 0 \ar[r] & V_0\otimes_KK' \ar[r] \ar[d] & V_0' \ar[r] \ar[d] & 0 & \\
    0 \ar[r] & N \ar[r] \ar[d] & \Omega^1_A\otimes_AK' \ar[r] \ar[d] &
      \Omega^1_{A'}\otimes_{A'}K' \ar[r] \ar[d] &
      \Omega^1_{A'/A}\otimes_{A'}K' \ar[r] \ar[d] & 0 \\
    0 \ar[r] & \Upsilon_{K'/K} \ar[r] \ar[d] &
      \Omega^1_K\otimes_KK' \ar[r] \ar[d] & \Omega^1_{K'} \ar[r] \ar[d] &
      \Omega^1_{K'/K} \ar[r] \ar[d] & 0 \\
    & 0 & 0 & 0 & 0 &
  }
  \tag{22.4.4.6}
\end{equation}
分别由中间横箭头的核和余核组成，这也最终证明了该图交换。此外：

\begin{lemma}[22.4.4.7]
\label{0.22.4.4.7-fr}
图 (22.4.4.6) 中的同态
\[
  N\longrightarrow\Upsilon_{K'/K}\qquad\text{和}\qquad
  \Omega^1_{A'/A}\otimes_{A'}K'\longrightarrow\Omega^1_{K'/K}
\]
都是双射。
\end{lemma}

\begin{proof}
事实上，把蛇形图（Bourbaki，\emph{Alg. comm.}，第 I 章，§~1，
n\textsuperscript{o}~4，命题 2）应用于 (22.4.4.6) 的两个中间列，
得到正合列
\[
  0\longrightarrow N\longrightarrow\Upsilon_{K'/K}\longrightarrow0
  \longrightarrow\Omega^1_{A'/A}\otimes_{A'}K'
  \longrightarrow\Omega^1_{K'/K}\longrightarrow0.
\]
而若 $t_i$ 是 $T_i$ 在 $A'$ 中的典范像，则关系
$F_i(t_i)=0$，连同在 $A'$ 中有
$p\cdot1\in\mathfrak mA'\subset\mathfrak m'$ 这一事实，立即表明
$d_{A'}f_i\in\mathfrak m'\Omega^1_{A'}$，所以
$d_{A'}f_i\otimes1_{K'}=0$；这证明诸 $d_Af_i\otimes1_{K'}$ 属于
$N$。此外，它们在 $\Upsilon_{K'/K}$ 中的像是诸
$d_Ka_i\otimes1_{K'}$；由 (21.4.7)，后者构成 $\Upsilon_{K'/K}$ 在
$K'$ 上的一个基。结合 (22.4.4.7)，这就完成了 (22.4.4) 的证明。
\end{proof}
\end{proof}

\begin{env}[22.4.5]
\label{0.22.4.5-fr}
现在考虑两个环 $A$、$B$，它们的\emph{特征}为 $p>0$；并考虑两个同态
\[
  A\xrightarrow{i}B\xrightarrow{j}A
\]

\oldpage[0\textsubscript{IV-1}]{197}

它们满足 (21.3.1) 的条件。再设 $\mathfrak m$ 是 $A$ 中一个
\emph{平方为零}的理想；置
\[
  K=A/\mathfrak m,\qquad
  B_{(K)}=B\otimes_AK=B/Bi(\mathfrak m)
\]
并令
\[
  \varphi:A\longrightarrow K,\qquad
  \psi:B\longrightarrow B_{(K)}
\]
为典范同态。

由于 $\varphi$ 是满射，$\Omega^1_{B_{(K)}/A}$ 与
$\Omega^1_{B_{(K)}/K}$ 典范认同（20.5.15）。注意，因
$p\geqslant2$，对每个 $x\in\mathfrak m$，有
$j(i(x))=x^p=0$；换言之，$j(i(\mathfrak m))=0$。所以从 $j$ 取商可得
一个同态
\[
  j':B_{(K)}\longrightarrow A
\]
使得，若 $i'=\psi\circ i:A\longrightarrow B_{(K)}$，则
\[
  j'(i'(x))=x^p\qquad\text{且}\qquad i'(j'(y))=y^p
  \qquad\text{对 }x\in A\text{ 和 }y\in B_{(K)}.
\]
按照 (21.3.2) 的记号，置
\[
  \Theta_{B_{(K)}/A}
  =\Omega^1_{B_{(K)}/A}\otimes_{B_{(K)}}A_{[j']}
  =\Omega^1_{B_{(K)}/K}\otimes_{B_{(K)}}A_{[j']},
\]
于是由 (21.3.2.2)，有典范同态
\begin{equation}
\label{0.22.4.5.1-fr}
  \pi_{B_{(K)}/A}:\Theta_{B_{(K)}/A}\longrightarrow\Omega^1_A.
  \tag{22.4.5.1}
\end{equation}

另一方面，从 $i$ 和 $j$ 取商（注意 $j(i(\mathfrak m))=0$），得到同态
\[
  K\xrightarrow{i_0}B_{(K)}\xrightarrow{j_0}K
\]
从而有交换图
\[
  \xymatrix{
    A \ar[r]^i \ar[d]_\varphi & B \ar[r]^j \ar[d]^\psi & A \ar[d]^\varphi \\
    K \ar[r]_{i_0} & B_{(K)} \ar[r]_{j_0} \ar[ur]_{j'} & K
  }
\]
而且显然有
\[
  j_0(i_0(s))=s^p\qquad\text{且}\qquad i_0(j_0(t))=t^p
  \qquad\text{对 }s\in K\text{ 和 }t\in B_{(K)}.
\]
因此同样定义
\[
  \Theta_{B_{(K)}/K}
  =\Omega^1_{B_{(K)}/K}\otimes_{B_{(K)}}K_{[j_0]}
  =\Theta_{B_{(K)}/A}\otimes_AK,
\]
所以把 (22.4.5.1) 与 $K$ 作张量，便得到一个 $K$-\emph{典范同态}
\begin{equation}
\label{0.22.4.5.2-fr}
  \pi'_{B_{(K)}/K}=\pi_{B_{(K)}/A}\otimes1_K:
  \Theta_{B_{(K)}/K}\longrightarrow\Omega^1_A\otimes_AK.
  \tag{22.4.5.2}
\end{equation}

\oldpage[0\textsubscript{IV-1}]{198}

由定义立即可知图
\[
  \xymatrix{
    \Theta_{B_{(K)}/K} \ar[r]^{\pi'_{B_{(K)}/K}} \ar[dr]_{\pi_{B_{(K)}/K}}
      & \Omega^1_A\otimes_AK \ar[d]^{\varphi_{K/A}} \\
    & \Omega^1_K
  }
\]
交换；其中 $\pi_{B_{(K)}/K}$ 是对应于 $i_0$ 和 $j_0$ 的典范同态
（21.3.2.2）。在
$\Xi_{B_{(K)}/K}=\operatorname{Ker}(\pi_{B_{(K)}/K})$ 上取
$\pi'_{B_{(K)}/K}$ 的\emph{限制}，就定义了一个
\emph{典范同态}
\begin{equation}
\label{0.22.4.5.3-fr}
  \Xi_{B_{(K)}/K}\longrightarrow
  \operatorname{Ker}(\Omega^1_A\otimes_AK\longrightarrow\Omega^1_K)=\Upsilon_{K/A}.
  \tag{22.4.5.3}
\end{equation}
\end{env}

\begin{env}[22.4.6]
\label{0.22.4.6-fr}
假设和记号同 (22.4.5)；再假设 $\mathfrak m$ 是一个\emph{平方为零的极大理想}，
所以 $K$ 是一个\emph{域}，并以 $V$ 表示理想 $\mathfrak m$，把它看成
$K$-向量空间。由于 $K$ 是其素域 $\mathbf F_p$ 上的形式平滑代数
（21.5.2），而 $A$ 是 $\mathbf F_p$ 上的代数（从而
$\Omega^1_A=\Omega^1_{A/\mathbf F_p}$），由 (20.5.14) 得到一个
\emph{正合列}
\begin{equation}
\label{0.22.4.6.1-fr}
  0\longrightarrow V\longrightarrow\Omega^1_A\otimes_AK
  \longrightarrow\Omega^1_K\longrightarrow0.
  \tag{22.4.6.1}
\end{equation}
所以由 (22.4.5.3)，有一个记作
\begin{equation}
\label{0.22.4.6.2-fr}
  \chi'_{B/A}:\Xi_{B_{(K)}/K}\longrightarrow V
  \tag{22.4.6.2}
\end{equation}
的同态；我们称它为 $A$-代数 $B$ 的\emph{特征同态}（相对于 $i$、
$j$ 和 $\mathfrak m$）。由定义可知，$\chi'_{B/A}$ 是\emph{单射}，
当且仅当 $\pi'_{B_{(K)}/K}$ 是单射，因为后者的核显然包含于
$\Xi_{B_{(K)}/K}$。
\end{env}

\begin{proposition}[22.4.7]
\label{0.22.4.7-fr}
设 $A$ 是一个 Artin 局部环，其极大理想 $\mathfrak m$ 平方为零，
$K=A/\mathfrak m$ 是其剩余域，$V$ 是理想 $\mathfrak m$，把它看成
$K$-向量空间；假设 $A$ 的特征为 $p>0$。设
\begin{equation}
\label{0.22.4.7.1-fr}
  F_i(T_i)=T_i^p-f_i\qquad\text{其中 }f_i\in A\quad(1\leqslant i\leqslant r)
  \tag{22.4.7.1}
\end{equation}
并用 $B$ 表示 $A[T_1,\ldots,T_r]$ 对由这 $r$ 个多项式 $F_i$ 生成的理想
所作的商环。于是：
\begin{enumerate}
  \item[\rm(i)] $B$ 是局部环。
  \item[\rm(ii)] 若 $\mathfrak m'$ 是 $B$ 的极大理想，
  $K'=B/\mathfrak m'$ 是其剩余域，$V'$ 是 $B$-模
  $\mathfrak m'/\mathfrak m'^2$，把它看成 $K'$-向量空间，则要使
  $\operatorname{rg}_K(V)=\operatorname{rg}_{K'}(V')$，当且仅当
  $K$-向量空间的特征同态
  \begin{equation}
  \label{0.22.4.7.2-fr}
    \chi'_{B/A}:\Xi_{B_{(K)}/K}\longrightarrow V
    \tag{22.4.7.2}
  \end{equation}
  （参见 (22.4.6.2)）是单射。
\end{enumerate}
\end{proposition}

\begin{proof}
仍以 $a_i$ $(1\leqslant i\leqslant r)$ 表示 $f_i$ 模
$\mathfrak m$ 的类；诸 $a_i$ 不再必然 $p$-无关（在 $K^p$ 上），但总可以假设
子族

\oldpage[0\textsubscript{IV-1}]{199}

$(a_i)_{s+1\leqslant i\leqslant r}$ 是 $p$-自由的（在 $K^p$ 上），并且
是一个 $p$-基底（在 $K^p$ 上），所生成的域为
$K^p(a_1,\ldots,a_r)$；从而有
\[
  a_i\in K^p(a_{s+1},\ldots,a_r)\qquad\text{对 }1\leqslant i\leqslant s.
\]
用 $A''$ 表示 $A[T_{s+1},\ldots,T_r]$ 对由诸 $F_i$ 生成的理想所作的
商环，其中指标满足 $i\geqslant s+1$。由 (22.4.4)，
$A''$ 是一个局部环，其极大理想
$\mathfrak m''=\mathfrak mA''$（因而平方为零），而剩余域
$K''=K(a_{s+1}^{1/p},\ldots,a_r^{1/p})$。所以有
$a_i\in K''^p$（$1\leqslant i\leqslant s$），从而在 $A''$ 中，元素
$f_i$（$1\leqslant i\leqslant s$）可写成
\begin{equation}
\label{0.22.4.7.3-fr}
  f_i=g_i^p+h_i
  \tag{22.4.7.3}
\end{equation}
其中 $g_i\in A''$ 且 $h_i\in\mathfrak m''$（因为
$\mathfrak m''^2=0$ 且 $p\mathfrak m''=0$，显然这些条件
\emph{唯一地}确定 $h_i$）。把 $T_i$ 换成 $T_i+g_i$，可见若置
\[
  G_i(T_i)=T_i^p+\sum_{k=1}^{p-1}\binom pk g_i^{p-k}T_i^k-h_i
  \qquad(1\leqslant i\leqslant s),
\]
则 $B$ 是 $A''[T_1,\ldots,T_s]$ 对由诸 $G_i$ 生成的理想所作的商环。
而 $G_i$ 除首项系数以外的所有系数都属于 $\mathfrak m''$，所以我们处于
(22.4.3) 的情形；$B$ 是局部环，且其剩余域 $K'$ \emph{同构于} $K''$。
这证明了 (i)（事实上也可从 $B^p\subset A$ 直接得出：因此 $B$ 中只有
一个理想位于 $\mathfrak m$ 之上）。

现在注意，(22.4.7.2) 是单射，当且仅当
\[
  \pi'_{B_{(K)}/K}:\Omega^1_{B_{(K)}/K}\otimes_{B_{(K)}}K
  \longrightarrow\Omega^1_A\otimes_AK
\]
是单射（由 (22.4.6)）。而 $B_{(K)}=B/\mathfrak mB$ 是代数
$C=K[T_1,\ldots,T_r]$ 对由多项式
$\overline F_i=T_i^p-a_i$ 生成的理想所作的商。由于
$d_{C/K}\overline F_i=0$，由 (20.5.13) 可知
$\Omega^1_{B_{(K)}/K}$ 是自由 $B_{(K)}$-模，其\emph{基}为诸
$d_{B_{(K)}/K}t_i$，其中 $t_i$ 是 $T_i$ 的典范像
$(1\leqslant i\leqslant r)$。但按定义 (22.4.5)，通过
$j':B_{(K)}\longrightarrow A$，一个模 $\mathfrak mB$ 的类（由元素
$x\in B$ 表示）的典范像，是元素 $x^p\in A$；因 $t_i^p=f_i$，立即可知
$d_{B_{(K)}/K}t_i\otimes1$ 通过 $\pi'_{B_{(K)}/K}$ 的像，是
$d_Af_i\otimes1$，这是 $\Omega^1_A\otimes_AK$ 中的元素。因此，
$\pi'_{B_{(K)}/K}$ 为单射，等价于诸 $d_Af_i\otimes1_K$ 在 $K$ 上
\emph{线性无关}；又等价于它们的像 $d_Af_i\otimes1_{K'}$ 在
$K'$ 上\emph{线性无关}（其所在空间为 $\Omega^1_A\otimes_AK'$）。

现在把 (22.4.4) 应用于 $A$-代数 $A''$，可知典范同态
$\Omega^1_A\otimes_AK'\longrightarrow\Omega^1_{A''}\otimes_{A''}K'$
的核，以诸 $d_Af_i\otimes1_{K'}$（其中
$s+1\leqslant i\leqslant r$）为基。因此，前一条件又等价于诸
$d_{A''}f_i\otimes1_{K'}$ 在 $\Omega^1_{A''}\otimes_{A''}K'$ 中线性无关
（其中 $1\leqslant i\leqslant s$）。而由 (22.4.7.3)，有
$d_{A''}f_i=d_{A''}h_i$。另一方面，若 $V''$ 是 $K'$-向量空间
$\mathfrak m''$，则
$\operatorname{rg}_{K'}(V'')=\operatorname{rg}_K(V)$
（22.4.1, (ii)）；所以条件
$\operatorname{rg}_{K'}(V')=\operatorname{rg}_K(V)$ 等价于
$\operatorname{rg}_{K'}(V'')=\operatorname{rg}_{K'}(V')$。但由
(22.4.3)，后一关系等价于诸 $h_i$ 的类

\oldpage[0\textsubscript{IV-1}]{200}

在 $V''$ 中在 $K'$ 上\emph{线性无关}。而在
$\Omega^1_{A''}\otimes_{A''}K'$ 中，诸
$d_{A''}h_i\otimes1_{K'}$ 是诸 $h_i$ 在单射
$V''\longrightarrow\Omega^1_{A''}\otimes_{A''}K'$ 下的典范像
（20.5.11.2；参见 (22.4.6.1)）。这就完成了 (22.4.7) 的证明。
\end{proof}

我们还证明了：

\begin{corollary}[22.4.7.4]
\label{0.22.4.7.4-fr}
要使 $\operatorname{rg}_K(V)=\operatorname{rg}_{K'}(V')$，当且仅当诸元素
$d_Af_i\otimes1_K$ $(1\leqslant i\leqslant r)$ 在 $K$-向量空间
$\Omega^1_A\otimes_AK$ 中线性无关。
\end{corollary}

\begin{remark}[22.4.8]
\label{0.22.4.8-fr}
可以在稍有不同的条件下定义同态 (22.4.5.2) 和 (22.4.6.2)。设 $A$ 是
一个环，$\mathfrak m$ 是 $A$ 中一个\emph{平方为零}的理想，
$K=A/\mathfrak m$；设 $p$ 是一个素数，满足 $p\cdot1\in\mathfrak m$
（换言之，$K$ 的特征为 $p$，但 $A$ 不必如此）。另一方面，设 $B$ 是
一个 $A$-代数，并且是\emph{忠实平坦} $A$-模（所以 $A\subset B$）；
假设对每个 $y\in B$ 都有
\begin{equation}
\label{0.22.4.8.1-fr}
  y^p\in A+pB\subset A+\mathfrak mB
  \tag{22.4.8.1}
\end{equation}
。如果 $y^p=x+pz$，其中 $x\in A$、$z\in B$，则模
$(A\cap pB)$ 所得的\emph{类}（由 $x$ 表示）被良好确定；并且有
$A\cap pB=pA$，因为 $B$ 是忠实平坦 $A$-模。所以由此定义了一个典范映射
$j:B\longrightarrow A/pA$，使得对 $x\in A$，$j(x)$ 是 $x^p$ 模
$pA$ 的类。再者，若 $y'\in B$ 且
$y'^p=x'+pz'$，其中 $x'\in A$、$z'\in B$，则立即可知元素
$(y+y')^p-x-x'$ 和 $(yy')^p-xx'$ 属于 $pB$；这是因为二项式系数
$\binom pi$ 是 $p$ 的倍数（$1\leqslant i\leqslant p-1$）。因此映射
$j$ 是一个\emph{环同态}。作复合，从 $j$ 得到同态
$j':B\xrightarrow{j}A/pA\longrightarrow A/\mathfrak m=K$，继而得到同态
$j_0:B_{(K)}\longrightarrow K$，使 $j'$ 是复合
$B\longrightarrow B_{(K)}=B\otimes_AK\xrightarrow{j_0}K$。此外，若
$i_0:K\longrightarrow B_{(K)}$ 是典范映射，则 $i_0$ 和 $j_0$ 对
\emph{特征为 $p$} 的环 $K$ 和 $B_{(K)}$ 满足 (21.3.1) 的条件。
在这种情况下，为定义一个 $K$-同态
$\pi'_{B_{(K)}/K}:\Omega^1_{B_{(K)}/K}\otimes_{B_{(K)}}K_{[j_0]}
\longrightarrow\Omega^1_A\otimes_AK$，由 (20.4.8)，只须定义一个
$K$-导射
\begin{equation}
\label{0.22.4.8.2-fr}
  D_0:B_{(K)}\longrightarrow\Omega^1_A\otimes_AK
  \tag{22.4.8.2}
\end{equation}
；这里把右端看成 $B_{(K)}$-模，其结构由 $j_0$ 给出。为此，只须定义一个
$A$-导射
\begin{equation}
\label{0.22.4.8.3-fr}
  D:B\longrightarrow\Omega^1_A\otimes_AK
  \tag{22.4.8.3}
\end{equation}
，使其\emph{在} $\mathfrak mB$ \emph{上为零}（这里把右端看成 $B$-模，
其结构由 $j'$ 给出）。若有关系 $y^p=x+pz$，其中 $y\in B$、$x\in A$、
$z\in B$，则元素 $d_Ax\otimes1_K$（属于
$\Omega^1_A\otimes_AK$）只依赖于类 $j(y)$；这里按模 $pA$ 取类，而该类
由 $x$ 表示（因而只依赖于 $y$）。这是因为对每个 $t\in A$，有
$d_A(pt)=p\,d_A(t)$，所以在
$d_A(pt)\otimes1=0$；此等式位于
$\Omega^1_A\otimes_AK=\Omega^1_A/\mathfrak m\Omega^1_A$ 中，因为
$p\cdot1\in\mathfrak m$。因此可取
\[
  D(y)=d_Ax\otimes1_K.
\]
$D$ 是导射这一事实，容易由 $j$ 为环同态推出。此外，对 $t\in A$，有
$d_A(t^px)=t^p d_A(x)+pt^{p-1}x d_A(t)$，由此推出
$D(ty)=tD(y)$，故 $D$ 是 $A$-导射。最后，若再有
$t\in\mathfrak m$，则

\oldpage[0\textsubscript{IV-1}]{201}

$D(ty)=tD(y)=0$，因为 $\Omega^1_A\otimes_AK$ 被 $\mathfrak m$ 零化。
这样便定义了所求的同态 (22.4.5.2)；并且立即可验证复合同态
\[
  \Theta_{B_{(K)}/K}\xrightarrow{\pi'_{B_{(K)}/K}}
  \Omega^1_A\otimes_AK\longrightarrow\Omega^1_K
\]
仍是典范同态 $\pi_{B_{(K)}/K}$（相对于 $i_0$ 和 $j_0$）。因此也有一个
类似于 (22.4.5.3) 的典范同态。

现在若 $K$ 是一个\emph{域}，则 $A/pA$ 是 $\mathbf F_p$-代数，并有
正合列 (20.5.14)
\[
  0\longrightarrow V/pA\longrightarrow
  \Omega^1_{A/pA}\otimes_{A/pA}K\longrightarrow\Omega^1_K\longrightarrow0
\]
；另一方面有正合列 (20.5.12.1)
\[
  pA\longrightarrow\Omega^1_A\otimes_A(A/pA)
  \longrightarrow\Omega^1_{A/pA}\longrightarrow0,
\]
与 $K$ 作张量后，它表明 $K$-向量空间
$\Omega^1_A\otimes_AK$ 与
$\Omega^1_{A/pA}\otimes_{A/pA}K$ 同构。因此，(22.4.6.2) 在这里的类似物
是同态
\begin{equation}
\label{0.22.4.8.4-fr}
  \Xi_{B_{(K)}/K}\longrightarrow V/pA=\mathfrak m/(\mathfrak m^2+pA).
  \tag{22.4.8.4}
\end{equation}

如此，\emph{判别法 (22.4.7) 仍然有效，只须把其中的 $V$ 和 $V'$ 分别换成
$V/pA$ 和 $V'/pB$，并且只假设 $K$ 的特征为 $p>0$}（换言之，有
$p\cdot1\in\mathfrak m$ 于 $A$ 中），\emph{以及诸 $F_i(T_i)$ 具有更一般的形式
(22.4.4.1)}。事实上，$B$ 是自由 $A$-模（因而忠实平坦）；只须重复
(22.4.7) 的论证，把其中的 (22.4.6.2) 换成 (22.4.8.4)，并把 $V''$ 换成
$V''/pA''$。
\end{remark}

\subsection{几何正则代数与形式平滑代数}
\label{subsection:0.22.5-fr}

\begin{proposition}[22.5.1]
\label{0.22.5.1-fr}
设 $A$ 是正则局部环，$A'$ 是包含 $A$ 的局部环，并且是有限
$A$-代数；$\mathfrak m$（相应地 $\mathfrak m'$）是 $A$（相应地
$A'$）的极大理想，$K=A/\mathfrak m$（相应地
$K'=A'/\mathfrak m'$）是其剩余域。则 $A'$ 为正则环，当且仅当
\begin{equation}
\label{0.22.5.1.1-fr}
  \operatorname{rg}_K(\mathfrak m/\mathfrak m^2)
  =\operatorname{rg}_{K'}(\mathfrak m'/\mathfrak m'^2).
  \tag{22.5.1.1}
\end{equation}
\end{proposition}

\begin{proof}
事实上，(22.5.1.1) 的左端是 $\dim(A)$（17.1.1）；又因
$A\subset A'$ 且 $A'$ 是有限 $A$-代数，有
$\dim(A')=\dim(A)$（16.1.5）。所以等式 (22.5.1.1) 正是 $A'$ 为正则环
的充要条件（17.1.1）。
\end{proof}

\begin{remark}[22.5.2]
\label{0.22.5.2-fr}
\begin{enumerate}
  \item[\rm(i)] 置 $A_1=A/\mathfrak m^2$，
  $A'_1=A'\otimes_AA_1=A'/\mathfrak m^2A'$。由于
  $\mathfrak m^2A'\subset\mathfrak m'^2$，
  $\mathfrak m'_1=\mathfrak m'/\mathfrak m^2A'$ 是 $A'_1$ 的极大理想，
  而 $\mathfrak m'_1/\mathfrak m_1'^2$ 作为 $K'$-向量空间同构于
  $\mathfrak m'/\mathfrak m'^2$。因此，$A'$ 的正则性条件只依赖于
  以 $A_1$ 为底的 $A'_1$-代数结构；这使我们能在下文应用 (22.4) 中关于
  Artin 环上有限代数的预备结果。
  \item[\rm(ii)] 由 (22.5.1) 的证明和 (16.2.6.1)，无论如何都有
  \begin{equation}
  \label{0.22.5.2.1-fr}
    \operatorname{rg}_K(\mathfrak m/\mathfrak m^2)
    \leqslant\operatorname{rg}_{K'}(\mathfrak m'/\mathfrak m'^2).
    \tag{22.5.2.1}
  \end{equation}

  \oldpage[0\textsubscript{IV-1}]{202}

  \item[\rm(iii)] 假设 $A$ 和 $A'$ 满足 (22.4.1) 的一般假设。由
  (19.8.10)，存在一个\emph{正则}局部环 $B$，其极大理想为
  $\mathfrak n$，使 $A$ 同构于 $B/\mathfrak n^2$。对每个
  $G_i\in B[T_i]$，取一个首一多项式，使它在 $A[T_i]$ 中的典范像为
  $F_i$ $(1\leqslant i\leqslant r)$。于是，若 $B'$ 是
  $B[T_1,\ldots,T_r]$ 对由诸 $G_i$ 生成的理想所作的商环，则显然
  $B'$ 是自由 $B$-模，秩为 $d$，并且
  $A'=B'\otimes_BA=B'/\mathfrak n^2B'$。因为假设 $A'$ 是局部环，
  $B'$ 也同样是局部环；若 $\mathfrak n'$ 是 $B'$ 的极大理想，则
  $B'/\mathfrak n'$ 同构于 $K'$，而
  $\mathfrak n'/\mathfrak n'^2$ 同构于 $\mathfrak m'/\mathfrak m'^2$，
  这里所指的是 $K'$-向量空间同构。因此，在 (22.4.1) 的假设下，
  不等式 (22.5.2.1) 表明
  \begin{equation}
  \label{0.22.5.2.2-fr}
    \operatorname{rg}_K(V)\leqslant\operatorname{rg}_{K'}(\mathfrak m'/\mathfrak m'^2).
    \tag{22.5.2.2}
  \end{equation}
\end{enumerate}
\end{remark}

\begin{corollary}[22.5.3]
\label{0.22.5.3-fr}
设 $A$ 是正则局部环，$\mathfrak m$ 是其极大理想，
$T_i$ $(1\leqslant i\leqslant r)$ 是未定元；对每个 $i$，设
\begin{equation}
\label{0.22.5.3.1-fr}
  F_i(T_i)=T_i^{d_i}+\sum_{1\leqslant k\leqslant d_i}c_{ik}T_i^{d_i-k}
  \tag{22.5.3.1}
\end{equation}
是 $A[T_i]$ 中的首一多项式，并且 $d_i\geqslant2$ 对每个 $i$
都成立，且 $c_{ik}\in\mathfrak m$。设 $A'$ 是多项式环
$A[T_1,\ldots,T_r]$ 对由这 $r$ 个多项式 $F_i$ 生成的理想所作的商环。
则：
\begin{enumerate}
  \item[\rm(i)] $A'$ 是局部环，而它的剩余域 $K'$ 同构于剩余域
  $K$，即 $A$ 的剩余域。
  \item[\rm(ii)] $A'$ 为正则环，当且仅当诸类 $\xi_i$ 模
  $\mathfrak m^2$（由元素 $c_{i,d_i}$ 给出）
  $(1\leqslant i\leqslant r)$ 在 $K$ 上线性无关。
\end{enumerate}
\end{corollary}

\begin{proof}
事实上，采用 (22.5.2, (i)) 的记号，并以 $\overline F_i$ 表示
$A_1[T_i]$ 中的多项式，该多项式由把 $F_i$ 的系数经同态
$A\longrightarrow A_1$ 映射而得；则 $A'_1$ 可由 $A_1$ 和诸多项式
$\overline F_i$ 按照 (22.4.1) 的方式定义。因此结论由 (22.5.1) 和
(22.4.3) 得出。
\end{proof}

\begin{theorem}[22.5.4]
\label{0.22.5.4-fr}
设 $A$ 是正则局部环，$K$ 是其剩余域，并假设 $K$ 的特征为
$p>0$；设
\begin{equation}
\label{0.22.5.4.1-fr}
  F_i(T_i)=T_i^p+p\sum_{k=1}^{p-1}c_{ik}T_i^{p-k}-f_i
  \tag{22.5.4.1}
\end{equation}
其中 $c_{ik}\in A$ 且 $f_i\in A$ $(1\leqslant i\leqslant r)$。
设 $B$ 是 $A[T_1,\ldots,T_r]$ 对由这 $r$ 个多项式 $F_i$ 生成的理想
所作的商环。则 $B$ 是局部环，并且下列条件等价：
\begin{enumerate}
  \item[a)] $B$ 是正则环；
  \item[b)] 诸元素 $d_Af_i\otimes1_K$ $(1\leqslant i\leqslant r)$
  在 $K$-向量空间 $\Omega^1_A\otimes_AK$ 中线性无关；
  \item[c)] 若 $B'=B/\mathfrak m^2B$，并置
  $\Xi_{B_{(K)}/K}=\Xi_{B'_{(K)}/K}$（22.4.5），则特征同态
  $\Xi_{B_{(K)}/K}\longrightarrow\mathfrak m/\mathfrak m^2$
  （22.4.6.2）是单射。
\end{enumerate}
\end{theorem}

\begin{proof}
事实上，$B'$ 仍可由 $A_1=A/\mathfrak m^2$ 和诸多项式
$\overline F_i$ 来定义；这些多项式由把诸 $F_i$ 的系数经同态
$A\longrightarrow A_1$ 映射而得。注意到 $\Omega^1_A\otimes_AK$ 典范同构于
$\Omega^1_{A_1}\otimes_{A_1}K$（20.5.12, (ii)），本定理由 (22.5.1)
与 (22.4.7.4) 得出；当 $A_1$ 的特征不为 $p$ 时，使用注 (22.4.8)。
\end{proof}

\oldpage[0\textsubscript{IV-1}]{203}

\begin{env}[22.5.5]
\label{0.22.5.5-fr}
设 $k$ 是特征为 $p>0$ 的域，$A$ 是局部环，$k'$ 是 $k$ 的一个
\emph{有限纯不可分扩张}，满足 $k'^p\subset k$；设 $\mathfrak m$ 是
$A$ 的极大理想，$K=A/\mathfrak m$ 是其剩余域。回忆
$A'=A\otimes_kk'$ 是局部环，其剩余域与局部环 $K\otimes_kk'$ 的
剩余域相同（Bourbaki, \emph{Alg. comm.}, chap.~V, §~2,
n\textsuperscript{o}~3, lemme~4）。注意
$A'_{(K)}=A'\otimes_AK=K\otimes_kk'$，从而由 (21.3.6)，在相差一个
典范同构的意义下，
\begin{equation}
\label{0.22.5.5.1-fr}
  \Theta_{A'_{(K)}/K}=\Theta_{k'/k}\otimes_kK
  \tag{22.5.5.1}
\end{equation}
。此外，由 (21.4.7)，有 $\Xi_{k'/k}=0$；换言之，典范同态
(21.3.2.2)
\[
  \pi_{k'/k}:\Theta_{k'/k}\longrightarrow\Omega^1_k
\]
是单射，因而有一个\emph{典范单射}
\begin{equation}
\label{0.22.5.5.2-fr}
  \Theta_{A'_{(K)}/K}\longrightarrow\Omega^1_k\otimes_kK
  \tag{22.5.5.2}
\end{equation}
；考虑到 (22.5.5.1)，它可具体描述如下：对每个 $x'\in k'$，把元素
$d_{A'_{(K)}/K}(x'\otimes1_K)\otimes1_K$（属于
$\Omega^1_{A'_{(K)}/K}\otimes_{A'_{(K)}}K$）映到元素
$d_k(x'^p)\otimes1_K$（属于 $\Omega^1_k\otimes_kK$）。
\end{env}

\begin{lemma}[22.5.6]
\label{0.22.5.6-fr}
采用 (22.5.5) 的记号：
\begin{enumerate}
  \item[\rm(i)] 通过 (22.5.5.2)，典范同态 $\pi_{A'_{(K)}/K}$
  （21.3.2.2）的核可等同于
  $(\Theta_{k'/k}\otimes_kK)\cap\Upsilon_{K/k}$。
  \item[\rm(ii)] 设 $A_1=A/\mathfrak m^2$，
  $A'_1=A_1\otimes_kk'=A'/\mathfrak m^2A'$；则特征同态
  $\chi'_{A'_1/A_1}$（22.4.6.2）可等同于特征同态
  $\chi_{A/k}:\Upsilon_{K/k}\longrightarrow V=\mathfrak m/\mathfrak m^2$
  （20.6.24）在
  $(\Theta_{k'/k}\otimes_kK)\cap\Upsilon_{K/k}$ 上的限制。
\end{enumerate}
\end{lemma}

\begin{proof}
(i) 在 (20.6.9) 中把 $A$ 换成素域，把 $B$ 换成 $k$，把 $C$ 换成
$K$，并把 $E$ 换成 $A_1=A/\mathfrak m^2$（它是 $K$ 由
$V=\mathfrak m/\mathfrak m^2$ 所作的扩张）。由该结果并参照 (22.5.5)，
只须验证：若 $x'\in k'$，而 $z\in A_1$ 的模
$\mathfrak m/\mathfrak m^2$ 类为 $x'^p\in k$，则
$d_k(x'^p)\otimes1_K$ 经典范同态
$\Omega^1_k\otimes_kK\longrightarrow\Omega^1_{A_1}\otimes_{A_1}K$
的像为 $d_{A_1}z\otimes1_K$；这由定义立即得出。

(ii) 因为 $A'_{(K)}=K\otimes_kk'=(A'_1)_{(K)}$（由关系
$K=A_1/(\mathfrak m/\mathfrak m^2)$），
$\Theta_{A'_{(K)}/K}$ 可等同于
$\Theta_{(A'_1)_{(K)}/K}$；验证由与 (i) 相同的计算以及
$\pi'_{(A'_1)_{(K)}/K}$ 的定义 (22.4.5.2) 得出。
\end{proof}

\begin{proposition}[22.5.7]
\label{0.22.5.7-fr}
设 $k$ 是特征为 $p>0$ 的域，$A$ 是作为正则局部环的 $k$-代数，
$\mathfrak m$ 是其极大理想，$K=A/\mathfrak m$ 是其剩余域；把
$V=\mathfrak m/\mathfrak m^2$ 看成 $K$-向量空间。设 $k'$ 是 $k$ 的
有限扩张，满足 $k'^p\subset k$，并置 $A'=A\otimes_kk'$。则局部环
$A'$ 为正则环，当且仅当特征同态
$\chi_{A/k}:\Upsilon_{K/k}\longrightarrow V$（20.6.24）在
$(\Theta_{k'/k}\otimes_kK)\cap\Upsilon_{K/k}$ 上的限制是单射。
\end{proposition}

\begin{proof}
需要证明命题中的条件等价于 (22.5.1.1)；这里以 $\mathfrak m'$ 表示
$A'$ 的极大理想，以 $K'=A'/\mathfrak m'$ 表示其剩余域。采用注
(22.5.2, (i)) 的记号，有 $A'_1=A_1\otimes_kk'$。现设
$(a_i)_{1\leqslant i\leqslant r}$ 是一个 $p$-基底（在域 $k'$ 中、
相对于 $k$），从而 $a_i^p=b_i\in k$；并且 $k'$ 同构于多项式环

\oldpage[0\textsubscript{IV-1}]{204}

$k[T_1,\ldots,T_r]$ 对由多项式 $F_i=T_i^p-b_i$
$(1\leqslant i\leqslant r)$ 生成的理想所作的商环。由此推出
$A'_1$ 同构于多项式环 $A_1[T_1,\ldots,T_r]$ 对由诸 $F_i$ 生成的
理想所作的商环。于是可以应用判别法 (22.4.7)，而命题中的条件便等价于
$\chi'_{A'_1/A_1}$ 为单射；最后用 (22.5.6) 即得结论。
\end{proof}

现在可以证明 (19.6.5) 的逆命题：

\begin{theorem}[22.5.8]
\label{0.22.5.8-fr}
设 $k$ 是域，$p$ 是它的特征指数，$A$ 是 Noether 局部 $k$-代数。
下列条件等价：
\begin{enumerate}
  \item[a)] $A$ 是形式平滑 $k$-代数（关于其预进拓扑）。
  \item[b)] $A$ 在 $k$ 上几何正则。
  \item[b')] 对每个有限扩张 $k'$（基域为 $k$），只要
  $k'^p\subset k$，$A'=A\otimes_kk'$ 就是正则环。
\end{enumerate}
\end{theorem}

\begin{proof}
考虑到 (19.6.6)，只须处理 $p>1$ 的情形。

a) 蕴含 b) 正是 (19.6.5)，而 b) 显然蕴含 b')。我们证明 b') 蕴含
(22.2.2) 的条件成立。第一个条件是显然的（17.1.1），因为 b') 首先蕴含
$A$ 是正则环。另一方面，设 $(x_\alpha)_{\alpha\in I}$ 是一个
$p$-基底（在域 $k^{1/p}$ 中、相对于 $k$）。由 (21.4.5)，诸元素
$d_k(x_\alpha^p)$ 构成 $k$-向量空间 $\Omega^1_k$ 的一个基，因而诸元素
$d_k(x_\alpha^p)\otimes1_K$ 构成 $K$-向量空间
$\Omega^1_k\otimes_kK$ 的一个基。于是（参见 (22.5.5)），当 $k'$ 遍历
$k^{1/p}$ 的所有在 $k$ 上有限的子扩张时，诸子空间
$\Theta_{k'/k}\otimes_kK$（位于 $\Omega^1_k\otimes_kK$ 中）
构成滤升族，其并为
$\Omega^1_k\otimes_kK$。由 (22.5.7)，条件 b') 遂蕴含
$\chi_{A/k}$ 为单射，而这正是 (22.2.2) 的条件 (ii)。
\end{proof}

\begin{corollary}[22.5.9]
\label{0.22.5.9-fr}
设 $k$ 是域，$A$ 是形式平滑的 Noether 局部 $k$-代数。则对每个素理想
$\mathfrak p$（属于 $A$），局部环 $A_{\mathfrak p}$ 都是形式平滑
$k$-代数（关于 $\mathfrak p$-预进拓扑）。
\end{corollary}

\begin{proof}
事实上，采用 (22.5.8, b')) 的记号，只须证明局部环
$A_{\mathfrak p}\otimes_kk'$ 是正则环；但它是局部环
$A\otimes_kk'$ 的分式环，而后者按假设是正则环，所以
$A_{\mathfrak p}\otimes_kk'$ 也是正则环（17.3.6）。
\end{proof}

\begin{remark}[22.5.10]
\label{0.22.5.10-fr}
\begin{enumerate}
  \item[\rm(i)] 若在条件 b') 中只考虑\emph{单生成}扩张
  $k'=k(x)$（基域为 $k$，且 $x^p\in k$），则 (22.5.8) 的结论不再成立。事实上，
  对所有这些扩张 $k'$，可能有
  $\Upsilon_{K/k}\cap(\Theta_{k'/k}\otimes_kK)=0$，尽管
  $\Upsilon_{K/k}\neq0$。这意味着对每个 $x\in k^{1/p}$，只要
  $x\notin k$，就必须有
  $d_k(x^p)\otimes1_K\notin\Upsilon_{K/k}$，等价地有
  $d_K(x^p)\neq0$，也就是 $x^p\notin K^p$；换言之，必须有
  $K^p\cap k=k^p$，尽管 $K$ 是 $k$ 的一个\emph{不可分}扩张。
  这种扩张很容易构造：从一个完全域 $k_0$ 出发，其特征为 $p>0$；取
  $X$、$Y$、$Z$ 三个未定元，并置 $k=k_0(X,Y)$ 以及
  $K=k(X^{1/p}+ZY^{1/p})$。容易验证 $K$ 满足上述条件，因而
  $\Upsilon_{K/k}\neq0$。现在应用引理 (22.2.5.2)，取 $V=K$ 且令
  $h$ 为零同态；此时 $A$ 是剩余域为 $K$ 的离散赋值环，并且
  $\chi_{A/k}=0$，所以由 (22.2.2)，它不是形式平滑 $k$-代数；然而对
  $A\otimes_kk'$ 都是正则环，其中 $k'\subset k^{1/p}$ 为
  $k$ 的每个\emph{单生成}扩张。

  \oldpage[0\textsubscript{IV-1}]{205}

  \item[\rm(ii)] 推论 (22.5.9) 引出如下问题：若 $A$ 是 Noether
  $k$-代数，则由满足下述性质的素理想
  $\mathfrak p\in\operatorname{Spec}(A)$ 所构成的集合在什么条件下是
  \emph{开集}：$A_{\mathfrak p}$ 是形式平滑 $k$-代数？我们以后将处理
  其中若干特殊情形。
\end{enumerate}
\end{remark}

\subsection{Zariski 的 Jacobian 判别法}
\label{subsection:0.22.6-fr}

\begin{theorem}[22.6.1]
\label{0.22.6.1-fr}
\textnormal{（形式平滑性的 Jacobian 判别法）。}
设 $A$、$B$ 是两个拓扑环，
$u:A\longrightarrow B$ 是连续同态，并使 $B$ 成为形式平滑
$A$-代数；设 $\mathfrak J$ 是 $B$ 的一个理想（不必闭），
$C=B/\mathfrak J$ 是商拓扑 $A$-代数。则 $C$ 为形式平滑 $A$-代数，
当且仅当典范同态（参见 (20.5.11.2)）
\begin{equation}
\label{0.22.6.1.1-fr}
  \delta_{C/B/A}:\mathfrak J/\mathfrak J^2
  \longrightarrow\Omega^1_{B/A}\otimes_BC
  \tag{22.6.1.1}
\end{equation}
是\emph{形式左可逆}的（参见 (19.1.5)）。
\end{theorem}

\begin{proof}
事实上，说 $B$（相应地 $C$）是形式平滑 $A$-代数，按照 (19.4.4)，
就是有
$\operatorname{Exalcotop}_A(B,L)=0$（相应地
$\operatorname{Exalcotop}_A(C,L)=0$）；这里所指的是离散
$B$-模（相应地 $C$-模）$L$，并被 $B$（相应地 $C$）的某个开理想零化。
按假设，$\operatorname{Exalcotop}_A(B,L)=0$ 对每个离散 $C$-模 $L$
都成立，只要它被 $C$ 的某个开理想零化。所以，说 $C$ 是形式平滑
$A$-代数，就是说典范同态
$\operatorname{Exalcotop}_A(C,L)\longrightarrow
\operatorname{Exalcotop}_A(B,L)$ 是单射。因此定理由 (20.7.8) 得出。
\end{proof}

回忆 (19.5.3)：当 $B$ 和 $C$ 都是形式平滑 $A$-代数时，
$\mathfrak J/\mathfrak J^2$ 是形式投射拓扑 $C$-模；此外，典范同态
$\varphi_n:\mathbf S_C^n(\mathfrak J/\mathfrak J^2)
\longrightarrow\mathfrak J^n/\mathfrak J^{n+1}$ 是\emph{形式双同态}，
这里假设 $B$ 是\emph{预可容许}环。
\begin{corollary}[22.6.2]
\label{0.22.6.2-fr}
假设和记号同 (22.6.1)。再假设在 $B$ 中每个开理想的平方仍是开理想，
并且在 $\mathfrak J$ 上，由 $B$ 的拓扑诱导的拓扑与由该拓扑导出的
拓扑 (19.0) 相同（注意，当 $B$ 为 Noether 环且 $B$ 的拓扑为预进拓扑
（$0_{\mathrm I}$, 7.3.2），或者 $B$ 的拓扑是 $\mathfrak J$-预进拓扑时，
这两个条件都成立）。设 $(\mathfrak b_\lambda)$ 是 $B$ 的开理想基本系，
并置 $B_\lambda=B/\mathfrak b_\lambda$，这对每个 $\lambda$。则：
\begin{enumerate}
  \item[\rm(i)] 下列条件等价：
  \begin{enumerate}
    \item[a)] $C$ 是形式平滑 $A$-代数。
    \item[b)] 对每个 $\lambda$，下列 $B_\lambda$-模同态（由
    $\delta_{C/B/A}$ 经张量积所得）
    \begin{equation}
    \label{0.22.6.2.1-fr}
      (\mathfrak J/\mathfrak J^2)\otimes_BB_\lambda
      \longrightarrow\Omega^1_{B/A}\otimes_BB_\lambda
      \tag{22.6.2.1}
    \end{equation}
    是左可逆的（换言之，它同构到
    $\Omega^1_{B/A}\otimes_BB_\lambda$ 的一个直因子上）。
  \end{enumerate}
  \item[\rm(ii)] 再假设拓扑环 $C$ 是预可容许的
  （$0_{\mathrm I}$, 7.1.2），并设 $\mathfrak L$ 是 $C$ 的一个定义理想；
  则 (i) 的条件 a) 和 b) 也等价于：

  \oldpage[0\textsubscript{IV-1}]{206}

  \begin{enumerate}
    \item[c)] 下列 $(C/\mathfrak L)$-模同态
    \[
      (\mathfrak J/\mathfrak J^2)\otimes_C(C/\mathfrak L)
      \longrightarrow\Omega^1_{B/A}\otimes_B(C/\mathfrak L)
    \]
    是左可逆的。
  \end{enumerate}
\end{enumerate}
\end{corollary}

\begin{proof}
这些假设蕴含 $\Omega^1_{B/A}$ 上的拓扑由 $B$ 的拓扑导出（20.4.5），
所以 (i) 由 (22.6.1) 和 (19.1.7) 得出。另一方面，回忆 (20.4.9)：
由于 $B$ 是形式平滑 $A$-代数，$\Omega^1_{B/A}$ 是形式投射
$B$-模，因而 $\Omega^1_{B/A}\otimes_BB_\lambda$ 是投射
$B_\lambda$-模；这对每个 $\lambda$ 都成立（19.2.4）。
当 $C$ 为预可容许环时，c) 与 a)、b) 的等价性遂由 (19.1.9) 得出。
\end{proof}

特别地：

\begin{corollary}[22.6.3]
\label{0.22.6.3-fr}
设 $A$ 是环，$B$ 是 $A$-代数，$\mathfrak J$ 是 $B$ 的一个理想，并置
$C=B/\mathfrak J$；假设 $A$、$B$、$C$ 都取离散拓扑，并且 $B$ 是形式平滑
$A$-代数。则 $C$ 为形式平滑 $A$-代数，当且仅当典范同态 (22.6.1.1)
是左可逆的。
\end{corollary}

注意，每个 $A$-代数 $C$ 都是某个多项式代数
$A[T_\alpha]_{\alpha\in I}$ 的商，而后者是形式平滑的（19.3.3）；
所以判别法 (22.6.3) 原则上可以判定 $C$ 是否为形式平滑 $A$-代数。

\begin{proposition}[22.6.4]
\label{0.22.6.4-fr}
设 $A$ 是环，$B$ 是形式平滑 $A$-代数（各取离散拓扑），
$\mathfrak J$ 是 $B$ 的一个理想，并置 $C=B/\mathfrak J$；假设
$\mathfrak J/\mathfrak J^2$ 是有限型 $C$-模。设 $\mathfrak p$ 是 $C$ 的
一个素理想，$k(\mathfrak p)$ 是 $C_{\mathfrak p}$ 的剩余域，
$\mathfrak p'$ 是 $\mathfrak p$ 在 $B$ 中的逆像，而 $\mathfrak q$ 是
$A$ 的素理想，即 $\mathfrak p'$ 的逆像。下列条件等价：
\begin{enumerate}
  \item[a)] $C_{\mathfrak p}$ 是形式平滑
  $A_{\mathfrak q}$-代数（各取离散拓扑）。
  \item[a')] $C_{\mathfrak p}$ 是形式平滑 $A$-代数（各取离散拓扑）。
  \item[b)] 典范同态
  \begin{equation}
  \label{0.22.6.4.1-fr}
    (\mathfrak J/\mathfrak J^2)\otimes_Ck(\mathfrak p)
    \longrightarrow\Omega^1_{B/A}\otimes_Bk(\mathfrak p)
    \tag{22.6.4.1}
  \end{equation}
  是单射。
  \item[c)] 存在 $f\in C-\mathfrak p$，使 $C_f$ 为形式平滑 $A$-代数。
\end{enumerate}
若进一步假设 $B$ 为 Noether 环，则上述条件还等价于：
\begin{enumerate}
  \item[d)] $C_{\mathfrak p}$ 是形式平滑 $A_{\mathfrak q}$-代数；其中
  $\mathfrak p$-预进拓扑置于 $C_{\mathfrak p}$，而离散拓扑（或
  $\mathfrak q$-预进拓扑）置于 $A_{\mathfrak q}$。
\end{enumerate}
\end{proposition}

\begin{proof}
因为 $A_{\mathfrak q}$ 是形式平滑 $A$-代数（19.3.5, (iv) 和 (ii)），
已知 a) 与 a') 等价（19.3.5, (ii)）。有
$C_{\mathfrak p}=B_{\mathfrak p'}/\mathfrak J B_{\mathfrak p'}$，并且
$B_{\mathfrak p'}$ 是关于离散拓扑的形式平滑 $A$-代数
（19.3.5, (iv)）；此外
$\Omega^1_{B_{\mathfrak p'}/A}=(\Omega^1_{B/A})_{\mathfrak p'}$
（20.5.9）。把 (22.6.3) 应用于 $B_{\mathfrak p'}$ 和
$C_{\mathfrak p}$，并注意 $\Omega^1_{B/A}$ 是投射 $B$-模
（20.4.9）、$\mathfrak J/\mathfrak J^2$ 是有限型 $B$-模，再结合
(19.2.11) 和 (19.1.12)，便得 a') 与 b) 等价。应用 (19.1.12) 也证明
b) 与 c) 等价。最后注意，因为 $B_{\mathfrak p'}$ 对离散拓扑是形式平滑
$A_{\mathfrak q}$-代数（19.3.5, (iv)），给 $B_{\mathfrak p'}$ 赋予
$\mathfrak p'$-预进拓扑后，它仍为形式平滑 $A_{\mathfrak q}$-代数
（19.3.8）；这里在 $A_{\mathfrak q}$ 上采用离散拓扑或

\oldpage[0\textsubscript{IV-1}]{207}

$\mathfrak q$-预进拓扑。当 $B$ 为 Noether 环时，为证明 b) 与 d) 等价，把 (22.6.2)
应用于环 $A_{\mathfrak q}$（取离散拓扑或 $\mathfrak q$-预进拓扑），
以及环 $B_{\mathfrak p'}$（取 $\mathfrak p'$-预进拓扑）。由于
$C_{\mathfrak p}$ 是预可容许的，并且
$\mathfrak pC_{\mathfrak p}$ 是其拓扑的定义理想，可以援用 (22.6.2) 中
c) 与 a) 的等价性。
\end{proof}

\begin{corollary}[22.6.5]
\label{0.22.6.5-fr}
在 (22.6.4) 的一般假设下，由所有
$\mathfrak p\in\operatorname{Spec}(C)$ 构成的下列集合是开集：那些使
$C_{\mathfrak p}$ 对离散拓扑为形式平滑 $A$-代数（或为形式平滑
$A_{\mathfrak q}$-代数，其中 $\mathfrak q$ 表示 $\mathfrak p$ 在
$A$ 中的逆像）的素理想。该开集位于 $\operatorname{Spec}(C)$ 中。
\end{corollary}

这由 (22.6.4) 的形式 c) 得出。注意，当 $B$ 为 Noether 环时，可以用预进
拓扑代替离散拓扑。

\begin{corollary}[22.6.6]
\label{0.22.6.6-fr}
在 (22.6.4) 的一般假设下，下列条件等价：
\begin{enumerate}
  \item[a)] $C$ 是形式平滑 $A$-代数（各取离散拓扑）。
  \item[b)] 对每个 $\mathfrak p\in\operatorname{Spec}(C)$（或者只对
  每个极大理想 $\mathfrak p$，其属于 $C$），$C_{\mathfrak p}$ 是形式平滑
  $A$-代数（或形式平滑 $A_{\mathfrak q}$-代数，其中 $\mathfrak q$ 表示
  $\mathfrak p$ 在 $A$ 中的逆像），各取离散拓扑。
\end{enumerate}
此外，当 $B$ 为 Noether 环时，可以在 b) 中把离散拓扑换成
$A_{\mathfrak q}$ 和 $C_{\mathfrak p}$ 上的预进拓扑。
\end{corollary}

\begin{proof}
最后一个断言由 (22.6.4) 得出。由 (22.6.3)，条件 a) 等价于同态
(22.6.1.1) 左可逆；注意 $\Omega^1_{B/A}$ 是投射 $B$-模，而
$\mathfrak J/\mathfrak J^2$ 是有限型 $C$-模，故 a) 与 b) 的等价性由
(19.1.14) 得出。
\end{proof}

\begin{proposition}[22.6.7]
\label{0.22.6.7-fr}
设 $k_0$ 是域，$k$ 是 $k_0$ 的可分扩张，$C$ 是有限型 $k$-代数。
\begin{enumerate}
  \item[\rm(i)] 设 $\mathfrak p$ 是 $C$ 的一个素理想。下列条件等价：
  \begin{enumerate}
    \item[a)] $C_{\mathfrak p}$ 对离散拓扑是形式平滑
    $k_0$-代数。
    \item[b)] $C_{\mathfrak p}$ 是形式平滑 $k_0$-代数（关于
    $\mathfrak p$-预进拓扑）。
    \item[c)] 存在 $f\in C-\mathfrak p$，使 $C_f$ 对离散拓扑是形式平滑
    $k_0$-代数。
    \item[d)] $C_{\mathfrak p}$ 在 $k_0$ 上是几何正则环（19.6.5）。
  \end{enumerate}
  此外，具有这些性质中任一性质的
  $\mathfrak p\in\operatorname{Spec}(C)$ 所成的集合，在
  $\operatorname{Spec}(C)$ 中是开集。
  \item[\rm(ii)] 下列条件等价：
  \begin{enumerate}
    \item[a)] $C$ 对离散拓扑是形式平滑 $k_0$-代数。
    \item[b)] 每个 $\mathfrak p\in\operatorname{Spec}(C)$ 都满足 (i) 的
    等价条件。
    \item[c)] 每个极大理想 $\mathfrak m$（属于 $C$）都满足 (i) 的等价条件。
    \item[d)] 对每个扩张 $k'$（底域为 $k_0$），
    $C\otimes_{k_0}k'$ 都是正则环（17.3.6）；也称 $C$ 在 $k_0$ 上
    几何正则。
  \end{enumerate}

  \oldpage[0\textsubscript{IV-1}]{208}

  \item[\rm(iii)] 设 $B=k[T_1,\ldots,T_r]$ 是多项式代数，
  $\mathfrak J$ 是 $B$ 的一个理想，使 $C$ 同构于
  $B/\mathfrak J$。设 $\mathfrak q$ 是 $B$ 的一个包含
  $\mathfrak J$ 的素理想，并置 $\mathfrak p=\mathfrak q/\mathfrak J$。
  此时 (i) 的条件还等价于下列条件（“Zariski 的 Jacobian 判别法”）：
  \begin{enumerate}
    \item[e)] 存在一个元素系 $(f_i)_{1\leq i\leq m}$，它属于
    $\mathfrak J$，并生成理想 $\mathfrak J_{\mathfrak q}$（位于
    $B_{\mathfrak q}$ 中）；并且存在 $k_0$-导射 $D_j$，它们从 $B$
    到自身，且满足 $(1\leq j\leq m)$，使
    $\det(D_j(f_i))\notin\mathfrak q$。
  \end{enumerate}
\end{enumerate}
\end{proposition}

\begin{proof}
已知 $C$ 总同构于一个 $k$-代数，它形如 $B/\mathfrak J$。由于 $B$ 是形式平滑
$k$-代数（19.3.3），而 $k$ 在 $k_0$ 上可分，故是形式平滑
$k_0$-代数（19.6.1）；从而 $B$ 是形式平滑 $k_0$-代数
（19.3.5, (ii)）。因此，(i) 中 a)、b)、c) 的等价性由 (22.6.4) 得出，
(ii) 中 a)、b)、c) 的等价性由 (22.6.6) 得出；(i) 中 a) 与 d) 的
等价性由 (22.5.8) 得出。由于 $C\otimes_{k_0}k'$ 的每个局部化，也都是
$C_{\mathfrak p}\otimes_{k_0}k'$ 的一个局部化，其中
$\mathfrak p\in\operatorname{Spec}(C)$ 适当，(ii) 中 d) 与 b) 的
等价性由 (i) 中 d) 与 a) 的等价性得出。最后，因为
$\Omega^1_{B/k_0}$ 是投射 $B$-模，Zariski 判别法 e) 与 (i) 中其他条件
的等价性，由 (19.1.12) 和 (22.6.3) 并结合 (20.4.8) 得出。
\end{proof}

Zariski 实际关心的是局部环 $C_{\mathfrak p}$ 的正则性微分判别法。对于
一个含域的 Noether 局部环，说它是正则环，等价于说它作为其素子域上的代数
形式平滑（19.6.4）；因此在 (22.6.7) 中把 $k_0$ 换成 $k$ 的素子域，
便立即得到这样的判别法。特别地，得到下面的结果；以后还会把它作为 Nagata
更一般结果的一个特殊情形重新得到（IV, 6.12.5）：

\begin{corollary}[22.6.8]
\label{0.22.6.8-fr}
\textnormal{（Zariski）。}
设 $C$ 是域上的有限型代数。则由
$\mathfrak p\in\operatorname{Spec}(C)$ 构成的下列集合是开集：那些使
$C_{\mathfrak p}$ 成为正则局部环的素理想。该开集位于
$\operatorname{Spec}(C)$ 中。
\end{corollary}

\begin{remark}[22.6.9]
\label{0.22.6.9-fr}
\begin{enumerate}
  \item[\rm(i)] 若不假设 $k$ 在 $k_0$ 上可分，而只假设它具有
  \emph{有限纯不可分重数}，则 (22.6.7) 的各断言仍然成立。换言之
  （19.6.6），存在有限纯不可分扩张 $k'_0$（底域为 $k_0$），使 Artin 局部环
  $k\otimes_{k_0}k'_0$ 的剩余域是可分扩张 $k'$（底域为 $k'_0$）。此时存在一个
  $k'_0$-单同态
  $k'\longrightarrow k\otimes_{k_0}k'_0$；它与典范同态
  $k\otimes_{k_0}k'_0\longrightarrow k'$ 的复合为恒等，因为 $k'$ 是形式平滑
  $k'_0$-代数（19.6.1）。因此
  $C\otimes_{k_0}k'_0$ 带有有限型 $k'$-代数结构；又因 $k'$ 在
  $k'_0$ 上可分，可以对这个 $k'$-代数应用 (22.6.7)，再用 (19.4.7)
  即得所述断言。我们以后不需要这一推广，故把细节留给读者。另一方面，
  我们不知道在对 $k$ 在 $k_0$ 上的扩张\emph{不作任何}假设时，
  (22.6.7) 的结果是否仍成立。
  \item[\rm(ii)] 在 (22.6.7) 的一般假设下，设 $(k_\alpha)$ 是 $k$ 的
  含 $k_0$ 的滤降子域族，并且
  $\bigcap_\alpha k_\alpha(k^p)=k_0(k^p)$（这里 $p$ 是 $k_0$ 的特征指数）。
  利用 Nagata 的维数计数方法，可以证明：在 Jacobian 判别法
  (22.6.7, (iii)) 中，只须考虑导射 $D_j$；它们属于 $B$，并且是
  $k_\alpha$-导射，其中 $\alpha$ 适当。

  \oldpage[0\textsubscript{IV-1}]{209}

  这个结果的意义在于，总存在这样的族 $(k_\alpha)$，使
  $[k:k_\alpha]$ 有限（21.8.6）。我们不证明 Zariski 判别法的这一精化，
  因为以后不会用到它。在 (22.7) 中，我们将对完备局部环给出 Zariski
  判别法的一个变体（同样由 Nagata 提出）；它基本上用相同的方法证明
  （技术困难略大）。
\end{enumerate}
\end{remark}

\subsection{Nagata 的 Jacobian 判别法}
\label{subsection:0.22.7-fr}

\begin{env}[22.7.1]
\label{0.22.7.1-fr}
Nagata 的 Jacobian 判别法是 Zariski 的 Jacobian 判别法的类似物，不过它适用
于域上的\emph{形式幂级数}环的商环。我们将像 Nagata [31] 那样给出两个版本，
并在这里把它们表述为形式平滑性判别法。
\end{env}

\begin{proposition}[22.7.2]
\label{0.22.7.2-fr}
设 $k$ 是域，$A=k[[T_1,\ldots,T_r]]$ 是 $k$ 上的形式幂级数环，并赋予其
通常的 $k$-代数结构；设 $\mathfrak q$ 是 $A$ 的一个理想，
$B=A/\mathfrak q$。设 $\mathfrak p$ 是 $A$ 的一个包含
$\mathfrak q$ 的素理想。假设存在一个 $k$-子代数 $C'$，它位于
$C=A/\mathfrak p$ 中，并同构于形式幂级数代数
$k[[X_1,\ldots,X_s]]$，使得 $C$ 在 $C'$ 上有限，并且分式域
$L$（属于 $C$）是 $L'=k((X_1,\ldots,X_s))$ 的可分扩张（若 $k$ 的特征为 0，
这一假设总是成立）。则下列条件等价：
\begin{enumerate}
  \item[a)] $B_{\mathfrak p}=A_{\mathfrak p}/\mathfrak qA_{\mathfrak p}$
  是形式平滑 $k$-代数（对 $\mathfrak p$-预进拓扑）。
  \item[b)] 存在 $k$-导射 $D_i$ $(1\leq i\leq m)$，它们从 $A$
  到自身；以及元素 $f_i$ $(1\leq i\leq m)$，它们属于
  $\mathfrak q$，使诸 $f_i$ 在 $\mathfrak qA_{\mathfrak p}$
  中的像生成 $A_{\mathfrak p}$ 的这个理想，并且
  $\det(D_i(f_j))\notin\mathfrak p$。
  \item[c)] $B_{\mathfrak p}$ 是正则环。
\end{enumerate}
\end{proposition}

\begin{proof}
$B_{\mathfrak p}$ 的剩余域是 $k(\mathfrak p)$。因为
$k((X_1,\ldots,X_s))$ 在 $k$ 上可分（21.9.6.4），按假设
$k(\mathfrak p)$ 在 $k$ 上可分；已知在这些条件下，性质 a) 与 c) 等价
（19.6.4）。此外，$A_{\mathfrak p}$ 是正则环（(17.1.4) 和 (17.3.2)）；
又因其剩余域 $k(\mathfrak p)$ 在 $k$ 上可分，$A_{\mathfrak p}$ 在
$k$ 上形式平滑（对 $\mathfrak p$-预进拓扑）（19.6.4）。因此由
(22.6.2, (ii))，条件 a) 等价于下列 $L$-向量空间同态为单射：
\begin{equation}
\label{0.22.7.2.1-fr}
  j_0:(\mathfrak q/\mathfrak q^2)\otimes_AL
  \longrightarrow\Omega^1_{A_{\mathfrak p}/k}\otimes_{A_{\mathfrak p}}L
  =\widehat\Omega^1_{A/k}\otimes_AL
  \tag{22.7.2.1}
\end{equation}
。

另一方面，考虑复合同态
\begin{equation}
\label{0.22.7.2.2-fr}
  j:(\mathfrak q/\mathfrak q^2)\otimes_AL
  \xrightarrow{j_0}\Omega^1_{A/k}\otimes_AL
  \longrightarrow\widehat\Omega^1_{A/k}\otimes_AL
  \tag{22.7.2.2}
\end{equation}
，并证明条件 b) 等价于 $j$ 为单射。事实上，注意
\[
  (\mathfrak q/\mathfrak q^2)\otimes_AL
  =(\mathfrak qA_{\mathfrak p}/\mathfrak q^2A_{\mathfrak p})
  \otimes_{A_{\mathfrak p}}L~;
\]
条件 b) 意味着：若 $F_i=j(f'_i\otimes1)$，其中 $f'_i$ 是 $f_i$ 在
$\mathfrak qA_{\mathfrak p}/\mathfrak q^2A_{\mathfrak p}$ 中的像，则矩阵
$(\langle F_i,D_j\otimes1\rangle)$ 可逆。这蕴含诸 $F_i$ 线性无关，因而
\emph{更有}诸 $f'_i\otimes1$ 线性无关；但后者生成
$(\mathfrak qA_{\mathfrak p}/\mathfrak q^2A_{\mathfrak p})
\otimes_{A_{\mathfrak p}}L$，故 $j$ 为单射。反过来，假设 $j$ 为单射，
并取诸 $f_i$，使诸 $f'_i\otimes1$ 构成

\oldpage[0\textsubscript{IV-1}]{210}

$(\mathfrak qA_{\mathfrak p}/\mathfrak q^2A_{\mathfrak p})
\otimes_{A_{\mathfrak p}}L$ 的一个基。于是诸 $F_i$ 构成 $j$ 的像的一个基。
但由 (21.9.3)，$\widehat\Omega^1_{A/k}$ 是自由 $A$-模，秩为 $r$；而
$k$-导射（从 $A$ 到自身）生成它的对偶。于是，诸 $F_i$ 线性无关蕴含存在
$k$-导射 $D_j$，它们从 $A$ 到自身，使矩阵
$(\langle F_i,D_j\otimes1\rangle)$ 可逆；这正是条件 b)。

这些说明已经表明 b) 蕴含 a)。反过来，为证明条件 a) 蕴含 $j$ 为单射，
考虑交换图
\begin{equation}
\label{0.22.7.2.3-fr}
\xymatrix{
  (\mathfrak q/\mathfrak q^2)\otimes_AL
    \ar[r]^j \ar[d]_\alpha &
  \widehat\Omega^1_{A/k}\otimes_AL \ar@{=}[d] \\
  (\mathfrak p/\mathfrak p^2)\otimes_AL
    \ar[r]_i &
  \widehat\Omega^1_{A/k}\otimes_AL
}
\tag{22.7.2.3}
\end{equation}
，并注意下列引理：
\end{proof}

\begin{lemma}[22.7.2.4]
\label{0.22.7.2.4-fr}
设 $R$ 是正则局部环，$\mathfrak m$ 是其极大理想，
$K=R/\mathfrak m$ 是其剩余域。若 $\mathfrak n$ 是 $R$ 的一个理想，
且 $R/\mathfrak n$ 为正则环，则典范同态
\begin{equation}
\label{0.22.7.2.5-fr}
  \alpha:(\mathfrak n/\mathfrak n^2)\otimes_RK
  \longrightarrow(\mathfrak m/\mathfrak m^2)\otimes_RK
  \tag{22.7.2.5}
\end{equation}
是单射。
\end{lemma}

\begin{proof}
事实上，由 (17.1.6)，$\mathfrak n$ 由一个序列
$(x_i)_{1\leq i\leq t}$ 生成，而该序列是一组正则参数系
$(x_i)_{1\leq i\leq n}$（属于 $R$）的一部分。诸 $x_i$ 的像
$(1\leq i\leq t)$ 构成
$(\mathfrak n/\mathfrak n^2)\otimes_RK$ 的一个基；它们在 $\alpha$ 下的
像，是 $(\mathfrak m/\mathfrak m^2)\otimes_RK$ 的一组基的一部分；这组基
由诸 $x_i$ 的像（$1\leq i\leq n$）构成。引理得证。
\end{proof}

\begin{proof}[命题~{\rm(22.7.2)} 证明的续篇]
由这个引理和图 (22.7.2.3)，考虑假设 a)，问题归结为证明 $i$ 是单射。
在正合列
\begin{equation}
\label{0.22.7.2.6-fr}
  \mathfrak p/\mathfrak p^2\xrightarrow{h}
  \widehat\Omega^1_{A/k}\widehat\otimes_A(A/\mathfrak p)
  \xrightarrow{g}\widehat\Omega^1_{(A/\mathfrak p)/k}\longrightarrow0
  \tag{22.7.2.6}
\end{equation}
中，由 (20.7.20)，$g$ 是满射，而 $h$ 的像在 $g$ 的核中对于
$\mathfrak m$-进拓扑是稠密的（$\mathfrak m$ 是 $A$ 的极大理想）。
因为 $\widehat\Omega^1_{A/k}$ 是有限型 $A$-模，
$(A/\mathfrak p)$-模
\[
  \widehat\Omega^1_{A/k}\otimes_A(A/\mathfrak p)
  =\widehat\Omega^1_{A/k}\widehat\otimes_A(A/\mathfrak p)
\]
的每个子模对 $\mathfrak m$-进拓扑都是闭的（$0_{\mathrm I}$, 7.3.5），
所以列 (22.7.2.6) 正合。与 $L$ 作张量，得到正合列
\[
  (\mathfrak p/\mathfrak p^2)\otimes_AL\xrightarrow{i}
  \widehat\Omega^1_{A/k}\otimes_AL\longrightarrow
  \widehat\Omega^1_{(A/\mathfrak p)/k}\otimes_{A/\mathfrak p}L
  \longrightarrow0.
\]

由对 $L$ 所作的可分性假设和 (21.9.5)，
$\widehat\Omega^1_{(A/\mathfrak p)/k}\otimes_{A/\mathfrak p}L$ 的秩为
$s$，底域为 $L$；又因 $\widehat\Omega^1_{A/k}\otimes_AL$ 的秩为
$r$，底域为 $L$（21.9.3），有 $\operatorname{rg}_L(\operatorname{Im}i)=r-s$。但是
$A/\mathfrak p$ 在一个同构于 $k[[T_1,\ldots,T_s]]$ 的子代数上有限，
所以 $\dim(A/\mathfrak p)=s$（(17.1.4) 和 (16.1.5)）；从而
\[
  r-s=\dim(A)-\dim(A/\mathfrak p)=\dim(A_{\mathfrak p})
\]

\oldpage[0\textsubscript{IV-1}]{211}

（由 (16.5.11)）。又因 $A_{\mathfrak p}$ 是正则环（17.3.2），
$\dim(A_{\mathfrak p})$ 等于以 $L$ 为底域时
$\mathfrak pA_{\mathfrak p}/\mathfrak p^2A_{\mathfrak p}$ 的秩
（17.1.1），因而等于以 $L$ 为底域时
$(\mathfrak p/\mathfrak p^2)\otimes_AL$ 的秩。这就完成了
$i$ 为单射的证明。
\end{proof}

\begin{theorem}[22.7.3]
\label{0.22.7.3-fr}
设 $k$ 是域，$A$ 是剩余域为 $K$ 的完备 Noether 局部环。假设：
\begin{enumerate}
  \item[$1^\circ$] $[k:k^p]<+\infty$（其中 $p$ 是 $k$ 的特征指数）；
  \item[$2^\circ$] $K$ 是一个可分扩张 $K_0$（在 $k$ 上）的有限扩张；
  \item[$3^\circ$] $A$ 带有形式平滑 $K_0$-代数结构（关于预进拓扑）。
\end{enumerate}
设 $\mathfrak q$ 是 $A$ 的一个理想，$B=A/\mathfrak q$；
$\mathfrak p$ 是 $A$ 的一个包含 $\mathfrak q$ 的素理想。下列条件等价：
\begin{enumerate}
  \item[a)] $B_{\mathfrak p}$ 是形式平滑 $k$-代数（关于
  $\mathfrak p$-预进拓扑）。
  \item[b)] 存在 $k$-导射 $D_i$，它们从 $A$ 到自身且满足
  $(1\leq i\leq m)$；以及元素 $f_i$ $(1\leq i\leq m)$，它们属于
  $\mathfrak q$，使诸 $f_i$ 在 $\mathfrak qA_{\mathfrak p}$ 中的像
  生成 $A_{\mathfrak p}$ 的这个理想，并且
  $\det(D_i(f_j))\notin\mathfrak p$。
  \item[b')] 存在一个子扩张 $k'$（属于 $K_0$），并含 $K_0^p$，使
  $[K_0:k']<+\infty$；存在 $k'$-导射 $D_i$，它们从 $A$ 到自身且满足
  $(1\leq i\leq m)$；以及元素 $f_i$，它们属于 $\mathfrak q$
  $(1\leq i\leq m)$，使诸 $f_i$ 在 $\mathfrak qA_{\mathfrak p}$ 中的像
  生成 $A_{\mathfrak p}$ 的这个理想，并且
  $\det(D_i(f_j))\notin\mathfrak p$。
\end{enumerate}
\end{theorem}

\begin{proof}
按 $p=1$ 或 $p>1$ 分为两种情形。

\emph{A)} $k$ \emph{的特征为} $0$。此时 $K$ 是 $K_0$ 的可分扩张，
由 (19.6.4)，$A$ 在 $K_0$ 上同构于形式幂级数环
$K[[T_1,\ldots,T_r]]$；该形式幂级数环取通常的 $K$-代数结构。把
(19.8.8, (ii)) 应用于
$A/\mathfrak p$，可见 (22.7.2) 的一般条件成立，只须把其中的 $k$ 换成
$K$。此外，由 (19.6.4)，说 $B_{\mathfrak p}$ 是形式平滑 $k$-代数，
或说它是形式平滑 $K$-代数，是一回事；两个条件都等价于
$B_{\mathfrak p}$ 为正则环。因此可以应用 (22.7.2) 的结论，并立即得到
a)、b) 与 b') 等价（在 b') 中取 $k'=K_0$）。

\emph{B)} $k$ \emph{的特征为} $p>0$。由于 $A$ 是形式平滑
$K_0$-代数，而 $K_0$ 在 $k$ 上可分，由 (19.3.5, (ii)) 和 (19.6.1)，
$A$ 是形式平滑 $k$-代数；再由 (22.5.9)，$A_{\mathfrak p}$ 也是形式平滑
$k$-代数（对 $\mathfrak p$-预进拓扑）。因此仍由
(22.6.2, (ii))，条件 a) 等价于 (22.7.2.1) 中定义的同态 $j_0$
（其中 $L=k(\mathfrak p)$）为单射。所以 (19.1.12, c)) 已经表明 b)
蕴含 a)；另一方面 b') 显然蕴含 b)，于是只须证明 $j_0$ 为单射这一假设
蕴含 b')。

先以 $k'$ 表示 $K_0$ 的任意子扩张，并满足
$[K_0:k']<+\infty$。注意 $\widehat\Omega^1_{A/k'}$ 是有限型自由
$A$-模。仿照 (22.7.2) 的论证，只须证明：适当选择 $k'$ 后，复合同态
\begin{equation}
\label{0.22.7.3.1-fr}
  j':(\mathfrak q/\mathfrak q^2)\otimes_AL
  \xrightarrow{j_0}\Omega^1_{A/k}\otimes_AL
  \longrightarrow\widehat\Omega^1_{A/k'}\otimes_AL
  \tag{22.7.3.1}
\end{equation}

\oldpage[0\textsubscript{IV-1}]{212}

是单射。再考虑交换图
\begin{equation}
\label{0.22.7.3.2-fr}
\xymatrix{
  (\mathfrak q/\mathfrak q^2)\otimes_AL
    \ar[r]^{j_0} \ar[d]_\alpha &
  \Omega^1_{A/k}\otimes_AL \ar[r] \ar@{=}[d] &
  \widehat\Omega^1_{A/k'}\otimes_AL \ar@{=}[d] \\
  (\mathfrak p/\mathfrak p^2)\otimes_AL
    \ar[r]_{i_0} &
  \Omega^1_{A/k}\otimes_AL \ar[r] &
  \widehat\Omega^1_{A/k'}\otimes_AL
}
\tag{22.7.3.2}
\end{equation}

因为 $j_0$ 是单射，有 $\operatorname{Ker}(i_0\circ\alpha)=0$。又因
$A_{\mathfrak p}$ 是正则环，而假设 a) 蕴含
$B_{\mathfrak p}=A_{\mathfrak p}/\mathfrak qA_{\mathfrak p}$ 也是正则环
（19.6.5），引理 (22.7.2.4) 表明 $\alpha$ 是单射，从而
$\alpha^{-1}(\operatorname{Ker} i_0)=0$。若 $i'$ 是 (22.7.3.2) 第二行
诸同态的复合，则同样有
$\operatorname{Ker}(j')=\alpha^{-1}(\operatorname{Ker}(i'))$；所以只须证明：
适当选择 $k'$ 后，有
\begin{equation}
\label{0.22.7.3.3-fr}
  \operatorname{Ker}(i')=\operatorname{Ker}(i).
  \tag{22.7.3.3}
\end{equation}

但与 (22.7.2) 中相同的论证给出正合列
\[
  (\mathfrak p/\mathfrak p^2)\otimes_AL\xrightarrow{i'}
  \widehat\Omega^1_{A/k'}\otimes_AL\longrightarrow
  \widehat\Omega^1_{(A/\mathfrak p)/k'}\otimes_{A/\mathfrak p}L
  \longrightarrow0.
\]
所作假设容许应用 (21.9.8)，所以存在一个子扩张 $k'$（属于 $K_0$），
并含 $K_0^p$，满足 $[K_0:k']<+\infty$，并且
\[
  \operatorname{rg}_L\bigl(\widehat\Omega^1_{(A/\mathfrak p)/k'}
  \otimes_{A/\mathfrak p}L\bigr)
  =\dim(A/\mathfrak p)+\operatorname{rg}_L\Upsilon_{L/k}
  +\operatorname{rg}_{K_0}\Omega^1_{K_0/k'}.
\]
另一方面，
\[
  \dim(A/\mathfrak p)=\dim(A)-\dim(A_{\mathfrak p})
\]
由 (16.5.11) 成立；又有
\[
  \dim(A_{\mathfrak p})
  =\operatorname{rg}_L\bigl((\mathfrak p/\mathfrak p^2)\otimes_AL\bigr)
\]
，因为 $A_{\mathfrak p}$ 是正则环（17.1.1）；最后，由 (21.9.2)，
\[
  \dim(A)+\operatorname{rg}_{K_0}\Omega^1_{K_0/k'}
  =\operatorname{rg}_L\bigl(\widehat\Omega^1_{A/k'}\otimes_AL\bigr).
\]
因此得到
\begin{equation}
\label{0.22.7.3.4-fr}
  \operatorname{rg}_L(\operatorname{Ker}(i'))=\operatorname{rg}_L\Upsilon_{L/k}.
  \tag{22.7.3.4}
\end{equation}

另一方面，$A_{\mathfrak p}$ 的剩余域在 $k$ 的素子域上形式平滑，因而有
正合列 (20.6.22.1)
\[
  \Upsilon_{L/k}\xrightarrow{\chi_{A_{\mathfrak p}/k}}
  (\mathfrak pA_{\mathfrak p})/(\mathfrak pA_{\mathfrak p})^2
  \xrightarrow{i}\Omega^1_{A_{\mathfrak p}/k}
  \otimes_{A_{\mathfrak p}}L,
\]
并且由于 $A_{\mathfrak p}$ 是形式平滑 $k$-代数（对
$\mathfrak p$-预进拓扑），由 (22.2.2)，$\chi_{A_{\mathfrak p}/k}$ 是单射。因此
$\operatorname{Ker}(i)$ 同构于 $\Upsilon_{L/k}$；结合 (22.7.3.4) 以及
$\operatorname{Ker}(i)\subset\operatorname{Ker}(i')$，便完成了
(22.7.3.3) 的证明，进而证明了定理。
\end{proof}

\oldpage[0\textsubscript{IV-1}]{213}

\begin{remark}[22.7.4]
\label{0.22.7.4-fr}
事实上，我们已经证明（其中使用了 (21.9.6)）：若再要求 $k'$ 是
一个递减有向族 $(k_\alpha)$ 中的域之一，其中该族的各个域都是
$K_0$ 的子域并包含 $K_0^p$，且 $[K_0:k_\alpha]<+\infty$ 对每个
$\alpha$ 都成立，并且
$\bigcap_\alpha k_\alpha=K_0^p$，则条件 b') 仍与 (22.7.3) 的其他条件等价。
\end{remark}

\begin{corollary}[22.7.5]
\label{0.22.7.5-fr}
在 (22.7.3) 的一般假设下，诸素理想
$\mathfrak n\in\operatorname{Spec}(B)$ 中，使 $B_{\mathfrak n}$ 成为
形式平滑 $k$-代数（关于 $\mathfrak n$-预进拓扑）者所成的集合，是
$\operatorname{Spec}(B)$ 中的开集。
\end{corollary}

\begin{proof}
采用 (22.7.3) 的记号，只须证明：若 $B_{\mathfrak p}$ 在 $k$ 上
形式平滑，则 $B_{\mathfrak p'}$ 也具有同一性质；这里
$\mathfrak p'\in V(\mathfrak q)$ 可取为充分接近 $\mathfrak p$ 的点，
而邻近性在 $\operatorname{Spec}(A)$ 中理解。事实上，若 $k$-导射 $D_i$
与元素 $f_i$（属于 $\mathfrak q$）满足 (22.7.3) 的判据 b)，则也有
$f=\det(D_if_j)\notin\mathfrak p'$，其中 $\mathfrak p'\in D(f)$；另一方面，
存在 $g\notin\mathfrak p$，使诸 $f_i$ 在
$\mathfrak qA_{\mathfrak p'}$ 中的像生成 $\mathfrak qA_{\mathfrak p'}$，
其中 $\mathfrak p'\in D(g)$（Bourbaki，
\emph{Alg. comm.}，第 II 章，§~5，第~1 号，命题~2）。这就完成了
推论的证明。
\end{proof}

\begin{corollary}[22.7.6]
\label{0.22.7.6-fr}
\textnormal{（Nagata）}设 $B$ 是含有一个域的完备 Noether 局部环。
则诸 $\mathfrak n\in\operatorname{Spec}(B)$ 中使
$B_{\mathfrak n}$ 为正则环者所成的集合，是
$\operatorname{Spec}(B)$ 中的开集。
\end{corollary}

\begin{proof}
事实上，$B$ 是某个素域 $k$ 上的代数，所以由 (19.8.8, (ii))，它具有
$A/\mathfrak q$ 的形式，其中
$A=K_0[[T_1,\ldots,T_n]]$ 是一个形式幂级数环，而 $K_0$ 是 $k$ 的一个
扩张；另一方面，说 $B_{\mathfrak n}$ 是正则环，等价于说它是形式平滑
$k$-代数（关于 $\mathfrak n$-预进拓扑）（19.6.4）。因此，应用
(22.7.5) 的全部条件都得到满足。
\end{proof}

\begin{remarks}[22.7.7]
\label{0.22.7.7-fr}
\begin{enumerate}
  \item[(i)] 我们稍后将同 Nagata 一起，利用 (22.7.6) 看到：可以把
  (22.7.6) 的结论推广到任意完备 Noether 局部环的情形（IV, 6.12.7）。

  \item[(ii)] 若不再假设 $[k:k^p]$ 有限，定理 (22.7.3) 的结论就未必
  成立。例如，取 $k=\mathbf F_p(X_n)_{n\geqslant0}$，
  $A=k[[T,U]]$，并令 $\mathfrak q$ 为主理想 $Af$，其中
  \[
    f=U^p-\sum_{n=0}^{\infty}X_nT^{np}~;
  \]
  $A$ 是唯一分解环，而 $f$ 在其中是不可约元，因为它在多项式环
  $k((T))[U]$ 中不可约，从而在形式幂级数环 $k((T))[[U]]$ 中也不可约
  （Bourbaki，\emph{Alg. comm.}，第 VII 章，§~3），因而在 $A$ 中
  \emph{更是}不可约。取 $\mathfrak p=\mathfrak q$，于是
  $B_{\mathfrak p}$ 是 (21.9.7) 中以 $E$ 表示的域；该域在 $k$ 上可分
  （\emph{同上}），所以它是形式平滑 $k$-代数。然而，在这里采用
  (22.7.3) 的记号，有 $k=K_0=K$，而 $k$-导射（从 $A$ 到自身）都是
  $\partial/\partial T$ 与 $\partial/\partial U$ 的线性组合；由于
  $\partial f/\partial T=\partial f/\partial U=0$，(22.7.3) 的判据 b)
  与 b') 都不成立。
\end{enumerate}
\end{remarks}

\section{日本型整环}
\label{section:0.23-fr}

本节的结果将在 IV, (7.6) 与 (7.7) 中得到补充。

\subsection{日本环}
\label{subsection:0.23.1-fr}

\begin{definition}[23.1.1]
\label{0.23.1.1-fr}
若整环 $A$ 满足下述条件，则称它为\emph{日本环}：对其分式域的每个
有限扩张 $K'$（该分式域记作 $K$），$A'$（即 $A$ 在 $K'$ 中的整闭包）
都是有限型 $A$-模（换言之，是有限 $A$-代数）。

\oldpage[0\textsubscript{IV-1}]{214}

若环 $A$ 的每个有限型整 $A$-代数都是日本环，则称这个环是
\emph{普遍日本的}。
\end{definition}

显然，若 $A$ 是日本环，则每个分式环 $S^{-1}A$ 也是日本环；事实上，
采用上述记号，$S^{-1}A'$ 是 $S^{-1}A$ 在 $K'$ 中的整闭包，并且显然是
有限型 $S^{-1}A$-模。

Noether 整环，甚至\emph{离散赋值环}，也未必是日本环 [30]。

\begin{proposition}[23.1.2]
\label{0.23.1.2-fr}
设 $A$ 是 Noether 整环，$K$ 是其分式域。若对每个有限纯不可分扩张
$K'$（底域为 $K$），$A$ 在 $K'$ 中的整闭包都是有限型 $A$-模，则
$A$ 是日本环。
\end{proposition}

\begin{proof}
由于 $A$ 是 Noether 环，为验证 $A$ 是日本环，只须证明：对每个有限
拟伽罗瓦扩张 $L$（底域为 $K$），整闭包 $B$（即 $A$ 在 $L$ 中的
整闭包）都是有限型 $A$-模。然而，$L$ 是如下域的伽罗瓦扩张：$K'$，
即 $K$ 在 $L$ 中所含的最大纯不可分扩张；
若 $A'$ 是 $A$ 在 $K'$ 中的整闭包，则 $B$ 是 $A'$ 在 $L$ 中的整闭包。
但已知，在 $K'$ 的一个可分扩张中，$A'$ 的整闭包是有限型 $A'$-模
（Bourbaki，\emph{Alg. comm.}，第 V 章，§~1，第~6 号，命题~20 的
推论~1）。命题因而得证。
\end{proof}

由 (23.1.2)，当 $K$ 的特征为 $0$ 时，说 $A$ 是日本环，就等价于说
它的\emph{整闭包} $A'$ 是有限型 $A$-模。

\begin{theorem}[23.1.3]
\label{0.23.1.3-fr}
\textnormal{（Tate）}设 $A$ 是 Noether 整环，$x\neq0$ 是 $A$ 的元素。
假设下列条件成立：
\begin{enumerate}
  \item[(i)] $A$ 是整闭的。
  \item[(ii)] $\mathfrak p=xA$ 是素理想，并且 $A$ 对
  $\mathfrak p$-预进拓扑是分离且完备的。
  \item[(iii)] $A/xA$ 是日本环。
\end{enumerate}
则 $A$ 是日本环。
\end{theorem}

我们将使用下列引理：

\begin{lemma}[23.1.3.1]
\label{0.23.1.3.1-fr}
设 $A$ 是环，$x$ 是 $A$ 的一个元素，且不为 $0$ 的因子（在 $A$ 中），
并使 $xA=\mathfrak p$ 为素理想。
则对每个整数 $n>0$，$x^nA_{\mathfrak p}$ 在 $A$ 中的原像是 $x^nA$。
\end{lemma}

\begin{proof}
事实上，设 $b\in A$ 满足 $b/1=x^na/s$ 在 $A_{\mathfrak p}$ 中成立，
其中 $a\in A$ 且 $s\notin\mathfrak p$；于是存在
$s'\notin\mathfrak p$，使 $s'sb=x^nas'$，从而
$s'sb\in\mathfrak p$。由于 $s's\notin\mathfrak p$，这蕴含
$b\in\mathfrak p$，换言之，$b=xb'$，其中 $b'\in A$；又因 $x$ 不为
$0$ 的因子，遂有 $s'sb'=x^{n-1}as'$，再对 $n$ 作归纳即可。
\SourceCorrectionNote{法文印本及外交式转录在消去非零因子后的末项漏掉了 第二个 s 的撇号。应为上式。}。
\end{proof}

\begin{proof}
为证明定理，可以只考虑如下情形：分式域 $K$ 的底环为 $A$，其特征为
$p>0$；这是因为 $A$ 是整闭的。设 $K'$ 是 $K$ 的有限纯不可分扩张，于是存在一个幂
$q=p^e$，使 $K'^q\subset K$；把 $K'$ 换成一个更大的纯不可分扩张后，
甚至可以假设存在 $y\in K'$，使 $y^q=x$。设 $A'$ 是 $A$ 在 $K'$ 中的
整闭包；由于 $A$ 是整闭的，$A'$ 就是所有 $x'\in K'$ 中满足
$x'^q\in A'\cap K=A$ 者所成的集合。置
$V=A_{\mathfrak p}$；由于其极大理想 $\mathfrak m=xV$（属于 $V$）是主理想，
已知 Noether 局部环 $V$ 是\emph{离散赋值环}（17.1.4）；又因 $V$ 是
整闭的，与上面相同的论证表明：

\oldpage[0\textsubscript{IV-1}]{215}

$V'$（即 $V$ 在 $K'$ 中的整闭包）是所有 $x'\in K'$ 中满足
$x'^q\in V$ 者所成的集合。已知（Bourbaki，\emph{Alg. comm.}，第 VI 章，
§~8，第~6 号，命题~6；以及第 V 章，§~2，第~3 号，引理~4）$V'$ 是
\emph{离散赋值环}；其极大理想 $\mathfrak m'$ 是所有 $x'\in K'$ 中满足
$x'^q\in\mathfrak m$ 者所成的集合，而且剩余域
$V'/\mathfrak m'$ 是 $V/\mathfrak m$ 的\emph{有限}扩张
（\emph{同上}，第 VI 章，§~8，第~1 号，引理~2）。我们来证明：对每个
整数 $n>0$，有
\begin{equation}
\label{0.23.1.3.2-fr}
  \mathfrak m'^n\cap A'=y^nA'.
  \tag{23.1.3.2}
\end{equation}

事实上，由于 $y^q=x$，有 $y\in\mathfrak m'$，所以
$y^nA'\subset\mathfrak m'^n\cap A'$。反过来，取
$x'\in\mathfrak m'^n\cap A'$，并写成 $x'=z'y^n$，其中 $z'\in K'$。
有 $x'^q\in A'$；另一方面，$x'$ 是若干乘积
$t'_1\cdots t'_n$ 的和，其中 $t'_i\in\mathfrak m'$，因而
$t_i'^q\in\mathfrak m$，由此得到 $x'^q\in\mathfrak m^n\cap A$。
但引理 (23.1.3.1) 表明 $\mathfrak m^n\cap A=x^nA$，所以有
$z'^qx^n\in x^nA$；这也就是 $z'^q\in A$，因为 $x\neq0$。从而证明了
(23.1.3.2)。

其次证明 $A'$ 上的 $xA'$-预进拓扑是\emph{分离的}；这个拓扑也就是
$yA'$-预进拓扑（在 $A'$ 上），而由 (23.1.3.2)，它是
$\mathfrak m'$-预进拓扑（在 $V'$ 上）所诱导的拓扑。后者是分离的，因为 $V'$ 是
离散赋值环。

再证明 $A'/xA'$ 是有限型 $A$-模。由于 $x=y^q$，并且
$y^kA'/y^{k+1}A'$ 同构于 $A'/yA'$，只须证明 $A'/yA'$ 是有限型
$A$-模。但把 (23.1.3.2) 用于 $n=1$，可见 $A'/yA'$ 是
$V'/\mathfrak m'$ 的子环；后者就是 $V'$ 的剩余域，并且是
$V/\mathfrak m$（即 $V$ 的剩余域）的有限扩张。又因 $V/\mathfrak m$ 是
$A/\mathfrak p$ 的分式域，而 $A'$ 在 $A$ 上整，所以 $A'/yA'$ 在
$A/\mathfrak p$ 上整，从而包含于 $A/\mathfrak p$ 在
$V'/\mathfrak m'$ 中的整闭包。由于按假设 $A/\mathfrak p$ 是 Noether
日本环，$A'/yA'$ 是有限型 $(A/\mathfrak p)$-模，因而\emph{更是}
有限型 $A$-模。

现在，$xA'$-预进拓扑（在 $A'$ 上）既然是分离的，$A'$ 就是它关于这个拓扑的
完备化 $\widehat{A'}$ 的一个子环；但因为 $A'/xA'$ 是有限型 $A$-模，
且 $A$ 对 $xA$-预进拓扑分离而完备，由 (0\textsubscript{I}, 7.2.9)，
$\widehat{A'}$ 是有限型 $A$-模，所以 $A'$ 也如此。证毕。
\end{proof}

\begin{corollary}[23.1.4]
\label{0.23.1.4-fr}
设 $A$ 是整闭的 Noether 日本环。则每个形式幂级数环
$A[[T_1,\ldots,T_r]]$ 都是日本环。
\end{corollary}

\begin{proof}
已知对每个整闭 Noether 环 $B$，形式幂级数环 $B[[T]]$ 都是 Noether
且整闭的（Bourbaki，\emph{Alg. comm.}，第 V 章，§~1，第~4 号，
命题~14）；因而对 $r$ 作归纳后，只须证明 $A[[T]]$ 是日本环。
而元素 $x=T$ 满足 (23.1.3) 的全部条件，结论随即成立。
\SourceCorrectionNote{法文印本及外交式转录把命题中的形式变量个数 r 误写为未约束的 n。应为上述归纳 变量。}。
\end{proof}

\begin{theorem}[23.1.5]
\label{0.23.1.5-fr}
\textnormal{（Nagata）}每个完备 Noether 局部整环 $A$ 都是日本环。
\end{theorem}

\begin{proof}
已知（19.8.8, (ii)），存在一个子环 $B$（属于 $A$），它是完备正则局部环，
并且 $A$ 是有限 $B$-代数；显然，只须证明

\oldpage[0\textsubscript{IV-1}]{216}

$B$ 是日本环。换言之，可以只考虑 $A$ 还是\emph{正则环}的情形。
对 $n=\dim(A)$ 作归纳；当 $n=0$ 时定理显然成立。设 $n>0$，以
$\mathfrak m$ 表示 $A$ 的极大理想，并取 $x$ 为
$\mathfrak m-\mathfrak m^2$ 中的元素；已知（17.1.8）$A/xA$ 是正则环，
而且由 (17.1.7) 与 (16.3.4)，有 $\dim(A/xA)=n-1$；此外，理想 $xA$
在 $A$ 中是闭的（0\textsubscript{I}, 7.3.5），故 $A/xA$ 是完备的。
由归纳假设，$A/xA$ 是日本环，再由 (23.1.3)，$A$ 也是日本环。
\end{proof}

\begin{corollary}[23.1.6]
\label{0.23.1.6-fr}
设 $A$ 是完备 Noether 局部整环，$K$ 是其分式域，$K'$ 是 $K$ 的有限
扩张；则整闭包 $A'$（即 $A$ 在 $K'$ 中的整闭包）是完备局部环。
\end{corollary}

\begin{proof}
由 (23.1.5)，已知 $A'$ 是有限 $A$-代数，因而是完备 Noether 半局部环，
所以是若干局部环的\emph{直积}；但 $A'$ 是整环，故它必为局部环。
\end{proof}

\begin{proposition}[23.1.7]
\label{0.23.1.7-fr}
设 $A$ 是 Noether 局部整环，$K$ 是其分式域，$K'$ 是 $K$ 的有限扩张，
$A'$ 是 $A$ 在 $K'$ 中的整闭包。假设完备化 $\widehat A$ 是约化的，并以
$R$ 表示它的全分式环。
\begin{enumerate}
  \item[(i)] 若环 $R\otimes_KK'$ 是约化的（特别地，当 $K'$ 是 $K$ 的
  可分扩张时，这个条件成立；这里 $R$
  是若干域的直积，这些域都是 $K$ 的扩张（Bourbaki，\emph{Alg.}，
  第 VIII 章，§~7，第~3 号，定理~1 的推论~1）），则 $A'$ 是有限型
  $A$-模。
  \item[(ii)] 特别地，若 $R$ 是可分 $K$-代数，则 $A$ 是日本环。
\end{enumerate}
\end{proposition}

\begin{proof}
\emph{(i)} 为简化记号，置 $A_1=\widehat A$、
$A'_1=A'\otimes_AA_1$、$K'_1=K'\otimes_AA_1$。由于 $A_1$ 是平坦
$A$-模（0\textsubscript{I}, 7.3.5），$A'_1$ 可视为 $K'_1$ 的一个子环，
并且显然在 $A_1$ 上整。另一方面，若置 $K_1=K\otimes_AA_1$，则可写成
$K'_1=K'\otimes_KK_1$；因为 $A$ 是整环且 $A_1$ 是平坦 $A$-模，
每个元素 $\neq0$（属于 $A$）都是 $A_1$-正则的（0\textsubscript{I}, 6.3.4）。由于
$K=S^{-1}A$，其中 $S=A-\{0\}$，有 $K_1=S^{-1}A_1$；上述说明表明
$K_1$ 可视为 $R$（即 $A_1$ 的全分式环）的一个子环。由于每个 $K$-模都是
平坦的，$K'_1$ 可视为 $K'\otimes_KR$ 的一个子环；按假设，后者是
约化的有限 $R$-代数。以 $\mathfrak q_i$（$1\leqslant i\leqslant n$）
表示 $A_1$ 的诸极小素理想，以 $B_i$ 表示整环
$A_1/\mathfrak q_i$，以 $L_i$ 表示 $B_i$ 的分式域；于是 $R$ 是诸
$L_i$ 的直积，而 $K'\otimes_KR$ 是诸 $K'\otimes_KL_i$ 的直积。
后面这些 $K$-代数既然是约化的，就是若干域的直积，而这些域是 $L_i$
的有限扩张。由于 $B_i$ 是\emph{完备的}局部整环，Nagata 定理
(23.1.5) 表明，$B_i$ 在 $K'\otimes_KL_i$ 中的整闭包是有限型
$B_i$-模，因而\emph{更是}有限型 $A_1$-模；又因
$K'\otimes_KL_i$ 中每个在 $A_1$ 上整的元素也在 $B_i$ 上整（因为它被
$\mathfrak q_i$ 零化），最终可知 $A_1$ 在 $K'\otimes_KR$ 中的整闭包
是有限型 $A_1$-模。但 $A'_1$ 包含于这个整闭包中，且 $A_1$ 是 Noether
环，所以 $A'_1$ 也是有限型 $A_1$-模。最后，由于 $A_1$ 是忠实平坦
$A$-模（0\textsubscript{I}, 7.3.5），遂知 $A'$ 是有限型 $A$-模
（Bourbaki，\emph{Alg. comm.}，第 I 章，§~3，第~6 号，命题~11）。

\oldpage[0\textsubscript{IV-1}]{217}

\emph{(ii)} 该假设蕴含诸 $L_i$ 都是 $K$ 的可分扩张，所以
$K'\otimes_KL_i$ 是约化的（Bourbaki，\emph{Alg.}，第 VIII 章，§~7，
第~3 号，定理~1），并且可以把 (i) 应用于每个有限扩张 $K'$（底域为 $K$）；
这就证明了断言。
\end{proof}

\subsection{Noether 局部整环的整闭包}
\label{subsection:0.23.2-fr}

\begin{env}[23.2.1]
\label{0.23.2.1-fr}
下文中，对每个环 $A$，为简便起见，我们以 $A_{\mathrm{red}}$ 表示 $A$
模其幂零根的商（因而，若 $X=\operatorname{Spec}(A)$，则有
$X_{\mathrm{red}}=\operatorname{Spec}(A_{\mathrm{red}})$）。

若\emph{局部}环 $A$ 满足：$A_{\mathrm{red}}$ 是整环，并且
$A_{\mathrm{red}}$ 的整闭包是\emph{局部}环，则称这个环是\emph{单支的}；
这推广了 (III, 4.3.6) 中给出的定义。若 $A$ 是单支的，而且局部环
$A_{\mathrm{red}}$ 的整闭包的剩余域，是 $A$ 的剩余域的
\emph{纯不可分扩张}，则称这个环是\emph{几何单支的}。显然，整闭局部环
是几何单支的；由 (23.1.6)，完备 Noether 局部整环是单支的。
\end{env}

\begin{lemma}[23.2.2]
\label{0.23.2.2-fr}
\textnormal{(*)}\footnote{在本号以下的部分，我们使用第 IV 章 §§~2 与 5
中陈述的概念和结果；由于本号的结果在第 IV 章 §~6 以前不会使用，故不
导致循环论证。} 设 $A$ 是整环，$A'$ 是其整闭包。典范态射
$f:\operatorname{Spec}(A')\to\operatorname{Spec}(A)$ 为根态射的充分必要
条件是：对每个 $\mathfrak p\in\operatorname{Spec}(A)$，
$A_{\mathfrak p}$ 都是几何单支的；在这种情形，$f$ 是同胚。
\end{lemma}

\begin{proof}
事实上，对每个 $\mathfrak p\in\operatorname{Spec}(A)$，
$A_{\mathfrak p}$ 的整闭包是 $A'_{\mathfrak p}$（Bourbaki，
\emph{Alg. comm.}，第 V 章，§~1，第~5 号，命题~16），并且
$A'_{\mathfrak p}$ 中位于 $\mathfrak pA_{\mathfrak p}$ 上方的所有素理想
都是极大理想（\emph{同上}，§~2，第~1 号，命题~1）。因此，说 $f$ 是
单射，等价于说对每个 $\mathfrak p\in\operatorname{Spec}(A)$，
$A'_{\mathfrak p}$ 都是\emph{局部}环，也就是说，诸
$A_{\mathfrak p}$ 都是单支的。于是，说每个 $A_{\mathfrak p}$ 都是
几何单支的，由 (I, 3.5.8) 等价于说 $f$ 是根态射。在这种情形，$f$ 是
满的且闭的（II, 6.1.10），因而是同胚。
\end{proof}

\begin{lemma}[23.2.3]
\label{0.23.2.3-fr}
设 $A$ 是整环，$K$ 是其分式域，$B$ 是 Noether 环，
$\varphi:A\to B$ 是使 $B$ 成为平坦 $A$-模的环同态。设 $K'$ 是 $K$ 的
一个扩张；置 $B_1=B_{\mathrm{red}}$，并以 $R$ 表示 $B_1$ 的全分式环。
则对每个 $A$-子代数 $A'$（属于 $K'$），典范同态
\begin{equation}
\label{0.23.1.8.1-fr}
  (A'\otimes_AB)_{\mathrm{red}}
  \longrightarrow(K'\otimes_AR)_{\mathrm{red}}
  \tag{23.1.8.1}
\end{equation}
是单射。
\end{lemma}

\begin{proof}
事实上，典范同态 $h:A'\otimes_AB\to K'\otimes_AR$ 可视为下列诸同态的
复合：
\[
  A'\otimes_AB\xrightarrow{u}K'\otimes_AB
  \xrightarrow{v}K'\otimes_AB_1\xrightarrow{w}K'\otimes_AR.
\]
由于 $B$ 是平坦 $A$-模，$u$ 是单射；同样，由于 $K'$ 是平坦 $K$-模、
$K$ 是平坦 $A$-模，$K'$ 是平坦 $A$-模。又有 $w$ 是单射，因为
$B_1$ 是约化的，故同态 $B_1\to R$ 是单射。最后，基于同一理由，若
$\mathfrak N$ 是

\oldpage[0\textsubscript{IV-1}]{218}

$B$ 的幂零根，则 $v$ 的核是 $K'\otimes_A\mathfrak N$，故包含于
$K'\otimes_AB$ 的幂零根中。由此立即可知：若在
$h=w\circ v\circ u$ 下，某个元素 $x\in A'\otimes_AB$ 的像是幂零的，则 $x$
也是幂零的。这就证明了引理。
\end{proof}

\begin{proposition}[23.2.4]
\label{0.23.2.4-fr}
设 $A$ 是整环，$K$ 是其分式域，$B$ 是 Noether 环；以
$\mathfrak q_i$（$1\leqslant i\leqslant n$）表示它的诸极小素理想，并设
$\varphi:A\to B$ 是环同态。假设诸 $B/\mathfrak q_i$ 都是日本环，并且
$B$ 是平坦 $A$-模。设 $K'$ 是 $K$ 的有限扩张，并以 $(C_\lambda)$ 表示
$K'$ 的诸子环所成的递增有向族；这些子环都是有限 $A$-代数，并以 $K'$
为分式域。因此，$A'$（诸 $C_\lambda$ 的并）是 $A$ 在 $K'$ 中的整闭包。
则：
\begin{enumerate}
  \item[(i)] 存在指标 $\alpha$，使得对 $\lambda\geqslant
  \alpha$，典范同态
  \begin{equation}
  \label{0.23.2.4.1-fr}
    (C_\alpha\otimes_AB)_{\mathrm{red}}
    \longrightarrow(C_\lambda\otimes_AB)_{\mathrm{red}}
    \tag{23.2.4.1}
  \end{equation}
  是双射。
  \item[(ii)] 再假设 $B$ 是忠实平坦 $A$-模，则态射
  $\operatorname{Spec}(A')\to\operatorname{Spec}(C_\alpha)$ 是根态射。
\end{enumerate}
\end{proposition}

\begin{proof}
\emph{(i)} 置 $B_1=B_{\mathrm{red}}$，并以 $R$ 表示 $B_1$ 的全分式环。
由于 $A$ 是整环且 $B$ 是平坦 $A$-模，每个元素 $\neq0$（属于 $A$）都是
$B$-正则的（0\textsubscript{I}, 6.3.4），故复合同态
$A\xrightarrow{\varphi}B\to B_1$ 延拓为同态 $K\to R$。于是可写成
$K'\otimes_AR=(K'\otimes_AK)\otimes_KR$；又因 $K'\otimes_AK$ 可与
$K'$ 等同，所以 $K'\otimes_AR=K'\otimes_KR$。由于 $R$ 是诸
$L_i$（即 $B/\mathfrak q_i$ 的分式域）的直积，$K'\otimes_AR$ 是诸
$K'\otimes_KL_i$ 的直积；因此 $(K'\otimes_AR)_{\mathrm{red}}$ 是
有限多个 $L_i$ 的有限扩张的直积。由假设，$B/\mathfrak q_i$ 在
$L_i$ 的一个有限扩张中的整闭包是有限型 $(B/\mathfrak q_i)$-模，
因而是有限型 $B$-模；由此可知，$B'$（即 $B$ 在
$(K'\otimes_AR)_{\mathrm{red}}$ 中的整闭包）是有限型 $B$-模。
由 (23.2.3)，诸 $(C_\lambda\otimes_AB)_{\mathrm{red}}$ 典范地等同于
$(K'\otimes_AR)_{\mathrm{red}}$ 的子环；它们都是有限 $B$-代数，故包含于
$B'$ 中。由于 $B$ 是 Noether 环且 $B'$ 是有限型 $B$-模，有向族
$(C_\lambda\otimes_AB)_{\mathrm{red}}$ 有一个\emph{最大元}
$(C_\alpha\otimes_AB)_{\mathrm{red}}$。这就证明了 (i)。

\emph{(ii)} 由于 $B$ 是忠实平坦 $A$-模，由 (IV, 2.6.1, (v))，只须证明
态射 $\operatorname{Spec}(A'\otimes_AB)\to
\operatorname{Spec}(C_\alpha\otimes_AB)$ 是根态射；或者等价地，由
(I, 5.1.6)，只须证明态射
$\operatorname{Spec}(A'\otimes_AB)_{\mathrm{red}}\to
\operatorname{Spec}(C_\alpha\otimes_AB)_{\mathrm{red}}$ 是根态射。但有
$A'\otimes_AB=\varinjlim_\lambda(C_\lambda\otimes_AB)$，所以由
(IV, 5.13.2)，
$(A'\otimes_AB)_{\mathrm{red}}=
\varinjlim_\lambda(C_\lambda\otimes_AB)_{\mathrm{red}}$；结论遂由 (i) 得出。
\end{proof}

\begin{corollary}[23.2.5]
\label{0.23.2.5-fr}
设 $A$ 是 Noether 局部整环，$K$ 是其分式域，$K'$ 是 $K$ 的有限扩张。
以 $(C_\lambda)$ 表示 $K'$ 的诸子环所成的递增有向族；这些子环都是有限
$A$-代数（因而是 Noether 半局部环（Bourbaki，\emph{Alg. comm.}，
第 IV 章，§~2，第~5 号，命题~9 的推论~3）），并以 $K'$ 为分式域。
则存在 $\alpha$，使得同态
$(\widehat C_\alpha)_{\mathrm{red}}\to
(\widehat C_\lambda)_{\mathrm{red}}$ 对 $\lambda\geqslant\alpha$ 是同构；并且，若 $A'$ 是 $A$ 在
$K'$ 中的整闭包，则态射
$\operatorname{Spec}(A')\to\operatorname{Spec}(C_\alpha)$ 是根态射
（参见 (23.2.2)）。
\end{corollary}

\oldpage[0\textsubscript{IV-1}]{219}

\begin{proof}
在 (23.2.4) 中取 $B=\widehat A$。诸 $B/\mathfrak q_i$ 都是完备 Noether
局部环，所以都是日本环（23.1.5）；$B$ 是忠实平坦 $A$-模
（0\textsubscript{I}, 7.3.5）；此外，还有
$C_\lambda\otimes_AB=\widehat C_\lambda$（Bourbaki，
\emph{Alg. comm.}，第 III 章，§~3，第~4 号，定理~3；以及第 IV 章，
§~2，第~5 号，命题~9 的推论~3）。
\end{proof}

\begin{corollary}[23.2.6]
\label{0.23.2.6-fr}
在 (23.2.5) 的假设下，$A'$（即 $A$ 在 $K'$ 中的整闭包）是半局部环；若
$\mathfrak m$ 是 $A$ 的极大理想，则对每个极大理想 $\mathfrak m'$
（属于 $A'$），$A'/\mathfrak m'$ 是 $A/\mathfrak m$ 的有限扩张。
\end{corollary}

\begin{proof}
同态 $(\widehat C_\alpha)_{\mathrm{red}}\to
(\widehat C_\lambda)_{\mathrm{red}}$ 对 $\lambda\geqslant\alpha$ 是双射；
这蕴含诸 $C_\lambda$ 的
极大理想个数在 $\lambda\geqslant\alpha$ 时恒定。若
$\mathfrak m_\alpha$ 是 $C_\alpha$ 的一个极大理想，而
$\mathfrak m_\lambda$ 是 $C_\lambda$ 中位于 $\mathfrak m_\alpha$ 上方的
唯一极大理想，则域 $C_\alpha/\mathfrak m_\alpha$ 与
$C_\lambda/\mathfrak m_\lambda$ 典范同构。结论来自
$\mathfrak m'=\varinjlim\mathfrak m_\lambda$ 与
$A'/\mathfrak m'=\varinjlim(C_\lambda/\mathfrak m_\lambda)$
（0\textsubscript{III}, 10.3.1.1.3）。
\end{proof}

\begin{proposition}[23.2.7]
\label{0.23.2.7-fr}
\textnormal{（Y. Mori）}设 $A$ 是 Noether 局部整环，$K$ 是其分式域，
$A'$ 是 $A$ 的整闭包。则 $A'$ 是半局部 Krull 环（Bourbaki，
\emph{Alg. comm.}，第 VII 章，§~1）；换言之，存在一族离散赋值环
$(V_\lambda)$，它们都以 $K$ 为分式域，并满足
$A'=\bigcap_\lambda V_\lambda$，而每个 $x\in K$ 只不属于有限多个
$V_\lambda$。此外，存在一个子环 $C$（属于 $K$），它是有限 $A$-代数，
使得诸 $V_\lambda$ 是诸环 $C_{\mathfrak p_\lambda}$ 的整闭包，其中
$(\mathfrak p_\lambda)$ 是所有高度为 $1$ 的素理想（属于局部环 $C$）所成的族。
\end{proposition}

先证明下列引理：

\begin{lemma}[23.2.7.1]
\label{0.23.2.7.1-fr}
设 $A$、$B$ 是两个环，$\varphi:A\to B$ 是使 $B$ 成为忠实平坦 $A$-模的
同态。假设 $A$ 是整环；若 $C=B_{\mathrm{red}}$，则复合同态
$\psi:A\to B\to B_{\mathrm{red}}=C$ 是单射，并延拓为从 $K$（$A$ 的
分式域）到 $R$（$C$ 的全分式环）的单同态。此外，若 $C'$ 是 $C$ 在 $R$ 中的
整闭包，则 $C'\cap K$ 是 $A$ 的整闭包。
\end{lemma}

\begin{proof}
由于 $A$ 是整环且 $\varphi$ 是单射，$\varphi(A)$ 与 $\mathfrak N$
（即 $B$ 的幂零根）的交只有 $0$；又因 $C=B/\mathfrak N$，$\psi$ 是单射。
每个 $a\neq0$（属于 $A$）在 $A$ 中都不是零因子；由平坦性，它在 $B$ 中
也不是零因子（0\textsubscript{I}, 6.3.4）。由此可知，$a$ 也不是
$0$ 的零因子（在 $C$ 中）：事实上，若有 $ax\in\mathfrak N$，其中
$x\notin\mathfrak N$（在 $B$ 中），则有 $a^nx^n=0$（对某个整数 $n$）；
但由于 $x^n\neq0$，这同上面的结论矛盾。因此，可以把
$\psi$ 延拓为从 $K$ 到 $R$ 的单同态。为证明最后一个断言，先注意
$A'$（即 $A$ 的整闭包）显然包含于 $C'\cap K$。反过来，取
$x\in C'\cap K$；于是 $x$ 在 $C$ 上整，因而\emph{更是}在 $B$ 上整，
换言之，$B[x]$ 是有限 $B$-代数。此外，$K$ 可视为
$B\otimes_AK$ 的一个子环，而 $B[x]$（$B\otimes_AK$ 的子环）由平坦性
可与 $B\otimes_AA[x]$ 等同；从而 $A[x]$ 是有限型 $A$-模
（Bourbaki，\emph{Alg. comm.}，第 I 章，§~3，第~6 号，命题~11），
所以 $x\in A'$。
\end{proof}

引理既已证明，把它应用于从 $A$ 到其完备化 $\widehat A$ 的典范单射；
后者是忠实平坦 $A$-模。环
$B=(\widehat A)_{\mathrm{red}}=\widehat A/\mathfrak N$（其中
$\mathfrak N$ 是 $\widehat A$ 的幂零根）是约化的完备 Noether 局部环；
其全分式

\oldpage[0\textsubscript{IV-1}]{220}

环 $R$ 是有限多个域 $L_i$ 的直积。$B_i$（$B$ 在 $L_i$ 中的典范像）是
\emph{完备的局部 Noether 整环}，其分式域为 $L_i$；而 $B'$（$B$ 在
$R$ 中的整闭包）是诸 $B'_i$ 的\emph{直积}，其中 $B'_i$ 是 $B_i$ 的整闭包。
但由 (23.1.5)，$B'_i$ 是有限 $B_i$-代数，因而是整闭的\emph{Noether}
局部环。已知（Bourbaki，\emph{Alg. comm.}，第 VII 章，§~1，第~3 号，
定理~2 的推论）$B'_i$ 是 Krull 环。现在，对每个 $i$，有单同态
$\varphi_i:K\to L_i$，从而 $\varphi_i^{-1}(B'_i)$ 是 $K$ 中的 Krull 环。
但是，由引理有 $A'=B'\cap K$，并且
$B'=\bigcap_i\varphi_i^{-1}(B'_i)$；所以 $A'$ 是有限多个 Krull 环的交，
因而是 Krull 环。

为证明命题的最后一个断言，注意存在一个有限 $A$-代数
$C\subset K$，使态射
$\operatorname{Spec}(A')\to\operatorname{Spec}(C)$ 是根态射兼同胚
（23.2.4）。对每个素理想 $\mathfrak p'\in\operatorname{Spec}(A')$，
$\mathfrak p'$ 是 $A'$ 中位于 $\mathfrak p=\mathfrak p'\cap C$ 上方的
唯一素理想，而 $A'_{\mathfrak p'}$ 是 $C_{\mathfrak p}$ 的整闭包；
另一方面，$\operatorname{Spec}(A')\to\operatorname{Spec}(C)$ 是同胚，
这蕴含映射 $\mathfrak p'\mapsto\mathfrak p'\cap C$ 把所有高度为
$1$ 的素理想（属于 $A'$）双射到所有高度为 $1$ 的素理想（属于 $C$）。
结合 Bourbaki，
\emph{Alg. comm.}，第 VII 章，§~1，第~6 号，定理~4，即得结论。

\begin{remarks}[23.2.8]
\label{0.23.2.8-fr}
\begin{enumerate}
  \item[(i)] 应当注意，在应用 (23.2.6) 与 (23.2.7) 时，环 $A'$
  \emph{未必是 Noether 环}；Nagata 在 $\dim(A)=3$ 的例子 [30] 说明了
  这一点。

  \item[(ii)] 若 $A'$ 是 $A$ 在一个有限扩张 $K'$（底域为 $K$）中的整闭包，
  (23.2.7) 的结论仍然成立。事实上，只须考虑一个有限 $A$-代数 $B$，
  其分式域为 $K'$。$B$ 是 Noether 半局部环，故为有限多个 Noether
  局部环 $B_{\mathfrak m_i}$ 的交，其中诸 $\mathfrak m_i$ 是 $B$ 的极大
  理想；它的整闭包（等于 $A'$）是诸 $A'_{\mathfrak m_i}$ 的交
  （Bourbaki，\emph{Alg. comm.}，第 II 章，§~3，第~3 号，定理~1 的
  推论~4）。这些环是诸 $B_{\mathfrak m_i}$ 的整闭包，因而由 (23.1.13)
  是 Krull 环；故 $A'$ 是 Krull 环。

  \item[(iii)] 可以证明（[30], 3.3.10），对任意 Noether 整环 $A$
  （不必是局部环），$A$ 的整闭包都是 Krull 环。
\end{enumerate}
\end{remarks}

\begin{corollary}[23.2.9]
\label{0.23.2.9-fr}
设 $A$ 是 Noether 整环，$A'$ 是其整闭包。假设对每个环 $C$，只要
$A\subset C\subset A'$ 且 $C$ 是有限 $A$-代数，并对每个素理想
$\mathfrak q$（高度为 $1$，属于 $C$），都有 $\mathfrak q\cap A$ 是
高度为 $1$ 的素理想（属于 $A$）。以 $P$ 表示所有高度为 $1$ 的素理想
（属于 $A$）所成的集合；则有
\[
  A'=\bigcap_{\mathfrak p\in P}A'_{\mathfrak p}.
\]
\end{corollary}

\begin{proof}
由于 $A'$ 是无挠 $A$-模，有 $A'=\bigcap A'_{\mathfrak m}$，其中
$\mathfrak m$ 遍历 $A$ 的所有极大理想（Bourbaki，\emph{Alg. comm.}，
第 II 章，§~3，第~3 号，定理~1 的推论~4）。

\oldpage[0\textsubscript{IV-1}]{221}

因此，只须证明 $A'_{\mathfrak m}$ 是诸 $A'_{\mathfrak p}$ 的交，其中
$\mathfrak p\subset\mathfrak m$ 遍历高度为 $1$ 的素理想。由于
$A'_{\mathfrak m}$ 是 $A_{\mathfrak m}$ 的整闭包，可见只须对
$A_{\mathfrak m}$ 证明本推论。此时由 (23.2.7)，存在有限
$A_{\mathfrak m}$-代数 $B$，且它包含于 $A'_{\mathfrak m}$ 中，使
$A'_{\mathfrak m}$ 是诸环 $B_{\mathfrak q}$ 的整闭包的交，其中
$\mathfrak q$ 遍历所有高度为 $1$ 的素理想（属于 $B$）。但 $B$ 具有
$C_{\mathfrak m}$ 的形式，其中 $C$ 是有限 $A$-代数；所以，对
每个素理想 $\mathfrak q$（高度为 $1$，属于 $B$），按假设，
$\mathfrak q\cap A$ 是高度为 $1$ 的素理想（属于 $A$）。由于 $B_{\mathfrak q}$ 的整闭包包含
$A_{\mathfrak q\cap A}$ 的整闭包，而后者又包含 $A'_{\mathfrak m}$，
遂知 $A'_{\mathfrak m}$ 确为所有 $A'_{\mathfrak p}$ 的交，其中
$\mathfrak p\in P$ 且 $\mathfrak p\subset\mathfrak m$。证明完毕。
\end{proof}

\begin{remarks}[23.2.10]
\label{0.23.2.10-fr}
\begin{enumerate}
  \item[(i)] 我们将看到（IV, 5.6.3 与 5.6.10），当 $A$ 是
  \emph{普遍链状的}时，(23.2.9) 的假设成立；特别地（IV, 5.10.17 与
  5.11.2），当 $A$ 是\emph{正则环的商}时亦然。最后，若
  $\operatorname{Spec}(A')\to\operatorname{Spec}(A)$ 是同胚，该假设显然
  成立；特别地，当这个态射是根态射时成立（23.2.2）。

  \item[(ii)] 若省去关于有限 $A$-子代数（属于 $A'$）的假设，(23.2.9) 的
  结论就未必成立。事实上，可以构造一个 Noether 局部整环 $A$，其维数为
  $2$，其整闭包 $A'$ 是有限 $A$-代数，但存在一个素理想
  $\mathfrak p'$（属于 $A'$，高度为 $1$），位于极大理想 $\mathfrak m$
  （高度为 $2$，属于 $A$）上方；此外，对每个素理想 $\mathfrak p$
  （高度为 $1$，属于 $A$），
  $A_{\mathfrak p}$ 都是整闭的（IV, 5.6.11）。这些环
  $A_{\mathfrak p}$ 的交不可能等于 $A'$，因为 $A'$ 是 Krull 环
  （Bourbaki，\emph{Alg. comm.}，第 VII 章，§~1，第~5 号，命题~9）。
\end{enumerate}
\end{remarks}

\begin{flushright}
\emph{（待续。）}
\end{flushright}

\cleardoublepage
\markboth{第 IV 章}{第 IV 章}
\oldpage[IV-1]{222}
\thispagestyle{plain}

\begin{center}
{\Large 第 IV 章}

\vspace{1.5cm}
{\LARGE\bfseries 概形与态射的\\[0.3em]
局部性质}

\vspace{1.5cm}
{\large\bfseries 目录}
\end{center}

\medskip
\begin{tabular}{@{}r p{0.84\textwidth}}
§ 1. & 相对有限性条件。预概形中的可构造集。\\
§ 2. & 基变换与平坦性。\\
§ 3. & 相伴素循环与准素分解。\\
§ 4. & 预概形的基域变换。\\
§ 5. & 预概形中的维数与深度。\\
§ 6. & 局部 Noether 预概形的平坦态射。\\
§ 7. & Noether 局部环与其完备化之间关系的应用。优环。\\
§ 8. & 预概形的逆极限。\\
§ 9. & 可构造性质。\\
§ 10. & Jacobson 预概形。\\
§ 11.\footnotemark[1] & 有限表示平坦态射的拓扑性质。平坦性的局部判别准则。\\
§ 12. & 有限表示平坦态射的纤维研究。\\
§ 13. & 等维态射。\\
§ 14. & 泛开态射。\\
§ 15. & 泛开态射的纤维研究。\\
§ 16. & 微分不变量。微分光滑态射。\\
§ 17. & 光滑、非分歧与平展态射。\\
§ 18. & 平展态射的补充。Hensel 局部环。\\
§ 19. & 正则嵌入与横截正则嵌入。\\
§ 20. & 超平面截面；一般投影。\\
§ 21. & 无穷小延拓。\\
\end{tabular}
\footnotetext[1]{§§ 11 至 21 的次序和内容仅供参考，在出版以前可能有所修改。}

\clearpage
\oldpage[IV-1]{223}

本章所研究的主题需要作如下说明：

\medskip
\noindent
\emph{a)} 它们的共同特征，是涉及预概形或态射的\emph{局部}性质，即在一点、
一条纤维的各点，或一点或一条纤维的某个（未指定的）邻域中考察的性质。这些
性质通常具有\emph{拓扑}、\emph{微分}或\emph{维数}性质（即涉及\emph{维数}
和\emph{深度}的概念），并同所考察各点处的\emph{局部环}的性质相联系。一类
典型问题是：对于给定态射 $f:X\to Y$ 和一点 $x\in X$，把 $X$ 在 $x$ 处的
性质同 $Y$ 在 $y=f(x)$ 处以及纤维 $X_y=f^{-1}(y)$ 在 $x$ 处的性质联系起来。
另一类问题是：确定点集的拓扑性质（例如可构造性，或者是否为开集或闭集）；
这里的点 $x\in X$ 是使得 $X$ 在该点具有某一性质，或使得纤维
$X_{f(x)}$ 通过 $x$ 并在 $x$ 处具有某一性质的点。同样，我们也要研究由如下
点 $y\in Y$ 所成集合的拓扑性质：$X$ 在纤维 $X_y$ 的每一点都具有某一性质，
或该纤维具有某一性质。

\medskip
\noindent
\emph{b)} 对后续各章最为重要的概念，是\emph{有限表示的平坦态射}，以及作为
其特殊情形的\emph{光滑态射}和\emph{平展态射}。对这些概念（连同有关问题）
的详细研究，严格说来从 §~11 开始。

\medskip
\noindent
\emph{c)} §§~1 至 10 可以视为预备性内容；其中发展三类方法，不但用于本章其余
各节，当然也用于后续各章：

\medskip
\noindent
\emph{c 1)} §§~1 至 4 从不同方面考察\emph{基变换}的概念，特别是它同
\emph{有限性}或\emph{平坦性}条件的关系；这里以最初等的形式开始发展
\emph{下降}方法（与这一方法有关的“有效性”问题将在第 V 章研究）。

\medskip
\noindent
\emph{c 2)} §§~5 至 7 围绕所谓的\emph{Noether 方法}展开：所考察的预概形
始终假设为局部 Noether，而对\emph{态射}通常反而不施加任何有限性条件；
这主要是因为维数和深度的概念，几乎只有对 Noether 局部环才便于处理。还要
指出，§~7 构成 Noether 局部环的一种“精细”理论，而在本章后文中很少使用。

\medskip
\noindent
\emph{c 3)} §§~8 至 10 尤其给出一种办法，可以\emph{消去}所研究预概形上的
\emph{Noether 假设}，代之以所考察\emph{态射}上的适当\emph{有限性}假设
（“有限表示”）；这种替换的优点在于，后一类假设在\emph{基变换下稳定}，而
预概形上的 Noether 假设并不具有这一性质。实现这种过渡的方法一方面依靠
预概形\emph{逆极限}的概念，借此可以把一个问题归约为 Noether 假设下的同一
问题；另一方面依靠系统地使用\emph{可构造集}，它们有双重优点：在任意态射
的逆像下仍为可构造集，在有限表示态射的正像下也仍为可构造集，并且在局部

\oldpage[IV-1]{224}

Noether 预概形中具有便于处理的拓扑性质。同样的方法甚至常能归约到更特殊的
Noether 环，例如有限型 $\mathbf Z$-\emph{代数}；正是在这里，“优”环（§~7
中研究）的性质起决定性作用。此外，即使不考虑消去 Noether 假设的问题，
§§~8 至 10 的方法本身也很初等，几乎在所有应用中都经常使用。

\setcounter{section}{0}

\section{相对有限性条件。预概形中的可构造集}
\label{section:IV.1-fr}

本节将补充并重新叙述 (I, 6.3 与 6.6) 中关于预概形态射 $f:X\to Y$ 的
“有限性条件”。本质上，有两种在 $X$ 上具有\emph{整体}性质的“有限性”概念：
\emph{拟紧}态射（定义见 (I, 6.6.1)）和\emph{拟分离}态射；另一方面，也有
两种在 $X$ 上具有\emph{局部}性质的“有限性”概念：\emph{局部有限型}态射
（定义见 (I, 6.6.2)）和\emph{局部有限表示}态射。把这些局部概念同前述整体
概念结合起来，就得到\emph{有限型}态射（定义见 (I, 6.3.1)）和
\emph{有限表示}态射的概念。为方便读者，我们将在本节重新给出 (I, 6.3 与
6.6) 中陈述的主要性质；其证明当然仍请参见第一章的相应各号。

在 (1.8) 和 (1.9) 中，我们将在预概形的框架内，利用上述有限性概念，补充
(0\textsubscript{III}, §~9) 给出的关于可构造集的结果。

\subsection{拟紧态射}
\label{subsection:IV.1.1-fr}

\begin{definition}[1.1.1]
\label{IV.1.1.1-fr}
若预概形态射 $f:X\to Y$ 所给出的连续映射 $f$（从拓扑空间 $X$ 到拓扑空间
$Y$）是拟紧的 (0\textsubscript{III}, 9.1.1)，则称该态射为\emph{拟紧}态射；
换言之，每个逆像 $f^{-1}(U)$ 都是拟紧的，其中 $U$ 是 $Y$ 的任意拟紧开集
（参见 (I, 6.6.1)）。
\end{definition}

若 $\mathfrak B$ 是由仿射开集组成的 $Y$ 的一个拓扑基，则 $f$ 拟紧，当且
仅当对每个 $V\in\mathfrak B$，$f^{-1}(V)$ 都是\emph{有限多个仿射开集的并}。
例如，若 $Y$ 仿射且 $X$ 拟紧，则\emph{每个}态射 $f:X\to Y$ 都是拟紧的
(I, 6.6.1)。

若 $f:X\to Y$ 是拟紧态射，则显然对每个开集 $V$（属于 $Y$），$f$ 在
$f^{-1}(V)$ 上的限制是拟紧态射 $f^{-1}(V)\to V$。反过来，若 $(U_\alpha)$
是 $Y$ 的一个开覆盖，而态射 $f:X\to Y$ 的各限制
$f^{-1}(U_\alpha)\to U_\alpha$ 都是拟紧的，则 $f$ 也是拟紧的。因此，若

\oldpage[IV-1]{225}

$f:X\to Y$ 是 $S$-预概形之间的一个 $S$-态射，并且存在一个开覆盖
$(S_\lambda)$（属于 $S$），使得各态射
$g^{-1}(S_\lambda)\to h^{-1}(S_\lambda)$（它们是 $f$ 的限制，其中 $g$ 与
$h$ 是结构态射）
都是拟紧的，则 $f$ 是拟紧的。

\begin{proposition}[1.1.2]
\label{IV.1.1.2-fr}
\begin{enumerate}
  \item[(i)] 嵌入 $X\to Y$ 在下列任一情形下都是拟紧的：它是闭嵌入；
  或 $Y$ 的底空间局部 Noether；或 $X$ 的底空间是 Noether 空间。
  \item[(ii)] 两个拟紧态射的复合是拟紧的。
  \item[(iii)] 若 $f:X\to Y$ 是拟紧 $S$-态射，则
  $f_{(S')}:X_{(S')}\to Y_{(S')}$ 对基预概形的每个扩张 $g:S'\to S$
  都是拟紧的。
  \item[(iv)] 若 $f:X\to X'$ 与 $g:Y\to Y'$ 是两个拟紧 $S$-态射，则
  \[
    f\times_S g:X\times_S Y\to X'\times_S Y'
  \]
  是拟紧的。
  \item[(v)] 若复合 $g\circ f$（其中 $f:X\to Y$、$g:Y\to Z$）
  是拟紧的，并且 $g$ 是分离态射，或 $X$ 的底空间局部 Noether，则 $f$
  是拟紧的。
  \item[(vi)] 态射 $f$ 是拟紧的，当且仅当 $f_{\mathrm{red}}$ 是拟紧的。
\end{enumerate}
\end{proposition}

证明见 (I, 6.6.4)。还应注意，断言 (vi) 也是下述更一般命题的推论：

\begin{proposition}[1.1.3]
\label{IV.1.1.3-fr}
设 $f:X\to Y$、$g:Y\to Z$ 是两个态射。若 $g\circ f$ 是拟紧的且 $f$
满射，则 $g$ 是拟紧的。
\end{proposition}

\begin{proof}
事实上，若 $V$ 是 $Z$ 中的拟紧开集，则依假设
$f^{-1}(g^{-1}(V))$ 是拟紧的；又有
$g^{-1}(V)=f(f^{-1}(g^{-1}(V)))$，因为 $f$ 满射，故 $g^{-1}(V)$ 是拟紧的。
\end{proof}

\begin{corollary}[1.1.4]
\label{IV.1.1.4-fr}
设 $f:X\to Y$、$g:Y'\to Y$ 是两个态射，$X'=X\times_Y Y'$，
$f'=f_{(Y')}:X'\to Y'$；假设 $g$ 拟紧且满射（相应地，只假设满射）。
则要使 $f$ 拟紧（相应地，支配），只须 $f'$ 具有这一性质。
\end{corollary}

\begin{proof}
若 $g':X'\to X$ 是典范投影，则 $g'$ 是满态射 (I, 3.5.2)，并且
$f\circ g'=g\circ f'$；若 $f'$ 拟紧（相应地，支配），则
$g\circ f'$ 也具有同一性质，因为 $g$ 拟紧 (1.1.2)（相应地，满射）；既然
$f\circ g'$ 拟紧（相应地，支配）且 $g'$ 满射，就可推出 $f$ 拟紧
(1.1.3)（相应地，支配）。
\end{proof}

\oldpage[IV-1]{226}

为简便起见，我们把预概形 $X$ 的任一不可约分支的泛点称为 $X$ 的
\emph{极大点}；若 $X=\operatorname{Spec}(A)$ 是仿射的，则 $X$ 的极大点
就是 $A$ 的\emph{极小素理想}。

\begin{proposition}[1.1.5]
\label{IV.1.1.5-fr}
设 $f:X\to Y$ 是拟紧态射。下列条件等价：
\begin{enumerate}
  \item[a)] $f$ 是支配态射；
  \item[b)] 对每个极大点 $y$（属于 $Y$），都有
  $f^{-1}(y)\neq\varnothing$；
  \item[c)] 对每个极大点 $y$（属于 $Y$），$f^{-1}(y)$ 含有 $X$ 的一个极大点。
\end{enumerate}
\end{proposition}

\emph{a)} 与 \emph{c)} 的等价性已在 (I, 6.6.5) 中证明。显然
\emph{c)} 蕴含 \emph{b)}；另一方面，若 $X'$ 是 $X$ 的一个不可约分支，
且同 $f^{-1}(y)$ 相交，则 $f(X')$ 不可约并含有 $y$，因而包含在不可约分支
$Y'=\overline{\{y\}}$（属于 $Y$）中 (0\textsubscript{I}, 2.1.6)；若 $x$
是 $X'$ 的泛点，就必有 $f(x)=y$ (0\textsubscript{I}, 2.1.5)；故
\emph{b)} 蕴含 \emph{c)}。

\begin{proposition}[1.1.6]
\label{IV.1.1.6-fr}
设 $f':X'\to Y$、$f'':X''\to Y$ 是两个态射，$X$ 是预概形的和
$X'\coprod X''$，而 $f:X\to Y$ 是等于 $f'$ 于 $X'$、等于 $f''$ 于 $X''$
的态射。则 $f$ 拟紧，当且仅当 $f'$ 与 $f''$ 都拟紧。
\end{proposition}

这由定义立即得出。

\subsection{拟分离态射}
\label{subsection:IV.1.2-fr}

\begin{definition}[1.2.1]
\label{IV.1.2.1-fr}
设 $X$、$Y$ 是两个预概形。称态射 $f:X\to Y$ 是\emph{拟分离}的
（或称 $X$ 在 $Y$ 上\emph{拟分离}），如果对角态射
$\Delta_f=(1_X,1_X)_Y:X\to X\times_Y X$ 是拟紧的。若预概形 $X$ 在
$\operatorname{Spec}(\mathbf Z)$ 上拟分离，则称该预概形是\emph{拟分离}的。
\end{definition}

由定义，每个分离态射都是拟分离的，因为 $\Delta_f$ 是闭嵌入
(1.1.2, (i))；每个\emph{概形} $X$ 都是拟分离预概形，因为它在
$\operatorname{Spec}(\mathbf Z)$ 上分离 (I, 5.4.1)。

\begin{proposition}[1.2.2]
\label{IV.1.2.2-fr}
\begin{enumerate}
  \item[(i)] 预概形的每个单态射 $f:X\to Y$（特别地，每个嵌入）都是拟分离的。
  \item[(ii)] 若 $f:X\to Y$ 与 $g:Y\to Z$ 是两个拟分离态射，则
  $g\circ f:X\to Z$ 是拟分离的。
  \item[(iii)] 设 $X$、$Y$ 是两个 $S$-预概形，而 $f:X\to Y$ 是拟分离
  $S$-态射。则对每个基变换 $g:S'\to S$，态射
  $f_{(S')}:X_{(S')}\to Y_{(S')}$ 都是拟分离的。
  \item[(iv)] 若 $f:X\to Y$、$f':X'\to Y'$ 是两个拟分离 $S$-态射，则
  \[
    f\times_S f':X\times_S X'\to Y\times_S Y'
  \]
  是拟分离的。
  \item[(v)] 若 $f:X\to Y$、$g:Y\to Z$ 是两个态射，并且 $g\circ f$
  是拟分离的，则 $f$ 是拟分离的。
  \item[(vi)] 若 $f:X\to Y$ 是拟分离态射，则
  $f_{\mathrm{red}}:X_{\mathrm{red}}\to Y_{\mathrm{red}}$ 是拟分离的。
\end{enumerate}
\end{proposition}

\begin{proof}
由 (I, 5.5.12)，只须证明 (i)、(ii) 和 (iii)。

断言 (i) 是平凡的，因为每个单态射都是分离的 (I, 5.5.1)。为证明 (iii)，
立即可归约到 $Y=S$ 的情形 (I, 3.3.11)，而断言由
$\Delta_{f_{(S')}}=(\Delta_f)_{(S')}$ (I, 5.3.4) 与 (1.1.2, (iii)) 得出。
为证明 (ii)，考虑两个投影 $p$ 与 $q$，它们从 $X\times_Y X$ 到 $X$；若
$i=(p,q)_Z$，已知图表
\[
\xymatrix{
  X\times_Y X \ar[r]^-{i} \ar[d]_-{\pi} & X\times_Z X
    \ar[d]^-{f\times_Z f} \\
  Y \ar[r]_-{\Delta_g} & Y\times_Z Y
}
\]
\oldpage[IV-1]{227}
（其中 $\pi$ 是结构态射）交换，并把 $X\times_Y X$ 等同于两个
$(Y\times_Z Y)$-预概形 $Y$ 与 $X\times_Z X$ 的积 (I, 5.3.5)。若 $g$
拟分离，则 $\Delta_g$ 拟紧，故 $i$ 也拟紧 (1.1.2, (iii))；若此外 $f$
拟分离，则 $\Delta_f$ 拟紧，故 $i\circ\Delta_f$ 也拟紧 (1.1.2, (ii))；
而它等于 $\Delta_{g\circ f}$。
\end{proof}

\begin{corollary}[1.2.3]
\label{IV.1.2.3-fr}
\begin{enumerate}
  \item[(i)] 设 $X$ 是拟分离预概形，则每个态射 $f:X\to Y$ 都是拟分离的。
  \item[(ii)] 设 $Y$ 是拟分离预概形。态射 $f:X\to Y$ 拟分离，当且仅当
  预概形 $X$ 拟分离。
\end{enumerate}
\end{corollary}

把 (1.2.2, (ii) 和 (v)) 应用于 $X$、$Y$ 和
$\operatorname{Spec}(\mathbf Z)$，立即得到此结论。

\begin{proposition}[1.2.4]
\label{IV.1.2.4-fr}
设 $f:X\to Y$ 是态射，$g:Y\to Z$ 是拟分离态射。若 $g\circ f$ 拟紧，
则 $f$ 也拟紧。
\end{proposition}

\begin{proof}
已知 $f$ 是复合态射
$X\xrightarrow{\Gamma_f}X\times_ZY\xrightarrow{\mathrm{pr}_2}Y$
(I, 5.3.13)；另一方面，$\mathrm{pr}_2$ 等同于
$(g\circ f)\times_Z1_Y$ (I, 3.3.4)，而若 $g\circ f$ 拟紧，则
$\mathrm{pr}_2$ 也拟紧 (1.1.2, (iv))。最后，有交换图
\[
\xymatrix{
  X \ar[r]^-{\Gamma_f} \ar[d]_-{f} & X\times_ZY
    \ar[d]^-{f\times 1_Y} \\
  Y \ar[r]_-{\Delta_g} & Y\times_ZY
}
\tag{1.2.4.1}
\label{IV.1.2.4.1-fr}
\]
它把 $X$ 等同于两个 $(Y\times_ZY)$-预概形 $Y$ 与 $X\times_ZY$ 的积
(I, 5.3.7)。依假设，$\Delta_g$ 是拟紧态射，故 $\Gamma_f$ 也拟紧
(1.1.2, (iii))；于是结论由 (1.1.2, (ii)) 得出。
\end{proof}

\begin{proposition}[1.2.5]
\label{IV.1.2.5-fr}
设 $f:X\to Y$、$g:Y'\to Y$ 是两个态射，
$X'=X\times_Y Y'$，$f'=f_{(Y')}:X'\to Y'$。若 $g$ 拟紧且满射，并且
$f'$ 拟分离，则 $f$ 拟分离。
\end{proposition}

\begin{proof}
事实上，有 $(X'\times_{Y'}X')=(X\times_YX)_{(Y')}$ 以及
$\Delta_{f'}=(\Delta_f)_{(Y')}$ (I, 5.3.4)；由于投影
$X'\times_{Y'}X'\to X\times_YX$ 拟紧 (1.1.2, (iii)) 且满射
(I, 3.5.2)，依 (I, 3.3.11) 可应用 (1.1.4)；它表明，既然
$\Delta_{f'}$ 是拟紧态射，$\Delta_f$ 也拟紧。
\end{proof}

\begin{proposition}[1.2.6]
\label{IV.1.2.6-fr}
设 $f:X\to Y$ 是态射，$(U_\alpha)$ 是 $Y$ 的一个开覆盖，并且各诱导预概形
$U_\alpha$ 都拟分离。则 $f$ 拟分离，当且仅当由 $X$ 在各开集
$f^{-1}(U_\alpha)$ 上诱导的预概形都拟分离。
\end{proposition}

\begin{proof}
在 $X\times_YX$ 中，$U_\alpha$ 的逆像是
$X_\alpha\times_{U_\alpha}X_\alpha$，其中令
$X_\alpha=f^{-1}(U_\alpha)$；而态射
$X_\alpha\to X_\alpha\times_{U_\alpha}X_\alpha$，即 $\Delta_f$ 的限制，正是
$\Delta_{f_\alpha}$，这里用 $f_\alpha$ 表示态射
$X_\alpha\to U_\alpha$，即 $f$ 的限制。由定义 (1.2.1) 以及拟紧态射概念的局部性
(1.1.1)，$f$ 拟分离，当且仅当每个态射 $f_\alpha$ 都拟分离。但依假设，
态射 $U_\alpha\to\operatorname{Spec}(\mathbf Z)$ 拟分离，所以
$f_\alpha$ 拟分离等价于复合
$X_\alpha\xrightarrow{f_\alpha}U_\alpha\to\operatorname{Spec}(\mathbf Z)$
拟分离 (1.2.2, (ii) 和 (v))，命题由此得证。
\end{proof}

\oldpage[IV-1]{228}

因此，为检验一个态射是否拟分离，问题归结为给出预概形拟分离的判别法：

\begin{proposition}[1.2.7]
\label{IV.1.2.7-fr}
设 $X$ 是预概形，$(U_\alpha)$ 是由拟紧开集组成的 $X$ 的一个覆盖。下列性质等价：
\begin{enumerate}
  \item[a)] $X$ 拟分离。
  \item[b)] 对每个拟紧开集 $U$（属于 $X$），典范嵌入
  $U\to X$ 都拟紧（换言之，$U$ 在 $X$ 中逆紧）。
  \item[b')] $X$ 的任意两个拟紧开集之交都拟紧。
  \item[c)] 对每对指标 $(\alpha,\beta)$，交集
  $U_\alpha\cap U_\beta$ 都拟紧。
\end{enumerate}
\end{proposition}

\begin{proof}
性质 \emph{b)} 与 \emph{b')} 由拟紧态射的定义显然等价。由于拟紧开集是有限个
仿射开集的并，若两个拟紧开集 $U$、$V$ 都属于 $X$，则
$U\times_{\mathbf Z}V$ 是 $X\times_{\mathbf Z}X$ 的拟紧开集
(I, 3.2.7)，而它在 $\Delta_X$ 下的逆像是 $U\cap V$，故 \emph{a)}
蕴含 \emph{b')}。\emph{b')} 蕴含 \emph{c)} 是平凡的；最后，若
\emph{c)} 成立，则 $U_\alpha\times_{\mathbf Z}U_\beta$ 组成
$X\times_{\mathbf Z}X$ 的一个拟紧开覆盖；我们知道，为使态射
$\Delta_X:X\to X\times_{\mathbf Z}X$ 拟紧，只须它在这些开集上的逆像
$U_\alpha\cap U_\beta$，即 $U_\alpha\times_{\mathbf Z}U_\beta$ 在
$\Delta_X$ 下的逆像，都拟紧 (1.1.1)。故 \emph{c)} 蕴含 \emph{a)}。
\end{proof}

\begin{corollary}[1.2.8]
\label{IV.1.2.8-fr}
底空间局部 Noether 的每个预概形 $X$（例如，每个局部 Noether 预概形）
都是拟分离的；于是每个态射 $f:X\to Y$ 都拟分离。
\end{corollary}

\begin{proposition}[1.2.9]
\label{IV.1.2.9-fr}
设 $X'$、$X''$ 是两个预概形，$X=X'\coprod X''$ 是它们的和，
$f':X'\to Y$、$f'':X''\to Y$ 是两个态射，而 $f:X\to Y$ 是等于
$f'$ 于 $X'$、等于 $f''$ 于 $X''$ 的态射。则 $f$ 拟分离，当且仅当
$f'$ 与 $f''$ 都拟分离。
\end{proposition}

\begin{proof}
事实上，$X\times_YX$ 是四个预概形
$X'\times_YX'$、$X'\times_YX''$、$X''\times_YX'$ 和 $X''\times_YX''$
的和，而 $\Delta_f$ 是等于 $\Delta_{f'}$ 于 $X'$、等于
$\Delta_{f''}$ 于 $X''$ 的态射；因此命题立即由定义得出。
\end{proof}

\subsection{局部有限型态射}
\label{subsection:IV.1.3-fr}

\begin{env}[1.3.1]
\label{IV.1.3.1-fr}
回忆：若 $A$ 是环，则称（交换）$A$-代数 $B$ 是\emph{有限型}的，如果它由
$B$ 的有限多个元素生成；等价地，如果它同多项式代数
$A[T_1,\ldots,T_n]$ 对该代数的一个理想的商同构。显然，对任意（交换）
$A$-代数 $A'$，$B\otimes_AA'$ 都是有限型 $A'$-代数。若 $B$ 是有限型
$A$-代数，而 $C$ 是有限型 $B$-代数，则 $C$ 是有限型 $A$-代数；事实上，
若 $C$ 是 $B[T_1,\ldots,T_n]$ 的商，而 $B$ 是
$A[S_1,\ldots,S_m]$ 的商，则
$B[T_1,\ldots,T_n]=B\otimes_AA[T_1,\ldots,T_n]$ 是
$A[S_1,\ldots,S_m,T_1,\ldots,T_n]$ 的商，因而 $C$ 也是它的商。
\end{env}

\begin{definition}[1.3.2]
\label{IV.1.3.2-fr}
设 $f:X\to Y$ 是预概形的态射，$x$ 是 $X$ 的一点，$y=f(x)$。称 $f$
\emph{在 $x$ 处为有限型}，如果存在一个仿射开邻域 $V$（属于 $y$）和
一个仿射开邻域 $U$（属于 $x$），使得 $f(U)\subset V$，并且环
$A(U)$ 是有限型 $A(V)$-代数。如果 $f$ 在 $X$ 的每一点处都为有限型，就称
该态射是\emph{局部有限型}的。
\end{definition}

\oldpage[IV-1]{229}

\begin{proposition}[1.3.3]
\label{IV.1.3.3-fr}
若 $Y$ 是局部 Noether 预概形，而 $f:X\to Y$ 是局部有限型态射，则
$X$ 是局部 Noether 的。
\end{proposition}

证明见 (I, 6.3.7)。

\begin{proposition}[1.3.4]
\label{IV.1.3.4-fr}
\begin{enumerate}
  \item[(i)] 每个局部嵌入都是局部有限型的。
  \item[(ii)] 若两个态射 $f:X\to Y$、$g:Y\to Z$ 都局部有限型，则
  $g\circ f$ 也局部有限型。
  \item[(iii)] 若 $f:X\to Y$ 是局部有限型 $S$-态射，则
  $f_{(S')}:X_{(S')}\to Y_{(S')}$ 对基预概形的每个扩张 $S'\to S$
  都局部有限型。
  \item[(iv)] 若 $f:X\to X'$ 与 $g:Y\to Y'$ 是两个局部有限型
  $S$-态射，则 $f\times_Sg$ 局部有限型。
  \item[(v)] 若两个态射的复合 $g\circ f$ 局部有限型，则 $f$ 局部有限型。
  \item[(vi)] 若态射 $f$ 局部有限型，则 $f_{\mathrm{red}}$ 也局部有限型。
\end{enumerate}
\end{proposition}

证明见 (I, 6.6.6)。

\begin{corollary}[1.3.5]
\label{IV.1.3.5-fr}
设 $f:X\to Y$ 是局部有限型态射。对每个态射 $Y'\to Y$，只要 $Y'$
局部 Noether，$X\times_YY'$ 就局部 Noether。
\end{corollary}

\begin{proof}
事实上，由 (1.3.4, (iii))，$f_{(Y')}:X\times_YY'\to Y'$ 局部有限型，
只须应用 (1.3.3)。
\end{proof}

\begin{proposition}[1.3.6]
\label{IV.1.3.6-fr}
设 $\varphi:A\to B$ 是环同态。相应态射
$f:\operatorname{Spec}(B)\to\operatorname{Spec}(A)$ 局部有限型，当且仅当
$B$ 是有限型 $A$-代数。
\end{proposition}

证明见 (I, 6.3.3)，其中要考虑到 $f$ 拟紧以及 (I, 6.6.3)。

\begin{proposition}[1.3.7]
\label{IV.1.3.7-fr}
局部有限型态射 $f:X\to Y$ 满射，当且仅当对每个代数闭域 $\Omega$，
映射 $X(\Omega)\to Y(\Omega)$（它对应于 $f$）(I, 3.4.1) 都满射。
\end{proposition}

证明见 (I, 6.3.10)。

\begin{env}[1.3.8]
\label{IV.1.3.8-fr}
设 $A$ 是环。称一个 $A$-代数 $B$ 是\emph{本质有限型}的，如果 $B$ 同
一个 $A$-代数 $A$-同构，而后者形如 $S^{-1}C$，其中 $C$ 是
\emph{有限型} $A$-代数，而 $S$ 是 $C$ 的乘法子集。
\end{env}

\par\medskip\noindent
\begin{proposition}[1.3.9]
\label{IV.1.3.9-fr}
\begin{enumerate}
  \item[(i)] 若 $B$ 是本质有限型 $A$-代数，而 $C$ 是本质有限型
  $B$-代数，则 $C$ 是本质有限型 $A$-代数。
  \item[(ii)] 设 $B$ 是本质有限型 $A$-代数，$A'$ 是一个 $A$-代数；则
  $B'=B\otimes_AA'$ 是本质有限型 $A'$-代数。
\end{enumerate}
\end{proposition}

\begin{proof}
\begin{enumerate}
  \item[(i)] 设 $B=S'^{-1}B'$ 且 $C=T'^{-1}C'$，其中 $B'$（相应地，
  $C'$）是有限型 $A$-代数（相应地，有限型 $B$-代数），而 $S'$（相应地，
  $T'$）是 $B'$（相应地，$C'$）的乘法子集；于是 $C'$ 形如
  $S'^{-1}C''$，其中 $C''$ 是有限型 $B'$-代数，故 $C$ 也是 $C''$
  的分式环；由于 $C''$ 是有限型 $A$-代数，这就证明了断言。
  \item[(ii)] 若 $B$ 是一个有限型 $A$-代数 $C$ 的分式环，则 $B'$ 是
  有限型 $A'$-代数 $C\otimes_AA'$ 的分式环，故得 (ii)。
\end{enumerate}
\end{proof}

\oldpage[IV-1]{230}

\begin{env}[1.3.10]
\label{IV.1.3.10-fr}
若 $B$ 是本质有限型 $A$-\emph{局部}代数，则 $B$ 形如
$C_{\mathfrak q}$，其中 $C$ 是有限型 $A$-代数，而 $\mathfrak q$ 是 $C$
的素理想（Bourbaki，\emph{Alg. comm.}，chap. II，\S~5，prop. 11）。设
$\mathfrak p$ 是在 $A$ 中理想 $\mathfrak q$ 的逆像；若令
$S=A-\mathfrak p$，则 $C_{\mathfrak q}$ 也是 $S^{-1}C$ 在一个素理想处的
局部环；由于 $S^{-1}C$ 是有限型
$A_{\mathfrak p}=S^{-1}A$-代数，可见 $B$ 也是本质有限型
$A_{\mathfrak p}$-代数，而且同态 $A_{\mathfrak p}\to B$ 是\emph{局部}同态。
\end{env}

\begin{proposition}[1.3.11]
\label{IV.1.3.11-fr}
若 $B$ 是本质有限型 $A$-局部代数，则 $B$ 同一个 $A$-代数
$A$-同构；该代数是一个 $A$-代数的商，后者形如 $C_{\mathfrak q}$，其中 $C$ 是多项式代数
$A[T_1,\ldots,T_n]$，而 $\mathfrak q$ 是 $C$ 的素理想。
\end{proposition}

\begin{proof}
事实上，按定义，$B$ 同 $C'_{\mathfrak q'}$ 同构，其中 $C'$ 是有限型
$A$-代数，而 $\mathfrak q'$ 是 $C'$ 的素理想；但
$C'=C/\mathfrak b$，其中 $C=A[T_1,\ldots,T_n]$，而 $\mathfrak b$ 是 $C$
的理想；故 $\mathfrak q'=\mathfrak q/\mathfrak b$，其中 $\mathfrak q$ 是
$C$ 的素理想，并且 $C'_{\mathfrak q'}$ 同
$C_{\mathfrak q}/\mathfrak bC_{\mathfrak q}$ 同构。
\end{proof}

\subsection{局部有限表示态射}
\label{subsection:IV.1.4-fr}

\begin{env}[1.4.1]
\label{IV.1.4.1-fr}
给定环 $A$，称（交换）$A$-代数 $B$ 是\emph{有限表示}的，如果它同
多项式代数对该代数的一个理想所成的商
$A[T_1,\ldots,T_n]/\mathfrak a$（作为 $A$-代数）同构，其中
$\mathfrak a$ 是 $A[T_1,\ldots,T_n]$ 的一个\emph{有限型}理想。对任意
（交换）$A$-代数 $A'$，$B\otimes_AA'$ 都是有限表示 $A'$-代数；事实上，
它同 $A'[T_1,\ldots,T_n]$ 对 $\mathfrak a\otimes_AA'$ 在该环中的典范像
所成的商同构，而该像显然是有限型 $A'[T_1,\ldots,T_n]$-模。若 $B$ 是
有限表示 $A$-代数，而 $C$ 是有限表示 $B$-代数，则 $C$ 是有限表示
$A$-代数。事实上，设
$B=A[S_1,\ldots,S_m]/\mathfrak a$，
$C=B[T_1,\ldots,T_n]/\mathfrak b$，其中 $\mathfrak a$（相应地，
$\mathfrak b$）是有限型理想，属于 $A[S_1,\ldots,S_m]$（相应地，
$B[T_1,\ldots,T_n]$）；环 $B[T_1,\ldots,T_n]$ 同
$A[S_1,\ldots,S_m,T_1,\ldots,T_n]/\mathfrak a'$ 同构，其中
$\mathfrak a'$ 是 $\mathfrak a\otimes_AA[T_1,\ldots,T_n]$ 的典范像，
因而是 $A[S_1,\ldots,S_m,T_1,\ldots,T_n]$ 的有限型理想；
理想 $\mathfrak b$（属于 $B[T_1,\ldots,T_n]$）形如
$\mathfrak b'/\mathfrak a'$，其中 $\mathfrak b'$ 是
$A[S_1,\ldots,S_m,T_1,\ldots,T_n]$ 的理想；由于 $\mathfrak a'$ 和
$\mathfrak b'/\mathfrak a'$ 都是有限型
$A[S_1,\ldots,S_m,T_1,\ldots,T_n]$-模，$\mathfrak b'$ 也如此；因此
$C$ 同 $A[S_1,\ldots,S_m,T_1,\ldots,T_n]/\mathfrak b'$ 同构，确实是
有限表示 $A$-代数。最后注意，若 $A$ 是\emph{Noether} 环，则
$A[T_1,\ldots,T_n]$ 也是 Noether 环，因而每个有限型 $A$-代数都是
有限表示的。
\end{env}

\begin{definition}[1.4.2]
\label{IV.1.4.2-fr}
设 $f:X\to Y$ 是预概形的态射，$x$ 是 $X$ 的一点，$y=f(x)$。称 $f$
\emph{在 $x$ 处为有限表示}，如果存在一个仿射开邻域 $V$（属于 $y$）和
一个仿射开邻域 $U$（属于 $x$），使得 $f(U)\subset V$，并且环
$A(U)$ 是有限表示 $A(V)$-代数。如果 $f$ 在 $X$ 的每一点处都为有限表示，
就称该态射是\emph{局部有限表示}的。
\end{definition}

若 $Y$ 局部 Noether，则 $f$ 局部有限型与局部有限表示是同一回事。

\begin{proposition}[1.4.3]
\label{IV.1.4.3-fr}
\begin{enumerate}
  \item[(i)] 每个局部同构都是局部有限表示的。
  \item[(ii)] 若两个态射 $f:X\to Y$、$g:Y\to Z$ 都局部有限表示，则
  $g\circ f$ 也局部有限表示。

\oldpage[IV-1]{231}

  \item[(iii)] 若 $f:X\to Y$ 是局部有限表示 $S$-态射，则
  $f_{(S')}:X_{(S')}\to Y_{(S')}$ 对基预概形的每个扩张 $S'\to S$
  都局部有限表示。
  \item[(iv)] 若 $f:X\to X'$ 与 $g:Y\to Y'$ 是两个局部有限表示
  $S$-态射，则 $f\times_Sg$ 局部有限表示。
  \item[(v)] 若两个态射的复合 $g\circ f$ 局部有限表示，并且 $g$
  局部有限型，则 $f$ 局部有限表示。
\end{enumerate}
\end{proposition}

\begin{proof}
断言 (i) 是平凡的。为证明 (iii)，可只考虑 $S=Y$ 的情形 (I, 3.3.11)；
设 $x'$ 是 $X_{(S')}$ 的一点，$s'$ 与 $x$ 分别是它在 $S'$ 与 $X$ 上的
投影，而 $s$ 是 $s'$ 与 $x$ 在 $S$ 上的共同投影。依假设，存在仿射开邻域
$V$（相应地，$U$），它是 $s$（相应地，$x$）的邻域，使 $U$
的像包含在 $V$ 中，并且 $A(U)$ 是有限表示 $A(V)$-代数；设 $V'$ 是
$s'$ 的一个仿射开邻域，其在 $S$ 中的像包含在 $V$ 中；则
$U\times_VV'$ 是 $x'$ 的一个仿射开邻域，它在 $S'$ 中的像包含在 $V'$
中，并且它的环 $A(U)\otimes_{A(V)}A(V')$ 是有限表示 $A(V')$-代数
(1.4.1)。

为证明 (ii)，考虑一点 $x\in X$，并令 $y=f(x)$、$z=g(f(x))$。依假设，
存在仿射开邻域 $W=\operatorname{Spec}(C)$（相应地，
$V=\operatorname{Spec}(B)$），它是 $z$（相应地，$y$）的邻域，使
$g(V)\subset W$，并且
$B$ 是有限表示 $C$-代数。另一方面，存在仿射开邻域
$V'=\operatorname{Spec}(B')$（相应地，$U=\operatorname{Spec}(A)$），
它是 $y$（相应地，$x$）的邻域，使 $f(U)\subset V'$，并且 $A$ 是有限表示
$B'$-代数。取 $s\in B$，使开集 $D(s)=\operatorname{Spec}(B_s)$（属于
$V$）是 $y$ 的一个包含在 $V'$ 中的邻域；并令
$U'=f^{-1}(D(s))\subset U$。有
$U'=\operatorname{Spec}(A\otimes_{B'}B_s)$，而 $A\otimes_{B'}B_s$ 是有限表示
$B_s$-代数 (1.4.1)；另一方面，$B_s=B[T]/(1-sT)$ 按定义是有限表示
$B$-代数；因此由 (1.4.1)，$A\otimes_{B'}B_s$ 是有限表示 $C$-代数，
这就证明了 (ii)。

已知 (iv) 由 (i)、(ii) 与 (iii) 得出 (I, 3.5.1)。为证明 (v)，考虑交换图
\[
\xymatrix{
X \ar[r]^{\Gamma_f} \ar[d] & X\times_ZY \ar[d]^{f\times 1} \\
Y \ar[r]_{\Delta_g} & Y\times_ZY
}
\]
它把 $X$ 等同于两个 $(Y\times_ZY)$-预概形 $Y$ 与 $X\times_ZY$ 的积
(I, 5.3.7，其中 $\Delta_g$ 是对角态射)。若已知 $\Delta_g$ 局部有限表示，
则由 (iii)，$\Gamma_f$ 也局部有限表示。另一方面，$f$ 有分解
$X\xrightarrow{\Gamma_f}X\times_ZY\xrightarrow{p_2}Y$ (I, 5.3.13)，其中投影
$p_2$ 等于 $(g\circ f)\times_Z1_Y$；因为假设 $g\circ f$ 局部有限表示，
由 (iv)，$p_2$ 也局部有限表示；最后由 (ii)，$f$ 也局部有限表示。
因此可归结为证明下面的
\end{proof}

\begin{corollary}[1.4.3.1]
\label{IV.1.4.3.1-fr}
设 $g:Y\to Z$ 是局部有限型态射；则对角态射
$\Delta_g:Y\to Y\times_ZY$ 局部有限表示。
\end{corollary}

\oldpage[IV-1]{232}

\begin{proof}
可只考虑 $Z=\operatorname{Spec}(A)$、$Y=\operatorname{Spec}(B)$ 的情形，
其中 $B$ 是有限型 $A$-代数；此时 $\Delta_g$ 对应于\emph{增广同态}
$\delta:B\otimes_AB\to B$，满足 $\delta(x\otimes y)=xy$。考虑定义 (1.4.2)，
只须证明核 $\mathfrak J$（即 $\delta$ 的核）是\emph{有限型}理想；这由
(0, 20.4.4) 得出。
\end{proof}

\begin{proposition}[1.4.4]
\label{IV.1.4.4-fr}
设 $A$ 是环，$B$ 是 $A$-代数，$B'$ 是多项式代数
$A[T_1,\ldots,T_n]$，而 $u:B'\to B$ 是满的 $A$-代数同态。则 $B$ 是
有限表示 $A$-代数，当且仅当核 $\mathfrak J$（属于 $u$）是 $B'$ 的有限型理想。
\end{proposition}

\begin{proof}
该条件由定义 (1.4.1) 显然充分。为说明其必要性，注意态射
$g:\operatorname{Spec}(B')\to\operatorname{Spec}(A)$ 局部有限型；若 $B$
是有限表示 $A$-代数，则态射
$f:\operatorname{Spec}(B)\to\operatorname{Spec}(B')$（对应于 $u$）使
$g\circ f:\operatorname{Spec}(B)\to\operatorname{Spec}(A)$ 局部有限表示，
故由 (1.4.3, (v))，$f$ 也局部有限表示。现在，为证明理想 $\mathfrak J$
有限型，只须证明 $\widetilde{\mathfrak J}$ 是有限型
$\widetilde{B'}$-模 (II, 6.1.4.1)。回到定义 (1.4.2)，最终归结为证明下面的
\end{proof}

\begin{lemma}[1.4.4.1]
\label{IV.1.4.4.1-fr}
设 $\mathfrak a$ 是环 $C$ 的理想，$s$ 是 $C$ 的一个元素，并且
$C_s/\mathfrak a_s$ 是有限表示 $C$-代数；则 $\mathfrak a_s$ 是环
$C_s$ 的有限型理想。
\end{lemma}

\begin{proof}
依假设，存在一个 $C$-多项式代数
$C'=C[T_1,\ldots,T_n]$ 以及一个满的 $C$-同态
$v:C'\to C_s/\mathfrak a_s$，其核 $\mathfrak b$ 有限型。设
$p:C_s\to C_s/\mathfrak a_s$ 是典范同态；对每个 $i$，存在
$t_i\in C$，使得 $p(t_i/s^k)=v(T_i)$（可假设指数 $k$ 对所有指标
$i$ 相同）。于是考虑 $C_s$-代数 $D=C_s[T_1,\ldots,T_n]$，以及
$C_s$-代数同态 $w:D\to C_s$，它满足 $w(T_i)=t_i/s^k$；$w$ 显然满射，
故 $v'=p\circ w:D\to C_s/\mathfrak a_s$ 也满射。另一方面，$D$ 的每个
多项式都可写成 $P/s^m$，其中 $P\in C'=C[T_1,\ldots,T_n]$，并且
$v'(P/s^m)=(1/s^m)v(P)$；由于 $1/s$ 在 $C_s$ 中可逆，关系
$v'(P/s^m)=0$ 在 $C_s/\mathfrak a_s$ 中等价于 $v(P)=0$，因此核
$\mathfrak b'$（属于 $v'$）由 $\mathfrak b$ 在 $D$ 中的典范像生成，从而
\emph{更不待言}是 $D$ 的有限型理想。但
$\mathfrak b'=v'^{-1}(0)=w^{-1}(p^{-1}(0))=w^{-1}(\mathfrak a_s)$；因为
$w$ 满射，所以 $\mathfrak a_s=w(\mathfrak b')$，这证明了 $\mathfrak a_s$
是 $C_s$ 的有限型理想。
\end{proof}

\begin{corollary}[1.4.5]
\label{IV.1.4.5-fr}
设 $X$、$Y$ 是两个预概形，$j:X\to Y$ 是浸入，$U$ 是 $Y$ 的一个开集，
使得 $j(X)$ 在 $U$ 中闭，并设 $\mathcal J$ 是 $\mathcal O_U$ 的拟凝聚理想，
它定义 $Y$ 的与 $j$ 相联系的闭子预概形。则 $j$ 局部有限表示，当且仅当
$\mathcal J$ 是有限型 $\mathcal O_U$-模。
\end{corollary}

\par\medskip\noindent
\begin{proposition}[1.4.6]
\label{IV.1.4.6-fr}
设 $\varphi:A\to B$ 是环同态。相应态射
$f:\operatorname{Spec}(B)\to\operatorname{Spec}(A)$ 局部有限表示，
当且仅当 $B$ 是有限表示 $A$-代数。
\end{proposition}

\begin{proof}
该条件显然充分，尚须证明其必要性。若 $f$ 局部有限表示，则已由 (1.3.6)
知道 $B$ 是有限型 $A$-代数，故存在满同态
$u:B'=A[T_1,\ldots,T_n]\to B$，它是 $A$-代数同态；由 (1.4.3, (v))，浸入
$j:\operatorname{Spec}(B)\to\operatorname{Spec}(B')$ 局部

\oldpage[IV-1]{233}

有限表示；因此由 (1.4.5)，若 $\mathfrak J$ 是 $u$ 的核，则
$\widetilde{B'}$-模 $\widetilde{\mathfrak J}$ 是有限型的，从而由
(II, 6.1.4.1)，$\mathfrak J$ 是有限型理想。
\end{proof}

\begin{proposition}[1.4.7]
\label{IV.1.4.7-fr}
设 $\rho:A\to B$ 是环同态，并使 $B$ 成为有限 $A$-代数。则 $B$ 是
有限表示 $A$-代数，当且仅当 $B$ 是有限表示 $A$-模。
\end{proposition}

我们将使用下面的引理。

\begin{lemma}[1.4.7.1]
\label{IV.1.4.7.1-fr}
若 $B$ 是有限 $A$-代数，则存在满的 $A$-代数同态 $u:B'\to B$，其中
$B'$ 是有限且有限表示的 $A$-代数，并且是自由 $A$-模。
\end{lemma}

\begin{proof}
事实上，有一个有限生成元系 $(a_i)_{1\leq i\leq m}$，它生成
$A$-代数 $B$，并使每个 $a_i$ 都满足关系 $F_i(a_i)=0$，其中
$F_i\in A[X]$ 是次数 $>0$ 的首一多项式；因此商 $A$-代数
$B'_i=A[X]/(F_i)$ 是有限秩自由 $A$-模；以 $c_i$ 表示 $X$ 在 $B'_i$
中的类。存在 $A$-代数同态 $u_i:B'_i\to B$，使 $u_i(c_i)=a_i$；此时只须
取 $B'$ 为张量积 $B'_1\otimes_AB'_2\otimes_A\cdots\otimes_AB'_m$，并考虑
同态 $u_1\otimes u_2\otimes\cdots\otimes u_m:B'\to B$；按构造它是满的。
此外，若 $B''=A[T_1,\ldots,T_m]$，并以 $b'_i$ 表示 $c_i$ 在 $B'$ 中的
典范像，则 $A$-代数同态 $v:B''\to B'$（满足 $v(T_i)=b'_i$，
$1\leq i\leq m$）是满的，而其核 $\mathfrak b'$ 是 $B''$ 的由
$F_i(T_i)$ 生成的理想，因而是有限型的；这证明了 $B'$ 是有限表示
$A$-代数。
\end{proof}

\begin{proof}
上述引理既已证明，沿用同样的记号，并令 $w=v\circ u:B''\to B$。若
$\mathfrak b$（相应地，$\mathfrak a$）是 $w$（相应地，$u$）的核，
则有 $\mathfrak a=v(\mathfrak b)$，因为 $v$ 满射。由 (1.4.4)，若 $B$ 是
有限表示 $A$-代数，则 $\mathfrak b$ 是 $B''$ 的有限型理想，故
$\mathfrak a$ 是 $B'$ 的有限型理想，因而是有限型 $A$-模，因为 $B'$ 是
有限 $A$-代数。由于 $B'$ 是自由 $A$-模，$B$ 是有限表示
$A$-模。反之，若 $B'$ 是有限表示 $A$-模，则 $\mathfrak a$ 是有限型
$A$-模（Bourbaki，\emph{Alg. comm.}，chap. I，\S~2，
n\textsuperscript{o}~8，lemme 9），从而\emph{更不待言}是 $B'$ 的有限型
理想；因此按定义，$B$ 是有限表示 $B'$-代数；又因 $B'$ 是有限表示
$A$-代数，所以 $B$ 是有限表示 $A$-代数。
\end{proof}

\subsection{有限型态射}
\label{subsection:IV.1.5-fr}

\begin{proposition}[1.5.1]
\label{IV.1.5.1-fr}
设 $f:X\to Y$ 是态射，$(U_\alpha)$ 是由 $Y$ 的仿射开集组成的覆盖。
下列条件等价：
\begin{enumerate}
  \item[a)] $f$ 局部有限型且拟紧。
  \item[b)] 对每个 $\alpha$，$f^{-1}(U_\alpha)$ 是有限多个仿射开集
  $V_{\alpha i}$ 的并，使得每个环 $A(V_{\alpha i})$ 都是有限型
  $A(U_\alpha)$-代数。
  \item[c)] 对每个仿射开集 $U$（属于 $Y$），$f^{-1}(U)$ 是有限多个仿射开集
  $V_j$ 的并，使得各环 $A(V_j)$ 都是有限型 $A(U)$-代数。
\end{enumerate}
\end{proposition}

证明见 (I, 6.3.2 与 6.6.3)。

\begin{definition}[1.5.2]
\label{IV.1.5.2-fr}
若态射 $f$ 满足命题 (1.5.1) 的各个等价条件，就称该态射是
\emph{有限型}的。
\end{definition}

\oldpage[IV-1]{234}

\begin{proposition}[1.5.3]
\label{IV.1.5.3-fr}
设 $f:X\to Y$ 是有限型态射；若 $Y$ 是 Noether 的，则 $X$ 也是
Noether 的。
\end{proposition}

证明见 (I, 6.3.7)。

\begin{proposition}[1.5.4]
\label{IV.1.5.4-fr}
\begin{enumerate}
  \item[(i)] 设 $j:X\to Y$ 是浸入态射。若 $j$ 拟紧（当 $j$ 是闭浸入，
  或 $X$ 的底空间是 Noether 的，或 $Y$ 的底空间局部 Noether 时，
  这总是成立），则 $j$ 是有限型的。
  \item[(ii)] 两个有限型态射的复合是有限型的。
  \item[(iii)] 若 $f:X\to Y$ 是有限型 $S$-态射，则
  $f_{(S')}:X_{(S')}\to Y_{(S')}$ 对基预概形的每个扩张 $S'\to S$
  都是有限型的。
  \item[(iv)] 若 $f:X\to X'$ 与 $g:Y\to Y'$ 是两个有限型 $S$-态射，
  则 $f\times_Sg$ 是有限型的。
  \item[(v)] 若两个态射的复合 $g\circ f$ 是有限型的，并且 $g$ 拟分离，
  或者 $X$ 是 Noether 的，则 $f$ 是有限型的。
  \item[(vi)] 若态射 $f$ 是有限型的，则 $f_{\mathrm{red}}$ 也是有限型的。
\end{enumerate}
\end{proposition}

除 (v) 中 $X$ 为 Noether 的情形（已在 (I, 6.3.6) 中证明）外，
这些断言都由定义以及 (1.1.2)、(1.2.4) 和 (1.3.4) 立即得出。

\begin{corollary}[1.5.5]
\label{IV.1.5.5-fr}
设 $f:X\to Y$ 是有限型态射。对每个态射 $Y'\to Y$，若 $Y'$ 是
Noether 预概形，则 $X\times_YY'$ 是 Noether 的。
\end{corollary}

这由 (1.5.4, (iii)) 和 (1.5.3) 得出。

\begin{corollary}[1.5.6]
\label{IV.1.5.6-fr}
设 $X$ 是局部 Noether 预概形 $S$ 上的有限型预概形。则每个 $S$-态射
$f:X\to Y$ 都是有限型的。
\end{corollary}

证明见 (I, 6.3.9)。

\begin{proposition}[1.5.7]
\label{IV.1.5.7-fr}
设 $\varphi:A\to B$ 是环同态。相应态射
$f:\operatorname{Spec}(B)\to\operatorname{Spec}(A)$ 是有限型的，
当且仅当 $\varphi$ 使 $B$ 成为有限型 $A$-代数。
\end{proposition}

证明见 (I, 6.3.3)。

\subsection{有限表示态射}
\label{subsection:IV.1.6-fr}

\begin{definition}[1.6.1]
\label{IV.1.6.1-fr}
设 $X$、$Y$ 是两个预概形。若态射 $f:X\to Y$ 满足下列三个条件，
就称它是\emph{有限表示}的：
\begin{enumerate}
  \item[(i)] $f$ 局部有限表示；
  \item[(ii)] $f$ 拟紧（当 (i) 成立时，这等价于说 $f$ 是有限型的
  (1.5.2)）；
  \item[(iii)] $f$ 拟分离。
\end{enumerate}
此时也称 $X$ \emph{在 $Y$ 上有限表示}，或称为一个
\emph{有限表示 $Y$-预概形}。
\end{definition}

若 $f$ 是分离的，或 $X$ 局部 Noether，则条件 (iii) 自动成立 (1.2.8)；
当 $Y$ 局部 Noether 时，说 $f$ \emph{有限表示}与说 $f$
\emph{有限型}是同一回事（后一条件蕴含 $X$ 局部 Noether (1.3.3)）。

\oldpage[IV-1]{235}

\begin{proposition}[1.6.2]
\label{IV.1.6.2-fr}
\begin{enumerate}
  \item[(i)] 每个拟紧且局部有限表示的浸入（特别地，每个拟紧开浸入）
  都是有限表示的。
  \item[(ii)] 两个有限表示态射的复合是有限表示的。
  \item[(iii)] 设 $X$、$Y$ 是两个 $S$-预概形，$f:X\to Y$ 是有限表示
  $S$-态射。则对每个态射 $S'\to S$，态射
  $f_{(S')}:X_{(S')}\to Y_{(S')}$ 都是有限表示的。
  \item[(iv)] 设 $f:X\to Y$、$f':X'\to Y'$ 是两个有限表示
  $S$-态射；则
  $f\times_Sf':X\times_SX'\to Y\times_SY'$ 是有限表示的。
  \item[(v)] 若两个态射的复合 $g\circ f$ 是有限表示的，并且 $g$
  拟分离且局部有限表示，则 $f$ 是有限表示的。
\end{enumerate}
\end{proposition}

这立即由 (1.1.2)、(1.2.2)、(1.2.4) 和 (1.4.3) 得出。

特别地，由 (1.6.2, (iii)) 可知，若 $f$ 是有限表示态射，$U$ 是 $Y$
的开集，则限制 $f^{-1}(U)\to U$（属于 $f$）是有限表示态射。反之，设
$(U_\alpha)$ 是由 $Y$ 的仿射开集组成的覆盖，并假设各个限制
$f^{-1}(U_\alpha)\to U_\alpha$（属于 $f$）都是有限表示态射；那么 $f$
是有限表示的，
因为 $f$ 显然局部有限表示，由 (1.1.1) 它拟紧，并由 (1.2.6) 它拟分离。
因此可以说，态射 $f:X\to Y$ 为有限表示这一性质是\emph{在 $Y$ 上局部}的。

若 $X$ 是拟分离预概形，则每个态射 $f:X\to Y$ 都拟分离 (1.2.3)。
因此，若 $f$ 拟紧且局部有限表示，则 $f$ 是有限表示的。

特别地：

\begin{corollary}[1.6.3]
\label{IV.1.6.3-fr}
设 $\varphi:A\to B$ 是环同态。相应态射
$f:\operatorname{Spec}(B)\to\operatorname{Spec}(A)$ 是有限表示的，
当且仅当 $\varphi$ 使 $B$ 成为有限表示 $A$-代数。
\end{corollary}

该条件的必要性确实由 (1.4.6) 得出。

\begin{remark}[1.6.4]
\label{IV.1.6.4-fr}
在定义 (1.6.1) 中，条件 (iii) 不是其余两个条件的推论。例如，设 $Y$
是一个底空间非 Noether 的仿射概形，并设 $U$ 是 $Y$ 中一个非拟紧开集。
设 $X$ 是把两个预概形 $Y_1$、$Y_2$（均与 $Y$ 同构）沿着开集
$U_1$、$U_2$（它们对应于 $U$）粘合所得的预概形（0$_{\mathrm I}$,
4.1.7），
使得 $X$ 是两个分别与 $Y_1$ 和 $Y_2$ 同构的开集之并（并将这些开集
分别同它们等同），而且 $Y_1\cap Y_2=U$。此外，有一个态射 $f:X\to Y$，
它在 $Y_i$ 上与给定的同构 $Y_i\to Y$ 相同（$i=1,2$）。显然，$f$
作为局部同构而满足 (1.6.1) 的条件 (i)；它也满足 (ii)，因为 $Y$ 的
拟紧开集的逆像是它在 $Y_1$ 与 $Y_2$ 中的像之并；但是由于
$Y_1\cap Y_2$ 非拟紧，$f$ 不是拟分离的 (1.2.7)。
\end{remark}

\begin{proposition}[1.6.5]
\label{IV.1.6.5-fr}
设 $X'$、$X''$ 是两个预概形，$X=X'\amalg X''$ 是它们的和，
$f':X'\to Y$、$f'':X''\to Y$ 是两个态射，$f:X\to Y$ 是这样一个态射：
它同 $f'$ 在 $X'$ 上相同，并同 $f''$ 在 $X''$ 上相同。则 $f$ 是有限表示的，
当且仅当 $f'$ 和 $f''$ 都是有限表示的。
\end{proposition}

\par\medskip\noindent
\begin{proof}
只须证明：$f$ 具有定义

\oldpage[IV-1]{236}

(1.6.1) 的三个性质中的任意一个，当且仅当 $f'$ 与 $f''$ 都具有该性质。
对局部有限表示这一性质，这是显然的，因为它在 $X$ 上局部；拟紧性质已在
(1.1.6) 中考察，拟分离性质则已在 (1.2.9) 中考察。
\end{proof}

\subsection{既有结果的改进}
\label{subsection:IV.1.7-fr}

本节将列出前几章中已经证明的一些命题；借助上面引入的新的有限性条件，
可以改进这些命题的陈述。

\begin{env}[1.7.1]
\label{IV.1.7.1-fr}
在 (I, 6.4.2、6.4.3 和 6.4.9) 的陈述中，可以不再假设 $X$ 和 $Y$
在域 $K$ 上有限型，而只假设它们在 $K$ 上\emph{局部有限型}；对于
(I, 6.4.2)，只须注意：预概形 $X$ 中若一个点在 $X$ 的每个仿射开集
中都是闭的，那么它在 $X$ 中也是闭的；并且，在 $K$ 上局部有限型的
仿射概形\emph{自动}是有限型的 (1.5.1)。对于 (I, 6.4.9)，使用 (1.3.7)
和 (1.3.4, (v)) 即可。
\end{env}

\begin{env}[1.7.2]
\label{IV.1.7.2-fr}
在 (I, 6.5.1, (ii)) 中，可以不再假设 $S$ 局部 Noether 且 $Y$ 在
$S$ 上有限型，而只假设 $Y$ 在 $S$ 上\emph{局部有限表示}；证明立即表明
这一点（应用定义 (1.4.2)）。同样，在 (I, 6.5.4, (ii)) 和
(I, 6.5.5, (ii)) 中，只须假设 $f$ 是一个 $S$-态射，并假设 $Y$ 在
$S$ 上局部有限表示。
\end{env}

\begin{env}[1.7.3]
\label{IV.1.7.3-fr}
在 (I, 7.1.11) 的第二个断言中，可以不再假设 $S$ 局部 Noether 且
$Y$ 在 $S$ 上有限型，而只假设 $Y$ 在 $S$ 上局部有限表示（证明依赖于
(I, 6.5.1, (ii))）；在 (I, 7.1.12 至 7.1.15) 中亦然。
\end{env}

\begin{env}[1.7.4]
\label{IV.1.7.4-fr}
在 (I, 9.2.2) 中，依据 (1.2.6) 和 (1.2.7)，可以把三个假设替换为
单一假设（其余每个假设都蕴含它）：$f$ \emph{拟紧且拟分离}。应注意，
在 (I, 9.2.1) 中，该假设所表达的恰好就是 $f$ 拟紧且拟分离。
\end{env}

\begin{env}[1.7.5]
\label{IV.1.7.5-fr}
在 (I, 9.3.1、9.3.2 和 9.3.3) 中，可以把关于 $X$ 的两个假设替换为
较弱的假设：$X$ 是\emph{拟紧且拟分离}的预概形；证明所用的恰好就是
这两个性质（依据 (1.2.7)）。
\end{env}

\begin{env}[1.7.6]
\label{IV.1.7.6-fr}
在 (I, 9.3.4 和 9.3.5) 中，只须假设 $X$ 拟紧，$\mathcal F$ 和
$\mathcal J$ 拟凝聚且有限型。
\end{env}

\begin{env}[1.7.7]
\label{IV.1.7.7-fr}
在 (I, 9.4.7) 中，可以将假设 b) 削弱为只假设 $X$ \emph{拟分离}，
因为只需典范浸入 $U\to X$ 拟紧即可 (1.2.7)。在 (I, 9.4.8、9.4.9
和 9.4.10) 中也可以作同样的修改。
\end{env}

\begin{env}[1.7.8]
\label{IV.1.7.8-fr}
在 (I, 9.5.1) 中，可以把陈述里括号内所列的条件替换为 $f$ 拟紧且拟分离
这一条件。
\end{env}

\begin{env}[1.7.9]
\label{IV.1.7.9-fr}
在 (I, 9.6.5) 中，可以把关于 $X$ 的假设替换为较弱的假设：$X$ 拟紧且
拟分离；证明本身表明了这一点，

\oldpage[IV-1]{237}

因为它只使用 (I, 9.4.7)。在 (I, 9.6.6) 中，假设“$X$ 是拟紧概形”
可以替换为“$X$ 是拟紧且拟分离的预概形”。
\end{env}

\begin{env}[1.7.10]
\label{IV.1.7.10-fr}
在 (II, 1.7.15) 中，只须假设 $Y$ 拟紧且拟分离，因为证明只使用
(I, 9.6.5)。
\end{env}

\begin{env}[1.7.11]
\label{IV.1.7.11-fr}
在 (II, 3.4.5) 的第二个断言中，可以把关于 $Y$ 的两个假设替换为较弱的
假设：$Y$ 拟紧且拟分离。在 (II, 3.4.8) 中亦然。
\end{env}

\begin{env}[1.7.12]
\label{IV.1.7.12-fr}
在 (II, 3.8.5) 中，也可以只假设 $Y$ 拟紧且拟分离；这里通过
(I, 9.4.7) 用到了该假设。
\end{env}

\begin{env}[1.7.13]
\label{IV.1.7.13-fr}
在 (II, 4.4.1、4.4.6 和 4.4.7) 中，只须假设 $Y$ \emph{拟紧且拟分离}；
在 (II, 4.4.7) 的证明中要结合 (1.2.3)。同样，在 (II, 4.4.10.1) 中，
只须假设 $Z$ 拟紧且拟分离。
\end{env}

\begin{env}[1.7.14]
\label{IV.1.7.14-fr}
在 (II, 4.5.2) 和 (II, 4.5.5) 中，只须假设 $X$ \emph{拟紧且拟分离}；
这里通过 (I, 9.3.1 和 9.3.2) 用到了该假设。
\end{env}

\begin{env}[1.7.15]
\label{IV.1.7.15-fr}
在 (II, 4.6.8) 中，仍只须假设 $X$ \emph{拟紧且拟分离}；这里通过
(I, 9.4.7) 用到了该假设。
\end{env}

\begin{env}[1.7.16]
\label{IV.1.7.16-fr}
在 (II, 5.1.2) 中，只须假设 $X$ 拟紧且拟分离（这里通过
(II, 4.5.2 和 4.5.5) 用到了该假设）。在使用 (I, 9.6.5) 的
(II, 5.1.9) 中亦然。
\end{env}

\begin{env}[1.7.17]
\label{IV.1.7.17-fr}
在 (II, 5.2.1) 中，仍只须假设 $X$ 拟紧且拟分离（这里通过
(II, 4.5.2) 用到了该假设）。
\end{env}

\begin{env}[1.7.18]
\label{IV.1.7.18-fr}
在 (II, 5.2.2) 中，必须假设 $f$ 拟紧且 $X$ 拟分离；事实上，按原有
假设，现有证明并不充分：由每个拟凝聚 $\mathcal O_X$-模 $\mathcal F$
都满足 $\mathrm R^1f_*(\mathcal F)=0$，并不必然能推出，对于仿射开集
$U$（它是 $Y$ 的开集），同一性质对拟凝聚
$(\mathcal O_X|f^{-1}(U))$-模也成立，除非已知这种模是某个拟凝聚
$\mathcal O_X$-模的\emph{限制}（同一说明也适用于条件 b) 和 c')）；当
$f$ 拟紧且 $X$ 拟分离时，依据上面经过修改的 (I, 9.4.7) (1.7.7)，
这种延拓是可能的。
\end{env}

\begin{env}[1.7.19]
\label{IV.1.7.19-fr}
在 (II, 5.3.2) 中，只须假设 $Y$ \emph{拟紧且拟分离}（这里通过
(II, 4.4.6 和 4.4.7) 用到了该假设）。在 (II, 5.5.3, (ii)) 中亦然。
\end{env}

\begin{env}[1.7.20]
\label{IV.1.7.20-fr}
在 (III, 1.2.2、1.2.3 和 1.2.4) 中，可以把关于 $X$ 的两个假设替换为
较弱的假设：$X$ 是\emph{拟紧且拟分离}的预概形；证明只使用
(I, 9.3.3)。
\end{env}

\begin{env}[1.7.21]
\label{IV.1.7.21-fr}
在 (III, 1.4.10 至 1.4.14) 中，可以把“分离”替换为“拟分离”，因为
该假设只用于应用 (I, 9.2.2)（见上文 (1.7.4)）。同样，在
(III, 1.4.15) 中，只须假设态射 $u$ \emph{拟分离且拟紧}。
\end{env}

\begin{env}[1.7.22]
\label{IV.1.7.22-fr}
在 (III, 2.1.3) 中，仍可以把各项假设替换为较弱的假设：$X$ 是拟紧且
拟分离的预概形；其论证与 (III, 1.2.2) 相同。
\end{env}

\subsection{有限表示态射与可构造集}
\label{subsection:IV.1.8-fr}

\oldpage[IV-1]{238}

\begin{env}[1.8.1]
\label{IV.1.8.1-fr}
设 $X$ 是预概形；已知 $X$ 中每个拟紧开集都是有限个仿射开集的并，反之亦然。
一个开集 $U$ 在 $X$ 中为\emph{逆紧} (0\textsubscript{III}, 9.1.1) 的
充分必要条件是：在 $X$ 中，对每个\emph{仿射开集} $V$（它是 $X$ 的开集），
$U\cap V$ 都拟紧。
当 $X$ \emph{拟分离}时 (1.2.1)，$X$ 中每个拟紧开集都在 $X$ 中逆紧
(1.2.7)，因而 $X$ 的每个局部可构造子集都在 $X$ 中逆紧
(0\textsubscript{III}, 9.1.13)；若\emph{再}假设 $X$ \emph{拟紧}，则
$X$ 中的可构造子集与局部可构造子集相同 (0\textsubscript{III}, 9.1.12)，
可构造开集与拟紧开集也相同 (0\textsubscript{III}, 9.1.5)。
\end{env}

\begin{proposition}[1.8.2]
\label{IV.1.8.2-fr}
设 $X$、$Y$ 是两个预概形，$f:X\to Y$ 是一个态射。对任意可构造
（分别地，局部可构造）子集 $Z$（它是 $Y$ 的子集），$f^{-1}(Z)$ 是 $X$ 的可构造
（分别地，局部可构造）子集。
\end{proposition}

\begin{proof}
假设 $Z$ 局部可构造；对每个 $x\in X$，存在仿射开邻域 $V$，它含有 $f(x)$，
使 $Z\cap V$ 在 $V$ 中可构造；若命题对可构造子集已经证明，则
$f^{-1}(Z)\cap f^{-1}(V)$ 在 $f^{-1}(V)$ 中可构造，因而 $f^{-1}(Z)$
局部可构造。因此只须考虑 $Z$ 可构造的情形；结合
(0\textsubscript{III}, 9.1.3)，可归结到 $Z$ 是 $Y$ 中的开集且\emph{逆紧}
的情形，换言之，典范注入 $Z\to Y$ 是拟紧态射；此时只须证明
$f^{-1}(Z)$ \emph{在 $X$ 中逆紧}，换言之，典范注入
$f^{-1}(Z)\to X$ 是拟紧态射；但由 (1.1.2, (iii)) 以及
$f^{-1}(Z)=Z\times_YX$ (I, 4.4.1)，这立即成立。
\end{proof}

\begin{lemma}[1.8.3]
\label{IV.1.8.3-fr}
设 $X$ 是拟紧且拟分离的预概形，$Z$ 是 $X$ 的一个可构造子集。则存在
仿射概形 $X'$ 和有限表示态射 $f:X'\to X$，使得 $f(X')=Z$。
\end{lemma}

\begin{proof}
由于 $X$ 拟紧，它是有限多个仿射开集 $X_i$（$i\in I$）的并；又由于
$X$ 拟分离，由 (1.2.7) 可知，各典范浸入 $g_i:X_i\to X$ 都是有限表示的。
因此，若 $f_i:X'_i\to X_i$ 是有限表示态射，则
$g_i\circ f_i:X'_i\to X$ 亦然 (1.6.2)；若 $X'$ 是各 $X'_i$ 的和预概形，
而 $h:X'\to X$ 与 $g_i\circ f_i$ 在相应的每个 $X'_i$ 上相同，则
该态射也是有限表示的 (1.6.5)。由于 $Z\cap X_i$ 在 $X_i$ 中可构造
(0\textsubscript{III}, 9.1.8)，故只须在 $X_i$ 是环 $A$ 的\emph{仿射}
概形时证明引理。此时 $Z$ 是有限多个形如 $U\cap\complement V$ 的集合之并，
其中 $U$ 和 $V$ 是拟紧开集 (0\textsubscript{III}, 9.1.3)；再取适当的
预概形之和，便可归结到 $Z=U\cap\complement V$ 的情形；又因为 $U$
是有限多个仿射开集的并，甚至可以只考虑 $U$ 为\emph{仿射}的情形。
由于 $V$ 拟紧，它是有限多个形如 $X_{f_i}$ 的开集之并，其中 $f_i\in A$；
设 $\mathfrak J$ 为 $A$ 中由各 $f_i$ 生成的理想，并设 $Z'$ 为它所定义的
$X$ 的仿射闭子概形（按构造，它在 $X$ 上有限表示）；依定义有
$V=\complement Z'$，故 $U\cap\complement V=U\cap Z'$。考虑仿射概形
$X'=Z'\times_XU$

\oldpage[IV-1]{239}

并令 $f:X'\to X$ 为结构态射；刚才已经看到，典范浸入 $Z'\to X$
是有限表示的；开浸入 $U\to X$ 亦然，因为 $U$ 和 $X$ 都是仿射的，
所以该浸入拟紧。因此由 (1.6.2, (iv)) 可知 $f$ 是有限表示的，并且
$f(X')=Z'\cap U$ (I, 3.4.8)，从而证明完毕。
\end{proof}

\begin{theorem}[1.8.4]
\label{IV.1.8.4-fr}
\textup{（Chevalley）}。--- 设 $f:X\to Y$ 是有限表示态射。对任意局部
可构造子集 $Z$（它是 $X$ 的子集），$f(Z)$ 在 $Y$ 中局部可构造。
\end{theorem}

\begin{proof}
设 $y\in Y$，$V$ 是 $y$ 的一个仿射开邻域；由于 $f$ 拟紧且拟分离，
它的限制 $f^{-1}(V)\to V$ 也拟紧且拟分离，故 $f^{-1}(V)$ 是拟紧且
拟分离的预概形 (1.2.3)。又因为 $Z\cap f^{-1}(V)$ 在 $f^{-1}(V)$ 中
可构造 (1.8.1)，所以只须考虑 $Y$ 仿射且 $Z$ 可构造的情形；此时 $X$
拟紧且拟分离，故由 (1.8.3)，存在有限表示态射 $g:X'\to X$，使
$Z=g(X')$。由于 $f\circ g$ 有限表示，问题便归结到 $Z=X$ 的情形；
换言之，只须证明下述引理。
\end{proof}

\begin{lemma}[1.8.4.1]
\label{IV.1.8.4.1-fr}
设 $Y$ 是仿射概形，$f:X\to Y$ 是拟紧且局部有限表示的态射；则 $f(X)$
是 $Y$ 的可构造子集。
\end{lemma}

\begin{proof}
由于 $X$ 拟紧，它是有限多个仿射开集的并；因此只须考虑
$Y=\operatorname{Spec}(A)$、$X=\operatorname{Spec}(B)$ 的情形，其中 $B$
是有限表示 $A$-代数 (1.4.6)。现有下述引理。
\end{proof}

\begin{lemma}[1.8.4.2]
\label{IV.1.8.4.2-fr}
设 $A_0$ 是一个环，$(A_\lambda)_{\lambda\in L}$ 是 $A_0$-代数的一个
归纳系统，$A=\varinjlim A_\lambda$；若 $B$ 是有限表示 $A$-代数，则存在
某个 $\lambda$ 以及有限表示 $A_\lambda$-代数 $B_\lambda$，使 $B$ 同构于
$B_\lambda\otimes_{A_\lambda}A$。
\end{lemma}

\begin{proof}
依假设，$B$ 同构于形如 $A[T_1,\ldots,T_n]/\mathfrak J$ 的代数，其中
各 $T_i$ 是不定元，而 $\mathfrak J$ 是有限型理想；设 $(F_j)$
（$1\leq j\leq m$）是 $\mathfrak J$ 的一组生成元。令
$\varphi_\lambda:A_\lambda\to A$ 为典范同态。由于 $L$ 是滤过的，存在
$\lambda$，使每个多项式 $F_j$ 的每个系数都是在 $\varphi_\lambda$ 下的
某个 $A_\lambda$ 元素之像。换言之，存在一族多项式
$(F_{j\lambda})_{1\leq j\leq m}$，它们属于
$A_\lambda[T_1,\ldots,T_n]$，并使得 $F_j=\varphi_\lambda(F_{j\lambda})$
对 $1\leq j\leq m$ 均成立。若 $\mathfrak J_\lambda$ 是
$A_\lambda[T_1,\ldots,T_n]$ 中由各 $F_{j\lambda}$ 生成的理想，则
$\mathfrak J$ 是 $\mathfrak J_\lambda\otimes_{A_\lambda}A$ 在
$A[T_1,\ldots,T_n]=A_\lambda[T_1,\ldots,T_n]\otimes_{A_\lambda}A$ 中的像；
环 $B_\lambda=A_\lambda[T_1,\ldots,T_n]/\mathfrak J_\lambda$ 即满足要求。
\end{proof}

这里把 $A$ 看作其各个\emph{有限型}子-$\mathbf Z$-代数的归纳极限，
并应用该引理。于是存在这样的子-$\mathbf Z$-代数 $A_0$，它包含于 $A$ 中；
还存在有限表示
$A_0$-代数 $B_0$，使得 $B$ 同构于 $B_0\otimes_{A_0}A$；置
$X_0=\operatorname{Spec}(B_0)$、$Y_0=\operatorname{Spec}(A_0)$，于是有
$X=X_0\times_{Y_0}Y$，且 $f$ 是投影 $X\to Y$；设
$f_0:X_0\to Y_0$、$g_0:Y\to Y_0$ 为结构态射。由 (I, 3.4.8) 得
$f(X)=g_0^{-1}(f_0(X_0))$；结合 (1.8.2)，只须证明 $f_0(X_0)$
可构造。换言之，最终归结为证明下述推论。

\begin{corollary}[1.8.5]
\label{IV.1.8.5-fr}
设 $Y$ 是 Noether 预概形，$f:X\to Y$ 是有限型态射。对任意可构造
子集 $Z$（它是 $X$ 的子集），$f(Z)$ 是 $Y$ 的可构造子集。
\end{corollary}

\oldpage[IV-1]{240}

\begin{proof}
同上，只须证明 $f(X)$ 可构造。应用判别准则 (0\textsubscript{III}, 9.2.3)：
须证明，对任意不可约闭子集 $T$（它是 $Y$ 的闭子集），$T\cap f(X)=f(f^{-1}(T))$
或者在 $T$ 中无处稠密，或者包含 $T$ 的一个非空开集。取 $Y$ 的一个
底空间为 $T$ 的子预概形 (I, 5.2.1)，并考虑它在 $f$ 下的逆像；结合
(1.5.4, (iii))，最终只须证明：若假设 $Y$ \emph{不可约且 $f$ 占优}，
则 $f(X)$ 包含 $Y$ 的一个非空开集。还可以假设 $Y$ 仿射；此时 $X$
是有限多个仿射开集 $V_i$ 的并，而由于 $f(X)$ 在 $Y$ 中稠密，至少一个
$f(V_i)$ 也稠密。因此还可以假设 $X$ 仿射；若 $(X_j)$ 是 $X$
的不可约分支所成的有限族，则如上可知，至少一个 $f(X_j)$ 在 $Y$ 中稠密；
故可假设 $X$ 不可约。最后，把 $f$ 替换为 $f_{\mathrm{red}}$
(1.5.4, (vi))，可以假设 $X$ 和 $Y$ \emph{都是整的}。于是
$X=\operatorname{Spec}(B)$、$Y=\operatorname{Spec}(A)$，其中 $A$ 是
Noether 整环，$B$ 是包含 $A$ 的有限型整 $A$-代数 (I, 1.2.7)。只须证明
存在 $g\in A$，使得对每个 $y\in D(g)$（即 $g\notin\mathfrak p_y$），
理想 $\mathfrak p_y$ 都是 $A$ 与 $B$ 的某个素理想的交；这将表明
$D(g)\subset f(X)$。最终只须证明下述引理。
\end{proof}

\begin{lemma}[1.8.5.1]
\label{IV.1.8.5.1-fr}
设 $A$ 是整环，$B$ 是有限型整 $A$-代数。存在 $g\in A$，使得 $A$
到代数闭域 $\Omega$ 的每个在 $g$ 上非零的同态，都能延拓为 $B$ 到
$\Omega$ 的同态。
\end{lemma}

\begin{proof}
这是交换代数中的一个经典结果（Bourbaki，\emph{Alg. comm.}，第 V 章，
\S~3，n\textsuperscript{o}~1，定理 1 的推论 3）。
\end{proof}

\begin{corollary}[1.8.6]
\label{IV.1.8.6-fr}
设 $Y$ 是不可约预概形，$\eta$ 是它的泛点，$f:X\to Y$ 是局部有限型态射。
若 $f^{-1}(\eta)$ 非空，则存在一个开邻域 $U$，它是 $\eta$ 在 $Y$ 中的
邻域，并使得 $f^{-1}(y)$ 对每个 $y\in U$ 都非空。若 $\mathscr F$ 是有限型拟凝聚
$\mathcal O_X$-模，并且 $\mathscr F_\eta=\mathscr F\otimes_{\mathcal O_Y}k(\eta)$
非零，则存在一个开邻域 $U'$，它是 $\eta$ 在 $Y$ 中的邻域，并使得
$\mathscr F_y=\mathscr F\otimes_{\mathcal O_Y}k(y)$ 对每个 $y\in U'$ 都非零。
\end{corollary}

\begin{proof}
若 $p_y$ 是纤维 $f^{-1}(y)=X\times_Y\operatorname{Spec}(k(y))$ 到 $X$
的典范投影，则 $\mathscr F_y=p_y^*(\mathscr F)$，因而有
$\operatorname{Supp}(\mathscr F_y)=p_y^{-1}(\operatorname{Supp}(\mathscr F))
=\operatorname{Supp}(\mathscr F)\cap f^{-1}(y)$；这是因为 (I, 9.1.13.1)、
(I, 3.6.1) 成立，且 $\mathscr F$ 有限型。此外，
$\operatorname{Supp}(\mathscr F)$ 在 $X$ 中是闭的 (0\textsubscript{I},
5.2.2)；若 $Z$ 是 $X$ 的一个闭子预概形，其底空间为
$\operatorname{Supp}(\mathscr F)$ (I, 5.2.1)，而 $j$ 是典范浸入
$Z\to X$，则 $f\circ j$ 局部有限型 (1.3.4)。这证明了第一个断言蕴含
第二个断言。为证明第一个断言，先注意可以假设 $X$ 和 $Y$ 都既约
(1.3.4, (vi))，换言之，$Y$ 是整的。设 $\xi$ 是 $f^{-1}(\eta)$ 中的一点；
可以用 $X$ 和 $Y$ 中分别包含 $\xi$ 和 $\eta$ 的仿射开邻域替换之，从而假设
$X=\operatorname{Spec}(B)$、$Y=\operatorname{Spec}(A)$，其中 $B$ 是
有限型 $A$-代数。此外，若 $Z$ 是 $X$ 的既约闭子预概形，其底集为
$\overline{\{\xi\}}$ (I, 5.2.1)，则可如上把 $X$ 替换为 $Z$，换言之，
可以假设 $B$ 为\emph{整环} (I, 5.1.4)。此时态射 $f$ 占优，所以相应同态
$\varphi:A\to B$ 是单射 (I, 1.2.7)；推论最终由引理 (1.8.5.1) 得出。
\end{proof}

\oldpage[IV-1]{241}

\begin{proposition}[1.8.7]
\label{IV.1.8.7-fr}
设 $S$ 是预概形，$g:X\to S$、$h:Y\to S$ 是两个有限表示态射，
$f:X\to Y$ 是一个 $S$-态射。对每个 $s\in S$，置
$X_s=g^{-1}(s)=X\times_S\operatorname{Spec}(k(s))$、
$Y_s=h^{-1}(s)=Y\times_S\operatorname{Spec}(k(s))$、
$f_s=f\times 1:X_s\to Y_s$。则由 $s\in S$ 中满足“$f_s$ 为满射
（分别地，为根态射）”者所成的集合是局部可构造的。
\end{proposition}

\begin{proof}
设 $E$ 为那些 $s\in S$ 所成的集合，其中 $f_s$ 满射；集合 $S-E$ 等于
$h(Y-f(X))$。由于 $f$ 是有限表示的 (1.6.2, (v))，$f(X)$ 在 $Y$ 中
局部可构造 (1.8.4)；又因为 $h$ 有限表示，$h(Y-f(X))$ 在 $S$ 中
局部可构造，因而 $E$ 也局部可构造。

为了证明集合 $E'$（由满足 $s\in S$ 且 $f_s$ 为根态射的点组成）局部可构造，
我们将使用下述引理。
\end{proof}

\begin{lemma}[1.8.7.1]
\label{IV.1.8.7.1-fr}
设 $U$、$V$ 是两个预概形；态射 $f:U\to V$ 为根态射的充分必要条件是
对角态射 $\Delta_f:U\to U\times_VU$ 为满射。因此，每个根态射都是
分离态射。
\end{lemma}

\begin{proof}
事实上，$\Delta_f$ 是浸入 (I, 5.3.9)，因而是局部有限型态射 (1.3.4)；
说 $\Delta_f$ 为满射，等价于说对每个代数闭域 $K$，相应映射
$U(K)\to(U\times_VU)(K)=U(K)\times_{V(K)}U(K)$ 为满射
(1.3.7 和 I, 3.4.2.1)；但由集合纤维积的定义，这意味着映射
$U(K)\to V(K)$（它与 $f$ 相应）是单射，而这等价于 $f$ 是根态射
(I, 3.5.5)。

现在，$\Delta_{f_s}$ 由 $\Delta_f:X\to X\times_YX$ 经基变换
$\operatorname{Spec}(k(s))\to S$ 得到 (I, 5.3.4)。由于 $f$ 有限表示，
结构态射 $X\times_YX\to Y$ 亦然 (1.6.2, (iv))，从而 $\Delta_f$
也有限表示 (1.6.2, (v))；因此只须把命题的第一部分应用于 $\Delta_f$，
并使用引理 (1.8.7.1)。
\end{proof}

\subsection{pro-可构造集与 ind-可构造集}
\label{subsection:IV.1.9-fr}

\begin{lemma}[1.9.1]
\label{IV.1.9.1-fr}
设 $(A_\alpha)_{\alpha\in I}$ 是环的一个归纳系统，其指标集 $I$ 向右滤过，
并设 $A=\varinjlim A_\alpha$。等式 $A=0$ 成立的充分必要条件是：存在一个
指标 $\gamma$，使得 $A_\gamma=0$（此时 $A_\beta=0$ 对每个
$\beta\geq\gamma$ 都成立）。
\end{lemma}

\begin{proof}
事实上，对每个 $\alpha\in I$，典范同态
$\varphi_\alpha:A_\alpha\to A$ 把单位元映为单位元。说 $A=0$，就是
$\varphi_\alpha(1)=0$，也就是
$\varphi_\alpha(1)=\varphi_\alpha(0)$。已知这等价于存在一个指标
$\gamma\geq\alpha$，使得同态
$\varphi_{\gamma\alpha}:A_\alpha\to A_\gamma$ 满足
$\varphi_{\gamma\alpha}(1)=\varphi_{\gamma\alpha}(0)$；而最后这个关系等价于
$A_\gamma=0$，引理得证。
\end{proof}

\begin{proposition}[1.9.2]
\label{IV.1.9.2-fr}
设 $B$ 是环，$(A_\alpha)_{\alpha\in I}$ 是 $B$-代数的一个归纳系统，
其指标集 $I$ 向右滤过，并设 $A=\varinjlim A_\alpha$；置
$Y=\operatorname{Spec}(B)$、$X_\alpha=\operatorname{Spec}(A_\alpha)$、
$X=\operatorname{Spec}(A)$，并令 $u:X\to Y$、$u_\alpha:X_\alpha\to Y$
为分别对应于结构同态 $B\to A$、$B\to A_\alpha$ 的态射。则：
\end{proposition}

\begin{env}[1.9.2.i]
\label{IV.1.9.2.i-fr}
\emph{$X=\varnothing$ 成立的充分必要条件是存在 $\gamma\in I$，使得
$X_\gamma=\varnothing$；此时 $X_\alpha=\varnothing$ 对
$\alpha\geq\gamma$ 均成立。}
\end{env}

\oldpage[IV-1]{242}

\begin{env}[1.9.2.ii]
\label{IV.1.9.2.ii-fr}
\emph{有}
\[
  u(X)=\bigcap_{\alpha\in I}u_\alpha(X_\alpha). \tag{1.9.2.1}
  \label{IV.1.9.2.1-fr}
\]
\end{env}

断言 (i) 正是 (1.9.1) 的翻译。为证明 (ii)，注意到 $u$ 分解为
$X\to X_\alpha\xrightarrow{u_\alpha}Y$，所以 (1.9.2.1) 的左端包含于右端。
反过来，设 $y\in Y-u(X)$；则 $u^{-1}(y)=\varnothing$，而 $u^{-1}(y)$ 是
$k(y)$-预概形 $\operatorname{Spec}(A\otimes_Bk(y))$ 的底空间。在 $B$-模范畴中，
函子 $\varinjlim$ 与张量积交换，所以 $A\otimes_Bk(y)$ 是环的归纳系统
$A_\alpha\otimes_Bk(y)$ 的归纳极限。由引理 (1.9.1)，存在 $\alpha$ 使
$A_\alpha\otimes_Bk(y)=0$，也就是 $u_\alpha^{-1}(y)=\varnothing$，从而
$y\notin u_\alpha(X_\alpha)$。证毕。

\begin{proposition}[1.9.3]
\label{IV.1.9.3-fr}
设 $X$ 是拟紧且拟分离的预概形，$E$ 是 $X$ 的一个子集，
$(X_\alpha)_{\alpha\in I}$ 是 $X$ 的一个仿射开覆盖，并且对每个
$\alpha\in I$，置 $E_\alpha=E\cap X_\alpha$。以下条件等价：
\begin{enumerate}
  \item[a)] 存在拟紧态射 $f:X'\to X$，使得 $E=f(X')$。
  \item[a$'$)] 对每个 $\alpha$，存在拟紧态射
  $f_\alpha:X'_\alpha\to X_\alpha$，使得
  $E_\alpha=f_\alpha(X'_\alpha)$。
  \item[a$''$)] 对每个 $\alpha$，存在仿射概形 $X''_\alpha$ 和态射
  $f'_\alpha:X''_\alpha\to X_\alpha$，使得
  $f'_\alpha(X''_\alpha)=E_\alpha$。
  \item[b)] $E$ 是 $X$ 的一些可构造子集之交。
  \item[b$'$)] $F=X-E$ 是 $X$ 的一些可构造子集之并。
\end{enumerate}
\end{proposition}

\begin{proof}
显然，$b)$ 与 $b')$ 等价。条件 $a)$ 蕴含 $a')$：取 $X'_\alpha$ 为
开集 $f^{-1}(X_\alpha)$ 上由 $X'$ 诱导的预概形，并取 $f_\alpha$ 为 $f$
的限制。为说明 $a')$ 蕴含 $a'')$，只须注意 $X'_\alpha$ 是一些仿射开集
$X''_{\alpha\lambda}$（$\lambda\in L$）的并；令 $X''_\alpha$ 为各
$X''_{\alpha\lambda}$（$\lambda\in L$）的预概形emph{之和}，它是仿射概形。
若 $g_\alpha:X''_\alpha\to X'_\alpha$ 是在每个 $X''_{\alpha\lambda}$ 上
都与恒等态射相同的态射，则置 $f'_\alpha=f_\alpha\circ g_\alpha$ 后，
显然有 $E_\alpha=f'_\alpha(X''_\alpha)$，因为 $g_\alpha$ 满射。为证明
$a'')$ 蕴含 $a)$，注意到存在一个有限子集 $J$，它包含于 $I$ 中，并使得
各 $X_\alpha$（$\alpha\in J$）覆盖 $X$。令 $X'$ 为各 $X''_\alpha$
（$\alpha\in J$）的和预概形，并令 $f:X'\to X$ 为与
$j_\alpha\circ f'_\alpha$ 对每个 $\alpha\in J$ 都相同的态射，其中
$j_\alpha:X_\alpha\to X$ 是典范嵌入。由于 $E$ 是各 $E\cap X_\alpha$
（$\alpha\in J$）的并，确有 $E=f(X')$；只剩证明 $f$ 拟紧。由于 $X$
拟分离，$j_\alpha$ 拟紧 (1.2.7)，所以 $j_\alpha\circ f'_\alpha$ 也拟紧
(1.1.1 和 1.1.2)，最终 $f$ 拟紧 (1.1.6)。

因此只须证明 $a'')$ 与 $b)$ 等价。先证 $a'')$ 蕴含 $b)$：取一个有限子集
$J$，它包含于 $I$ 中，并使得各 $X_\alpha$（$\alpha\in J$）覆盖 $X$。
只须证明，对每个 $\alpha\in J$，$E_\alpha$ 是一些可构造子集
$F_{\alpha\lambda}$ 之交；这些子集位于 $X_\alpha$ 中
（$\lambda\in L_\alpha$）。事实上，由
(0\textsubscript{III}, 9.1.8, (ii))，每个 $F_{\alpha\lambda}$ 也是 $X$
的可构造子集，因为 $X$ 拟分离这一假设蕴含每个 $X_\alpha$ 在 $X$ 中
逆紧 (1.2.7)。由于 $E$ 是各 $E_\alpha$（$\alpha\in J$）的并，它等于
所有有限并 $\bigcup_{\alpha\in J}F_{\alpha,\lambda(\alpha)}$ 的交，其中
$\lambda(\alpha)\in L_\alpha$ 取遍所有选择。这些并在 $X$ 中可构造，
因而断言成立。于是只须考虑如下情形：

\oldpage[IV-1]{243}

$X=\operatorname{Spec}(A)$，且 $E=f(X')$，其中
$X'=\operatorname{Spec}(A')$ 为仿射概形，$f:X'\to X$ 是与环同态
$A\to A'$ 对应的态射。现注意到，$A'$ 是族 $(A'_\lambda)$ 按包含关系
排序所得的归纳极限；其成员是\emph{有限型}子-$A$-代数。若
$X'_\lambda=\operatorname{Spec}(A'_\lambda)$，则 $f$ 分解为
$X'\to X'_\lambda\xrightarrow{f_\lambda}X$；由 (1.9.2, (ii)) 得
$f(X')=\bigcap_\lambda f_\lambda(X'_\lambda)$。若能证明
$f_\lambda(X'_\lambda)$ 是 $X$ 的一些可构造子集之交，则 $f(X')$ 亦然。
因此可归结到 $A'$ 是\emph{有限型} $A$-代数的情形。现有下述引理。
\end{proof}

\begin{lemma}[1.9.3.1]
\label{IV.1.9.3.1-fr}
设 $A$ 是环，$A'$ 是有限型 $A$-代数。则 $A'$ 在 $A$ 上与一个滤过归纳极限
$\varinjlim A''_\mu$ 同构，其中各 $A''_\mu$ 都是有限表示
$A$-代数 (1.4.1)。
\end{lemma}

\begin{proof}
事实上，可以写成 $A'=B/\mathfrak J$，其中 $B=A[T_1,\ldots,T_n]$，
$\mathfrak J$ 是 $B$ 的一个理想。但 $\mathfrak J$ 是各有限型理想
$\mathfrak J'_\mu$ 的归纳极限；这些是 $B$ 的理想，并包含于 $\mathfrak J$ 中，
按包含关系排序。由于在 $B$-模范畴中函子 $\varinjlim$ 正合，遂得
$A'=\varinjlim B/\mathfrak J'_\mu$，相差一个 $A$-同构。
\end{proof}

于是用与上面相同的论证，可以归结到 $A'$ 是\emph{有限表示} $A$-代数的情形；
但此时由 Chevalley 定理 (1.8.5)，$f(X')$ 在 $X$ 中可构造，从而证明了
$b)$。

最后证明 $b)$ 蕴含 $a'')$。若 $E$ 是一族可构造子集
$(G_\lambda)_{\lambda\in L}$ 之交，这族子集位于 $X$ 中，则每个 $E_\alpha$ 都是各
$G_\lambda\cap X_\alpha$ 之交；后者在 $X_\alpha$ 中可构造
(0\textsubscript{III}, 9.1.8, (i))。因此可归结到
$X=\operatorname{Spec}(A)$ 仿射的情形。

这时由 (1.8.3) 知，对每个 $\lambda\in L$，存在态射
$f_\lambda:X'_\lambda\to X$，其中
$X'_\lambda=\operatorname{Spec}(A'_\lambda)$ 为仿射概形，并满足
$f_\lambda(X'_\lambda)=G_\lambda$。因此只须证明下述引理。

\begin{lemma}[1.9.3.2]
\label{IV.1.9.3.2-fr}
设 $(A'_\lambda)_{\lambda\in L}$ 是一族 $A$-代数，
$X=\operatorname{Spec}(A)$，$X'_\lambda=\operatorname{Spec}(A'_\lambda)$，
并设 $f_\lambda:X'_\lambda\to X$ 为结构态射。对任意有限子集
$J$（它是 $L$ 的子集），置 $A'_J=\bigotimes_{\lambda\in J}A'_\lambda$（$A$-代数的张量积）、
$X'_J=\operatorname{Spec}(A'_J)$，并设 $f_J:X'_J\to X$ 为结构态射。则
$f_J(X'_J)=\bigcap_{\lambda\in J}f_\lambda(X'_\lambda)$。若
$A'=\varinjlim A'_J$，其中 $J$ 遍历 $L$ 的有限子集所成的滤过集，且
$X'=\operatorname{Spec}(A')$，并且 $f:X'\to X$ 为结构态射，则
$f(X')=\bigcap_{\lambda\in L}f_\lambda(X'_\lambda)$。
\end{lemma}

\begin{proof}
第一个断言由 (I, 3.4.7) 得出；第二个断言由 (1.9.2, (ii)) 得出，因为后者给出
$f(X')=\bigcap_Jf_J(X'_J)$。
\end{proof}

\begin{definition}[1.9.4]
\label{IV.1.9.4-fr}
设 $X$ 是拓扑空间。称一个子集 $E$ 在 $X$ 中为
\emph{pro-可构造的}（分别地，\emph{ind-可构造的}），如果对 $X$ 中每个
$x\in X$，存在开邻域 $U$，它是 $x$ 在 $X$ 中的邻域，并使得
$E\cap U$ 是 $U$ 的局部可构造子集之交（分别地，之并）。
\end{definition}

考虑到 (1.2.7) 以及 (0\textsubscript{III}, 9.1.8 和 9.1.12)，命题
(1.9.3) 的诸等价条件因而可表述为：$E$ 在 $X$ 中是
\emph{pro-可构造的}；在 (1.9.3) 的假设下，$X$ 中的局部可构造集与
$X$ 中的可构造集相同。

\oldpage[IV-1]{244}

\begin{proposition}[1.9.5]
\label{IV.1.9.5-fr}
设 $X$ 是预概形。
\begin{enumerate}
  \item[(i)] 子集 $E$ 在 $X$ 中为 pro-可构造，当且仅当 $X-E$ 为
  ind-可构造。
  \item[(ii)] $X$ 的 pro-可构造子集的任意有限并以及任意交都是
  pro-可构造的。$X$ 的 ind-可构造子集的任意有限交以及任意并都是
  ind-可构造的。
  \item[(iii)] $X$ 的局部可构造子集的任意交（分别地，任意并）是
  pro-可构造的（分别地，ind-可构造的）。反过来，若 $X$ 拟紧且
  拟分离，则 $X$ 的每个 pro-可构造（分别地，ind-可构造）子集都是
  $X$ 的可构造子集之交（分别地，之并）。
  \item[(iv)] 设 $(U_\alpha)$ 是 $X$ 的开覆盖。子集 $E$ 在 $X$ 中为
  pro-可构造的（分别地，ind-可构造的），当且仅当在 $X$ 中，对每个 $\alpha$，
  $E\cap U_\alpha$ 在 $U_\alpha$ 中为 pro-可构造的（分别地，
  ind-可构造的）。
  \item[(v)] 每个 pro-可构造子集 $E$（属于 $X$）在 $X$ 中都是逆紧的。
  一个局部闭子集 $E$（属于 $X$）在 $X$ 中为逆紧，当且仅当它在
  $X$ 中为 pro-可构造的。
  \item[(vi)] 设 $f:X'\to X$ 是预概形态射；对每个 pro-可构造
  （分别地，ind-可构造）子集 $E$（属于 $X$），$f^{-1}(E)$ 在 $X'$ 中
  为 pro-可构造的（分别地，ind-可构造的）。
  \item[(vii)] 设 $f:X'\to X$ 是拟紧态射；对每个 pro-可构造子集
  $E'$（属于 $X'$），$f(E')$ 在 $X$ 中为 pro-可构造的；特别地，
  $f(X')$ 在 $X$ 中为 pro-可构造的。
  \item[(viii)] 设 $f:X'\to X$ 是局部有限表示的态射；对每个
  ind-可构造子集 $E'$（属于 $X'$），$f(E')$ 在 $X$ 中为 ind-可构造的；特别地，
  $f(X')$ 在 $X$ 中为 ind-可构造的。
  \item[(ix)] 假设 $X$ 拟紧且拟分离；则子集 $E$ 在 $X$ 中为
  pro-可构造的，当且仅当存在仿射概形 $X'$ 以及态射（必然拟紧）
  $f:X'\to X$，使得 $E=f(X')$。
\end{enumerate}
\end{proposition}

\begin{proof}
性质 (i)、(ii)、(iv) 以及 (iii) 的第一个断言都由定义 (1.9.4) 得出，
并且利用交与并的相互分配律以及下述事实，它们对任意拓扑空间都成立：
若 $T$ 在 $X$ 中局部可构造，则 $T\cap U$ 在 $U$ 中局部可构造，其中
$U$ 是 $X$ 的任意开集。若 $X$ 拟紧且拟分离，并且 $E$ 在 $X$ 中为
pro-可构造的，则存在有限开覆盖 $(U_i)$，它由 $X$ 的仿射开集组成，
使得 $E\cap U_i$ 是各可构造子集 $M_{i\lambda}$（属于 $U_i$，
$\lambda\in L_i$）之交；由 (1.2.7) 和
(0\textsubscript{III}, 9.1.8, (ii))，各 $M_{i\lambda}$ 在 $X$ 中也可构造，
而 $E$ 是诸有限并 $\bigcup_iM_{i,\lambda(i)}$（其中
$\lambda(i)\in L_i$ 对每个 $i$）之交；这些有限并都是 $X$ 的可构造子集。这证明了
(iii) 的第二个断言。(ix) 由 (iii) 和 (1.9.3) 得出。为证明 (v) 的
第一个断言，可以归结到 $X$ 仿射的情形，并在此情形证明 $E$ 拟紧
(0\textsubscript{III}, 9.1.1)；但由于 $X$ 此时拟分离，由 (ix) 存在拟紧态射
$f:X'\to X$，使得 $E=f(X')$；由于 $X'$ 拟紧，它在连续映射下的像
$E$ 也拟紧。

为证明 (vi)，由 (iv) 可以归结到 $X$ 仿射的情形；此时结论由 (iii) 和
(1.8.2) 得出。

\oldpage[IV-1]{245}

同样，为证明 (vii)，可以归结到 $X$ 仿射的情形；此时 $X'$ 是有限个
仿射开集 $X'_i$ 的并，而 $f(E')$ 是各 $f(E'\cap X'_i)$ 的并，故由 (ii)
也可假设 $X'$ 仿射；但这时由 (ix)，$E'=g(X'')$，其中
$g:X''\to X'$ 是拟紧态射，并且 $f(E')=f(g(X''))$，所以 $f(E')$ 为
pro-可构造的，因为 $g\circ f$ 是拟紧的。

由 (vii) 可推出 (v) 的第二个断言：事实上，设 $E$ 是 $X$ 的局部闭且
逆紧的子集，并考虑 $X$ 的一个以 $E$ 为底空间的子预概形 (I, 5.2.1)；
典范注入 $j:E\to X$ 依假设是拟紧态射，故由 (vii)，$E=j(E)$ 在
$X$ 中为 pro-可构造的。

最后，为证明 (viii)，仍可假设 $X$ 仿射；此外，若 $(X'_\alpha)$ 是
$X'$ 的仿射开覆盖，并且 $E'\cap X'_\alpha$ 是 $X'_\alpha$ 的可构造子集之并
（(iii) 和 (iv)），则 $f(E')$ 是各 $f(E'\cap X'_\alpha)$ 之并；由 Chevalley
定理 (1.8.4)，每个 $f(E'\cap X'_\alpha)$ 都是 $X$ 的可构造子集之并，
结论遂得。
\end{proof}

\begin{env}[1.9.6]
\label{IV.1.9.6-fr}
\textbf{例} (1.9.6). --- 预概形 $X$ 的每个\emph{有限}子集都是
pro-可构造的：事实上，只须考虑只含一个点的子集 $\{x\}$
(1.9.5, (ii))；而 $\{x\}$ 是 $\operatorname{Spec}(k(x))$ 在典范态射
$\operatorname{Spec}(k(x))\to X$ 下的像，此态射拟紧；结论由
(1.9.5, (vii)) 得出。反之，只含一个点的子集 $\{x\}$ 未必是
ind-可构造的；例如，设 $A$ 是具有无穷多个极大理想的 Noether 整环，并设
$x$ 是 $X=\operatorname{Spec}(A)$ 的泛点；若 $\{x\}$ 在 $X$ 中为
ind-可构造的，则由 (1.9.5, (iii)) 它是可构造的，因而包含 $X$ 的一个
非空开集；但由 Artin--Tate 定理 (0, 16.3.3)，这与 $A$ 具有无穷多个
极大理想的假设矛盾。

每个闭子集 $Y$（属于预概形 $X$）由 (1.9.5, (v)) 都是 pro-可构造的。
因此 $X$ 的每个开子集都是 ind-可构造的；但开子集 $U$（属于 $X$）只有在
\emph{逆紧}时才是 pro-可构造的，这仍由 (1.9.5, (v)) 得出。

最后，若 $X$ 是 Noether 预概形，因而拟分离 (1.2.8)，则由
(1.9.5, (iii)) 和 (0\textsubscript{III}, 9.1.7)，$X$ 的 ind-可构造子集
正是 $X$ 的\emph{局部闭子集之并}。
\end{env}

\begin{theorem}[1.9.7]
\label{IV.1.9.7-fr}
设 $X$ 是拟紧预概形，$E$ 是 $X$ 的 ind-可构造子集，
$(F_\lambda)_{\lambda\in L}$ 是 $X$ 的一族 pro-可构造子集，并且
$\bigcap_{\lambda\in L}F_\lambda\subset E$；则存在有限子集 $J$（它是 $L$ 的子集），
使得 $\bigcap_{\lambda\in J}F_\lambda\subset E$。
\end{theorem}

\begin{proof}
注意子集 $F=X-E$ 和 $F\cap F_\lambda$ 都是 pro-可构造的，故可归结到
下述推论。
\end{proof}

\begin{corollary}[1.9.8]
\label{IV.1.9.8-fr}
设 $X$ 是拟紧预概形，$(F_\lambda)_{\lambda\in L}$ 是 $X$ 的一族
pro-可构造子集，且其交为空；则存在交为空的有限子族
$(F_\lambda)_{\lambda\in J}$。
\end{corollary}

\begin{proof}
以有限多个仿射开集覆盖 $X$，并使用 (1.9.5, (iv))，可归结到 $X$ 仿射的
情形；由 (1.9.5, (ix))，有 $F_\lambda=f_\lambda(X'_\lambda)$，

\oldpage[IV-1]{246}

其中 $X'_\lambda=\operatorname{Spec}(A'_\lambda)$ 是仿射的，而 $f_\lambda$
是态射 $X'_\lambda\to X$；于是采用 (1.9.3.2) 的记号，由假设有
$f(X')=\varnothing$，故 $X'=\varnothing$；但由 (1.9.2, (i))，这蕴含存在
有限子集 $J$（它是 $L$ 的子集），使得 $X'_J=\varnothing$，从而
$f_J(X'_J)=\bigcap_{\lambda\in J}F_\lambda=\varnothing$。
\end{proof}

\begin{corollary}[1.9.9]
\label{IV.1.9.9-fr}
设 $X$ 是拟紧预概形，$F$ 是 $X$ 的 pro-可构造子集，
$(E_\lambda)_{\lambda\in L}$ 是 $X$ 的一族 ind-可构造子集，并且
$F\subset\bigcup_{\lambda\in L}E_\lambda$；则存在有限子集 $J$（它是 $L$ 的子集），
使得 $F\subset\bigcup_{\lambda\in J}E_\lambda$。特别地，从任意由
ind-可构造子集组成的 $X$ 的覆盖中，都可抽取出 $X$ 的有限覆盖。
\end{corollary}

\begin{proof}
对补集应用 (1.9.7)，即得此结论。
\end{proof}

\begin{proposition}[1.9.10]
\label{IV.1.9.10-fr}
设 $X$ 是预概形。子集 $E$（属于 $X$）为 ind-可构造的一个必要条件是：
对每个 $x\in X$，交 $E\cap\overline{\{x\}}$ 是 $x$ 在
$\overline{\{x\}}$ 中的一个邻域。若 $X$ 局部 Noether，则此条件也充分。
\end{proposition}

\begin{proof}
假设 $E$ 为 ind-可构造的，并设 $Y$ 是 $\overline{\{x\}}$ 与 $X$ 中一个包含
$x$ 的仿射开集之交；因而存在 $X$ 的一个以 $Y$ 为底空间的子预概形
(I, 5.2.1)。若 $j:Y\to X$ 是典范注入，则 $E\cap Y=j^{-1}(E)$ 在 $Y$ 中
为 ind-可构造的 (1.9.5, (vi))。由于 $Y$ 拟紧且分离，$E\cap Y$ 是 $Y$ 的
可构造子集之并 (1.9.5, (iii))；因此存在开集 $U$、$V$（属于 $Y$），使得
$x\in U\cap\complement V\subset E\cap Y$
(0\textsubscript{III}, 9.1.3)。但由于 $x$ 是不可约空间 $Y$ 的泛点，$V$
必然为空，并且有 $U\subset E\cap Y$。

现假设 $X$ 局部 Noether，并设 $E$ 是 $X$ 的一个满足命题所述条件的子集。
由定义 (1.9.4)，可以归结到 $X$ 为 Noether 的情形。于是，对每个
$x\in E$，$E\cap\overline{\{x\}}$ 包含一个形如
$U\cap\overline{\{x\}}$ 的非空部分，其中 $U$ 在 $X$ 中开；这表明 $E$
是 $X$ 的局部闭子集之并，因而是 ind-可构造的 (1.9.6)。
\end{proof}

\begin{proposition}[1.9.11]
\label{IV.1.9.11-fr}
设 $X$ 是预概形，$E$ 是 $X$ 的一个子集。下列性质等价：
\begin{enumerate}
  \item[a)] $E$ 局部可构造。
  \item[b)] $E$ 既 ind-可构造又 pro-可构造。
  \item[c)] $E$ 与 $X-E$ 都 pro-可构造。
  \item[d)] $E$ 与 $X-E$ 都 ind-可构造。
\end{enumerate}
\end{proposition}

\begin{proof}
显然，只须证明 $d)$ 蕴含 $a)$，并且可以归结到 $X$ 仿射的情形。此时由
(1.9.5, (iii))，$E$（分别地，$X-E$）是某族 $(E_\lambda)$（分别地，
$(E'_\mu)$）可构造子集之并，这些子集属于 $X$；由于各 $E_\lambda$ 与各
$E'_\mu$ 构成 $X$ 的一个覆盖，由 (1.9.9) 存在有限多个指标 $\lambda_i$、
$\mu_j$，使得各 $E_{\lambda_i}$ 与各 $E'_{\mu_j}$ 构成 $X$ 的一个覆盖；
这蕴含 $E$ 是各 $E_{\lambda_i}$ 之并，因而可构造。
\end{proof}

\begin{corollary}[1.9.12]
\label{IV.1.9.12-fr}
设 $f:X\to Y$ 是满态射，它或者拟紧，或者局部有限表示。子集 $E$（属于
$Y$）为局部可构造的（分别地，pro-可构造的，分别地，ind-可构造的），
当且仅当 $f^{-1}(E)$ 在 $X$ 中具有相应性质。
\end{corollary}

\oldpage[IV-1]{247}

\begin{proof}
已知此条件是必要的（(1.8.2) 和 (1.9.5, (vi))）；为证明其充分性，由
(1.9.5, (iv)) 可归结到 $X$ 仿射的情形；此外，由 (1.9.11)，只须考虑
$f^{-1}(E)$ 为 pro-可构造的情形，\emph{或者} $f^{-1}(E)$ 为
ind-可构造的情形。若 $f$ 满且拟紧，并且 $f^{-1}(E)$ pro-可构造，则由
(1.9.5, (vii))，$E=f(f^{-1}(E))$ pro-可构造。若 $f$ 满且局部有限表示，
并且 $f^{-1}(E)$ ind-可构造，则由 (1.9.5, (viii))，
$E=f(f^{-1}(E))$ ind-可构造，证明完毕。
\end{proof}

\begin{env}[1.9.13]
\label{IV.1.9.13-fr}
设 $X$ 是预概形，$\mathscr T$ 是它的拓扑；由 (1.9.5, (i) 和 (ii))，
$X$ 的 ind-可构造子集（分别地，pro-可构造子集）是 $X$ 上某个拓扑的
开子集（分别地，闭子集）。这个拓扑称为\emph{可构造拓扑}，记为
$\mathscr T^{\mathrm{cons}}$；以 $X^{\mathrm{cons}}$ 表示集合 $X$ 配备拓扑
$\mathscr T^{\mathrm{cons}}$ 后所得的空间。
\end{env}

\begin{proposition}[1.9.14]
\label{IV.1.9.14-fr}
设 $X$ 是预概形，$\mathscr T$ 是它的拓扑，$\mathscr T^{\mathrm{cons}}$
是 $X$ 上的可构造拓扑。
\begin{enumerate}
  \item[(i)] 拓扑 $\mathscr T^{\mathrm{cons}}$ 比 $\mathscr T$ 更细。
  \item[(ii)] $X$ 的局部可构造子集恰好是空间 $X^{\mathrm{cons}}$ 中既开又闭的子集。
  \item[(iii)] 对每个态射 $f:X\to Y$，从 $X^{\mathrm{cons}}$ 到
  $Y^{\mathrm{cons}}$ 的底层映射连续；记为 $f^{\mathrm{cons}}$。
  \item[(iv)] 若态射 $f:X\to Y$ 拟紧，则 $f^{\mathrm{cons}}$ 是闭映射；
  特别地，若 $f$ 拟紧且双射，则 $f^{\mathrm{cons}}$ 是同胚。
  \item[(v)] 若态射 $f:X\to Y$ 局部有限表示，则 $f^{\mathrm{cons}}$ 是开映射；
  特别地，若 $f$ 双射且局部有限表示，则 $f^{\mathrm{cons}}$ 是同胚。
  \item[(vi)] 对每个开集 $U$（属于 $X$），$\mathscr T^{\mathrm{cons}}$ 在 $U$ 上
  诱导的拓扑等于 $U^{\mathrm{cons}}$ 的拓扑。
\end{enumerate}
\end{proposition}

\begin{proof}
事实上，(i) 来自如下事实：$\mathscr T$ 的每个开集都是
$\mathscr T^{\mathrm{cons}}$ 的开集 (1.9.6)；(ii) 是 (1.9.11) 的改写；
(iii)、(iv) 和 (v) 分别改写了 (1.9.5, (vi)、(vii) 和 (viii))。最后，为证明
(vi)，只须注意典范注入 $j:U\to X$ 局部有限表示 (1.4.3)，故
$j^{\mathrm{cons}}:U^{\mathrm{cons}}\to X^{\mathrm{cons}}$ 是连续开映射。
\end{proof}

\begin{proposition}[1.9.15]
\label{IV.1.9.15-fr}
设 $X$ 是预概形。
\begin{enumerate}
  \item[(i)] $X^{\mathrm{cons}}$ 拟紧，当且仅当 $X$ 拟紧。
  \item[(ii)] $X^{\mathrm{cons}}$ 分离，当且仅当 $X$ 拟分离；此时
  $X^{\mathrm{cons}}$ 局部紧且全不连通。
  \item[(iii)] $X^{\mathrm{cons}}$ 紧，当且仅当 $X$ 拟紧且拟分离。
  \item[(iv)] $X$ 的每个点对于拓扑 $\mathscr T^{\mathrm{cons}}$ 都有一个
  既开又紧的邻域。
  \item[(v)] 态射 $f:X\to Y$ 拟紧，当且仅当连续映射
  $f^{\mathrm{cons}}:X^{\mathrm{cons}}\to Y^{\mathrm{cons}}$ 是固有映射。
\end{enumerate}
\end{proposition}

\begin{proof}
(i) 由于 $\mathscr T^{\mathrm{cons}}$ 比 $\mathscr T$ 更细，显然若
$X^{\mathrm{cons}}$ 拟紧，则 $X$ 也拟紧；反命题由 (1.9.9) 得出。

(ii) 假设 $X$ 拟分离，并证明 $X^{\mathrm{cons}}$ 全不连通：事实上，若
$x$、$y$ 是 $X$ 的两个不同点，并且例如
$x\notin\overline{\{y\}}$，则存在一个
\oldpage[IV-1]{248}
仿射开邻域 $U$（属于 $X$），它不含 $y$；由于 $X$ 拟分离，$U$ 在 $X$ 中
逆紧 (1.2.7)，故 $U$ 与 $X-U$ 在 $X$ 中局部可构造，因而由 (1.9.11)
在 $X^{\mathrm{cons}}$ 中开，断言遂得。对每个仿射开集 $U$（属于 $X$），
$\mathscr T^{\mathrm{cons}}$ 在其上诱导的拓扑由 (1.9.14, (vi)) 等于
$U^{\mathrm{cons}}$ 的拓扑；由上述结果，$X^{\mathrm{cons}}$ 局部紧，因为
$U$ 在 $X^{\mathrm{cons}}$ 中开，而 $U^{\mathrm{cons}}$ 紧。尚须证明：若
$X^{\mathrm{cons}}$ 分离，则 $X$ 拟分离。事实上，取一个仿射开集 $U$
（属于 $X$）；$U$ 上由 $\mathscr T^{\mathrm{cons}}$ 诱导的拓扑等于
$U^{\mathrm{cons}}$ 的拓扑 (1.9.14, (vi))，故 $U$ 对此诱导拓扑是紧的，
因为论证的前半部分已表明 $U^{\mathrm{cons}}$ 紧。若 $V$ 是 $X$ 中第二个
仿射开集，则 $U\cap V$ 是分离空间 $X^{\mathrm{cons}}$ 的紧子集，因为它是
此空间中两个紧子集的交；由于 $\mathscr T^{\mathrm{cons}}$ 在 $U\cap V$ 上
诱导的拓扑等于 $(U\cap V)^{\mathrm{cons}}$ 的拓扑 (1.9.14, (vi))，并且恒等映射
$(U\cap V)^{\mathrm{cons}}\to U\cap V$ 连续，故 $U\cap V$ 对
$\mathscr T$ 诱导的拓扑拟紧；于是由 (1.2.7)，$X$ 拟分离。

(iii) 是 (i) 与 (ii) 的直接推论。

(iv) 对每个 $x\in X$，一个仿射开邻域 $U$（它是 $x$ 对 $\mathscr T$ 的邻域）
也是 $x$ 对 $\mathscr T^{\mathrm{cons}}$ 的邻域，并且由 (iii) 和
(1.9.14, (vi))，$U$ 是紧的，因为 $\mathscr T^{\mathrm{cons}}$ 在其上诱导的
拓扑等于 $U^{\mathrm{cons}}$ 的拓扑。

(v) 假设 $f$ 拟紧；则已知 (1.9.14, (iv)) $f^{\mathrm{cons}}$ 是闭映射。
另一方面，设 $y$ 是 $Y$ 的一点，并置
\[
  Z=f^{-1}(y)=X\times_Y\operatorname{Spec}(k(y));
\]
$Z$ 拟紧 (1.1.2, (iii))；又因典范态射 $p:Z\to X$ 为单射，由 (i) 以及
$p^{\mathrm{cons}}$ 连续可知，$f^{-1}(y)$ 上由 $X^{\mathrm{cons}}$ 的拓扑
诱导的拓扑使 $f^{-1}(y)$ 成为拟紧空间。这证明 $f^{\mathrm{cons}}$ 是
固有映射（Bourbaki, \emph{Top. gén.}, chap. I\textsuperscript{er},
3\textsuperscript{e} éd., \S~10, n\textsuperscript{o}~2, th.~1）。反过来，假设
连续映射 $f^{\mathrm{cons}}$ 固有，并设 $V$ 是 $Y$ 的拟紧开集；若
$h:V\to Y$ 是典范注入，则 $h^{\mathrm{cons}}:V^{\mathrm{cons}}\to Y^{\mathrm{cons}}$
连续且单射，并且由 (i)，$V^{\mathrm{cons}}$ 拟紧，故 $V$ 上由
$Y^{\mathrm{cons}}$ 的拓扑诱导的拓扑使 $V$ 成为拟紧空间。$f^{\mathrm{cons}}$ 固有这一
假设继而蕴含：$f^{-1}(V)$ 上由 $X^{\mathrm{cons}}$ 的拓扑诱导的拓扑使
$f^{-1}(V)$ 成为拟紧空间（\emph{loc. cit.}, prop.~6），故 $f^{-1}(V)$ 也是
$X$ 的拟紧子空间，这表明态射 $f$ 拟紧。
\end{proof}

\begin{env}[1.9.16]
\label{IV.1.9.16-fr}
稍后我们将说明，对每个预概形 $X$，如何给空间 $X^{\mathrm{cons}}$ 配备环层，
使其成为一个局部环全为域的预概形；这些域就是 $X$ 各点的剩余域。例如，
在用概形语言转述 Néron [32] 的 Abel 簇约化理论中的构造时，这类预概形会
自然出现。
\end{env}

\oldpage[IV-1]{249}

\subsection{开态射的应用}
\label{subsection:IV.1.10-fr}

\begin{theorem}[1.10.1]
\label{IV.1.10.1-fr}
设 $X$ 是预概形，$E$ 是 $X$ 的 ind-可构造子集 (1.9.4)，$x$ 是 $X$ 的一点。
$x$ 是 $E$ 的内点，当且仅当每个一般化点 ($0_{\mathrm I}$, 2.1.2)
$x'$（它是 $x$ 的一般化）都属于 $E$；换言之 (I, 2.4.2)，当且仅当
$\operatorname{Spec}(\mathcal O_{X,x})\subset E$。
\end{theorem}

\begin{proof}
必要性显然，因为 $x$ 的每个邻域都包含 $x$ 的一般化点。为证明充分性，显然
可以归结到 $X=\operatorname{Spec}(A)$ 仿射的情形，而 $x$ 是素理想
$\mathfrak p$（属于 $A$）。此时由 (1.9.5, (ix))，存在仿射概形
$X'=\operatorname{Spec}(A')$ 和态射 $f:X'\to X$，使得
$f(X')=X-E$。置 $Y=\operatorname{Spec}(\mathcal O_{X,x})$ 以及
$Y'=X'\times_XY$；由 (I, 3.6.5)，假设意味着 $Y'=\varnothing$，换言之
$A'\otimes_AA_{\mathfrak p}=0$。由于
$A_{\mathfrak p}=\varinjlim A_t$，其中 $t$ 遍历 $A-\mathfrak p$
($0_{\mathrm I}$, 1.4.5)，又有 $\varinjlim A_t'=0$，因为函子
$\varinjlim$ 与张量积可交换。于是由 (1.9.1)，存在 $t\in A-\mathfrak p$，使得
$A_t'=0$；此时 $\mathrm D(t)$ 是 $x$ 在 $X$ 中的开邻域，并包含于 $E$。
\end{proof}

\begin{corollary}[1.10.2]
\label{IV.1.10.2-fr}
设 $f:X\to Y$ 是局部有限表示态射，并设 $y\in Y$。$y$ 是 $f(X)$ 的内点，
当且仅当每个一般化点 $y'$（它是 $y$ 的一般化）都属于 $f(X)$。
\end{corollary}

\begin{proof}
事实上，已知 (1.9.5, (viii)) $f(X)$ 在 $Y$ 中 ind-可构造，只须应用
(1.10.1)。
\end{proof}

称连续映射 $f:X\to Y$ 在一点 $x\in X$ 处\emph{为开映射}，如果 $f$ 将
$x$ 在 $X$ 中的\emph{每个}邻域都映成 $f(x)$ 在 $Y$ 中的邻域。

\begin{proposition}[1.10.3]
\label{IV.1.10.3-fr}
设 $f:X\to Y$ 是态射，$x$ 是 $X$ 的一点，$y=f(x)$。考虑下列条件：
\begin{enumerate}
  \item[a)] $f$ 在点 $x$ 处为开映射。
  \item[b)] 对每个一般化点 $y'$（它是 $y$ 的一般化），存在一般化点 $x'$（它是 $x$ 的一般化），
  使得 $f(x')=y'$。
  \item[b')] $f$ 在 $\operatorname{Spec}(\mathcal O_{X,x})$ 上的像是
  $\operatorname{Spec}(\mathcal O_{Y,y})$。
  \item[c)] 对每个不可约闭子集 $Y'$（属于 $Y$ 且包含 $y$），
  $X'=f^{-1}(Y')$ 有一个包含 $x$ 且支配 $Y'$ 的不可约分支。
\end{enumerate}
则有蕴含关系
$\mathrm{a)}\Rightarrow\mathrm{b)}\Leftrightarrow\mathrm{b')}
\Leftrightarrow\mathrm{c)}$。此外，若 $f$ 局部有限表示，则四个条件等价。
\end{proposition}

\begin{proof}
由一般化点的定义，$f$ 在 $\operatorname{Spec}(\mathcal O_{X,x})$ 上的像总包含于
$\operatorname{Spec}(\mathcal O_{Y,y})$，故 b) 与 b') 等价。c) 蕴含 b) 是立即的：
若 $y'$ 是 $y$ 的一个一般化点，则 $Y'=\overline{\{y'\}}$ 是 $Y$ 的一个包含
$y$ 的不可约子集；若 $x'$ 是 $f^{-1}(Y')$ 的某个包含 $x$ 且支配 $Y'$ 的
不可约分支的泛点，则 $f(x')=y'$ ($0_{\mathrm I}$, 2.1.5)，并且 $x'$ 是 $x$ 的
一般化点。反过来，b) 蕴含 c)：事实上，设 $y'$ 是 $Y'$ 的泛点，并设 $x'$
是 $x$ 的一个一般化点，使得 $f(x')=y'$；设 $Z'$ 是 $f^{-1}(Y')$ 的一个包含
$x'$（因而包含 $x$）的不可约分支；由于 $f(x')$ 已经在 $Y'$ 中稠密，
$f(Z')$ 更是如此。

为证明 a) 蕴含 b')，显然可以归结到 $X=\operatorname{Spec}(B)$ 和
$Y=\operatorname{Spec}(A)$ 都仿射的情形，其中 $x$ 是素理想
$\mathfrak q$（属于 $B$），从而
$\mathcal O_{X,x}=B_{\mathfrak q}=\varinjlim B_t$，
\oldpage[IV-1]{250}
其中 $t$ 遍历 $B-\mathfrak q$。由 (1.9.2)，有
\[
  f(\operatorname{Spec}(\mathcal O_{X,x}))
  =\bigcap_t f(\operatorname{Spec}(B_t))=\bigcap_t f(\mathrm D(t)),
\]
并且由假设，对每个 $t\in B-\mathfrak q$，$y$ 是 $f(\mathrm D(t))$ 的内点，
故 $f(\mathrm D(t))$ 包含 $\operatorname{Spec}(\mathcal O_{Y,y})$；于是
$\operatorname{Spec}(\mathcal O_{Y,y})
\subset f(\operatorname{Spec}(\mathcal O_{X,x}))$，这证明了断言。最后，若 $f$
局部有限表示，则由 (1.10.2) 应用于限制 $U\to Y$；这是 $f$ 在任意开邻域
$U$（它是 $x$ 的邻域）上的限制，故 b) 蕴含 a)。
\end{proof}

\begin{corollary}[1.10.4]
\label{IV.1.10.4-fr}
设 $f:X\to Y$ 是态射。考虑下列条件：
\begin{enumerate}
  \item[a)] $f$ 是开映射。
  \item[b)] 对每个 $x\in X$ 以及每个一般化点 $y'$（它是 $y=f(x)$ 的一般化），
  存在一般化点 $x'$（它是 $x$ 的一般化），使得 $f(x')=y'$。
  \item[b')] 对每个 $x\in X$，$f$ 在
  $\operatorname{Spec}(\mathcal O_{X,x})$ 上的像是
  $\operatorname{Spec}(\mathcal O_{Y,f(x)})$。
  \item[c)] 对每个不可约闭子集 $Y'$（属于 $Y$），$f^{-1}(Y')$ 的每个不可约分支
  都支配 $Y'$。
\end{enumerate}
则有蕴含关系
$\mathrm{a)}\Rightarrow\mathrm{b)}\Leftrightarrow\mathrm{b')}
\Leftrightarrow\mathrm{c)}$。此外，若 $f$ 局部有限表示，则四个条件等价。
\end{corollary}

\begin{proof}
说 $f$ 是开映射，就是说 $f$ 在每个点 $x\in X$ 处为开映射；因此蕴含关系
$\mathrm{a)}\Rightarrow\mathrm{b)}\Leftrightarrow\mathrm{b')}$ 来自 (1.10.3) 中
相应的蕴含关系；蕴含关系 $\mathrm{b)}\Rightarrow\mathrm{a)}$ 在 $f$ 局部有限表示时
也同样如此。蕴含关系
$\mathrm{c)}\Rightarrow\mathrm{b)}$ 也来自 (1.10.3) 中相应的蕴含关系；最后证明
b) 蕴含 c)。事实上，设 $x'$ 是 $f^{-1}(Y')$ 的某个不可约分支的泛点，并证明
$y'=f(x')$ 是 $Y'$ 的泛点。设 $y''$ 是 $y'$ 的一个一般化点；由假设，存在
一般化点 $x''$（它是 $x'$ 的一般化），使得 $f(x'')=y''$；又因
$x''\in f^{-1}(Y')$，必有 $x''=x'$，从而 $y''=y'$，证明完毕。
\end{proof}

\begin{center}
\emph{（待续。）}
\end{center}

\clearpage
\oldpage[IV-1]{251}
\phantomsection
\addcontentsline{toc}{section}{书目（续）}
\section*{书目（\emph{续}）}

\begingroup
\small
\setlength{\parindent}{0pt}
\setlength{\parskip}{3pt}
[30] M.~Nagata, \emph{Local rings}, \emph{Interscience Tracts}, 13 (1962),
Interscience, New York.

[31] M.~Nagata, \emph{A Jacobian criterion of simple points},
\emph{Ill. Journ. of Math.}, 1 (1957), p.~427--432.

[32] A.~Néron, \emph{Modèles minimaux des variétés abéliennes sur les corps
locaux et globaux}, \emph{Publ. Inst. Hautes Études Sci.}, n\textsuperscript{o}
21 (1964).

\clearpage
\oldpage[IV-1]{252}
\phantomsection
\addcontentsline{toc}{section}{记号索引}
\section*{记号索引}

$\dim_c(X)$，$\dim(X)$（$X$ 为拓扑空间）：0, 14.1.2.\par
$\dim_x(X)$（$X$ 为拓扑空间，$x$ 为 $X$ 的点）：0, 14.1.2.\par
$\operatorname{codim}(Y,X)$（$X$ 为拓扑空间，$Y$ 为 $X$ 的闭子集）：0, 14.2.1.\par
$\operatorname{codim}_x(Y,X)$（$X$ 为拓扑空间，$Y$ 为 $X$ 的闭子集，$x$ 为 $X$ 的点）：0, 14.2.4.\par
$N[T_1,\ldots,T_n]$（$N$ 为 $A$-模）：0, 15.1.1.\par
$\dim(A)$（$A$ 为环）：0, 16.1.1.\par
$\operatorname{ht}(\mathfrak a)$（$\mathfrak a$ 为理想）：0, 16.1.3.\par
$\dim_A(M)$，$\dim(M)$（$M$ 为 $A$-模）：0, 16.1.7.\par
$P_{\mathfrak q}(M,n)$（$M$ 是 Noether 半局部环 $A$ 上的有限型模，$\mathfrak q$ 是 $A$ 的定义理想）：0, 16.2.1.\par
$\operatorname{prof}_A(M)$，$\operatorname{prof}(M)$（$M$ 为 $A$-模）：0, 16.4.5.\par
$\operatorname{coprof}_A(M)$，$\operatorname{coprof}(M)$（$M$ 为 $A$-模）：0, 16.4.9.\par
$\operatorname{dim.proj}_A(M)$，$\operatorname{dim.proj}(M)$，
$\operatorname{dim.inj}_A(M)$，$\operatorname{dim.inj}(M)$（$M$ 为 $A$-模）：0, 17.2.1.\par
$\operatorname{dim.coh}(A)$（$A$ 为环）：0, 17.2.8.\par
$\operatorname{dim.proj}(\mathscr F)$，$\operatorname{dim.inj}(\mathscr F)$
（$\mathscr F$ 为环化空间上的模）：0, 17.2.14.\par
$\operatorname{dim.coh}(X)$（$X$ 为环化空间）：0, 17.2.14.\par
$E\times_B F$（$B$、$E$、$F$ 为 $A$-环）：0, 18.1.2.\par
$D_B(L)$（$B$ 为环，$L$ 为 $B$-双模）：0, 18.2.3.\par
$Y\oplus_X Z$（$X$、$Y$、$Z$ 为 $A$-双模）：0, 18.2.7.\par
$\operatorname{Exan}_A(B,L)$（$B$ 为 $A$-环，$L$ 为 $B$-双模）：0, 18.3.4.\par
$\operatorname{Exan}_{A/A'}(B,L)$：0, 18.3.7.\par
$\operatorname{Exal}_A(B,L)$，$\operatorname{Exal}_{A/A'}(B,L)$（$A$ 为交换环，$B$ 为 $A$-代数，$L$ 为 $B$-双模）：0, 18.4.1.\par
$\operatorname{Exalcom}_A(B,L)$，$\operatorname{Exalcom}_{A/A'}(B,L)$（$A$ 为交换环，$B$ 为交换 $A$-代数，$L$ 为 $B$-模）：0, 18.4.2.\par
$\mathrm H_A^2(B,L)$，$\mathrm H_A^2(B,L)^s$：0, 18.4.3.\par
$\operatorname{Exantop}_A(B,L)$，$\operatorname{Exaltop}_A(B,L)$，
$\operatorname{Exalcotop}_A(B,L)$：0, 18.5.1.\par
$\operatorname{Exantop}_{A/A'}(B,L)$，$\operatorname{Exaltop}_{A/A'}(B,L)$，
$\operatorname{Exalcotop}_{A/A'}(B,L)$：0, 18.5.2.\par
$\operatorname{Hom.\ cont}_C(M,N)$（$M$、$N$ 为拓扑 $C$-双模）：0, 18.5.3.\par
$\operatorname{Dér}_A(B,L)$（$B$ 为 $A$-环，$L$ 为 $B$-双模）：0, 20.1.2.\par
$\operatorname{Dér.\ cont}_A(B,L)$：0, 20.3.1.\par
$p_{B/A}$，$\mathfrak J_{B/A}$（$B$ 为交换拓扑 $A$-代数）：0, 20.4.1.\par
$P^1_{B/A}$，$\varepsilon_{B/A}$（$B$ 为交换拓扑 $A$-代数）：0, 20.4.2.\par
$\Omega^1_{B/A}$，$\widehat{\Omega}^1_{B/A}$（$B$ 为交换拓扑 $A$-代数）：0, 20.4.3.\par
$\Omega^1_B$（$B$ 为拓扑环）：0, 20.4.3.\par
$d_{B/A}$，$d(x)$，$dx$，$d_B$：0, 20.4.6.\par
$\Upsilon_{C/B/A}$，$\Upsilon_{C/B}$（$A$、$B$、$C$ 为环）：0, 20.6.1.\par
$\chi_E$（$E$ 是一个 $B$-扩张，它在 $A$ 上平凡，并且是 $C$ 由 $L$ 给出的扩张）：0, 20.6.8.\par

\clearpage
\oldpage[IV-1]{253}
$\Upsilon^C_{B/A/\Lambda}$，$\Upsilon^C_{B/A}$（$\Lambda$、$A$、$B$、$C$ 为环）：0, 20.6.14.\par
$\chi_{B/A}$，$\chi_B$（$B$ 为 $A$-代数）：0, 20.6.24.\par
$\chi_E$（$E$ 为拓扑 $B$-代数）：0, 20.7.10.\par
$F_A$，$A^p$，$A^{(p)}$，$E^{(p)}$（$A$ 为特征 $p>0$ 的环，$E$ 为 $A$-模）：0, 21.1.4.\par
$\pi_{B/A}$，$\Theta_{B/A}$，$\Xi_{B/A}$（$A$、$B$ 为特征 $p>0$ 的环）：0, 21.3.2.\par
$\Delta(L/K,k/k_0)$，$d(L/K,k/k_0)$，$\Delta(L/K,k)$，$d(L/K,k)$（$k_0$、$k$、$K$、$L$ 为域）：0, 21.6.1.\par
$\chi'_{B/A}$（$A$、$B$ 为特征 $p>0$ 的环）：0, 22.4.6.\par
$\mathscr T^{\mathrm{cons}}$，$X^{\mathrm{cons}}$（$X$ 为预概形，$\mathscr T$ 为 $X$ 的拓扑）：IV, 1.9.12.\par
$f^{\mathrm{cons}}$（$f$ 为预概形态射）：IV, 1.9.13.\par

\clearpage
\oldpage[IV-1]{254}
\phantomsection
\addcontentsline{toc}{section}{术语索引}
\section*{术语索引}

一阶主部增广代数：0, 20.4.2.\par
Cohen $A$-代数：0, 19.8.1.\par
有限表示代数：IV, 1.4.1.\par
本质有限型 $A$-代数：IV, 1.3.8.\par
$A$-环：0, 18.1.1.\par
$A$-环，增广于 $B$：0, 18.1.4.\par
双等维环：0, 16.1.4.\par
悬链环：0, 16.1.4.\par
特征 $p$ 的环：0, 21.1.1.\par
Cohen 环：0, 19.8.4.\par
Cohen--Macaulay 环：0, 16.5.1 和 16.5.13.\par
等余维环、等维环：0, 16.1.4.\par
几何单支局部环：0, 23.1.7.\par
日本环：0, 23.1.1.\par
素局部环、完备素局部环：0, 19.8.3.\par
正则环：0, 17.1.1 和 17.3.6.\par
单支局部环：0, 23.1.7.\par
泛日本环：0, 23.1.1.\par
在一点处为开映射：IV, 1.10.2.\par
增广：0, 18.1.4.\par

形式双态射：0, 19.1.2.\par

有序集中的链、拓扑空间的不可约闭子集链：0, 14.1.1.\par
饱和链：0, 14.3.1.\par
组合余维数、拓扑空间闭子集的余维数：0, 14.2.1.\par
拓扑空间闭子集在一点处的余维数：0, 14.2.4.\par
链条件：0, 14.3.2.\par
Noether 局部环上模的余深度：0, 16.4.9.\par
$k_0$-容许域、扩张的容许域：0, 21.6.1.\par
相对于子域的有限根重数域：0, 19.6.6.\par

扩张的 $k_0$-容许性亏格：0, 21.6.1.\par
导子、$A$-导子：0, 20.1.2.\par
微分、一阶微分、绝对微分：0, 20.4.3.\par
外微分：0, 20.4.6.\par
组合维数、拓扑空间的维数：0, 14.1.2.\par
拓扑空间在一点处的维数：0, 14.1.2.\par
Krull 维数、环的维数：0, 16.1.1.\par
模的维数：0, 16.1.7.\par
上同调维数、环的整体上同调维数：0, 17.2.8.\par
环化空间在一点处的上同调维数、环化空间的上同调维数：0, 17.2.14.\par

\clearpage
\oldpage[IV-1]{255}
模的内射维数、投射维数：0, 17.2.1.\par
在一点处的内射维数、层的内射维数：0, 17.2.14.\par
在一点处的投射维数、层的投射维数：0, 17.2.14.\par

$A$ 中的元素，它在某个 $A$-模中不是零因子：0, 15.1.1.\par
$M$-正则元素、$M$-拟正则元素：0, 15.1.4.\par
ind-可构造集、pro-可构造集：IV, 1.9.4.\par
形式满态射：0, 19.1.2.\par
$A$-扩张的 $A$-等价：0, 18.2.2.\par
双等维空间：0, 14.3.3.\par
悬链空间：0, 14.3.2.\par
等余维空间：0, 14.2.1.\par
等维空间：0, 14.1.3.\par
一个 $A$-扩张：它扩张的 $A$-环为 $B$，所用模是 $B$-双模：0, 18.2.2.\par
经由双模同态从另一扩张导出的扩张：0, 18.2.8.\par
Hochschild 扩张：0, 18.4.3.\par
$B$-扩张，$A$-平凡：0, 18.3.7.\par
典型平凡扩张：0, 18.2.3.\par
$A$-等价的 $A$-扩张：0, 18.2.2.\par

$p$-自由族、绝对 $p$-自由族：0, 21.1.9.\par
形式平 étale（代数）：0, 19.10.2.\par
形式右可逆（同态）：0, 19.1.15.\par
形式左可逆（同态）：0, 19.1.5.\par
形式非分歧（代数）：0, 19.10.2.\par
形式投射（模）：0, 19.2.1.\par
形式平滑（代数）：0, 19.3.1.\par
相对于一个环形式平滑（代数）：0, 19.9.1.\par

在域上几何正则（局部环）：0, 19.6.5.\par
$A$-扩张类群：0, 18.3.4.\par

理想的高度：0, 16.1.3.\par
$A$-环的 $A$-同态：0, 18.1.1.\par
$B$-扩张（$A$-平凡）的特征同态：0, 20.6.8.\par
相对于一个环与一个理想的 $A$-代数特征同态：0, 20.6.24.\par
特征 $p>0$ 代数的特征同态：0, 22.4.6.\par

增广理想：0, 18.1.4.\par
增广环的逆像：0, 18.1.5.\par
$A$-扩张的逆像：0, 18.2.5.\par

链的长度：0, 14.1.1.\par

双等维、悬链、等余维、等维模：0, 16.1.7.\par
Cohen--Macaulay 模：0, 16.5.1 和 16.5.13.\par
不可完备度模：0, 20.6.1.\par
形式单态射：0, 19.1.2.\par
$A$-扩张的态射：0, 18.2.4.\par
在一点处有限表示的态射、局部有限表示态射：IV, 1.4.2.\par
有限表示态射：IV, 1.6.1.\par
拟分离态射：IV, 1.2.1.\par

Cohen $p$-环：0, 19.8.4.\par
$p$-基、绝对 $p$-基：0, 21.1.9.\par
预概形的极大点：IV, 1.1.4.\par

\clearpage
\oldpage[IV-1]{256}
Hilbert--Samuel 多项式：0, 16.2.1.\par
在另一预概形上有限表示的预概形：IV, 1.6.1.\par
拟分离预概形：IV, 1.2.1.\par
$A$-环的纤维积：0, 18.1.2.\par
模的深度：0, 16.4.5.\par

双模的拼合和：0, 18.2.7.\par
严格形式投射（模）：0, 19.2.3.\par
$M$-正则序列、$M$-拟正则序列：0, 15.1.7.\par
$\mathscr F$-正则序列、$\mathscr F$-拟正则序列：0, 15.2.2.\par
Noether 半局部环上模的参数系：0, 16.3.6.\par
$p$-生成元系、绝对 $p$-生成元系：0, 21.1.9.\par
开理想、开子模的基本系：0, 19.0.2.\par
正则局部环的正则参数系：0, 17.1.6.\par

预概形上的可构造拓扑：IV, 1.9.12.\par
由 $A$ 的拓扑导出的 $A$-模拓扑：0, 19.0.2.\par
平凡（增广 $A$-环）：0, 18.1.4.\par
$A$-平凡（扩张）：0, 18.2.3 和 18.3.7.\par

\clearpage
\oldpage[IV-1]{257}
\phantomsection
\addcontentsline{toc}{section}{目录}
\section*{目录}

\newcommand{\ivonecontentsline}[3]{\noindent\makebox[17mm][l]{#1}\parbox[t]{0.77\linewidth}{#2\dotfill}\makebox[12mm][r]{#3}\par}
\ivonecontentsline{第 0 章}{\textbf{预备知识}（\emph{续}）}{5}
\ivonecontentsline{\S~14.}{拓扑空间的组合维数}{6}
\ivonecontentsline{14.1.}{拓扑空间的组合维数}{6}
\ivonecontentsline{14.2.}{闭子集的余维数}{8}
\ivonecontentsline{14.3.}{链条件}{10}
\ivonecontentsline{\S~15.}{$M$-正则序列与 $\mathscr F$-正则序列}{12}
\ivonecontentsline{15.1.}{$M$-正则序列与 $M$-拟正则序列}{12}
\ivonecontentsline{15.2.}{$\mathscr F$-正则序列}{20}
\ivonecontentsline{\S~16.}{Noether 局部环中的维数与深度}{22}
\ivonecontentsline{16.1.}{环的维数}{22}
\ivonecontentsline{16.2.}{Noether 半局部环的维数}{25}
\ivonecontentsline{16.3.}{Noether 局部环中的参数系}{28}
\ivonecontentsline{16.4.}{深度与余深度}{32}
\ivonecontentsline{16.5.}{Cohen--Macaulay 模}{36}
\ivonecontentsline{\S~17.}{正则环}{39}
\ivonecontentsline{17.1.}{正则环的定义}{39}
\ivonecontentsline{17.2.}{模的投射维数与内射维数回顾}{42}
\ivonecontentsline{17.3.}{正则环的上同调理论}{46}
\ivonecontentsline{\S~18.}{代数扩张的补充}{51}
\ivonecontentsline{18.1.}{增广环的逆像}{51}
\ivonecontentsline{18.2.}{环由双模给出的扩张}{54}
\ivonecontentsline{18.3.}{$A$-扩张类群}{59}
\ivonecontentsline{18.4.}{代数的扩张}{64}
\ivonecontentsline{18.5.}{拓扑环的情形}{66}
\ivonecontentsline{\S~19.}{形式平滑代数与 Cohen 环}{69}
\ivonecontentsline{19.0.}{引言}{69}
\ivonecontentsline{19.1.}{形式满态射与形式单态射}{71}
\ivonecontentsline{19.2.}{形式投射模}{78}

\clearpage
\oldpage[IV-1]{258}
\ivonecontentsline{19.3.}{形式平滑代数}{79}
\ivonecontentsline{19.4.}{形式平滑性的初步判别准则}{86}
\ivonecontentsline{19.5.}{形式平滑性与相伴分次环}{90}
\ivonecontentsline{19.6.}{域上代数的情形}{100}
\ivonecontentsline{19.7.}{局部同态的情形：存在性与唯一性定理}{104}
\ivonecontentsline{19.8.}{Cohen 代数与 Cohen $p$-环；对完备局部环结构的应用}{109}
\ivonecontentsline{19.9.}{相对形式平滑代数}{114}
\ivonecontentsline{19.10.}{形式非分歧代数与形式平 étale 代数}{115}
\ivonecontentsline{\S~20.}{导子与微分}{116}
\ivonecontentsline{20.1.}{导子与代数扩张}{117}
\ivonecontentsline{20.2.}{导子的函子性质}{119}
\ivonecontentsline{20.3.}{拓扑环中的连续导子}{121}
\ivonecontentsline{20.4.}{主部与微分}{123}
\ivonecontentsline{20.5.}{$\Omega^1_{B/A}$ 的基本函子性质}{128}
\ivonecontentsline{20.6.}{不可完备度模与特征同态}{136}
\ivonecontentsline{20.7.}{推广到拓扑环上}{147}
\ivonecontentsline{\S~21.}{特征 $p$ 环中的微分}{153}
\ivonecontentsline{21.1.}{$p$-生成元系与 $p$-基}{154}
\ivonecontentsline{21.2.}{$p$-基与形式平滑性}{157}
\ivonecontentsline{21.3.}{$p$-基与不可完备度模}{160}
\ivonecontentsline{21.4.}{域扩张的情形}{162}
\ivonecontentsline{21.5.}{应用：可分性判别准则}{164}
\ivonecontentsline{21.6.}{扩张的容许域}{167}
\ivonecontentsline{21.7.}{Cartier 等式}{169}
\ivonecontentsline{21.8.}{容许性判别准则}{171}
\ivonecontentsline{21.9.}{形式幂级数环中的完备化微分模}{176}
\ivonecontentsline{\S~22.}{形式平滑性与正则性的微分判别准则}{182}
\ivonecontentsline{22.1.}{形式平滑性的提升}{183}
\ivonecontentsline{22.2.}{域上形式平滑局部代数的微分刻画}{186}
\ivonecontentsline{22.3.}{对若干局部环及其完备化之间关系的应用}{191}
\ivonecontentsline{22.4.}{极大理想平方为零的局部环有限扩张的初步结果}{193}

\clearpage
\oldpage[IV-1]{259}
\ivonecontentsline{22.5.}{几何正则代数与形式平滑代数}{201}
\ivonecontentsline{22.6.}{Zariski 的 Jacobian 判别准则}{205}
\ivonecontentsline{22.7.}{Nagata 的 Jacobian 判别准则}{209}
\ivonecontentsline{\S~23.}{日本环}{213}
\ivonecontentsline{23.1.}{日本环}{213}
\ivonecontentsline{23.2.}{Noether 整局部环的整闭包}{217}
\ivonecontentsline{}{\textsc{第 IV 章}——\textbf{概形及概形态射的局部研究}}{222}
\ivonecontentsline{\S~1.}{相对有限性条件。预概形中的可构造集}{224}
\ivonecontentsline{1.1.}{拟紧态射}{224}
\ivonecontentsline{1.2.}{拟分离态射}{226}
\ivonecontentsline{1.3.}{局部有限型态射}{228}
\ivonecontentsline{1.4.}{局部有限表示态射}{230}
\ivonecontentsline{1.5.}{有限型态射}{233}
\ivonecontentsline{1.6.}{有限表示态射}{234}
\ivonecontentsline{1.7.}{先前结果的改进}{236}
\ivonecontentsline{1.8.}{有限表示态射与可构造集}{238}
\ivonecontentsline{1.9.}{pro-可构造集与 ind-可构造集}{241}
\ivonecontentsline{1.10.}{开态射的应用}{249}
\ivonecontentsline{}{书目（\emph{续}）}{251}
\ivonecontentsline{}{记号索引}{252}
\ivonecontentsline{}{术语索引}{254}

\vspace{2em}
\begin{flushright}
\emph{1963 年 3 月 15 日收到。}
\end{flushright}
\let\ivonecontentsline\relax
\endgroup
